\documentclass[11pt]{article}

\usepackage[T1]{fontenc}
\usepackage[margin=1in]{geometry}
\usepackage{amsmath,amssymb,amsthm}
\usepackage{booktabs}
\usepackage{hyperref}

\hypersetup{
  colorlinks=true,
  linkcolor=blue,
  urlcolor=blue
}

\newcommand{\R}{\mathbb{R}}
\newcommand{\E}{\mathbb{E}}
\newcommand{\Prob}{\mathbb{P}}
\newcommand{\fl}{\operatorname{fl}}
\newcommand{\Frob}[1]{\left\lVert #1\right\rVert_F}
\newcommand{\Norm}[1]{\left\lVert #1\right\rVert_2}
\newcommand{\ActNorm}[1]{\left\lVert #1\right\rVert_{2,\mathrm{act}}}
\newcommand{\Atil}{\widetilde A}
\newcommand{\Util}{\widetilde U}
\newcommand{\Ahat}{\widehat{\widetilde A}}
\newcommand{\Uhat}{\widehat{\widetilde U}}
\newcommand{\uunit}{u}
\newcommand{\lean}[1]{\texttt{\detokenize{#1}}}
\newcommand{\leanfit}[1]{\noindent{\scriptsize\nolinkurl{#1}}\par}

\newtheorem{theorem}{Theorem}
\newtheorem{corollary}[theorem]{Corollary}
\newtheorem{lemma}[theorem]{Lemma}
\theoremstyle{definition}
\newtheorem{definition}[theorem]{Definition}
\theoremstyle{remark}
\newtheorem*{remark}{Remark}
\newenvironment{formalized}
  {\begin{quote}\small\noindent\textbf{Lean formalization.}\par\smallskip}
  {\end{quote}}
\title{Lean Formalization Summary for the Explicit RandNLA CACM Algorithms}
\author{LeanFpAnalysis repository}
\date{}

\begin{document}
\sloppy
\maketitle

\begin{abstract}
This note summarizes the Lean formalization attached to the three explicit
meta-algorithms in Drineas and Mahoney's CACM survey,
\href{https://dl.acm.org/doi/10.1145/2842602}{RandNLA: Randomized Numerical
Linear Algebra}.  The formalized results focus on stability and on the
paper-level approximation guarantees that are currently present in the
repository: Algorithm 1 entrywise sampling stability, a hard-thresholded
element-wise spectral theorem used as a source-aligned support variant,
and a faithful nontruncated Algorithm~1 Frobenius/Markov floating-point
corollary with an explicit scalar radius, plus a literal Algorithm~1
source-rate specialization under a no-small-entry condition and a
nonconditional literal-law support-radius theorem with all reciprocal-entry
and floating-point terms expanded,
Algorithm 2 row-sampling Gram approximation and its leverage-score
specialization, the Algorithm 2 leverage-score
least-squares consequence behind equation~(8), and Algorithm 3
matrix-preconditioning stability.  It also records the remaining accuracy
backlog: the source-uniform sharp literal Algorithm 1 equation~(2)
matrix-Bernstein theorem for \(p_{ij}=A_{ij}^2/\Frob{A}^2\), with no
reciprocal-entry or no-small-entry dependence, its fully rectangular
CACM-prose variant,
the implementation-backed rounded least-squares sketch and
solver pipeline for equation~(8), and the later paper-level randomized
approximation claims.  Each mathematical statement below is paired with its
Lean theorem names.
\end{abstract}

\section{Common Floating-Point Model}

The repository uses an axiomatic floating-point model with unit roundoff
\(\uunit\).  Each basic arithmetic operation satisfies the standard relative
error model
\[
  \fl(x \circ y)=(x\circ y)(1+\delta),
  \qquad |\delta|\le \uunit .
\]
Accumulated dot-product and matrix-product errors use Higham's
\[
  \gamma_k=\frac{k\uunit}{1-k\uunit},
\]
with Lean side condition \(\lean{gammaValid fp k}\).

\section{Algorithm 1: Element-Wise Sampling}

Algorithm 1 samples entries \((i,j)\) independently with probabilities
\[
  p_{ij}=\frac{A_{ij}^2}{\sum_{k,\ell} A_{k\ell}^2}
  =\frac{A_{ij}^2}{\Frob{A}^2}.
\]
The current formalization studies exact unbiasedness, the one-step
mean-zero residual increments, the floating-point stability of the sampled
entry update, and the random hit counter of a fixed entry.

By current convention, the probabilities \(p_{ij}\) are exact mathematical
sampling laws.  The implementation-facing floating-point statements therefore
charge the arithmetic performed after sampling: the sampled-entry rescaling
division, repeated accumulation into the sketch, and residual formation.  They
do not charge floating-point error for constructing, normalizing, repairing, or
sampling from the probability table.  The `WithProb` Lean functions are exact
supplied-law adapters, useful for deterministic transfers, not certificates for
probability construction.

\begin{formalized}
\leanfit{sqMagProb}
\leanfit{sqMagProb_sum_eq_one}
\leanfit{elementwiseSampleUpdateWithProb}
\leanfit{fl_elementwiseSampleUpdateWithProb_error_bound}
\leanfit{fl_elementwiseTraceSketchWithProb_error_bound}
\end{formalized}

\begin{theorem}[Exact unbiasedness]
Assume the trace length equals the Algorithm 1 sample count, \(steps=s\), and
\(s\ne 0\).  Under the independent squared-magnitude trace law, the exact
zero-initialized estimator is unbiased entrywise:
\[
  \mathbb E\!\left[\Atil_{ij}\right]=A_{ij}.
\]
Equivalently, as a matrix-valued statement,
\[
  \bigl(\mathbb E[\Atil_{ij}]\bigr)_{i,j}=A .
\]
\end{theorem}

\begin{formalized}
\leanfit{sqMagTraceProbMass_marginal_one}
\leanfit{sqMagTraceProbability_eventProb_sampleHits}
\leanfit{sqMagTraceProbability_expectationReal_hitCount_eq}
\leanfit{elementwiseTraceSketch_zero_init_of_entry_eq_zero}
\leanfit{sqMagTraceProbability_expectationReal_elementwiseTraceSketch_nonzero_entry}
\leanfit{sqMagTraceProbability_expectationReal_elementwiseTraceSketch_zero_entry}
\leanfit{sqMagTraceProbability_expectationReal_elementwiseTraceSketch_entry}
\leanfit{sqMagTraceProbability_expectationReal_elementwiseTraceSketch_matrix}
\end{formalized}

\begin{lemma}[Zero-mean residual increments]
For one squared-magnitude sample \(X\), define the residual increment
\[
  Z_X(i,j)=\frac{A_{ij}}{s}
  -
  \mathbf 1_{\{X=(i,j)\}}\frac{A_{ij}}{s p_{ij}}.
\]
If \(s\ne0\) and \(\Frob{A}^2>0\), then
\[
  \mathbb E[Z_X(i,j)] = 0 .
\]
For an \(s\)-sample trace \(X_1,\ldots,X_s\), the exact residual decomposes as
\[
  A-\Atil = \sum_{t=1}^s Z_{X_t}.
\]
Consequently, for every fixed vector \(x\),
\[
  \mathbb E[Z_X x]=0,
  \qquad
  (A-\Atil)x=\sum_{t=1}^s Z_{X_t}x.
\]
\end{lemma}


\begin{formalized}
\leanfit{elementwiseSampleContribution}
\leanfit{elementwiseTraceContribution_eq_sampleContribution}
\leanfit{elementwiseSampleResidualIncrement}
\leanfit{elementwiseTraceResidual_eq_sum_sampleResidualIncrement}
\leanfit{rectMatMulVec_elementwiseTraceResidual_eq_sum_sampleResidualIncrement}
\leanfit{sqMagProb_sum_elementwiseSampleContribution_entry}
\leanfit{sqMagProb_sum_elementwiseSampleResidualIncrement_entry_eq_zero}
\leanfit{sqMagProb_sum_rectMatMulVec_elementwiseSampleResidualIncrement_eq_zero}
\leanfit{sqMagTraceProbability_expectationReal_elementwiseTraceResidual_entry_eq_zero}
\leanfit{sqMagTraceProbability_expectationReal_rectMatMulVec_elementwiseTraceResidual_eq_zero}
\end{formalized}

\begin{theorem}[Deterministic entrywise stability]
Fix an entry \((i,j)\), and let \(q_{ij}\) be the number of times this entry is
sampled in a deterministic Algorithm 1 trace.  For nonzero \(A_{ij}\) and
valid \(\gamma_{q_{ij}}\), the computed entry satisfies
\[
\left|
  \widehat{\Atil}_{ij}
  -
  \left(
    \Atil^{(0)}_{ij}
    +q_{ij}\frac{\Frob{A}^2}{sA_{ij}}
  \right)
\right|
\le
\left|\Atil^{(0)}_{ij}\right|\gamma_{q_{ij}}
+
q_{ij}
\left|\frac{\Frob{A}^2}{sA_{ij}}\right|
\gamma_{q_{ij}+1}.
\]
\end{theorem}

\begin{formalized}
\leanfit{fl_elementwiseTraceSketch_sqMag_error_bound}
\leanfit{fl_elementwiseTraceSketch_sqMag_error_bound_exact}
\leanfit{fl_elementwiseSketch_zero_init_sqMag_error_bound}
\end{formalized}

\begin{corollary}[High-probability stability from hit-count concentration]
If a proved concentration theorem gives
\[
  \Prob[q_{ij}\le Q]\ge 1-\delta,
\]
then the same event implies the deterministic stability bound with the budget
\[
  B(Q)=
  |\Atil^{(0)}_{ij}|\gamma_Q+
  Q\left|\frac{\Frob{A}^2}{sA_{ij}}\right|\gamma_{Q+1}.
\]
The repository instantiates this transfer with Markov, pairwise-Chebyshev, and
independent Chernoff concentration for the Algorithm 1 product trace law.
\end{corollary}

\begin{formalized}
\leanfit{highProbability_sqMagTraceStability_of_markov_budget}
\leanfit{highProbability_sqMagTraceStability_of_pairwise_chebyshev_budget}
\leanfit{highProbability_sqMagTraceStability_of_independent_chernoff_budget}
\leanfit{highProbability_sqMagTraceStability_of_independent_chernoff_optimized_tail_budget}
\end{formalized}

\section{Algorithm 1: Spectral Transfer for Equation (2)}

The CACM survey's equation~(2) is a high-probability spectral-norm
concentration theorem for the exact sampled residual \(A-\Atil\).  The
repository contains both the deterministic floating-point transfer layer and a
closed source-aligned square theorem for the truncated squared-magnitude route
used in Drineas--Zouzias.  The operator norm is stated in rectangular
vector-action form:
\[
  \ActNorm{M}\le c
  \quad\Longleftrightarrow\quad
  \Norm{Mx}\le c\Norm{x}\quad\text{for all }x .
\]
The first few statements in this section are support adapters: they expose the
generic perturbation matrix used internally by the Lean library.  They are not
the advertised equation~(2) result.  The end-of-section corollary eliminates
that matrix and gives explicit \(\gamma_{s+1}\)-radii: first for the
hard-thresholded source variant, and then for a faithful nontruncated
Frobenius/Markov Algorithm~1 corollary.  It also records a literal
source-rate corollary under the explicit condition that the source threshold
\(\varepsilon/(2n)\) would not remove any nonzero entry.  The fully
unconditional sharp literal CACM equation~(2) matrix-Bernstein theorem remains
open.

\begin{theorem}[Support adapter: floating-point transfer for an exact spectral event]
Let \(R=A-\Atil\) be the exact Algorithm 1 residual for a fixed trace, and let
\(\widehat R=A-\widehat{\Atil}\) be the floating-point residual for the same
trace.  Suppose
\[
  \Norm{Rx}\le \varepsilon\Norm{x}\quad\text{for all }x,
  \qquad
  |\widehat{\Atil}_{ij}-\Atil_{ij}|\le B_{ij},
  \qquad B_{ij}\ge 0 .
\]
Then
\[
  \Norm{\widehat R x}
  \le
  \left(\varepsilon+\Frob{B}\right)\Norm{x}
  \quad\text{for all }x .
\]
\end{theorem}

\begin{formalized}
\leanfit{frobNormRect}
\leanfit{rectOpNorm2Le}
\leanfit{rectOpNorm2Le_add_of_rectOpNorm2Le_of_frobNormRect_le}
\leanfit{fl_elementwiseTraceResidual_rectOpNorm2Le_of_exact}
\end{formalized}

\begin{corollary}[Support adapter: hit-count-budget transfer]
Assume the sampled-entry denominators are nonzero for every entry considered
(the Lean corollary uses the sufficient hypothesis \(A_{ij}\ne 0\) for all
\((i,j)\)), and assume the required \(\gamma_Q\) and \(\gamma_{Q+1}\) validity
conditions.  If every entry hit count in a fixed trace is at most \(Q\), then
the entrywise Algorithm 1 stability budget \(B(Q)\) can be used as the matrix
\(B\) above.  Consequently an exact spectral residual bound transfers to the
floating-point residual with additive perturbation \(\Frob{B(Q)}\).
\end{corollary}

\begin{formalized}
\leanfit{fl_elementwiseTraceResidual_rectOpNorm2Le_of_exact_and_hitCount_le}
\end{formalized}

\begin{corollary}[Support adapter: probability transfer]
On a finite probability space, if the exact spectral event has probability at
least \(\rho\), and the entrywise floating-point budget holds for every
outcome, then the floating-point spectral event has probability at least
\(\rho\).
\end{corollary}

\begin{formalized}
\leanfit{algorithm1ExactSpectralEvent_subset_flSpectralEvent}
\leanfit{probability_algorithm1_fl_spectral_of_exact_spectral}
\end{formalized}

\begin{theorem}[Self-adjoint dilation support]
Let \(R=A-\Atil\) be the exact rectangular residual, and define its
self-adjoint dilation by
\[
  \mathcal D(R)=
  \begin{bmatrix}
    0 & R\\
    R^T & 0
  \end{bmatrix}.
\]
If the square dilation satisfies the vector-action bound
\[
  \Norm{\mathcal D(R)z}\le \varepsilon\Norm{z}
  \quad\text{for all block vectors }z,
\]
then the rectangular residual satisfies
\[
  \Norm{Rx}\le \varepsilon\Norm{x}
  \quad\text{for all }x .
\]
Consequently, any proved high-probability event for the dilation transfers to
the exact rectangular residual event.  The generic floating-point handoff is
kept in Lean as a support adapter; the displayed equation~(2) endpoint below
uses the explicit \(\gamma_{s+1}\)-radius instead.
\end{theorem}

\begin{lemma}[Dilation increment foundations]
For an \(s\)-sample Algorithm 1 trace,
\[
  \mathcal D(A-\Atil)
  =
  \sum_{t=1}^s \mathcal D(Z_{X_t}),
\]
where \(Z_X\) is the one-sample residual increment.  Under the one-step
squared-magnitude law, every entry of \(\mathcal D(Z_X)\) has mean zero, and
\[
  \mathbb E\Frob{\mathcal D(Z_X)}^2
  \le
  2mn\,\frac{\Frob{A}^2}{s^2}.
\]
The finite-probability library also now contains the real-valued exponential
Markov kernel
\[
  \Prob[T\le X]\le
  \exp(-\lambda T)\,
  \mathbb E\exp(\lambda X),
  \qquad \lambda>0.
\]
The square finite matrix layer contains product and trace operations, trace
additivity, trace homogeneity, cyclicity \(\operatorname{tr}(MN)=\operatorname{tr}(NM)\),
finite identity and transpose operations, quadratic forms, PSD and Loewner-order
predicates, quadratic-form add/neg/sub algebra, two-sided Loewner-to-absolute
quadratic-form control, symmetric-square PSD facts, the deterministic implication
\[
  \Norm{Mx}\le L\Norm{x}\ \forall x
  \quad\Longrightarrow\quad
  M^2 \preceq L^2 I
\]
for symmetric finite matrices, and the dilation trace-of-square identity
\[
  \operatorname{tr}\!\left(\mathcal D(M)^2\right)
  =
  2\Frob{M}^2 .
\]
Conversely, for a symmetric finite matrix and \(L\ge0\),
\[
  M^2\preceq L^2I
  \quad\Longrightarrow\quad
  \Norm{Mx}\le L\Norm{x}\ \text{for every }x .
\]
Therefore, for an Algorithm 1 residual \(R\), the squared dilation event
\[
  \mathcal D(R)^2\preceq \varepsilon^2 I
\]
implies both the exact dilation operator event and the rectangular event
\[
  \Norm{Rx}\le \varepsilon\Norm{x}\qquad\text{for every }x .
\]
The same event transfers to the floating-point residual after adding the
proved entrywise perturbation budget.
It also proves the quadratic-form variance proxy
\[
  \mathbb E\,
  z^T\mathcal D(Z_X)^2z
  \le
  2mn\,\frac{\Frob{A}^2}{s^2}\Norm{z}^2
\]
for every block vector \(z\).  Equivalently, the one-step variance matrix
\[
  V_1=\mathbb E\,\mathcal D(Z_X)^2
\]
is positive semidefinite and satisfies the Loewner bound
\[
  0\preceq V_1
  \preceq
  2mn\,\frac{\Frob{A}^2}{s^2}I .
\]
After summing the \(s\) independent one-step proxies, the variance matrix
\[
  V_s=\sum_{t=1}^s\mathbb E\,\mathcal D(Z_{X_t})^2
\]
is positive semidefinite and satisfies
\[
  \sum_{t=1}^s
  \mathbb E\,
  z^T\mathcal D(Z_{X_t})^2z
  \le
  2mn\,\frac{\Frob{A}^2}{s}\Norm{z}^2 .
\]
Equivalently, \(0\preceq V_s\preceq
2mn\,\Frob{A}^2 I/s\).
The trace form is also formalized:
\[
  \operatorname{tr}(V_1)
  \le
  2mn\,\frac{\Frob{A}^2}{s^2}(m+n),
  \qquad
  \operatorname{tr}(V_s)
  \le
  2mn\,\frac{\Frob{A}^2}{s}(m+n).
\]
The product-law bridge is now formalized explicitly as well.  If
\(X_1,\ldots,X_s\) are the independent Algorithm 1 samples and
\(R=\sum_t Z_{X_t}\), then every entry of the full dilation residual is
mean-zero:
\[
  \mathbb E\,\mathcal D(R)_{ab}=0 .
\]
Moreover the trace-law expected variance matrix
\[
  \sum_{t=1}^s \mathbb E\,\mathcal D(Z_{X_t})^2
\]
is the same object bounded above by \(2mn\,\Frob{A}^2 I/s\) in Loewner order.
The scalar product-law MGF bridge is also formalized.  For any one-step
statistic \(f\colon (i,j)\mapsto\mathbb R\),
\[
  \mathbb E\exp\!\left(\lambda\sum_{t=1}^s f(X_t)\right)
  =
  \left(
    \sum_{i,j}p_{ij}\exp(\lambda f(i,j))
  \right)^s .
\]
Consequently the repository has exponential-Markov upper-tail and complement
forms for \(\sum_t f(X_t)\), including variants that consume an explicit
one-step bound
\(\sum_{i,j}p_{ij}\exp(\lambda f(i,j))\le \exp(\psi)\) and return the tail
\(\exp(s\psi-\lambda T)\).  This is useful scalar concentration
infrastructure for fixed-projection or covering arguments.  These results are
not by themselves the noncommutative trace-exponential theorem needed for
matrix Bernstein, but the later trace-MGF and Bernstein sections compose the
source-aligned route into a closed high-probability spectral event and
floating-point transfer.
\end{lemma}

\begin{corollary}[Finite test-vector quadratic-form MGF support]
Let \(\mathcal N\) be a finite index set of deterministic test vectors
\(z_a\in\mathbb R^{m+n}\).  For each \(a\in\mathcal N\), define the one-step
scalar statistic
\[
  f_a(i,j)
  =
  z_a^T D\!\left(\Delta_{ij}\right)z_a,
\]
where \(\Delta_{ij}\) is the one-sample residual increment of Algorithm~1 and
\(D(\cdot)\) is the self-adjoint dilation.  The formalization proves the
identity
\[
  z_a^T D(R) z_a
  =
  \sum_{t=1}^s f_a(X_t).
\]
Consequently, if for every \(a\in\mathcal N\)
\[
  \sum_{i,j} p_{ij}\exp(\lambda_a f_a(i,j))
  \le \exp(\psi_a),
  \qquad \lambda_a>0,
\]
then
\[
  \mathbb P\!\left[
    z_a^T D(R)z_a\le T_a
    \text{ for all }a\in\mathcal N
  \right]
  \ge
  1-\sum_{a\in\mathcal N}\exp(s\psi_a-\lambda_aT_a).
\]
The one-step MGF hypothesis can itself be discharged from a pointwise bound:
if \(f_a(i,j)\le B_a\) for every sampled entry, then the same conclusion holds
with \(\psi_a=\lambda_aB_a\).  There is also a support-aware form: it is enough
to prove \(f_a(i,j)\le B_a\) only when \(p_{ij}>0\), since the zero-probability
entries contribute zero mass to the one-step MGF.  This support-aware statement
is the version used by the truncated source-aligned route, where retained-entry
side conditions are available exactly on positive-probability samples.
This is a finite-family scalar-MGF theorem for supplied test vectors.  It does
not construct a net, prove sharp Bernstein/Hoeffding-type one-step MGF bounds,
or prove a largest-eigenvalue or matrix-Bernstein theorem.
\end{corollary}

\begin{corollary}[One-sided dilation event transfer]
Let \(R=A-\widetilde A\) be the exact Algorithm~1 residual and let
\[
  D(R)=
  \begin{bmatrix}
    0 & R\\
    R^T & 0
  \end{bmatrix}.
\]
If \(\varepsilon\ge 0\) and
\[
  D(R)\preceq \varepsilon I,
\]
then
\[
  \Norm{R}\le \varepsilon.
\]
Consequently, any high-probability largest-eigenvalue theorem proving the
one-sided dilation event transfers to the repository's exact rectangular
operator event, and then to the floating-point event after adding the proved
entrywise FP perturbation budget.  This is an event-transfer theorem only: it
does not prove that \(D(R)\preceq\varepsilon I\) occurs with high probability.
\end{corollary}


\begin{formalized}
\leanfit{finiteVecNorm2}
\leanfit{finiteMatVec}
\leanfit{finiteFrobNormSq}
\leanfit{finiteMatMul}
\leanfit{finiteIdMatrix}
\leanfit{finiteBasisVec}
\leanfit{finiteTranspose}
\leanfit{finiteTrace}
\leanfit{finiteTrace_add}
\leanfit{finiteTrace_smul}
\leanfit{finiteTrace_neg}
\leanfit{finiteTrace_sub}
\leanfit{finiteTrace_finiteIdMatrix}
\leanfit{finiteTrace_smul_finiteIdMatrix}
\leanfit{finiteTrace_fintype_sum}
\leanfit{finiteTrace_fintype_sum_smul}
\leanfit{finiteTrace_nonneg_of_finitePSD}
\leanfit{finiteTrace_mono_of_finiteLoewnerLe}
\leanfit{finiteTrace_le_of_finiteLoewnerLe_smul_id}
\leanfit{finiteTrace_finiteMatMul_comm}
\leanfit{finiteMatVec_finiteMatMul}
\leanfit{finiteMatVec_finiteIdMatrix}
\leanfit{finiteMatVec_finiteBasisVec}
\leanfit{finiteVecNorm2Sq_finiteMatVec_le_finiteFrobNormSq_mul}
\leanfit{finiteOpNorm2Le}
\leanfit{finiteQuadraticForm}
\leanfit{finitePSD}
\leanfit{finiteLoewnerLe}
\leanfit{finitePSD_iff_zero_loewnerLe}
\leanfit{finiteQuadraticForm_finiteIdMatrix}
\leanfit{finiteQuadraticForm_finiteBasisVec}
\leanfit{finiteQuadraticForm_smul_finiteIdMatrix}
\leanfit{finiteQuadraticForm_add}
\leanfit{finiteQuadraticForm_neg}
\leanfit{finiteQuadraticForm_sub}
\leanfit{finiteLoewnerLe_iff_sub_finitePSD}
\leanfit{finiteLoewnerLe_add}
\leanfit{finiteLoewnerLe_smul_of_nonneg}
\leanfit{finiteQuadraticForm_finset_sum}
\leanfit{finiteQuadraticForm_finset_sum_smul}
\leanfit{finiteQuadraticForm_fintype_sum}
\leanfit{finiteQuadraticForm_fintype_sum_smul}
\leanfit{finiteLoewnerLe_fintype_sum}
\leanfit{finiteLoewnerLe_fintype_sum_smul_of_nonneg}
\leanfit{finitePSD_fintype_sum_of_finitePSD}
\leanfit{finitePSD_fintype_sum_smul_of_nonneg}
\leanfit{abs_finiteQuadraticForm_le_of_loewnerLe_neg}
\leanfit{finiteVecInnerProduct_sq_le}
\leanfit{abs_finiteVecInnerProduct_le_finiteVecNorm2_mul}
\leanfit{abs_finiteVecInnerProduct_finiteMatVec_le_of_finiteOpNorm2Le}
\leanfit{abs_finiteQuadraticForm_le_of_finiteOpNorm2Le}
\leanfit{finiteQuadraticForm_le_of_finiteOpNorm2Le}
\leanfit{finiteLoewnerLe_smul_id_of_finiteOpNorm2Le}
\leanfit{finiteLoewnerLe_neg_smul_id_of_finiteOpNorm2Le}
\leanfit{finiteVecInnerProduct_finiteMatVec_eq_transpose}
\leanfit{finiteVecInnerProduct_finiteMatVec_left_eq_right_of_symmetric}
\leanfit{finiteQuadraticForm_finiteMatMul_self_of_symmetric}
\leanfit{finitePSD_finiteMatMul_self_of_symmetric}
\leanfit{finiteQuadraticForm_finiteMatMul_self_le_finiteFrobNormSq_mul_of_symmetric}
\leanfit{finiteMatMul_self_loewnerLe_scalar_id_of_finiteOpNorm2Le}
\leanfit{finiteTrace_finiteMatMul_self_le_card_mul_sq_of_finiteOpNorm2Le}
\leanfit{finiteOpNorm2Le_of_finiteMatMul_self_loewnerLe_scalar_id}
\leanfit{IsSymmetricFiniteMatrix}
\leanfit{IsSymmetricFiniteMatrix.to_matrix_isSymm}
\leanfit{Matrix_isSymm.to_IsSymmetricFiniteMatrix}
\leanfit{IsSymmetricFiniteMatrix.to_matrix_isHermitian}
\leanfit{Matrix_isHermitian.to_IsSymmetricFiniteMatrix}
\leanfit{finitePSD.to_matrix_posSemidef}
\leanfit{Matrix_posSemidef.to_finitePSD}
\leanfit{finitePSD_iff_matrix_posSemidef_of_symmetric}
\leanfit{finiteLoewnerLe.to_matrix_posSemidef_sub}
\leanfit{Matrix_posSemidef_sub.to_finiteLoewnerLe}
\leanfit{finiteLoewnerLe_iff_matrix_posSemidef_sub_of_symmetric}
\leanfit{finiteHermitianEigenvalues}
\leanfit{finiteHermitianEigenvalues_mem_spectrum_real}
\leanfit{finiteTrace_eq_sum_finiteHermitianEigenvalues}
\leanfit{finiteDiagonal}
\leanfit{finiteMatrixExp}
\leanfit{finiteMatrixExp_smul_finiteIdMatrix}
\leanfit{finiteTrace_finiteMatrixExp_smul_finiteIdMatrix}
\leanfit{finiteMatrixExp_symmetric}
\leanfit{finiteMatrixExp_finiteDiagonal}
\leanfit{finiteTrace_finiteMatrixExp_finiteDiagonal}
\leanfit{finiteHermitianCfcExp}
\leanfit{finiteTrace_finiteHermitianCfcExp_eq_sum_exp_finiteHermitianEigenvalues}
\leanfit{finiteTrace_finiteMatrixExp_eq_sum_exp_finiteHermitianEigenvalues}
\leanfit{finiteTrace_finiteMatrixExp_neg_eq_sum_exp_neg_finiteHermitianEigenvalues}
\leanfit{exp_le_finiteTrace_finiteMatrixExp_of_finiteHermitianEigenvalue_ge}
\leanfit{finiteTrace_finiteMatrixExp_nonneg}
\leanfit{exp_neg_le_finiteTrace_finiteMatrixExp_neg_of_finiteHermitianEigenvalue_le}
\leanfit{finiteTrace_finiteMatrixExp_neg_nonneg}
\leanfit{FiniteProbability.eventProb_exists_finiteHermitianEigenvalue_ge_le_exp_neg_mul_expected_trace_exp}
\leanfit{FiniteProbability.eventProb_exists_finiteHermitianEigenvalue_le_le_exp_mul_expected_trace_exp_neg}
\leanfit{FiniteProbability.eventProb_forall_abs_finiteHermitianEigenvalue_lt_ge_one_sub_exp_neg_mul_expected_trace_exp_add}
\leanfit{finitePSD_iff_finiteHermitianEigenvalues_nonneg}
\leanfit{finiteLoewnerLe_iff_sub_finiteHermitianEigenvalues_nonneg}
\leanfit{smulFiniteIdMatrix_symmetric}
\leanfit{finiteScalarUpperDiff_symmetric}
\leanfit{finiteScalarUpperDiffEigenvalues}
\leanfit{finiteLoewnerLe_smul_id_iff_finiteScalarUpperDiffEigenvalues_nonneg}
\leanfit{finiteLoewnerLe_smul_id_iff_sub_finiteHermitianEigenvalues_nonneg}
\leanfit{rectSelfAdjointDilation}
\leanfit{rectSelfAdjointDilation_symmetric}
\leanfit{finitePSD_rectSelfAdjointDilation_square}
\leanfit{rectSelfAdjointDilation_square_loewnerLe_scalar_id_of_finiteOpNorm2Le}
\leanfit{rectSelfAdjointDilation_opNorm2Le_of_square_loewnerLe_scalar_id}
\leanfit{sumBothVec}
\leanfit{finiteVecNorm2Sq_sumBothVec}
\leanfit{finiteMatVec_rectSelfAdjointDilation_sumBothVec}
\leanfit{finiteQuadraticForm_rectSelfAdjointDilation_sumBothVec}
\leanfit{rectOpNorm2Le_of_selfAdjointDilation_loewnerLe_scalar_id}
\leanfit{rectOpNorm2Le_of_selfAdjointDilation_square_loewnerLe_scalar_id}
\leanfit{finiteFrobNormSq_rectSelfAdjointDilation}
\leanfit{finiteTrace_finiteMatMul_rectSelfAdjointDilation_self}
\leanfit{finiteMatVec_rectSelfAdjointDilation_sumInr}
\leanfit{rectOpNorm2Le_of_selfAdjointDilation}
\leanfit{algorithm1ExactDilationEvent}
\leanfit{algorithm1ExactDilationUpperEvent}
\leanfit{algorithm1ExactDilationEigenUpperEvent}
\leanfit{algorithm1ExactDilationEigenUpperIndexEvent}
\leanfit{algorithm1ExactDilationSquareEvent}
\leanfit{algorithm1ExactDilationEvent_subset_exactSpectralEvent}
\leanfit{algorithm1ExactDilationUpperEvent_subset_exactSpectralEvent}
\leanfit{algorithm1ExactDilationEigenUpperEvent_subset_exactDilationUpperEvent}
\leanfit{algorithm1ExactDilationEigenUpperEvent_subset_exactSpectralEvent}
\leanfit{algorithm1ExactDilationSquareEvent_subset_exactDilationEvent}
\leanfit{algorithm1ExactDilationSquareEvent_subset_exactSpectralEvent}
\leanfit{probability_algorithm1_exact_spectral_of_dilation}
\leanfit{probability_algorithm1_exact_spectral_of_dilation_upper}
\leanfit{probability_algorithm1_exact_spectral_of_dilation_eigen_upper}
\leanfit{probability_algorithm1_exact_dilation_eigen_upper_of_index_bounds}
\leanfit{probability_algorithm1_exact_spectral_of_dilation_eigen_upper_index_bounds}
\leanfit{probability_algorithm1_exact_spectral_of_dilation_square}
\leanfit{probability_algorithm1_fl_spectral_of_exact_dilation}
\leanfit{probability_algorithm1_fl_spectral_of_exact_dilation_upper}
\leanfit{probability_algorithm1_fl_spectral_of_exact_dilation_eigen_upper}
\leanfit{probability_algorithm1_fl_spectral_of_exact_dilation_eigen_upper_index_bounds}
\leanfit{probability_algorithm1_fl_spectral_of_exact_dilation_square}
\leanfit{rectSelfAdjointDilation_elementwiseTraceResidual_eq_sum_sampleResidualIncrement}
\leanfit{sqMagProb_sum_rectSelfAdjointDilation_elementwiseSampleResidualIncrement_eq_zero}
\leanfit{sqMagProb_sum_finiteFrobNormSq_rectSelfAdjointDilation_elementwiseSampleResidualIncrement_le}
\leanfit{sqMagProb_sum_finiteQuadraticForm_rectSelfAdjointDilation_square_le}
\leanfit{sqMagProb_sum_steps_finiteQuadraticForm_rectSelfAdjointDilation_square_le}
\leanfit{sqMagProb_sum_rectSelfAdjointDilation_square_loewnerLe_scalar_id}
\leanfit{sqMagProb_sum_rectSelfAdjointDilation_square_psd}
\leanfit{sqMagProb_sum_steps_rectSelfAdjointDilation_square_loewnerLe_scalar_id}
\leanfit{sqMagProb_sum_steps_rectSelfAdjointDilation_square_psd}
\leanfit{sqMagProb_sum_finiteTrace_rectSelfAdjointDilation_square_le}
\leanfit{sqMagProb_sum_steps_finiteTrace_rectSelfAdjointDilation_square_le}
\leanfit{sqMagTraceProbability_expectationReal_step_eq}
\leanfit{exp_sum_stepFunction_eq_prod}
\leanfit{sqMagTraceProbMass_exp_sum_stepFunction_eq}
\leanfit{sqMagTraceProbability_expectationReal_exp_sum_stepFunction_eq}
\leanfit{sqMagTraceProbability_eventProb_sum_stepFunction_ge_le_exp_mul_mgf}
\leanfit{sqMagTraceProbability_eventProb_sum_stepFunction_le_ge_one_sub_exp_mul_mgf}
\leanfit{sqMagTraceProbability_eventProb_sum_stepFunction_ge_le_exp_of_one_step_mgf_bound}
\leanfit{sqMagTraceProbability_eventProb_sum_stepFunction_le_ge_one_sub_exp_of_one_step_mgf_bound}
\leanfit{sqMagTraceProbability_eventProb_forall_sum_stepFunction_le_ge_one_sub_sum_exp_of_one_step_mgf_bound}
\leanfit{sqMagProb_sum_exp_stepFunction_le_exp_of_forall_le}
\leanfit{sqMagProb_sum_exp_stepFunction_le_exp_of_support_forall_le}
\leanfit{sqMagTraceProbability_eventProb_forall_sum_stepFunction_le_ge_one_sub_sum_exp_of_pointwise_bound}
\leanfit{sqMagTraceProbability_eventProb_forall_sum_stepFunction_le_ge_one_sub_sum_exp_of_support_pointwise_bound}
\leanfit{finiteQuadraticForm_rectSelfAdjointDilation_elementwiseTraceResidual_eq_sum_sampleResidualIncrement}
\leanfit{sqMagTraceProbability_eventProb_forall_finiteQuadraticForm_rectSelfAdjointDilation_elementwiseTraceResidual_le_ge_one_sub_sum_exp_of_one_step_mgf_bound}
\leanfit{sqMagTraceProbability_eventProb_forall_finiteQuadraticForm_rectSelfAdjointDilation_elementwiseTraceResidual_le_ge_one_sub_sum_exp_of_pointwise_bound}
\leanfit{sqMagTraceProbability_eventProb_forall_finiteQuadraticForm_rectSelfAdjointDilation_elementwiseTraceResidual_le_ge_one_sub_sum_exp_of_support_pointwise_bound}
\leanfit{sqMagTraceProbability_expectationReal_rectSelfAdjointDilation_sampleResidualIncrement_eq_zero}
\leanfit{sqMagTraceProbability_expectationReal_rectSelfAdjointDilation_elementwiseTraceResidual_eq_zero}
\leanfit{sqMagTraceProbability_expectationReal_finiteQuadraticForm_rectSelfAdjointDilation_sampleResidualIncrement_square_le}
\leanfit{sqMagTraceProbability_expectationReal_sum_finiteQuadraticForm_rectSelfAdjointDilation_sampleResidualIncrement_square_le}
\leanfit{sqMagTraceProbability_sum_expectationReal_rectSelfAdjointDilation_square_loewnerLe_scalar_id}
\leanfit{sqMagTraceProbability_sum_expectationReal_rectSelfAdjointDilation_square_psd}
\leanfit{FiniteProbability.eventProb_real_ge_le_exp_mul_mgf}
\leanfit{FiniteProbability.eventProb_real_le_ge_one_sub_exp_mul_mgf}
\leanfit{FiniteProbability.eventProb_real_le_ge_one_sub_exp_of_mgf_bound}
\leanfit{FiniteProbability.eventProb_real_le_mean_add_ge_one_sub_exp_sq_of_subgaussian_mgf}
\leanfit{FiniteProbability.log_mgf_differential_le_of_entropyReal_exp_mul_le}
\leanfit{FiniteProbability.log_mgf_div_sub_quadratic_antitoneOn_of_differential_le}
\leanfit{FiniteProbability.tendsto_log_mgf_div_nhdsGT_zero}
\leanfit{FiniteProbability.log_mgf_le_mean_add_quadratic_of_differential_le}
\leanfit{FiniteProbability.log_mgf_le_mean_add_quadratic_of_entropyReal_exp_mul_le}
\leanfit{FiniteProbability.expectationReal_exp_centered_le_exp_of_log_mgf_le}
\leanfit{FiniteProbability.eventProb_real_le_mean_add_ge_one_sub_exp_sq_of_log_mgf_bound}
\leanfit{FiniteProbability.eventProb_real_le_mean_add_ge_one_sub_exp_sq_of_entropyReal_exp_mul_le}
\end{formalized}

\begin{lemma}[Finite MGF, entropy, and Laplace calculus]
Let \(X\) be a real random variable on a finite probability space.  The finite
moment-generating function
\[
  M(\lambda)=\mathbb E\exp\{\lambda X\}
\]
is strictly positive, and
\[
  M'(\lambda)=\mathbb E\!\left[X\exp\{\lambda X\}\right],
  \qquad
  (\log M)'(\lambda)
    =\frac{\mathbb E[X\exp\{\lambda X\}]}{\mathbb E[\exp\{\lambda X\}]}.
\]
For the real entropy functional
\[
  \operatorname{Ent}(Z)=
  \mathbb E[Z\log Z]-\mathbb E[Z]\log(\mathbb E[Z]),
\]
the exponential tilt satisfies
\[
  \operatorname{Ent}(e^{\lambda X})
  =
  \lambda\,\mathbb E[Xe^{\lambda X}]
  -\mathbb E[e^{\lambda X}]\log\mathbb E[e^{\lambda X}].
\]
Consequently, any explicit entropy bound
\[
  \operatorname{Ent}(e^{\lambda X})
    \le c\lambda^2 M(\lambda)
\]
implies the pointwise Herbst differential inequality
\[
  \lambda\,\frac{M'(\lambda)}{M(\lambda)}-\log M(\lambda)
    \le c\lambda^2.
\]
If this differential inequality holds for every \(\lambda>0\), then the
corrected quotient
\[
  \lambda\longmapsto \frac{\log M(\lambda)}{\lambda}-c\lambda
\]
is antitone on \((0,\infty)\).
The endpoint value is formalized as
\[
  \frac{\log M(\lambda)}{\lambda}\longrightarrow \mathbb E X
  \qquad(\lambda\downarrow 0),
\]
and therefore the same differential inequality implies the finite Herbst
log-Laplace estimate
\[
  \log M(\lambda)\le \lambda\,\mathbb E X+c\lambda^2
  \qquad(\lambda>0).
\]
Equivalently, the formalization packages the direct entropy-to-log-Laplace
wrapper: the entropy bound
\[
  \operatorname{Ent}(e^{\lambda X})
  \le c\lambda^2 M(\lambda) \qquad(\lambda>0)
\]
itself implies the same log-Laplace estimate.
Also, any visible log-Laplace estimate
\[
  \log M(\lambda)\le \lambda\mu+R
\]
implies the centered MGF estimate
\[
  \mathbb E\exp\{\lambda(X-\mu)\}\le e^R.
\]
In particular, if \(\sigma>0\), \(t>0\), and
\[
  \log M(\lambda)\le \lambda\mu+\lambda^2\sigma^2/2
  \qquad(\lambda>0),
\]
then
\[
  \mathbb P\{X\le \mu+t\}
  \ge 1-\exp\!\left(-\frac{t^2}{2\sigma^2}\right).
\]
In the special case where the entropy bound has the Ledoux/Herbst normalization
\[
  \operatorname{Ent}(e^{\lambda X})
  \le \frac{\sigma^2}{2}\lambda^2 M(\lambda)
  \qquad(\lambda>0),
\]
the same tail bound follows directly.
\end{lemma}


\begin{formalized}
\lean{Analysis/FiniteProbability.lean}
\leanfit{FiniteProbability.expectationReal_exp_pos}
\leanfit{FiniteProbability.hasDerivAt_expectationReal_exp_mul}
\leanfit{FiniteProbability.hasDerivAt_log_expectationReal_exp_mul}
\leanfit{FiniteProbability.entropyReal}
\leanfit{FiniteProbability.entropyReal_exp_mul_eq}
\leanfit{FiniteProbability.log_mgf_differential_le_of_entropyReal_exp_mul_le}
\leanfit{FiniteProbability.log_mgf_div_sub_quadratic_antitoneOn_of_differential_le}
\leanfit{FiniteProbability.tendsto_log_mgf_div_nhdsGT_zero}
\leanfit{FiniteProbability.log_mgf_le_mean_add_quadratic_of_differential_le}
\leanfit{FiniteProbability.log_mgf_le_mean_add_quadratic_of_entropyReal_exp_mul_le}
\leanfit{FiniteProbability.expectationReal_exp_centered_le_exp_of_log_mgf_le}
\leanfit{FiniteProbability.eventProb_real_le_mean_add_ge_one_sub_exp_sq_of_log_mgf_bound}
\leanfit{FiniteProbability.eventProb_real_le_mean_add_ge_one_sub_exp_sq_of_entropyReal_exp_mul_le}
\end{formalized}

\begin{lemma}[Finite product entropy chain rule]
Let \(P\) and \(Q\) be finite probability laws on finite sets
\(\Omega\) and \(\Lambda\).  For any real function
\(Z:\Omega\times\Lambda\to\mathbb R\), define
\[
  \bar Z(\omega):=\mathbb E_Q Z(\omega,\cdot).
\]
Then the repository-native finite entropy functional satisfies the exact
product decomposition
\[
  \operatorname{Ent}_{P\otimes Q}(Z)
  =
  \mathbb E_P\!\left[
    \operatorname{Ent}_Q\bigl(Z(\omega,\cdot)\bigr)
  \right]
  +
  \operatorname{Ent}_P(\bar Z).
\]
\end{lemma}


\begin{formalized}
\lean{Analysis/FiniteProbability.lean}
\leanfit{FiniteProbability.prod_expectationReal_eq}
\leanfit{FiniteProbability.prod_expectationReal_fst_eq}
\leanfit{FiniteProbability.prod_expectationReal_snd_eq}
\leanfit{FiniteProbability.entropyReal_prod_eq_expectation_entropyReal_add_entropyReal_expectation}
\end{formalized}

\begin{lemma}[Scalar and fixed-vector independent-copy symmetrization]
Let \(P\) be a finite probability law on a finite set \(\Omega\).  Let
\((\omega,\omega')\) be distributed according to the product law
\(P\otimes P\).  For a scalar random variable \(X:\Omega\to\R\),
\[
  \E_P\!\left|X-\E_P X\right|
  \le
  \E_{P\otimes P}\!\left|X(\omega)-X(\omega')\right|.
\]
Consequently, if \(\E_P X=0\), then
\[
  \E_P |X|
  \le
  \E_{P\otimes P}\!\left|X(\omega)-X(\omega')\right|.
\]

Let \(d\in\mathbb N\), and let \(Y:\Omega\to\R^d\).  Write
\[
  \mu_i=\E_P Y_i
  \qquad (1\le i\le d),
  \qquad
  \|z\|_2=\left(\sum_{i=1}^d z_i^2\right)^{1/2}.
\]
Then the Euclidean vector form is
\[
  \E_P\!\left\|Y-\mu\right\|_2
  \le
  \E_{P\otimes P}\!\left\|Y(\omega)-Y(\omega')\right\|_2.
\]
In particular, if \(\E_PY_i=0\) for every coordinate \(i\), then
\[
  \E_P\|Y\|_2
  \le
  \E_{P\otimes P}\!\left\|Y(\omega)-Y(\omega')\right\|_2.
\]

For Algorithm~1, let \(A\in\R^{m\times n}\), let \(\widetilde A(\omega)\)
be the exact trace sketch produced from the exact squared-magnitude product
law, and fix \(x\in\R^n\).  With
\[
  Y(\omega)=(A-\widetilde A(\omega))x\in\R^m,
\]
the exact mean-zero residual identity gives
\[
  \E\!\left\|(A-\widetilde A)x\right\|_2
  \le
  \E\!\left\|
     (A-\widetilde A(\omega))x
     -(A-\widetilde A(\omega'))x
  \right\|_2.
\]
The law \(P\) is exact by convention.  All quantities in this lemma are exact
analysis objects; no floating-point operation or computed non-probability
quantity is introduced at this layer.  The result is fixed-vector
Khintchine-route infrastructure and is not a uniform spectral-norm tail for
equation~(2).
\end{lemma}

\begin{formalized}
\leanfit{FiniteProbability.expectationReal_abs_sub_mean_le_prod_expectationReal_abs_sub}
\leanfit{FiniteProbability.expectationReal_abs_le_prod_expectationReal_abs_sub_of_expectation_eq_zero}
\leanfit{sqMagTraceProbability_expectationReal_abs_rectSelfAdjointDilation_elementwiseTraceResidual_le_prod_abs_sub}
\leanfit{FiniteProbability.expectationReal_vecNorm2_mean_le_expectationReal_vecNorm2}
\leanfit{FiniteProbability.expectationReal_vecNorm2_sub_mean_le_prod_expectationReal_vecNorm2_sub}
\leanfit{FiniteProbability.expectationReal_vecNorm2_le_prod_expectationReal_vecNorm2_sub_of_expectation_eq_zero}
\leanfit{sqMagTraceProbability_expectationReal_vecNorm2_rectMatMulVec_elementwiseTraceResidual_le_prod_vecNorm2_sub}
\end{formalized}

\begin{lemma}[Operator copy-difference symmetrization bridge]
Let \(P\) be a finite probability law on a finite set \(\Omega\).  Let
\(m,n\in\mathbb N\), and let \(X:\Omega\to\R^{m\times n}\) be a
matrix-valued random variable satisfying
\[
  \E_P X_{ij}=0
  \qquad (1\le i\le m,\ 1\le j\le n).
\]
Fix an outcome \(\omega\in\Omega\) and a radius \(L\in\R\).  Suppose that
for every independent-copy outcome \(\eta\in\Omega\),
\[
  \|(X(\omega)-X(\eta))z\|_2
  \le L\|z\|_2
  \qquad\text{for every } z\in\R^n.
\]
Then
\[
  \|X(\omega)z\|_2\le L\|z\|_2
  \qquad\text{for every } z\in\R^n.
\]

For Algorithm~1, let \(P\) be the exact squared-magnitude product law and
write
\[
  R(\omega)=A-\widetilde A(\omega).
\]
Define the exact all-copy-differences event
\[
  \mathcal E_{\rm diff}(L)
  =
  \left\{\omega:
    \|(R(\omega)-R(\eta))z\|_2\le L\|z\|_2
    \text{ for every } \eta \text{ and every } z\in\R^n
  \right\}.
\]
If \(P(\mathcal E_{\rm diff}(L))\ge\rho\), then the exact Algorithm~1
spectral event
\[
  \left\{\omega:
    \|R(\omega)z\|_2\le L\|z\|_2
    \text{ for every } z\in\R^n
  \right\}
\]
also has probability at least \(\rho\).  This is exact-probability and
exact-arithmetic infrastructure.  No floating-point computation is modeled in
this bridge, and the probability tail for \(\mathcal E_{\rm diff}(L)\) is a
separate Khintchine-route target.
\end{lemma}

\begin{formalized}
\leanfit{FiniteProbability.rectOpNorm2Le_of_entrywise_mean_zero_of_copy_diff_rectOpNorm2Le}
\leanfit{sqMagTraceProbability_rectOpNorm2Le_elementwiseTraceResidual_of_all_copy_diffs}
\leanfit{algorithm1ExactAllCopyDiffSpectralEvent}
\leanfit{algorithm1ExactAllCopyDiffSpectralEvent_subset_exactSpectralEvent}
\leanfit{sqMagTraceProbability_eventProb_algorithm1ExactSpectralEvent_ge_of_all_copy_diff}
\end{formalized}

\begin{lemma}[Finite-test Rademacher signed-matrix quadratic forms]
Let \(m,n\in\mathbb N\).  Let
\[
  \Omega_m=\{-1,1\}^m
\]
with the exact uniform Rademacher law, and for \(\omega\in\Omega_m\) write
\(\varepsilon_k(\omega)\in\{-1,1\}\) for the \(k\)-th exact sign.  Let
\[
  M_k\in\R^{n\times n}\qquad (1\le k\le m)
\]
be a fixed exact matrix family.  For \(z\in\R^n\), define the exact quadratic
form
\[
  q(M,z)=z^\mathsf T Mz .
\]
Then, for every \(z\in\R^n\) and every exact sign outcome \(\omega\),
\[
  q\!\left(\sum_{k=1}^m\varepsilon_k(\omega)M_k,z\right)
  =
  \sum_{k=1}^m q(M_k,z)\,\varepsilon_k(\omega).
\]

Let \(\mathcal A\) be a finite index set.  For each \(a\in\mathcal A\), let
\(z_a\in\R^n\), \(T_a>0\), and \(\sigma_a>0\).  Assume the exact variance
proxy
\[
  \sum_{k=1}^m q(M_k,z_a)^2\le \sigma_a^2
  \qquad (a\in\mathcal A).
\]
Then the exact Rademacher law satisfies
\[
  \mathbb P\!\left\{
    \forall a\in\mathcal A,\quad
    \left|
      q\!\left(\sum_{k=1}^m\varepsilon_k M_k,z_a\right)
    \right|
    \le T_a
  \right\}
  \ge
  1-\sum_{a\in\mathcal A}
    2\exp\!\left(-\frac{T_a^2}{2\sigma_a^2}\right).
\]
This is the finite-test scalar-Hoeffding layer for a later matrix
Khintchine/Bernstein route.  It is exact-probability and exact-arithmetic
only: the signs are drawn from the exact law, the matrices and test vectors
are exact analysis objects, and no floating-point computation or computed
non-probability quantity appears in this lemma.  It does not by itself prove
an operator-norm or spectral event for Algorithm~1.
\end{lemma}

\begin{formalized}
\lean{Algorithms/RandNLA/Preconditioning.lean}
\leanfit{finiteQuadraticForm_rademacher_signed_matrix_sum_eq_sum}
\leanfit{rademacherTraceProbability_eventProb_forall_abs_finiteQuadraticForm_signed_matrix_sum_le_ge_one_sub_sum_two_mul_exp_neg_sq_div}
\end{formalized}

\begin{corollary}[Algorithm 1 Rademacherized stepwise copy-difference tests]
Let \(A\in\R^{m\times n}\), let \(s\in\mathbb N\), and let
\[
  S=(S_1,\ldots,S_s),\qquad S'=(S'_1,\ldots,S'_s)
\]
be two fixed exact Algorithm~1 traces.  For a one-step sample \(S_t=(i_t,j_t)\),
let \(Z_t(S_t)\in\R^{m\times n}\) denote the exact one-sample residual
increment
\[
  Z_t(S_t)=s^{-1}A-\frac{A_{i_tj_t}}{s\,p_{i_tj_t}}e_{i_t}e_{j_t}^{\mathsf T},
\]
with the formalized zero-probability branch used when the exact sampling
probability \(p_{i_tj_t}\) is zero.  Define the exact self-adjoint dilation
difference
\[
  \Delta_t
  =
  \begin{bmatrix}
  0&Z_t(S_t)-Z_t(S'_t)\\
  (Z_t(S_t)-Z_t(S'_t))^{\mathsf T}&0
  \end{bmatrix}
  \in\R^{(m+n)\times(m+n)}.
\]
For \(w\in\R^{m+n}\), define
\[
  q(\Delta,w)=w^{\mathsf T}\Delta w .
\]
Let \(\mathcal A\) be a finite test set.  For each \(a\in\mathcal A\), let
\(w_a\in\R^{m+n}\), \(T_a>0\), and \(\sigma_a>0\).  Assume
\[
  \sum_{t=1}^s q(\Delta_t,w_a)^2\le\sigma_a^2
  \qquad (a\in\mathcal A).
\]
Let \(\varepsilon_1,\ldots,\varepsilon_s\) be exact independent Rademacher
signs.  No additional computed quantity is randomized or charged in this
statement.  Then, under the exact Rademacher law,
\[
  \mathbb P\!\left\{
    \forall a\in\mathcal A,\quad
    \left|
      q\!\left(\sum_{t=1}^s\varepsilon_t\Delta_t,w_a\right)
    \right|
    \le T_a
  \right\}
  \ge
  1-\sum_{a\in\mathcal A}
    2\exp\!\left(-\frac{T_a^2}{2\sigma_a^2}\right).
\]
This corollary is exact-probability and exact-arithmetic only.  The traces
\(S,S'\), the matrices \(\Delta_t\), and the test vectors \(w_a\) are fixed
analysis objects; the only randomness in the statement is the exact
Rademacher sign law.  No floating-point computation or computed
non-probability quantity is introduced.  The corollary is the Algorithm~1
finite-test Rademacherization adapter for future copy-difference
Khintchine/Bernstein work; it is not yet the spectral-norm probability tail
for the all-copy-differences event.
\end{corollary}

\begin{formalized}
\lean{Algorithms/RandNLA/Preconditioning.lean}
\leanfit{finiteQuadraticForm_rademacher_signed_matrix_sum_eq_sum_fintype}
\leanfit{rademacherTraceProbability_eventProb_forall_abs_finiteQuadraticForm_signed_matrix_sum_fintype_le_ge_one_sub_sum_two_mul_exp_neg_sq_div}
\lean{Algorithms/RandNLA/ElementwiseSpectral.lean}
\leanfit{sqMagTraceProbability_eventProb_forall_abs_finiteQuadraticForm_rademacher_signed_rectSelfAdjointDilation_sampleResidualIncrement_diff_le_ge_one_sub_sum_two_mul_exp_neg_sq_div}
\end{formalized}

\begin{corollary}[Finite-cover Rademacher operator event for fixed Algorithm 1 copy differences]
Use the notation of the preceding corollary.  Let
\(\mathcal K=\{1,\ldots,m\}\sqcup\{1,\ldots,n\}\), and view each
\(\Delta_t\) as an exact self-adjoint matrix on \(\R^{\mathcal K}\).  Let
\(\mathcal A\) be a finite index set and let \(w_a\in\R^{\mathcal K}\)
\((a\in\mathcal A)\) be an exact finite cover of the Euclidean unit ball at
radius \(\rho\ge0\): for every \(x\in\R^{\mathcal K}\) with
\(\|x\|_2\le1\), some \(a\in\mathcal A\) satisfies
\(\|x-w_a\|_2\le\rho\).  Let \(\eta>0\), let \(L\ge0\), and let
\(\sigma_a>0\) for \(a\in\mathcal A\).  Define
\[
  D(\varepsilon)=\sum_{t=1}^s \varepsilon_t\Delta_t,
  \qquad
  R(\varepsilon)=
  \sum_{t=1}^s\varepsilon_t\bigl(Z_t(S_t)-Z_t(S'_t)\bigr).
\]
Thus \(D(\varepsilon)\) is the exact self-adjoint dilation of
\(R(\varepsilon)\).  Assume the exact deterministic coarse radius
\[
  \|D(\varepsilon)\|_{2\to2}\le L
  \qquad\text{for every exact sign outcome }\varepsilon,
\]
and the exact finite-test variance proxies
\[
  \sum_{t=1}^s q(\Delta_t,w_a)^2\le\sigma_a^2
  \qquad (a\in\mathcal A).
\]
Then, under the exact Rademacher sign law,
\[
  \mathbb P\!\left\{
    \|R(\varepsilon)\|_{2\to2}
    \le
    \eta+L(2\rho+\rho^2)
  \right\}
  \ge
  1-\sum_{a\in\mathcal A}
    2\exp\!\left(-\frac{\eta^2}{2\sigma_a^2}\right).
\]
This is an exact-law and exact-arithmetic cover upgrade from finite
quadratic-form tests to a rectangular operator event for fixed signed
copy-difference sums.  No floating-point computation, stored matrix, computed
basis, or probability-construction error appears.  The coarse radius \(L\) is
a visible deterministic exact hypothesis, not a computed certificate and not a
probability event.  Consequently this corollary is intermediate
infrastructure; it is not the final source-uniform Algorithm~1 spectral tail
or the CACM equation~(2) theorem.
\end{corollary}

\begin{formalized}
\lean{Analysis/MatrixAlgebra.lean}
\leanfit{finiteUnitBallCover}
\leanfit{finiteLoewnerLe_of_finite_unit_ball_cover_quadraticForm}
\lean{Algorithms/RandNLA/ElementwiseSpectral.lean}
\leanfit{rademacherTraceProbability_eventProb_rectOpNorm2Le_signed_sampleResidualIncrement_diff_ge_one_sub_sum_two_mul_exp_neg_sq_div_of_finiteUnitBallCover}
\end{formalized}

\begin{theorem}[Small-entry obstruction for the literal all-copy-differences event]
Let \(L\in\R\) be any proposed rectangular operator radius and define the exact
literal Algorithm~1 input
\[
  A_L=\begin{bmatrix}1&\alpha\end{bmatrix}\in\R^{1\times2},
  \qquad
  \alpha=(|L|+2)^{-1}.
\]
Let \(p_{01}\) denote the exact squared-magnitude sampling probability of the
tiny entry \(A_{0,1}=\alpha\).  Then \(p_{01}>0\).  Let \(S_{\rm tiny}\) be
the exact one-step trace that samples entry \((0,1)\), and let \(S_{\rm unit}\)
be the exact one-step trace that samples entry \((0,0)\).  If
\(R(S)=A_L-\widetilde A(S)\) denotes the exact one-step residual, then
\[
  \neg\bigl(\|R(S_{\rm tiny})-R(S_{\rm unit})\|_{2\to2}\le L\bigr).
\]
Equivalently, \(S_{\rm tiny}\) is not in the all-copy-differences event
\[
  \mathcal E_{\rm allcopy}(A_L,L)
  =
  \left\{
    S:\ \forall S',\
    \|R(S)-R(S')\|_{2\to2}\le L
  \right\}.
\]
Consequently, under the exact one-step squared-magnitude law for \(A_L\),
\[
  p_{01}
  \le
  \mathbb P\{S:S\notin \mathcal E_{\rm allcopy}(A_L,L)\}.
\]
Equivalently, for every \(\delta\in\R\), any lower bound
\[
  \mathbb P\{S:S\in\mathcal E_{\rm allcopy}(A_L,L)\}\ge1-\delta
\]
forces
\[
  p_{01}\le\delta .
\]
In particular, the all-copy-differences support event has positive failure
probability for this literal untruncated input.

This is an exact-law and exact-arithmetic obstruction result.  It introduces
no floating-point computation and no computed non-probability quantity.  It
does not say that every possible high-probability spectral tail is false; it
says that the literal untruncated all-copy support event cannot be treated as a
probability-one prerequisite, nor can the missing Algorithm~1 equation~(2)
tail be closed by hiding a uniform deterministic copy-difference radius.
\end{theorem}

\begin{formalized}
\lean{Algorithms/RandNLA/ElementwiseSpectral.lean}
\leanfit{algorithm1SmallEntrySupportMatrix_trace_residual_small_unit_diff_not_rectOpNorm2Le}
\leanfit{sqMagTraceProbability_eventProb_not_algorithm1ExactAllCopyDiffSpectralEvent_smallEntry_ge}
\leanfit{sqMagTraceProbability_eventProb_not_algorithm1ExactAllCopyDiffSpectralEvent_smallEntry_pos}
\leanfit{sqMagTraceProbability_algorithm1ExactAllCopyDiffSpectralEvent_smallEntry_delta_ge}
\end{formalized}

\begin{theorem}[Sample-count obstruction for the literal exact spectral event]
Let \(s\ge1\) be the Algorithm~1 sample count and also the update denominator
used in the exact rescaling.  Let \(L\in\R\) be any proposed rectangular
operator radius.  Define
\[
  A_L=\begin{bmatrix}1&\alpha\end{bmatrix}\in\R^{1\times2},
  \qquad
  \alpha=(|L|+2)^{-1},
\]
and let
\[
  p_{01}
  =
  \frac{\alpha^2}{1+\alpha^2}
\]
be the exact squared-magnitude probability of the tiny entry \(A_{0,1}\).
Let \(S_{\rm tiny}^{(s)}\) be the exact length-\(s\) trace that samples
\((0,1)\) at every step.  For any exact length-\(s\) trace \(S\), write
\[
  R_s(S)=A_L-\widetilde A_s(S)
\]
for the exact Algorithm~1 residual formed with denominator \(s\).  Then
\[
  \neg\bigl(\|R_s(S_{\rm tiny}^{(s)})\|_{2\to2}\le L\bigr).
\]
Consequently, under the exact independent length-\(s\) squared-magnitude law,
\[
  p_{01}^{\,s}
  \le
  \mathbb P\{S:\|R_s(S)\|_{2\to2}>L\}.
\]
Equivalently, the exact success probability satisfies the upper bound
\[
  \mathbb P\{S:\|R_s(S)\|_{2\to2}\le L\}
  \le
  1-p_{01}^{\,s}
  <
  1.
\]
Equivalently, if the exact spectral event
\[
  \mathcal E_{\rm spec}(A_L,L,s)
  =
  \{S:\|R_s(S)\|_{2\to2}\le L\}
\]
is claimed to satisfy
\[
  \mathbb P\{S:S\in\mathcal E_{\rm spec}(A_L,L,s)\}\ge1-\delta ,
\]
then the failure budget must obey
\[
  p_{01}^{\,s}\le\delta .
\]
For \(\delta>0\), this is equivalently bounded in logarithmic sample-count
form by the proved inequality
\[
  \log(1/\delta)
  \le
  s\,\log(1/p_{01})
  =
  s\,\log\!\left(\frac{1+\alpha^2}{\alpha^2}\right).
\]
Thus any such claimed literal exact spectral theorem must, at minimum, have
\[
  s
  \ge
  \frac{\log(1/\delta)}
       {\log\!\left((1+\alpha^2)/\alpha^2\right)}
  =
  \Omega_L(\log(1/\delta)),
\]
where the subscript means \(L\), hence \(\alpha=(|L|+2)^{-1}\), is fixed.
The divided inequality is formalized directly: \(0<p_{01}<1\), so the
logarithmic denominator is positive and no additional analytic assumption is
introduced.

Equivalently, the same Lean development proves the contrapositive forms:
if
\[
  \delta < p_{01}^{\,s},
\]
then the displayed \(1-\delta\) lower bound for
\(\mathcal E_{\rm spec}(A_L,L,s)\) cannot hold.  If \(\delta>0\) and
\[
  s
  <
  \frac{\log(1/\delta)}
       {\log\!\left((1+\alpha^2)/\alpha^2\right)},
\]
then the same \(1-\delta\) lower bound also cannot hold.  These are exact
arithmetic contradictions derived from the exact product-law mass of the
all-tiny trace; no floating-point operation or probability-construction error
appears in the statement.

The same exact product-mass argument also gives a fixed-sample existential
form: for every positive \(s\) and every \(L\), there is a positive
\(\delta\), for instance \(\delta=p_{01}^{\,s}/2\), for which the claimed
\(1-\delta\) lower bound cannot hold.

The formalization also includes a concrete source-budget witness.  For
\[
  A=\begin{bmatrix}1&1/102\end{bmatrix},\qquad
  s=1,\qquad L=100,\qquad \delta=1/30000,
\]
with \(m=1\), \(n=2\), and
\[
  C=2\max\{1,2\}+\frac{4\sqrt2}{3}\sqrt2,
\]
the rectangular source-style budget
\[
  4C\|A\|_F^2
    \log\!\left(\frac{2(1+2)}{1/30000}\right)
  \le 100^2
\]
is true in exact arithmetic, while the exact lower bound
\[
  \mathbb P\{S:\|R_1(S)\|_{2\to2}\le100\}
  \ge 1-\frac{1}{30000}
\]
is formally impossible for the literal squared-magnitude law.
Writing \(\mathcal B(A)\) for the displayed source-style budget at these
fixed parameters with \(A\) in place of the concrete witness, the
formalization also proves
\[
  \neg\Big[
    \forall A,
    \operatorname{sqMagProbDen}(A)>0
    \ \wedge\
    \mathcal B(A)
    \Longrightarrow
    \mathbb P\{S:\|R_1(S)\|_{2\to2}\le100\}
      \ge 1-\frac{1}{30000}
  \Big].
\]
Thus the failure is a formal refutation of the source-budget-only implication
at these fixed parameters, not just a warning about the displayed witness.

This is still an exact-law and exact-arithmetic obstruction, not a
floating-point theorem.  It is stronger than the previous all-copy support
obstruction because it applies directly to the exact spectral event itself.
It does not rule out a different matrix-tail proof, but it shows that any
literal untruncated source-uniform theorem must pay at least this all-tiny
trace mass, or else expose reciprocal-entry, truncation, or no-small-entry
dependence.
\end{theorem}

\begin{formalized}
\lean{Analysis/FiniteProbability.lean}
\leanfit{FiniteProbability.prob_le_eventProb_of_mem}
\lean{Algorithms/RandNLA/ElementwiseSpectral.lean}
\leanfit{sqMagTraceProbMass_algorithm1SmallEntrySupportMatrix_all_tiny}
\leanfit{algorithm1SmallEntrySupportMatrix_all_tiny_trace_residual_not_rectOpNorm2Le}
\leanfit{sqMagTraceProbability_eventProb_not_algorithm1ExactSpectralEvent_all_tiny_smallEntry_ge}
\leanfit{sqMagTraceProbability_algorithm1ExactSpectralEvent_all_tiny_smallEntry_delta_ge}
\leanfit{sqMagProb_algorithm1SmallEntrySupportMatrix_small_lt_one}
\leanfit{log_inv_delta_le_nat_mul_log_inv_of_pow_le}
\leanfit{log_inv_delta_div_log_inv_le_nat_of_pow_le}
\leanfit{sqMagTraceProbability_algorithm1ExactSpectralEvent_all_tiny_smallEntry_log_delta_le}
\leanfit{sqMagTraceProbability_algorithm1ExactSpectralEvent_all_tiny_smallEntry_sample_count_ge}
\leanfit{sqMagTraceProbability_not_algorithm1ExactSpectralEvent_all_tiny_smallEntry_of_delta_lt_pow}
\leanfit{sqMagTraceProbability_not_algorithm1ExactSpectralEvent_all_tiny_smallEntry_of_sample_count_lt}
\leanfit{sqMagTraceProbability_eventProb_algorithm1ExactSpectralEvent_all_tiny_smallEntry_le_one_sub_pow}
\leanfit{sqMagTraceProbability_eventProb_algorithm1ExactSpectralEvent_all_tiny_smallEntry_lt_one}
\leanfit{exists_delta_not_sqMagTraceProbability_algorithm1ExactSpectralEvent_all_tiny_smallEntry}
\leanfit{real_log_180000_le_18}
\leanfit{algorithm1RectSourceBudget_1_2_100_one_div_30000}
\leanfit{algorithm1SmallEntrySupportMatrix_100_small_prob_gt_one_div_30000}
\leanfit{algorithm1SmallEntrySupportMatrix_100_rect_source_budget_one_div_30000}
\leanfit{sqMagTraceProbability_not_algorithm1ExactSpectralEvent_rect_source_budget_witness}
\leanfit{not_forall_algorithm1ExactSpectralEvent_of_rect_source_budget_one_div_30000}
\end{formalized}

\begin{lemma}[Unbiased Bernoulli coordinate law]
Let \(B=\{0,1\}\) and let \(\nu\) be the unbiased probability law on \(B\):
\(\nu(\{0\})=\nu(\{1\})=1/2\).  For every \(X:B\to\mathbb R\),
\[
  \mathbb E_{\nu} X
    = \frac{X(0)+X(1)}{2}.
\]
For every \(Z:B\to\mathbb R\), the finite entropy functional satisfies
\[
  \operatorname{Ent}_{\nu}(Z)
    =
    \frac{Z(0)\log Z(0)+Z(1)\log Z(1)}{2}
    -
    \frac{Z(0)+Z(1)}{2}
      \log\!\left(\frac{Z(0)+Z(1)}{2}\right).
\]
\end{lemma}


\begin{formalized}
\lean{Analysis/FiniteProbability.lean}
\leanfit{FiniteProbability.boolUniformProbability}
\leanfit{FiniteProbability.boolUniformProbability_prob}
\leanfit{FiniteProbability.boolUniformProbability_expectationReal}
\leanfit{FiniteProbability.entropyReal_boolUniformProbability_eq}
\end{formalized}

\begin{lemma}[Two-point Bernoulli log-Sobolev inequality]
For \(a,b>0\), the two-point entropy satisfies
\[
  \frac{a\log a+b\log b}{2}
  - \frac{a+b}{2}\log\!\left(\frac{a+b}{2}\right)
  \le
  \frac{(a-b)^2}{2(a+b)}.
\]
Consequently, if \(g:B\to(0,\infty)\), then
\[
  \operatorname{Ent}_{\nu}(g^2)
    \le \bigl(g(1)-g(0)\bigr)^2 .
\]
\end{lemma}


\begin{formalized}
\lean{Analysis/FiniteProbability.lean}
\leanfit{FiniteProbability.twoPointEntropy_le_sq_sub_div_of_pos}
\leanfit{FiniteProbability.entropyReal_boolUniformProbability_sq_le_sq_sub_of_pos}
\end{formalized}

\begin{lemma}[One-coordinate Bernoulli tensorization peel-off]
Let \(P\) be any finite probability law on \(\Omega\), let
\(\nu\) be the unbiased Bernoulli law on \(B=\{0,1\}\), and let
\(g:\Omega\times B\to(0,\infty)\).  Then
\[
  \operatorname{Ent}_{P\otimes\nu}(g^2)
  \le
  \mathbb E_P\!\left[
    \bigl(g(\omega,1)-g(\omega,0)\bigr)^2
  \right]
  +
  \operatorname{Ent}_P\!\left(
    \mathbb E_{\nu} g(\omega,\cdot)^2
  \right).
\]
\end{lemma}


\begin{formalized}
\lean{Analysis/FiniteProbability.lean}
\leanfit{FiniteProbability.entropyReal_prod_boolUniformProbability_sq_le_coordinate_add_entropy}
\end{formalized}

\begin{lemma}[Finite \(L^2\) section-norm bridge]
Let \(P\) be a finite probability law on \(\Omega\).  For real-valued
functions \(X,Y:\Omega\to\mathbb R\), the finite \(L^2(P)\) seminorm satisfies
the reverse triangle inequality
\[
  \left|
    \sqrt{\mathbb E_P X^2}
      - \sqrt{\mathbb E_P Y^2}
  \right|
  \le
  \sqrt{\mathbb E_P (X-Y)^2}.
\]
In particular, for the unbiased Bernoulli coordinate law \(\nu\),
\[
  \left|
    \sqrt{\mathbb E_{\nu} u^2}
      - \sqrt{\mathbb E_{\nu} v^2}
  \right|
  \le
  \sqrt{\mathbb E_{\nu}(u-v)^2}.
\]
\end{lemma}


\begin{formalized}
\lean{Analysis/FiniteProbability.lean}
\leanfit{FiniteProbability.expectationReal_sq_nonneg}
\leanfit{FiniteProbability.abs_expectationReal_mul_le_sqrt_mul_sqrt}
\leanfit{FiniteProbability.sqrt_expectationReal_sq_add_le}
\leanfit{FiniteProbability.abs_sqrt_expectationReal_sq_sub_le_sqrt_expectationReal_sub_sq}
\leanfit{FiniteProbability.boolUniformProbability_abs_sqrt_expectationReal_sq_sub_le_sqrt_expectationReal_sub_sq}
\end{formalized}

\begin{lemma}[Abstract Bernoulli-product induction lift]
Let \(P\) be a finite probability law on \(\Omega\), let
\(\nu\) be the unbiased Bernoulli law, and let
\(\{T_i:\Omega\to\Omega\}_{i\in I}\) be a finite family of coordinate moves.
Assume that \(P\) satisfies the entropy-gradient bound
\[
  \operatorname{Ent}_P(h^2)
  \le
  \sum_{i\in I}\mathbb E_P\!\left[(h(\omega)-h(T_i\omega))^2\right]
\]
for every positive \(h:\Omega\to(0,\infty)\).  Then every positive
\(g:\Omega\times\{0,1\}\to(0,\infty)\) satisfies
\[
\begin{aligned}
  \operatorname{Ent}_{P\otimes\nu}(g^2)
  &\le
  \mathbb E_P\!\left[
    \bigl(g(\omega,1)-g(\omega,0)\bigr)^2
  \right] \\
  &\quad+
  \sum_{i\in I}\mathbb E_{P\otimes\nu}\!\left[
    \bigl(g(\omega,b)-g(T_i\omega,b)\bigr)^2
  \right].
\end{aligned}
\]
\end{lemma}


\begin{formalized}
\lean{Analysis/FiniteProbability.lean}
\leanfit{FiniteProbability.entropyReal_prod_boolUniformProbability_sq_le_lifted_diff_sum_add}
\end{formalized}

\begin{theorem}[Concrete Rademacher-cube entropy-gradient inequality]
\label{thm:concrete-rademacher-cube-entropy}
Let \(\Omega_m=\{0,1\}^m\) with the uniform Rademacher law \(\mu_m\).  For
\(\omega\in\Omega_m\), let \(\omega^{(i)}\) denote the trace obtained by
flipping coordinate \(i\).  If \(g:\Omega_m\to(0,\infty)\), then
\[
  \operatorname{Ent}_{\mu_m}(g^2)
  \le
  \sum_{i=1}^m
    \mathbb E_{\mu_m}\!\left[
      \bigl(g(\omega)-g(\omega^{(i)})\bigr)^2
    \right].
\]
\end{theorem}


\begin{formalized}
\lean{Algorithms/RandNLA/Preconditioning.lean}
\leanfit{rademacherTraceFlip}
\leanfit{rademacherTraceProbMass_snoc}
\leanfit{rademacherTraceProbability_expectationReal_succ_last_eq}
\leanfit{rademacherTraceProbability_expectationReal_succ_eq_prod}
\leanfit{rademacherTraceProbability_entropyReal_succ_eq_prod}
\leanfit{rademacherTraceFlip_castSucc_snoc}
\leanfit{rademacherTraceFlip_last_snoc}
\leanfit{rademacherTraceProbability_entropyReal_sq_le_sum_flip}
\end{formalized}

\begin{corollary}[Conditional exponential-tilt reduction on the Rademacher cube]
\label{cor:conditional-tilt}
Let \(X:\Omega_m\to\mathbb R\), let \(\lambda\in\mathbb R\), and let
\(c\in\mathbb R\).  Assume the deterministic tilt flip-gradient estimate
\[
  \sum_{i=1}^m
    \mathbb E_{\mu_m}\!\left[
      \left(
        \exp\{\lambda X(\omega)/2\}
        -
        \exp\{\lambda X(\omega^{(i)})/2\}
      \right)^2
    \right]
  \le
  c\lambda^2\,
  \mathbb E_{\mu_m}\exp\{\lambda X(\omega)\}.
\]
Then
\[
  \operatorname{Ent}_{\mu_m}\!\bigl(\exp\{\lambda X\}\bigr)
  \le
  c\lambda^2\,
  \mathbb E_{\mu_m}\exp\{\lambda X\}.
\]
Consequently, if \(\sigma>0\) and, for every \(\lambda>0\), the displayed
flip-gradient estimate holds with \(c=\sigma^2/2\), then for every \(t>0\),
\[
  \mu_m\{\,X\le \mathbb E_{\mu_m}X+t\,\}
  \ge
  1-\exp\!\left\{-\frac{t^2}{2\sigma^2}\right\}.
\]
\end{corollary}


\begin{formalized}
\lean{Algorithms/RandNLA/Preconditioning.lean}
\leanfit{rademacherTraceProbability_entropyReal_exp_mul_le_of_flip_tilt_sq_sum_bound}
\leanfit{rademacherTraceProbability_eventProb_real_le_mean_add_ge_one_sub_exp_sq_of_flip_tilt_sq_sum_bound}
\end{formalized}

\begin{lemma}[Scalar and finite-cube tilt-gradient symmetrization]
For real \(a,b\),
\[
  \left(e^{a/2}-e^{b/2}\right)^2
  \le
  2\left(\frac a2-\frac b2\right)^2
  \left(e^a+e^b\right).
\]
Consequently, on the finite Rademacher cube, if coefficients \(L_i\ge0\)
satisfy
\[
  |X(\omega)-X(\omega^{(i)})|\le L_i
\]
for every coordinate flip, and if
\[
  2\sum_i \frac{\lambda^2}{2}L_i^2 \le c\lambda^2,
\]
then the flip-gradient hypothesis in Corollary~\ref{cor:conditional-tilt}
holds with that \(c\).
\end{lemma}


\begin{formalized}
\lean{Analysis/FiniteProbability.lean}
\leanfit{real_exp_sub_one_le_mul_exp}
\leanfit{real_abs_exp_sub_exp_le_abs_sub_mul_exp_add_exp}
\leanfit{real_exp_half_sub_sq_le_two_mul_half_diff_sq}
\lean{Algorithms/RandNLA/Preconditioning.lean}
\leanfit{rademacherTraceProbability_expectationReal_flip}
\leanfit{rademacherTraceProbability_flip_tilt_sq_sum_bound_of_pointwise_halfdiff_sq_le}
\leanfit{rademacherTraceProbability_flip_tilt_sq_sum_bound_of_pointwise_absdiff_le}
\end{formalized}

\begin{lemma}[Finite subgaussian Chernoff tail]
Let \(X\) be a real random variable on a finite probability space.  Suppose
that for every \(\lambda>0\),
\[
  \mathbb E\exp\{\lambda(X-\mu)\}
    \le \exp\{\lambda^2\sigma^2/2\},
  \qquad \sigma>0.
\]
Then for every \(t>0\),
\[
  \mathbb P\{X\le \mu+t\}
    \ge 1-\exp\!\left(-\frac{t^2}{2\sigma^2}\right).
\]
\end{lemma}


\begin{formalized}
\lean{Analysis/FiniteProbability.lean}
\leanfit{FiniteProbability.eventProb_real_le_ge_one_sub_exp_of_mgf_bound}
\leanfit{FiniteProbability.eventProb_real_le_mean_add_ge_one_sub_exp_sq_of_subgaussian_mgf}
\end{formalized}

\begin{lemma}[Finite spectral Loewner bridge]
Let \(M\in\mathbb R^{d\times d}\) be symmetric, and let
\(\lambda_1(M),\ldots,\lambda_d(M)\) denote the Hermitian eigenvalues supplied
by mathlib.  The formalization proves
\[
  M\succeq 0
  \quad\Longleftrightarrow\quad
  \lambda_i(M)\ge 0 \quad\text{for all } i .
\]
More generally, for symmetric \(M,N\),
\[
  M\preceq N
  \quad\Longleftrightarrow\quad
  \lambda_i(N-M)\ge 0 \quad\text{for all } i .
\]
In particular,
\[
  M\preceq L I
  \quad\Longleftrightarrow\quad
  \lambda_i(LI-M)\ge 0 \quad\text{for all } i .
\]
This is the deterministic endpoint needed by a future largest-eigenvalue tail
bound.  It does not give a probability estimate for the event \(M\preceq LI\).
\end{lemma}

\begin{formalized}
\leanfit{finitePSD_iff_finiteHermitianEigenvalues_nonneg}
\leanfit{finiteLoewnerLe_iff_sub_finiteHermitianEigenvalues_nonneg}
\leanfit{finiteScalarUpperDiffEigenvalues}
\leanfit{finiteLoewnerLe_smul_id_iff_finiteScalarUpperDiffEigenvalues_nonneg}
\leanfit{finiteLoewnerLe_smul_id_iff_sub_finiteHermitianEigenvalues_nonneg}
\end{formalized}

\begin{lemma}[Scalar matrix-exponential normalization]
Let \(I_d\) be the \(d\times d\) identity matrix and let \(L\in\mathbb R\).
The formalization defines a repository-native matrix exponential
\(\operatorname{Exp}(M)\) by routing finite matrices through mathlib's matrix
exponential, and proves
\[
  \operatorname{Exp}(L I_d)=e^L I_d,
  \qquad
  \operatorname{tr}(\operatorname{Exp}(L I_d))=d e^L .
\]
For a diagonal matrix \(\operatorname{diag}(v)\), it also proves
\[
  \operatorname{Exp}(\operatorname{diag}(v))
    =\operatorname{diag}(i\mapsto e^{v_i}),
  \qquad
  \operatorname{tr}(\operatorname{Exp}(\operatorname{diag}(v)))
    =\sum_i e^{v_i}.
\]
The same formal layer proves that \(\operatorname{Exp}(M)\) is symmetric when
\(M\) is symmetric.
This is the scalar normalization term that occurs in trace-MGF and
matrix-Bernstein proofs.  It is not a trace-exponential domination theorem and
does not by itself imply a largest-eigenvalue tail bound.
\end{lemma}

\begin{formalized}
\leanfit{finiteDiagonal}
\leanfit{finiteMatrixExp}
\leanfit{finiteMatrixExp_smul_finiteIdMatrix}
\leanfit{finiteTrace_finiteMatrixExp_smul_finiteIdMatrix}
\leanfit{finiteMatrixExp_symmetric}
\leanfit{finiteMatrixExp_finiteDiagonal}
\leanfit{finiteTrace_finiteMatrixExp_finiteDiagonal}
\end{formalized}

\begin{lemma}[Hermitian spectral-calculus trace diagonalization]
Let \(M\in\mathbb R^{d\times d}\) be symmetric, and let
\(\lambda_1(M),\ldots,\lambda_d(M)\) be the Hermitian eigenvalues used by the
repository.  The formalization defines the Hermitian continuous-functional-
calculus exponential \(E_{\mathrm{cfc}}(M)\) and proves
\[
  \operatorname{tr}(E_{\mathrm{cfc}}(M))
  =
  \sum_{i=1}^{d} e^{\lambda_i(M)} .
\]
This is the diagonalized trace identity normally used inside trace-MGF
arguments.  It is still not a trace-exponential domination theorem for random
matrix sums, and it does not close equation~(2).
\end{lemma}

\begin{formalized}
\leanfit{finiteHermitianCfcExp}
\leanfit{finiteTrace_finiteHermitianCfcExp_eq_sum_exp_finiteHermitianEigenvalues}
\end{formalized}

\begin{lemma}[Power-series trace diagonalization]
Let \(M\in\mathbb R^{d\times d}\) be symmetric, and let
\(\operatorname{Exp}(M)\) denote the repository-native matrix exponential
defined by mathlib's power-series matrix exponential.  Then
\[
  \operatorname{tr}(\operatorname{Exp}(M))
  =
  \sum_{i=1}^{d} e^{\lambda_i(M)} .
\]
It also proves the lower-tail companion
\[
  \operatorname{tr}(\operatorname{Exp}(-M))
  =
  \sum_{i=1}^{d} e^{-\lambda_i(M)} .
\]
The proof diagonalizes \(M\) by the Hermitian spectral theorem, applies
invariance of the matrix exponential under unitary conjugation, and then uses
trace cyclicity and the diagonal exponential formula.  This closes the
power-series trace-diagonalization dependency.  It is still not a
trace-exponential domination theorem for random matrix sums, and it does not
close equation~(2).
\end{lemma}

\begin{formalized}
\leanfit{finiteTrace_finiteMatrixExp_eq_sum_exp_finiteHermitianEigenvalues}
\leanfit{finiteTrace_finiteMatrixExp_neg_eq_sum_exp_neg_finiteHermitianEigenvalues}
\end{formalized}

\begin{lemma}[Scalar Loewner bound to trace-exponential bound]
If \(M\in\mathbb R^{d\times d}\) is symmetric and
\[
  M\preceq cI,
\]
then every Hermitian eigenvalue satisfies \(\lambda_i(M)\le c\), and hence
\[
  \operatorname{tr}(\operatorname{Exp}(M))\le d e^c .
\]
The lower-tail companion is stated as
\[
  -M\preceq cI
  \quad\Longrightarrow\quad
  \operatorname{tr}(\operatorname{Exp}(-M))\le d e^c .
\]
This is the deterministic scalarization step that a future matrix-CGF
Loewner estimate can feed into the trace-MGF interface.
\end{lemma}

\begin{formalized}
\leanfit{finiteQuadraticForm_finiteHermitianEigenvector_eq_eigenvalue_mul_norm_sq}
\leanfit{finiteVecNorm2Sq_finiteHermitianEigenvector_eq_one}
\leanfit{finiteHermitianEigenvalues_le_of_finiteLoewnerLe_smul_id}
\leanfit{finiteTrace_finiteMatrixExp_le_card_mul_exp_of_finiteLoewnerLe_smul_id}
\leanfit{finiteTrace_finiteMatrixExp_neg_le_card_mul_exp_of_neg_finiteLoewnerLe_smul_id}
\end{formalized}

\begin{lemma}[Trace-exponential eigenvalue tail interface]
Let \(M(\omega)\in\mathbb R^{d\times d}\) be a finite random symmetric matrix
family and let \(T\in\mathbb R\).  The formalization proves
\[
  \Prob\{\exists i:\lambda_i(M(\omega))\ge T\}
  \le
  e^{-T}\,
  \E\!\left[\operatorname{tr}(\operatorname{Exp}(M(\omega)))\right].
\]
Consequently, if
\[
  \E\!\left[\operatorname{tr}(\operatorname{Exp}(M(\omega)))\right]\le B,
\]
then
\[
  \Prob\{\exists i:\lambda_i(M(\omega))\ge T\}\le e^{-T}B.
\]
The same formal layer also proves the high-probability complement
\[
  \Prob\{\forall i:\lambda_i(M(\omega))<T\}
  \ge
  1-e^{-T}\,
  \E\!\left[\operatorname{tr}(\operatorname{Exp}(M(\omega)))\right],
\]
and its scalar-bound version with \(B\).  These are useful final-mile forms
for any future trace-MGF domination theorem.
The lower-tail companion is also formalized:
\[
  \Prob\{\exists i:\lambda_i(M(\omega))\le T\}
  \le
  e^{T}\,
  \E\!\left[\operatorname{tr}(\operatorname{Exp}(-M(\omega)))\right],
\]
and therefore
\[
  \Prob\{\forall i:T<\lambda_i(M(\omega))\}
  \ge
  1-e^{T}\,
  \E\!\left[\operatorname{tr}(\operatorname{Exp}(-M(\omega)))\right],
\]
again with scalar-bound versions.
Combining the upper and lower complements gives the two-sided event
\[
  \Prob\{\forall i:|\lambda_i(M(\omega))|<T\}
  \ge
  1-e^{-T}\E[\operatorname{tr}(\operatorname{Exp}(-M(\omega)))]
    -e^{-T}\E[\operatorname{tr}(\operatorname{Exp}(M(\omega)))] ,
\]
with a scalar-bound version obtained by replacing the two expectations by
proved bounds \(B_{-}\) and \(B_{+}\).
This is scalar Markov applied to the nonnegative random variable
\(\operatorname{tr}(\operatorname{Exp}(\pm M(\omega)))\), using the
deterministic witnesses
\[
  \lambda_i(M)\ge T
  \quad\Longrightarrow\quad
  e^T \le \operatorname{tr}(\operatorname{Exp}(M)).
\]
and
\[
  \lambda_i(M)\le T
  \quad\Longrightarrow\quad
  e^{-T} \le \operatorname{tr}(\operatorname{Exp}(-M)).
\]
It is a genuine MGF-to-eigenvalue bridge, but it still requires a separate
trace-MGF domination theorem to bound the expectation on the right.
\end{lemma}

\begin{formalized}
\leanfit{exp_le_finiteTrace_finiteMatrixExp_of_finiteHermitianEigenvalue_ge}
\leanfit{finiteTrace_finiteMatrixExp_nonneg}
\leanfit{exp_neg_le_finiteTrace_finiteMatrixExp_neg_of_finiteHermitianEigenvalue_le}
\leanfit{finiteTrace_finiteMatrixExp_neg_nonneg}
\leanfit{FiniteProbability.eventProb_exists_finiteHermitianEigenvalue_ge_le_exp_neg_mul_expected_trace_exp}
\leanfit{FiniteProbability.eventProb_exists_finiteHermitianEigenvalue_ge_le_exp_neg_mul_trace_bound}
\leanfit{FiniteProbability.eventProb_forall_finiteHermitianEigenvalue_lt_ge_one_sub_exp_neg_mul_expected_trace_exp}
\leanfit{FiniteProbability.eventProb_forall_finiteHermitianEigenvalue_lt_ge_one_sub_exp_neg_mul_trace_bound}
\leanfit{FiniteProbability.eventProb_exists_finiteHermitianEigenvalue_le_le_exp_mul_expected_trace_exp_neg}
\leanfit{FiniteProbability.eventProb_exists_finiteHermitianEigenvalue_le_le_exp_mul_trace_bound}
\leanfit{FiniteProbability.eventProb_forall_finiteHermitianEigenvalue_gt_ge_one_sub_exp_mul_expected_trace_exp_neg}
\leanfit{FiniteProbability.eventProb_forall_finiteHermitianEigenvalue_gt_ge_one_sub_exp_mul_trace_bound}
\leanfit{FiniteProbability.eventProb_forall_abs_finiteHermitianEigenvalue_lt_ge_one_sub_exp_neg_mul_expected_trace_exp_add}
\leanfit{FiniteProbability.eventProb_forall_abs_finiteHermitianEigenvalue_lt_ge_one_sub_exp_neg_mul_trace_bound_add}
\end{formalized}

\begin{corollary}[Eigenvalue form of the dilation upper event]
Let \(R=A-\widetilde A\) be the exact Algorithm~1 residual and let
\[
  D(R)=
  \begin{bmatrix}
    0 & R\\
    R^{T} & 0
  \end{bmatrix}.
\]
If \(\varepsilon\ge 0\) and every Hermitian eigenvalue of
\(\varepsilon I-D(R)\) is nonnegative, then
\[
  D(R)\preceq \varepsilon I
  \qquad\Longrightarrow\qquad
  \|R\|_2\le \varepsilon .
\]
Consequently, any probability lower bound for this eigenvalue event transfers
to the exact Algorithm~1 spectral residual event and, after the entrywise
floating-point perturbation budget is supplied, to the floating-point spectral
event.  The formalization also includes a finite union-bound form: if each
single eigenvalue inequality has failure budget \(\delta_a\), then the full
eigenvalue event has failure budget \(\sum_a\delta_a\).  This corollary is
only an event adapter: it does not prove the high-probability eigenvalue event
or the single-eigenvalue scalar tail bounds.
\end{corollary}

\begin{formalized}
\leanfit{algorithm1ExactDilationEigenUpperEvent}
\leanfit{algorithm1ExactDilationEigenUpperIndexEvent}
\leanfit{algorithm1ExactDilationEigenUpperEvent_subset_exactDilationUpperEvent}
\leanfit{algorithm1ExactDilationEigenUpperEvent_subset_exactSpectralEvent}
\leanfit{probability_algorithm1_exact_spectral_of_dilation_eigen_upper}
\leanfit{probability_algorithm1_exact_dilation_eigen_upper_of_index_bounds}
\leanfit{probability_algorithm1_exact_spectral_of_dilation_eigen_upper_index_bounds}
\leanfit{probability_algorithm1_fl_spectral_of_exact_dilation_eigen_upper}
\leanfit{probability_algorithm1_fl_spectral_of_exact_dilation_eigen_upper_index_bounds}
\end{formalized}

\begin{remark}[External concentration source chain]
The next matrix-concentration target is guided by Tropp,
\emph{User-friendly tail bounds for sums of random matrices},
Foundations of Computational Mathematics 12 (2012), 389--434,
\url{https://arxiv.org/abs/1004.4389}.  The relevant source locations are
Theorem 3.2 (Lieb trace concavity), Corollary 3.3 (the one-step
trace-expectation inequality), Theorem 3.6 (master tail bound for independent
self-adjoint sums), Corollary 3.7 (trace-exponential tail bound),
Theorem 1.4 (matrix Bernstein), and Theorem 1.6 (rectangular matrix
Bernstein via self-adjoint dilation).  Tropp's monograph
\emph{An Introduction to Matrix Concentration Inequalities},
\url{https://tropp.caltech.edu/books/Tro14-Introduction-Matrix-preprint.pdf},
Theorem 8.1.1, records the same Lieb concavity theorem and expands the
relative-entropy proof route.
These references guided theorem design only: no Lean theorem in this
repository cites them as a hypothesis.  The source-aligned square theorem below
formalizes the needed Tropp/Lieb-style trace-MGF, Bernstein, and truncation
steps locally.  An older Ahlswede--Winter operator Chernoff route remains
recorded in the ledger as an advisory alternative, but it is no longer the
active route for the closed hard-thresholded source-variant theorem.
The local symmetry, PSD, and Loewner predicates for finite square matrices are
now connected to mathlib's \(\texttt{Matrix.IsSymm}\),
\(\texttt{Matrix.IsHermitian}\), and \(\texttt{Matrix.PosSemidef}\) APIs, so
the next formal trace-exponential or largest-eigenvalue proof can use mathlib's
spectral and matrix-order vocabulary without changing the RandNLA algorithm
statements.  A small spectral bridge also names the Hermitian eigenvalues of a
local symmetric finite matrix, proves trace equals their sum, and identifies
local PSD and scalar-identity Loewner upper bounds with nonnegativity of the
Hermitian eigenvalues of the corresponding difference matrices.  The same
bridge now includes the scalar matrix-exponential normalization
\(\operatorname{tr}(\exp(LI))=d e^L\) and preservation of symmetry by matrix
exponentiation, plus the diagonal identity
\(\operatorname{tr}(\exp(\operatorname{diag}(v)))=\sum_i e^{v_i}\).  It also
includes the Hermitian spectral-calculus identity
\(\operatorname{tr}(E_{\mathrm{cfc}}(M))=\sum_i e^{\lambda_i(M)}\), and the
repository-native matrix-exponential trace identity
\(\operatorname{tr}(\operatorname{Exp}(M))=\sum_i e^{\lambda_i(M)}\).  These
bridges are foundational only; they are not themselves concentration theorems.
For source alignment, Drineas--Zouzias,
\emph{A Note on Element-wise Matrix Sparsification via a Matrix-valued
Bernstein Inequality}, arXiv:1006.0407, Algorithm~1 and Theorem~1, uses a
truncated matrix \(\hat A\): small entries are zeroed out before sampling from
\(\hat A_{ij}^2/\|\hat A\|_F^2\).  The deterministic truncation transfer layer,
the matrix-Bernstein half-budget event for sampling \(\hat A\), and the
support-aware floating-point transfer are formalized below.  The untruncated
user-requested law
\(A_{ij}^2/\|A\|_F^2\) remains a separate theorem target.
\end{remark}

\begin{lemma}[Operator-log monotonicity bridge]
For complex \(d\times d\) C\(^*\)-matrices \(X,Y\), if \(0<X\preceq Y\), then
\[
  \log X \preceq \log Y .
\]
\end{lemma}


\begin{formalized}
\lean{Analysis/OperatorLog.lean}
\leanfit{cstarMatrix_log_le_log}
\end{formalized}

\begin{lemma}[Self-adjoint logarithm--exponential inverse]
Let \(X\) be a self-adjoint complex \(d\times d\) C\(^*\)-matrix.  Then
\[
  \log(\exp X)=X .
\]
The same identity is formalized both for the normed-algebra exponential and
for the real and complex continuous-functional-calculus exponentials:
\[
  \log(\exp_{\mathrm{cfc},\mathbb C}X)=X,
  \qquad
  \log(\exp_{\mathrm{cfc},\mathbb R}X)=X .
\]
\end{lemma}


\begin{formalized}
\lean{Analysis/OperatorLog.lean}
\leanfit{cstarMatrix_normedSpaceExp_isTopologicalRing}
\leanfit{cstarMatrix_normedRingExp_isTopologicalRing}
\leanfit{cstarMatrix_realContinuousFunctionalCalculus}
\leanfit{cstarMatrix_log_normedSpace_exp_of_isSelfAdjoint}
\leanfit{cstarMatrix_log_cfc_complex_exp_of_isSelfAdjoint}
\leanfit{cstarMatrix_log_cfc_real_exp_of_isSelfAdjoint}
\end{formalized}

\begin{lemma}[Self-adjoint matrix exponentials are strictly positive]
Let \(X\) be a self-adjoint complex \(d\times d\) C\(^*\)-matrix.  Then
\[
  \exp X \succ 0 .
\]
The same strict-positivity statement is formalized for the normed-algebra
exponential and for the real and complex continuous-functional-calculus
exponentials:
\[
  \exp_{\mathrm{cfc},\mathbb C}(X)\succ 0,
  \qquad
  \exp_{\mathrm{cfc},\mathbb R}(X)\succ 0 .
\]
Consequently, for self-adjoint \(H\), the matrix
\[
  \exp_{\mathrm{cfc},\mathbb C}(H+\log A)
\]
appearing in the local Lieb trace functional is strictly positive.
\end{lemma}


\begin{formalized}
\lean{Analysis/OperatorLog.lean}
\leanfit{cstarMatrix_normedSpace_exp_isStrictlyPositive_of_isSelfAdjoint}
\leanfit{cstarMatrix_cfc_complex_exp_isStrictlyPositive_of_isSelfAdjoint}
\leanfit{cstarMatrix_cfc_real_exp_isStrictlyPositive_of_isSelfAdjoint}
\lean{Analysis/LiebTrace.lean}
\leanfit{liebTraceCfcExp_isStrictlyPositive}
\end{formalized}

\begin{theorem}[Scalar Bernstein parabola]
Let \(a\ge 0\).  For every \(x\le 1\),
\[
  e^{a x}\le 1+a x+(e^a-a-1)x^2 .
\]
Equivalently, if \(\theta\ge 0\), \(R>0\), and \(x\le R\), then
\[
  e^{\theta x}
  \le
  1+\theta x+
  \frac{e^{\theta R}-\theta R-1}{R^2}\,x^2 .
\]
\end{theorem}


\begin{formalized}
\lean{Analysis/LiebTrace.lean}
\leanfit{real_exp_quadratic_remainder_monotone}
\leanfit{real_exp_tail_two_hasSum}
\leanfit{real_exp_mul_le_quadratic_of_nonneg_of_le_one}
\leanfit{real_exp_mul_le_quadratic_scaled_of_nonneg_of_pos_of_le}
\end{formalized}

\begin{lemma}[CFC Bernstein parabola lift]
Let \(X\) be a self-adjoint finite complex C\(^*\)-matrix and let
\(\theta,\beta\in\mathbb R\).  Suppose the scalar inequality
\[
  \exp(\theta x)\le 1+\theta x+\beta x^2
  \qquad\text{for every }x\in\sigma_{\mathbb R}(X)
\]
holds on the real spectrum of \(X\).  Then the same inequality lifts to
Loewner order:
\[
  \exp_{\mathrm{cfc},\mathbb R}(\theta X)
  \preceq
  I+\theta X+\beta X^2 .
\]
The formalization also proves the normalization identity
\[
  \operatorname{cfc}\!\left(x\mapsto 1+\theta x+\beta x^2,X\right)
  =
  I+\theta X+\beta X^2 .
\]
\end{lemma}


\begin{formalized}
\lean{Analysis/LiebTrace.lean}
\leanfit{cstarMatrix_cfc_quadratic_eq}
\leanfit{cstarMatrix_cfc_real_exp_mul_le_quadratic_of_spectrum}
\leanfit{cstarMatrix_cfc_real_exp_mul_le_bernstein_quadratic_of_spectrum_le}
\end{formalized}

\begin{theorem}[One-sample Bernstein log-CGF variance proxy]
Let \(X\) be a self-adjoint finite complex C\(^*\)-matrix-valued random
variable on a finite probability space.  Assume
\[
  \mathbb E X = 0,
  \qquad
  x\le R \quad\text{for all }x\in\sigma_{\mathbb R}(X(\omega)),
  \qquad
  \theta\ge0,\quad R>0 .
\]
Define
\[
  g(\theta,R)=\frac{e^{\theta R}-\theta R-1}{R^2},
  \qquad
  V=\mathbb E[X^2].
\]
Then
\[
  \log\mathbb E\exp(\theta X)
  \preceq
  g(\theta,R)\,V .
\]
Equivalently, before taking logarithms,
\[
  \mathbb E\exp_{\mathrm{cfc}}(\theta X)
  \preceq
  I+g(\theta,R)\mathbb E[X^2].
\]
\end{theorem}


\begin{formalized}
\lean{Analysis/CStarMatrixExpectation.lean}
\leanfit{FiniteProbability.expectationCStarMatrix_real_smul}
\lean{Analysis/LiebTrace.lean}
\leanfit{cstarMatrix_real_smul_isSelfAdjoint}
\leanfit{cstarMatrix_cfc_real_exp_mul_eq_normedSpace_exp_real_smul}
\leanfit{cstarMatrix_selfAdjoint_mul_self_nonneg}
\leanfit{cstarMatrix_one_add_le_normedSpace_exp_of_nonneg}
\leanfit{cstarMatrix_log_one_add_le_self_of_nonneg}
\leanfit{FiniteProbability.expectationCStarMatrix_nonneg_of_prob_pos}
\leanfit{FiniteProbability.expectationCStarMatrix_mono_of_prob_pos}
\leanfit{FiniteProbability.expectationCStarMatrix_cfc_real_exp_mul_le_bernstein_variance_proxy}
\leanfit{FiniteProbability.cstarMatrix_log_expectationCStarMatrix_cfc_real_exp_mul_le_bernstein_variance_proxy}
\leanfit{FiniteProbability.cstarMatrix_log_expectationCStarMatrix_normed_exp_real_smul_le_bernstein_variance_proxy}
\leanfit{FiniteProbability.expectationCStarMatrix_cfc_real_exp_mul_le_bernstein_variance_proxy_of_prob_pos}
\leanfit{FiniteProbability.cstarMatrix_log_expectationCStarMatrix_cfc_real_exp_mul_le_bernstein_variance_proxy_of_prob_pos}
\leanfit{FiniteProbability.cstarMatrix_log_expectationCStarMatrix_normed_exp_real_smul_le_bernstein_variance_proxy_of_prob_pos}
\end{formalized}

\begin{corollary}[Algorithm 1 truncated one-sample log-CGF]
Let \(\widehat A\) be the hard-thresholded matrix used in the
Drineas--Zouzias truncated element-wise sampler, let \(S\) be one
squared-magnitude sample from \(\widehat A\), and define the one-sample
residual increment
\[
  Z(S)=s^{-1}\widehat A-C(S).
\]
Let \(D(Z(S))\) be its self-adjoint dilation and assume \(s>0\), \(\tau>0\),
and \(\|\widehat A\|_F^2>0\).  Put
\[
  R_{\tau,s}(A)=
  \sqrt2\left(s^{-1}\|\widehat A\|_F+
    \frac{\|\widehat A\|_F^2}{s\tau}\right).
\]
Then, for every \(\theta\ge0\),
\[
  \log\mathbb E\exp(\theta D(Z(S)))
  \preceq
  \frac{e^{\theta R_{\tau,s}(A)}-\theta R_{\tau,s}(A)-1}
       {R_{\tau,s}(A)^2}
  \,\mathbb E[D(Z(S))^2].
\]
\end{corollary}


\begin{formalized}
\lean{Analysis/OperatorLog.lean}
\leanfit{cstarMatrix_spectrum_le_of_le_real_smul_one}
\lean{Algorithms/RandNLA/ElementwiseSpectral.lean}
\leanfit{sqMagSampleProbability_expectationCStarMatrix_rectSelfAdjointDilation_sampleResidualIncrement_eq_zero}
\leanfit{sqMagSampleProbability_finiteComplex_rectSelfAdjointDilation_sampleResidualIncrement_truncated_spectrum_le}
\leanfit{sqMagSampleProbability_cstarMatrix_log_expectationCStarMatrix_normed_exp_rectSelfAdjointDilation_truncated_sampleResidualIncrement_le}
\end{formalized}

\begin{lemma}[Algorithm 1 one-sample dilation variance proxy]
For the same one-sample residual increment \(Z(S)\),
\[
  \mathbb E\,D(Z(S))^2
  \preceq
  2mn\,\frac{\|\widehat A\|_F^2}{s^2}\,I_{m+n}.
\]
Here the expectation is under the squared-magnitude one-sample law for
\(\widehat A\), and the order is the complex C\(^*\)-matrix Loewner order
after embedding the real self-adjoint dilation.
\end{lemma}


\begin{formalized}
\lean{Analysis/CStarMatrixBridge.lean}
\leanfit{finiteComplexCStarMatrix_mul}
\lean{Analysis/MatrixAlgebra.lean}
\leanfit{finiteMatMul_self_symmetric_of_symmetric}
\leanfit{IsSymmetricFiniteMatrix.sum_smul}
\lean{Algorithms/RandNLA/ElementwiseSpectral.lean}
\leanfit{sqMagSampleProbability_expectationCStarMatrix_rectSelfAdjointDilation_sampleResidualIncrement_square_le}
\end{formalized}

\begin{lemma}[Source-sharp square dilation variance proxy]
For square \(n\times n\) input matrices, the Drineas--Zouzias route now has the
sharper one-step proxy
\[
  \mathbb E\,D(Z(S))^2
  \preceq
  n\,\frac{\|\widehat A\|_F^2}{s^2}\,I_{2n}.
\]
This theorem avoids both the earlier Frobenius-detour factor \(2n^2\) and the
extra \(\sqrt2\) support-radius loss.
\end{lemma}


\begin{formalized}
\lean{Algorithms/RandNLA/ElementwiseSpectral.lean}
\leanfit{sqMagProb_sum_vecNorm2Sq_rectMatMulVec_elementwiseSampleResidualIncrement_le_sharp}
\leanfit{sqMagProb_sum_vecNorm2Sq_transposeRectMatMulVec_elementwiseSampleResidualIncrement_le_sharp}
\leanfit{sqMagProb_sum_finiteQuadraticForm_rectSelfAdjointDilation_square_le_sharp_square}
\leanfit{sqMagProb_sum_rectSelfAdjointDilation_square_loewnerLe_scalar_id_sharp_square}
\leanfit{sqMagSampleProbability_expectationCStarMatrix_rectSelfAdjointDilation_sampleResidualIncrement_square_le_sharp_square}
\end{formalized}

\begin{lemma}[Source-aligned rectangular dilation variance proxy]
Let \(A\in\R^{m\times n}\), let \(s>0\), and let one Algorithm~1 sample be
drawn from the exact squared-magnitude law
\[
  p_{ij}=\frac{A_{ij}^2}{\Frob{A}^2}.
\]
For a sample \(S=(i,j)\), define the exact residual increment
\[
  Z(S)=s^{-1}A-\frac{A_{ij}}{s\,p_{ij}}e_i e_j^\mathsf T,
\]
with the usual zero-probability branch used by the formalized sampler.  Let
\[
  D(Z)=\begin{bmatrix}0&Z\\ Z^\mathsf T&0\end{bmatrix}
\]
be the self-adjoint dilation.  Then, in exact arithmetic,
\[
  \mathbb E\,D(Z(S))^2
  \preceq
  \max\{m,n\}\,\frac{\Frob{A}^2}{s^2}\,I_{m+n}.
\]
Equivalently, for every \(x\in\R^{m+n}\),
\[
  \mathbb E\,x^\mathsf T D(Z(S))^2x
  \le
  \max\{m,n\}\,\frac{\Frob{A}^2}{s^2}\,\|x\|_2^2 .
\]
The same result is also formalized in the C-star matrix order used by the
trace-MGF/log-CGF infrastructure.  This closes the unconditional rectangular
variance-proxy dependency for the literal Algorithm~1 law; it is not yet the
final matrix-Bernstein/Khintchine tail theorem for CACM equation~(2).
\end{lemma}

\begin{formalized}
\lean{Algorithms/RandNLA/ElementwiseSpectral.lean}
\leanfit{sqMagProb_sum_finiteQuadraticForm_rectSelfAdjointDilation_square_le_sharp_rect}
\leanfit{sqMagProb_sum_rectSelfAdjointDilation_square_loewnerLe_scalar_id_sharp_rect}
\leanfit{sqMagSampleProbability_expectationCStarMatrix_rectSelfAdjointDilation_sampleResidualIncrement_square_le_sharp_rect}
\end{formalized}

\begin{lemma}[Source-aligned rectangular summed variance proxy]
Let \(Z_1,\ldots,Z_s\) be the exact independent Algorithm~1 residual
increments under the squared-magnitude trace law above, and let \(D(Z_t)\) be
their self-adjoint dilations.  In exact arithmetic,
\[
  \mathbb E\sum_{t=1}^s D(Z_t)^2
  \preceq
  \max\{m,n\}\,\frac{\Frob{A}^2}{s}\,I_{m+n}.
\]
Equivalently, for every \(x\in\R^{m+n}\),
\[
  \mathbb E\sum_{t=1}^s x^\mathsf T D(Z_t)^2x
  \le
  \max\{m,n\}\,\frac{\Frob{A}^2}{s}\,\|x\|_2^2.
\]
The trace form is also proved:
\[
  \operatorname{tr}\!\left(\mathbb E\sum_{t=1}^s D(Z_t)^2\right)
  \le
  \max\{m,n\}\,\frac{\Frob{A}^2}{s}\,(m+n).
\]
This is the exact-law variance parameter expected by a rectangular
self-adjoint-dilation Bernstein route.  It still does not supply the missing
unconditional largest-eigenvalue tail for CACM equation~(2).
\end{lemma}

\begin{formalized}
\lean{Algorithms/RandNLA/ElementwiseSpectral.lean}
\leanfit{sqMagProb_sum_steps_finiteQuadraticForm_rectSelfAdjointDilation_square_le_sharp_rect}
\leanfit{sqMagProb_sum_steps_rectSelfAdjointDilation_square_loewnerLe_scalar_id_sharp_rect}
\leanfit{sqMagProb_sum_finiteTrace_rectSelfAdjointDilation_square_le_sharp_rect}
\leanfit{sqMagProb_sum_steps_finiteTrace_rectSelfAdjointDilation_square_le_sharp_rect}
\leanfit{sqMagTraceProbability_expectationReal_finiteQuadraticForm_rectSelfAdjointDilation_sampleResidualIncrement_square_le_sharp_rect}
\leanfit{sqMagTraceProbability_expectationReal_sum_finiteQuadraticForm_rectSelfAdjointDilation_sampleResidualIncrement_square_le_sharp_rect}
\leanfit{sqMagTraceProbability_sum_expectationReal_rectSelfAdjointDilation_square_loewnerLe_scalar_id_sharp_rect}
\end{formalized}

\begin{lemma}[Literal untruncated bounded-increment obstruction]
The exact literal squared-magnitude law does not give a uniform bounded
one-step support radius for all untruncated inputs.  For every proposed scalar
radius \(L\), define the \(1\times2\) exact matrix
\[
  A_L=\begin{bmatrix}1&(|L|+2)^{-1}\end{bmatrix}.
\]
Its squared-magnitude denominator is positive, and the small entry is sampled
with positive probability.  If \(S\) is that small-entry sample, then the exact
one-step residual increment with \(s=1\) satisfies
\[
  \left|Z(S)_{1,2}\right|=|L|+2>L,
\]
and consequently the repository's rectangular operator support predicate
\[
  \mathrm{rectOpNorm2Le}(Z(S),L)
\]
is false.  Moreover, under the exact one-step squared-magnitude product law
for \(A_L\), the event
\[
  \bigl\{\omega:
    \neg\,\mathrm{rectOpNorm2Le}
      (A_L-\widetilde A_L(\omega),L)\bigr\}
\]
has strictly positive probability.
\end{lemma}

\begin{remark}
This is an exact-arithmetic route-elimination theorem, not a floating-point
theorem and not the final CACM equation~(2) tail bound.  Sampling laws remain
exact mathematical inputs.  The result shows that the currently formalized
support-aware Bernstein API cannot close the source-uniform literal
untruncated Algorithm~1 theorem by inserting a uniform deterministic
bounded-increment radius.  The closed source-rate path must therefore keep its
truncation/no-small-entry hypothesis, or the literal theorem must use a
different concentration route.  The nonconditional support-radius theorem
below follows the latter principle by replacing the impossible uniform radius
with an explicit reciprocal-entry radius.
\end{remark}

\begin{formalized}
\lean{Algorithms/RandNLA/ElementwiseSpectral.lean}
\leanfit{algorithm1SmallEntrySupportMatrix}
\leanfit{sqMagProbDen_algorithm1SmallEntrySupportMatrix_pos}
\leanfit{sqMagProb_algorithm1SmallEntrySupportMatrix_small_pos}
\leanfit{algorithm1SmallEntrySupportMatrix_residual_increment_abs_eq}
\leanfit{algorithm1SmallEntrySupportMatrix_residual_increment_not_rectOpNorm2Le}
\leanfit{exists_sqMagPositive_sampleResidualIncrement_entry_abs_gt}
\leanfit{exists_sqMagPositive_sampleResidualIncrement_not_rectOpNorm2Le}
\leanfit{sqMagTraceProbability_eventProb_not_algorithm1ExactSpectralEvent_smallEntry_pos}
\end{formalized}

\begin{corollary}[Truncated trace-MGF scalarization]
Let
\[
  L_{\tau,s}(A)
  =
  \sqrt2\left(s^{-1}\|\widehat A\|_F+
    \frac{\|\widehat A\|_F^2}{s\tau}\right),
  \qquad
  V_{\tau,s}(A)=
  2mn\,\frac{\|\widehat A\|_F^2}{s^2}.
\]
For every \(\theta\ge0\), define
\[
  g(\theta,L)=\frac{\exp(\theta L)-\theta L-1}{L^2}.
\]
Then the iid trace-MGF expression for the positive dilation increments
\(\theta D(Z_t)\) satisfies
\[
  \operatorname{tr}\exp\!\left(\sum_{t=1}^s K_t^+\right)
  \le
  (m+n)\exp\!\left(s\,g(\theta,L_{\tau,s}(A))V_{\tau,s}(A)\right),
\]
where \(K_t^+=\log\mathbb E\exp(\theta D(Z_t))\).  The same bound holds for
the negative increments \(-\theta D(Z_t)\), with
\(K_t^-=\log\mathbb E\exp(-\theta D(Z_t))\).
\end{corollary}


\begin{formalized}
\lean{Analysis/CStarMatrixTrace.lean}
\leanfit{cstarMatrix_normedSpace_exp_le_real_smul_one_of_le_real_smul_one}
\leanfit{cstarMatrixTrace_normedSpace_exp_re_le_card_mul_exp_of_le_real_smul_one}
\lean{Algorithms/RandNLA/ElementwiseTraceMGF.lean}
\leanfit{sqMagTraceProbabilityFiniteRealTraceMGFLogBound_zero_le_card_mul_exp_of_log_le_real_smul_one}
\lean{Algorithms/RandNLA/ElementwiseSpectral.lean}
\leanfit{sqMagSampleProbability_neg_finiteComplex_rectSelfAdjointDilation_sampleResidualIncrement_truncated_spectrum_le}
\leanfit{sqMagTraceProbabilityFiniteRealTraceMGFLogBound_rectSelfAdjointDilation_truncated_sampleResidualIncrement_le}
\leanfit{sqMagTraceProbabilityFiniteRealTraceMGFLogBound_neg_rectSelfAdjointDilation_truncated_sampleResidualIncrement_le}
\end{formalized}

\begin{corollary}[Source-sharp rectangular truncated trace-MGF skeleton]
For rectangular \(m\times n\) matrices, the positive and negative
truncated trace-MGF scalarizations are also formalized with the sharper
rectangular variance proxy
\[
  V_{\tau,s}^{\rm rect}(A)
  =
  \max\{m,n\}\,
  \frac{\|\widehat A\|_F^2}{s^2},
  \qquad
  \widehat A=\operatorname{trunc}_\tau(A).
\]
The support radius remains the generic truncated one
\[
  L_{\tau,s}^{\rm rect}(A)
  =
  \sqrt2\left(
    s^{-1}\|\widehat A\|_F+
    \frac{\|\widehat A\|_F^2}{s\tau}
  \right).
\]
Thus the two-sided scaled dilation eigenvalue event has an explicit
\(1-\delta\) form with
\[
  B_{\rm rect}=(m+n)\,
  \exp\!\left(s\,g(\theta,L_{\tau,s}^{\rm rect}(A))
    V_{\tau,s}^{\rm rect}(A)\right),
  \qquad
  T=\log(2B_{\rm rect}/\delta).
\]
This is exact-arithmetic, exact-probability matrix-concentration
infrastructure.  It is consumed below by the rectangular source-budget and
floating-point gamma-transfer corollaries; it is still distinct from the
literal untruncated CACM law when tiny nonzero entries are not hard-thresholded.
\end{corollary}

\begin{formalized}
\lean{Algorithms/RandNLA/ElementwiseSpectral.lean}
\leanfit{sqMagTraceProbabilityFiniteRealTraceMGFLogBound_rectSelfAdjointDilation_truncated_sampleResidualIncrement_le_sharp_rect}
\leanfit{sqMagTraceProbabilityFiniteRealTraceMGFLogBound_neg_rectSelfAdjointDilation_truncated_sampleResidualIncrement_le_sharp_rect}
\leanfit{sqMagTraceProbability_eventProb_forall_abs_finiteHermitianEigenvalue_scaled_rectSelfAdjointDilation_truncatedTraceResidual_lt_ge_exp_sharp_rect}
\leanfit{sqMagTraceProbability_eventProb_forall_abs_finiteHermitianEigenvalue_scaled_rectSelfAdjointDilation_truncatedTraceResidual_lt_ge_one_sub_delta_sharp_rect}
\end{formalized}

\begin{corollary}[Source-sharp square trace-MGF scalarization]
For square \(n\times n\) matrices, the same positive and negative
trace-MGF scalarizations are now formalized with
\[
  L_{\tau,s}^{\sharp}(A)
  =
  s^{-1}\|\widehat A\|_F+
  \frac{\|\widehat A\|_F^2}{s\tau},
  \qquad
  V_{\tau,s}^{\sharp}(A)=
  n\,\frac{\|\widehat A\|_F^2}{s^2}.
\]
Consequently the two-sided truncated dilation eigenvalue event has an
explicit \(1-\delta\) form with
\[
  B_{\sharp}=2n\,
  \exp\!\left(s\,g(\theta,L_{\tau,s}^{\sharp}(A))
    V_{\tau,s}^{\sharp}(A)\right),
  \qquad
  T=\log(2B_{\sharp}/\delta).
\]
\end{corollary}

\begin{formalized}
\lean{Algorithms/RandNLA/ElementwiseSpectral.lean}
\leanfit{sqMagTraceProbabilityFiniteRealTraceMGFLogBound_rectSelfAdjointDilation_truncated_sampleResidualIncrement_le_sharp_square}
\leanfit{sqMagTraceProbabilityFiniteRealTraceMGFLogBound_neg_rectSelfAdjointDilation_truncated_sampleResidualIncrement_le_sharp_square}
\leanfit{sqMagTraceProbability_eventProb_forall_abs_finiteHermitianEigenvalue_scaled_rectSelfAdjointDilation_truncatedTraceResidual_lt_ge_exp_sharp_square}
\leanfit{sqMagTraceProbability_eventProb_forall_abs_finiteHermitianEigenvalue_scaled_rectSelfAdjointDilation_truncatedTraceResidual_lt_ge_one_sub_delta_sharp_square}
\end{formalized}

\begin{corollary}[Parameterized two-sided Bernstein tail skeleton]
For \(T\in\mathbb R\) and \(\theta\ge0\),
\[
  \mathbb P\!\left[
    \lambda_{\max}\bigl(\theta D(A-\widetilde A)\bigr)\ge T
  \right]
  \le
  e^{-T}(m+n)\exp\!\left(s\,g(\theta,L_{\tau,s}(A))V_{\tau,s}(A)\right).
\]
Moreover,
\[
  \mathbb P\!\left[
    \forall r,\,
    \left|\lambda_r\bigl(\theta D(A-\widetilde A)\bigr)\right|<T
  \right]
  \ge
  1-2e^{-T}(m+n)
    \exp\!\left(s\,g(\theta,L_{\tau,s}(A))V_{\tau,s}(A)\right).
\]
\end{corollary}


\begin{formalized}
\lean{Algorithms/RandNLA/ElementwiseSpectral.lean}
\leanfit{sqMagTraceProbability_eventProb_exists_finiteHermitianEigenvalue_scaled_rectSelfAdjointDilation_truncatedTraceResidual_ge_le_exp}
\leanfit{sqMagTraceProbability_eventProb_forall_abs_finiteHermitianEigenvalue_scaled_rectSelfAdjointDilation_truncatedTraceResidual_lt_ge_exp}
\end{formalized}

\begin{corollary}[Explicit \(1-\delta\) scaled dilation tail]
Let
\[
  B_{\theta,\tau,s}(A)
  =
  (m+n)\exp\!\left(s\,g(\theta,L_{\tau,s}(A))V_{\tau,s}(A)\right).
\]
Assume \(B_{\theta,\tau,s}(A)>0\) and \(\delta>0\).  With the choice
\[
  T=\log\!\left(\frac{2B_{\theta,\tau,s}(A)}{\delta}\right),
\]
the formalization proves
\[
  \mathbb P\!\left[
    \forall r,\,
    \left|\lambda_r\bigl(\theta D(A-\widetilde A)\bigr)\right|
      <
    \log\!\left(\frac{2B_{\theta,\tau,s}(A)}{\delta}\right)
  \right]
  \ge 1-\delta .
\]
This is still a scaled dilation-eigenvalue statement.  It removes the free tail
parameter \(T\), but it does not yet optimize \(\theta\) or convert the result
to the final CACM equation (2) rectangular spectral-norm constants.
\end{corollary}


\begin{formalized}
\lean{Analysis/MatrixConcentration.lean}
\leanfit{real_exp_neg_log_two_mul_div_mul_self_add}
\lean{Algorithms/RandNLA/ElementwiseSpectral.lean}
\leanfit{sqMagTraceProbability_eventProb_forall_abs_finiteHermitianEigenvalue_scaled_rectSelfAdjointDilation_truncatedTraceResidual_lt_ge_one_sub_delta}
\end{formalized}

\begin{corollary}[Explicit \(1-\delta\) spectral event with unoptimized radius]
Assume additionally that \(\theta>0\) and \(0<\delta\le 1\).  Then the
formalization proves the rectangular spectral event
\[
  \mathbb P\!\left[
    \left\|\widehat A-\widetilde A\right\|_2
      \le
    \frac{\log\!\left(2B_{\theta,\tau,s}(A)/\delta\right)}{\theta}
  \right]\ge 1-\delta .
\]
Here \(\widehat A\) is the hard-thresholded matrix used by the truncated
Algorithm~1 route, and \(\widetilde A\) is the exact sampled sketch from
that truncated distribution.  This is now a genuine spectral-norm event, not
only a scaled eigenvalue event.  Its radius is still unoptimized; the CACM
equation~(2) constant form requires choosing \(\theta\) and simplifying the
resulting expression.
\end{corollary}


\begin{formalized}
\lean{Analysis/MatrixAlgebra.lean}
\leanfit{finiteLoewnerLe_of_smul_left_le_smul_id}
\lean{Analysis/MatrixSpectral.lean}
\leanfit{finiteLoewnerLe_smul_id_of_finiteHermitianEigenvalues_le}
\lean{Algorithms/RandNLA/ElementwiseSpectral.lean}
\leanfit{algorithm1ScaledDilationAbsEigenvalueEvent}
\leanfit{algorithm1ScaledDilationAbsEigenvalueEvent_subset_exactSpectralEvent}
\leanfit{sqMagTraceProbability_eventProb_algorithm1ExactSpectralEvent_truncatedTraceResidual_ge_one_sub_delta_scaled_radius}
\end{formalized}

\begin{corollary}[Bennett-optimized \(1-\delta\) spectral event]
Let
\[
  W_{\tau,s}(A)=s\,V_{\tau,s}(A),\qquad
  q=\log\!\left(\frac{2(m+n)}{\delta}\right),
\]
and fix a desired radius \(r>0\).  Assume \(L_{\tau,s}(A)>0\),
\(W_{\tau,s}(A)>0\), and the Bennett budget
\[
  q\le
  \frac{W_{\tau,s}(A)}{L_{\tau,s}(A)^2}
  \left[
    \left(1+\frac{L_{\tau,s}(A)r}{W_{\tau,s}(A)}\right)
    \log\!\left(1+\frac{L_{\tau,s}(A)r}{W_{\tau,s}(A)}\right)
    -
    \frac{L_{\tau,s}(A)r}{W_{\tau,s}(A)}
  \right].
\]
With
\[
  \theta=\frac{\log(1+L_{\tau,s}(A)r/W_{\tau,s}(A))}
               {L_{\tau,s}(A)},
\]
the formalization proves
\[
  \mathbb P\!\left[
    \left\|\widehat A-\widetilde A\right\|_2 \le r
  \right]\ge 1-\delta .
\]
In Lean, the positivity of \(L_{\tau,s}(A)\) and \(W_{\tau,s}(A)\) is derived
from \(s>0\), \(\tau>0\), a positive squared-magnitude denominator, and the
domain assumption \(mn>0\).  The theorem is still the truncated exact
Algorithm~1 route; it does not yet simplify the Bennett budget to the final
Drineas--Zouzias/CACM sample-complexity constants or compose with
floating-point perturbation.
\end{corollary}


\begin{formalized}
\lean{Analysis/MatrixConcentration.lean}
\leanfit{real_bernstein_exact_radius_le_of_log_le}
\lean{Algorithms/RandNLA/ElementwiseSpectral.lean}
\leanfit{algorithm1ExactSpectralEvent_mono}
\leanfit{sqMagTraceProbability_eventProb_algorithm1ExactSpectralEvent_truncatedTraceResidual_ge_one_sub_delta_bennett_radius}
\end{formalized}

\begin{corollary}[Source-sharp Bennett spectral event]
In the square \(n\times n\) source-sharp route, set
\[
  L_{\tau,s}^{\sharp}(A)
  =
  s^{-1}\|\widehat A\|_F+
  \frac{\|\widehat A\|_F^2}{s\tau},
  \qquad
  V_{\tau,s}^{\sharp}(A)
  =
  n\,\frac{\|\widehat A\|_F^2}{s^2},
  \qquad
  W_{\tau,s}^{\sharp}(A)=sV_{\tau,s}^{\sharp}(A).
\]
Let
\[
  q_{\sharp}=\log\!\left(\frac{2(2n)}{\delta}\right),
\]
and fix \(r>0\).  If
\[
  q_{\sharp}\le
  \frac{W_{\tau,s}^{\sharp}(A)}{(L_{\tau,s}^{\sharp}(A))^2}
  \left[
    \left(1+\frac{L_{\tau,s}^{\sharp}(A)r}
                    {W_{\tau,s}^{\sharp}(A)}\right)
    \log\!\left(1+\frac{L_{\tau,s}^{\sharp}(A)r}
                    {W_{\tau,s}^{\sharp}(A)}\right)
    -
    \frac{L_{\tau,s}^{\sharp}(A)r}{W_{\tau,s}^{\sharp}(A)}
  \right],
\]
then the exact truncated Algorithm~1 sketch satisfies
\[
  \mathbb P\!\left[
    \left\|\widehat A-\widetilde A\right\|_2\le r
  \right]\ge 1-\delta .
\]
This is now the source-aligned Bennett step: it uses
\(V_{\tau,s}^{\sharp}(A)\) and the no-\(\sqrt2\) support radius.  It still
leaves open the algebraic simplification of this budget to the final
Drineas--Zouzias/CACM equation~(2) sample-complexity constants, followed by
truncation and floating-point transfer at those constants.
\end{corollary}


\begin{formalized}
\lean{Algorithms/RandNLA/ElementwiseSpectral.lean}
\leanfit{sqMagTraceProbability_eventProb_algorithm1ExactSpectralEvent_truncatedTraceResidual_ge_one_sub_delta_scaled_radius_sharp_square}
\leanfit{sqMagTraceProbability_eventProb_algorithm1ExactSpectralEvent_truncatedTraceResidual_ge_one_sub_delta_bennett_radius_sharp_square}
\end{formalized}

\begin{corollary}[Conservative source-sharp denominator form]
With the notation of the previous corollary, suppose instead of the Bennett
transform budget that
\[
  q_{\sharp}\le
  \frac{r^2}{2W_{\tau,s}^{\sharp}(A)+L_{\tau,s}^{\sharp}(A)r}.
\]
Then
\[
  \mathbb P\!\left[
    \left\|\widehat A-\widetilde A\right\|_2\le r
  \right]\ge 1-\delta .
\]
This statement is intentionally marked as a conservative fallback: the
denominator \(2W+Lr\) is weaker than the sharper Bernstein denominator
\(2W+\frac23 Lr\) used in the Drineas--Zouzias final constant simplification.
It is nevertheless a closed local theorem, because the scalar inequality
\[
  \frac{x^2}{2+x}\le (1+x)\log(1+x)-x,\qquad x\ge 0,
\]
is formalized and then composed with the source-sharp Bennett spectral event.
\end{corollary}


\begin{formalized}
\lean{Analysis/MatrixConcentration.lean}
\leanfit{real_bennett_transform_lower_bound_two_add}
\leanfit{real_bennett_budget_of_quadratic_denominator_two_add}
\lean{Algorithms/RandNLA/ElementwiseSpectral.lean}
\leanfit{sqMagTraceProbability_eventProb_algorithm1ExactSpectralEvent_truncatedTraceResidual_ge_one_sub_delta_bernstein_denominator_sharp_square}
\end{formalized}

\begin{corollary}[Source sample budget and truncation transfer]
Let \(A\in\mathbb R^{n\times n}\), \(n>0\), \(s>0\),
\(\varepsilon>0\), \(0<\delta\le1\), and set
\[
  \tau=\frac{\varepsilon}{2n},\qquad
  \widehat A=\operatorname{trunc}_{\tau}(A).
\]
Assume \(\|\widehat A\|_F>0\) and
\[
  14\,n\,\|A\|_F^2
  \log\!\left(\frac{2(2n)}{\delta}\right)
  \le s\varepsilon^2 .
\]
Then the exact truncated Algorithm~1 sampler satisfies
\[
  \mathbb P\!\left[
    \left\|\widehat A-\widetilde A\right\|_2
    \le \frac{\varepsilon}{2}
  \right]\ge 1-\delta .
\]
Moreover, after adding the deterministic truncation error,
\[
  \mathbb P\!\left[
    \left\|A-\widetilde A\right\|_2
    \le \varepsilon
  \right]\ge 1-\delta .
\]
\end{corollary}


\begin{formalized}
\lean{Analysis/MatrixConcentration.lean}
\leanfit{real_bennett_transform_lower_bound_two_add_two_thirds}
\leanfit{real_bennett_budget_of_quadratic_denominator_two_add_two_thirds}
\lean{Algorithms/RandNLA/ElementwiseSpectral.lean}
\leanfit{elementwiseTruncate_tau_le_frobNormRect_of_sqMagProbDen_pos}
\leanfit{sqMagTraceProbability_eventProb_algorithm1ExactSpectralEvent_truncatedTraceResidual_ge_one_sub_delta_source_sample_budget_sharp_square}
\leanfit{sqMagTraceProbability_eventProb_algorithm1ExactTruncatedSpectralEvent_ge_one_sub_delta_source_sample_budget_sharp_square}
\end{formalized}

\begin{corollary}[Floating-point transfer with an explicit perturbation budget]
Under the hypotheses of the previous corollary, let \(B\in\mathbb R^{n\times n}\)
be an entrywise nonnegative perturbation budget.  If for every trace sample
\(\omega\) and every entry \((i,j)\),
\[
  \left|
    \operatorname{fl}(\widetilde A_{\tau})(\omega)_{ij}
    -\widetilde A_{\tau}(\omega)_{ij}
  \right|\le B_{ij},
\]
then the rounded truncated Algorithm~1 sketch satisfies
\[
  \mathbb P\!\left[
    \left\|A-\operatorname{fl}(\widetilde A_{\tau})\right\|_2
    \le \varepsilon+\|B\|_F
  \right]\ge 1-\delta .
\]
This is a genuine FP transfer theorem, but it still records the entrywise
perturbation budget as an explicit hypothesis.  The next corollary discharges
that budget for the canonical squared-magnitude trace law by using the
probability-one positive-support event.
\end{corollary}

\begin{formalized}
\lean{Algorithms/RandNLA/ElementwiseSpectral.lean}
\leanfit{sqMagTraceProbability_eventProb_algorithm1FlTruncatedSpectralEvent_ge_one_sub_delta_source_sample_budget_sharp_square}
\end{formalized}

\begin{corollary}[Floating-point source budget with derived gamma budget]
Under the source sample-budget hypotheses above, assume additionally
\(\gamma_s\) and \(\gamma_{s+1}\) are valid for the floating-point model.  Let
\[
  B_{ij}
  =
  \operatorname{budget}_{ij}
  =
  \gamma_s |0|
  + s\left|\frac{\|\widehat A\|_F^2}{s\,\widehat A_{ij}}\right|\gamma_{s+1},
\]
where Lean writes this as
\(\operatorname{sqMagTraceErrorBudget}(\mathrm{fp},s,s,\widehat A,0)_{ij}\).
Then
\[
  \mathbb P\!\left[
    \left\|A-\operatorname{fl}(\widetilde A_{\tau})\right\|_2
    \le \varepsilon+\|B\|_F
  \right]\ge 1-\delta .
\]
The proof does not assume the entrywise perturbation event.  It intersects the
exact source-budget event with the event that every sampled entry of
\(\widehat A\) has positive squared-magnitude probability.  That support event
has probability one under the canonical product law.  On the support event, all
actually used denominators are nonzero, each hit count is at most \(s\), and
the local Algorithm~1 gamma stability theorem supplies the displayed budget.
Traces that hit zero-probability entries are irrelevant because they have zero
mass, and the floating-point model intentionally makes no claim about division
by a zero denominator on those impossible traces.
\end{corollary}

\begin{formalized}
\lean{Algorithms/RandNLA/ElementwiseSampling.lean}
\leanfit{hitCount_le_steps}
\leanfit{hitCount_eq_zero_of_forall_not_hit}
\leanfit{fl_elementwiseTraceSketch_zero_init_eq_zero_of_forall_not_hit}
\leanfit{sqMagTraceErrorBudget_nonneg}
\lean{Algorithms/RandNLA/ElementwiseSpectral.lean}
\leanfit{fl_elementwiseTraceSketch_zero_init_sqMag_error_bound_of_positiveProb}
\leanfit{sqMagTraceProbability_eventProb_algorithm1FlTruncatedSpectralEvent_ge_one_sub_delta_source_sample_budget_gamma_square}
\end{formalized}

\begin{corollary}[Intermediate Algorithm 1 floating-point radius expansion]
Assume the source sample-budget hypotheses of the previous corollaries:
\[
  0<n,\qquad 0<s,\qquad 0<\varepsilon,\qquad 0<\delta\le 1,
\]
\[
  14n\Frob{A}^2
  \log\!\left(\frac{2(2n)}{\delta}\right)
  \le s\varepsilon^2.
\]
Let
\[
  \tau=\frac{\varepsilon}{2n},
  \qquad
  \widehat A=\operatorname{trunc}_{\tau}(A),
\]
and assume the squared-magnitude denominator of \(\widehat A\) is positive.
Assume also that the floating-point gamma factors \(\gamma_s\) and
\(\gamma_{s+1}\) are valid.  Under the canonical product law that samples
entries from \(\widehat A_{ij}^2/\Frob{\widehat A}^2\),
\[
  \mathbb P\!\left[
    \left\|A-\operatorname{fl}(\widetilde A_{\tau})\right\|_2
    \le
    \varepsilon+
    n\,\frac{\Frob{\widehat A}^2}{\tau}\,\gamma_{s+1}
  \right]\ge 1-\delta .
\]
Since \(\Frob{\widehat A}\le\Frob{A}\) and
\(\tau=\varepsilon/(2n)\), the fully source-level version is
\[
  \mathbb P\!\left[
    \left\|A-\operatorname{fl}(\widetilde A_{\tau})\right\|_2
    \le
    \varepsilon+
    \frac{2n^2\Frob{A}^2}{\varepsilon}\,\gamma_{s+1}
  \right]\ge 1-\delta .
\]
Here
\[
  \gamma_{s+1}
  =
  \frac{(s+1)u}{1-(s+1)u},
\]
with the usual validity condition \((s+1)u<1\).  Thus the final FP error term
is displayed in the dimension \(n\), the sample count \(s\), the unit roundoff
\(u\), \(\Frob{A}\), and \(\varepsilon\), rather than hidden inside a budget
matrix.  The same statement is restated as the final Algorithm~1 equation~(2)
corollary at the end of the Algorithm~1 section.
\end{corollary}


\begin{formalized}
\lean{Algorithms/RandNLA/ElementwiseSpectral.lean}
\leanfit{sqMagTraceErrorBudget_zero_init_truncated_le_const}
\leanfit{frobNormRect_sqMagTraceErrorBudget_zero_init_truncated_le_const_square}
\leanfit{sqMagTraceProbability_eventProb_algorithm1FlTruncatedSpectralRadius_ge_one_sub_delta_source_sample_budget_gamma_square}
\leanfit{sqMagTraceProbability_eventProb_algorithm1FlTruncatedSpectralRadius_ge_one_sub_delta_source_sample_budget_explicit_gamma_square}
\end{formalized}

\begin{corollary}[Rectangular source-budget and floating-point gamma radius]
Let \(A\in\mathbb R^{m\times n}\), set
\[
  R=\sqrt{mn},\qquad
  M=\max\{m,n\},\qquad
  C_{\rm rect}=2M+\frac{4\sqrt2}{3}R,
\]
and choose the rectangular truncation threshold
\[
  \tau=\frac{\varepsilon}{2R},
  \qquad
  \widehat A=\operatorname{trunc}_{\tau}(A).
\]
Assume \(mn>0\), \(s>0\), \(0<\delta\le 1\), \(\varepsilon>0\), and
\(\|\widehat A\|_F^2>0\).  If
\[
  4\,C_{\rm rect}\,\|A\|_F^2
  \log\!\left(\frac{2(m+n)}{\delta}\right)
  \le s\varepsilon^2,
\]
then under the exact product law that samples entries from
\(\widehat A_{ij}^2/\|\widehat A\|_F^2\),
\[
  \mathbb P\!\left[
    \left\|A-\widetilde A_{\tau}\right\|_2\le \varepsilon
  \right]\ge 1-\delta .
\]
If \(\gamma_s\) and \(\gamma_{s+1}\) are valid for the floating-point model,
then the rounded sampled-update computation satisfies
\[
  \mathbb P\!\left[
    \left\|A-\operatorname{fl}(\widetilde A_{\tau})\right\|_2
    \le
    \varepsilon+
    R\,\frac{\|\widehat A\|_F^2}{\tau}\,\gamma_{s+1}
  \right]\ge 1-\delta .
\]
Using \(\|\widehat A\|_F\le \|A\|_F\) and
\(\tau=\varepsilon/(2R)\), the fully source-level displayed radius is
\[
  \varepsilon+
  \frac{2mn\,\|A\|_F^2}{\varepsilon}\,\gamma_{s+1}.
\]
Equivalently, since
\(\gamma_{s+1}=(s+1)u/(1-(s+1)u)\), the FP additive term is
\[
  \Theta\!\left(
    \frac{mn\,\|A\|_F^2}{\varepsilon}\,s u
  \right)
\]
in the regime of fixed \(m,n,A,\varepsilon\) and small \((s+1)u\).
The sample condition is
\[
  s
  =
  \Omega\!\left(
    \frac{(M+R)\|A\|_F^2}{\varepsilon^2}
    \log\!\frac{m+n}{\delta}
  \right)
  =
  \Theta\!\left(
    \frac{\max\{m,n\}\|A\|_F^2}{\varepsilon^2}
    \log\!\frac{m+n}{\delta}
  \right),
\]
because \(R=\sqrt{mn}\le M\).  The result is nonvacuous: for fixed dimensions,
source parameters, and a fixed sample count satisfying the displayed budget,
the floating-point radius tends to \(\varepsilon\) as \(u\to0\).
\end{corollary}

\begin{formalized}
\lean{Algorithms/RandNLA/ElementwiseSpectral.lean}
\leanfit{elementwiseTruncate_rect_error_frobNormRect_le_half}
\leanfit{sqMagTraceProbability_eventProb_algorithm1ExactSpectralEvent_truncatedTraceResidual_ge_one_sub_delta_source_sample_budget_sharp_rect}
\leanfit{sqMagTraceProbability_eventProb_algorithm1ExactTruncatedSpectralEvent_ge_one_sub_delta_source_sample_budget_sharp_rect}
\leanfit{sqMagTraceProbability_eventProb_algorithm1FlTruncatedSpectralEvent_ge_one_sub_delta_source_sample_budget_gamma_rect}
\leanfit{frobNormRect_sqMagTraceErrorBudget_zero_init_truncated_le_const_rect}
\leanfit{sqMagTraceProbability_eventProb_algorithm1FlTruncatedSpectralRadius_ge_one_sub_delta_source_sample_budget_gamma_rect}
\leanfit{sqMagTraceProbability_eventProb_algorithm1FlTruncatedSpectralRadius_ge_one_sub_delta_source_sample_budget_explicit_gamma_rect}
\end{formalized}

\begin{lemma}[Strictly positive exponential--logarithm inverse]
Let \(A\) be a strictly positive complex \(d\times d\) C\(^*\)-matrix.  Then
\[
  \exp(\log A)=A .
\]
The same identity is formalized both for the normed-algebra exponential and
for the real and complex continuous-functional-calculus exponentials:
\[
  \exp_{\mathrm{cfc},\mathbb C}(\log A)=A,
  \qquad
  \exp_{\mathrm{cfc},\mathbb R}(\log A)=A .
\]
\end{lemma}


\begin{formalized}
\lean{Analysis/OperatorLog.lean}
\leanfit{cstarMatrix_normedRingExp_nonnegSpectrumClass}
\leanfit{cstarMatrix_normedSpace_exp_log_of_isStrictlyPositive}
\leanfit{cstarMatrix_cfc_complex_exp_log_of_isStrictlyPositive}
\leanfit{cstarMatrix_cfc_real_exp_log_of_isStrictlyPositive}
\end{formalized}

\begin{lemma}[Finite-real to complex C\(^*\)-matrix algebraic and order bridge]
Let \(M,N\in\mathbb R^{d\times d}\).  Define
\[
  M_{\mathbb C}=(M_{ij}: \mathbb C)_{i,j}.
\]
Then entrywise embedding into complex C\(^*\)-matrices preserves subtraction,
the identity, and scalar identities:
\[
  (M-N)_{\mathbb C}=M_{\mathbb C}-N_{\mathbb C},\qquad
  I_{\mathbb C}=I,\qquad
  (aI)_{\mathbb C}=aI .
\]
Moreover, if \(M\) is symmetric, then \(M_{\mathbb C}\) is self-adjoint.  If
\(M\succeq 0\) in the local finite quadratic-form sense, then
\(M_{\mathbb C}\succeq 0\) in complex C\(^*\)-matrix order.  If \(M,N\) are
symmetric and \(M\preceq N\) locally, then
\[
  M_{\mathbb C}\preceq N_{\mathbb C}
\]
in the same complex C\(^*\)-matrix order.
\end{lemma}


\begin{formalized}
\lean{Analysis/CStarMatrixBridge.lean}
\leanfit{finiteComplexCStarMatrix}
\leanfit{finiteComplexCStarMatrix_isSelfAdjoint_of_symmetric}
\leanfit{finiteComplexCStarMatrix_sub}
\leanfit{finiteComplexCStarMatrix_finiteIdMatrix}
\leanfit{finiteComplexCStarMatrix_smul_finiteIdMatrix}
\leanfit{finiteComplexCStarMatrix_nonneg_of_finitePSD}
\leanfit{finiteComplexCStarMatrix_le_of_finiteLoewnerLe}
\end{formalized}

\begin{lemma}[Regularized finite-real operator-log bridge]
Let \(M,N\in\mathbb R^{d\times d}\) be symmetric, assume \(M\succeq 0\),
and assume \(M\preceq N\) in local finite Loewner order.  For every
\(\eta>0\),
\[
  M_{\mathbb C}+\eta I \quad\text{is strictly positive}
\]
as a complex C\(^*\)-matrix, and
\[
  M_{\mathbb C}+\eta I \preceq N_{\mathbb C}+\eta I .
\]
Consequently,
\[
  \log(M_{\mathbb C}+\eta I)
  \preceq
  \log(N_{\mathbb C}+\eta I).
\]
\end{lemma}


\begin{formalized}
\lean{Analysis/CStarMatrixBridge.lean}
\leanfit{cstarMatrix_pos_real_smul_one_isStrictlyPositive}
\leanfit{finiteComplexCStarMatrix_add_pos_smul_one_isStrictlyPositive_of_finitePSD}
\leanfit{finiteComplexCStarMatrix_add_smul_one_le_of_finiteLoewnerLe}
\lean{Analysis/OperatorLog.lean}
\leanfit{finiteComplexCStarMatrix_regularized_log_le_log_of_finiteLoewnerLe}
\end{formalized}

\begin{lemma}[Complex C\(^*\)-matrix trace bridge]
For a complex \(d\times d\) C\(^*\)-matrix \(X\), define
\[
  \operatorname{tr}_{\mathbb C}(X)=\sum_{i=1}^d X_{ii}.
\]
This trace is additive and homogeneous, and
\[
  \operatorname{tr}_{\mathbb C}(I)=d,\qquad
  \operatorname{tr}_{\mathbb C}(aI)=ad .
\]
It is also cyclic:
\[
  \operatorname{tr}_{\mathbb C}(XY)
  =
  \operatorname{tr}_{\mathbb C}(YX).
\]
The trace is positive and monotone in real part for the C\(^*\)-spectral
order:
\[
  0\preceq X
  \quad\Longrightarrow\quad
  0\le \Re\,\operatorname{tr}_{\mathbb C}(X),
\]
and
\[
  X\preceq Y
  \quad\Longrightarrow\quad
  \Re\,\operatorname{tr}_{\mathbb C}(X)
  \le
  \Re\,\operatorname{tr}_{\mathbb C}(Y).
\]
If \(M\in\mathbb R^{d\times d}\) is embedded as \(M_{\mathbb C}\), then
\[
  \operatorname{tr}_{\mathbb C}(M_{\mathbb C})
  =
  \operatorname{tr}(M).
\]
Consequently, if \(M\succeq0\) in local finite PSD order, then
\[
  0\le \Re\,\operatorname{tr}_{\mathbb C}(M_{\mathbb C}),
\]
and if \(M\preceq N\) in local finite Loewner order, then
\[
  \Re\,\operatorname{tr}_{\mathbb C}(M_{\mathbb C})
  \le
  \Re\,\operatorname{tr}_{\mathbb C}(N_{\mathbb C}).
\]
If \(M_{\mathbb C}=M_{\mathbb C}^*\), then
\[
  \Im\,\operatorname{tr}_{\mathbb C}(M_{\mathbb C})=0 .
\]
Finally, for scalar identities, the complex continuous functional calculus
normalizes the exponential trace:
\[
  \operatorname{tr}_{\mathbb C}\!\left(\exp_{\mathrm{cfc}}(aI)\right)
  =
  e^a d .
\]
\end{lemma}


\begin{formalized}
\lean{Analysis/CStarMatrixTrace.lean}
\leanfit{cstarMatrixTrace}
\leanfit{cstarMatrixTrace_neg}
\leanfit{cstarMatrixTrace_sub}
\leanfit{cstarMatrixTrace_mul_comm}
\leanfit{cstarMatrixTrace_star_mul_self_re_nonneg}
\leanfit{cstarMatrixTrace_re_nonneg_of_nonneg}
\leanfit{cstarMatrixTrace_re_mono}
\leanfit{cstarMatrixTrace_im_eq_zero_of_isSelfAdjoint}
\leanfit{cstarMatrixTrace_finiteComplexCStarMatrix}
\leanfit{cstarMatrixTrace_finiteComplexCStarMatrix_re_nonneg_of_finitePSD}
\leanfit{cstarMatrixTrace_finiteComplexCStarMatrix_re_mono_of_finiteLoewnerLe}
\leanfit{cstarMatrixTrace_cfc_exp_algebraMap}
\leanfit{cstarMatrixTrace_cfc_exp_real_smul_one}
\end{formalized}

\begin{lemma}[Complex C\(^*\)-matrix expectation, order, and Jensen bridge]
Let \((\Omega,\mathbb P)\) be a finite probability space.  For a
complex-valued random variable \(Z:\Omega\to\mathbb C\), define
\[
  \mathbb E_{\mathbb C} Z
  =
  \sum_{\omega\in\Omega} \mathbb P(\omega) Z(\omega).
\]
For a complex C\(^*\)-matrix-valued random variable
\(X:\Omega\to\mathbb C^{d\times d}\), define expectation entrywise:
\[
  (\mathbb E_{\mathbb C}^{\mathrm{mat}} X)_{ij}
  =
  \mathbb E_{\mathbb C}\bigl(X_{ij}\bigr).
\]
Then complex expectation is linear, agrees with the repository's real
expectation after coercion \(\mathbb R\hookrightarrow\mathbb C\), satisfies
\[
  \operatorname{Re}\bigl(\mathbb E_{\mathbb C}Z\bigr)
  =
  \mathbb E_{\mathbb R}\operatorname{Re}(Z),
\]
and
\[
  \operatorname{tr}_{\mathbb C}
    \bigl(\mathbb E_{\mathbb C}^{\mathrm{mat}} X\bigr)
  =
  \mathbb E_{\mathbb C}
    \bigl(\operatorname{tr}_{\mathbb C}(X)\bigr).
\]
Moreover, if \(M(\omega)\) is a repository-native finite real matrix, then
\[
  \mathbb E_{\mathbb C}^{\mathrm{mat}}\bigl(M(\omega)_{\mathbb C}\bigr)
  =
  \left(\mathbb E_{\mathbb R} M(\omega)\right)_{\mathbb C}.
\]
The matrix expectation has both complex and real weighted-sum forms
\[
  \mathbb E_{\mathbb C}^{\mathrm{mat}} X
  =
  \sum_{\omega\in\Omega}\mathbb P(\omega)X(\omega).
\]
Consequently, if \(X(\omega)\succeq 0\) for every \(\omega\), then
\[
  \mathbb E_{\mathbb C}^{\mathrm{mat}} X \succeq 0.
\]
If \(X(\omega)\succ 0\) for every \(\omega\), then
\[
  \mathbb E_{\mathbb C}^{\mathrm{mat}} X \succ 0.
\]
If \(X(\omega)\preceq Y(\omega)\) for every \(\omega\), then
\[
  \mathbb E_{\mathbb C}^{\mathrm{mat}} X
  \preceq
  \mathbb E_{\mathbb C}^{\mathrm{mat}} Y.
\]
Finally, if \(X(\omega)\succeq 0\) and \(\varepsilon>0\), then
\[
  \mathbb E_{\mathbb C}^{\mathrm{mat}} X+\varepsilon I
\]
is strictly positive.

In addition, if \(f:S\to\mathbb R\) is concave on a convex set
\(S\subseteq \mathbb C^{d\times d}\) and \(X(\omega)\in S\), then finite
Jensen gives
\[
  \mathbb E_{\mathbb R}\bigl[f(X)\bigr]
  \le
  f\!\left(\mathbb E_{\mathbb C}^{\mathrm{mat}}X\right).
\]
\end{lemma}


\begin{formalized}
\lean{Analysis/CStarMatrixExpectation.lean}
\leanfit{FiniteProbability.expectationComplex}
\leanfit{FiniteProbability.expectationComplex_sum}
\leanfit{FiniteProbability.expectationComplex_ofReal}
\leanfit{FiniteProbability.expectationComplex_re}
\leanfit{FiniteProbability.exists_prob_pos}
\leanfit{FiniteProbability.expectationReal_le_of_concaveOn}
\leanfit{FiniteProbability.expectationComplex_add}
\leanfit{FiniteProbability.expectationComplex_smul}
\leanfit{FiniteProbability.expectationCStarMatrix}
\leanfit{FiniteProbability.cstarMatrixTrace_expectationCStarMatrix}
\leanfit{FiniteProbability.expectationCStarMatrix_finiteComplexCStarMatrix}
\leanfit{FiniteProbability.expectationCStarMatrix_add}
\leanfit{FiniteProbability.expectationCStarMatrix_smul}
\leanfit{FiniteProbability.expectationCStarMatrix_eq_sum_smul}
\leanfit{FiniteProbability.expectationCStarMatrix_eq_sum_real_smul}
\leanfit{FiniteProbability.expectationReal_le_of_concaveOn_expectationCStarMatrix}
\leanfit{FiniteProbability.expectationReal_trace_cfc_exp_add_log_le_of_concaveOn}
\leanfit{FiniteProbability.expectationCStarMatrix_nonneg}
\leanfit{FiniteProbability.expectationCStarMatrix_isStrictlyPositive}
\leanfit{FiniteProbability.expectationCStarMatrix_mono}
\leanfit{FiniteProbability.expectationCStarMatrix_add_pos_smul_one_isStrictlyPositive}
\end{formalized}

\begin{lemma}[Strictly positive cone vocabulary for Lieb]
Let
\[
  \mathcal P_{++}
  =
  \{A\in\mathbb C^{d\times d}: A \text{ is strictly positive}\}
\]
in the complex C\(^*\)-matrix spectral order.  This cone is convex over real
scalars.  Positive real scalar multiples preserve strict positivity, and
nonnegative real scalar multiples preserve nonnegativity.  The local Lieb
trace functional is named as
\[
  \Phi_H(A)
  =
  \operatorname{Re}\operatorname{tr}
  \exp_{\mathrm{cfc}}\!\bigl(H+\log A\bigr).
\]
For self-adjoint \(H\), the matrix
\(\exp_{\mathrm{cfc}}(H+\log A)\) is nonnegative in the complex
C\(^*\)-spectral order, its trace is real-valued, and
\[
  \Phi_H(A)\ge0 .
\]
It also agrees with the standard normed-algebra exponential notation used in
Tropp's trace-MGF theorem:
\[
  \Phi_H(A)
  =
  \operatorname{Re}\operatorname{tr}
  \exp\!\bigl(H+\log A\bigr),
\]
and on the strictly positive cone it satisfies the normalization
\[
  \Phi_0(A)=\operatorname{Re}\operatorname{tr}(A).
\]
Consequently the \(H=0\) special case of the local Lieb concavity target is
proved: \(\Phi_0\) is affine, hence concave, on \(\mathcal P_{++}\).
The remaining Lieb theorem target is the assertion that \(\Phi_H\) is concave
on \(\mathcal P_{++}\) for fixed self-adjoint \(H\).
\end{lemma}


\begin{corollary}[Scalar normalized-kernel representation of \(x\log x\)]
For \(x>0\), let
\[
  g(x)=
  \begin{cases}
    1, & x=1,\\
    x\log x/(x-1), & x\ne 1.
  \end{cases}
\]
Then
\[
  (x-1)g(x)=x\log x
\]
and, using the already formalized unit-interval representation of \(g\),
\[
  (x-1)\int_0^1 \frac{x}{u+(1-u)x}\,du=x\log x .
\]
\end{corollary}

\begin{theorem}[Correct unit-interval kernel for \(x\log x\)]
For \(x>0\),
\[
  \int_0^1 \frac{x(x-1)}{u+(1-u)x}\,du=x\log x .
\]
Moreover, for \(0\le u<1\),
\[
  \frac{x(x-1)}{u+(1-u)x}
  =
  \frac{x}{1-u}
  -
  \frac{1}{(1-u)^2}
  +
  \frac{u}{(1-u)^3}
  \left(x+\frac{u}{1-u}\right)^{-1}.
\]
\end{theorem}


\begin{remark}
The auxiliary square-kernel identity below is still true and useful, but it is
not the kernel whose integral reconstructs \(x\log x\).  The reconstruction
kernel is the \(x(x-1)\) kernel in the theorem above.
\end{remark}

\begin{formalized}
\leanfit{real_xlog_eq_sub_one_mul_realNormalizedLogKernel}
\leanfit{real_xlog_eq_sub_one_mul_normalizedKernel_setIntegral}
\leanfit{real_xlog_eq_unit_interval_xlog_kernel_integral}
\leanfit{real_unit_interval_xlog_kernel_integrand_eq_affine_add_shifted_inv}
\leanfit{cstarMatrix_cfc_unit_interval_xlog_kernel_integrand_eq_affine_add_shifted_inv}
\end{formalized}

\begin{lemma}[Auxiliary square-kernel decomposition]
For \(x>0\) and \(0\le u<1\),
\[
  \frac{(x-1)^2}{u+(1-u)x}
  =
  \frac{x}{1-u}
  -
  \frac{2-u}{(1-u)^2}
  +
  \frac{1}{(1-u)^3}
  \left(x+\frac{u}{1-u}\right)^{-1}.
\]
\end{lemma}


\begin{formalized}
\leanfit{real_unit_interval_xlog_integrand_eq_affine_add_shifted_inv}
\leanfit{cstarMatrix_cfc_unit_interval_xlog_integrand_eq_affine_add_shifted_inv}
\end{formalized}

\begin{formalized}
\lean{Analysis/LiebTrace.lean}
\leanfit{strictPositiveCStarMatrixCone}
\leanfit{mem_strictPositiveCStarMatrixCone}
\leanfit{cstarMatrix_isStrictlyPositive_pos_real_smul}
\leanfit{cstarMatrix_nonneg_nonneg_real_smul}
\leanfit{cstarMatrix_isStrictlyPositive_pos_nonneg_real_smul_add}
\leanfit{strictPositiveCStarMatrixCone_convex}
\leanfit{cstarMatrix_log_isSelfAdjoint}
\leanfit{liebTraceArgument_isSelfAdjoint}
\leanfit{liebTraceArgument_isStarNormal}
\leanfit{liebTraceCfcExp_nonneg}
\leanfit{liebTraceCfcExp_isStrictlyPositive}
\leanfit{liebTraceFunctional_trace_im_eq_zero}
\leanfit{liebTraceFunctional}
\leanfit{liebTraceFunctional_nonneg}
\leanfit{liebTraceFunctional_eq_normedSpace_exp}
\leanfit{liebTraceFunctional_zero_eq_trace}
\leanfit{liebTraceConcavityTarget}
\leanfit{liebTraceConcavityTarget_zero}
\end{formalized}

\begin{corollary}[Conditional one-step Tropp--Jensen trace-MGF bound]
Assume the Lieb concavity statement
\[
  \Phi_H(A)=\operatorname{Re}\operatorname{tr}
    \exp\!\bigl(H+\log A\bigr)
  \quad\text{is concave on }\mathcal P_{++}
\]
for the fixed self-adjoint matrix \(H\).  Let \(X\) be a finite-valued
self-adjoint matrix random variable.  Then the local formalization proves
\[
  \mathbb E\,\operatorname{Re}\operatorname{tr}\exp(H+X)
  \le
  \operatorname{Re}\operatorname{tr}
  \exp\!\left(H+\log\mathbb E\,\exp X\right).
\]
\end{corollary}


\begin{formalized}
\lean{Analysis/LiebTrace.lean}
\leanfit{FiniteProbability.expectationReal_trace_normed_exp_add_le_of_liebTraceConcavityTarget}
\end{formalized}

\begin{definition}[Scalar and finite-vector relative entropy]
For positive scalars \(a,b\), define
\[
  d(a;b)=a(\log a-\log b)-(a-b).
\]
For finite real vectors \(x,y\) indexed by a finite type \(I\), define
\[
  D_{\mathrm{vec}}(x;y)=\sum_{i\in I} d(x_i;y_i).
\]
\end{definition}

\begin{lemma}[Scalar relative-entropy nonnegativity]
If \(a>0\) and \(b>0\), then
\[
  d(a;b)\ge 0.
\]
\end{lemma}

\begin{corollary}[Finite-vector relative-entropy nonnegativity]
If \(x_i>0\) and \(y_i>0\) for every coordinate \(i\), then
\[
  D_{\mathrm{vec}}(x;y)\ge 0.
\]
\end{corollary}


\begin{formalized}
\lean{Analysis/LiebTrace.lean}
\leanfit{realRelativeEntropy}
\leanfit{realRelativeEntropy_nonneg}
\leanfit{finiteRealRelativeEntropy}
\leanfit{finiteRealRelativeEntropy_nonneg}
\end{formalized}

\begin{theorem}[Finite log-sum inequality]
Let \(p_i>0\) and \(q_i>0\) for all \(i\in I\), with \(I\) finite and
nonempty.  Then
\[
  \left(\sum_i p_i\right)
  \log\!\left(\frac{\sum_i p_i}{\sum_i q_i}\right)
  \le
  \sum_i p_i\log\!\left(\frac{p_i}{q_i}\right).
\]
\end{theorem}

\begin{corollary}[Commutative joint convexity of relative entropy]
If \(x,y,a,b>0\), \(\alpha,\beta\ge0\), and \(\alpha+\beta=1\), then
\[
  d(\alpha x+\beta y;\alpha a+\beta b)
  \le
  \alpha d(x;a)+\beta d(y;b).
\]
Consequently, if all vector coordinates are positive, then
\[
  D_{\mathrm{vec}}(\alpha x+\beta y;\alpha a+\beta b)
  \le
  \alpha D_{\mathrm{vec}}(x;a)+
  \beta D_{\mathrm{vec}}(y;b).
\]
\end{corollary}


\begin{formalized}
\lean{Analysis/LiebTrace.lean}
\leanfit{finite_log_sum_inequality}
\leanfit{realRelativeEntropy_jointConvex_of_pos_weights}
\leanfit{realRelativeEntropy_jointConvex}
\leanfit{finiteRealRelativeEntropy_jointConvex}
\end{formalized}

\begin{definition}[C-star matrix relative entropy vocabulary]
For finite complex C\(^*\)-matrices \(A\) and \(B\), the current
Tropp-route vocabulary defines
\[
  D(A;B)
  =
  \operatorname{Re}\operatorname{tr}
  \left(A(\log A-\log B)-(A-B)\right).
\]
This is the finite-dimensional relative-entropy expression used by the
chosen route toward Lieb trace concavity.
\end{definition}

\begin{lemma}[Diagonal normalization of relative entropy]
For every finite complex C\(^*\)-matrix \(A\),
\[
  D(A;A)=0.
\]
\end{lemma}


\begin{formalized}
\lean{Analysis/LiebTrace.lean}
\leanfit{cstarMatrixRelativeEntropy}
\leanfit{cstarMatrixRelativeEntropy_self}
\end{formalized}

\begin{definition}[Real diagonal C-star matrices]
For a finite index set \(I\) and a real vector \(a:I\to\mathbb R\), write
\[
  \Delta(a)=\operatorname{diag}(a_i)_{i\in I}
\]
for the corresponding complex C\(^*\)-matrix.
\end{definition}

\begin{lemma}[Diagonal logarithm]
If \(a_i\ne 0\) for every \(i\in I\), then
\[
  \log \Delta(a)=\Delta\bigl((\log a_i)_{i\in I}\bigr).
\]
\end{lemma}

\begin{corollary}[Diagonal matrix relative entropy]
If \(a_i\ne 0\) and \(b_i\ne 0\) for every \(i\in I\), then
\[
  D(\Delta(a);\Delta(b))
  =
  \sum_{i\in I}
  \left(a_i(\log a_i-\log b_i)-(a_i-b_i)\right).
\]
Consequently, if \(a_i>0\) and \(b_i>0\) for every \(i\), then
\[
  D(\Delta(a);\Delta(b))\ge 0.
\]
\end{corollary}


\begin{formalized}
\lean{Analysis/OperatorLog.lean}
\leanfit{complex_realContinuousFunctionalCalculus}
\leanfit{piComplex_realContinuousFunctionalCalculus}
\lean{Analysis/LiebTrace.lean}
\leanfit{cstarMatrixDiagonalStarAlgHom}
\leanfit{cstarMatrixDiagonalStarAlgHom_continuous}
\leanfit{cstarMatrixRealDiagonal}
\leanfit{cstarMatrixTrace_realDiagonal}
\leanfit{cstarMatrix_log_realDiagonal}
\leanfit{cstarMatrixRelativeEntropy_realDiagonal}
\leanfit{cstarMatrixRelativeEntropy_realDiagonal_nonneg}
\end{formalized}

\begin{corollary}[Real-diagonal matrix relative entropy is jointly convex]
Let \(x_i,y_i,a_i,b_i>0\) for all \(i\), and let
\(\alpha,\beta\ge0\) with \(\alpha+\beta=1\).  Then
\[
\begin{aligned}
&D\!\left(\alpha\Delta(x)+\beta\Delta(y);
         \alpha\Delta(a)+\beta\Delta(b)\right)\\
&\qquad\le
\alpha D(\Delta(x);\Delta(a))+
\beta D(\Delta(y);\Delta(b)).
\end{aligned}
\]
\end{corollary}


\begin{formalized}
\lean{Analysis/LiebTrace.lean}
\leanfit{cstarMatrixRealDiagonal_smul_add}
\leanfit{positive_weighted_sum_pos}
\leanfit{cstarMatrixRelativeEntropy_realDiagonal_jointConvex}
\end{formalized}

\begin{lemma}[Scalar-identity matrix relative entropy]
Let \(I_d\) denote the identity matrix in a finite \(d\)-dimensional
C\(^*\)-matrix algebra.  For real scalars \(a,b\),
\[
  D(aI_d;bI_d)=d\,d(a;b).
\]
Consequently, if \(a>0\) and \(b>0\), then
\[
  D(aI_d;bI_d)\ge 0.
\]
\end{lemma}


\begin{formalized}
\lean{Analysis/LiebTrace.lean}
\leanfit{cstarMatrixRelativeEntropy_algebraMap_real}
\leanfit{cstarMatrixRelativeEntropy_algebraMap_real_nonneg}
\end{formalized}

\begin{lemma}[Left/right multiplication and the ratio endomorphism]
For a finite complex C\(^*\)-matrix \(A\), let
\[
  L_A(Z)=AZ,\qquad R_A(Z)=ZA
\]
be endomorphisms of the complex vector space of finite C\(^*\)-matrices.
Then \(A\mapsto L_A\) and \(A\mapsto R_A\) are affine for real weighted sums:
\[
  L_{\alpha A+\beta B}=\alpha L_A+\beta L_B,\qquad
  R_{\alpha A+\beta B}=\alpha R_A+\beta R_B.
\]
Moreover left and right multiplication commute,
\[
  L_A R_B = R_B L_A,
\]
and they obey the product rules
\[
  L_{AB}=L_A L_B,\qquad R_{AB}=R_B R_A.
\]
Consequently, for every natural number \(k\),
\[
  L_{A^k}=L_A^k,\qquad R_{A^k}=R_A^k.
\]
and if \(A\) is a unit, then \(L_A\) and \(R_A\) are units in the
endomorphism ring.  In particular this holds for strictly positive \(A\).

If \(A\) is supplied as a unit, define the ratio endomorphism
\[
  T_{X,A}:=L_X R_{A^{-1}} .
\]
Then
\[
  T_{X,A}(Z)=XZA^{-1},
  \qquad
  T_{X,A}(A)=X.
\]
The second identity also holds when \(A\) is strictly positive, using the unit
coming from strict positivity.
\end{lemma}


\begin{formalized}
\lean{Analysis/LiebTrace.lean}
\leanfit{cstarMatrixLeftMul}
\leanfit{cstarMatrixRightMul}
\leanfit{cstarMatrixLeftMul_mul}
\leanfit{cstarMatrixRightMul_mul}
\leanfit{cstarMatrixLeftMul_pow}
\leanfit{cstarMatrixRightMul_pow}
\leanfit{cstarMatrixLeftMul_real_smul_add}
\leanfit{cstarMatrixRightMul_real_smul_add}
\leanfit{cstarMatrixLeftRightMul_commute}
\leanfit{cstarMatrixLeftMul_isUnit_of_isStrictlyPositive}
\leanfit{cstarMatrixRightMul_isUnit_of_isStrictlyPositive}
\leanfit{cstarMatrixLeftRightRatio}
\leanfit{cstarMatrixLeftRightRatio_apply}
\leanfit{cstarMatrixLeftRightRatio_apply_unit}
\leanfit{cstarMatrixLeftRightRatio_apply_of_unit_eq}
\leanfit{cstarMatrixLeftRightRatio_apply_of_isStrictlyPositive}
\end{formalized}

\begin{lemma}[Kronecker lifts for the perspective route]
Let \(A,B,H\in\mathbb C^{n\times n}\), and define
\[
  K_L(A)=A\otimes I_n,\qquad K_R(H)=I_n\otimes H.
\]
For real scalars \(\alpha,\beta\), the lifts are affine:
\[
  K_L(\alpha A+\beta B)=\alpha K_L(A)+\beta K_L(B),
  \qquad
  K_R(\alpha A+\beta B)=\alpha K_R(A)+\beta K_R(B).
\]
They also commute and multiply to the ordinary Kronecker product:
\[
  K_L(A)K_R(H)=A\otimes H=K_R(H)K_L(A).
\]
If \(A\) is positive definite, then \(K_L(A)\) and \(K_R(A)\) are positive
definite.
Under product-index vectorization,
\[
  (A\otimes B^{\mathsf T})\operatorname{vec}(M)=\operatorname{vec}(AMB).
\]
For every integer \(k\ge 0\),
\[
  (A\otimes B^{\mathsf T})^k\operatorname{vec}(M)
    =\operatorname{vec}(A^k M B^k).
\]
Finally,
\[
  \operatorname{tr}(A\otimes H)=\operatorname{tr}(A)\operatorname{tr}(H),
  \qquad
  \operatorname{tr}(K_L(A))=\operatorname{tr}(K_R(A))=
  n\,\operatorname{tr}(A).
\]
If \(v_I=\operatorname{vec}(I_n)\), then the vectorized-identity trace pairing is
\[
  v_I^*(A\otimes B^{\mathsf T})v_I=\operatorname{tr}(AB).
\]
The powered trace-pairing identity is
\[
  v_I^*(A\otimes B^{\mathsf T})^k v_I=\operatorname{tr}(A^kB^k).
\]
Consequently, for any finite \(S\subset\mathbb N\) and coefficients
\(c_k\in\mathbb C\),
\[
  v_I^*\left(\sum_{k\in S} c_k(A\otimes B^{\mathsf T})^k\right)v_I
  =
  \sum_{k\in S} c_k\,\operatorname{tr}(A^kB^k).
\]
Equivalently, Lean's polynomial-evaluation form gives the same identity for
\(p(A\otimes B^{\mathsf T})\), with the right-hand side expanded over the
finite support of \(p\).
The domain side needed for real functional calculus is also closed: if
\(A=A^*\) and \(B=B^*\), then \(A\otimes B^{\mathsf T}\) is self-adjoint, and
positive semidefiniteness or positive definiteness is preserved by the same
Kronecker-transpose lift.  In particular, if \(X\succ0\) and \(A\succ0\), then
\[
  X\otimes (A^{-1})^{\mathsf T}\succ0.
\]
Mathlib's Weierstrass approximation theorem supplies real polynomials
approximating \(\log x\), \(x\log x\), and \(x\log x-(x-1)\) uniformly on every
positive compact interval.

The source-faithful polynomial form of the Effros superoperator trace term is
now formalized.  For \(R_C=I_n\otimes C^{\mathsf T}\),
\[
  v_I^*(A\otimes B^{\mathsf T})^k R_C v_I
  =
  \operatorname{tr}(A^k C B^k),
\]
and hence, for a polynomial \(p\),
\[
  v_I^*p(A\otimes B^{\mathsf T})R_Cv_I
  =
  \sum_k p_k\,\operatorname{tr}(A^k C B^k).
\]
Specializing to the left-right ratio \(X\otimes(A^{-1})^{\mathsf T}\) and
\(C=A\), the formalized real-polynomial CFC statement is
\[
  v_I^*p(L_XR_A^{-1})R_Av_I
  =
  \sum_k p_k\,\operatorname{tr}\!\bigl(X^kA(A^{-1})^k\bigr).
\]
The same finite-polynomial trace terms now have the explicit spectral-overlap
form
\[
  \operatorname{Re}\operatorname{tr}\!\bigl(X^kA(A^{-1})^k\bigr)
  =
  \sum_{i,j}\lambda_i^k\,\mu_j\,\mu_j^{-k}
       |\langle u_i,v_j\rangle|^2,
\]
and hence
\[
  \operatorname{Re}\sum_k p_k
       \operatorname{tr}\!\bigl(X^kA(A^{-1})^k\bigr)
  =
  \sum_{i,j}\mu_j\,p(\lambda_i/\mu_j)
       |\langle u_i,v_j\rangle|^2.
\]
\end{lemma}


\begin{formalized}
\lean{Analysis/LiebTrace.lean}
\leanfit{matrixVecId}
\leanfit{matrixVec}
\leanfit{matrixComplexQuadraticForm}
\leanfit{continuous_matrixComplexQuadraticForm}
\leanfit{finset_sum_product_diagonal}
\leanfit{matrixVecId_inner_matrixVec}
\leanfit{matrix_transpose_conjTranspose_eq_self_of_isSelfAdjoint}
\leanfit{matrix_kronecker_transpose_isSelfAdjoint_of_isSelfAdjoint}
\leanfit{matrix_kronecker_transpose_posSemidef}
\leanfit{matrix_kronecker_transpose_posDef}
\leanfit{matrix_kronecker_inv_transpose_posDef}
\leanfit{matrixSelfAdjointCfc}
\leanfit{matrixSelfAdjointCfc_polynomial}
\leanfit{exists_realPolynomial_near_log_on_Icc}
\leanfit{exists_realPolynomial_near_xlog_on_Icc}
\leanfit{exists_realPolynomial_near_realEntropyKernel_on_Icc}
\leanfit{matrix_posDef_spectrum_real_pos}
\leanfit{matrix_posDef_spectrum_real_subset_Icc}
\leanfit{exists_realPolynomial_tendstoUniformlyOn_realEntropyKernel_on_subset_Icc}
\leanfit{exists_realPolynomial_tendstoUniformlyOn_realEntropyKernel_spectrum_of_posDef}
\leanfit{matrix_kronecker_transpose_mulVec_matrixVec}
\leanfit{matrix_kronecker_transpose_pow}
\leanfit{matrix_kronecker_transpose_pow_mulVec_matrixVec}
\leanfit{matrixVec_one}
\leanfit{matrixComplexQuadraticForm_vecId_kronecker_transpose_pow}
\leanfit{matrixComplexQuadraticForm_vecId_kronecker_transpose_pow_right}
\leanfit{matrixComplexQuadraticForm_sum}
\leanfit{matrixComplexQuadraticForm_smul}
\leanfit{matrixComplexQuadraticForm_vecId_kronecker_transpose_polynomial}
\leanfit{matrixComplexQuadraticForm_vecId_kronecker_transpose_polynomial_right}
\leanfit{matrixComplexQuadraticForm_vecId_kronecker_transpose_aeval}
\leanfit{matrixComplexQuadraticForm_vecId_kronecker_transpose_aeval_right}
\leanfit{matrixComplexQuadraticForm_vecId_matrixSelfAdjointCfc_kronecker_transpose_realPolynomial}
\leanfit{matrixComplexQuadraticForm_vecId_matrixSelfAdjointCfc_kronecker_inv_transpose_realPolynomial_right}
\leanfit{tendsto_matrixComplexQuadraticForm_matrixSelfAdjointCfc_mul}
\leanfit{tendsto_superoperator_entropyKernel_of_realPolynomial_uniform_approx}
\leanfit{exists_realPolynomial_tendsto_superoperator_entropyKernel_trace_of_posDef}
\leanfit{matrixComplexQuadraticForm_re_nonneg_of_posSemidef}
\leanfit{matrixComplexQuadraticForm_re_mono_of_posSemidef_sub}
\leanfit{matrixComplexQuadraticForm_re_mono_of_cstarMatrix_le}
\leanfit{matrixComplexQuadraticForm_add}
\leanfit{exists_realPolynomial_tendsto_superoperator_entropyKernel_trace_and_overlap_of_posDef}
\leanfit{matrixTrace_pow_mul_inv_pow_re_eq_sum}
\leanfit{matrixPolynomialTraceRatio_re_eq_sum}
\leanfit{tendsto_matrixPolynomialTraceRatio_overlap_sum_of_uniform_approx}
\leanfit{matrixSuperoperatorEntropyKernelTrace}
\leanfit{matrixSuperoperatorEntropyKernelOverlapExpansion}
\leanfit{matrixSuperoperatorEntropyKernelOverlapExpansion_all}
\leanfit{matrixRelativeEntropyTraceRepresentation_of_superoperator_overlap}
\leanfit{cstarMatrixRelativeEntropyTraceRepresentationBySuperoperator}
\leanfit{cstarMatrixRelativeEntropyTraceRepresentationBySuperoperator_all}
\leanfit{continuous_matrix_polynomial_aeval}
\leanfit{matrixComplexQuadraticForm_vecId_kronecker_transpose}
\leanfit{matrix_kronecker_left_identity_real_smul_add}
\leanfit{matrix_kronecker_right_identity_real_smul_add}
\leanfit{matrix_kronecker_right_identity_transpose_real_smul_add}
\leanfit{cstarMatrixSuperoperatorLeftLift}
\leanfit{cstarMatrixSuperoperatorRightLift}
\leanfit{cstarMatrixSuperoperatorLeftLift_real_smul_add}
\leanfit{cstarMatrixSuperoperatorRightLift_real_smul_add}
\leanfit{cstarMatrixSuperoperatorLeftLift_isStrictlyPositive}
\leanfit{cstarMatrixSuperoperatorRightLift_isStrictlyPositive}
\leanfit{cstarMatrixSuperoperatorLeftLift_rightLift_commute}
\leanfit{cstarMatrixSuperoperatorPositiveInvSqrtRightLift_commute_leftLift}
\leanfit{cstarMatrixSuperoperatorPositiveSqrtRightLift_commute_leftLift}
\leanfit{cstarMatrixSuperoperatorPerspective_normalizedArgument_reorder}
\leanfit{cstarMatrixSuperoperatorPerspective_normalizedArgument_eq_leftLift_mul_rightLift_unit_inv}
\leanfit{cstarMatrixSuperoperatorLeftLift_mul_rightLift_unit_inv_commute_positiveSqrtRightLift}
\leanfit{cstarMatrixSuperoperatorPerspectiveTrace}
\leanfit{cstarMatrixSuperoperatorPerspectiveTrace_jointConvex}
\leanfit{cstarMatrixRelativeEntropyPerspectiveTraceRepresentation}
\leanfit{cstarMatrixRelativeEntropyJointConvexOnStrictPositive_of_perspectiveTraceRepresentation}
\leanfit{matrix_kronecker_left_identity_mul_right_identity}
\leanfit{matrix_kronecker_right_identity_mul_left_identity}
\leanfit{matrix_kronecker_left_right_commute}
\leanfit{matrix_kronecker_posDef_left_identity}
\leanfit{matrix_kronecker_posDef_right_identity}
\leanfit{matrix_trace_kronecker}
\leanfit{matrix_trace_kronecker_left_identity}
\leanfit{matrix_trace_kronecker_right_identity}
\end{formalized}

\begin{definition}[Positive operator-convexity source target]
For a real function \(f\), the local ordinary positive-cone matrix-convexity
target is
\[
  f(\alpha T_1+\beta T_2)
  \preceq
  \alpha f(T_1)+\beta f(T_2)
\]
whenever \(T_1,T_2\) are strictly positive C\(^*\)-matrices,
\(\alpha,\beta\ge 0\), and \(\alpha+\beta=1\).  The Effros route needs this
target for
\[
  f(x)=x\log x.
\]
\end{definition}

\begin{lemma}[Identity sanity case for ordinary convexity]
The positive-cone matrix-convexity target is formalized and proved for
\(f(x)=x\).  This checks the scalar-weighted functional-calculus statement in
the local C\(^*\)-matrix vocabulary.  It is not a proof of nonlinear operator
convexity.
\end{lemma}

\begin{definition}[Hansen--Pedersen transfer target]
The Hansen--Pedersen source theorem is represented as the transfer
\[
  \begin{aligned}
  &\hbox{ordinary positive-cone operator convexity of } f \\
  &\qquad\Longrightarrow\quad
  f(A^*T_1A+B^*T_2B)
  \preceq
  A^*f(T_1)A+B^*f(T_2)B .
  \end{aligned}
\]
for strictly positive \(T_1,T_2\) and \(A^*A+B^*B=I\).  The generic
all-functions transfer theorem is still not formalized.  For the concrete
function \(x\log x\), however, the repository now proves the two-point Jensen
conclusion later in this section by a direct reflection-average argument, so no
final theorem uses this transfer target as a hidden hypothesis.
\end{definition}

\begin{lemma}[Identity sanity case for two-point Jensen]
The assembled two-point Hansen--Pedersen target is also formalized and proved
for \(f(x)=x\).  For the identity function, continuous functional calculus
reduces both sides to
\[
  A^*T_1A+B^*T_2B.
\]
The nonlinear \(x\log x\) Jensen statement is named as the downstream target and
proved later in this section by the unit-interval kernel and reflection-average
corner route.
\end{lemma}


\begin{lemma}[Derivative monotonicity for the Bendat--Sherman route]
On strictly positive C\(^*\)-matrices, the formal derivative of
\(x\log x\) is operator-monotone:
\[
  A\preceq B,\quad A\succ0,\quad B\succ0
  \quad\Longrightarrow\quad
  I+\log A \preceq I+\log B .
\]
\end{lemma}


\begin{definition}[Bendat--Sherman divided-difference route]
For \(c>0\), let the first divided difference of \(f(x)=x\log x\) at the base
point \(c\) be
\[
  f_c(x)=
  \begin{cases}
  1+\log c, & x=c,\\[2mm]
  \dfrac{x\log x-c\log c}{x-c}, & x\ne c .
  \end{cases}
\]
The source-faithful finite route separates into a monotonicity theorem for
every \(f_c\) and the Bendat--Sherman bridge:
\[
  \begin{aligned}
  &\hbox{operator monotonicity of every } f_c,\\
  &\hbox{plus the Bendat--Sherman divided-difference bridge}\\
  &\qquad\Longrightarrow
    \hbox{operator convexity of } x\log x .
  \end{aligned}
\]
on the strictly positive finite C\(^*\)-matrix cone.  The monotonicity part is
now proved below; the Bendat--Sherman bridge itself remains a local open target,
not a hidden hypothesis used by a final paper theorem.
\end{definition}

\begin{lemma}[Scalar divided-difference normalization]
For positive \(c,x\) with \(x\ne c\), the repository proves the scalar
identities
\[
  f_c(x)
  =
  \log c + \frac{x\log(x/c)}{x-c}
  =
  \log c + \frac{(x/c)\log(x/c)}{x/c-1}.
\]
The diagonal value \(f_c(c)=1+\log c\) is also formalized.
\end{lemma}


\begin{lemma}[Scalar logarithmic integral identity]
Define the normalized scalar kernel
\[
  g(t)=
  \begin{cases}
    1, & t=1,\\[2mm]
    \dfrac{t\log t}{t-1}, & t\ne1.
  \end{cases}
\]
Then the first divided difference satisfies
\[
  f_c(x)=\log c+g(x/c)
  \qquad (c>0,\ x>0).
\]
Moreover, for \(t>0\) and \(t\ne1\),
\[
  \int_0^1 \frac{t}{u+(1-u)t}\,du
  =
  \frac{t\log t}{t-1}.
\]
\end{lemma}


\begin{lemma}[Fractional-kernel monotonicity]
Let \(0\preceq A\preceq B\) be finite complex C\(^*\)-matrices.  Then the
unital CFC kernels
\[
  x\longmapsto 1-(1+x)^{-1},
  \qquad
  x\longmapsto \frac{x}{1+x}
\]
are operator-monotone on this nonnegative cone:
\[
  \left(I-(I+A)^{-1}\right)\preceq
  \left(I-(I+B)^{-1}\right),
  \qquad
  \frac{A}{I+A}\preceq \frac{B}{I+B}
\]
where the fractions are interpreted by continuous functional calculus.  The
repository also proves the spectrum wrapper needed to use nonnegativity of
\(A\) and \(B\) inside these CFC congruence arguments.  Moreover, for every
\(s>0\), the scaled fractional kernel
\[
  x\longmapsto \frac{x}{s+x}
\]
is operator-monotone on the same nonnegative cone.  Consequently, every finite
nonnegative linear combination
\[
  x\longmapsto
  \sum_{k\in F} w_k\,\frac{x}{\sigma_k+x},
  \qquad
  w_k\ge0,\quad \sigma_k>0,
\]
is operator-monotone on the same cone.
For \(0<u<1\), the unit-interval integrand
\[
  x\longmapsto \frac{x}{u+(1-u)x}
\]
is also operator-monotone, by rewriting it as a positive multiple of the
scaled kernel with \(s=u/(1-u)\).  On the strictly positive cone this is
extended to every \(u\in[0,1]\): the endpoint \(u=0\) is the constant-one CFC
kernel and the endpoint \(u=1\) is the identity kernel.
\end{lemma}


\begin{lemma}[Concrete interval-integral side conditions]
Let \(A\succ0\), and let \(M\) be any real number with
\(\lambda\le M\) for every \(\lambda\in\sigma(A)\).  For \(u\in[0,1]\) and
\(\lambda\in\sigma(A)\),
\[
  \left|\frac{\lambda}{u+(1-u)\lambda}\right|
  \le \max\{1,M\}.
\]
The same estimate is formalized both for Lebesgue measure restricted to
\([0,1]\) and for the subtype \(\{u:\mathbb R\mid u\in[0,1]\}\).  The subtype
version is the one aligned with the current \texttt{cfc\_integral} API, because
the API asks for continuity on the whole integration type.  The repository also
proves the corresponding joint-continuity statement on the interval subtype and
finite integrability of the constant bound for any finite measure on that
subtype.
\end{lemma}


\begin{theorem}[Normalized logarithmic-kernel monotonicity]
Define
\[
  g(t)=
  \begin{cases}
    1, & t=1,\\[2mm]
    \dfrac{t\log t}{t-1}, & t\ne1.
  \end{cases}
\]
For every \(t>0\),
\[
  g(t)=\int_0^1 \frac{t}{u+(1-u)t}\,du.
\]
Consequently, if \(0\prec A\preceq B\) are finite complex C\(^*\)-matrices,
then
\[
  g(A)\preceq g(B),
\]
where \(g(A)\) and \(g(B)\) are interpreted by real continuous functional
calculus.
\end{theorem}


\begin{theorem}[Divided-difference operator monotonicity for \(x\log x\)]
For every \(c>0\), the first divided difference
\[
  f_c(x)=
  \begin{cases}
  1+\log c, & x=c,\\[2mm]
  \dfrac{x\log x-c\log c}{x-c}, & x\ne c
  \end{cases}
\]
is operator-monotone on the strictly positive finite C\(^*\)-matrix cone:
\[
  0\prec A\preceq B
  \quad\Longrightarrow\quad
  f_c(A)\preceq f_c(B).
\]
\end{theorem}


\begin{corollary}[Conditional Bendat--Sherman closure]
If the finite Bendat--Sherman divided-difference bridge is proved locally,
then the concrete positive-cone operator-convexity target for \(x\log x\)
follows.  This is still only a route adapter: the bridge theorem remains open,
but the divided-difference monotonicity input is now proved.
\end{corollary}

\begin{theorem}[Finite inverse convexity by Schur complement]
Let \(A,B\in\mathbb C^{n\times n}\) be positive definite and let
\(\alpha,\beta\ge 0\) with \(\alpha+\beta=1\). Then
\[
  (\alpha A+\beta B)^{-1}
  \preceq
  \alpha A^{-1}+\beta B^{-1}.
\]
\end{theorem}


\begin{corollary}[Finite C\(^*\)-matrix inverse-kernel convexity]
Let \(A,B\) be strictly positive finite complex C\(^*\)-matrices, and let
\(\alpha,\beta\ge 0\) with \(\alpha+\beta=1\). Then the real continuous
functional calculus satisfies
\[
  f(\alpha A+\beta B)
  \preceq
  \alpha f(A)+\beta f(B),
  \qquad f(x)=x^{-1}.
\]
Equivalently, the unshifted inverse kernel is operator-convex in the finite
C\(^*\)-matrix vocabulary used by the Lieb/Tropp route.
\end{corollary}


\begin{corollary}[Shifted inverse-kernel convexity]
Let \(A,B\) be nonnegative finite complex C\(^*\)-matrices, let \(s>0\), and
let \(\alpha,\beta\ge0\) with \(\alpha+\beta=1\). Then
\[
  g_s(\alpha A+\beta B)
  \preceq
  \alpha g_s(A)+\beta g_s(B),
  \qquad g_s(x)=(s+x)^{-1}.
\]
\end{corollary}


\begin{theorem}[Direct operator convexity of \(x\log x\)]
Let \(A,B\) be strictly positive finite complex C\(^*\)-matrices and let
\(\alpha,\beta\ge0\) with \(\alpha+\beta=1\). Then
\[
  (\alpha A+\beta B)\log(\alpha A+\beta B)
  \preceq
  \alpha\,A\log A+\beta\,B\log B,
\]
where the products are interpreted through the real continuous functional
calculus \(f(x)=x\log x\).
\end{theorem}


\begin{corollary}[Closed ordinary \(x\log x\) operator-convexity target]
The Lean target \lean{cstarMatrixXLogXPositiveOperatorConvexTarget} is now
closed locally by the direct unit-interval kernel route.  The next open
Hansen--Pedersen/Effros dependency is the perspective/relative-entropy
joint-convexity layer, not ordinary positive-cone operator convexity of
\(x\log x\).
\end{corollary}

\begin{corollary}[All-finite source correction for Hansen--Pedersen]
The direct unit-interval proof also gives
\lean{cstarMatrixXLogXPositiveOperatorConvexAllFiniteTarget}: ordinary
positive-cone operator convexity of \(x\log x\) at every finite matrix size.
This is the hypothesis shape needed by the standard Hansen--Pedersen
block-matrix proof.  Consequently,
\lean{cstarMatrixXLogXHansenPedersenJensenTarget_of_allFiniteTransfer} reduces
the assembled two-point Jensen target to the source-faithful all-finite
transfer theorem
\lean{cstarMatrixXLogXHansenPedersenTransferAllFiniteTarget}.  The concrete
two-point Jensen target is closed below by a direct reflection-average proof, so
the all-finite transfer theorem remains only a generic source-route target.
\end{corollary}

\begin{lemma}[Block-compression algebra for Hansen--Pedersen]
Let
\[
  V=\begin{bmatrix}A\\ B\end{bmatrix},\qquad
  D=\begin{bmatrix}T_1&0\\0&T_2\end{bmatrix}.
\]
Then, in the finite C\(^*\)-matrix model,
\[
  V^*V=A^*A+B^*B,\qquad
  DV=\begin{bmatrix}T_1A\\T_2B\end{bmatrix},
\]
and
\[
  V^*DV=A^*T_1A+B^*T_2B.
\]
In particular, if \(A^*A+B^*B=I\), then \(V^*V=I\).  In addition, the block
diagonal constructor preserves the star-algebra operations used by the proof:
zero, identity, addition, subtraction, star, multiplication, and units.
It also preserves positivity and strict positivity:
if \(T_1,T_2\ge 0\), then
\[
  \operatorname{diag}(T_1,T_2)\ge 0,
\]
and if \(T_1,T_2\) are strictly positive, then
\[
  \operatorname{diag}(T_1,T_2)
\]
is strictly positive.  Thus the block algebra and positive-cone bookkeeping
used by the standard Hansen--Pedersen proof are now formalized; the nonlinear
functional-calculus corner and concrete Jensen transfer are handled in the
following lemmas.

The block-isometry range projection is also formalized.  If \(V=[A;B]\) and
\(P=VV^*\), then under \(V^*V=I\),
\[
  P=P^*,\qquad P^2=P,\qquad PV=V,\qquad V^*P=V^* .
\]
These are the algebraic projection identities used by the pinching/compression
route to Hansen--Pedersen Jensen.  The associated range reflection
\[
  R = 2P-I
\]
is also formalized:
\[
  R=R^*,\qquad R^2=I,\qquad R \text{ is a unitary unit},\qquad RV=V .
\]
The adjoint block row is also fixed by the same reflection, \(V^*R=V^*\), and
therefore the algebraic reflection average has the same compression as the
original block matrix:
\[
  V^*\,\frac{D+RDR}{2}\,V = V^*DV .
\]
The same averaged block is invariant under the reflection and commutes with it:
\[
  R\frac{D+RDR}{2}R=\frac{D+RDR}{2},
  \qquad
  R\frac{D+RDR}{2}=\frac{D+RDR}{2}R .
\]
It also commutes with the range projection \(P=VV^*\).
Consequently its action on the block column and adjoint block row factors
through the same compressed corner as \(D\):
\[
  \frac{D+RDR}{2}\,V = V(V^*DV),
  \qquad
  V^*\frac{D+RDR}{2}=(V^*DV)V^* .
\]
This closes the linear projection/reflection algebra and the algebraic
pinching-average compression/invariance/projection-commutation/range-reduction
identities.
The nonlinear inequality side and the \(x\log x\) corner of the pinching route
are handled below.
\end{lemma}

\begin{formalized}
\lean{Analysis/CStarMatrixBridge.lean}
\leanfit{cstarMatrixBlockDiagonal}
\leanfit{cstarMatrixBlockDiagonal_zero_zero}
\leanfit{cstarMatrixBlockDiagonal_one_one}
\leanfit{cstarMatrixBlockDiagonal_add}
\leanfit{cstarMatrixBlockDiagonal_neg}
\leanfit{cstarMatrixBlockDiagonal_sub}
\leanfit{cstarMatrixBlockDiagonal_star}
\leanfit{cstarMatrixBlockDiagonal_isSelfAdjoint}
\leanfit{cstarMatrixBlockDiagonal_mul}
\leanfit{cstarMatrixBlockDiagonal_isUnit}
\leanfit{cstarMatrixBlockDiagonal_left_nonneg}
\leanfit{cstarMatrixBlockDiagonal_right_nonneg}
\leanfit{cstarMatrixBlockDiagonal_nonneg}
\leanfit{cstarMatrixBlockDiagonal_isStrictlyPositive}
\leanfit{cstarMatrixBlockDiagonalStarAlgHom}
\leanfit{cstarMatrixBlockDiagonalStarAlgHom_continuous}
\leanfit{cstarMatrix_mul_assoc_rect}
\leanfit{cstarMatrix_mul_add_rect}
\leanfit{cstarMatrix_add_mul_rect}
\leanfit{cstarMatrix_mul_smul_rect}
\leanfit{cstarMatrix_smul_mul_rect}
\leanfit{cstarMatrix_mul_one_rect}
\leanfit{cstarMatrix_one_mul_rect}
\leanfit{cstarMatrixColumnPair}
\leanfit{cstarMatrixColumnPair_conjTranspose_mul_columnPair}
\leanfit{cstarMatrixColumnPair_conjTranspose_mul_self}
\leanfit{cstarMatrixBlockDiagonal_mul_columnPair}
\leanfit{cstarMatrixColumnPair_conjTranspose_mul_blockDiagonal_mul_columnPair}
\leanfit{cstarMatrixColumnPair_conjTranspose_mul_self_eq_one_of_sum}
\leanfit{cstarMatrixColumnPairRangeProjection}
\leanfit{cstarMatrixColumnPairRangeProjection_isSelfAdjoint}
\leanfit{cstarMatrixColumnPairRangeProjection_mul_self_of_sum}
\leanfit{cstarMatrixColumnPairRangeProjection_mul_columnPair_of_sum}
\leanfit{cstarMatrixColumnPair_conjTranspose_mul_rangeProjection_of_sum}
\leanfit{cstarMatrixProjectionReflection}
\leanfit{cstarMatrixProjectionReflection_isSelfAdjoint_of_isSelfAdjoint}
\leanfit{cstarMatrixProjectionReflection_mul_self_of_idempotent}
\leanfit{cstarMatrixProjectionReflection_isUnit_of_idempotent}
\leanfit{cstarMatrixProjectionReflection_mem_unitary_of_isSelfAdjoint_of_idempotent}
\leanfit{cstarMatrixProjectionReflection_mul_of_mul_eq_self}
\leanfit{cstarMatrix_mul_projectionReflection_of_mul_eq_self}
\leanfit{cstarMatrix_reflectionAverage_compression_of_fixed}
\leanfit{cstarMatrix_reflectionAverage_conj_of_involutive}
\leanfit{cstarMatrix_reflectionAverage_commute_of_involutive}
\leanfit{cstarMatrix_commute_projection_of_commute_reflection}
\leanfit{cstarMatrixColumnPairRangeReflection}
\leanfit{cstarMatrixColumnPairRangeReflection_isSelfAdjoint}
\leanfit{cstarMatrixColumnPairRangeReflection_mul_self_of_sum}
\leanfit{cstarMatrixColumnPairRangeReflection_isUnit_of_sum}
\leanfit{cstarMatrixColumnPairRangeReflection_mem_unitary_of_sum}
\leanfit{cstarMatrixColumnPairRangeReflection_mul_columnPair_of_sum}
\leanfit{cstarMatrixColumnPair_conjTranspose_mul_rangeReflection_of_sum}
\leanfit{cstarMatrixColumnPair_reflectionAverage_compression_of_sum}
\leanfit{cstarMatrixColumnPair_reflectionAverage_conj_rangeReflection_of_sum}
\leanfit{cstarMatrixColumnPair_reflectionAverage_commute_rangeReflection_of_sum}
\leanfit{cstarMatrixColumnPair_reflectionAverage_commute_rangeProjection_of_sum}
\leanfit{cstarMatrixColumnPair_mul_columnPair_eq_columnPair_compression_of_commute}
\leanfit{cstarMatrixColumnPair_conjTranspose_mul_eq_compression_mul_conjTranspose_of_commute}
\leanfit{cstarMatrixColumnPair_reflectionAverage_mul_columnPair_of_sum}
\leanfit{cstarMatrixColumnPair_conjTranspose_mul_reflectionAverage_of_sum}
\end{formalized}

\begin{lemma}[Block-diagonal functional calculus]
Let \(T_1,T_2\) be self-adjoint finite C\(^*\)-matrices, and let
\(f:\mathbb{R}\to\mathbb{R}\) be continuous on
\(\sigma(T_1)\cup\sigma(T_2)\).  Then continuous functional calculus commutes
with the block diagonal:
\[
  f\!\left(\begin{bmatrix}T_1&0\\0&T_2\end{bmatrix}\right)
  =
  \begin{bmatrix}f(T_1)&0\\0&f(T_2)\end{bmatrix}.
\]
This closes the block-diagonal CFC part of the Hansen--Pedersen block proof.
It does not by itself prove the compression inequality
\[
  f(V^*DV)\le V^*f(D)V.
\]
\end{lemma}

\begin{formalized}
\lean{Analysis/LiebTrace.lean}
\leanfit{cstarMatrixBlockDiagonal_cfc}
\end{formalized}

\begin{lemma}[Unitary-conjugation functional calculus]
Let \(U\) be a unitary finite C\(^*\)-matrix, let \(T=T^*\), and let
\(f:\mathbb{R}\to\mathbb{R}\) be continuous on \(\sigma(T)\).  Then
\[
  f(UTU^*) = U f(T) U^* .
\]
This closes the CFC-conjugation part of the range-reflection route.  It does
not prove the nonlinear CFC pinching inequality or the compression Jensen
inequality.
\end{lemma}

\begin{formalized}
\lean{Analysis/LiebTrace.lean}
\leanfit{cstarMatrix_cfc_unitary_conj}
\end{formalized}

\begin{lemma}[Reflection-average CFC pinching inequality]
Let \(V=[A;B]\), assume \(V^*V=I\), and put
\[
  R=2VV^*-I,\qquad E=\frac{D+RDR}{2}.
\]
Assume \(D\) is strictly positive and \(f\) is continuous on \(\sigma(D)\).
If \(f\) is ordinary operator-convex on the positive cone at every finite
matrix size, then
\[
  f(E)\le \frac{f(D)+Rf(D)R}{2}.
\]
Moreover, after compression by \(V\),
\[
  V^* f(E) V \le V^* f(D) V .
\]
\end{lemma}


\begin{formalized}
\lean{Analysis/LiebTrace.lean}
\leanfit{cstarMatrix_compression_nonneg}
\leanfit{cstarMatrix_compression_mono}
\leanfit{cstarMatrixColumnPair_reflectionAverage_cfc_le_average_of_sum}
\leanfit{cstarMatrixColumnPair_reflectionAverage_compressed_cfc_le_compressed_of_sum}
\end{formalized}

\begin{remark}
This is a genuine new closure in the Hansen--Pedersen route, but it is still
not the final concrete \(x\log x\) Jensen theorem.  The next nonlinear corner
step is
\[
  V^*f(E)V=f(V^*DV),
\]
first for the shifted-inverse kernels and then for \(x\log x\), or an
alternative Effros/perspective proof that avoids this corner identity.
\end{remark}

\begin{lemma}[Shifted-inverse corner identity]
Let \(V=[A;B]\), assume \(V^*V=I\), and put
\[
  R=2VV^*-I,\qquad E=\frac{D+RDR}{2},\qquad X=V^*DV .
\]
If \(D\) is strictly positive and \(s>0\), then the shifted inverse kernel
satisfies the nonlinear corner identity
\[
  V^*(sI+E)^{-1}V=(sI+X)^{-1}.
\]
Equivalently, in real continuous-functional-calculus notation,
\[
  V^*\,\bigl[x\mapsto(s+x)^{-1}\bigr](E)\,V
  =
  \bigl[x\mapsto(s+x)^{-1}\bigr](V^*DV).
\]
\end{lemma}


\begin{formalized}
\lean{Analysis/CStarMatrixBridge.lean}
\leanfit{cstarMatrix_units_inv_mul_rect_eq_mul_units_inv_of_mul_eq}
\lean{Analysis/LiebTrace.lean}
\leanfit{cstarMatrix_cfc_shifted_inv_mul_rect_eq_mul_cfc_shifted_inv_of_mul_eq}
\leanfit{cstarMatrixColumnPair_reflectionAverage_shifted_inv_corner_of_sum}
\end{formalized}

\begin{remark}
This closes the nonlinear corner identity for the shifted-inverse kernels used
in the direct \(x\log x\) integral representation.  The next theorem assembles
the affine and Bochner-integral pieces into the full \(x\log x\) corner
identity and then the concrete two-point Hansen--Pedersen Jensen statement.
\end{remark}

\begin{theorem}[Concrete Hansen--Pedersen Jensen for \(x\log x\)]
Let \(A,B,T_1,T_2\) be finite complex C\(^*\)-matrices.  Assume
\[
  T_1\succ0,\qquad T_2\succ0,\qquad A^*A+B^*B=I .
\]
Then
\[
  f(A^*T_1A+B^*T_2B)
  \preceq
  A^*f(T_1)A+B^*f(T_2)B,
  \qquad f(x)=x\log x .
\]
Equivalently, the concrete target
\lean{cstarMatrixXLogXHansenPedersenJensenTarget} is now proved locally.
\end{theorem}


\begin{formalized}
\lean{Analysis/CStarMatrixBridge.lean}
\leanfit{cstarMatrix_complex_finiteDimensional}
\leanfit{cstarMatrix_compression_add}
\leanfit{cstarMatrix_compression_sub}
\leanfit{cstarMatrix_compression_smul}
\leanfit{cstarMatrix_compression_real_smul}
\leanfit{cstarMatrixCompressionCLM}
\leanfit{cstarMatrixCompressionCLM_apply}
\leanfit{cstarMatrix_compression_one_of_conjTranspose_mul_self_eq_one}
\lean{Analysis/LiebTrace.lean}
\leanfit{cstarMatrix_compression_setIntegral}
\leanfit{cstarMatrix_compression_isStrictlyPositive_of_injective_mulVec}
\leanfit{cstarMatrixColumnPair_compression_isStrictlyPositive_of_sum}
\leanfit{cstarMatrixColumnPair_reflectionAverage_xlog_kernel_corner_of_sum}
\leanfit{cstarMatrixColumnPair_reflectionAverage_xlog_corner_of_sum}
\leanfit{cstarMatrixXLogXHansenPedersenJensenTarget_of_reflectionAverage_xlog_corner}
\leanfit{cstarMatrixXLogXHansenPedersenJensenTarget_of_unit_interval_kernel}
\end{formalized}

\begin{remark}
This closes the concrete two-point \(x\log x\) Jensen dependency needed by the
chosen Effros/Tropp source route.  This lemma alone does not prove the
perspective theorem or the trace representation of matrix relative entropy.
The ordinary finite perspective theorem for the normalized entropy kernel is
stated below; the source-faithful superoperator trace representation, matrix
relative-entropy joint convexity, arbitrary-\(H\) Lieb trace concavity,
trace-MGF domination, and matrix Bernstein remain open.
\end{remark}

\begin{corollary}[Normalized entropy-kernel Jensen]
Let
\[
  \varphi(x)=x\log x-(x-1).
\]
The ordinary positive-cone operator-convexity target for \(\varphi\) holds at
every finite matrix size.  In addition, under the same hypotheses as the
preceding theorem,
\[
  \varphi(A^*T_1A+B^*T_2B) \preceq
  A^*\varphi(T_1)A+B^*\varphi(T_2)B .
\]
\end{corollary}


\begin{formalized}
\lean{Analysis/LiebTrace.lean}
\leanfit{realEntropyKernel}
\leanfit{cstarMatrixEntropyKernelPositiveOperatorConvexTarget}
\leanfit{cstarMatrixEntropyKernelPositiveOperatorConvexAllFiniteTarget}
\leanfit{cstarMatrixEntropyKernelPositiveOperatorConvexTarget_of_xlog}
\leanfit{cstarMatrixEntropyKernelPositiveOperatorConvexTarget_of_unit_interval_kernel}
\leanfit{cstarMatrixEntropyKernelPositiveOperatorConvexAllFiniteTarget_of_unit_interval_kernel}
\leanfit{cstarMatrixEntropyKernelHansenPedersenJensenTarget}
\leanfit{cstarMatrix_cfc_realEntropyKernel_eq_xlog_sub_id_add_one}
\leanfit{cstarMatrixEntropyKernelHansenPedersenJensenTarget_of_xlog}
\leanfit{cstarMatrixEntropyKernelHansenPedersenJensenTarget_of_unit_interval_kernel}
\end{formalized}

\begin{remark}
This corollary still does not prove the finite operator-perspective theorem,
the trace representation of matrix relative entropy, or the noncommutative
joint-convexity theorem.  It closes only the affine-corrected scalar kernel
needed by that perspective route.
\end{remark}

\begin{lemma}[Perspective square-root substrate]
Let \(A\) be a strictly positive finite C\(^*\)-matrix.  The real continuous
functional calculus defines
\[
  A^{1/2}:=\sqrt{\cdot}(A), \qquad
  A^{-1/2}:=(\sqrt{\cdot})^{-1}(A).
\]
Then
\[
  A^{1/2}A^{1/2}=A,\qquad
  A^{-1/2}A^{1/2}=A^{1/2}A^{-1/2}=I,\qquad
  A^{-1/2}AA^{-1/2}=I.
\]
Moreover, if \(X\succ0\), then \(A^{-1/2}XA^{-1/2}\succ0\).
\end{lemma}


\begin{formalized}
\lean{Analysis/LiebTrace.lean}
\leanfit{cstarMatrixPositiveSqrt}
\leanfit{cstarMatrixPositiveInvSqrt}
\leanfit{cstarMatrixPositiveSqrt_isSelfAdjoint}
\leanfit{cstarMatrixPositiveInvSqrt_isSelfAdjoint}
\leanfit{continuousOn_real_inv_sqrt_spectrum_of_isStrictlyPositive}
\leanfit{cstarMatrixPositiveSqrt_mul_self}
\leanfit{cstarMatrixPositiveInvSqrt_mul_sqrt}
\leanfit{cstarMatrixPositiveSqrt_mul_invSqrt}
\leanfit{cstarMatrixPositiveInvSqrt_isUnit}
\leanfit{complex_ofReal_sqrt_mul_self_of_nonneg}
\leanfit{cstarMatrixPositiveSqrt_isStrictlyPositive}
\leanfit{cstarMatrixPositiveInvSqrt_mul_self_mul}
\leanfit{cstarMatrixPositiveInvSqrt_mul_self_eq_unit_inv}
\leanfit{cstarMatrixPositiveSqrt_commute_unit_inv}
\leanfit{cstarMatrixPositiveInvSqrt_conj_isStrictlyPositive}
\end{formalized}

\begin{theorem}[Ordinary finite entropy-kernel perspective]
Let \(f(x)=x\log x-(x-1)\), and define the finite C\(^*\)-matrix
perspective
\[
  P_f(X,A)=A^{1/2}f\!\left(A^{-1/2}XA^{-1/2}\right)A^{1/2}.
\]
If \(X,Y,A,B\succ0\), \(a,b\ge0\), and \(a+b=1\), then
\[
  P_f(aX+bY,aA+bB)
  \preceq
  aP_f(X,A)+bP_f(Y,B).
\]
\end{theorem}


\begin{formalized}
\lean{Analysis/LiebTrace.lean}
\leanfit{cstarMatrixPerspective}
\leanfit{cstarMatrixPerspectiveWeight}
\leanfit{cstarMatrixPerspectiveWeight_star_mul_self}
\leanfit{cstarMatrixPerspectiveWeights_star_mul_self_add}
\leanfit{cstarMatrixPerspectiveWeight_compress_normalized}
\leanfit{cstarMatrixPerspectiveWeight_mul_positiveSqrt}
\leanfit{cstarMatrixPositiveSqrt_mul_perspectiveWeight_star}
\leanfit{cstarMatrixPerspectiveWeight_value_uncompress}
\leanfit{cstarMatrixEntropyKernelPerspective_jointConvex}
\end{formalized}

\begin{remark}
This is a dependency closure for the Effros perspective route.  It proves the
ordinary finite perspective theorem for the normalized entropy kernel, but it
does not by itself prove the Umegaki relative-entropy trace representation.
The source-faithful route for \(D(X;A)\) still requires the superoperator
perspective built from \(L_XR_A^{-1}\) and the trace identity connecting that
superoperator perspective to matrix relative entropy.
\end{remark}

\begin{lemma}[Unitary-conjugation domain preservation]
Let \(U\) be a unitary finite C\(^*\)-matrix and let \(T\) be strictly
positive.  Then \(UTU^*\) is strictly positive.  In particular, if
\(V^*V=I\) and \(R=2VV^*-I\) is the range reflection, then \(RTR^*\) is
strictly positive whenever \(T\) is strictly positive.

This closes the strict-positive-domain side condition needed before applying
continuous functional calculus along the reflection/pinching route.  It does
not prove the nonlinear corner CFC identity or compression Jensen
inequality.
\end{lemma}

\begin{formalized}
\lean{Analysis/CStarMatrixBridge.lean}
\leanfit{cstarMatrix_unitary_conj_isStrictlyPositive}
\leanfit{cstarMatrixColumnPairRangeReflection_conj_isStrictlyPositive_of_sum}
\end{formalized}

\begin{lemma}[Strict positivity of the compressed block argument]
Let \(V=[A;B]\) and assume \(V^*V=A^*A+B^*B=I\).  If \(T_1\) and \(T_2\)
are strictly positive, then
\[
  V^*
  \begin{bmatrix}T_1&0\\0&T_2\end{bmatrix}
  V
  =
  A^*T_1A+B^*T_2B
\]
is strictly positive.

The proof bridges strict positivity to ordinary finite matrix positive
definiteness, uses \(V^*V=I\) to prove that the matrix action of \(V\) is
injective, applies preservation of positive definiteness under injective
compression, and bridges back to the C\(^*\)-matrix order.  This is a domain
lemma only; it does not prove the nonlinear Jensen inequality.
\end{lemma}

\begin{formalized}
\lean{Analysis/LiebTrace.lean}
\leanfit{cstarMatrix_isStrictlyPositive_of_matrix_posDef}
\leanfit{cstarMatrixColumnPair_mulVec_injective_of_sum}
\leanfit{cstarMatrixColumnPair_compress_blockDiagonal_isStrictlyPositive_of_sum}
\leanfit{cstarMatrixHansenPedersenCompression_isStrictlyPositive_of_sum}
\end{formalized}

\begin{formalized}
\lean{Analysis/LiebTrace.lean}
\leanfit{cstarMatrixPositiveOperatorConvexTarget}
\leanfit{cstarMatrixPositiveOperatorConvexAllFiniteTarget}
\leanfit{cstarMatrixPositiveOperatorConvexTarget_id}
\leanfit{cstarMatrixPositiveOperatorConvexAllFiniteTarget_id}
\leanfit{cstarMatrixPositiveHansenPedersenTransferTarget}
\leanfit{cstarMatrixPositiveHansenPedersenTransferAllFiniteTarget}
\leanfit{cstarMatrixXLogXPositiveOperatorConvexTarget}
\leanfit{cstarMatrixXLogXHansenPedersenTransferTarget}
\leanfit{cstarMatrixXLogXPositiveOperatorConvexAllFiniteTarget}
\leanfit{cstarMatrixXLogXHansenPedersenTransferAllFiniteTarget}
\leanfit{cstarMatrixHansenPedersenJensenTwoPointTarget}
\leanfit{cstarMatrixHansenPedersenJensenTwoPointTarget_id}
\leanfit{cstarMatrixXLogXHansenPedersenJensenTarget}
\leanfit{cstarMatrixXLogXHansenPedersenJensenTarget_of_positiveOperatorConvex_of_transfer}
\leanfit{cstarMatrixXLogXHansenPedersenJensenTarget_of_transfer}
\leanfit{cstarMatrixXLogXPositiveOperatorConvexAllFiniteTarget_of_unit_interval_kernel}
\leanfit{cstarMatrixXLogXHansenPedersenJensenTarget_of_allFiniteTransfer}
\leanfit{cstarMatrixXLogXPositiveOperatorConvexTarget_of_unit_interval_kernel}
\leanfit{cstarMatrix_cfc_one_add_log_eq_one_add_log}
\leanfit{cstarMatrixXLogXDerivativeMonotoneTarget}
\leanfit{cstarMatrixXLogXDerivativeMonotoneTarget_of_log_monotone}
\leanfit{cstarMatrixStrictPositiveOperatorMonotoneTarget}
\leanfit{realXLogXDividedDifference}
\leanfit{realXLogXDividedDifference_self}
\leanfit{realXLogXDividedDifference_eq_log_add_ratio}
\leanfit{realXLogXDividedDifference_eq_log_add_normalized}
\leanfit{realNormalizedLogKernel}
\leanfit{realNormalizedLogKernel_eq_of_ne_one}
\leanfit{realNormalizedLogKernel_eq_mul_dslope_log}
\leanfit{continuousOn_realNormalizedLogKernel_Ioi}
\leanfit{realXLogXDividedDifference_eq_log_add_normalizedKernel}
\leanfit{real_normalizedLogKernel_offdiag_intervalIntegral}
\leanfit{realNormalizedLogKernel_setIntegral}
\leanfit{cstarMatrix_spectrum_nonneg_of_nonneg}
\leanfit{cstarMatrix_cfc_one_sub_one_add_inv_monotone}
\leanfit{cstarMatrix_cfc_pos_over_one_add_monotone}
\leanfit{cstarMatrix_cfc_pos_over_pos_add_monotone}
\leanfit{cstarMatrix_cfc_unit_interval_fractional_kernel_monotone}
\leanfit{cstarMatrix_cfc_unit_interval_fractional_kernel_monotone_of_mem_Icc}
\leanfit{continuousOn_uncurry_unit_interval_fractional_kernel_spectrum}
\leanfit{real_unit_interval_fractional_kernel_abs_le_max_of_le}
\leanfit{real_unit_interval_fractional_kernel_spectrum_norm_le_max}
\leanfit{ae_unit_interval_fractional_kernel_spectrum_norm_le_max}
\leanfit{hasFiniteIntegral_const_max_one_spectrum_bound}
\leanfit{continuousOn_uncurry_unit_interval_subtype_fractional_kernel_spectrum}
\leanfit{ae_unit_interval_subtype_fractional_kernel_spectrum_norm_le_max}
\leanfit{hasFiniteIntegral_unit_interval_subtype_const_max_one_spectrum_bound}
\leanfit{cstarMatrix_cfc_realNormalizedLogKernel_eq_unit_interval_integral}
\leanfit{cstarMatrix_setIntegral_mono_on}
\leanfit{cstarMatrix_cfc_realNormalizedLogKernel_monotone_of_spectrum_bound}
\leanfit{cstarMatrix_cfc_realNormalizedLogKernel_monotone}
\leanfit{cstarMatrix_cfc_realXLogXDividedDifference_eq_log_add_scaled_normalizedKernel}
\leanfit{cstarMatrix_cfc_finset_sum_nonneg_mul_pos_over_pos_add_monotone}
\leanfit{cfc_integral_mono_of_forall_of_bound}
\leanfit{cstarMatrixXLogXDividedDifferenceMonotoneTarget}
\leanfit{cstarMatrixXLogXDividedDifferenceMonotoneTarget_of_normalizedLogKernel}
\leanfit{cstarMatrixBendatShermanDividedDifferenceBridgeTarget}
\leanfit{cstarMatrixXLogXPositiveOperatorConvexTarget_of_bendatShermanDividedDifferenceBridge}
\leanfit{cstarMatrixBendatShermanDerivativeBridgeTarget}
\leanfit{cstarMatrixXLogXPositiveOperatorConvexTarget_of_bendatShermanDerivativeBridge}
\leanfit{matrix_posDef_inverse_schur_block}
\leanfit{matrix_weighted_inverse_schur_block}
\leanfit{matrix_posDef_weighted_sum}
\leanfit{matrix_inv_convex_posDef}
\leanfit{cstarMatrix_nonneg_of_matrix_posSemidef}
\leanfit{cstarMatrix_le_of_matrix_le}
\leanfit{cstarMatrix_cfc_inv_convex_isStrictlyPositive}
\leanfit{cstarMatrix_cfc_shifted_inv_eq_cfc_inv_add_smul_one}
\leanfit{cstarMatrix_cfc_shifted_inv_convex_nonneg}
\leanfit{real_xlog_eq_unit_interval_xlog_kernel_integral}
\leanfit{real_unit_interval_xlog_kernel_integrand_eq_affine_add_shifted_inv}
\leanfit{cstarMatrix_cfc_unit_interval_xlog_kernel_integrand_eq_affine_add_shifted_inv}
\leanfit{continuousOn_uncurry_unit_interval_xlog_kernel_spectrum}
\leanfit{real_unit_interval_xlog_kernel_spectrum_norm_le_max_sq}
\leanfit{ae_unit_interval_xlog_kernel_spectrum_norm_le_max_sq}
\leanfit{hasFiniteIntegral_const_max_one_spectrum_bound_sq}
\leanfit{cstarMatrix_cfc_xlog_eq_unit_interval_xlog_kernel_integral}
\leanfit{cstarMatrix_cfc_unit_interval_xlog_integrand_convex_of_pos_lt_one}
\leanfit{cstarMatrix_cfc_unit_interval_xlog_kernel_integrand_convex_of_pos_lt_one}
\end{formalized}

\begin{definition}[Relative-entropy variational route]
For fixed \(H\), define the variational objective
\[
  \Psi_H(X,A)
  =
  \operatorname{Re}\operatorname{tr}(XH)-D(X;A)
  +\operatorname{Re}\operatorname{tr}(A).
\]
The extra trace term matches the repository normalization
\[
  D(X;A)=
  \operatorname{Re}\operatorname{tr}
  \bigl(X(\log X-\log A)-(X-A)\bigr).
\]
At this point in the development, the generalized Klein,
variational-formula, relative-entropy joint-convexity, finite-dimensional
Lieb, one-step Tropp trace-MGF, independent-sum trace-MGF iteration, and
finite-real trace-exponential transport foundations are closed locally.  The
remaining concentration foundation is the scalar matrix-CGF/log-MGF bound that
feeds the matrix Bernstein/Khintchine tail conversion.
\end{definition}

\begin{lemma}[Optimizer-candidate value]
With the normalized objective above, the candidate
\[
  X_\star=\exp(H+\log A)
\]
satisfies
\[
  \Psi_H(X_\star,A)=\Phi_H(A).
\]
This is only the attainment algebra for the expected optimizer.  The global
upper bound
\(\Psi_H(X,A)\le \Phi_H(A)\) for all \(X\succ0\) is supplied below through
the local generalized Klein and relative-entropy nonnegativity theorem.
\end{lemma}


\begin{lemma}[Variational formula from relative-entropy nonnegativity]
If matrix relative entropy is nonnegative on the strictly positive cone, then
the normalized variational formula holds.  Indeed, for
\(X_\star=\exp(H+\log A)\),
\[
  \Psi_H(Y,A)=\Phi_H(A)-D(Y;X_\star).
\]
Thus the missing maximality inequality
\(\Psi_H(Y,A)\le\Phi_H(A)\) is exactly the nonnegativity of
\(D(Y;X_\star)\).
\end{lemma}


\begin{lemma}[Relative-entropy nonnegativity from generalized Klein]
Let
\[
  \Phi(X)=\operatorname{Re}\operatorname{tr}(X\log X-X).
\]
Assume the generalized Klein first-order trace inequality on the strictly
positive cone:
\[
  \Phi(A)
  +\operatorname{Re}\operatorname{tr}\bigl((\log A)(X-A)\bigr)
  \le \Phi(X)
  \qquad (X,A\succ0).
\]
Then the local matrix relative entropy is nonnegative:
\[
  D(X;A)
  =
  \operatorname{Re}\operatorname{tr}
  \bigl(X(\log X-\log A)-(X-A)\bigr)
  \ge 0.
\]
\end{lemma}


\begin{corollary}[Klein and local relative entropy are equivalent]
In the repository normalization
\[
  D(X;A)=\operatorname{Re}\operatorname{tr}
  \bigl(X(\log X-\log A)-(X-A)\bigr),
\]
the generalized Klein first-order trace inequality above is equivalent to
nonnegativity of \(D(X;A)\) on the strictly positive cone.
\end{corollary}


\begin{lemma}[Hermitian spectral-overlap expansion]
Let \(A\) and \(B\) be Hermitian \(d\times d\) complex matrices.  For finite
families of scalar functions \(f_r,g_r:\mathbb R\to\mathbb R\),
\[
  \operatorname{Re}\sum_r
  \operatorname{tr}\!\left(f_r(A)g_r(B)\right)
  =
  \sum_{j,k}
  \left(\sum_r f_r(\lambda_j(A))g_r(\lambda_k(B))\right)
  |\langle u_j,v_k\rangle|^2,
\]
where \(\{u_j\}\) and \(\{v_k\}\) are Hermitian eigenbases for \(A\) and
\(B\).  Hence a nonnegative scalar kernel gives a nonnegative trace sum.
\end{lemma}


\begin{corollary}[Entropy first-order kernel on positive spectra]
For \(\phi(t)=t\log t-t\) and \(\psi(t)=\log t\), if the eigenvalues of
Hermitian \(A\) and \(B\) are positive, then
\[
\begin{aligned}
0 \le \operatorname{Re}\operatorname{tr}\big(&
  \phi(A)I
  - I\phi(B)
  - A\psi(B)
  + I(\psi(B)B)\big).
\end{aligned}
\]
Equivalently, this is the separated-kernel Hermitian form of
\(\operatorname{tr}\phi(A)-\operatorname{tr}\phi(B)
-\operatorname{tr}(\psi(B)(A-B))\ge0\).
\end{corollary}


\begin{theorem}[Compact Hermitian generalized Klein inequality]
Let \(X,A\) be Hermitian complex matrices whose eigenvalues are all strictly
positive.  Then
\[
  \operatorname{Re}\operatorname{tr}
  \bigl(A\log A-A+\log(A)(X-A)\bigr)
  \le
  \operatorname{Re}\operatorname{tr}(X\log X-X).
\]
\end{theorem}


\begin{lemma}[Strictly positive C-star matrices have positive spectra]
If \(M\) is nonnegative as a finite complex C\(^*\)-matrix, then its underlying
Hermitian matrix is positive semidefinite.  If \(M\) is strictly positive,
then every Hermitian eigenvalue of the underlying matrix is strictly positive.
\end{lemma}


\begin{theorem}[Generalized Klein inequality on the strict cone]
For all strictly positive finite complex C\(^*\)-matrices \(X,A\),
\[
  \Phi(A)
  +\operatorname{Re}\operatorname{tr}\bigl(\log(A)(X-A)\bigr)
  \le
  \Phi(X),
  \qquad
  \Phi(Y)=\operatorname{Re}\operatorname{tr}(Y\log Y-Y).
\]
\end{theorem}


\begin{corollary}[Matrix relative entropy is nonnegative]
For all strictly positive finite complex C\(^*\)-matrices \(X,A\),
\[
  D(X;A)=
  \operatorname{Re}\operatorname{tr}
  \bigl(X(\log X-\log A)-(X-A)\bigr)
  \ge 0.
\]
\end{corollary}


\begin{corollary}[Normalized entropy variational formula]
For self-adjoint \(H\), the normalized variational formula holds:
\[
  \Phi_H(A)
  =
  \sup_{X\succ0}\Psi_H(X,A),
  \qquad
  \Psi_H(X,A)=
  \operatorname{Re}\operatorname{tr}(XH)-D(X;A)
  +\operatorname{Re}\operatorname{tr}(A).
\]
\end{corollary}


\begin{corollary}[Joint convexity alone now implies the local Lieb target]
Assume matrix relative entropy is jointly convex on the strictly positive
cone.  Then, for every self-adjoint \(H\),
\[
  A\mapsto
  \operatorname{Re}\operatorname{tr}\exp(H+\log A)
\]
is concave on the strictly positive cone.
\end{corollary}


\begin{theorem}[Finite-dimensional relative entropy is jointly convex]
For finite complex matrices \(X_1,X_2,A_1,A_2\succ0\) and
\(\lambda,\mu\ge0\) with \(\lambda+\mu=1\),
\[
  D(\lambda X_1+\mu X_2;\lambda A_1+\mu A_2)
  \le
  \lambda D(X_1;A_1)+\mu D(X_2;A_2).
\]
\end{theorem}


\begin{corollary}[Finite-dimensional Lieb trace concavity]
If \(H=H^\ast\), then
\[
  A\mapsto
  \operatorname{Re}\operatorname{tr}\exp(H+\log A)
\]
is concave on the strictly positive cone.
\end{corollary}


\begin{corollary}[One-step Tropp trace-MGF inequality]
For a finite probability space, a self-adjoint \(H\), and a self-adjoint
matrix random variable \(X\),
\[
  \mathbb E\,\operatorname{Re}\operatorname{tr}\exp(H+X)
  \le
  \operatorname{Re}\operatorname{tr}
  \exp\!\left(H+\log \mathbb E\,\exp X\right).
\]
\end{corollary}


\begin{theorem}[Relative-entropy route implies Lieb concavity]
Assume \(H\) is self-adjoint.  If matrix relative entropy is jointly convex
on the strictly positive cone and the variational formula above holds for
this \(H\), then
\[
  A\mapsto
  \operatorname{Re}\operatorname{tr}\exp(H+\log A)
\]
is concave on the strictly positive cone.
\end{theorem}


\begin{lemma}[Spectral-overlap form of the compact relative-entropy trace]
Let \(X,A\in\mathbb C^{d\times d}\) be Hermitian.  Write
\(\lambda_i\) and \(\mu_j\) for their eigenvalues, and write
\[
  W=U_X^\ast U_A
\]
for the unitary overlap between their eigenvector bases.  Then the repository's
compact trace expression satisfies
\[
  \operatorname{Re}\operatorname{tr}\!\left(
    X\log X-X\log A-(X-A)\right)
  =
  \sum_{i,j}
    \bigl(\lambda_i(\log\lambda_i-\log\mu_j)
      -(\lambda_i-\mu_j)\bigr)|W_{ij}|^2 .
\]
\end{lemma}


\begin{formalized}
\lean{Analysis/LiebTrace.lean}
\leanfit{cstarMatrixEntropyVariationalObjective}
\leanfit{cstarMatrixEntropyVariationalObjective_liebOptimizer}
\leanfit{cstarMatrixRelativeEntropyNonnegOnStrictPositive}
\leanfit{cstarMatrixEntropyTraceFirstOrderConvexityOnStrictPositive}
\leanfit{cstarMatrixRelativeEntropyNonnegOnStrictPositive_of_entropyTraceFirstOrder}
\leanfit{cstarMatrixEntropyTraceFirstOrderConvexityOnStrictPositive_of_relativeEntropy_nonneg}
\leanfit{cstarMatrixEntropyTraceFirstOrderConvexityOnStrictPositive_iff_relativeEntropy_nonneg}
\leanfit{matrixTrace_diagonal_mul_mul_diagonal_mul_star}
\leanfit{matrixTrace_sum_diagonal_mul_mul_diagonal_mul_star_re}
\leanfit{matrixTrace_sum_hermitianCfc_mul_cfc_re}
\leanfit{matrixTrace_sum_hermitianCfc_mul_cfc_nonneg_of_kernel_nonneg}
\leanfit{matrixTrace_sum_hermitianCfc_mul_cfc_nonneg_of_eigen_kernel_nonneg}
\leanfit{matrixTrace_hermitianCfc_firstOrderKernel_sum_nonneg}
\leanfit{matrixTrace_hermitianCfc_firstOrderKernel_sum_nonneg_of_eigen}
\leanfit{realEntropy_firstOrderKernel_nonneg}
\leanfit{matrixTrace_hermitianCfc_entropy_firstOrder_sum_nonneg}
\leanfit{matrix_isHermitian_cfc_const_one}
\leanfit{matrix_isHermitian_cfc_const_neg_one}
\leanfit{matrix_isHermitian_cfc_neg_id}
\leanfit{matrix_isHermitian_cfc_id}
\leanfit{matrix_isHermitian_cfc_entropy}
\leanfit{matrix_isHermitian_cfc_xlog}
\leanfit{matrix_isHermitian_cfc_log_mul_id}
\leanfit{matrixTrace_hermitianCfc_relativeEntropy_re_eq_sum}
\leanfit{matrixTrace_hermitianCfc_entropy_firstOrder_compact_nonneg}
\leanfit{cstarMatrix_nonneg_to_matrix_posSemidef}
\leanfit{cstarMatrix_isStrictlyPositive_to_matrix_posDef}
\leanfit{cstarMatrixEntropyTraceFirstOrderConvexityOnStrictPositive_of_hermitianCfc}
\leanfit{cstarMatrixRelativeEntropyNonnegOnStrictPositive_of_hermitianCfc}
\leanfit{cstarMatrixEntropyVariationalFormula_of_relativeEntropy_nonneg}
\leanfit{cstarMatrixEntropyVariationalFormula_of_hermitianCfc}
\leanfit{cstarMatrixRelativeEntropyJointConvexOnStrictPositive}
\leanfit{cstarMatrixSuperoperatorEntropyKernelCfc_ratio_commute_positiveSqrtRightLift}
\leanfit{cstarMatrixSuperoperatorPerspective_outerSqrt_cfc_ratio_mul_outerSqrt}
\leanfit{cstarMatrixSuperoperatorPerspective_eq_cfc_ratio_mul_rightLift}
\leanfit{cstarMatrix_unit_inv_to_matrix}
\leanfit{cstarMatrixSuperoperatorRightLift_unit_inv_to_matrix}
\leanfit{cstarMatrixSuperoperatorPerspectiveTrace_eq_matrixSuperoperatorEntropyKernelTrace}
\leanfit{cstarMatrixRelativeEntropyPerspectiveTraceRepresentation_all}
\leanfit{cstarMatrixRelativeEntropyJointConvexOnStrictPositive_all}
\leanfit{cstarMatrixEntropyVariationalFormula}
\leanfit{liebTraceConcavityTarget_of_relativeEntropy_route}
\leanfit{liebTraceConcavityTarget_of_relativeEntropy_jointConvex}
\leanfit{liebTraceConcavityTarget_all}
\leanfit{FiniteProbability.expectationReal_trace_normed_exp_add_le}
\end{formalized}

\begin{lemma}[Algorithm 1 one-step product-law adapter]
Let \(P_s\) be the canonical independent squared-magnitude trace law for
Algorithm~1, and let \(P_1\) be the corresponding one-sample law on sampled
entries.  For any fixed trace coordinate \(t<s\), any complex-valued
one-sample statistic \(f\), and any C\(^*\)-matrix-valued one-sample statistic
\(F\),
\[
  \mathbb E_{P_s} f(\omega_t)
  =
  \mathbb E_{P_1} f,\qquad
  \mathbb E_{P_s} F(\omega_t)
  =
  \mathbb E_{P_1} F.
\]
Consequently, if \(H\) and every \(X(x)\) are self-adjoint, then
\[
  \mathbb E_{P_s}
  \operatorname{Re}\operatorname{tr}
  \exp\!\bigl(H+X(\omega_t)\bigr)
  \le
  \operatorname{Re}\operatorname{tr}
  \exp\!\left(
    H+
    \log\mathbb E_{P_1}\exp(X)
  \right).
\]
This is the one-coordinate Algorithm~1 specialization of the closed
finite-dimensional Tropp/Lieb one-step trace-MGF theorem.
\end{lemma}

\begin{formalized}
\leanfit{sqMagSampleProbability}
\leanfit{sqMagTraceProbability_expectationComplex_step_eq}
\leanfit{sqMagTraceProbability_expectationCStarMatrix_step_eq}
\leanfit{sqMagTraceProbability_expectationCStarMatrix_step_eq_sampleExpectation}
\leanfit{sqMagTraceProbability_expectationReal_trace_normed_exp_add_step_le}
\end{formalized}

\begin{theorem}[Algorithm 1 iid trace-MGF domination]
Let \(P_s\) be the \(s\)-fold independent squared-magnitude trace law and
let \(P_1\) be its one-sample marginal.  Let
\[
  X:\{1,\ldots,m\}\times\{1,\ldots,n\}
  \longrightarrow \mathbb C^{d\times d}
\]
be self-adjoint valued, and let \(H\) be self-adjoint.  Define
\[
  K=\log \mathbb E_{P_1}\exp(X).
\]
Then
\[
  \mathbb E_{P_s}
  \operatorname{Re}\operatorname{tr}
  \exp\!\left(H+\sum_{t=1}^s X(\omega_t)\right)
  \le
  \operatorname{Re}\operatorname{tr}
  \exp\!\left(H+\sum_{t=1}^s K\right).
\]
This is the independent-sum lift of the one-step Tropp/Lieb theorem for the
canonical Algorithm~1 product sampler.  The proof first factors the trace mass
after appending a last sample, then conditions on the last coordinate, applies
the one-step theorem, and iterates.
\end{theorem}

\begin{formalized}
\leanfit{sqMagTraceProbMass_snoc}
\leanfit{sqMagTraceProbability_expectationReal_succ_last_eq}
\leanfit{cstarMatrix_finset_sum_isSelfAdjoint}
\leanfit{sqMagTraceProbability_expectationReal_trace_normed_exp_add_sum_le}
\end{formalized}

\begin{corollary}[Algorithm 1 finite-real trace-MGF bridge]
Let \(H\in\mathbb R^{d\times d}\) and every one-sample increment
\(X(x)\in\mathbb R^{d\times d}\) be symmetric.  Under the same independent
Algorithm~1 trace law \(P_s\), define
\[
  K =
  \log\mathbb E_{P_1}
    \exp\!\bigl(X_{\mathbb C}\bigr),
\]
where \(X_{\mathbb C}\) denotes the complex C\(^*\)-matrix embedding of the
real matrix \(X\).  Then
\[
  \mathbb E_{P_s}
    \operatorname{tr}_{\mathbb R}
    \exp_{\mathrm{mat}}\!\left(
      H+\sum_{t=1}^s X(\omega_t)
    \right)
  \le
  \operatorname{Re}\operatorname{tr}_{\mathbb C}
    \exp\!\left(
      H_{\mathbb C}+\sum_{t=1}^s K
    \right).
\]
The bridge uses the newly formalized identity
\[
  \operatorname{Re}\operatorname{tr}_{\mathbb C}
    \exp(M_{\mathbb C})
  =
  \operatorname{tr}_{\mathbb R}\bigl(\exp_{\mathrm{mat}}(M)\bigr),
\]
proved by showing that the finite-real embedding is a continuous ring
homomorphism and therefore commutes with the power-series matrix exponential.
This is a genuine interface theorem between the C\(^*\)-matrix Lieb iteration
and the repository's finite-real trace-exponential concentration layer.  It
does not yet bound \(K\) by a scalar variance proxy and therefore is not by
itself the CACM equation~(2) Bernstein/Khintchine tail bound.
\end{corollary}

\begin{formalized}
\leanfit{finiteComplexCStarMatrixRingHom}
\leanfit{finiteComplexCStarMatrixRingHom_continuous}
\leanfit{finiteComplexCStarMatrix_finiteMatrixExp}
\leanfit{cstarMatrixTrace_normed_exp_finiteComplexCStarMatrix}
\leanfit{cstarMatrixTrace_normed_exp_finiteComplexCStarMatrix_re}
\leanfit{finiteComplexCStarMatrix_add}
\leanfit{finiteComplexCStarMatrix_finset_sum}
\leanfit{sqMagTraceProbabilityFiniteRealTraceMGFLogBound}
\leanfit{sqMagTraceProbability_expectationReal_finiteTrace_finiteMatrixExp_add_sum_le}
\end{formalized}

\begin{corollary}[Algorithm 1 residual-dilation trace-MGF target]
Let \(Z(x)\) be the one-sample Algorithm~1 residual increment and let
\[
  D(Z)=\begin{bmatrix}0&Z\\ Z^\mathsf T&0\end{bmatrix}
\]
be its self-adjoint dilation.  For every scalar \(\theta\), the finite-real
trace-MGF bridge applies to the symmetric increments \(\theta D(Z(x))\).
Consequently,
\[
  \mathbb E_{P_s}
  \operatorname{tr}_{\mathbb R}
  \exp_{\mathrm{mat}}\!\left(
    \theta D(A-\widetilde A)
  \right)
  \le
  \operatorname{Re}\operatorname{tr}_{\mathbb C}
  \exp\!\left(
    \sum_{t=1}^s
      \log\mathbb E_{P_1}
      \exp\!\bigl(\theta D(Z)_{\mathbb C}\bigr)
  \right).
\]
The proof first instantiates the finite-real trace-MGF theorem with
\(\theta D(Z(x))\), then rewrites
\(\sum_t\theta D(Z(\omega_t))=\theta D(A-\widetilde A)\) using the formal
residual decomposition.  This is the exact trace-exponential object needed by
the matrix Bernstein route.  The remaining unproved step is to bound the
logarithmic mean increment by a scalar/variance proxy.
\end{corollary}

\begin{formalized}
\lean{Algorithms/RandNLA/ElementwiseSpectral.lean}
\leanfit{rectSelfAdjointDilation_elementwiseSampleResidualIncrement_smul_symmetric}
\leanfit{sqMagTraceProbability_expectationReal_finiteTrace_finiteMatrixExp_sum_rectSelfAdjointDilation_sampleResidualIncrement_le}
\leanfit{sqMagTraceProbability_expectationReal_finiteTrace_finiteMatrixExp_rectSelfAdjointDilation_elementwiseTraceResidual_le}
\end{formalized}

\begin{corollary}[Algorithm 1 eigenvalue Markov step after trace-MGF]
With the notation of the preceding corollary, let \(T\in\mathbb R\).  The
repository now also specializes the generic trace-exponential Markov interface
to the scaled dilation residual:
\[
  \mathbb P\!\left\{
    \exists a:\,
    T\le \lambda_a\!\left(\theta D(A-\widetilde A)\right)
  \right\}
  \le
  e^{-T}\,
  \operatorname{Re}\operatorname{tr}_{\mathbb C}
  \exp\!\left(
    \sum_{t=1}^s
      \log\mathbb E_{P_1}
      \exp\!\bigl(\theta D(Z)_{\mathbb C}\bigr)
  \right).
\]
It also proves the corresponding two-sided absolute-eigenvalue event by using
the same trace-MGF theorem with \(\theta\) and \(-\theta\).  This closes the
trace-exponential-to-eigenvalue Markov dependency for the actual Algorithm 1
dilation residual.  The later source-aligned sections add the scalar
matrix-CGF/Bernstein constants and convert this dependency into the closed
equation~(2)-style theorem for the truncated squared-magnitude law.
\end{corollary}

\begin{formalized}
\lean{Algorithms/RandNLA/ElementwiseSpectral.lean}
\leanfit{sqMagTraceProbability_eventProb_exists_finiteHermitianEigenvalue_scaled_rectSelfAdjointDilation_elementwiseTraceResidual_ge_le}
\leanfit{sqMagTraceProbability_eventProb_forall_abs_finiteHermitianEigenvalue_scaled_rectSelfAdjointDilation_elementwiseTraceResidual_lt_ge}
\end{formalized}

\begin{definition}[Truncated elementwise variant]
For a threshold \(\tau\ge 0\), define the hard-thresholded matrix
\[
  \hat A_{ij} =
  \begin{cases}
    0, & |A_{ij}| < \tau,\\
    A_{ij}, & \text{otherwise}.
  \end{cases}
\]
The source-aligned truncated sketch samples from \(\hat A\) but measures the
final residual against the original matrix \(A\).
\end{definition}

\begin{theorem}[Deterministic truncation cost]
Let \(A\in\mathbb R^{n\times n}\), \(n>0\), and \(\varepsilon\ge 0\).  With
\(\tau=\varepsilon/(2n)\),
\[
  \Norm{A-\hat A} \le \Frob{A-\hat A} \le \frac{\varepsilon}{2}.
\]
Thus, if the exact sampled residual of the truncated matrix satisfies
\[
  \Norm{\hat A-\widetilde{\hat A}} \le \frac{\varepsilon}{2},
\]
then the truncated sketch is within \(\varepsilon\) of the original matrix:
\[
  \Norm{A-\widetilde{\hat A}} \le \varepsilon.
\]
\end{theorem}

\begin{corollary}[Exact truncated transfer]
Any high-probability half-budget theorem for sampling the truncated matrix
\(\hat A\) transfers immediately to the original-matrix event
\(\Norm{A-\widetilde{\hat A}}\le\varepsilon\).  This is the exact
source-alignment layer.  The rounded endpoint is not stated here with a
generic perturbation matrix; the final Algorithm~1 corollary below expands the
internal floating-point error into the explicit
\(\gamma_{s+1}\)-radius.
\end{corollary}

\begin{theorem}[Support and bounded increments for the truncated route]
Assume \(\tau>0\), \(s>0\), and
\(\|\hat A\|_F^2>0\).  Under the independent product law that samples entries
with probabilities \(\hat A_{ij}^2/\|\hat A\|_F^2\), the event that every
sampled entry has positive probability has probability one.  On that event,
each one-step residual increment
\[
  Z_X=\frac{\hat A}{s}-C_X
\]
satisfies
\[
  \Norm{Z_X}\le \Frob{Z_X}
  \le
  \frac{1}{s}\Frob{\hat A}
  +\frac{\Frob{\hat A}^2}{s\tau}.
\]
Equivalently, its self-adjoint dilation \(D(Z_X)\) satisfies the finite
operator bound
\[
  \Norm{D(Z_X)}
  \le
  \sqrt{2}\left(
    \frac{1}{s}\Frob{\hat A}
    +\frac{\Frob{\hat A}^2}{s\tau}
  \right),
\]
and therefore the squared-Loewner bound
\[
  D(Z_X)^2
  \preceq
  2\left(
    \frac{1}{s}\Frob{\hat A}
    +\frac{\Frob{\hat A}^2}{s\tau}
  \right)^2 I.
\]
The operator bound is also exposed in the one-sided and lower Loewner forms
\[
  D(Z_X)\preceq L I,
  \qquad
  -D(Z_X)\preceq L I,
  \qquad
  L=\sqrt{2}\left(
    \frac{1}{s}\Frob{\hat A}
    +\frac{\Frob{\hat A}^2}{s\tau}
  \right).
\]
The corresponding ``all increments are bounded'', ``all increments are
two-sided Loewner-bounded'', and ``all squared increments are Loewner-bounded''
trace events have probability one.  Their intersection, the simultaneous
Bernstein-boundedness event needed by a future matrix Bernstein theorem, also
has probability one.
\end{theorem}

\begin{corollary}[Support-aware truncated finite-test MGF]
Let \(\mathcal N\) be a finite family of deterministic test vectors
\(z_a\in\mathbb R^{m+n}\), and set
\[
  L_\tau
  =
  \sqrt{2}\left(
    \frac{1}{s}\Frob{\hat A}
    +\frac{\Frob{\hat A}^2}{s\tau}
  \right).
\]
For every positive-probability sampled entry of the truncated law,
\[
  z_a^T D(Z_X)z_a
  \le L_\tau\|z_a\|_2^2.
\]
Therefore, for any \(\lambda_a>0\),
\[
  \mathbb P\!\left[
    z_a^T D(R_{\hat A})z_a\le T_a
    \text{ for all }a\in\mathcal N
  \right]
  \ge
  1-\sum_{a\in\mathcal N}
    \exp\!\left(s\lambda_a L_\tau\|z_a\|_2^2-\lambda_aT_a\right).
\]
This is a support-aware finite-test scalar tail.  It removes the hidden
``the sampled entry is retained'' side condition from finite test-vector
arguments, but it is not the matrix Bernstein half-budget event.
\end{corollary}


\begin{formalized}
\leanfit{elementwiseTruncate}
\leanfit{elementwiseTruncatedTraceResidual}
\leanfit{fl_elementwiseTruncatedTraceResidual}
\leanfit{elementwiseTruncate_error_entry_abs_le}
\leanfit{elementwiseTruncate_abs_le}
\leanfit{frobNormSqRect_elementwiseTruncate_le}
\leanfit{frobNormRect_elementwiseTruncate_le}
\leanfit{elementwiseTruncate_nonzero_abs_ge}
\leanfit{frobNormRect_const_square}
\leanfit{elementwiseTruncate_square_error_frobNormRect_le_half}
\leanfit{elementwiseTruncate_square_error_rectOpNorm2Le_half}
\leanfit{elementwiseTruncatedTraceResidual_rectOpNorm2Le_of_truncated}
\leanfit{elementwiseTruncatedTraceResidual_square_rectOpNorm2Le_of_half}
\leanfit{algorithm1ExactTruncatedSpectralEvent}
\leanfit{algorithm1ExactSpectralEvent_truncated_half_subset_original}
\leanfit{probability_algorithm1_exact_truncated_spectral_of_sampled_half}
\leanfit{fl_elementwiseTruncatedTraceResidual_rectOpNorm2Le_of_truncated}
\leanfit{algorithm1FlTruncatedSpectralEvent}
\leanfit{algorithm1ExactSpectralEvent_truncated_half_subset_fl_original}
\leanfit{probability_algorithm1_fl_truncated_spectral_of_sampled_half}
\leanfit{entry_ne_zero_of_sqMagProb_pos}
\leanfit{FiniteProbability.eventProb_le_one}
\leanfit{elementwiseTracePositiveProb}
\leanfit{sqMagTraceProbability_eventProb_elementwiseTracePositiveProb}
\leanfit{frobNormRect_elementwiseSampleContribution_truncated_le}
\leanfit{frobNormRect_elementwiseSampleResidualIncrement_truncated_le}
\leanfit{rectOpNorm2Le_elementwiseSampleResidualIncrement_truncated}
\leanfit{sqMagTraceProbability_eventProb_truncatedResidualIncrementsBoundedEvent_eq_one}
\leanfit{finiteOpNorm2Le_rectSelfAdjointDilation_elementwiseSampleResidualIncrement_truncated}
\leanfit{finiteQuadraticForm_rectSelfAdjointDilation_elementwiseSampleResidualIncrement_truncated_le}
\leanfit{sqMagTraceProbability_eventProb_truncatedDilationIncrementsBoundedEvent_eq_one}
\leanfit{finiteLoewnerLe_rectSelfAdjointDilation_elementwiseSampleResidualIncrement_truncated}
\leanfit{finiteLoewnerLe_neg_rectSelfAdjointDilation_elementwiseSampleResidualIncrement_truncated}
\leanfit{truncatedDilationIncrementLoewnerBoundedEvent}
\leanfit{sqMagTraceProbability_eventProb_truncatedDilationIncrementLoewnerBoundedEvent_eq_one}
\leanfit{truncatedDilationIncrementLoewnerBoundedEvent_subset_exactDilationUpperEvent_sum_bound}
\leanfit{sqMagTraceProbability_eventProb_algorithm1ExactDilationUpperEvent_truncated_sum_bound_eq_one}
\leanfit{sqMagTraceProbability_eventProb_algorithm1ExactSpectralEvent_truncated_sum_bound_eq_one}
\leanfit{finiteLoewnerLe_rectSelfAdjointDilation_square_elementwiseSampleResidualIncrement_truncated}
\leanfit{truncatedDilationIncrementSquaresBoundedEvent}
\leanfit{sqMagTraceProbability_eventProb_truncatedDilationIncrementSquaresBoundedEvent_eq_one}
\leanfit{FiniteProbability.eventProb_inter_eq_one_of_eq_one}
\leanfit{truncatedDilationBernsteinBoundedEvent}
\leanfit{sqMagTraceProbability_eventProb_truncatedDilationBernsteinBoundedEvent_eq_one}
\leanfit{sqMagTraceProbability_eventProb_forall_finiteQuadraticForm_rectSelfAdjointDilation_truncatedTraceResidual_le_ge_one_sub_sum_exp_of_support_bound}
\end{formalized}

\begin{theorem}[Accumulated bounded-increment spectral bound]
Let \(\hat A=T_\tau(A)\), and suppose each retained truncated self-adjoint
dilation increment \(X_t\) satisfies the one-sided Loewner bound
\[
  X_t\preceq L I .
\]
Then the full exact truncated residual dilation satisfies
\[
  \sum_{t=1}^{s} X_t \preceq sL\, I .
\]
Consequently, under the truncated squared-magnitude product law, the exact
truncated sampled residual satisfies
\[
  \Norm{\hat A-\widetilde{\hat A}}\le sL
\]
with probability \(1\), for the explicit retained-increment bound
\[
  L=\sqrt2\left(\frac{\Frob{\hat A}}{s}
      +\frac{\Frob{\hat A}^2}{s\tau}\right).
\]
\end{theorem}


\begin{formalized}
\leanfit{truncatedDilationIncrementLoewnerBoundedEvent_subset_exactDilationUpperEvent_sum_bound}
\leanfit{sqMagTraceProbability_eventProb_algorithm1ExactDilationUpperEvent_truncated_sum_bound_eq_one}
\leanfit{sqMagTraceProbability_eventProb_algorithm1ExactSpectralEvent_truncated_sum_bound_eq_one}
\end{formalized}

\begin{theorem}[One-step residual variance proxies]
Let \(X\) be one Algorithm 1 sampled entry drawn from the squared-magnitude
law, and define the one-step residual increment
\[
  Z_X = \frac{A}{s}-C_X,
\]
where \(C_X\) is the one-sample contribution matrix.  Then, for every entry
\((i,j)\),
\[
  \mathbb E (Z_X)_{ij}^2
  \le
  \frac{\Frob{A}^2}{s^2}.
\]
Consequently,
\[
  \mathbb E\Frob{Z_X}^2
  \le
  mn\,\frac{\Frob{A}^2}{s^2},
  \qquad
  \mathbb E\Norm{Z_Xx}^2
  \le
  mn\,\frac{\Frob{A}^2}{s^2}\Norm{x}^2 .
\]
\end{theorem}

\begin{corollary}[Fixed-vector residual tail]
For the full \(s\)-sample residual \(R=A-\Atil\), every fixed vector \(x\), and
every \(\eta>0\),
\[
  \mathbb P\!\left[\Norm{Rx}\le\eta\right]
  \ge
  1-
  \frac{(mn/s)\Frob{A}^2\Norm{x}^2}{\eta^2}.
\]
This is a fixed-vector Markov theorem.  It is not uniform over all vectors and
therefore is not the CACM equation~(2) spectral-norm concentration theorem.
\end{corollary}

\begin{corollary}[Finite test set residual tail]
Let \(\mathcal N\) be any finite family of test vectors, with thresholds
\(\eta_x>0\).  Then the same proof plus the finite union bound gives
\[
  \mathbb P\!\left[
    \Norm{Rx}\le \eta_x\ \text{for every }x\in\mathcal N
  \right]
  \ge
  1-\sum_{x\in\mathcal N}
  \frac{(mn/s)\Frob{A}^2\Norm{x}^2}{\eta_x^2}.
\]
This is a future net-argument support theorem; it does not assert that
\(\mathcal N\) covers the unit sphere.
\end{corollary}

\begin{corollary}[Finite-cover operator support]
Suppose now that the finite family \(\mathcal N\) is supplied together with a
radius \(\rho\) such that every vector in the unit ball is within Euclidean
distance \(\rho\) of some vector in \(\mathcal N\).  Let \(\eta>0\) and
\(L>0\).  Combining the finite-test theorem with the Frobenius residual event
gives
\[
  \mathbb P\!\left[
    \Norm{Rx}\le(\eta+L\rho)\Norm{x}\ \text{for every }x
  \right]
  \ge
  1-
  \sum_{x\in\mathcal N}
    \frac{(mn/s)\Frob{A}^2\Norm{x}^2}{\eta^2}
  -
    \frac{(mn/s)\Frob{A}^2}{L^2}.
\]
This is a useful covering-route theorem.  It assumes the finite cover and does
not itself construct one, so it is not the route used for the closed
source-aligned equation~(2) theorem below.
\end{corollary}

\begin{lemma}[Product-grid cover]
Let \(G=\{g_a:a\in\alpha\}\subset\mathbb R\) be a finite one-dimensional grid
such that for every \(t\in[-1,1]\) there is \(a\in\alpha\) with
\[
  |t-g_a|\le\delta .
\]
Assume \(\delta\ge0\).  For dimension \(n\), define the product grid
\[
  \mathcal G_n
  =
  \{(g_{a_1},\ldots,g_{a_n}): (a_1,\ldots,a_n)\in \alpha^n\}.
\]
Then every \(x\in\mathbb R^n\) with \(\|x\|_2\le1\) is within Euclidean
distance \(\sqrt n\,\delta\) of some point of \(\mathcal G_n\).  Moreover the
formal index type has cardinality
\[
  |\alpha^n|=|\alpha|^n .
\]
This closes the product-cover bookkeeping dependency for a covering-net route.
It is not a concentration theorem and does not supply the sharp scalar tail
bounds needed for equation~(2).
\end{lemma}


\begin{formalized}
\leanfit{abs_coord_le_vecNorm2}
\leanfit{realUnitIntervalCover}
\leanfit{rectUnitBallCover_product_grid}
\leanfit{fintype_card_product_grid_index}
\end{formalized}

\begin{formalized}
\leanfit{vecNorm2_smul}
\leanfit{rectMatMulVec_smul}
\leanfit{rectMatMulVec_add}
\leanfit{rectMatMulVec_sub}
\leanfit{rectUnitBallCover}
\leanfit{rectOpNorm2Le_of_unit_ball_bound}
\leanfit{rectOpNorm2Le_of_unit_ball_cover}
\leanfit{sqMagProb_sum_elementwiseSampleContribution_entry_sq_le}
\leanfit{sqMagProb_sum_elementwiseSampleResidualIncrement_entry_sq_le}
\leanfit{sqMagProb_sum_elementwiseSampleResidualIncrement_frob_sq_le}
\leanfit{sqMagProb_sum_vecNorm2Sq_rectMatMulVec_elementwiseSampleResidualIncrement_le}
\leanfit{sqMagTraceProbability_expectationReal_vecNorm2Sq_rectMatMulVec_elementwiseTraceResidual_le}
\leanfit{sqMagTraceProbability_eventProb_vecNorm2_rectMatMulVec_elementwiseTraceResidual_le_ge_one_sub}
\leanfit{fl_elementwiseTraceResidual_vecNorm2_le_of_exact_fixed_vector}
\leanfit{sqMagTraceProbability_eventProb_fl_vecNorm2_rectMatMulVec_elementwiseTraceResidual_le_ge_one_sub}
\leanfit{sqMagTraceProbability_eventProb_forall_vecNorm2_rectMatMulVec_elementwiseTraceResidual_le_ge_one_sub_sum}
\leanfit{sqMagTraceProbability_eventProb_forall_fl_vecNorm2_rectMatMulVec_elementwiseTraceResidual_le_ge_one_sub_sum}
\leanfit{sqMagTraceProbability_eventProb_algorithm1ExactSpectralEvent_ge_one_sub_cover}
\leanfit{sqMagTraceProbability_eventProb_algorithm1FlSpectralEvent_ge_one_sub_cover}
\end{formalized}

\begin{theorem}[Exact Frobenius residual moment]
Under the canonical independent squared-magnitude product trace law for
Algorithm 1, with \(s>0\),
\[
  \mathbb E\Frob{A-\Atil}^2
  \le
  \frac{mn}{s}\Frob{A}^2 .
\]
Equivalently, for every threshold \(\eta>0\),
\[
  \mathbb P\!\left[\Frob{A-\Atil}\le \eta\right]
  \ge
  1-\frac{(mn/s)\Frob{A}^2}{\eta^2}.
\]
\end{theorem}


\begin{formalized}
\leanfit{sqMagTraceProbMass_marginal_two_ne}
\leanfit{sqMagTraceProbability_eventProb_sampleHits_pair_ne}
\leanfit{hitCountPairwiseCenteredMoment_le_steps_mul}
\leanfit{sqMagTraceProbability_expectationReal_elementwiseTraceResidual_entry_sq_le}
\leanfit{sqMagTraceProbability_expectationReal_elementwiseTraceResidual_frob_sq_le}
\leanfit{sqMagTraceProbability_eventProb_algorithm1ExactFrobEvent_ge_one_sub}
\end{formalized}

\begin{corollary}[Frobenius residual bridge]
If the exact Algorithm 1 residual satisfies
\[
  \Frob{R}\le \varepsilon,
\]
then it also satisfies the rectangular vector-action operator event.  The
generic floating-point adapter for this weaker route remains a Lean support
lemma, but the PDF does not advertise it as an Algorithm~1 endpoint because its
radius is phrased through an internal perturbation matrix.
\end{corollary}

\begin{formalized}
\leanfit{algorithm1ExactFrobEvent}
\leanfit{algorithm1ExactFrobEvent_subset_exactSpectralEvent}
\leanfit{probability_algorithm1_exact_spectral_of_frob}
\end{formalized}

\begin{corollary}[Nonconditional Frobenius-to-operator probability]
Combining the exact Frobenius residual theorem with the deterministic bridge,
the canonical product trace law satisfies
\[
  \mathbb P\!\left[
    \Norm{(A-\Atil)x}\le \eta\Norm{x}\ \text{for all }x
  \right]
  \ge
  1-\frac{(mn/s)\Frob{A}^2}{\eta^2}.
\]
\end{corollary}

\begin{formalized}
\leanfit{sqMagTraceProbability_eventProb_algorithm1ExactSpectralEvent_ge_one_sub_frob}
\end{formalized}

\begin{theorem}[Exact entrywise union-bound fallback, not equation (2)]
Let \(\tau>0\).  Under the same canonical product trace law, each residual
entry satisfies
\[
  \mathbb P\!\left[
    |A_{ij}-\Atil_{ij}|\le \tau
  \right]
  \ge
  1-\frac{\Frob{A}^2}{s\tau^2}.
\]
Consequently, by a finite union bound over all \(mn\) entries,
\[
  \mathbb P\!\left[
    |A_{ij}-\Atil_{ij}|\le \tau
    \quad\text{for all } i,j
  \right]
  \ge
  1-mn\,\frac{\Frob{A}^2}{s\tau^2}.
\]
On this event,
\[
  \Norm{(A-\Atil)x}
  \le
  \Frob{\tau {\bf 1}}\,\Norm{x}
  \quad\text{for all }x,
\]
where \(\tau{\bf 1}\) denotes the \(m\times n\) matrix whose entries are all
\(\tau\).  This is an exact-arithmetic fallback theorem.  It is deliberately
not the Algorithm~1 equation~(2) theorem, and the PDF does not advertise the
internal floating-point adapter for this weaker route because that adapter is
phrased through a generic perturbation budget.
\end{theorem}


\begin{formalized}
\leanfit{FiniteProbability.eventProb_finset_forall_ge_one_sub_sum}
\leanfit{FiniteProbability.eventProb_forall_ge_one_sub_sum}
\leanfit{algorithm1ExactEntrywiseEvent}
\leanfit{sqMagTraceProbability_eventProb_elementwiseTraceResidual_entry_abs_le_ge_one_sub}
\leanfit{sqMagTraceProbability_eventProb_algorithm1ExactEntrywiseEvent_ge_one_sub}
\leanfit{algorithm1ExactEntrywiseEvent_subset_exactSpectralEvent_const}
\leanfit{probability_algorithm1_exact_spectral_of_entrywise_const}
\leanfit{sqMagTraceProbability_eventProb_algorithm1ExactSpectralEvent_ge_one_sub_entrywise}
\end{formalized}

\begin{remark}
This section is not the proof of equation~(2) itself.  The Frobenius and
entrywise probability theorems are nonconditional Markov or union-bound routes
with weaker dimension factors.  The sharper source-aligned
Bernstein/Khintchine-type theorem is proved later by the trace-MGF route.
\end{remark}

\begin{corollary}[Truncated element-wise sampler with explicit FP radius]
Let \(A\in\R^{n\times n}\), \(n>0\), \(s>0\), \(\varepsilon>0\), and
\(0<\delta\le 1\).  Set
\[
  \tau=\frac{\varepsilon}{2n},
  \qquad
  \widehat A=\operatorname{trunc}_{\tau}(A),
\]
where the truncation is the hard-thresholding map
\[
  \widehat A_{ij}
  =
  \begin{cases}
    0, & |A_{ij}|<\tau,\\
    A_{ij}, & |A_{ij}|\ge\tau .
  \end{cases}
\]
The value \(\tau=\varepsilon/(2n)\) is chosen because \(A\) is \(n\times n\):
every discarded entry has magnitude at most \(\tau\), hence
\[
  \Frob{A-\widehat A}
  \le \sqrt{n^2\tau^2}
  = n\tau
  = \frac{\varepsilon}{2}.
\]
Thus deterministic truncation uses one half of the final
\(\varepsilon\)-radius, leaving the other half for the sampled residual of
\(\widehat A\).  The same choice also gives
\(|\widehat A_{ij}|\ge\tau\) for every retained nonzero sampled entry, which
is what turns the internal division-by-\(\widehat A_{ij}\) stability budget
into the explicit floating-point term below.
Finally, assume the squared-magnitude denominator of \(\widehat A\) is positive.
Assume the source sample budget
\[
  14n\Frob{A}^2
  \log\!\left(\frac{2(2n)}{\delta}\right)
  \le s\varepsilon^2 .
\]
Then exact source-aligned truncated element-wise sampling from
\(\widehat A_{ij}^2/\Frob{\widehat A}^2\) satisfies
\[
  \mathbb P\!\left[
    \left\|A-\widetilde A_{\tau}\right\|_2
    \le \varepsilon
  \right]\ge 1-\delta .
\]
If additionally \(\gamma_s\) and \(\gamma_{s+1}\) are valid for the
floating-point model, then the rounded sketch satisfies
\[
  \mathbb P\!\left[
    \left\|A-\operatorname{fl}(\widetilde A_{\tau})\right\|_2
    \le
    \varepsilon+
    \frac{2n^2\Frob{A}^2}{\varepsilon}\,\gamma_{s+1}
  \right]\ge 1-\delta ,
\]
where
\[
  \gamma_{s+1}=\frac{(s+1)u}{1-(s+1)u}.
\]
This is a theorem for the hard-thresholded element-wise sampler above.  It is
not the literal CACM Algorithm~1 sampler with distribution
\(p_{ij}=A_{ij}^2/\Frob{A}^2\), because the sampling law has been changed to
use \(\widehat A\).  The term after \(\varepsilon\) is the explicit
floating-point contribution; no budget matrix is hidden in the statement.
\end{corollary}

\begin{formalized}
\leanfit{sqMagTraceProbability_eventProb_algorithm1ExactSpectralEvent_truncatedTraceResidual_ge_one_sub_delta_source_sample_budget_sharp_square}
\leanfit{sqMagTraceProbability_eventProb_algorithm1ExactTruncatedSpectralEvent_ge_one_sub_delta_source_sample_budget_sharp_square}
\leanfit{sqMagTraceProbability_eventProb_algorithm1FlTruncatedSpectralRadius_ge_one_sub_delta_source_sample_budget_explicit_gamma_square}
\end{formalized}

\begin{remark}[Exact probability convention for Corollary 143]
This corollary is an exact-law result: the sampler draws from the truncated
squared-magnitude law
\(\widehat A_{ij}^2/\Frob{\widehat A}^2\).  Under the current convention no
floating-point error is charged for computing that law.  The FP radius in the
formal theorem charges only the rounded sketch arithmetic after the sample has
been drawn.
\end{remark}

\begin{corollary}[Literal Algorithm 1 without truncation]
Let \(A\in\R^{n\times n}\), let \(s>0\), let \(\varepsilon>0\), and assume
\[
  \Frob{A}^2>0.
\]
Run Algorithm~1 with the literal squared-magnitude sampling distribution
\[
  p_{ij}=\frac{A_{ij}^2}{\Frob{A}^2}.
\]
Assume the Frobenius/Markov sample budget
\[
  \frac{n^2\Frob{A}^2}{s\varepsilon^2}\le \delta .
\]
For the floating-point part, assume there is a displayed nonzero-entry floor
\(\alpha>0\) such that
\[
  A_{ij}\ne0 \quad\Longrightarrow\quad \alpha\le |A_{ij}|
  \qquad\text{for all }i,j.
\]
This is not a truncation: it is only a stated condition on the actual input
matrix \(A\), used to make the division-by-\(A_{ij}\) floating-point term
readable.  If \(\gamma_{s+1}\) is valid for the floating-point model, then
\[
  \mathbb P\!\left[
    \left\|A-\operatorname{fl}(\widetilde A)\right\|_2
    \le
    \varepsilon+
    n\,\frac{\Frob{A}^2}{\alpha}\,\gamma_{s+1}
  \right]\ge 1-\delta ,
\]
where
\[
  \gamma_{s+1}=\frac{(s+1)u}{1-(s+1)u}.
\]
The probability is taken under the literal Algorithm~1 trace law above.  The
proof intersects the exact Frobenius/Markov spectral event with the
probability-one support event for \(p_{ij}\).  Hence zero entries of \(A\) are
never sampled on the event where the local FP division theorem is used, and no
hard-thresholded matrix \(\widehat A\) appears.
\end{corollary}

\begin{formalized}
\leanfit{sqMagTraceErrorBudget_zero_init_le_const_of_entry_abs_ge}
\leanfit{frobNormRect_sqMagTraceErrorBudget_zero_init_le_const_square_of_entry_abs_ge}
\leanfit{sqMagTraceProbability_eventProb_algorithm1FlSpectralRadius_ge_one_sub_delta_frob_explicit_gamma_square}
\end{formalized}

\begin{corollary}[Nonconditional literal Algorithm 1 support-radius theorem]
Let \(A\in\R^{m\times n}\), \(s>0\), \(\Frob{A}^2>0\), and
\(0<\delta\le1\).  Run Algorithm~1 with the literal exact squared-magnitude
law
\[
  p_{ij}=\frac{A_{ij}^2}{\Frob{A}^2}.
\]
The probability law is exact by convention.  The computed non-probability
quantities are the sampled-entry rescaling, the additions into the sketch, and
the final rounded residual; their floating-point contribution is charged
below.  The threshold quantities in the display are exact analysis parameters,
not additional algorithm outputs.  If an implementation also computes these
thresholds for a stopping or certification rule, that separate computation is
not part of this theorem.

Define the exact reciprocal-entry quantity
\[
  G(A)=
  \sum_{\{(i,j):A_{ij}\ne0\}}
    \frac{\Frob{A}^2}{|A_{ij}|},
  \qquad
  H(A)=\Frob{A}+G(A).
\]
For a chosen analysis parameter \(\theta>0\), define
\[
  L_{\rm lit}=\frac{H(A)}{s},
  \qquad
  V_{\rm lit}=\max\{m,n\}\frac{\Frob{A}^2}{s^2},
\]
and
\[
  \beta_{\rm lit}(\theta)
  =
  \frac{\exp(\theta L_{\rm lit})-\theta L_{\rm lit}-1}
       {L_{\rm lit}^2}.
\]
The exact spectral radius appearing in the formal theorem is
\[
  \varepsilon_{\rm lit}(\theta,\delta)
  =
  \frac{
    \log\!\left(
      \frac{2(m+n)\exp(s\beta_{\rm lit}(\theta)V_{\rm lit})}{\delta}
    \right)}
  {\theta}.
\]
Substituting all definitions gives the irreducible expanded form
\[
  \varepsilon_{\rm lit}(\theta,\delta)
  =
  \frac{\log(2(m+n)/\delta)}{\theta}
  +
  \frac{\max\{m,n\}\Frob{A}^2\,s}{\theta H(A)^2}
  \left(
    \exp\!\left(\frac{\theta H(A)}{s}\right)
    -\frac{\theta H(A)}{s}
    -1
  \right).
\]
If \(\gamma_s\) and \(\gamma_{s+1}\) are valid for the floating-point model,
then the rounded literal Algorithm~1 sketch satisfies
\[
  \Prob\!\left[
    \left\|A-\fl(\widetilde A)\right\|_2
    \le
    T_{\rm lit}^{\rm fp}(\theta,\delta)
  \right]\ge 1-\delta ,
\]
where the fully expanded final radius is
\[
  T_{\rm lit}^{\rm fp}(\theta,\delta)
  =
  \frac{\log(2(m+n)/\delta)}{\theta}
  +
  \frac{\max\{m,n\}\Frob{A}^2\,s}{\theta H(A)^2}
  \left(
    \exp\!\left(\frac{\theta H(A)}{s}\right)
    -\frac{\theta H(A)}{s}
    -1
  \right)
  +
  \sqrt{mn}\,\gamma_{s+1}\,G(A),
\]
with
\[
  \gamma_{s+1}=\frac{(s+1)u}{1-(s+1)u}.
\]
Thus, in the regime where \(\theta>0\) is fixed and
\(\theta H(A)/s=O(1)\), the same bound has order
\[
  T_{\rm lit}^{\rm fp}(\theta,\delta)
  =
  \Theta\!\left(
    \frac{\log((m+n)/\delta)}{\theta}
    +
    \frac{\max\{m,n\}\Frob{A}^2\,s}{\theta H(A)^2}
    \left(
      \exp\!\left(\frac{\theta H(A)}{s}\right)
      -\frac{\theta H(A)}{s}
      -1
    \right)
    +
    \sqrt{mn}\,\gamma_{s+1}G(A)
  \right).
\]
In the small-argument subregime \(\theta H(A)/s\to0\), this simplifies to
\[
  T_{\rm lit}^{\rm fp}(\theta,\delta)
  =
  \Theta\!\left(
    \frac{\log((m+n)/\delta)}{\theta}
    +
    \frac{\theta\,\max\{m,n\}\Frob{A}^2}{s}
    +
    \sqrt{mn}\,\gamma_{s+1}G(A)
  \right).
\]
For fixed \(A,m,n,s,\theta,\delta\), the floating-point part tends to zero as
\(\gamma_{s+1}\to0\).  The theorem is nonconditional and literal, but it is
not the source-uniform CACM equation~(2) rate: tiny nonzero entries enlarge
\(G(A)\), exactly as the formal support obstruction predicts.

The formalization now also removes the free parameter \(\theta\) by applying
the proved Bennett optimizer.  Let
\[
  q_\delta=\log\!\left(\frac{2(m+n)}{\delta}\right),
  \qquad
  W_{\rm lit}=sV_{\rm lit}=\frac{\max\{m,n\}\Frob{A}^2}{s}.
\]
For any target exact radius \(r>0\), if
\[
  q_\delta
  \le
  \frac{r^2}{2W_{\rm lit}+\frac23L_{\rm lit}r},
\]
equivalently, after substituting \(L_{\rm lit}=H(A)/s\),
\[
  s
  \ge
  \frac{q_\delta\left(2\max\{m,n\}\Frob{A}^2+\frac23H(A)r\right)}{r^2},
\]
then the exact literal sketch satisfies \(\|A-\widetilde A\|_2\le r\) with
probability at least \(1-\delta\).  The rounded implementation-facing version
has the fully expanded scalar radius
\[
  T_{\rm lit,B}^{\rm fp}(r)
  =
  r+\sqrt{mn}\,\gamma_{s+1}G(A).
\]
Thus the direct Bennett sample requirement has order
\[
  s=
  \Theta\!\left(
    \frac{\left(\max\{m,n\}\Frob{A}^2+H(A)r\right)
      \log((m+n)/\delta)}{r^2}
  \right),
\]
and the final FP radius has order
\[
  T_{\rm lit,B}^{\rm fp}(r)
  =
  \Theta\!\left(r+\sqrt{mn}\,\gamma_{s+1}G(A)\right).
\]
This form is nonvacuous for fixed input and target radius as
\(\gamma_{s+1}\to0\), and the sample condition is nonvacuous whenever the
displayed lower bound on \(s\) can be met.  The dependence on \(G(A)\) and
\(H(A)\) is deliberate and irreducible for this support-radius route.

The formalization also gives a simpler floor-bound surface for readers who
prefer not to inspect \(G(A)\) directly.  Suppose that for some
\(\alpha>0\),
\[
  A_{ij}\ne0\quad\Longrightarrow\quad \alpha\le |A_{ij}|
  \qquad\text{for all }i,j .
\]
Then the exact deterministic simplifications
\[
  G(A)\le \frac{mn\,\Frob{A}^2}{\alpha},
  \qquad
  H(A)\le
  H_\alpha(A):=\Frob{A}+\frac{mn\,\Frob{A}^2}{\alpha}
\]
are proved in Lean.  Consequently, the stronger but more readable budget
\[
  q_\delta
  \le
  \frac{r^2}
       {2W_{\rm lit}+\frac23\,\frac{H_\alpha(A)}{s}\,r}
\]
implies the exact event
\[
  \Prob\!\left[
    \|A-\widetilde A\|_2\le r
  \right]\ge1-\delta ,
\]
and the rounded implementation-facing event
\[
  \Prob\!\left[
    \|A-\fl(\widetilde A)\|_2
    \le
    r+\sqrt{mn}\,
      \frac{mn\,\Frob{A}^2}{\alpha}\,
      \gamma_{s+1}
  \right]\ge1-\delta .
\]
Equivalently, the alpha-floor sample budget has order
\[
  s=
  \Theta\!\left(
    \frac{
      \left(\max\{m,n\}\Frob{A}^2+H_\alpha(A)r\right)
      \log((m+n)/\delta)}
      {r^2}
  \right),
\]
and the displayed rounded radius has order
\[
  \Theta\!\left(
    r+\sqrt{mn}\,
      \frac{mn\,\Frob{A}^2}{\alpha}\,
      \gamma_{s+1}
  \right).
\]
For fixed \(A,m,n,s,r,\alpha,\delta\), the floating-point additive term tends
to zero as \(\gamma_{s+1}\to0\).  This entry-floor corollary is still not
source-uniform, but it is often easier to interpret than the irreducible
reciprocal-entry formula.
\end{corollary}

\begin{formalized}
\leanfit{elementwiseLiteralContributionRadius}
\leanfit{elementwiseLiteralResidualSupportRadius}
\leanfit{elementwiseLiteralContributionRadius_le_of_entry_abs_ge}
\leanfit{elementwiseLiteralResidualSupportRadius_le_of_entry_abs_ge}
\leanfit{smul_elementwiseLiteralContributionRadius_le_of_entry_abs_ge}
\leanfit{sqMagTraceProbability_eventProb_algorithm1ExactSpectralEvent_literalTraceResidual_ge_one_sub_delta_scaled_radius_supportRadius}
\leanfit{sqMagTraceProbability_eventProb_algorithm1FlSpectralEvent_literal_scaled_radius_gamma_supportRadius}
\leanfit{sqMagTraceErrorBudget_zero_init_le_literalContributionRadius}
\leanfit{frobNormRect_sqMagTraceErrorBudget_zero_init_le_literalContributionRadius}
\leanfit{sqMagTraceProbability_eventProb_algorithm1FlSpectralRadius_literal_scaled_radius_gamma_supportRadius}
\leanfit{sqMagTraceProbability_eventProb_algorithm1ExactSpectralEvent_literalTraceResidual_ge_one_sub_delta_bennett_radius_supportRadius}
\leanfit{sqMagTraceProbability_eventProb_algorithm1ExactSpectralEvent_literalTraceResidual_ge_one_sub_delta_bernstein_denominator_two_thirds_supportRadius}
\leanfit{sqMagTraceProbability_eventProb_algorithm1FlSpectralRadius_literal_ge_one_sub_delta_bernstein_denominator_gamma_supportRadius}
\leanfit{algorithm1LiteralBernsteinDenominatorBudget_of_entry_floor}
\leanfit{sqMagTraceProbability_eventProb_algorithm1ExactSpectralEvent_literalTraceResidual_ge_one_sub_delta_bernstein_denominator_entry_floor}
\leanfit{sqMagTraceProbability_eventProb_algorithm1FlSpectralRadius_literal_ge_one_sub_delta_bernstein_denominator_gamma_entry_floor}
\end{formalized}

\begin{corollary}[Literal Algorithm 1 at the source spectral rate]
Let \(A\in\R^{n\times n}\), \(n>0\), \(s>0\), \(\varepsilon>0\),
\(0<\delta\le1\), and \(\Frob{A}^2>0\).  Run Algorithm~1 with the literal
distribution
\[
  p_{ij}=\frac{A_{ij}^2}{\Frob{A}^2}.
\]
Assume the source-rate sample budget
\[
  14n\Frob{A}^2
  \log\!\left(\frac{2(2n)}{\delta}\right)
  \le s\varepsilon^2 .
\]
Assume also that the source threshold would not remove any nonzero entry:
\[
  A_{ij}\ne0
  \quad\Longrightarrow\quad
  \frac{\varepsilon}{2n}\le |A_{ij}|
  \qquad\text{for all } i,j .
\]
Equivalently, \(\operatorname{trunc}_{\varepsilon/(2n)}(A)=A\).  Therefore
the theorem below is a literal Algorithm~1 statement; the sampling distribution
is not changed.

Then the exact sketch satisfies
\[
  \mathbb P\!\left[
    \left\|A-\widetilde A\right\|_2\le \varepsilon
  \right]\ge 1-\delta .
\]
If \(\gamma_s\) and \(\gamma_{s+1}\) are valid for the floating-point model,
then the rounded sketch satisfies
\[
  \mathbb P\!\left[
    \left\|A-\operatorname{fl}(\widetilde A)\right\|_2
    \le
    \varepsilon+
    \frac{2n^2\Frob{A}^2}{\varepsilon}\,\gamma_{s+1}
  \right]\ge 1-\delta ,
\]
where
\[
  \gamma_{s+1}=\frac{(s+1)u}{1-(s+1)u}.
\]
Thus, when \(\varepsilon=\eta\Frob{A}\), the sampling condition has the
paper-style size \(s\gtrsim n\log(n/\delta)/\eta^2\), plus the displayed
no-small-entry condition needed to make the source truncation identity and the
floating-point denominator bound explicit.

The same specialization is formalized for rectangular
\(A\in\R^{m\times n}\).  Define
\[
  R=\sqrt{mn},\qquad
  M=\max\{m,n\},\qquad
  C_{\rm rect}=2M+\frac{4\sqrt2}{3}R .
\]
Assume \(m n>0\), \(\Frob{A}^2>0\), and
\[
  4C_{\rm rect}\Frob{A}^2
  \log\!\left(\frac{2(m+n)}{\delta}\right)
  \le s\varepsilon^2 .
\]
Assume also the rectangular source threshold is the identity on the literal
input:
\[
  A_{ij}\ne0
  \quad\Longrightarrow\quad
  \frac{\varepsilon}{2R}\le |A_{ij}|
  \qquad\text{for all } i,j .
\]
Equivalently, \(\operatorname{trunc}_{\varepsilon/(2R)}(A)=A\), so the
probability law remains the exact literal law
\(p_{ij}=A_{ij}^2/\Frob{A}^2\).  In exact arithmetic,
\[
  \mathbb P\!\left[
    \left\|A-\widetilde A\right\|_2\le \varepsilon
  \right]\ge 1-\delta .
\]
With valid \(\gamma_s\) and \(\gamma_{s+1}\), the computed non-probability
operations in the rounded sampled updates and residual satisfy
\[
  \mathbb P\!\left[
    \left\|A-\operatorname{fl}(\widetilde A)\right\|_2
    \le
    \varepsilon+
    \frac{2mn\,\Frob{A}^2}{\varepsilon}\,\gamma_{s+1}
  \right]\ge 1-\delta .
\]
In the small-\((s+1)u\) regime the rectangular FP additive term is
\[
  \Theta\!\left(
    \frac{mn\,\Frob{A}^2}{\varepsilon}\,s u
  \right),
\]
and, since \(R\le M\), the rectangular sample budget has order
\[
  \Theta\!\left(
    \frac{\max\{m,n\}\Frob{A}^2}{\varepsilon^2}
    \log\!\frac{m+n}{\delta}
  \right).
\]
\end{corollary}

\begin{formalized}
\leanfit{elementwiseTruncate_eq_self_of_forall_nonzero_entry_abs_ge}
\leanfit{sqMagTraceProbability_eventProb_algorithm1ExactSpectralEvent_ge_one_sub_delta_source_sample_budget_no_small_entries_square}
\leanfit{sqMagTraceProbability_eventProb_algorithm1FlSpectralRadius_ge_one_sub_delta_source_sample_budget_no_small_entries_square}
\leanfit{sqMagTraceProbability_eventProb_algorithm1ExactSpectralEvent_ge_one_sub_delta_source_sample_budget_no_small_entries_rect}
\leanfit{sqMagTraceProbability_eventProb_algorithm1FlSpectralRadius_ge_one_sub_delta_source_sample_budget_no_small_entries_rect}
\end{formalized}

\begin{remark}[Exact probability convention for Corollaries 144 and 145]
The preceding literal Algorithm~1 corollaries use the exact squared-magnitude
sampling law
\[
  p_{ij}=\frac{A_{ij}^2}{\Frob{A}^2}.
\]
Under the current convention there is no separate perturbation theorem for a
rounded probability table.  The implementation-facing FP radius is the radius
already displayed in the literal corollaries: it charges the rounded sampled
updates and residual arithmetic while keeping the probability law exact.
\end{remark}

\begin{remark}[Literal Algorithm 1 status]
The literal corollaries above are faithful to the Algorithm~1
distribution
\[
  p_{ij}=\frac{A_{ij}^2}{\Frob{A}^2}
\]
and contain no changed sampling law.  The Frobenius/Markov corollary is
unconditional but weaker.  The support-radius corollary is also
unconditional and nonconditional, with no nonzero-entry floor, but it exposes
the reciprocal-entry term \(G(A)\).  The source-rate corollary has the
paper-style \(n\log n\) sample budget, but it assumes that the source threshold
\(\varepsilon/(2n)\) would not remove any nonzero entry.  The source-uniform
literal matrix-Bernstein/Khintchine version of equation~(2), with neither
reciprocal-entry nor no-small-entry dependence, remains a separate theorem
target if kept in scope.
\end{remark}

\section{Algorithm 2: Row Sampling with Equation (4)}

For \(A\in\R^{m\times n}\), define
\[
  r_i(A)=\Norm{A_{i*}}^2,
  \qquad
  D(A)=\Frob{A}^2=\sum_i r_i(A).
\]
Equation (4) samples row \(i\) with probability
\[
  p_i=\frac{r_i(A)}{D(A)}.
\]
If \(I_t\) is the row sampled at trial \(t\), Algorithm 2 returns
\[
  \Atil_{tj}=\frac{A_{I_tj}}{\sqrt{s\,p_{I_t}}}.
\]

By current convention the row probabilities \(p_i\) are exact mathematical
sampling laws.  The implementation-facing formalization still allows computed
scale denominators
\(\widehat d_i\approx\sqrt{s\,p_i}\), with explicit denominator errors.  The
rounded-row sketch is then
\[
  \widehat{\Atil}_{tj}
  =
  \fl\!\left(
    \frac{A_{I_tj}}{\widehat d_{I_t}}
  \right),
\]
so square-root denominator computation and row-rescaling division are visible
in the theorem surface.  No FP error is charged for row-probability
construction, normalization, or sampling-law perturbation.

\begin{formalized}
\leanfit{ComputedRowScaleDen}
\leanfit{rowSampleSketchWithComputedDen}
\leanfit{fl_rowSampleSketchWithComputedDen_error_bound}
\end{formalized}

\begin{theorem}[Unbiased sampled Gram matrix]
Assume \(D(A)>0\) and \(s>0\).  Under the independent row sampler with
probabilities \(p_i=r_i(A)/D(A)\),
\[
  \E[(\Atil^T\Atil)_{jk}]=(A^TA)_{jk}
  \qquad\text{for all }j,k.
\]
\end{theorem}

\begin{formalized}
\leanfit{rowSqNormTraceProbability_expectationReal_rowSampleGram_entry}
\end{formalized}

\begin{theorem}[Equation (5)]
Assume \(D(A)>0\), \(s>0\), and \(\varepsilon>0\).  Then
\[
  \Prob\!\left[
    \Frob{\Atil^T\Atil-A^TA}
    \le \varepsilon D(A)
  \right]
  \ge 1-\frac{1}{s\varepsilon^2}.
\]
\end{theorem}

\begin{formalized}
\leanfit{rowSqNormTraceProbability_expectationReal_rowSampleGram_frob_error_sq_le}
\leanfit{rowSqNormTraceProbability_eventProb_rowSampleGram_frob_error_le_epsilon}
\end{formalized}

\begin{corollary}[Computed-denominator floating-point equation (5)]
Let
\[
  d_i=\sqrt{s\,p_i},\qquad
  \widehat d_i=\operatorname{flsqrt}(\operatorname{flmul}(s,p_i)),
  \qquad
  \alpha(\uunit)=\sqrt{1+\uunit}\,\uunit+\uunit .
\]
The exact probabilities \(p_i=\Norm{A_{i*}}^2/D(A)\) are not rounded, but the
non-probability denominator table is computed.  On positive-probability rows,
Lean proves
\[
  |\widehat d_i-d_i|\le d_i\,\alpha(\uunit),
  \qquad
  \widehat d_i\ne0.
\]
Define
\[
  \rho_A=
  \sum_{i:\,p_i>0}\frac{d_i\,\alpha(\uunit)}{|\widehat d_i|},
  \qquad
  \eta_A=\uunit+(1+\uunit)\rho_A .
\]
If \(\widehat A_{\rm comp}\) is the sampled sketch with entries
\[
  (\widehat A_{\rm comp})_{tj}
  =
  \operatorname{fldiv}(A_{I_tj},\widehat d_{I_t}),
\]
then, under the same sampling assumptions,
\[
  \Prob\!\left[
    \Frob{\widehat A_{\rm comp}^{T}\widehat A_{\rm comp}-A^TA}
    \le
    \varepsilon D(A)+n(2\eta_A+\eta_A^2)D(A)
  \right]
  \ge 1-\frac{1}{s\varepsilon^2}.
\]
If each Gram entry is also computed by the repository length-\(s\)
dot-product kernel, then
\[
  \Prob\!\left[
    \Frob{\widehat G_{\rm dot}-A^TA}
    \le
    \varepsilon D(A)
    +n\!\left((2\eta_A+\eta_A^2)
      +\gamma_s(1+\eta_A)^2\right)D(A)
  \right]
  \ge 1-\frac{1}{s\varepsilon^2}.
\]
The additive fully computed FP term has order
\[
  \Theta\!\left(
    nD(A)\bigl(\eta_A+\eta_A^2+\gamma_s(1+\eta_A)^2\bigr)
  \right),
\]
and, in the small-roundoff regime \(\eta_A,\gamma_s\to0\),
\[
  \Theta\!\left(nD(A)(\eta_A+\gamma_s)\right).
\]
For fixed \(A,m,n,s\) and nonzero computed positive-probability denominators,
the additive FP term tends to zero as \(\uunit,\gamma_s\to0\), so the result is
non-vacuous.  The exact-denominator theorems are recovered by setting
\(\rho_A=0\) and \(\eta_A=\uunit\).
\end{corollary}

\begin{formalized}
\leanfit{ComputedRowScaleDen.flMulThenSqrt}
\leanfit{fl_rowSampleSketchWithComputedDen_total_error_bound_le_budget}
\leanfit{rowSampleGramRelPerturbBudget_eq_nat_mul}
\leanfit{rowSampleGramDotProductRelBudget_eq_nat_mul}
\leanfit{fl_rowSampleGramWithComputedDen_perturb_bound}
\leanfit{fl_rowSampleGramDotWithComputedDen_perturb_bound}
\leanfit{rowSqNormTraceProbability_eventProb_fl_rowSampleGramWithComputedDen_frob_error_le_epsilon_add_explicit_budget}
\leanfit{rowSqNormTraceProbability_eventProb_fl_rowSampleGramDotWithComputedDen_frob_error_le_epsilon_add_explicit_budget}
\end{formalized}

\begin{remark}[Exact row probabilities for computed-denominator equation (5)]
The displayed bound uses the exact row-norm sampling law
\(p_i=\Norm{A_{i*}}^2/D(A)\). Under the current convention this law is exact,
so no probability-construction or sampling-law transfer term appears.  The
computed non-probability denominator, final row-scaling division, and optional
Gram dot-product kernel are all charged explicitly.
\end{remark}

\section{Algorithm 2: Leverage Scores and Equation (7)}

Let \(U\in\R^{m\times n}\) have orthonormal columns, so \(U^TU=I\).  Equation
(6) is the specialization of equation (4) to \(U\):
\[
  p_i=\frac{\Norm{U_{i*}}^2}{n}.
\]

By current convention, the leverage probabilities in equation~(6) are exact
sampling laws.  Floating-point construction of probabilities is not modeled in
the Algorithm~2 leverage theorems.  If a future implementation theorem needs a
computed SVD/QR basis, that basis perturbation should be stated separately from
probability construction and then transferred to the exact-law sampling
theorem.

\begin{theorem}[Equation (7), exact arithmetic]
Assume \(U^TU=I\), \(n>0\), \(s>0\), and \(\varepsilon>0\).  The leverage-score
row sampler satisfies
\[
  \Prob\!\left[
    \forall x,
    \Norm{(\Util^T\Util-I)x}
    \le \varepsilon\Norm{x}
  \right]
  \ge
  1-\frac{1}{s(\varepsilon/n)^2}.
\]
\end{theorem}

\begin{formalized}
\leanfit{leverageScoreProb_eq_rowNormSq_div_nat}
\leanfit{leverageTraceProbability_eventProb_rowSampleGram_opNorm2_error_le_epsilon}
\end{formalized}

\begin{lemma}[One-step leverage covariance facts]
Let \(u_i\in\R^n\) be row \(i\) of \(U\), and let
\[
  X_i=\begin{cases}
    u_i^Tu_i/p_i, & p_i>0,\\
    0, & p_i=0.
  \end{cases}
\]
For the leverage law \(p_i=\Norm{u_i}^2/n\), the one-sample estimator
satisfies
\[
  X_i\succeq0,\qquad
  \sum_i p_iX_i=I_n,\qquad
  X_i\preceq nI_n .
\]
\end{lemma}


\begin{formalized}
\leanfit{rowOuterGramSample_eq_zero_of_prob_zero}
\leanfit{finiteQuadraticForm_rowOuterGramSample_eq_sq_div}
\leanfit{finitePSD_rowOuterGramSample}
\leanfit{leverage_rowOuterGramSample_finitePSD}
\leanfit{leverage_rowOuterGramSample_mean_eq_id}
\leanfit{leverage_rowOuterGramSample_finiteLoewnerLe_nat}
\end{formalized}

\begin{lemma}[Row-trace MGF product-law adapter]
Let \(X_i\) be any self-adjoint matrix observable attached to one sampled row,
and let \(\omega=(i_1,\ldots,i_s)\) be drawn from the Algorithm~2 independent
row-trace law.  The no-hidden-Lieb one-step trace-MGF inequality iterates to
\[
  \E\operatorname{tr}\exp\!\left(
    H+\sum_{t=1}^s X_{i_t}
  \right)
  \le
  \operatorname{tr}\exp\!\left(
    H+s\log\E\exp(X_i)
  \right).
\]
For finite real matrices this is also connected to the repository's
`finiteMatrixExp` trace interface.  If the one-step logarithmic mean increment
is bounded by \(cI\), the right-hand side is at most
\[
  n\exp(sc).
\]
\end{lemma}


\begin{formalized}
\leanfit{rowSqNormSampleProbability}
\leanfit{rowSqNormTraceProbability_expectationReal_trace_normed_exp_add_step_le}
\leanfit{rowSqNormTraceProbability_expectationReal_succ_last_eq}
\leanfit{rowSqNormTraceProbability_expectationReal_trace_normed_exp_add_sum_le}
\leanfit{rowSqNormTraceProbability_expectationReal_finiteTrace_finiteMatrixExp_add_sum_le}
\leanfit{rowSqNormTraceProbabilityFiniteRealTraceMGFLogBound_zero_le_card_mul_exp_of_log_le_real_smul_one}
\end{formalized}

\begin{lemma}[Centered leverage covariance log-CGF]
Let
\[
  Y_i=X_i,\qquad Z_i=Y_i-I_n,
\]
where \(Y_i\) is the one-sample leverage outer-product estimator from the
previous lemma.  Under the leverage law \(p_i=\Norm{u_i}^2/n\),
\[
  \E Z_i=0,\qquad Z_i=Z_i^T,\qquad \lambda_{\max}(Z_i)\le n.
\]
The exact second-moment proxy is
\[
  \E[Z_i^2]=(n-1)I_n.
\]
Consequently, for every \(\theta\ge0\),
\[
  \log\E\exp(\theta Z_i)
  \preceq
  \frac{\exp(\theta n)-\theta n-1}{n^2}\,
  \E[Z_i^2].
\]
The lower-tail observable \( -Z_i=I_n-Y_i\) is handled separately:
\[
  \lambda_{\max}(-Z_i)\le 1,\qquad
  \log\E\exp(\theta(-Z_i))
  \preceq
  (\exp\theta-\theta-1)\E[Z_i^2].
\]
\end{lemma}


\begin{formalized}
\leanfit{Algorithms/RandNLA/RowSamplingLeverageMGF.lean}
\leanfit{rowOuterGramSample_symmetric}
\leanfit{rowOuterGramSample_centered_symmetric}
\leanfit{leverage_rowOuterGramSample_centered_cstar_selfAdjoint}
\leanfit{leverage_rowOuterGramSample_centered_expectationCStarMatrix_eq_zero}
\leanfit{leverage_rowOuterGramSample_centered_finiteLoewnerLe_nat}
\leanfit{leverage_rowOuterGramSample_centered_spectrum_le_nat}
\leanfit{leverage_rowOuterGramSample_centered_square_expectationCStarMatrix_eq}
\leanfit{leverage_rowOuterGramSample_centered_log_cgf_le}
\leanfit{leverage_rowOuterGramSample_centered_log_cgf_le_scalar}
\leanfit{leverage_rowOuterGramSample_neg_centered_log_cgf_le}
\leanfit{leverage_rowOuterGramSample_neg_centered_log_cgf_le_scalar}
\end{formalized}

\begin{theorem}[Two-sided leverage row-sampling concentration]
Let \(i_1,\ldots,i_s\) be independent samples from the leverage law
\(p_i=\Norm{u_i}^2/n\), and set
\[
  \widehat G=\frac1s\sum_{t=1}^sY_{i_t}.
\]
For every \(\theta>0\) and every \(\varepsilon\in\R\), define
\[
  \beta_+(\theta)=
  \frac{\exp(\theta n)-\theta n-1}{n^2},
  \qquad
  \beta_-(\theta)=\exp\theta-\theta-1.
\]
Then the repository proves
\[
\begin{aligned}
&\Prob\!\left[
  \widehat G-I_n\preceq\varepsilon I_n
  \ \hbox{and}\
  I_n-\widehat G\preceq\varepsilon I_n
\right] \\
&\qquad\ge
1-
n\exp\!\left(-\theta s\varepsilon+s\beta_+(\theta)(n-1)\right)
-
n\exp\!\left(-\theta s\varepsilon+s\beta_-(\theta)(n-1)\right).
\end{aligned}
\]
Equivalently, if the two displayed scalar tails sum to at most \(\delta\),
then the same event holds with probability at least \(1-\delta\).
\end{theorem}


\begin{formalized}
\leanfit{leverage_rowSqNormTraceProbabilityFiniteRealTraceMGFLogBound_centered_le}
\leanfit{leverage_rowSqNormTraceProbabilityFiniteRealTraceMGFLogBound_neg_centered_le}
\leanfit{leverage_rowSqNormTraceProbability_eventProb_exists_finiteHermitianEigenvalue_centered_sum_ge_le_exp}
\leanfit{leverageTraceProbability_eventProb_rowSampleGram_finiteLoewnerLe_upper_lt_ge_one_sub_exp}
\leanfit{leverageTraceProbability_eventProb_rowSampleGram_finiteLoewnerLe_lower_lt_ge_one_sub_exp}
\leanfit{leverageTraceProbability_eventProb_rowSampleGram_two_sided_finiteLoewnerLe_ge_one_sub_exp}
\leanfit{leverageTraceProbability_eventProb_rowSampleGram_two_sided_finiteLoewnerLe_ge_one_sub_delta_of_tail_budget}
\end{formalized}

\begin{corollary}[Bennett sample-budget equation (7)]
Assume \(n>1\), \(\varepsilon>0\), \(\delta>0\), and
\[
  \log\frac{2n}{\delta}
  \le
  s\,\frac{\varepsilon^2}{2(n-1)+(2/3)n\varepsilon},
  \qquad
  \log\frac{2n}{\delta}
  \le
  s\,\frac{\varepsilon^2}{2(n-1)+(2/3)\varepsilon}.
\]
Then
\[
  \Prob\!\left[
  \widehat G-I_n\preceq\varepsilon I_n
  \ \hbox{and}\
  I_n-\widehat G\preceq\varepsilon I_n
  \right]\ge 1-\delta .
\]
\end{corollary}


\begin{formalized}
\leanfit{real_bernstein_tail_le_half_delta_of_quadratic_budget}
\leanfit{leverageTraceProbability_eventProb_rowSampleGram_finiteLoewnerLe_upper_ge_one_sub_delta_half_of_sample_budget}
\leanfit{leverageTraceProbability_eventProb_rowSampleGram_finiteLoewnerLe_lower_ge_one_sub_delta_half_of_sample_budget}
\leanfit{leverageTraceProbability_eventProb_rowSampleGram_two_sided_finiteLoewnerLe_ge_one_sub_delta_of_sample_budget}
\end{formalized}

\begin{corollary}[Computed-denominator floating-point Bennett equation (7)]
With the same sample-budget assumptions and \(\gamma_s\) validity, let
\(\widetilde G_{\rm comp}\) be the Gram matrix computed from the leverage
sketch using the concrete denominator routine
\[
  \widehat d_i
  =
  \operatorname{flsqrt}\!\left(
    \operatorname{flmul}\!\left(s,\frac{\Norm{U_{i*}}^2}{n}\right)
  \right),
\]
rounded row scaling, and rounded length-\(s\) Gram dot products.  Define
\[
  \eta_U=\uunit+(1+\uunit)
  \sum_{i:\,\Norm{U_{i*}}^2>0}
  \frac{\sqrt{s\Norm{U_{i*}}^2/n}\,
        (\sqrt{1+\uunit}\,\uunit+\uunit)}
       {\left|
        \operatorname{flsqrt}\!\left(
          \operatorname{flmul}\!\left(s,\Norm{U_{i*}}^2/n\right)
        \right)\right|}.
\]
Since \(U^TU=I_n\), \(D(U)=n\).  Hence the fully computed FP radius is
\[
  \tau_{\rm comp}(U;s)
  =
  n^2\left((2\eta_U+\eta_U^2)+\gamma_s(1+\eta_U)^2\right).
\]
The theorem proved in Lean is
\[
  \Prob\!\left[
  \widetilde G_{\rm comp}-I_n\preceq
    \bigl(\varepsilon+\tau_{\rm comp}(U;s)\bigr) I_n
  \ \hbox{and}\
  I_n-\widetilde G_{\rm comp}\preceq
    \bigl(\varepsilon+\tau_{\rm comp}(U;s)\bigr) I_n
  \right]\ge 1-\delta .
\]
In the small-roundoff regime,
\[
  \tau_{\rm comp}(U;s)
  =
  \Theta\!\left(n^2(\eta_U+\gamma_s)\right).
\]
The exact-denominator comparison theorem is recovered by setting
\(\eta_U=\uunit\).
\end{corollary}

\begin{formalized}
\leanfit{unitRoundoff_lt_one_of_pos_gammaValid}
\leanfit{leverageFlMulThenSqrtRowScaleDen}
\leanfit{leverageFlMulThenSqrtRowScaleDen_den}
\leanfit{leverage_fl_rowSampleGramDotWithComputedDen_perturb_bound}
\leanfit{leverageTraceProbability_eventProb_fl_rowSampleGramDotWithComputedDen_two_sided_finiteLoewnerLe_ge_one_sub_delta_of_sample_budget}
\leanfit{leverageTraceProbability_eventProb_fl_rowSampleGramDotWithFlMulThenSqrtDen_two_sided_finiteLoewnerLe_ge_one_sub_delta_of_sample_budget}
\end{formalized}

\begin{corollary}[Actual-input Bennett equation (7)]
Let \(U\in\R^{m\times r}\) have orthonormal columns, let
\(C\in\R^{r\times q}\), and set \(A=UC\).  The exact leverage law is
\[
  p_i^U=\Norm{U_{i*}}^2/r.
\]
The factors \(U\) and \(C\) are exact analysis witnesses in this theorem, not
computed outputs.  The implemented row sketch samples rows of the actual input
\(A\) from this exact law.  Under the Bennett sample-budget hypotheses with
\(r\) replacing \(n\), the exact sampled actual-input Gram satisfies
\[
  \Prob\!\left[
    \widetilde A_U^T\widetilde A_U-A^TA\preceq\varepsilon A^TA
    \ \hbox{and}\
    A^TA-\widetilde A_U^T\widetilde A_U\preceq\varepsilon A^TA
  \right]\ge 1-\delta .
\]
The concrete floating-point theorem computes
\[
  \widehat d_i^U=\operatorname{flsqrt}(\operatorname{flmul}(s,p_i^U)),
\]
rounds the sampled-row divisions on entries of \(A\), and computes the Gram
entries by length-\(s\) floating-point dot products.  Define
\[
  \eta_U^A
  =
  \uunit+(1+\uunit)
  \sum_{i:\,p_i^U>0}
  \frac{\sqrt{s p_i^U}\,(\sqrt{1+\uunit}\,\uunit+\uunit)}
       {|\operatorname{flsqrt}(\operatorname{flmul}(s,p_i^U))|}
\]
and
\[
  T_A^{\rm fp}
  =
  \left(\gamma_s(1+\eta_U^A)^2+2\eta_U^A+(\eta_U^A)^2\right)
  \sum_{j=1}^q
  \left(\sum_{i:\,p_i^U>0}\frac{|A_{ij}|}{\sqrt{p_i^U}}\right)^2 .
\]
Then the computed Gram \(\widehat G_{\rm dot}(A;U)\) satisfies
\[
  \Prob\!\left[
    \widehat G_{\rm dot}(A;U)-A^TA
      \preceq \varepsilon A^TA+T_A^{\rm fp}I_q
    \ \hbox{and}\
    A^TA-\widehat G_{\rm dot}(A;U)
      \preceq \varepsilon A^TA+T_A^{\rm fp}I_q
  \right]\ge 1-\delta .
\]
With \(k_U=|\{i:p_i^U>0\}|\), fixed \(U,C,m,r,q,s\), and nonzero
positive-probability computed denominators,
\[
  T_A^{\rm fp}
  =
  \Theta\!\left(
    (\gamma_s+\uunit(1+k_U))
    \sum_{j=1}^q
    \left(\sum_{i:\,p_i^U>0}\frac{|A_{ij}|}{\sqrt{p_i^U}}\right)^2
  \right),
\]
so the floating-point additive radius tends to zero as
\(\uunit,\gamma_s\to0\).  A concrete QR, SVD, or rank-revealing routine that
produces \(U\) or \(C\) must still be certified separately before it is composed
with this endpoint.
\end{corollary}

\begin{formalized}
\leanfit{leverageFactoredInputSampleGram}
\leanfit{leverageTraceProbability_eventProb_factoredInputSampleGram_two_sided_finiteLoewnerLe_ge_one_sub_delta_of_sample_budget}
\leanfit{leverageFactoredInputDenGramBudget}
\leanfit{leverage_fl_rowSampleGramDotWithComputedDen_factoredInput_perturb_bound}
\leanfit{leverageTraceProbability_eventProb_factoredInput_fl_rowSampleGramDotWithFlMulThenSqrtDen_two_sided_finiteLoewnerLe_ge_one_sub_delta_of_sample_budget}
\end{formalized}

\begin{corollary}[Stored-basis floating-point Bennett equation (7)]
The exact leverage law is still \(p_i=\Norm{U_{i*}}^2/n\), and it remains an
exact sampling law.  The computed sketch may instead use a stored basis table
\[
  \widehat U^{(\mathrm{mul1})}_{ij}=\operatorname{flmul}(U_{ij},1),
  \qquad
  \widehat U^{(\mathrm{add0})}_{ij}=\operatorname{fladd}(U_{ij},0),
  \qquad
  \widehat U^{(\mathrm{sub0})}_{ij}=\operatorname{flsub}(U_{ij},0),
\]
each satisfying
\[
  |\widehat U^{(\mu)}_{ij}-U_{ij}|\le \uunit |U_{ij}|.
\]
Let
\[
  \rho_U=
  \sum_{i:\,p_i>0}
  \frac{\sqrt{s p_i}\,(\sqrt{1+\uunit}\,\uunit+\uunit)}
       {|\operatorname{flsqrt}(\operatorname{flmul}(s,p_i))|},
  \qquad
  \lambda_U=2\uunit+\uunit^2+(1+\uunit)^2\rho_U .
\]
The fully expanded stored-basis radius can be written as
\[
  T_U^{\rm fp}
  =
  \left(\gamma_s(1+\lambda_U)^2+2\lambda_U+\lambda_U^2\right)
  \sum_{j=1}^n
  \left(\sum_{i:\,p_i>0}\frac{|U_{ij}|}{\sqrt{p_i}}\right)^2 .
\]
For each \(\mu\in\{\mathrm{mul1},\mathrm{add0},\mathrm{sub0}\}\), the Gram
matrix computed from \(\widehat U^{(\mu)}\), the concrete denominator
\(\operatorname{flsqrt}(\operatorname{flmul}(s,p_i))\), rounded row divisions,
and rounded length-\(s\) dot products satisfies the same Bennett probability
lower bound with event radius \(\varepsilon+T_U^{\rm fp}\).  In the fixed
problem small-roundoff regime,
\[
  T_U^{\rm fp}
  =
  \Theta\!\left(
    (\gamma_s+\uunit(1+|\{i:p_i>0\}|))
    \sum_{j=1}^n
    \left(\sum_{i:\,p_i>0}\frac{|U_{ij}|}{\sqrt{p_i}}\right)^2
  \right),
\]
and the radius tends to zero as \(\uunit,\gamma_s\to0\).  This closes the
three modeled stored-table paths; computing \(U\) from a data matrix by QR,
SVD, or another routine remains a separate computed-basis routine obligation.
\end{corollary}

\begin{formalized}
\leanfit{RowSamplingLeverageComputedBasis.lean}
\leanfit{leverageExactBasisSampleColumnAbsBudget}
\leanfit{leverageComputedBasisSampleColumnErrorBudget}
\leanfit{leverageComputedBasisDenGramBudget}
\leanfit{leverageTraceProbability_eventProb_fl_rowSampleGramDotWithStoredBasisMulOneAndFlMulThenSqrtDen_two_sided_finiteLoewnerLe_ge_one_sub_delta_of_sample_budget}
\leanfit{leverageTraceProbability_eventProb_fl_rowSampleGramDotWithStoredBasisAddZeroRightAndFlMulThenSqrtDen_two_sided_finiteLoewnerLe_ge_one_sub_delta_of_sample_budget}
\leanfit{leverageTraceProbability_eventProb_fl_rowSampleGramDotWithStoredBasisSubZeroRightAndFlMulThenSqrtDen_two_sided_finiteLoewnerLe_ge_one_sub_delta_of_sample_budget}
\end{formalized}

\begin{remark}[Exact leverage probabilities for computed-denominator equation (7)]
The displayed Bennett FP theorem samples from the exact leverage-score law
defined by the ideal orthonormal basis \(U\).  Under the current convention,
there is no rounded leverage-probability table and no sampling-law transfer
loss.  The computed-denominator theorem charges the non-probability
denominator, rounded row scaling, and rounded Gram dot products.  The
stored-basis theorem additionally charges the three modeled rounded-copy basis
tables.  Any future QR/SVD-generated basis theorem must instantiate a concrete
computed-basis routine rather than being hidden inside the leverage law.
\end{remark}

\section{Algorithm 3: Random-Projection Preconditioning}

Algorithm 3 applies preprocessing matrices to an input matrix \(A\):
\[
  \Pi_LA,\qquad A\Pi_R,\qquad \Pi_LA\Pi_R.
\]
The random projection law is intentionally not fixed in the meta-algorithm.
The Lean results are therefore deterministic after \(\Pi_L\) and \(\Pi_R\) are
chosen.

For an implementation, the preprocessing matrices may themselves be generated
or stored in floating point before the matrix products are evaluated.  The
implementation-facing surface records computed matrices
\(\widehat\Pi_L,\widehat\Pi_R\) with entrywise errors against the ideal
preconditioners, and separates those transform-entry errors from the rounded
matrix-multiplication error in \(\fl(\widehat\Pi_L A)\),
\(\fl(A\widehat\Pi_R)\), and
\(\fl(\widehat\Pi_L A\widehat\Pi_R)\).  The uniform row sampler used after
preconditioning has a separate certificate for the computed
\(\widehat d\approx\sqrt{s/m}\) scale denominator.

\begin{formalized}
\leanfit{ComputedPreconditioner}
\leanfit{fl_preconditionRowsWithComputedLeft_total_error_bound}
\leanfit{fl_preconditionColumnsWithComputedRight_total_error_bound}
\leanfit{fl_preconditionElementsWithComputed_error_bound}
\leanfit{fl_preconditionElementsWithComputed_total_error_bound}
\leanfit{ComputedUniformRowScaleDen}
\leanfit{fl_uniformRowSampleSketch_computedDen_error_bound}
\leanfit{uniformRowSampleIncrementWithComputedDen_ideal_error_bound}
\leanfit{fl_uniformRowSampleSketch_computedDen_total_error_bound}
\leanfit{fl_uniformRowSampleSketch_computedDen_total_error_bound_exact}
\leanfit{ComputedUniformRowScaleDen.flDivThenSqrt}
\leanfit{ComputedUniformRowScaleDen.flInvMulThenSqrt}
\leanfit{ComputedUniformRowScaleDen.flSqrtDivSqrt}
\leanfit{ComputedUniformRowScaleDen.flSqrtMulInvSqrt}
\end{formalized}

\begin{theorem}[Computed uniform-row denominator certificate]
Let
\[
  d=\sqrt{s/m},\qquad \widehat d\ne0,\qquad |\widehat d-d|\le \eta_d,
\]
where \(d\ne0\).  For an exact uniform row sample \(i\) and an implemented row
scaling division using the computed denominator \(\widehat d\),
\[
  B_{ij}=U_{ij}/d,\qquad
  \widehat B_{ij}^{\,\widehat d}=\fl(U_{ij}/\widehat d),
\]
one has
\[
  \left|\widehat B_{ij}^{\,\widehat d}-B_{ij}\right|
  \le
  \left|U_{ij}/\widehat d\right|\,u
  +
  \frac{|U_{ij}|\,\eta_d}{|\widehat d|\,|d|}.
\]
Thus the row law remains the exact uniform law \(1/m\), while the
non-probability scalar computed by the algorithm is charged explicitly.  In the
exact-denominator specialization \(\widehat d=d\) and \(\eta_d=0\), this reduces
to the usual rounded-division bound
\[
  \left|\fl(U_{ij}/d)-U_{ij}/d\right|\le |U_{ij}/d|\,u.
\]
\end{theorem}

\begin{formalized}
\leanfit{ComputedUniformRowScaleDen.exact}
\leanfit{uniformRowSampleIncrementWithComputedDen_ideal_error_bound}
\leanfit{fl_uniformRowSampleSketch_computedDen_total_error_bound}
\leanfit{fl_uniformRowSampleSketch_computedDen_total_error_bound_exact}
\end{formalized}

\begin{corollary}[Rounded-sqrt exact-input denominator certificate]
Assume \(0<u<1\) in the floating-point model and that the scalar
\((s/m)\) is supplied exactly to the square-root primitive.  If
\[
  \widehat d=\fl(\sqrt{s/m}),
\]
then the square-root model gives
\[
  \widehat d\ne0,\qquad |\widehat d-d|\le d\,u,
  \qquad d=\sqrt{s/m}.
\]
Consequently, the computed-denominator theorem applies with
\(\eta_d=d\,u\):
\[
  \left|\widehat B_{ij}^{\,\widehat d}-B_{ij}\right|
  \le
  \left|U_{ij}/\widehat d\right|u
  +
  \frac{|U_{ij}|\,d\,u}{|\widehat d|\,|d|}.
\]
This corollary charges the rounded square-root primitive.  If the implementation
also computes \(s/m\) by rounded division or multiplication, that earlier scalar
routine must instantiate a separate certificate before this exact-input
constructor is used.
\end{corollary}

\begin{formalized}
\leanfit{ComputedUniformRowScaleDen.flSqrtExactInput}
\leanfit{ComputedUniformRowScaleDen.flSqrtExactInput_den}
\leanfit{ComputedUniformRowScaleDen.flSqrtExactInput_den_abs_error}
\end{formalized}

\begin{corollary}[Rounded division-then-square-root denominator certificate]
Assume \(0<u<1\), \(m>0\), and \(s>0\).  Let the implementation compute
\[
  \widehat x=\fl(s/m),\qquad
  \widehat d=\fl(\sqrt{\widehat x}),\qquad
  d=\sqrt{s/m}.
\]
Then \(\widehat d\ne0\) and the denominator error satisfies
\[
  |\widehat d-d|
  \le
  d\bigl(\sqrt{1+u}\,u+u\bigr).
\]
Consequently, for an exact uniform row sample \(i\),
\[
  \left|\fl(U_{ij}/\widehat d)-U_{ij}/d\right|
  \le
  \left|U_{ij}/\widehat d\right|u
  +
  \frac{|U_{ij}|\,d\bigl(\sqrt{1+u}\,u+u\bigr)}
       {|\widehat d|\,|d|}.
\]
The row-sampling law remains the exact uniform law \(1/m\).  The two displayed
losses charge only non-probability computations performed by the implementation:
rounded scalar division in \(s/m\), rounded square root, and the final rounded
row-scaling division.
\end{corollary}

\begin{formalized}
\leanfit{ComputedUniformRowScaleDen.abs_sqrt_one_add_sub_one_le_abs}
\leanfit{ComputedUniformRowScaleDen.flDivThenSqrt}
\leanfit{ComputedUniformRowScaleDen.flDivThenSqrt_den}
\leanfit{ComputedUniformRowScaleDen.flDivThenSqrt_den_abs_error}
\end{formalized}

\begin{corollary}[Rounded reciprocal-multiply-square-root denominator certificate]
Assume \(0<u<1\), \(m>0\), and \(s>0\).  Let the implementation compute
\[
  \widehat r=\fl(1/m),\qquad
  \widehat x=\fl(s\widehat r),\qquad
  \widehat d=\fl(\sqrt{\widehat x}),\qquad
  d=\sqrt{s/m}.
\]
Then \(\widehat d\ne0\) and
\[
  |\widehat d-d|
  \le
  d\left(\sqrt{(1+u)(1+u)}\,u+2u+u^2\right).
\]
Consequently, for an exact uniform row sample \(i\),
\[
  \left|\fl(U_{ij}/\widehat d)-U_{ij}/d\right|
  \le
  \left|U_{ij}/\widehat d\right|u
  +
  \frac{|U_{ij}|\,d\left(\sqrt{(1+u)(1+u)}\,u+2u+u^2\right)}
       {|\widehat d|\,|d|}.
\]
This corollary charges the reciprocal-as-division \(1/m\), the rounded
multiplication by \(s\), the rounded square root, and the final rounded
row-scaling division.  It introduces no probability-construction,
probability-normalization, or sampling-law transfer term: the row law remains
the exact uniform law \(1/m\).
\end{corollary}

\begin{formalized}
\leanfit{ComputedUniformRowScaleDen.flInvMulThenSqrt}
\leanfit{ComputedUniformRowScaleDen.flInvMulThenSqrt_den}
\leanfit{ComputedUniformRowScaleDen.flInvMulThenSqrt_den_abs_error}
\end{formalized}

\begin{corollary}[Rounded split-square-root denominator certificate]
Assume \(0<u<1\), \(m>0\), and \(s>0\).  Let the implementation compute
\[
  \widehat s_0=\fl(\sqrt{s}),\qquad
  \widehat m_0=\fl(\sqrt{m}),\qquad
  \widehat d=\fl(\widehat s_0/\widehat m_0),\qquad
  d=\sqrt{s/m}.
\]
Then \(\widehat d\ne0\) and
\[
  |\widehat d-d|
  \le
  d\,\frac{3u+u^2}{1-u}.
\]
Consequently, for an exact uniform row sample \(i\),
\[
  \left|\fl(U_{ij}/\widehat d)-U_{ij}/d\right|
  \le
  \left|U_{ij}/\widehat d\right|u
  +
  \frac{|U_{ij}|\,d(3u+u^2)}
       {|\widehat d|\,|d|\,(1-u)}.
\]
This corollary charges the two rounded square roots, the rounded scalar
division used to form \(\widehat d\), and the final rounded row-scaling
division.  The row-sampling law remains the exact uniform law \(1/m\).
\end{corollary}

\begin{formalized}
\leanfit{ComputedUniformRowScaleDen.flSqrtDivSqrt}
\leanfit{ComputedUniformRowScaleDen.flSqrtDivSqrt_den}
\leanfit{ComputedUniformRowScaleDen.flSqrtDivSqrt_den_abs_error}
\end{formalized}

\begin{corollary}[Rounded square-root times reciprocal-square-root denominator certificate]
Assume \(0<u<1\), \(m>0\), and \(s>0\).  Let the implementation compute
\[
  \widehat s_0=\fl(\sqrt{s}),\qquad
  \widehat m_0=\fl(\sqrt{m}),\qquad
  \widehat r_m=\fl(1/\widehat m_0),\qquad
  \widehat d=\fl(\widehat s_0\,\widehat r_m),\qquad
  d=\sqrt{s/m}.
\]
Then \(\widehat d\ne0\) and
\[
  |\widehat d-d|
  \le
  d\,\frac{4u+3u^2+u^3}{1-u}.
\]
Consequently, for an exact uniform row sample \(i\),
\[
  \left|\fl(U_{ij}/\widehat d)-U_{ij}/d\right|
  \le
  \left|U_{ij}/\widehat d\right|u
  +
  \frac{|U_{ij}|\,d(4u+3u^2+u^3)}
       {|\widehat d|\,|d|\,(1-u)}.
\]
This corollary charges the two rounded square roots, the rounded reciprocal of
\(\widehat m_0\), the final rounded multiplication used to form
\(\widehat d\), and the final rounded row-scaling division.  It does not charge
the construction of the uniform sampling law, which remains the exact law
\(1/m\).
\end{corollary}

\begin{formalized}
\leanfit{ComputedUniformRowScaleDen.flSqrtMulInvSqrt}
\leanfit{ComputedUniformRowScaleDen.flSqrtMulInvSqrt_den}
\leanfit{ComputedUniformRowScaleDen.flSqrtMulInvSqrt_den_abs_error}
\end{formalized}

\begin{theorem}[Exact Frobenius stability for orthogonal preprocessing]
If \(A,\Pi_L,\Pi_R\in\R^{n\times n}\) and the preprocessing matrices are
orthogonal, then
\[
  \Frob{\Pi_L A}=\Frob{A},\qquad
  \Frob{A\Pi_R}=\Frob{A},\qquad
  \Frob{\Pi_LA\Pi_R}=\Frob{A}.
\]
\end{theorem}

\begin{formalized}
\leanfit{preconditionRows_frobNorm_orthogonal}
\leanfit{preconditionColumns_frobNorm_orthogonal}
\leanfit{preconditionElements_frobNorm_orthogonal}
\end{formalized}

\begin{theorem}[Orthonormal-column invariance under orthogonal preprocessing]
Let \(U\in\R^{m\times n}\) satisfy
\[
  U^TU=I_n,
  \qquad\text{i.e.}\qquad
  \sum_{i=1}^m U_{ij}U_{ik}=\delta_{jk}.
\]
If \(\Pi_L\in\R^{m\times m}\) and \(\Pi_R\in\R^{n\times n}\) are
orthogonal, then
\[
  (\Pi_LU)^T(\Pi_LU)=I_n,\qquad
  (U\Pi_R)^T(U\Pi_R)=I_n,\qquad
  (\Pi_LU\Pi_R)^T(\Pi_LU\Pi_R)=I_n .
\]
\end{theorem}

\begin{corollary}[Leverage denominator after Algorithm 3 preprocessing]
Under the same hypotheses, the equation~(6) leverage denominator remains
the number of columns:
\[
  \sum_{i=1}^m \Norm{(\Pi_LU)_{i,*}}^2
  =
  \sum_{i=1}^m \Norm{(U\Pi_R)_{i,*}}^2
  =
  \sum_{i=1}^m \Norm{(\Pi_LU\Pi_R)_{i,*}}^2
  =
  n.
\]
This is a deterministic basis-invariance statement.  It does not say that a
particular random projection law makes the individual leverage scores nearly
uniform.
\end{corollary}

\begin{formalized}
\leanfit{preconditionRows_hasOrthonormalColumns_of_orthogonal}
\leanfit{preconditionColumns_hasOrthonormalColumns_of_orthogonal}
\leanfit{preconditionElements_hasOrthonormalColumns_of_orthogonal}
\leanfit{rowSqNormProbDen_preconditionRows_eq_nat_of_orthogonal}
\leanfit{rowSqNormProbDen_preconditionColumns_eq_nat_of_orthogonal}
\leanfit{rowSqNormProbDen_preconditionElements_eq_nat_of_orthogonal}
\end{formalized}

\begin{theorem}[Signed orthogonal preprocessing]
Let \(H\in\R^{m\times m}\) be orthogonal and let
\(d\in\R^m\) satisfy \(d_i^2=1\) for every \(i\).  Write
\(D=\operatorname{diag}(d)\).  Then \(D\) is orthogonal, \(HD\) is
orthogonal, and for every \(U\in\R^{m\times n}\) with \(U^TU=I_n\),
\[
  (HDU)^T(HDU)=I_n,
  \qquad
  \sum_{i=1}^m \Norm{(HDU)_{i,*}}^2=n .
\]
\end{theorem}


\begin{formalized}
\leanfit{IsOrthogonal.diagMatrix_of_sq_eq_one}
\leanfit{signedOrthogonalPreconditioner_isOrthogonal}
\leanfit{signedOrthogonalPreconditionRows_hasOrthonormalColumns}
\leanfit{rowSqNormProbDen_signedOrthogonalPreconditionRows_eq_nat}
\end{formalized}

This closes the deterministic signed-orthogonal part of the SRHT proof-source
route.  The finite Rademacher support and moment facts are stated next.  Later
sections close the source-sharp signed-Hadamard row-norm flattening theorem in
explicit-\(t\) and logarithmic forms; the remaining Algorithm~3 obligations are
computed-object instantiations for concrete transform/basis/projector routines,
not probability-construction errors.

\begin{corollary}[Probability-one signed preprocessing under Rademacher signs]
Let \(\omega\) be uniformly distributed over sign traces in \(\{\pm1\}^m\),
and let \(D_\omega=\operatorname{diag}(\omega)\).  If \(H\in\R^{m\times m}\)
is orthogonal and \(U^TU=I_n\), then
\[
  \mathbb P\!\left[
    HD_\omega\ \text{is orthogonal},\quad
    (HD_\omega U)^T(HD_\omega U)=I_n,\quad
    \sum_{i=1}^m \Norm{(HD_\omega U)_{i,*}}^2=n
  \right]=1 .
\]
\end{corollary}


\begin{formalized}
\leanfit{FiniteProbability.eventProb_eq_one_of_forall}
\leanfit{RademacherTrace}
\leanfit{rademacherSignVector_sq}
\leanfit{rademacherTraceProbability}
\leanfit{rademacherTraceProbability_eventProb_signedOrthogonalPreconditionRows_eq_one}
\end{formalized}

\begin{theorem}[Finite Rademacher moments]
Let \(\omega\) be uniformly distributed over \(\{\pm1\}^m\).  Then for all
indices \(a,b\),
\[
  \mathbb E\,\omega_a = 0,
  \qquad
  \mathbb E\,\omega_a\omega_b=\delta_{ab}.
\]
Consequently, for every deterministic vector \(c\in\R^m\),
\[
  \mathbb E\left(\sum_{k=1}^m c_k\omega_k\right)^2
  =
  \sum_{k=1}^m c_k^2 .
\]
\end{theorem}

\begin{lemma}[Exact Rademacher signed-linear-form MGF]
Let \(\omega\) be uniformly distributed over \(\{\pm1\}^m\).  For every
deterministic vector \(a\in\R^m\),
\[
  \mathbb E \exp\left(\sum_{k=1}^m a_k\omega_k\right)
  =
  \prod_{k=1}^m \frac{\exp(a_k)+\exp(-a_k)}{2}.
\]
Consequently, for every \(\lambda>0\) and threshold \(T\),
\[
  \mathbb P\left\{
    \sum_{k=1}^m a_k\omega_k \le T
  \right\}
  \ge
  1-\exp(-\lambda T)
    \prod_{k=1}^m
      \frac{\exp(\lambda a_k)+\exp(-\lambda a_k)}2 .
\]
\end{lemma}


\begin{theorem}[Scalar Rademacher Hoeffding tail]
Let \(\omega\) be uniformly distributed over \(\{\pm1\}^m\).  For every
\(a\in\R^m\), every \(\lambda\in\R\), and every \(T\in\R\),
\[
  \mathbb E\exp\left(
    \lambda\sum_{k=1}^m a_k\omega_k
  \right)
  \le
  \exp\left(
    \frac{\lambda^2}{2}\sum_{k=1}^m a_k^2
  \right).
\]
Consequently, for \(\lambda>0\),
\[
  \mathbb P\left\{
    \sum_{k=1}^m a_k\omega_k\ge T
  \right\}
  \le
  \exp(-\lambda T)
  \exp\left(
    \frac{\lambda^2}{2}\sum_{k=1}^m a_k^2
  \right),
\]
and
\[
  \mathbb P\left\{
    \left|\sum_{k=1}^m a_k\omega_k\right|\le T
  \right\}
  \ge
  1-2\exp(-\lambda T)
  \exp\left(
    \frac{\lambda^2}{2}\sum_{k=1}^m a_k^2
  \right).
\]
\end{theorem}


\begin{formalized}
\leanfit{real_rademacher_cosh_factor_le_exp_sq_div_two}
\leanfit{rademacherTraceProbability_expectationReal_exp_sum_mul_sign_le_exp_sum_sq_div_two}
\leanfit{rademacherTraceProbability_expectationReal_exp_lam_sum_mul_sign_le_exp_lam_sq_sum_sq_div_two}
\leanfit{rademacherTraceProbability_eventProb_sum_mul_sign_ge_le_exp_sq_bound}
\leanfit{rademacherTraceProbability_eventProb_sum_mul_sign_le_ge_one_sub_exp_sq_bound}
\leanfit{rademacherTraceProbability_eventProb_abs_sum_mul_sign_le_ge_one_sub_two_mul_exp_sq_bound}
\end{formalized}

\begin{corollary}[Optimized scalar Rademacher Hoeffding tail]
Let \(\omega\) be uniformly distributed over the exact finite Rademacher cube
\(\{\pm1\}^m\).  Let \(a\in\R^m\), \(T>0\), and \(\sigma>0\).  If
\[
  \sum_{k=1}^m a_k^2 \le \sigma^2 ,
\]
then
\[
  \mathbb P\left\{
    \left|\sum_{k=1}^m a_k\omega_k\right|\le T
  \right\}
  \ge
  1-2\exp\left(-\frac{T^2}{2\sigma^2}\right).
\]
The statement optimizes the preceding auxiliary-parameter bound with
\(\lambda=T/\sigma^2\).  It is an exact-probability concentration primitive:
no floating-point computation is involved, and it does not by itself close the
Gaussian, FJLT, or input-sparsity Algorithm 3 transform laws.
\end{corollary}

\begin{formalized}
\leanfit{rademacherTraceProbability_eventProb_abs_sum_mul_sign_le_ge_one_sub_two_mul_exp_neg_sq_div}
\end{formalized}

\begin{corollary}[Finite-family Rademacher Hoeffding tail]
Let \(\omega\) be uniformly distributed over the exact finite Rademacher cube
\(\{\pm1\}^m\).  Let \(\mathcal I\) be a finite index set.  For each
\(\alpha\in\mathcal I\), let \(a_\alpha\in\R^m\), \(T_\alpha>0\), and
\(\sigma_\alpha>0\).  If
\[
  \sum_{k=1}^m a_{\alpha,k}^2 \le \sigma_\alpha^2
  \qquad\text{for every }\alpha\in\mathcal I,
\]
then
\[
  \mathbb P\left\{
    \left|\sum_{k=1}^m a_{\alpha,k}\omega_k\right|\le T_\alpha
    \text{ for every }\alpha\in\mathcal I
  \right\}
  \ge
  1-\sum_{\alpha\in\mathcal I}
    2\exp\left(-\frac{T_\alpha^2}{2\sigma_\alpha^2}\right).
\]
This is the finite union-bound packaging of the preceding exact-probability
Rademacher tail.  It is intended as a reusable primitive for later finite
Rademacher/FJLT preprocessing branches, not as a completed transform law.
\end{corollary}

\begin{formalized}
\leanfit{rademacherTraceProbability_eventProb_forall_abs_sum_mul_sign_le_ge_one_sub_sum_two_mul_exp_neg_sq_div}
\end{formalized}

\begin{corollary}[Rademacher signed-linear vector norm bound]
Let \(\omega\) be uniformly distributed over the exact finite Rademacher cube
\(\{\pm1\}^m\).  For \(j=1,\ldots,n\), let
\[
  X_j(\omega)=\sum_{k=1}^m a_{j,k}\omega_k .
\]
Let \(T_j>0\), \(\sigma_j>0\), and assume
\[
  \sum_{k=1}^m a_{j,k}^2 \le \sigma_j^2
  \qquad\text{for every }j .
\]
Then
\[
  \mathbb P\left\{
    \sum_{j=1}^n X_j(\omega)^2
    \le
    \sum_{j=1}^n T_j^2
  \right\}
  \ge
  1-\sum_{j=1}^n
    2\exp\left(-\frac{T_j^2}{2\sigma_j^2}\right).
\]
The proof uses the finite-family coordinate event and the deterministic
implication \(|X_j|\le T_j\) for all \(j\Rightarrow \sum_j X_j^2\le\sum_j
T_j^2\).  This is a row-norm-shaped exact-probability primitive, still not a
completed Algorithm 3 transform law.
\end{corollary}

\begin{formalized}
\leanfit{rademacherTraceProbability_eventProb_vecNorm2Sq_sum_mul_sign_le_ge_one_sub_sum_two_mul_exp_neg_sq_div}
\end{formalized}

\begin{corollary}[Exact signed-mixing row-norm bound]
Let \(G\in\R^{r\times m}\) be a deterministic exact mixing table and let
\(U\in\R^{m\times n}\) have orthonormal columns.  Let
\(\omega\in\{\pm1\}^m\) be an exact Rademacher sign vector and define the exact
signed-mixing preconditioner
\[
  \Pi_\omega(i,k)=G_{ik}\omega_k .
\]
For row thresholds \(T_{ij}>0\) and row coefficient caps
\(\alpha_i>0\), assume the deterministic exact coefficient bound
\[
  G_{ik}^2\le \alpha_i^2
  \qquad\text{for every }i,k .
\]
Then
\[
  \mathbb P\left\{
    \forall i,\quad
    \Norm{(\Pi_\omega U)_{i,*}}^2
      \le \sum_{j=1}^n T_{ij}^2
  \right\}
  \ge
  1-\sum_{i=1}^r\sum_{j=1}^n
    2\exp\left(-\frac{T_{ij}^2}{2\alpha_i^2}\right).
\]
In particular, if \(T_{ij}=B_i>0\) for all coordinates in row \(i\), then
\[
  \mathbb P\left\{
    \forall i,\quad
    \Norm{(\Pi_\omega U)_{i,*}}^2
      \le n B_i^2
  \right\}
  \ge
  1-\sum_{i=1}^r\sum_{j=1}^n
    2\exp\left(-\frac{B_i^2}{2\alpha_i^2}\right).
\]
When \(\alpha_i=\alpha\) and \(B_i=B\) are uniform, the displayed failure
term is exactly \(2rn\exp(-B^2/(2\alpha^2))\), hence has order
\[
  \Theta\!\left(rn\exp\left(-\frac{B^2}{2\alpha^2}\right)\right)
\]
in the regime where \(r,n,B/\alpha\) vary and the exact sampling law is fixed.
This is an exact-probability and exact-arithmetic preprocessing result:
the Rademacher law is exact by project convention, \(G\) and \(U\) are exact
analysis objects in this corollary, and no floating-point storage or transform
application is charged here.  Floating-point use of a concrete \(G\) remains a
separate computed-preconditioner obligation.
\end{corollary}

\begin{formalized}
\leanfit{signedMixingRows}
\leanfit{signedMixingRows_preconditionRows_entry}
\leanfit{signedMixingRows_coeff_sq_sum_le_of_entry_sq_le}
\leanfit{rademacherTraceProbability_eventProb_forall_rowNormSq_signedMixingRows_entry_sq_le_ge_one_sub_sum_sum_two_mul_exp_neg_sq_div}
\leanfit{rademacherTraceProbability_eventProb_forall_rowNormSq_signedMixingRows_entry_sq_le_uniform_ge_one_sub_sum_sum_two_mul_exp_neg_sq_div}
\end{formalized}

\begin{corollary}[Exact signed-mixing one-row mean identity]
Let \(G\in\R^{r\times m}\) be a deterministic exact mixing table and let
\(U\in\R^{m\times n}\) have orthonormal columns.  Let
\(\omega\in\{\pm1\}^m\) be an exact Rademacher sign vector and let
\(i\) be an exact uniform row index in \(\{1,\ldots,r\}\), independent of
\(\omega\).  Define \(\Pi_\omega(i,k)=G_{ik}\omega_k\).  Assume the exact
column-square normalization
\[
  \sum_{i=1}^r G_{ik}^2=1
  \qquad\text{for every }k=1,\ldots,m .
\]
This normalization is computable in exact arithmetic as an analysis
assumption on \(G\).  The formalization also proves the direct adapter from
the stronger rectangular-isometry hypothesis \(G^TG=I_m\), since that
hypothesis implies the displayed column-square identities.
For the one-row estimator \(r(\Pi_\omega U)_{i,*}^T(\Pi_\omega U)_{i,*}\),
\[
  \mathbb E_{\omega,i}\!
  \left[
    r(\Pi_\omega U)_{i,j}(\Pi_\omega U)_{i,\ell}
  \right]
  =
  \delta_{j\ell}.
\]
Equivalently, the expected one-row outer product is \(I_n\).  This closes the
exact first-moment/normalization layer for finite signed-mixing rows under
column-square normalization, and directly under \(G^TG=I_m\) by the adapter.
It is not a concentration theorem, not a subspace-embedding theorem, and not a
floating-point theorem: the law of
\((\omega,i)\) is exact by project convention, and applying or storing \(G\)
in floating point remains a separate computed-preconditioner obligation.
\end{corollary}

\begin{formalized}
\leanfit{rademacherTraceProbability_expectationReal_sum_mul_sign_mul_sum_mul_sign_eq_sum_mul}
\leanfit{signedMixingUniformRowSampleProbability}
\leanfit{rademacherTraceProbability_expectationReal_signedMixingRows_entry_mul_eq}
\leanfit{signedMixingUniformRowSampleProbability_expectationReal_uniformRowOuterGramSample_eq_id}
\leanfit{signedMixingUniformRowSampleProbability_expectationReal_uniformRowOuterGramSample_eq_id_of_hasOrthonormalColumns}
\end{formalized}

\begin{corollary}[Exact CountSketch first moment]
Let \(r>0\), and let
\[
  \mathcal H_{r,m}=\{h:\{1,\ldots,m\}\to\{1,\ldots,r\}\}
\]
be the finite set of exact CountSketch hash maps.  Let \(h\) be drawn from
the exact uniform law on \(\mathcal H_{r,m}\), and let
\(\omega\in\{\pm1\}^m\) be drawn from the exact Rademacher product law,
independently of \(h\).  For exact arithmetic, define the sparse
CountSketch preconditioner \(S_{h,\omega}\in\R^{r\times m}\) by
\[
  (S_{h,\omega})_{ik}
  =
  \begin{cases}
    \omega_k, & h(k)=i,\\
    0, & h(k)\ne i .
  \end{cases}
\]
Thus every input row \(k\) contributes to exactly one output bucket
\(h(k)\), with exact sign \(\omega_k\).

For any exact matrix \(U\in\R^{m\times n}\),
\[
  \mathbb E_{h,\omega}
    \left[(S_{h,\omega}U)^T(S_{h,\omega}U)\right]
  =
  U^T U .
\]
Equivalently, entrywise,
\[
  \mathbb E_{h,\omega}
    \left[
      {\rm rowGram}(S_{h,\omega}U)_{j\ell}
    \right]
  =
  {\rm rowGram}(U)_{j\ell}.
\]
If \(U^TU=I_n\), then the expected sketched Gram is \(I_n\).

The proof first fixes the exact hash \(h\).  Rademacher signs make all
collision cross terms have zero expectation, and summing over the \(r\)
buckets counts each input row once.  The final Lean theorem then packages the
exact uniform hash law and exact Rademacher sign law.

The formalization also proves the deterministic collision-free foundation.
If \(h\) is injective and the exact signs satisfy \(\omega_k^2=1\), then
\[
  S_{h,\omega}^T S_{h,\omega}=I_m,
\]
so for every exact matrix \(A\in\R^{m\times n}\),
\[
  (S_{h,\omega}A)^T(S_{h,\omega}A)=A^TA.
\]
Consequently, if \(U^TU=I_n\), then \(S_{h,\omega}U\) has orthonormal
columns.  This collision-free statement is exact arithmetic and deterministic;
it does not itself supply the probability of the injective event.

The formalization now also proves the exact pair-collision probability for
the hash law.  For distinct input rows \(a,b\in\{1,\ldots,m\}\), define
\[
  E_{a,b}=\{h\in\mathcal H_{r,m}:h(a)=h(b)\},
  \qquad
  E_{a,b}^{\rm nc}=\{h\in\mathcal H_{r,m}:h(a)\ne h(b)\}.
\]
Under the exact iid uniform hash law,
\[
  \Pr_h(E_{a,b})=\frac1r,
  \qquad
  \Pr_h(E_{a,b}^{\rm nc})=1-\frac1r .
\]
Thus the exact pair-collision probability has order
\[
  \Pr_h(E_{a,b})=\Theta(r^{-1})
\]
as \(r\to\infty\) with the two input rows fixed.  This is non-vacuous for
\(r>1\), where \(0<1-1/r<1\), and the collision probability tends to zero
as the number of buckets grows.  This probability statement concerns only
the exact sampling law; it does not compute, store, round, or apply the
sparse matrix \(S_{h,\omega}\).

The non-injective CountSketch second-moment route now has its exact
sign/hash-moment foundation.  The first new ingredient is the indicator
identity
\[
  \mathbb E_h {\bf 1}_{\{h(a)=h(b)\}}=\frac1r
  \qquad(a\ne b),
\]
and the same identity after marginalizing the independent exact Rademacher
sign law:
\[
  \mathbb E_{h,\Omega}{\bf 1}_{\{h(a)=h(b)\}}=\frac1r .
\]
The second ingredient is the exact fourth-moment classifier for ordered
distinct Rademacher sign pairs.  If \(a\ne b\) and \(c\ne d\), then
\[
\begin{aligned}
&\mathbb E_\Omega[
  \omega_a\omega_b\omega_c\omega_d]\\
&\quad =
\begin{cases}
1, & a=c \hbox{ and } b=d,\\
1, & a=d \hbox{ and } b=c,\\
0, & \hbox{otherwise}.
\end{cases}
\end{aligned}
\]
Equivalently, for any exact coefficient table
\(c:\mathcal P_m\to\R\), where
\(\mathcal P_m=\{(a,b):a\ne b\}\),
\[
\begin{aligned}
&\mathbb E_\Omega
  \left(
    \sum_{(a,b)\in\mathcal P_m} c_{ab}\omega_a\omega_b
  \right)^2\\
&=
\sum_{(a,b)\in\mathcal P_m}
\sum_{(c,d)\in\mathcal P_m}
c_{ab}c_{cd}
\left[
  {\bf 1}_{\{a=c,\ b=d\}}
  +
  {\bf 1}_{\{a=d,\ b=c\}}
\right].
\end{aligned}
\]
This exact formula is the reusable local substrate for a standard
non-collision-free CountSketch variance or Frobenius embedding proof: after
the CountSketch Gram error is expanded as an ordered collision sum, the hash
indicator mean contributes the \(1/r\) factor and the Rademacher
fourth-moment formula leaves only identical and reversed ordered pairs.  No
floating-point quantity is introduced in this layer, and no concentration
event is claimed yet.  The primitive probability factor has order
\(\Theta(r^{-1})\) for fixed input-row pair; the summed sign identity is exact
and remains non-vacuous for arbitrary exact coefficients, including the
degenerate zero-coefficient case.

The formalization also proves a global injectivity lower bound for the exact
hash law.  Let
\[
  I_{r,m}=\{h\in\mathcal H_{r,m}: h \hbox{ is injective}\},
  \qquad
  \mathcal P_m=\{(a,b):a,b\in\{1,\ldots,m\},\ a\ne b\}.
\]
The event \(I_{r,m}\) is exact and analysis-level: it is not computed by the
algorithm.  Lean proves that \(I_{r,m}\) is the same as the simultaneous
pairwise no-collision event over all \((a,b)\in\mathcal P_m\), and therefore
the pair-collision law and the finite union bound give
\[
  \Pr_h(I_{r,m})
  \ge
  1-\sum_{(a,b)\in\mathcal P_m}\frac1r .
\]
Using only the cardinality bound \(|\mathcal P_m|\le m^2\), Lean also proves
the simplified closed form
\[
  \Pr_h(I_{r,m})\ge 1-\frac{m^2}{r}.
\]
Equivalently, the proven failure radius is
\[
  1-\Pr_h(I_{r,m})\le \sum_{(a,b)\in\mathcal P_m}\frac1r
  \le \frac{m^2}{r}.
\]
In order notation, for varying \(m,r\) this union-bound failure radius is
\[
  O\!\left(\frac{m^2}{r}\right),
\]
and for fixed \(m\) it is \(O(r^{-1})\) as \(r\to\infty\).  This lower bound
is non-vacuous when \(r>m^2\), and it tends to one in the asymptotic regime
\(m^2/r\to0\).  The formalized theorem is one-sided, so the order statement is
big-\(O\), not a matching \(\Theta\) lower-and-upper estimate.

The formalization also proves the first concrete sparse floating-point apply
bound for CountSketch.  This is no longer an exact-arithmetic statement about
the table alone; it is an implementation-facing statement for one computed
output entry.  The hash \(h\), signs \(\omega\), and integer bucket selection
are exact, while the selected signed products and the subsequent bucket
accumulation are computed in floating point.

For an exact hash \(h\), define the exact bucket
\[
  B_i(h)=\{k\in\{1,\ldots,m\}:h(k)=i\},
  \qquad b_i=|B_i(h)|.
\]
Let \(k_1,\ldots,k_{b_i}\) be the canonical finite enumeration of this
bucket used by Lean.  Exact arithmetic gives
\[
  (S_{h,\omega}A)_{ij}
  =
  \sum_{t=1}^{b_i}\omega_{k_t}A_{k_tj}.
\]
The modeled sparse floating-point routine computes
\[
  \widehat z_t=\operatorname{fl}(\omega_{k_t}A_{k_tj})
\]
for each selected row, initializes \(\widehat y^{(0)}=0\), and then performs
left-to-right floating-point addition
\[
  \widehat y^{(t)}
  =
  \operatorname{fl}\!\left(\widehat y^{(t-1)}+\widehat z_t\right),
  \qquad t=1,\ldots,b_i,
  \qquad
  \widehat y_{ij}=\widehat y^{(b_i)}.
\]
When \(b_i=0\), the routine returns \(\widehat y_{ij}=0\).  If
\(\gamma_{b_i}=b_i u/(1-b_i u)\) is valid, Lean proves
\[
  |\widehat y_{ij}-(S_{h,\omega}A)_{ij}|
  \le
  \left(u+\gamma_{b_i}+u\gamma_{b_i}\right)
  \sum_{t=1}^{b_i}|\omega_{k_t}|\,|A_{k_tj}|.
\]
Under the exact sign-magnitude condition \(|\omega_k|\le 1\), this simplifies
to the readable column-one-norm bound
\[
  |\widehat y_{ij}-(S_{h,\omega}A)_{ij}|
  \le
  \left(u+\gamma_{b_i}+u\gamma_{b_i}\right)\|A_{:j}\|_1 .
\]
The coefficient has the irreducible expansion
\[
  u+\gamma_{b_i}+u\gamma_{b_i}
  =
  u+\frac{b_i u}{1-b_i u}
    +u\frac{b_i u}{1-b_i u}
  =
  u+\frac{b_i u(1+u)}{1-b_i u}.
\]
For \(b_i\ge1\), \(b_i u<1\), and small unit roundoff \(u\), this coefficient
has order
\[
  u+\frac{b_i u(1+u)}{1-b_i u}
  =
  \Theta(b_i u),
\]
so the sparse-entry error bound is
\[
  \Theta\!\left(
    b_i u
    \sum_{t=1}^{b_i}|\omega_{k_t}|\,|A_{k_tj}|
  \right),
\]
or, under \(|\omega_k|\le1\),
\[
  O\!\left(b_i u\,\|A_{:j}\|_1\right).
\]
For the empty-bucket case \(b_i=0\), the exact and computed entries are both
zero and the displayed sum-bound is exactly zero.  The result is non-vacuous:
for fixed bucket size and fixed exact input column, the right-hand side tends
to zero as \(u\to0\) while \(b_i u<1\).

At matrix level, Lean defines the computed sparse apply
\[
  \widehat Y=\operatorname{flCountSketchApply}_{h,\omega}(A),
  \qquad
  \widehat Y_{ij}=\widehat y_{ij},
\]
entrywise by the same bucket routine.  Therefore, if
\(\gamma_{b_i}\) is valid for every output bucket \(i\), Lean proves for all
\(i,j\)
\[
  |\widehat Y_{ij}-(S_{h,\omega}A)_{ij}|
  \le
  \left(u+\gamma_{b_i}+u\gamma_{b_i}\right)
  \sum_{t=1}^{b_i}|\omega_{k_t}|\,|A_{k_tj}|,
\]
and, under \(|\omega_k|\le1\),
\[
  |\widehat Y_{ij}-(S_{h,\omega}A)_{ij}|
  \le
  \left(u+\gamma_{b_i}+u\gamma_{b_i}\right)\|A_{:j}\|_1 .
\]
Thus the whole computed sparse matrix is now exposed entrywise; no additional
matrix-apply perturbation event is assumed.

Lean now also closes the deterministic floating-point arithmetic for the
sparse CountSketch Gram formed from this computed matrix.  Define the exact
sparse sketch
\[
  Y=S_{h,\omega}A,\qquad Y_{ij}=(S_{h,\omega}A)_{ij},
\]
and define the exact, nonnegative sparse-apply error radius
\[
  E_{ij}
  =
  \left(
    u+\frac{b_i u}{1-b_i u}
      +u\frac{b_i u}{1-b_i u}
  \right)
  \sum_{t=1}^{b_i}|\omega_{k_t}|\,|A_{k_tj}|.
\]
Equivalently,
\[
  E_{ij}
  =
  \left(
    u+\frac{b_i u(1+u)}{1-b_i u}
  \right)
  \sum_{t=1}^{b_i}|\omega_{k_t}|\,|A_{k_tj}|.
\]
The computation first forms
\(\widehat Y=\operatorname{flCountSketchApply}_{h,\omega}(A)\), as above, and
then forms the displayed Gram matrix by rounded length-\(r\) dot products:
\[
  \widehat G_{jk}
  =
  \operatorname{flDot}_r
    \left((\widehat Y_{1j},\ldots,\widehat Y_{rj}),
          (\widehat Y_{1k},\ldots,\widehat Y_{rk})\right).
\]
The exact reference matrix is
\[
  G_{jk}
  =
  (Y^T Y)_{jk}
  =
  \sum_{i=1}^r Y_{ij}Y_{ik}.
\]
If \(\gamma_r\) and all \(\gamma_{b_i}\) are valid, Lean proves
\[
  \|\widehat G-G\|_F
  \le
  T_{\rm CSGram}^{\rm fp},
\]
where the exact-only budget is
\[
\begin{aligned}
T_{\rm CSGram}^{\rm fp}
  &=
  \left\|
    \left[
      \gamma_r
      \sum_{i=1}^r
        (|Y_{ij}|+E_{ij})(|Y_{ik}|+E_{ik})
    \right]_{j,k}
  \right\|_F \\
  &\quad+
  \left\|
    \left[
      \sum_{i=1}^r
        \left(
          E_{ij}|Y_{ik}|+|Y_{ij}|E_{ik}+E_{ij}E_{ik}
        \right)
    \right]_{j,k}
  \right\|_F .
\end{aligned}
\]
In compact matrix notation, with absolute values taken entrywise, this is
\[
  T_{\rm CSGram}^{\rm fp}
  =
  \left\|\gamma_r(|Y|+E)^T(|Y|+E)\right\|_F
  +
  \left\|E^T|Y|+|Y|^TE+E^TE\right\|_F .
\]
This expression exposes the two floating-point sources separately: the first
term is rounded Gram-dot-product error, and the second term propagates the
already proved sparse-apply error into the Gram matrix.

For interpretation, define the exact nonnegative matrix
\[
  Q_{ij}
  =
  b_i\sum_{t=1}^{b_i}|\omega_{k_t}|\,|A_{k_tj}|.
\]
When \(r u<1\), \(b_i u<1\) for all buckets, \(u\) is small, and the displayed
leading matrix is not identically zero, the exact budget has leading order
\[
  T_{\rm CSGram}^{\rm fp}
  =
  \Theta\!\left(
    r u\,\bigl\||Y|^T|Y|\bigr\|_F
    +
    u\,\bigl\|Q^T|Y|+|Y|^TQ\bigr\|_F
  \right),
\]
with higher-order terms of order
\[
  O\!\left(
    r u^2\bigl\||Y|^TQ+Q^T|Y|\bigr\|_F
    +r u^3\|Q^TQ\|_F
    +u^2\|Q^TQ\|_F
  \right).
\]
The exact formula, not the asymptotic display, handles degenerate cases such
as \(Y=0\) or empty buckets.  The result is non-vacuous: for fixed exact
\(A,h,\omega\), fixed dimensions, and valid \(\gamma\)-values, the right-hand
side tends to zero as \(u\to0\).

The non-injective CountSketch route now has an actual Frobenius/Markov sparse
Gram endpoint, rather than only a moment substrate.  Let \(x=(h,\Omega)\) be
drawn from the exact CountSketch hash-sign product law, let
\(\omega=\operatorname{rademacherSignVector}(\Omega)\), and define the exact
sketched Gram error
\[
  \Delta_{\rm CS}(h,\omega;A)
  =
  (S_{h,\omega}A)^T(S_{h,\omega}A)-A^TA .
\]
For each Gram entry, Lean first proves the exact identity
\[
\begin{aligned}
  (\Delta_{\rm CS})_{jl}
  &=
  \sum_{(a,b)\in\mathcal P_m}
    {\bf 1}_{\{h(a)=h(b)\}}\,
    \omega_a\omega_b\,A_{aj}A_{bl},
  \qquad
  \mathcal P_m=\{(a,b):a\ne b\}.
\end{aligned}
\]
Using the fourth-moment classifier and exact hash collision law gives
\[
  \mathbb E_{h,\Omega}
  \bigl[(\Delta_{\rm CS})_{jl}^2\bigr]
  \le
  \frac{2}{r}
  \sum_{(a,b)\in\mathcal P_m}(A_{aj}A_{bl})^2 .
\]
Summing entries yields the exact Frobenius second-moment bound
\[
  \mathbb E_{h,\Omega}\|\Delta_{\rm CS}(h,\omega;A)\|_F^2
  \le
  \frac{2}{r}
  \sum_{j,l=1}^n
  \sum_{(a,b)\in\mathcal P_m}
    (A_{aj}A_{bl})^2 .
\]
The same exact sign/hash argument is now available at the fixed-vector
quadratic-form level.  For an exact test vector \(z\in\R^n\), define
\[
  v:=Az\in\R^m,\qquad v_a=\sum_{j=1}^n A_{aj}z_j .
\]
The formalization first proves the algebraic bridge
\[
  z^T\Delta_{\rm CS}(h,\omega;A)z
  =
  {\rm rowGram}\!\left(S_{h,\omega}v\right)_{11}
  -
  {\rm rowGram}(v)_{11},
\]
where \(v\) is regarded as a one-column matrix.  Therefore the already
formalized CountSketch entry second-moment theorem applies to the concrete
one-column input \(v=Az\), giving
\[
  \mathbb E_{h,\Omega}
  \left[
    \left(z^T\Delta_{\rm CS}(h,\omega;A)z\right)^2
  \right]
  \le
  \frac{2}{r}
  \sum_{(a,b)\in\mathcal P_m}(v_av_b)^2 .
\]
This is the exact coefficient form: every term is determined by the exact
input matrix \(A\), exact test vector \(z\), and exact hash/sign laws.  There
is no computed perturbation object and no certificate-existence assumption in
this exact-law layer.

Lean also proves the readable coefficient simplifications
\[
  \sum_{(a,b)\in\mathcal P_m}(v_av_b)^2
  \le
  \|v\|_2^4
  =
  \|Az\|_2^4
  \le
  \|A\|_F^4\|z\|_2^4.
\]
Consequently,
\[
  \mathbb E_{h,\Omega}
  \left[
    \left(z^T\Delta_{\rm CS}(h,\omega;A)z\right)^2
  \right]
  \le
  \frac{2\|Az\|_2^4}{r}
  \le
  \frac{2\|A\|_F^4\|z\|_2^4}{r}.
\]
For any exact threshold \(\eta>0\),
\[
\begin{aligned}
&\Pr_{h,\Omega}
  \left\{
    \left|z^T\Delta_{\rm CS}(h,\omega;A)z\right|
    \le \eta
  \right\}\\
&\qquad\ge
  1-
  \frac{2}{r\eta^2}
  \sum_{(a,b)\in\mathcal P_m}(v_av_b)^2 .
\end{aligned}
\]
Equivalently, if
\[
  \frac{2}{r\eta^2}
  \sum_{(a,b)\in\mathcal P_m}(v_av_b)^2
  \le \delta,
\]
then this fixed-vector quadratic-form event has probability at least
\(1-\delta\).  The irreducible order is
\[
  \Theta\!\left(
    \frac{\sum_{a\ne b}((Az)_a(Az)_b)^2}{r\eta^2}
  \right),
\]
and the readable orders are
\[
  O\!\left(\frac{\|Az\|_2^4}{r\eta^2}\right),
  \qquad
  O\!\left(\frac{\|A\|_F^4\|z\|_2^4}{r\eta^2}\right).
\]
When the displayed coefficient sum is nonzero and \(A,z,\eta\) are fixed,
the first expression is \(\Theta(r^{-1})\) as \(r\to\infty\); in the zero
case the exact proved loss is zero.  This fixed-vector theorem is a sharper
foundation for later finite-cover or subspace-embedding work, but it is not
yet a uniform-over-all-vectors CountSketch embedding theorem.

The fixed-vector statement is also lifted to an arbitrary finite family of
exact test vectors.  Let \(\mathcal T\) be a finite index set, let
\(z_\alpha\in\R^n\) and \(\eta_\alpha>0\) for \(\alpha\in\mathcal T\), and set
\[
  v^\alpha=A z_\alpha,\qquad
  B_\alpha
  =
  \frac{2}{r\eta_\alpha^2}
  \sum_{a\ne b}
    \bigl(v^\alpha_a v^\alpha_b\bigr)^2 .
\]
Then Lean proves
\[
\begin{aligned}
&\Pr_{h,\Omega}
  \left\{
    \forall \alpha\in\mathcal T,\quad
    \left|z_\alpha^T\Delta_{\rm CS}(h,\omega;A)z_\alpha\right|
    \le\eta_\alpha
  \right\}\\
&\qquad\ge
  1-\sum_{\alpha\in\mathcal T} B_\alpha .
\end{aligned}
\]
Equivalently, if \(\sum_{\alpha}B_\alpha\le\delta\), then the same finite-test
event has probability at least \(1-\delta\).  This is a finite union-bound
theorem, not a new probability assumption.  Its irreducible loss is
\[
  \Theta\!\left(
    \sum_{\alpha\in\mathcal T}
      \frac{\sum_{a\ne b}((Az_\alpha)_a(Az_\alpha)_b)^2}
           {r\eta_\alpha^2}
  \right)
\]
when the displayed numerator is nonzero, with readable upper forms
\[
  O\!\left(
    \sum_{\alpha\in\mathcal T}
      \frac{\|Az_\alpha\|_2^4}{r\eta_\alpha^2}
  \right),
  \qquad
  O\!\left(
    \sum_{\alpha\in\mathcal T}
      \frac{\|A\|_F^4\|z_\alpha\|_2^4}{r\eta_\alpha^2}
  \right).
\]
The event remains exact-law and exact-arithmetic.  The finite set
\(\mathcal T\) is an analysis test set; applying this theorem to a net or cover
requires a separate deterministic cover-to-Loewner step.

That deterministic cover step is now formalized without a supplied
operator-radius certificate.  Suppose the same finite family
\(\{z_\alpha\}_{\alpha\in\mathcal T}\) is a Euclidean unit-ball cover of
\(\R^n\) with radius \(\rho\ge0\).  Choose common thresholds \(\eta>0\) and
\(L>0\), and define
\[
  \tau_{\rm cover}=\eta+L(2\rho+\rho^2).
\]
The event
\[
  -\tau_{\rm cover}I
  \preceq
  \Delta_{\rm CS}(h,\omega;A)
  \preceq
  \tau_{\rm cover}I
\]
has probability at least
\[
\begin{aligned}
1
&-
\sum_{\alpha\in\mathcal T}
  \frac{2}{r\eta^2}
  \sum_{a\ne b}
    \bigl((Az_\alpha)_a(Az_\alpha)_b\bigr)^2\\
&-
\frac{2}{rL^2}
  \sum_{j,\ell}
  \sum_{a\ne b}
    \bigl(A_{aj}A_{b\ell}\bigr)^2 .
\end{aligned}
\]
Equivalently, if the two displayed losses sum to at most \(\delta\), then the
same two-sided finite-Loewner event has probability at least \(1-\delta\).
Lean also proves the readable sufficient exact-cover bound
\[
  \Pr\{-\tau_{\rm cover}I\preceq\Delta_{\rm CS}\preceq\tau_{\rm cover}I\}
  \ge
  1-
  \left(
    \sum_{\alpha\in\mathcal T}
      \frac{2\|Az_\alpha\|_2^4}{r\eta^2}
    +
      \frac{2\|A\|_F^4}{rL^2}
  \right),
\]
plus the corresponding direct \(1-\delta\) wrapper for that readable loss.
The Lean names are
\[
\texttt{countSketchProbability\_eventProb\_rowGram\_twoSidedLoewner\_cover\_ge\_one\_sub\_sum\_vecNorm\_add\_frobNorm}
\]
and
\[
\texttt{countSketchProbability\_eventProb\_rowGram\_twoSidedLoewner\_cover\_ge\_one\_sub\_delta\_of\_sum\_vecNorm\_add\_frobNorm\_budget}.
\]
The proof intersects the finite-test event with the exact Frobenius event
\(\|\Delta_{\rm CS}(h,\omega;A)\|_F\le L\); the Frobenius event supplies the
coarse radius required by the deterministic cover lemma.  Thus no hidden
``there exists a coarse operator certificate'' assumption is introduced.

The irreducible failure order is
\[
  \Theta\!\left(
    \sum_{\alpha\in\mathcal T}
      \frac{\sum_{a\ne b}((Az_\alpha)_a(Az_\alpha)_b)^2}
           {r\eta^2}
    +
      \frac{\sum_{j,\ell}\sum_{a\ne b}(A_{aj}A_{b\ell})^2}
           {rL^2}
  \right)
\]
when the displayed numerator is nonzero, with readable upper order
\[
  O\!\left(
    \sum_{\alpha\in\mathcal T}
      \frac{\|Az_\alpha\|_2^4}{r\eta^2}
    +
      \frac{\|A\|_F^4}{rL^2}
  \right).
\]
For fixed \(A\), fixed cover, and fixed positive \(\eta,L\), this loss tends to
zero like \(\Theta(r^{-1})\) whenever the numerator is nonzero, and is exactly
zero when both coefficient sums vanish.  This is still exact-law and
exact-arithmetic infrastructure; it is not a floating-point endpoint.

The supplied-cover version has a constructive product-grid specialization.
Let \(\mathcal I\) be a finite exact index set and let
\(g:\mathcal I\to\R\) be an exact one-dimensional grid.  Suppose \(g\) covers
the interval \([-1,1]\) with radius \(\delta_{\rm grid}\), with
\(\delta_{\rm grid}\ge0\), and suppose
\[
  \sqrt n\,\delta_{\rm grid}\le\rho .
\]
For each \(a:\{1,\ldots,n\}\to\mathcal I\), define the exact product-grid
vector
\[
  z_a(j)=g(a(j)),\qquad j=1,\ldots,n .
\]
The Lean theorems
\[
\texttt{finiteUnitBallCover\_of\_rectUnitBallCover},
\qquad
\texttt{finiteUnitBallCover\_product\_grid}
\]
prove that \(\{z_a\}_{a:\{1,\ldots,n\}\to\mathcal I}\) is an exact
unit-ball cover of radius \(\rho\).  Its index cardinality is
\[
  \#\{a:\{1,\ldots,n\}\to\mathcal I\}=|\mathcal I|^n,
\]
as formalized by
\[
\texttt{fintype\_card\_product\_grid\_index}.
\]
Instantiating the CountSketch finite-cover theorem with this explicit grid
gives
\[
  \Pr\{-\tau_{\rm cover}I\preceq\Delta_{\rm CS}\preceq\tau_{\rm cover}I\}
  \ge
  1-
  \left(
    \sum_{a:\{1,\ldots,n\}\to\mathcal I}
      \frac{2\|Az_a\|_2^4}{r\eta^2}
    +
      \frac{2\|A\|_F^4}{rL^2}
  \right),
\]
where \(\tau_{\rm cover}=\eta+L(2\rho+\rho^2)\).  The corresponding direct
target-budget theorem replaces the parenthesized loss by any
\(\delta\) that upper-bounds it.  The Lean theorem names are
\[
\texttt{countSketchProbability\_eventProb\_rowGram\_twoSidedLoewner\_productGrid\_ge\_one\_sub\_sum\_vecNorm\_add\_frobNorm}
\]
and
\[
\texttt{countSketchProbability\_eventProb\_rowGram\_twoSidedLoewner\_productGrid\_ge\_one\_sub\_delta\_of\_sum\_vecNorm\_add\_frobNorm\_budget}.
\]
For fixed exact \(A\), fixed grid \(g\), and fixed positive \(\eta,L\), the
product-grid loss has readable order
\[
  \Theta\!\left(
    \frac{
      \sum_{a:\{1,\ldots,n\}\to\mathcal I}\|Az_a\|_2^4/\eta^2
      +
      \|A\|_F^4/L^2
    }{r}
  \right)
\]
when the numerator is nonzero, and is exactly zero when the numerator is zero.
The grid and its cardinality are exact analysis objects; no probability table
or floating-point cover construction is computed by Algorithm 3.

Therefore, for any exact threshold \(\eta>0\),
\[
\begin{aligned}
&\Pr_{h,\Omega}
  \left\{
    \|\Delta_{\rm CS}(h,\omega;A)\|_F\le\eta
  \right\} \\
&\qquad\ge
  1-
  \frac{2}{r\eta^2}
  \sum_{j,l=1}^n
  \sum_{(a,b)\in\mathcal P_m}
    (A_{aj}A_{bl})^2 .
\end{aligned}
\]
This exact-object result is then transferred to the computed sparse Gram
\[
  \widehat G(h,\Omega)
  =
  \operatorname{fl}\!\left(
    \operatorname{flCountSketchApply}_{h,\omega}(A)^T
    \operatorname{flCountSketchApply}_{h,\omega}(A)
  \right).
\]
The computed event is
\[
  \left\{
    (h,\Omega):
    \|\widehat G(h,\Omega)-A^TA\|_F
    \le
    \eta+
    T_{\rm CSGram}^{\rm fp}(h,\omega,A)
  \right\}.
\]
It has the same probability lower bound as the exact event.  The proof is a
deterministic triangle inequality plus the already formalized concrete sparse
Gram arithmetic theorem
\(\|\widehat G-(S_{h,\omega}A)^T(S_{h,\omega}A)\|_F
\le T_{\rm CSGram}^{\rm fp}(h,\omega,A)\); no perturbation-existence
assumption or collision-free event is used.

The CountSketch sign-storage arithmetic is also now formalized as a
deterministic implementation substrate.  Let \(\widehat\omega\) be any computed
table certified against the exact realized signs \(\omega\), with entrywise
storage radius \(e_k\).  If the sparse apply uses \(\widehat\omega_k\) in the
rounded products, then the entrywise sparse-apply radius is
\[
  E^{{\rm CS},{\rm store}}_{ij}
  =
  (u+\gamma_{b_i}+u\gamma_{b_i})
  \sum_{k:h(k)=i}|\widehat\omega_k|\,|A_{kj}|
  +
  \sum_{k:h(k)=i} e_k |A_{kj}|.
\]
The second term is the exact storage deviation back to the ideal
\(S_{h,\omega}A\).  Theorems
\[
\texttt{preconditionRows\_countSketchRows\_storedSign\_entry\_error\_bound},
\quad
\texttt{fl\_countSketchSparseApplyEntry\_withStoredSign\_error\_bound},
\]
\[
\texttt{fl\_countSketchSparseApplyWithStoredSign\_entry\_error\_bound},
\quad
\texttt{fl\_countSketchSparseGramDotWithStoredSign\_perturb\_bound}
\]
thread this radius through the same rounded length-\(r\) Gram-dot interface.
For the concrete rounded-copy sign-storage modes
\[
  \widehat\omega_k\in
  \{\operatorname{fl}(\omega_k\cdot1),
    \operatorname{fl}(\omega_k+0),
    \operatorname{fl}(\omega_k-0)\},
\]
theorems
\[
\texttt{fl\_countSketchSparseGramDotWithFlStoredSign\_perturb\_bound},
\quad
\texttt{fl\_countSketchSparseGramDotWithFlStoredSignAddZeroRight\_perturb\_bound},
\quad
\texttt{fl\_countSketchSparseGramDotWithFlStoredSignSubZeroRight\_perturb\_bound}
\]
instantiate \(e_k=u\) using \(|\omega_k|=1\).  In the regime of fixed realized
\(h,\omega,A\) and small roundoff, this stored-sign radius has order
\[
  \Theta\!\left(
    \gamma_r\||Y|^T|Y|\|_F+
    \|(E^{{\rm CS},{\rm store}})^T|Y|
      +|Y|^TE^{{\rm CS},{\rm store}}\|_F
  \right),
  \qquad
  E^{{\rm CS},{\rm store}}_{ij}
  =
  \Theta\!\left((u+\gamma_{b_i})
  \sum_{k:h(k)=i}|A_{kj}|\right)
\]
for the three displayed storage modes, and it tends to zero as
\(u,\gamma_r,\gamma_{b_i}\to0\).

The stored-sign radius has now also been threaded through the product-grid
finite-cover probability theorem.  The event
\[
  -\bigl(\eta+L(2\rho+\rho^2)+
    T_{{\rm CSGram},{\rm store}}^{\rm fp}(h,\omega,A)\bigr)I
  \preceq
  \widehat G^{\rm store}(h,\omega,A)-A^TA
  \preceq
  \bigl(\eta+L(2\rho+\rho^2)+
    T_{{\rm CSGram},{\rm store}}^{\rm fp}(h,\omega,A)\bigr)I
\]
has the same exact-law product-grid lower bound as the exact-sign computed
sparse-Gram event:
\[
  1-
  \left(
    \sum_{a:\mathrm{Fin}\,n\to\alpha}
      \frac{2\|A z_a\|_2^4}{r\eta^2}
    +
    \frac{2\|A\|_F^4}{rL^2}
  \right).
\]
Lean records the generic transfer as
\texttt{countSketchRowGramTwoSidedLoewnerCoverEvent\_subset\_flSparseGramDotWithStoredSignRowGramTwoSidedLoewnerEvent}
and the product-grid wrappers as
\texttt{countSketchProbability\_eventProb\_flSparseGramDotWithStoredSign\_rowGram\_twoSidedLoewner\_productGrid\_ge\_one\_sub\_sum\_vecNorm\_add\_frobNorm}
and
\texttt{countSketchProbability\_eventProb\_flSparseGramDotWithStoredSign\_rowGram\_twoSidedLoewner\_productGrid\_ge\_one\_sub\_delta\_of\_sum\_vecNorm\_add\_frobNorm\_budget}.
The final concrete sign-copy endpoints are
\texttt{countSketchProbability\_eventProb\_flSparseGramDotWithFlStoredSign\_rowGram\_twoSidedLoewner\_productGrid\_ge\_one\_sub\_delta\_of\_sum\_vecNorm\_add\_frobNorm\_budget},
\texttt{countSketchProbability\_eventProb\_flSparseGramDotWithFlStoredSignAddZeroRight\_rowGram\_twoSidedLoewner\_productGrid\_ge\_one\_sub\_delta\_of\_sum\_vecNorm\_add\_frobNorm\_budget},
and
\texttt{countSketchProbability\_eventProb\_flSparseGramDotWithFlStoredSignSubZeroRight\_rowGram\_twoSidedLoewner\_productGrid\_ge\_one\_sub\_delta\_of\_sum\_vecNorm\_add\_frobNorm\_budget}.
The probability-loss order remains
\[
  \Theta\!\left(
    \frac{
      \sum_a\|A z_a\|_2^4/\eta^2+\|A\|_F^4/L^2
    }{r}
  \right)
\]
in the fixed exact-input/grid/threshold regime when the numerator is nonzero;
stored sign copies add only the deterministic roundoff radius, not a new
probability-loss term.

The same product-grid stored-sign theorem is now closed for any fixed
per-bucket memory order.  To avoid overloading notation in this paragraph,
write \(\rho_{\rm grid}\) for the product-grid cover radius and
\(\rho_{\rm sig}\) for the selected sign-copy mode.  For each realized exact
hash \(h\), write
\[
  B_i(h)=\{k:h(k)=i\},\qquad b_i(h)=|B_i(h)|,
\]
and let
\[
  k_{i,1},\ldots,k_{i,b_i(h)}
\]
be an exact ordering of this bucket.  This order is a discrete index choice;
it is not a floating-point real computation.  For a generic rounded-copy mode
\(\rho\), the permuted sparse-apply entry radius is
\[
\begin{aligned}
  E^{{\rm CS},\rho,\pi}_{ij}
  &=
  (u+\gamma_{b_i(h)}+u\gamma_{b_i(h)})
  \sum_{t=1}^{b_i(h)}
    |\widehat\omega^\rho_{k_{i,t}}|\,|A_{k_{i,t}j}|
  +
  u\sum_{t=1}^{b_i(h)}|A_{k_{i,t}j}|.
\end{aligned}
\]
The exact reference remains
\[
  Y_{ij}=\sum_{k:h(k)=i}\omega_kA_{kj}.
\]
The resulting deterministic Gram radius is
\[
\begin{aligned}
  T_{{\rm CSGram},\rho,\pi}^{\rm fp}
  &:=
  \left\|
    \left(
      \gamma_r\sum_{i=1}^r
      (|Y_{ij}|+E^{{\rm CS},\rho,\pi}_{ij})
      (|Y_{ik}|+E^{{\rm CS},\rho,\pi}_{ik})
    \right)_{j,k}
  \right\|_F \\
  &\quad+
  \left\|
    \left(
      \sum_{i=1}^r
      \left[
        E^{{\rm CS},\rho,\pi}_{ij}|Y_{ik}|
        +|Y_{ij}|E^{{\rm CS},\rho,\pi}_{ik}
        +E^{{\rm CS},\rho,\pi}_{ij}E^{{\rm CS},\rho,\pi}_{ik}
      \right]
    \right)_{j,k}
  \right\|_F .
\end{aligned}
\]
Consequently the event
\[
  -\bigl(\eta+L(2\rho_{\rm grid}+\rho_{\rm grid}^2)+
    T_{{\rm CSGram},\rho_{\rm sig},\pi}^{\rm fp}\bigr)I
  \preceq
  \widehat G^{\rho_{\rm sig},\pi}-A^TA
  \preceq
  \bigl(\eta+L(2\rho_{\rm grid}+\rho_{\rm grid}^2)+
    T_{{\rm CSGram},\rho_{\rm sig},\pi}^{\rm fp}\bigr)I
\]
has the same exact product-grid lower bound
\[
  1-
  \left(
    \sum_{a:\mathrm{Fin}\,n\to\alpha}
      \frac{2\|A z_a\|_2^4}{r\eta^2}
    +
    \frac{2\|A\|_F^4}{rL^2}
  \right).
\]
Lean records this endpoint through
\texttt{countSketchFlSparseGramDotWithStoredSignPermutedRowGramTwoSidedLoewnerEvent},
\texttt{countSketchRowGramTwoSidedLoewnerCoverEvent\_subset\_flSparseGramDotWithStoredSignPermutedRowGramTwoSidedLoewnerEvent},
\texttt{countSketchProbability\_eventProb\_flSparseGramDotWithStoredSignPermuted\_rowGram\_twoSidedLoewner\_productGrid\_ge\_one\_sub\_sum\_vecNorm\_add\_frobNorm},
and its target-budget and three concrete stored-sign copy wrappers.  For fixed
realized \(h,\omega,\pi,A\),
\[
  E^{{\rm CS},\rho,\pi}_{ij}
  =
  \Theta\!\left(
    (u+\gamma_{b_i(h)})
    \sum_{t=1}^{b_i(h)}|A_{k_{i,t}j}|
  \right),
\]
and, whenever the displayed leading term is nonzero,
\[
  T_{{\rm CSGram},\rho,\pi}^{\rm fp}
  =
  \Theta\!\left(
    \gamma_r\||Y|^T|Y|\|_F+
    \|(E^{{\rm CS},\rho,\pi})^T|Y|
      +|Y|^TE^{{\rm CS},\rho,\pi}\|_F
  \right).
\]
The radius tends to zero with \(u,\gamma_r,\gamma_{b_i}\) for fixed realized
data.  This closes fixed per-bucket permutation layouts; parallel
reassociation and aliasing semantics remain separate implementation routines.

The same stored-sign product-grid sparse-Gram endpoint is now closed for
explicit tree-reduced bucket accumulation.  For each realized hash \(h\), the
tree shapes \(\mathcal T_i(h)\) are exact discrete algorithm choices, with
depths \(d_i(h)\), and the computed sparse sketch stores the realized signs,
rounds each signed product, reduces each bucket through its supplied binary
tree, and finally forms rounded length-\(r\) Gram dot products.  The event
radius is
\[
  \eta+L(2\rho+\rho^2)
  +T_{{\rm CSGram},\rho,\mathcal T}^{\rm fp}(h,\omega,A),
\]
where the entry budget expands to
\[
  E^{{\rm CS},\rho,\mathcal T}_{ij}
  =
  (u+\gamma_{d_i(h)}+u\gamma_{d_i(h)})
  \sum_{t=1}^{b_i(h)}
    |\widehat\omega^\rho_{k_{i,t}}|\,|A_{k_{i,t}j}|
  +
  \sum_{t=1}^{b_i(h)}e^\rho_{k_{i,t}}|A_{k_{i,t}j}|,
\]
with \(e^\rho_k\le u\) for the three concrete sign-copy modes
\(\fl(\omega_k\cdot1)\), \(\fl(\omega_k+0)\), and \(\fl(\omega_k-0)\).  The
Gram radius keeps the same expanded shape
\[
\begin{aligned}
  T_{{\rm CSGram},\rho,\mathcal T}^{\rm fp}
  &=
  \left\|
    \left(
      \gamma_r
      \sum_i(|Y_{ij}|+E^{{\rm CS},\rho,\mathcal T}_{ij})
             (|Y_{ik}|+E^{{\rm CS},\rho,\mathcal T}_{ik})
    \right)_{j,k}
  \right\|_F\\
  &\quad+
  \left\|
    \left(
      \sum_i
      E^{{\rm CS},\rho,\mathcal T}_{ij}|Y_{ik}|
      +|Y_{ij}|E^{{\rm CS},\rho,\mathcal T}_{ik}
      +E^{{\rm CS},\rho,\mathcal T}_{ij}
       E^{{\rm CS},\rho,\mathcal T}_{ik}
    \right)_{j,k}
  \right\|_F .
\end{aligned}
\]
Lean records this endpoint through
\texttt{fl\_countSketchSparseApplyEntryTree\_error\_bound},
\texttt{fl\_countSketchSparseGramDotTree\_perturb\_bound},
\texttt{fl\_countSketchSparseApplyWithStoredSignTree\_entry\_error\_bound},
\texttt{fl\_countSketchSparseGramDotWithStoredSignTree\_perturb\_bound},
\texttt{countSketchFlSparseGramDotWithStoredSignTreeRowGramTwoSidedLoewnerEvent},
\texttt{countSketchRowGramTwoSidedLoewnerCoverEvent\_subset\_flSparseGramDotWithStoredSignTreeRowGramTwoSidedLoewnerEvent},
\texttt{countSketchProbability\_eventProb\_flSparseGramDotWithStoredSignTree\_rowGram\_twoSidedLoewner\_productGrid\_ge\_one\_sub\_sum\_vecNorm\_add\_frobNorm},
and its target-budget and three concrete stored-sign copy wrappers.  For fixed
realized data,
\[
  E^{{\rm CS},\rho,\mathcal T}_{ij}
  =
  \Theta\!\left(
    (u+\gamma_{d_i(h)})
    \sum_{t=1}^{b_i(h)}|A_{k_{i,t}j}|
  \right),
\]
and balanced trees give the interpretable specialization
\[
  E^{{\rm CS},\rho,\mathcal T}_{ij}
  =
  \Theta\!\left(
    u(1+\log(b_i(h)+1))
    \sum_{t=1}^{b_i(h)}|A_{k_{i,t}j}|
  \right).
\]
The probability loss is unchanged from the stored-sign product-grid theorem,
\(\Theta((\sum_a\|Az_a\|_2^4/\eta^2+\|A\|_F^4/L^2)/r)\), and the deterministic
radius tends to zero with
\(u,\gamma_r,\gamma_{d_i(h)},e_k^\rho\).  This closes tree-reduced sparse-Gram
CountSketch layouts; the downstream uniform-row tree-reduced composition is
closed below.  Aliasing or parallel sparse-sketch memory layouts remain
separate implementation routines.

The same stored-sign sparse-apply radius has now been lifted through the
downstream iid uniform-row sampling stage.  The final event uses the concrete
computed denominator for \(\sqrt{s/r}\), rounded row divisions, and rounded
length-\(s\) sampled-Gram dot products, with radius
\[
  \eta+L(2\rho+\rho^2)+\eta_{\rm row}
  +T_{{\rm CS,row},{\rm store}}^{\rm fp}(h,\omega,\sigma,A).
\]
The exact product-grid plus row-sampling loss is unchanged from the exact-sign
downstream theorem:
\[
  \sum_a
    \frac{2r^{-1}\sum_{p\ne q}((Az_a)_p(Az_a)_q)^2}{\eta^2}
  +
  \frac{2r^{-1}\sum_{j,k}\sum_{p\ne q}(A_{pj}A_{qk})^2}{L^2}
  +
  \frac{r m^2\|A\|_F^4}{s\eta_{\rm row}^2}.
\]
Lean records the generic stored-sign downstream wrappers as
\texttt{countSketchUniformRowTraceProbability\_eventProb\_sparseStoredSignComputedPreconditioned\_fl\_uniformRowSampleGramDotWithComputedDen\_rowGram\_twoSidedLoewner\_productGrid\_ge\_one\_sub\_sum\_coeff\_add\_frob\_add\_row}
and its target-budget variant, and the concrete final endpoints for
\(\fl(\omega_i\cdot1)\), \(\fl(\omega_i+0)\), and \(\fl(\omega_i-0)\) are the
corresponding
\texttt{sparseFlStoredSignComputedPreconditioned},
\texttt{sparseFlStoredSignAddZeroRightComputedPreconditioned}, and
\texttt{sparseFlStoredSignSubZeroRightComputedPreconditioned} theorems.  These
are final implementation-facing results for the modeled stored-sign
CountSketch-plus-uniform-row path; no probability-computation term,
perturbation event, or certificate-existence assumption is added.

The stored-sign downstream theorem is also closed for fixed exact per-bucket
permutation layouts.  For each realized hash \(h\), the exact bucket order
\(\pi\) is a discrete index operation, while the computed sparse sketch visits
the bucket in that order.  The final event radius is
\[
  \eta+L(2\rho+\rho^2)+\eta_{\rm row}
  +T_{{\rm CS,row},{\rm store},\pi}^{\rm fp}(h,\omega,\sigma,A),
\]
and the exact product-grid plus row-sampling loss remains the same as in the
canonical stored-sign downstream theorem.  Lean records the new basis and
budget through
\texttt{countSketchSparseComputedPreconditionedBasisWithStoredSignPermuted}
and
\texttt{countSketchSparseUniformRowComputedDenStoredSignPermutedPerturbBudget},
the probability-one perturbation theorem through
\texttt{countSketchUniformRowTraceProbability\_eventProb\_sparseStoredSignPermutedComputedPreconditioned\_fl\_uniformRowComputedDenPerturbEvent\_eq\_one},
and the final product-grid wrappers through the
\texttt{sparseStoredSignPermutedComputedPreconditioned} and concrete
\texttt{sparseFlStoredSignPermutedComputedPreconditioned},
\texttt{sparseFlStoredSignAddZeroRightPermutedComputedPreconditioned}, and
\texttt{sparseFlStoredSignSubZeroRightPermutedComputedPreconditioned}
theorems.  For fixed realized \(h,\omega,\sigma,\pi,A\) and \(d>0\),
\[
  T_{{\rm CS,row},{\rm store},\pi}^{\rm fp}
  =
  \Theta\!\left(
    \gamma_s\Frob{|\widehat B|^T|\widehat B|}
    +
    \Frob{(E^{\rm den,\pi})^T|\widehat B|
      +|\widetilde B|^TE^{\rm den,\pi}}
    +
    \Frob{(F^\pi)^T|\widetilde B|+|B|^TF^\pi}
  \right),
\]
with
\[
  E^{\rm den,\pi}_{qj}
  =
  \Theta\!\left(
    u|\widehat B_{qj}|
    +
    \frac{|\widehat Y^\pi_{\sigma_qj}|}{|\widehat d|}
    \frac{4u+3u^2+u^3}{1-u}
  \right),
  \qquad
  F^\pi_{qj}
  =
  \Theta\!\left(
    \frac{u+\gamma_{b_{\sigma_q}(h)}}{d}
    \sum_t |A_{k_{\sigma_q,t}j}|
  \right).
\]
Thus the exact bucket order adds no probability loss and no probability-law
perturbation; it only selects the concrete order charged by the deterministic
floating-point accumulation radius.  Parallel reassociation and aliasing
sparse-sketch layouts remain separate implementation routines.

The stored-sign downstream theorem is also closed for explicit binary
tree-reduced bucket accumulation.  For each realized hash \(h\), the exact
tree shapes \(\mathcal T_i(h)\) and their trailing zero leaves are discrete
implementation choices, while the computed sparse sketch stores signs, rounds
signed products, reduces each bucket through the selected tree, forms the
concrete denominator \(\widehat d_\star=\fl(\fl(\sqrt{s})\fl(1/\fl(\sqrt r)))\),
rounds sampled-row divisions, and rounds length-\(s\) sampled-Gram dot
products.  Lean records the basis and budget through
\texttt{countSketchSparseComputedPreconditionedBasisWithStoredSignTree}
and
\texttt{countSketchSparseUniformRowComputedDenStoredSignTreePerturbBudget},
the probability-one perturbation theorem through
\texttt{countSketchUniformRowTraceProbability\_eventProb\_sparseStoredSignTreeComputedPreconditioned\_fl\_uniformRowComputedDenPerturbEvent\_eq\_one},
and the final product-grid wrappers through the
\texttt{sparseStoredSignTreeComputedPreconditioned} and concrete
\texttt{sparseFlStoredSignTreeComputedPreconditioned},
\texttt{sparseFlStoredSignAddZeroRightTreeComputedPreconditioned}, and
\texttt{sparseFlStoredSignSubZeroRightTreeComputedPreconditioned} theorems.
The final radius is
\[
  \eta+L(2\rho+\rho^2)+\eta_{\rm row}
  +T_{{\rm CS,row},{\rm store},\mathcal T}^{\rm fp}(h,\omega,\sigma,A),
\]
and the exact product-grid plus row-sampling loss is unchanged.  For fixed
realized data and \(d=\sqrt{s/r}>0\),
\[
  T_{{\rm CS,row},{\rm store},\mathcal T}^{\rm fp}
  =
  \Theta\!\left(
    \gamma_s\Frob{|\widehat B|^T|\widehat B|}
    +
    \Frob{(E^{\rm den,\mathcal T})^T|\widehat B|
      +|\widetilde B|^TE^{\rm den,\mathcal T}}
    +
    \Frob{(F^{\mathcal T})^T|\widetilde B|+|B|^TF^{\mathcal T}}
  \right),
\]
where
\[
  E^{\rm den,\mathcal T}_{qj}
  =
  \Theta\!\left(
    u|\widehat B_{qj}|
    +
    \frac{|\widehat Y^{\mathcal T}_{\sigma_qj}|}{|\widehat d_\star|}
    \frac{4u+3u^2+u^3}{1-u}
  \right),
\]
and
\[
  F^{\mathcal T}_{qj}
  =
  \Theta\!\left(
    \frac{u+\gamma_{d_{\sigma_q}(h)}}{d}
    \sum_t |A_{k_{\sigma_q,t}j}|
  \right).
\]
Balanced bucket trees give
\[
  F^{\mathcal T}_{qj}
  =
  \Theta\!\left(
    \frac{u(1+\log(b_{\sigma_q}(h)+1))}{d}
    \sum_t |A_{k_{\sigma_q,t}j}|
  \right).
\]
Exact hash/sign laws, exact row sampling, exact product-grid vectors, tree
shapes, and \(A^TA\) remain exact analysis/probability or discrete objects.
All computed non-probability objects on this modeled path are charged, and no
perturbation-event or certificate-existence assumption is used.

The same stored-sign downstream product-grid theorem is now instantiated with
four additional concrete denominator routines for \(d=\sqrt{s/r}\).  The
routes are
\[
\begin{array}{@{}l@{}}
\widehat d=\fl_{\rm sqrt}(s/r)
  \quad\text{when the ratio \(s/r\) is supplied exactly},\\
\widehat d=\fl_{\rm sqrt}(\fl_{\rm div}(s,r)),\\
\widehat d=\fl_{\rm sqrt}(\fl_{\rm mul}(s,\fl_{\rm div}(1,r))),\\
\widehat d=\fl_{\rm div}(\fl_{\rm sqrt}(s),\fl_{\rm sqrt}(r)).
\end{array}
\]
Their denominator relative radii are respectively
\[
  u,\qquad
  \sqrt{1+u}\,u+u,\qquad
  \sqrt{(1+u)(1+u)}\,u+2u+u^2,\qquad
  \frac{3u+u^2}{1-u}.
\]
Thus the exact product-grid plus downstream row probability loss is unchanged,
and the deterministic stored-sign row radius is obtained by substituting the
chosen route-specific denominator factor.  For fixed realized
\((h,\omega,\sigma,A)\) and \(d>0\), the stored-sign downstream floating-point
radius has order
\[
  \Theta\!\left(
    \gamma_s\Frob{|\widehat B|^T|\widehat B|}
    +
    \Frob{(E^{\rm den})^T|\widehat B|
      +|\widetilde B|^TE^{\rm den}}
    +
    \Frob{F^T|\widetilde B|+|B|^TF}
  \right),
\]
with
\[
  E^{\rm den}_{qj}
  =
  \Theta\!\left(
    u|\widehat B_{qj}|
    +
    \frac{|\widehat Y_{\sigma_qj}|}{|\widehat d|}\Delta_d
  \right),
\]
where \(\Delta_d\) is the selected relative denominator factor.  The radius
tends to zero as \(u,\gamma_s,\gamma_{b_i(h)},\Delta_d\to0\); denominator
computation changes only the deterministic radius, not the exact probability
law.

The coefficient sum is also simplified in Lean:
\[
  \sum_{j,l=1}^n
  \sum_{(a,b)\in\mathcal P_m}
    (A_{aj}A_{bl})^2
  \le
  \|A\|_F^4 .
\]
Thus the readable computed sparse Gram endpoint is
\[
  \Pr_{h,\Omega}
  \left\{
    \|\widehat G(h,\Omega)-A^TA\|_F
    \le
    \eta+T_{\rm CSGram}^{\rm fp}(h,\omega,A)
  \right\}
  \ge
  1-\frac{2\|A\|_F^4}{r\eta^2}.
\]
In the asymptotic regime with exact \(A\) and \(\eta>0\) fixed and
\(r\to\infty\), the failure term has order
\[
  \Theta\!\left(\frac{\|A\|_F^4}{r\eta^2}\right).
\]
For an exact orthonormal analysis basis \(U\in\R^{m\times n}\), \(U^TU=I_n\)
implies \(\|U\|_F^2=n\), and the bound becomes
\[
  \Pr_{h,\Omega}
  \left\{
    \|\widehat G(h,\Omega)-I_n\|_F
    \le
    \eta+T_{\rm CSGram}^{\rm fp}(h,\omega,U)
  \right\}
  \ge
  1-\frac{2n^2}{r\eta^2},
\]
with failure order
\[
  \Theta\!\left(\frac{n^2}{r\eta^2}\right)
\]
when \(n\) and \(\eta\) are fixed and \(r\to\infty\).  The result is
non-vacuous when \(r>2\|A\|_F^4/\eta^2\), or \(r>2n^2/\eta^2\) in the
orthonormal case, and the failure term tends to zero as \(r\to\infty\).
This is a fully computed sparse-Gram Frobenius/Markov endpoint for the
non-injective route.  It is still weaker than the standard CountSketch
finite-Loewner subspace embedding and has not yet been composed with the
downstream uniform-row sampling theorem.

The same non-injective theorem now has the finite-Loewner form needed by the
geometric Algorithm 3 statements.  For each realized \(x=(h,\Omega)\), set
\[
  \tau_{\rm CS}^{\rm fp}(x;A,\eta)
  :=
  \eta +
  T_{\rm CSGram}^{\rm fp}
    \bigl(h,\operatorname{rademacherSignVector}(\Omega),A\bigr),
\]
and define the computed two-sided Gram event
\[
  \mathcal E_{\rm CS,L}^{\rm fp}(\eta)
  =
  \left\{x:
  \begin{array}{l}
  \widehat G(x)-A^TA
  \preceq
  \tau_{\rm CS}^{\rm fp}(x;A,\eta)I_n,\\[2pt]
  -(\widehat G(x)-A^TA)
  \preceq
  \tau_{\rm CS}^{\rm fp}(x;A,\eta)I_n
  \end{array}
  \right\}.
\]
Here \(\widehat G(x)\) is the computed sparse CountSketch Gram matrix above,
while \(A^TA\) is the exact analysis reference Gram.  Lean proves the
deterministic inclusion from the Frobenius event by the local bridge
\[
  \|M\|_F\le\tau \quad\Longrightarrow\quad -\tau I\preceq M\preceq \tau I .
\]
Consequently,
\[
  \Pr_{h,\Omega}\{\mathcal E_{\rm CS,L}^{\rm fp}(\eta)\}
  \ge
  1-\frac{2}{r\eta^2}
    \sum_{j,l=1}^{n}\sum_{(a,b)\in\mathcal P_m}(A_{aj}A_{bl})^2
  \ge
  1-\frac{2\|A\|_F^4}{r\eta^2}.
\]
The failure term has the same readable order
\[
  \Theta\!\left(\frac{\|A\|_F^4}{r\eta^2}\right)
\]
for fixed exact \(A\) and \(\eta>0\) as \(r\to\infty\).  If \(U^TU=I_n\), then
\[
  \Pr_{h,\Omega}\{\mathcal E_{\rm CS,L}^{\rm fp}(\eta)\}
  \ge
  1-\frac{2n^2}{r\eta^2},
  \qquad
  \hbox{failure order }
  \Theta\!\left(\frac{n^2}{r\eta^2}\right).
\]
This Loewner statement charges the same computed sparse apply and rounded
Gram-dot arithmetic as the Frobenius statement.  It does not add a
collision-free event or an external certificate assumption; it is still a
Markov/Frobenius-derived Loewner endpoint, not the optimal CountSketch
subspace-embedding concentration theorem.

The exact finite-cover CountSketch Loewner theorem is now also transferred to
the computed sparse Gram object.  Let \(\mathcal A\) be a finite index set and
let \(z_\alpha\in\R^n\), \(\alpha\in\mathcal A\), be an exact unit-ball cover
of radius \(\rho\).  For exact thresholds \(\eta>0\) and \(L>0\), define
\[
  \tau_{\rm cover}(A;\rho,\eta,L)
  :=
  \eta+L(2\rho+\rho^2)
\]
and the exact probability loss
\[
\begin{aligned}
  \mathcal L_{\rm cover}(A,Z;\eta,L,r)
  :=
  &\sum_{\alpha\in\mathcal A}
    \frac{2}{r\eta^2}
    \sum_{(a,b)\in\mathcal P_m}
      \bigl((Az_\alpha)_a(Az_\alpha)_b\bigr)^2       \\
  &+
    \frac{2}{rL^2}
    \sum_{j,l=1}^{n}\sum_{(a,b)\in\mathcal P_m}
      (A_{aj}A_{bl})^2 .
\end{aligned}
\]
For each exact CountSketch realization \(x=(h,\Omega)\), set
\[
  \tau_{\rm cover}^{\rm fp}(x;A,\rho,\eta,L)
  :=
  \eta+L(2\rho+\rho^2)
  +
  T_{\rm CSGram}^{\rm fp}
    \bigl(h,\operatorname{rademacherSignVector}(\Omega),A\bigr).
\]
The theorem
\[
\texttt{countSketchRowGramTwoSidedLoewnerCoverEvent\_subset\_flSparseGramDotRowGramTwoSidedLoewnerEvent}
\]
proves the deterministic transfer
\[
  -\tau_{\rm cover}^{\rm fp}(x)I_n
  \preceq
  \widehat G(x)-A^TA
  \preceq
  \tau_{\rm cover}^{\rm fp}(x)I_n
\]
from the exact finite-cover event.  Here \(\widehat G(x)\) is computed by
\(\operatorname{flCountSketchSparseGramDot}(fp,h,\omega,A)\); the exact
reference \(A^TA\), the finite cover, and the hash/sign probability laws remain
analysis objects.  Sparse signed products, bucket accumulation, and rounded
Gram dot products are charged through the concrete
\(T_{\rm CSGram}^{\rm fp}\) budget, not by an assumed perturbation certificate.
Consequently Lean proves
\[
  \Pr_{h,\Omega}
  \left\{
    -\tau_{\rm cover}^{\rm fp}(x)I_n
    \preceq
    \widehat G(x)-A^TA
    \preceq
    \tau_{\rm cover}^{\rm fp}(x)I_n
  \right\}
  \ge
  1-\mathcal L_{\rm cover}(A,Z;\eta,L,r),
\]
with direct target-budget theorem
\[
  \mathcal L_{\rm cover}(A,Z;\eta,L,r)\le\delta
  \quad\Longrightarrow\quad
  \Pr\{\hbox{the displayed computed-cover event}\}\ge 1-\delta .
\]
The Lean names are
\[
\texttt{countSketchProbability\_eventProb\_flSparseGramDot\_rowGram\_twoSidedLoewner\_cover\_ge\_one\_sub\_sum\_coeff\_add\_frob}
\]
and
\[
\texttt{countSketchProbability\_eventProb\_flSparseGramDot\_rowGram\_twoSidedLoewner\_cover\_ge\_one\_sub\_delta\_of\_sum\_coeff\_add\_frob\_budget}.
\]
Lean also proves the readable sufficient version
\[
  \Pr\{\hbox{the displayed computed-cover event}\}
  \ge
  1-
  \left(
    \sum_{\alpha\in\mathcal A}
      \frac{2\|Az_\alpha\|_2^4}{r\eta^2}
    +
      \frac{2\|A\|_F^4}{rL^2}
  \right),
\]
and its target-budget wrapper
\[
  \sum_{\alpha\in\mathcal A}
      \frac{2\|Az_\alpha\|_2^4}{r\eta^2}
    +
      \frac{2\|A\|_F^4}{rL^2}
  \le\delta
  \quad\Longrightarrow\quad
  \Pr\{\hbox{the displayed computed-cover event}\}\ge1-\delta .
\]
The Lean names are
\[
\texttt{countSketchProbability\_eventProb\_flSparseGramDot\_rowGram\_twoSidedLoewner\_cover\_ge\_one\_sub\_sum\_vecNorm\_add\_frobNorm}
\]
and
\[
\texttt{countSketchProbability\_eventProb\_flSparseGramDot\_rowGram\_twoSidedLoewner\_cover\_ge\_one\_sub\_delta\_of\_sum\_vecNorm\_add\_frobNorm\_budget}.
\]
In the regime of fixed exact \(A\), fixed cover \(Z\), fixed positive
\(\eta,L\), and \(r\to\infty\), the irreducible failure order is
\[
  \Theta\!\left(
    \sum_{\alpha\in\mathcal A}
      \frac{\sum_{(a,b)\in\mathcal P_m}
        ((Az_\alpha)_a(Az_\alpha)_b)^2}{r\eta^2}
    +
      \frac{\sum_{j,l=1}^{n}\sum_{(a,b)\in\mathcal P_m}
        (A_{aj}A_{bl})^2}{rL^2}
  \right)
\]
when the numerator is nonzero.  A readable upper order is
\[
  O\!\left(
    \sum_{\alpha\in\mathcal A}
      \frac{\|Az_\alpha\|_2^4}{r\eta^2}
    +
      \frac{\|A\|_F^4}{rL^2}
  \right).
\]
The probability term is non-vacuous whenever
\(\mathcal L_{\rm cover}(A,Z;\eta,L,r)<1\), and tends to zero as
\(r\to\infty\) for fixed exact inputs and thresholds.  The additional
floating-point term is deterministic radius inflation and vanishes when the
charged sparse-apply and Gram-dot roundoff budgets vanish.

The same product grid also has a computed sparse-Gram endpoint.  Keep the
finite exact grid \(g:\mathcal I\to\R\), product vectors \(z_a\), and
\(\sqrt n\,\delta_{\rm grid}\le\rho\) hypotheses from the preceding
product-grid paragraph, and assume the floating-point model has valid
\(\gamma_m\) and \(\gamma_r\).  Then Lean proves
\[
  \Pr_{h,\Omega}
  \left\{
    -\tau_{\rm pg}^{\rm fp}(x)I_n
    \preceq
    \widehat G(x)-A^TA
    \preceq
    \tau_{\rm pg}^{\rm fp}(x)I_n
  \right\}
  \ge
  1-
  \left(
    \sum_{a:\{1,\ldots,n\}\to\mathcal I}
      \frac{2\|Az_a\|_2^4}{r\eta^2}
    +
      \frac{2\|A\|_F^4}{rL^2}
  \right),
\]
where
\[
  \tau_{\rm pg}^{\rm fp}(x)
  =
  \eta+L(2\rho+\rho^2)
  +
  T_{\rm CSGram}^{\rm fp}
    \bigl(h,\operatorname{rademacherSignVector}(\Omega),A\bigr).
\]
The direct \(1-\delta\) wrapper uses the same displayed product-grid loss.
The Lean theorem names are
\[
\texttt{countSketchProbability\_eventProb\_flSparseGramDot\_rowGram\_twoSidedLoewner\_productGrid\_ge\_one\_sub\_sum\_vecNorm\_add\_frobNorm}
\]
and
\[
\texttt{countSketchProbability\_eventProb\_flSparseGramDot\_rowGram\_twoSidedLoewner\_productGrid\_ge\_one\_sub\_delta\_of\_sum\_vecNorm\_add\_frobNorm\_budget}.
\]
The probability order is the same product-grid
\(\Theta(((\sum_a\|Az_a\|_2^4/\eta^2)+\|A\|_F^4/L^2)/r)\) order above in
the fixed exact-input/grid/threshold regime.  The added
\(T_{\rm CSGram}^{\rm fp}\) term is deterministic radius inflation from the
computed sparse signed products, bucket accumulation, and rounded Gram dot
products; it is not a probability loss and vanishes as those charged
roundoff budgets vanish.

The downstream iid uniform-row sampling foundation is now available without
assuming that the matrix being sampled already has orthonormal columns.  Let
\(M\in\R^{m\times n}\) be an exact matrix, let \(s>0\), and let
\(\sigma=(\sigma_1,\ldots,\sigma_s)\) be exact iid uniform row indices in
\(\{1,\ldots,m\}\).  Define the exact sampled Gram average
\[
  \widehat G_{\rm unif}(M,\sigma)
  :=
  \frac1s\sum_{q=1}^{s}
    m\,M_{\sigma_q,*}^{T}M_{\sigma_q,*}.
\]
For each row \(i\), write
\[
  R_i(M):=\|M_{i,*}\|_2^2
  =
  \sum_{j=1}^{n}M_{ij}^2,
  \qquad
  \|M\|_F^2=\sum_{i=1}^{m}R_i(M).
\]
All objects in this paragraph are exact arithmetic objects; the only random
object is the exact uniform row law.  No floating-point sparse apply,
denominator, row scaling, or Gram-dot routine is being charged here, because
none is computed in this exact foundation theorem.

Lean first proves the one-row mean identity
\[
  \mathbb E_i\!\left[m\,M_{i,*}^{T}M_{i,*}\right]=M^TM,
\]
and the iid sample-average variance identity for each Gram entry.  Summing
entries and bounding the centered second moment by the raw second moment gives
\[
  \mathbb E_\sigma
  \left\|\widehat G_{\rm unif}(M,\sigma)-M^TM\right\|_F^2
  \le
  \frac{m}{s}\sum_{i=1}^{m}R_i(M)^2.
\]
Since
\[
  \sum_{i=1}^{m}R_i(M)^2
  \le
  \left(\sum_{i=1}^{m}R_i(M)\right)^2
  =
  \|M\|_F^4,
\]
the readable bound is
\[
  \mathbb E_\sigma
  \left\|\widehat G_{\rm unif}(M,\sigma)-M^TM\right\|_F^2
  \le
  \frac{m}{s}\|M\|_F^4.
\]
Therefore, for any exact threshold \(\eta>0\),
\[
\begin{aligned}
  \Pr_\sigma
  \left\{
    \left\|\widehat G_{\rm unif}(M,\sigma)-M^TM\right\|_F
    \le \eta
  \right\}
  &\ge
  1-\frac{m\sum_{i=1}^{m}R_i(M)^2}{s\eta^2} \\
  &\ge
  1-\frac{m\|M\|_F^4}{s\eta^2}.
\end{aligned}
\]
The exact row-fourth failure term has order
\[
  \Theta\!\left(
    \frac{m\sum_{i=1}^{m}R_i(M)^2}{s\eta^2}
  \right)
\]
when \(\sum_iR_i(M)^2\ne0\), and the simplified Frobenius upper radius has
order
\[
  \Theta\!\left(\frac{m\|M\|_F^4}{s\eta^2}\right)
\]
when \(\|M\|_F\ne0\), in the regime with exact \(M\), \(m\), and \(\eta>0\)
fixed and \(s\to\infty\).  The bound is non-vacuous when
\(s>m\sum_iR_i(M)^2/\eta^2\), or under the simplified form when
\(s>m\|M\|_F^4/\eta^2\), and the failure term tends to zero as
\(s\to\infty\).

This exact theorem is not, by itself, a final implementation-facing
CountSketch-plus-uniform-row theorem: it is the exact-law foundation that the
computed sampled-Gram bridge below uses.  The Lean theorem surface is:
\leanfit{uniform\_rowOuterGramSample\_mean\_eq\_rowGram}
\leanfit{uniformRowTraceProbability\_expectationReal\_sampleAverage\_sub\_mean\_sq}
\leanfit{uniformRowTraceProbability\_expectationReal\_uniformRowSampleGram\_entry\_error\_sq}
\leanfit{uniformRowTraceProbability\_expectationReal\_uniformRowSampleGram\_frob\_error\_sq\_le}
\leanfit{uniformRowTraceProbability\_expectationReal\_uniformRowSampleGram\_frob\_error\_sq\_le\_frobNorm}
\leanfit{uniformRowTraceProbability\_eventProb\_uniformRowSampleGram\_frob\_error\_le\_ge\_one\_sub}
\leanfit{uniformRowTraceProbability\_eventProb\_uniformRowSampleGram\_frob\_error\_le\_ge\_one\_sub\_frobNorm}

The arbitrary-matrix uniform-row sampled Gram now also has a concrete
floating-point bridge.  Keep the same exact matrix \(M\), exact iid uniform row
samples \(\sigma\), sample count \(s>0\), and exact reference Gram \(M^TM\).
The probability law for \(\sigma\) is exact by project convention.  The
non-probability operations charged in this theorem are: row scaling by the
uniform sampling denominator, optional computation/storage of that denominator,
and the length-\(s\) dot products used to form the sampled Gram.

Let
\[
  d:=\sqrt{s/m},\qquad
  B_{qj}:=\frac{M_{\sigma_q j}}{d},
  \qquad
  C_B:=\bigl(\sum_{q=1}^{s}|B_{qj}|\,|B_{qk}|\bigr)_{j,k}.
\]
Here \(B\) is the exact sampled-and-rescaled sketch used only as the
reference object in the analysis.  If the implementation divides each sampled
entry by the exact denominator \(d\) in floating point and then computes every
Gram entry by a rounded length-\(s\) dot product, Lean proves
\[
\begin{aligned}
  \Pr_\sigma\Bigl\{
    \bigl\|\widehat G_{\rm unif}^{\rm fp}(M,\sigma)-M^TM\bigr\|_F
    \le \eta + T_{\rm unif}^{\rm fp}(M,\sigma)
  \Bigr\}
  &\ge
  1-\frac{m\sum_{i=1}^{m}R_i(M)^2}{s\eta^2} \\
  &\ge
  1-\frac{m\|M\|_F^4}{s\eta^2},
\end{aligned}
\]
where the exact expanded floating-point radius is
\[
\begin{aligned}
  T_{\rm unif}^{\rm fp}(M,\sigma)
  &=
  \left\|
    \bigl((2u+u^2)\sum_{q=1}^{s}|B_{qj}|\,|B_{qk}|\bigr)_{j,k}
  \right\|_F \\
  &\quad+
  \left\|
    \bigl(\gamma_s(1+u)^2
      \sum_{q=1}^{s}|B_{qj}|\,|B_{qk}|\bigr)_{j,k}
  \right\|_F \\
  &=
  \bigl(2u+u^2+\gamma_s(1+u)^2\bigr)\,\|C_B\|_F.
\end{aligned}
\]
The first term is the rounded row-division contribution; the second is the
rounded dot-product contribution.  In the usual small-roundoff regime with
\(su\) bounded away from \(1\), \(\gamma_s=\Theta(su)\), and therefore
\[
  T_{\rm unif}^{\rm fp}(M,\sigma)
  =
  \Theta\!\left((u+\gamma_s)\|C_B\|_F\right)
  =
  \Theta\!\left(su\,\|C_B\|_F\right)
\]
when \(C_B\ne0\).  For fixed \(s\), exact \(M\), and fixed realized samples,
the same expression is \(\Theta(u\|C_B\|_F)\).  It tends to zero as the unit
roundoff tends to zero.

Lean also proves the implementation-facing variant where the denominator is
itself computed.  Let \(\widehat d\) be the computed nonzero denominator, let
\(\Delta_d:=|\widehat d-d|\), define
\[
  \widehat B_{qj}:=\operatorname{fl}\!\left(
    \frac{M_{\sigma_q j}}{\widehat d}
  \right),
  \qquad
  E_{qj}:=
    |\widehat B_{qj}|\,u+
    |M_{\sigma_qj}|\frac{\Delta_d}{|\widehat d|\,|d|}.
\]
The quantity \(E_{qj}\) is fully computable from the modeled rounded sketch,
the exact sampled input entry used in the analysis, and the proved denominator
absolute-error certificate.  With this computed denominator, Lean proves
\[
\begin{aligned}
  \Pr_\sigma\Bigl\{
    \bigl\|\widehat G_{\rm unif,den}^{\rm fp}(M,\widehat d,\sigma)
      -M^TM\bigr\|_F
    \le
      \eta + T_{\rm unif,den}^{\rm fp}(M,\widehat d,\sigma)
  \Bigr\}
  &\ge
  1-\frac{m\sum_{i=1}^{m}R_i(M)^2}{s\eta^2} \\
  &\ge
  1-\frac{m\|M\|_F^4}{s\eta^2},
\end{aligned}
\]
with exact expanded radius
\[
\begin{aligned}
  T_{\rm unif,den}^{\rm fp}
  &=
  \left\|
    \bigl(\sum_{q=1}^{s}
      (E_{qj}|\widehat B_{qk}|+|B_{qj}|E_{qk})\bigr)_{j,k}
  \right\|_F \\
  &\quad+
  \left\|
    \bigl(\gamma_s
      \sum_{q=1}^{s}|\widehat B_{qj}|\,|\widehat B_{qk}|\bigr)_{j,k}
  \right\|_F.
\end{aligned}
\]
Equivalently, using matrix notation with entrywise absolute values,
\[
  T_{\rm unif,den}^{\rm fp}
  =
  \|E^T|\widehat B|+|B|^TE\|_F
  +
  \gamma_s\||\widehat B|^T|\widehat B|\|_F,
\]
and its interpretable order is
\[
  \Theta\!\left(
    \|E^T|\widehat B|+|B|^TE\|_F
    +
    \gamma_s\||\widehat B|^T|\widehat B|\|_F
  \right).
\]
Substituting \(E_{qj}\) gives the irreducible order form
\[
  \Theta\!\left(
    \left\|
      (u|\widehat B|+\delta_d|M_\sigma|)^T|\widehat B|
      +
      |B|^T(u|\widehat B|+\delta_d|M_\sigma|)
    \right\|_F
    +
    \gamma_s\||\widehat B|^T|\widehat B|\|_F
  \right),
\]
where \(\delta_d:=\Delta_d/(|\widehat d|\,|d|)\) and
\((M_\sigma)_{qj}:=M_{\sigma_qj}\).  This radius is non-vacuous in the
floating-point sense: for fixed exact data and samples, it tends to zero as
\(u\to0\), \(\gamma_s\to0\), and the denominator error
\(\Delta_d\to0\).

The Lean theorem surface for this floating-point bridge is:
\leanfit{uniformRowSampleGramRowGramFrobErrorEvent}
\leanfit{uniformRowFlSampleGramDotRowGramFrobErrorEvent}
\leanfit{uniformRowFlSampleGramDotWithComputedDenRowGramFrobErrorEvent}
\leanfit{uniformRowSampleGramRowGramFrobErrorEvent\_subset\_flSampleGramDot}
\leanfit{uniformRowSampleGramRowGramFrobErrorEvent\_subset\_flSampleGramDotWithComputedDen}
\leanfit{uniformRowTraceProbability\_eventProb\_fl\_uniformRowSampleGramDot\_rowGram\_frob\_error\_le\_ge\_one\_sub}
\leanfit{uniformRowTraceProbability\_eventProb\_fl\_uniformRowSampleGramDot\_rowGram\_frob\_error\_le\_ge\_one\_sub\_frobNorm}
\leanfit{uniformRowTraceProbability\_eventProb\_fl\_uniformRowSampleGramDotWithComputedDen\_rowGram\_frob\_error\_le\_ge\_one\_sub}
\leanfit{uniformRowTraceProbability\_eventProb\_fl\_uniformRowSampleGramDotWithComputedDen\_rowGram\_frob\_error\_le\_ge\_one\_sub\_frobNorm}

The non-injective CountSketch branch is now composed with this downstream
uniform-row machinery.  Draw \(x=(h,\Omega)\) from the exact CountSketch
hash-sign law, set
\(\omega=\operatorname{rademacherSignVector}(\Omega)\), and let
\(\sigma\) be exact iid uniform samples from the \(r\) CountSketch output rows.
Define
\[
  V(h,\Omega):=S_{h,\omega}A,
  \qquad
  \widehat V(h,\Omega):=\operatorname{flCountSketchSparseApply}(h,\omega,A).
\]
The new deterministic CountSketch bucket identity proves
\[
  \sum_{i=1}^{r}\sum_{k:h(k)=i} f(k)=\sum_{k=1}^{m}f(k),
\]
and, by Cauchy--Schwarz, the exact CountSketch output satisfies
\[
  \|V(h,\Omega)\|_F^2\le m\|A\|_F^2
\]
for every realized hash and every exact sign vector with \(|\omega_k|\le1\).
This supplies the uniform, nonconditional row-sampling loss needed after
non-injective CountSketch.

Let \(\eta_{\rm CS}>0\) be the exact CountSketch Gram threshold and
\(\eta_{\rm row}>0\) the exact downstream row-sampling threshold.  The exact
product event is the intersection
\[
\begin{aligned}
  \|V(h,\Omega)^TV(h,\Omega)-A^TA\|_F &\le \eta_{\rm CS},\\
  \|\widehat G_{\rm unif}(V(h,\Omega),\sigma)
    -V(h,\Omega)^TV(h,\Omega)\|_F &\le \eta_{\rm row}.
\end{aligned}
\]
Lean proves
\[
\begin{aligned}
  \Pr\{\mathcal E_{\rm exact}^{\rm CS+row}\}
  \ge
  1
  &-
  \frac{2}{r\eta_{\rm CS}^2}
    \sum_{j,k=1}^{n}\sum_{(a,b)\in\mathcal P_m}
      (A_{aj}A_{bk})^2  \\
  &-
  \frac{r\,m^2\|A\|_F^4}{s\eta_{\rm row}^2},
  \qquad
  \mathcal P_m=\{(a,b):a\ne b\}.
\end{aligned}
\]
Using the already-proved coefficient simplification, the readable form is
\[
  \Pr\{\mathcal E_{\rm exact}^{\rm CS+row}\}
  \ge
  1
  -
  \frac{2\|A\|_F^4}{r\eta_{\rm CS}^2}
  -
  \frac{r\,m^2\|A\|_F^4}{s\eta_{\rm row}^2}.
\]
The corresponding computed event forms \(\widehat V\), computes a denominator
\(\widehat d\) for \(\sqrt{s/r}\), rounds sampled-row divisions, and rounds
the length-\(s\) sampled-Gram dot products:
\[
  \left\|
    \widehat G_{\rm CS+row}^{\rm fp}(h,\Omega,\sigma)
    -
    A^TA
  \right\|_F
  \le
  \eta_{\rm CS}+\eta_{\rm row}
  +
  T_{\rm CS,row}^{\rm fp}(h,\Omega,\sigma).
\]
The same probability lower bound holds for this computed event.  The realized
floating-point radius is not an opaque certificate:
\[
  T_{\rm CS,row}^{\rm fp}
  =
  T_{\rm unif,den}^{\rm fp}(\widehat V,\widehat d,\sigma)
  +
  T_{\rm basis}^{\rm fp}(V,\widehat V,E^{\rm CS},\sigma),
\]
where the first term is the expanded computed-denominator sampled-Gram radius
displayed above with \(M=\widehat V\), and
\[
  T_{\rm basis}^{\rm fp}
  =
  \left\|
    \left(
      \sum_{q=1}^{s}
      \left(
        \frac{E^{\rm CS}_{\sigma_q j}}{d}
        |\widehat B_{qk}|
        +
        |B_{qj}|
        \frac{E^{\rm CS}_{\sigma_q k}}{d}
      \right)
    \right)_{j,k}
  \right\|_F,
  \qquad d=\sqrt{s/r}.
\]
Here \(B\) is the exact sampled sketch from \(V\), \(\widehat B\) is the
sampled sketch from \(\widehat V\), and
\[
  E^{\rm CS}_{ij}
  =
  (u+\gamma_{b_i}+u\gamma_{b_i})
  \sum_{k:h(k)=i}|\omega_k|\,|A_{kj}|,
  \qquad b_i=|\{k:h(k)=i\}|.
\]
For fixed exact \(A,m,\eta_{\rm CS},\eta_{\rm row}\), the simplified failure
loss has order
\[
  \Theta\!\left(
    \frac{\|A\|_F^4}{r\eta_{\rm CS}^2}
    +
    \frac{r\,m^2\|A\|_F^4}{s\eta_{\rm row}^2}
  \right).
\]
It is non-vacuous when the displayed loss is below one, and it tends to zero
in the regime \(r\to\infty\), \(s/r\to\infty\).  For fixed realized
hashes, signs, and row samples, the FP radius tends to zero as \(u\to0\),
\(\gamma_s\to0\), every \(\gamma_{b_i}\to0\), and
the denominator error \(|\widehat d-d|\to0\).  This closes a genuine
non-injective CountSketch plus downstream row-sampling implementation endpoint;
it is still Markov/Frobenius-derived and weaker than the optimal CountSketch
subspace-embedding theorem.

The Lean theorem surface is:
\leanfit{countSketchBucket\_sum\_sum\_eq}
\leanfit{frobNormSqRect\_preconditionRows\_countSketchRows\_le}
\leanfit{countSketchUniformRowSampleGramRowGramFrobEvent}
\leanfit{countSketchSparseComputedPreconditionedFlUniformRowSampleGramWithComputedDenRowGramFrobEvent}
\leanfit{countSketchUniformRowTraceProbability\_eventProb\_uniformRowSampleGram\_rowGram\_frob\_error\_le\_ge\_one\_sub}
\leanfit{countSketchUniformRowTraceProbability\_eventProb\_sparseComputedPreconditioned\_fl\_uniformRowSampleGramDotWithComputedDen\_rowGram\_frob\_error\_le\_ge\_one\_sub}
\leanfit{countSketchUniformRowTraceProbability\_eventProb\_sparseComputedPreconditioned\_fl\_uniformRowSampleGramDotWithComputedDen\_rowGram\_frob\_error\_le\_ge\_one\_sub\_frobNorm}

The same fully computed non-injective CountSketch-plus-uniform-row endpoint now
has a two-sided finite-Loewner form.  For each realization, set
\[
  \tau_{\rm CS,row}^{\rm fp}
  :=
  \eta_{\rm CS}+\eta_{\rm row}
  +
  T_{\rm CS,row}^{\rm fp}(h,\Omega,\sigma).
\]
The computed two-sided event is
\[
  -\tau_{\rm CS,row}^{\rm fp}I_n
  \preceq
  \widehat G_{\rm CS+row}^{\rm fp}(h,\Omega,\sigma)-A^TA
  \preceq
  \tau_{\rm CS,row}^{\rm fp}I_n .
\]
It is obtained deterministically from the Frobenius event by the locally proved
bridge
\[
  \|M\|_F\le\tau \quad\Longrightarrow\quad -\tau I\preceq M\preceq\tau I.
\]
Therefore the exact coefficient and readable probability losses are unchanged:
\[
\begin{aligned}
  \Pr\{\mathcal E_{\rm CS+row,L}^{\rm fp}\}
  \ge
  1
  &-
  \frac{2}{r\eta_{\rm CS}^2}
    \sum_{j,k=1}^{n}\sum_{(a,b)\in\mathcal P_m}
      (A_{aj}A_{bk})^2  \\
  &-
  \frac{r\,m^2\|A\|_F^4}{s\eta_{\rm row}^2} \\
  \ge
  1
  &-
  \frac{2\|A\|_F^4}{r\eta_{\rm CS}^2}
  -
  \frac{r\,m^2\|A\|_F^4}{s\eta_{\rm row}^2}.
\end{aligned}
\]
The simplified failure loss again has order
\[
  \Theta\!\left(
    \frac{\|A\|_F^4}{r\eta_{\rm CS}^2}
    +
    \frac{r\,m^2\|A\|_F^4}{s\eta_{\rm row}^2}
  \right).
\]
The finite-Loewner theorem is non-vacuous in the same regime: the displayed
loss is below one, \(r\to\infty\), \(s/r\to\infty\), and the charged
roundoff/denominator errors vanish.  This is a geometry-shaped theorem for
the fully computed downstream row-sampling path, but its probability rate is
still inherited from Markov/Frobenius estimates rather than from the optimal
CountSketch subspace-embedding concentration theorem.

The Lean theorem surface is:\par
\noindent{\scriptsize\ttfamily countSketchSparseComputedPreconditionedFlUniformRowSampleGramWithComputedDen}\par
\noindent{\scriptsize\ttfamily RowGramTwoSidedLoewnerEvent}\par
\leanfit{countSketchSparseComputedPreconditionedFlUniformRowSampleGramWithComputedDenRowGramFrobEvent\_subset\_twoSidedLoewnerEvent}
\leanfit{countSketchUniformRowTraceProbability\_eventProb\_sparseComputedPreconditioned\_fl\_uniformRowSampleGramDotWithComputedDen\_rowGram\_twoSidedLoewner\_ge\_one\_sub}
\leanfit{countSketchUniformRowTraceProbability\_eventProb\_sparseComputedPreconditioned\_fl\_uniformRowSampleGramDotWithComputedDen\_rowGram\_twoSidedLoewner\_ge\_one\_sub\_frobNorm}

The corresponding sample-budget corollaries are also formalized.  The sharper
exact-coefficient wrapper states that if
\[
  \frac{2}{r\eta_{\rm CS}^2}
    \sum_{j,k=1}^{n}\sum_{(a,b)\in\mathcal P_m}
      (A_{aj}A_{bk})^2
  +
  \frac{r\,m^2\|A\|_F^4}{s\eta_{\rm row}^2}
  \le \delta,
\]
then the same fully computed two-sided event has probability at least
\(1-\delta\).  The first term is the exact CountSketch collision/sign
coefficient loss; the second term is the proved downstream uniform-row
Frobenius-growth loss for \(S_{h,\Omega}A\).

The readable Frobenius-norm wrapper follows by bounding the coefficient sum by
\(\|A\|_F^4\).  If
\[
  \frac{2\|A\|_F^4}{r\eta_{\rm CS}^2}
  +
  \frac{r\,m^2\|A\|_F^4}{s\eta_{\rm row}^2}
  \le \delta,
\]
then the same fully computed two-sided event has probability at least
\(1-\delta\).  With the equal split
\(\eta_{\rm CS}=\eta_{\rm row}=\varepsilon/2\), the exact
non-floating-point part of the Loewner radius is \(\varepsilon\), so the
realized radius is
\[
  \varepsilon + T_{\rm CS,row}^{\rm fp}(h,\Omega,\sigma).
\]
The readable sufficient loss condition expands irreducibly to
\[
  \frac{8\|A\|_F^4}{r\varepsilon^2}
  +
  \frac{4r\,m^2\|A\|_F^4}{s\varepsilon^2}
  \le \delta.
\]
The equal-radius exact-coefficient sufficient condition is
\[
  \frac{8}{r\varepsilon^2}
    \sum_{j,k=1}^{n}\sum_{(a,b)\in\mathcal P_m}
      (A_{aj}A_{bk})^2
  +
  \frac{4r\,m^2\|A\|_F^4}{s\varepsilon^2}
  \le \delta.
\]
Lean now exposes these two expanded equal-radius sufficient conditions as
checked theorem surfaces, not only as prose simplifications of the
\(\eta=\varepsilon/2\) wrappers.
Thus, for fixed exact \(A,m,\varepsilon>0\) in the small-roundoff regime and
varying \(r,s\), the equal-radius failure loss has order
\[
  \Theta\!\left(
    \frac{\|A\|_F^4}{r\varepsilon^2}
    +
    \frac{r\,m^2\|A\|_F^4}{s\varepsilon^2}
  \right).
\]
Keeping the exact coefficient sum gives the irreducible order
\[
  \Theta\!\left(
    \frac{1}{r\varepsilon^2}
      \sum_{j,k=1}^{n}\sum_{(a,b)\in\mathcal P_m}
        (A_{aj}A_{bk})^2
    +
    \frac{r\,m^2\|A\|_F^4}{s\varepsilon^2}
  \right),
\]
with the displayed coefficient sum treated as fixed exact input data.
The statement is non-vacuous when the displayed loss is below one; it tends to
zero when \(r\to\infty\), \(s/r\to\infty\), and the charged roundoff and
denominator errors vanish.

The finite-cover CountSketch downstream wrappers replace the Frobenius
preprocessing radius \(\eta_{\rm CS}\) by the exact finite-cover Loewner radius
\[
  \tau_{\rm CS}=\eta+L(2\rho+\rho^2).
\]
Let \(z_\alpha\) be the exact finite unit-ball cover vectors.  If
\[
  \sum_\alpha
    \frac{2r^{-1}}{\eta^2}
      \sum_{(a,b)\in\mathcal P_m}
        \big((Az_\alpha)_a(Az_\alpha)_b\big)^2
  +
  \frac{2r^{-1}}{L^2}
    \sum_{j,k=1}^{n}\sum_{(a,b)\in\mathcal P_m}
      (A_{aj}A_{bk})^2
  +
  \frac{r\,m^2\|A\|_F^4}{s\eta_{\rm row}^2}
  \le \delta,
\]
then the fully computed downstream event has probability at least
\(1-\delta\) and radius
\[
  \tau_{\rm CS}+\eta_{\rm row}
  +T_{\rm CS,row}^{\rm fp}(h,\Omega,\sigma).
\]
Thus, in the fixed-cover, fixed-\(A\), small-roundoff regime, the irreducible
loss order is
\[
  \Theta\!\left(
    \sum_\alpha
      \frac{r^{-1}}{\eta^2}
        \sum_{(a,b)\in\mathcal P_m}
          \big((Az_\alpha)_a(Az_\alpha)_b\big)^2
    +
    \frac{r^{-1}}{L^2}
      \sum_{j,k=1}^{n}\sum_{(a,b)\in\mathcal P_m}
        (A_{aj}A_{bk})^2
    +
    \frac{r\,m^2\|A\|_F^4}{s\eta_{\rm row}^2}
  \right).
\]
This is still finite-cover Markov/Frobenius-derived, not the optimal
CountSketch subspace-embedding concentration theorem, but it is a checked
end-to-end implementation-facing theorem for the modeled CountSketch plus
uniform-row path.

The same downstream finite-cover theorem is now instantiated by the exact
product grid.  Let \(\mathcal I\) be a finite exact index set, let
\(g:\mathcal I\to\R\) cover \([-1,1]\) with radius
\(\delta_{\rm grid}\ge0\), and assume
\[
  \sqrt n\,\delta_{\rm grid}\le\rho .
\]
For \(a:\{1,\ldots,n\}\to\mathcal I\), set
\[
  z_a(j)=g(a(j)).
\]
The product-grid downstream theorem has the same computed event radius
\[
  \eta+L(2\rho+\rho^2)+\eta_{\rm row}
  +T_{\rm CS,row}^{\rm fp}(h,\Omega,\sigma),
\]
but its exact target loss is displayed without a supplied cover:
\[
\begin{aligned}
  &
  \sum_{a:\{1,\ldots,n\}\to\mathcal I}
    \frac{2r^{-1}}{\eta^2}
      \sum_{(p,q)\in\mathcal P_m}
        \big((Az_a)_p(Az_a)_q\big)^2
  \\
  &\quad+
  \frac{2r^{-1}}{L^2}
    \sum_{j,k=1}^{n}\sum_{(p,q)\in\mathcal P_m}
      (A_{pj}A_{qk})^2
  +
  \frac{r\,m^2\|A\|_F^4}{s\eta_{\rm row}^2}.
\end{aligned}
\]
If this loss is at most \(\delta\), then the fully computed downstream event
has probability at least \(1-\delta\).  In the fixed exact-input/grid/threshold
regime, the irreducible order is
\[
  \Theta\!\left(
    \sum_{a:\{1,\ldots,n\}\to\mathcal I}
      \frac{r^{-1}}{\eta^2}
        \sum_{(p,q)\in\mathcal P_m}
          \big((Az_a)_p(Az_a)_q\big)^2
    +
    \frac{r^{-1}}{L^2}
      \sum_{j,k=1}^{n}\sum_{(p,q)\in\mathcal P_m}
        (A_{pj}A_{qk})^2
    +
    \frac{r\,m^2\|A\|_F^4}{s\eta_{\rm row}^2}
  \right).
\]
The concrete product-grid denominator theorem selects the same
\[
  \operatorname{fl\_mul}\!\left(
    \operatorname{fl\_sqrt}(s),
    \operatorname{fl\_div}\!\left(1,\operatorname{fl\_sqrt}(r)\right)
  \right)
\]
routine for \(d=\sqrt{s/r}\).  The grid and product index set are exact
analysis objects; sparse apply, denominator formation, row divisions, and
sampled-Gram dot products are the computed non-probability quantities charged
by \(T_{\rm CS,row}^{\rm fp}\).

The corresponding Lean theorem surfaces are:\par
\noindent{\scriptsize\ttfamily\detokenize{countSketchUniformRowTraceProbability_eventProb_sparseComputedPreconditioned_fl}}\par
\noindent{\scriptsize\ttfamily\detokenize{_uniformRowSampleGramDotWithComputedDen_rowGram_twoSidedLoewner_cover_ge_one_sub}}\par
\noindent{\scriptsize\ttfamily\detokenize{_sum_coeff_add_frob_add_row}}\par
\noindent{\scriptsize\ttfamily\detokenize{countSketchUniformRowTraceProbability_eventProb_sparseComputedPreconditioned_fl}}\par
\noindent{\scriptsize\ttfamily\detokenize{_uniformRowSampleGramDotWithComputedDen_rowGram_twoSidedLoewner_cover_ge_one_sub}}\par
\noindent{\scriptsize\ttfamily\detokenize{_delta_of_sum_coeff_add_frob_add_row_budget}}\par
\noindent{\scriptsize\ttfamily\detokenize{countSketchUniformRowTraceProbability_eventProb_sparseComputedPreconditioned_fl}}\par
\noindent{\scriptsize\ttfamily\detokenize{_uniformRowSampleGramDotWithFlSqrtMulInvSqrtDen_rowGram_twoSidedLoewner_cover}}\par
\noindent{\scriptsize\ttfamily\detokenize{_ge_one_sub_delta_of_sum_coeff_add_frob_add_row_budget}}\par
\noindent{\scriptsize\ttfamily\detokenize{countSketchUniformRowTraceProbability_eventProb_sparseComputedPreconditioned_fl}}\par
\noindent{\scriptsize\ttfamily\detokenize{_uniformRowSampleGramDotWithComputedDen_rowGram_twoSidedLoewner_productGrid}}\par
\noindent{\scriptsize\ttfamily\detokenize{_ge_one_sub_sum_coeff_add_frob_add_row}}\par
\noindent{\scriptsize\ttfamily\detokenize{countSketchUniformRowTraceProbability_eventProb_sparseComputedPreconditioned_fl}}\par
\noindent{\scriptsize\ttfamily\detokenize{_uniformRowSampleGramDotWithComputedDen_rowGram_twoSidedLoewner_productGrid}}\par
\noindent{\scriptsize\ttfamily\detokenize{_ge_one_sub_delta_of_sum_coeff_add_frob_add_row_budget}}\par
\noindent{\scriptsize\ttfamily\detokenize{countSketchUniformRowTraceProbability_eventProb_sparseComputedPreconditioned_fl}}\par
\noindent{\scriptsize\ttfamily\detokenize{_uniformRowSampleGramDotWithFlSqrtMulInvSqrtDen_rowGram_twoSidedLoewner_productGrid}}\par
\noindent{\scriptsize\ttfamily\detokenize{_ge_one_sub_delta_of_sum_coeff_add_frob_add_row_budget}}\par

The sample-budget Lean theorem surface is:\par
\noindent{\scriptsize\ttfamily\detokenize{countSketchUniformRowTraceProbability_eventProb_sparseComputedPreconditioned_fl}}\par
\noindent{\scriptsize\ttfamily\detokenize{_uniformRowSampleGramDotWithComputedDen_rowGram_twoSidedLoewner_ge_one_sub}}\par
\noindent{\scriptsize\ttfamily\detokenize{_delta_of_coeff_budget}}\par
\noindent{\scriptsize\ttfamily\detokenize{countSketchUniformRowTraceProbability_eventProb_sparseComputedPreconditioned_fl}}\par
\noindent{\scriptsize\ttfamily\detokenize{_uniformRowSampleGramDotWithComputedDen_rowGram_twoSidedLoewner_ge_one_sub}}\par
\noindent{\scriptsize\ttfamily\detokenize{_delta_of_frobNorm_budget}}\par
\noindent{\scriptsize\ttfamily\detokenize{countSketchUniformRowTraceProbability_eventProb_sparseComputedPreconditioned_fl}}\par
\noindent{\scriptsize\ttfamily\detokenize{_uniformRowSampleGramDotWithComputedDen_rowGram_twoSidedLoewner_ge_one_sub}}\par
\noindent{\scriptsize\ttfamily\detokenize{_delta_of_equal_radius_coeff_budget}}\par
\noindent{\scriptsize\ttfamily\detokenize{countSketchUniformRowTraceProbability_eventProb_sparseComputedPreconditioned_fl}}\par
\noindent{\scriptsize\ttfamily\detokenize{_uniformRowSampleGramDotWithComputedDen_rowGram_twoSidedLoewner_ge_one_sub}}\par
\noindent{\scriptsize\ttfamily\detokenize{_delta_of_equal_radius_frobNorm_budget}}\par
\noindent{\scriptsize\ttfamily\detokenize{countSketchUniformRowTraceProbability_eventProb_sparseComputedPreconditioned_fl}}\par
\noindent{\scriptsize\ttfamily\detokenize{_uniformRowSampleGramDotWithComputedDen_rowGram_twoSidedLoewner_ge_one_sub}}\par
\noindent{\scriptsize\ttfamily\detokenize{_delta_of_equal_radius_coeff_budget_expanded}}\par
\noindent{\scriptsize\ttfamily\detokenize{countSketchUniformRowTraceProbability_eventProb_sparseComputedPreconditioned_fl}}\par
\noindent{\scriptsize\ttfamily\detokenize{_uniformRowSampleGramDotWithComputedDen_rowGram_twoSidedLoewner_ge_one_sub}}\par
\noindent{\scriptsize\ttfamily\detokenize{_delta_of_equal_radius_frobNorm_budget_expanded}}\par

The final concrete-denominator wrappers for this non-injective finite-Loewner
path instantiate
\[
  \widehat d
  =
  \operatorname{fl\_mul}\!\left(
    \operatorname{fl\_sqrt}(s),
    \operatorname{fl\_div}\!\left(1,\operatorname{fl\_sqrt}(r)\right)
  \right),
\]
so the exact denominator is \(d=\sqrt{s/r}\) and the charged denominator
radius is
\[
  \sqrt{s/r}\,
  \frac{4u+3u^2+u^3}{1-u}.
\]
The generic \texttt{WithComputedDen} statements remain reusable
infrastructure for other denominator routines; the concrete theorem surface is:
\par
\noindent{\scriptsize\ttfamily\detokenize{countSketchUniformRowTraceProbability_eventProb_sparseComputedPreconditioned_fl}}\par
\noindent{\scriptsize\ttfamily\detokenize{_uniformRowSampleGramDotWithFlSqrtMulInvSqrtDen_rowGram_twoSidedLoewner_ge_one_sub}}\par
\noindent{\scriptsize\ttfamily\detokenize{_delta_of_coeff_budget}}\par
\noindent{\scriptsize\ttfamily\detokenize{countSketchUniformRowTraceProbability_eventProb_sparseComputedPreconditioned_fl}}\par
\noindent{\scriptsize\ttfamily\detokenize{_uniformRowSampleGramDotWithFlSqrtMulInvSqrtDen_rowGram_twoSidedLoewner_ge_one_sub}}\par
\noindent{\scriptsize\ttfamily\detokenize{_delta_of_equal_radius_coeff_budget}}\par
\noindent{\scriptsize\ttfamily\detokenize{countSketchUniformRowTraceProbability_eventProb_sparseComputedPreconditioned_fl}}\par
\noindent{\scriptsize\ttfamily\detokenize{_uniformRowSampleGramDotWithFlSqrtMulInvSqrtDen_rowGram_twoSidedLoewner_ge_one_sub}}\par
\noindent{\scriptsize\ttfamily\detokenize{_delta_of_frobNorm_budget}}\par
\noindent{\scriptsize\ttfamily\detokenize{countSketchUniformRowTraceProbability_eventProb_sparseComputedPreconditioned_fl}}\par
\noindent{\scriptsize\ttfamily\detokenize{_uniformRowSampleGramDotWithFlSqrtMulInvSqrtDen_rowGram_twoSidedLoewner_ge_one_sub}}\par
\noindent{\scriptsize\ttfamily\detokenize{_delta_of_equal_radius_frobNorm_budget}}\par
\noindent{\scriptsize\ttfamily\detokenize{countSketchUniformRowTraceProbability_eventProb_sparseComputedPreconditioned_fl}}\par
\noindent{\scriptsize\ttfamily\detokenize{_uniformRowSampleGramDotWithFlSqrtMulInvSqrtDen_rowGram_twoSidedLoewner_ge_one_sub}}\par
\noindent{\scriptsize\ttfamily\detokenize{_delta_of_equal_radius_coeff_budget_expanded}}\par
\noindent{\scriptsize\ttfamily\detokenize{countSketchUniformRowTraceProbability_eventProb_sparseComputedPreconditioned_fl}}\par
\noindent{\scriptsize\ttfamily\detokenize{_uniformRowSampleGramDotWithFlSqrtMulInvSqrtDen_rowGram_twoSidedLoewner_ge_one_sub}}\par
\noindent{\scriptsize\ttfamily\detokenize{_delta_of_equal_radius_frobNorm_budget_expanded}}\par

The collision-free exact hash route now also has a computed sparse-Gram
probability endpoint.  Assume the exact signs satisfy \(\omega_k^2=1\) for all
input rows and assume \(\gamma_m\) and \(\gamma_r\) are valid.  The lemma
\(b_i\le m\) lets the single assumption \(\gamma_m\) cover every bucket
coefficient \(\gamma_{b_i}\).  Define the exact hash event
\[
  \mathcal E_{\rm fp}
  =
  \left\{
    h:
    \|\widehat G(h)-A^TA\|_F
    \le T_{\rm CSGram}^{\rm fp}(h)
  \right\},
\]
where \(\widehat G(h)\) is computed by the sparse apply and rounded Gram-dot
routine above, while \(A^TA\) and \(T_{\rm CSGram}^{\rm fp}(h)\) are exact
mathematical quantities in the theorem statement.  If \(h\) is injective, then
the exact CountSketch table has orthonormal columns and therefore
\[
  (S_{h,\omega}A)^T(S_{h,\omega}A)=A^TA.
\]
Combining this exact identity with the sparse-Gram FP bound gives
\[
  \{h:h\hbox{ is injective}\}\subseteq \mathcal E_{\rm fp}.
\]
The exact hash-law union bound then yields
\[
  \Pr_h(\mathcal E_{\rm fp})
  \ge
  \Pr_h(h\hbox{ is injective})
  \ge
  1-\frac{m^2}{r}.
\]
This probability statement has no floating-point probability computation in
it: the probability is the exact hash law, while the event describes the
deterministic rounded sparse apply and rounded Gram dot products.  It is a
collision-free sufficient route and is non-vacuous when \(r>m^2\).  It is
deliberately weaker than the usual CountSketch concentration theorem.

The same collision-free route is now closed for the full exact CountSketch
hash-sign product law.  Draw \(x=(h,\Omega)\) from
\(\texttt{countSketchProbability}\), and let
\(\omega=\operatorname{rademacherSignVector}(\Omega)\).  Define
\[
  \mathcal E_{\rm fp}^{\rm prod}
  =
  \left\{
    (h,\Omega):
    \|\widehat G(h,\Omega)-A^TA\|_F
    \le
    T_{\rm CSGram}^{\rm fp}
      \bigl(h,\operatorname{rademacherSignVector}(\Omega)\bigr)
  \right\},
\]
where \(\widehat G(h,\Omega)\) is computed by the same sparse apply and rounded
Gram-dot routine using the realized exact Rademacher signs.  The Rademacher
sign theorem gives \(\omega_k^2=1\) for every realized sign vector, so
\[
  \{(h,\Omega):h\hbox{ is injective}\}
  \subseteq
  \mathcal E_{\rm fp}^{\rm prod}.
\]
The product-law marginal identity for first-coordinate events and the exact
hash union bound give the fully realized CountSketch statement
\[
  \Pr_{h,\Omega}(\mathcal E_{\rm fp}^{\rm prod})
  \ge
  \Pr_h(h\hbox{ is injective})
  \ge
  1-\frac{m^2}{r}.
\]
Again, probabilities and sign/hash draws are exact mathematical laws.  The
only floating-point terms are the deterministic sparse signed products, bucket
accumulations, and Gram dot products inside \(\widehat G(h,\Omega)\).  The
bound is non-vacuous on the same collision-free sufficient regime \(r>m^2\)
and remains weaker than the usual CountSketch concentration/subspace-embedding
theorem.

The collision-free CountSketch route now also composes with iid exact uniform
row sampling and the computed uniform-row sample-Gram routine.  Let
\(U\in\R^{m\times n}\) be an exact analysis basis with \(U^TU=I_n\), and draw
\((h,\Omega,\sigma)\) from the exact product law consisting of the exact
CountSketch hash/sign law and \(s\) exact iid uniform samples from the \(r\)
CountSketch output rows.  Put
\[
  V(h,\Omega)=S_{h,\operatorname{rademacherSignVector}(\Omega)}U.
\]
On the injective-hash event, \(V(h,\Omega)\) has orthonormal columns; hence
each exact preconditioned row has squared norm at most one.  The row-sampling
Loewner radius in the local uniform-row MGF theorem is therefore the explicit
quantity
\[
  L_{\rm CS}=r.
\]
For \(\theta>0\), sample size \(s>0\), and sample failure budget
\(\delta_{\rm sample}\ge0\), assume the exact tail-budget inequality
\[
\begin{aligned}
&n\exp\!\left(
    -\theta s\varepsilon
    +s\frac{\exp(\theta r)-\theta r-1}{r^2}(r^2+1)
  \right)\\
&\quad+
 n\exp\!\left(
    -\theta s\varepsilon
    +s(\exp\theta-\theta-1)(r^2+1)
  \right)
 \le \delta_{\rm sample}.
\end{aligned}
\]
Then the exact sample-Gram event
\[
  -\varepsilon I_n
  \preceq
  \operatorname{Gram}_{\rm unif}(V(h,\Omega),\sigma)-I_n
  \preceq
  \varepsilon I_n
\]
has exact product-law probability at least
\[
  1-\left(\frac{m^2}{r}+\delta_{\rm sample}\right).
\]

The computed sparse CountSketch/sample-Gram endpoint replaces \(V\) by
\[
  \widehat V(h,\Omega)
  =
  \operatorname{flCountSketchSparseApply}
  \bigl(h,\operatorname{rademacherSignVector}(\Omega),U\bigr),
\]
and uses a computed uniform denominator \(\widehat d\) for
\(\sqrt{s/r}\).  Its visible perturbation radius is
\[
\begin{aligned}
  T_{\rm CS,row}^{\rm fp}(h,\Omega,\sigma)
  &=
  \mathcal B_{\rm den}^{\rm unif}
  \bigl(fp,\widehat V,\widehat d,\sigma\bigr)\\
  &\quad+
  \mathcal B_{\rm basis}^{\rm unif}
  \bigl(V,\widehat V,
    \mathcal B_{\rm CS}^{\rm sparse}(fp,h,\omega,U),
    \sigma\bigr),
\end{aligned}
\]
where \(\omega=\operatorname{rademacherSignVector}(\Omega)\), and the budget
symbols are exactly the local Lean definitions
\leanfit{uniformRowSampleGramComputedDenFullFpPerturbBudget}
\leanfit{uniformRowSampleGramBasisPerturbBudget}
\leanfit{countSketchSparseApplyFpAbsBudget}
with the arguments displayed in the formula above.  The first term
charges the computed denominator, rounded row divisions, and rounded
length-\(s\) Gram dot products.  The second term propagates the already proved
sparse CountSketch application entry radii, which charge rounded signed
products and rounded bucket accumulation.  The final fully computed-object
event is
\[
  -(\varepsilon+T_{\rm CS,row}^{\rm fp})I_n
  \preceq
  \widehat G_{\rm row}(h,\Omega,\sigma)-I_n
  \preceq
  (\varepsilon+T_{\rm CS,row}^{\rm fp})I_n,
\]
where \(\widehat G_{\rm row}\) is computed from
\(\widehat V(h,\Omega)\), \(\widehat d\), the sampled rows, and rounded Gram
dot products.  The Lean theorem proves
\[
  \Pr_{h,\Omega,\sigma}
  \!\left[
    -(\varepsilon+T_{\rm CS,row}^{\rm fp})I_n
    \preceq
    \widehat G_{\rm row}-I_n
    \preceq
    (\varepsilon+T_{\rm CS,row}^{\rm fp})I_n
  \right]
  \ge
  1-\left(\frac{m^2}{r}+\delta_{\rm sample}\right).
\]

The actual-input version is now closed for the same collision-free route.  Let
\(A=UC\), where \(U\in\R^{m\times q}\) is an exact analysis witness with
\(U^TU=I_q\) and \(C\in\R^{q\times n}\) is an exact right factor.  The factors
\(U,C\) are not computed by this theorem; the implemented non-probability
quantity is the sparse CountSketch of the actual input
\[
  \widehat Y(h,\Omega)
  =
  \operatorname{flCountSketchSparseApply}
  \bigl(h,\operatorname{rademacherSignVector}(\Omega),A\bigr).
\]
With the same computed denominator \(\widehat d\), define
\[
\begin{aligned}
  T_{\rm CS,A}^{\rm fp}(h,\Omega,\sigma)
  &=
  \mathcal B_{\rm den}^{\rm unif}
  \bigl(fp,\widehat Y,\widehat d,\sigma\bigr)\\
  &\quad+
  \mathcal B_{\rm basis}^{\rm unif}
  \bigl(S_{h,\omega}A,\widehat Y,
    \mathcal B_{\rm CS}^{\rm sparse}(fp,h,\omega,A),
    \sigma\bigr),
\end{aligned}
\]
where the three \(\mathcal B\)-terms are the same concrete Lean budget
definitions displayed above, now instantiated on the actual input matrix
\(A\).  The final actual-input event is
\[
  -\bigl(\varepsilon\,A^TA+T_{\rm CS,A}^{\rm fp}I_n\bigr)
  \preceq
  \widehat G_A(h,\Omega,\sigma)-A^TA
  \preceq
  \varepsilon\,A^TA+T_{\rm CS,A}^{\rm fp}I_n,
\]
where \(\widehat G_A\) is computed from \(\widehat Y(h,\Omega)\),
\(\widehat d\), the sampled rows, and rounded Gram dot products.  Lean proves
\[
  \Pr_{h,\Omega,\sigma}\!\left[\text{the displayed actual-input event}\right]
  \ge
  1-\left(\frac{m^2}{r}+\delta_{\rm sample}\right).
\]
Equivalently, if
\[
  \frac{m^2}{r}+\delta_{\rm sample}\le \delta,
\]
the same fully computed actual-input event has probability at least
\(1-\delta\).  This target-budget form is a separate Lean theorem, so later
input-sparsity theorem surfaces do not need to manually recombine the
collision-free and downstream row-sampling losses.  The order of the
collision-free preprocessing loss is \(\Theta(m^2/r)\), and the combined
failure budget is non-vacuous whenever \(m^2/r+\delta_{\rm sample}<1\).
The exact right-congruence proof transfers the orthonormal-basis event to
\(A^TA\); the floating-point perturbation event is still proved with
probability one from the concrete sparse-apply, computed-denominator,
row-division, and Gram-dot budgets.
There is no additional perturbation-event probability or certificate-existence
assumption: the concrete computed-denominator perturbation event is proved
with probability one from the displayed local budgets.  The probability laws
remain exact mathematical laws.

The same actual-input collision-free endpoint is now closed when the realized
CountSketch signs are stored before sparse application.  For storage mode
\(\rho_{\rm sig}\in\{\mathrm{mul1},\mathrm{add0},\mathrm{sub0}\}\), the
computed sign table satisfies
\[
  |\widehat\omega_k^{\rho_{\rm sig}}-\omega_k|\le u,
\]
and the sparse-apply radius is replaced by
\[
\begin{aligned}
  E^{{\rm CS},\rho_{\rm sig}}_{ij}
  &=
  \bigl(u+\gamma_{b_i(h)}+u\gamma_{b_i(h)}\bigr)
  \sum_{k:h(k)=i}
      |\widehat\omega_k^{\rho_{\rm sig}}|\,|A_{kj}|\\
  &\quad+
  \sum_{k:h(k)=i}
      |\widehat\omega_k^{\rho_{\rm sig}}-\omega_k|\,|A_{kj}|.
\end{aligned}
\]
This term then propagates through the same concrete denominator,
sampled-row division, and sampled-Gram dot-product budget, giving a stored-sign
radius \(T_{{\rm CS,A},\rho_{\rm sig}}^{\rm fp}\).  Lean proves the same
two-sided event centered at \(A^TA\) with probability at least
\[
  1-\left(\frac{m^2}{r}+\delta_{\rm sample}\right),
\]
and the three final wrappers instantiate the sign storage as
\(\fl(\omega_k\cdot1)\), \(\fl(\omega_k+0)\), or \(\fl(\omega_k-0)\).
The probability loss and non-vacuity condition are unchanged; sign storage
affects only the deterministic floating-point radius.

The actual-input collision-free route is also closed for fixed per-bucket
stored-sign traversal orders.  For each realized hash \(h\), the order
\(\pi_i(h)\) for bucket \(i\) is exact discrete implementation data.  Let
\[
  k_{i,t}^{\pi}(h)=\mathrm{bucketIndex}(h,i,\pi_i(h)(t)).
\]
The stored-sign sparse-apply radius becomes
\[
\begin{aligned}
  E^{{\rm CS},\rho_{\rm sig},\pi}_{ij}
  &=
  \bigl(u+\gamma_{b_i(h)}+u\gamma_{b_i(h)}\bigr)
  \sum_{t=0}^{b_i(h)-1}
      |\widehat\omega_{k_{i,t}^{\pi}(h)}^{\rho_{\rm sig}}|\,
      |A_{k_{i,t}^{\pi}(h),j}|\\
  &\quad+
  \sum_{t=0}^{b_i(h)-1}
      |\widehat\omega_{k_{i,t}^{\pi}(h)}^{\rho_{\rm sig}}
        -\omega_{k_{i,t}^{\pi}(h)}|\,
      |A_{k_{i,t}^{\pi}(h),j}|.
\end{aligned}
\]
This radius is propagated through the same concrete denominator,
row-division, and sampled-Gram dot-product budget.  The final event is centered
at \(A^TA\), has probability at least
\[
  1-\left(\frac{m^2}{r}+\delta_{\rm sample}\right),
\]
and the three final wrappers instantiate the sign storage as
\(\fl(\omega_k\cdot1)\), \(\fl(\omega_k+0)\), or \(\fl(\omega_k-0)\).  Exact
objects are \(U,C,A=UC\), the hash/sign/row-sampling laws, \(A^TA\), and the
bucket orders.  Computed objects are the stored signs, sparse signed products,
bucket accumulations in the selected order, denominator, row divisions, and
sampled-Gram dot products.  In fixed realized data regimes the sampled-basis
part has order
\[
  \Theta\!\left(
    \frac{u+\gamma_{b_{\sigma_t}(h)}}{\sqrt{s/r}}
    \sum_{v=0}^{b_{\sigma_t}(h)-1}
      |A_{k_{\sigma_t,v}^{\pi}(h),j}|
  \right),
\]
and the full deterministic radius tends to zero as
\(u,\gamma_s,\gamma_{b_i(h)}\to0\).  No additional \(\delta_{\rm comp}\),
perturbation-event assumption, or certificate-existence assumption is used.

The actual-input collision-free route is also closed for tree-reduced
stored-sign bucket accumulation.  The tree shape \(\mathcal T_i(h)\) for each
realized bucket is an exact discrete implementation choice, while arithmetic
along that tree is computed.  If \(d_i^{\mathcal T}(h)\) is the tree depth, the
stored-sign sparse-apply radius becomes
\[
\begin{aligned}
  E^{{\rm CS},\rho_{\rm sig},\mathcal T}_{ij}
  &=
  \bigl(u+\gamma_{d_i^{\mathcal T}(h)}
        +u\gamma_{d_i^{\mathcal T}(h)}\bigr)
  \sum_{k:h(k)=i}
      |\widehat\omega_k^{\rho_{\rm sig}}|\,|A_{kj}|\\
  &\quad+
  \sum_{k:h(k)=i}
      |\widehat\omega_k^{\rho_{\rm sig}}-\omega_k|\,|A_{kj}|.
\end{aligned}
\]
This radius is propagated through the same concrete denominator, row-division,
and sampled-Gram dot-product budget.  The final event is again centered at
\(A^TA\), has probability at least
\[
  1-\left(\frac{m^2}{r}+\delta_{\rm sample}\right),
\]
and adds no probability loss for choosing the exact tree shapes.  Balanced
bucket trees give the readable sampled-basis order
\[
  \Theta\!\left(
    \frac{u(1+\log(b_i(h)+1))}{\sqrt{s/r}}
    \sum_{k:h(k)=i}|A_{kj}|
  \right)
\]
inside the deterministic radius.

Separately, the non-injective sparse-Gram branch is now closed at the
computed finite-Loewner Gram level.  For an exact input matrix
\(A\in\R^{m\times n}\), an exact CountSketch hash \(h\), and exact
Rademacher trace \(\Omega\), write
\(\omega=\operatorname{rademacherSignVector}(\Omega)\).  The computed
non-probability quantity is
\[
  \widehat G_{\rm CS}^{\rm fp}(h,\Omega)
  =
  \operatorname{flCountSketchSparseGramDot}(fp,h,\omega,A),
\]
formed from rounded sparse signed products, rounded bucket accumulation, and
rounded Gram dot products.  The exact reference object is the analysis Gram
\(A^TA\).  For \(\eta>0\), define the realized radius
\[
  \tau_{\rm CS}^{\rm fp}(h,\Omega,A,\eta)
  =
  \eta+T_{\rm CSGram}^{\rm fp}(h,\Omega,A),
  \qquad
  T_{\rm CSGram}^{\rm fp}
  =
  \operatorname{countSketchSparseGramFullFpPerturbBudget}
  (fp,h,\omega,A).
\]
Lean proves the event
\[
  -\tau_{\rm CS}^{\rm fp} I_n
  \preceq
  \widehat G_{\rm CS}^{\rm fp}-A^TA
  \preceq
  \tau_{\rm CS}^{\rm fp} I_n
\]
under the exact CountSketch hash/sign product law with probability at least
\[
  1-
  \frac{2}{r\eta^2}
  \sum_{j=1}^n\sum_{\ell=1}^n
  \sum_{\substack{a,b=1\\a\ne b}}^m
  (A_{aj}A_{b\ell})^2.
\]
The readable Frobenius upper form is
\[
  \Pr[\text{displayed finite-Loewner event}]
  \ge
  1-\frac{2\|A\|_F^4}{r\eta^2},
\]
and if \(A=U\) with \(U^TU=I_n\), the orthonormal specialization is
\[
  \Pr[\text{displayed finite-Loewner event}]
  \ge
  1-\frac{2n^2}{r\eta^2}.
\]
The target-budget wrappers make these non-vacuity conditions direct theorem
surfaces: for \(0<\delta<1\), if the exact coefficient loss, the Frobenius
loss \(2\|A\|_F^4/(r\eta^2)\), or the orthonormal loss
\(2n^2/(r\eta^2)\) is at most \(\delta\), respectively, then the same
computed finite-Loewner event has probability at least \(1-\delta\).  The
irreducible loss orders are
\[
  \Theta\!\left(
    \frac{\sum_{j,\ell}\sum_{a\ne b}(A_{aj}A_{b\ell})^2}
         {r\eta^2}
  \right),
  \qquad
  \Theta\!\left(\frac{\|A\|_F^4}{r\eta^2}\right),
  \qquad
  \Theta\!\left(\frac{n^2}{r\eta^2}\right)
\]
in the coefficient, Frobenius, and orthonormal forms when the displayed
numerators are nonzero and \(\eta\) is fixed.  These bounds are non-vacuous
whenever the chosen loss is \(<1\), and the loss tends to zero as
\(r\eta^2\) grows with fixed exact input.

The product-grid finite-cover branch now also has an orthonormal-basis
simplification.  Let \(U^TU=I_n\), let
\(\mathcal I\) be the exact one-dimensional grid index set, let
\(g:\mathcal I\to\mathbb R\) be the exact grid, and write
\(z_a(j)=g(a(j))\) for each product-grid vector
\(a:\{1,\ldots,n\}\to\mathcal I\).  Since
\(\|Uz_a\|_2^2=\|z_a\|_2^2\) and \(\|U\|_F^2=n\), the exact product-grid loss
for the computed sparse-Gram event reduces to
\[
  \mathcal L_{\rm pg}^{\rm orth}
  =
  \sum_a\frac{2\|z_a\|_2^4}{r\eta^2}
  +
  \frac{2n^2}{rL^2}.
\]
Lean proves both the exact-law sparse-Gram theorem and the implementation
facing stored-sign theorem with this same probability loss.  In the
stored-sign path the computed non-probability objects are the stored realized
signs, sparse signed products, bucket accumulations, and rounded Gram dot
products; their deterministic contribution remains the already proved
\(T_{{\rm CSGram},\rho_{\rm sig}}^{\rm fp}\) radius.  The irreducible order is
\[
  \Theta\!\left(
    \frac{\sum_a\|z_a\|_2^4/\eta^2+n^2/L^2}{r}
  \right)
\]
for fixed exact grid and positive \(\eta,L\) when the numerator is nonzero,
and the bound is non-vacuous whenever
\(\mathcal L_{\rm pg}^{\rm orth}<1\).

The same orthonormal specialization is now composed through downstream
uniform row sampling for the ordinary stored-sign CountSketch path.  With
exact iid downstream row indices, concrete denominator
\(\widehat d=\fl(\sqrt{s}\,\fl(1/\sqrt r))\), and
\(\eta_{\rm row}>0\), the final event is centered at \(I_n\) and has
computed radius
\[
  \eta+L(2\rho+\rho^2)+\eta_{\rm row}
  +T_{{\rm CS,row},\rho_{\rm sig}}^{\rm fp}(h,\omega,\sigma;U).
\]
Its simplified sufficient loss is
\[
  \mathcal L_{\rm pg,row}^{\rm orth}
  =
  \sum_a\frac{2\|z_a\|_2^4}{r\eta^2}
  +
  \frac{2n^2}{rL^2}
  +
  \frac{r(mn)^2}{s\eta_{\rm row}^2},
\]
with order
\[
  \Theta\!\left(
    \frac{\sum_a\|z_a\|_2^4/\eta^2+n^2/L^2}{r}
    +
    \frac{r m^2 n^2}{s\eta_{\rm row}^2}
  \right)
\]
for fixed exact grid and positive thresholds when the numerator is nonzero.
The exact objects are \(U\), the product grid, exact hash/sign laws, and exact
downstream row laws.  The computed non-probability objects are the stored
signs, sparse signed products, bucket accumulations, the concrete
\(\sqrt{s/r}\) denominator, sampled-row divisions, and sampled-Gram dot
products.  No probability-computation term, perturbation event, auxiliary
certificate, or generic probability-loss assumption is introduced; the result
is non-vacuous whenever \(\mathcal L_{\rm pg,row}^{\rm orth}<1\) and the
deterministic floating-point radius is small.

The same orthonormal downstream product-grid simplification is now closed for
fixed per-bucket traversal orders and for exact supplied bucket trees.  For the
fixed-order route, the exact order is discrete implementation data and the
computed radius charges stored signs, rounded sparse products, accumulation in
that chosen order, the concrete \(\sqrt{s/r}\) denominator, sampled-row
divisions, and sampled-Gram dot products.  For the tree-reduced route, the
exact tree shapes are discrete implementation data and the computed radius
replaces the flat bucket-accumulation factor by the tree-depth
\(\gamma_{d_i^{\mathcal T}(h)}\) factors, under explicit validity
hypotheses.  Both routes use the same sufficient exact loss
\[
  \mathcal L_{\rm pg,row}^{\rm orth}
  =
  \sum_a\frac{2\|z_a\|_2^4}{r\eta^2}
  +\frac{2n^2}{rL^2}
  +\frac{r(mn)^2}{s\eta_{\rm row}^2},
\]
and the same irreducible order
\[
  \Theta\!\left(
    \frac{\sum_a\|z_a\|_2^4/\eta^2+n^2/L^2}{r}
    +
    \frac{r m^2n^2}{s\eta_{\rm row}^2}
  \right).
\]
No probability-computation term, perturbation event, auxiliary certificate, or
generic probability-loss assumption is introduced; the only change from the
ordinary orthonormal downstream theorem is the layout-specific deterministic
floating-point radius.

The collision-free row-sampling results above remain collision-free
sufficient results.  The non-injective sparse-Gram result is
Markov/Frobenius-derived; it does not yet prove the standard optimal
CountSketch concentration/subspace-embedding theorem or the complete
input-sparsity Algorithm 3 theorem.
\end{corollary}

\begin{formalized}
\leanfit{CountSketchHash}
\leanfit{rowGram_preconditionRows_eq_of_left_hasOrthonormalColumns}
\leanfit{preconditionRows_hasOrthonormalColumns_of_left_hasOrthonormalColumns}
\leanfit{countSketchRows}
\leanfit{countSketchHashProbability}
\leanfit{countSketchHashPairCollision}
\leanfit{countSketchHashPairNoCollision}
\leanfit{countSketchProbability}
\leanfit{uniformRowTraceProbability_eventProb_pair_collision_eq_inv}
\leanfit{countSketchHashProbability_eventProb_pairCollision_eq_inv}
\leanfit{FiniteProbability.expectationReal_indicator_eq_eventProb}
\leanfit{countSketchHashProbability_expectationReal_pairCollisionIndicator_eq_inv}
\leanfit{countSketchProbability_expectationReal_pairCollisionIndicator_eq_inv}
\leanfit{countSketchHashProbability_eventProb_pairNoCollision_eq_one_sub_inv}
\leanfit{CountSketchDistinctPair}
\leanfit{rademacherTraceProbability_expectationReal_eq_zero_of_flip_neg}
\leanfit{rademacherTraceProbability_expectationReal_sign_four_eq_zero_of_left_unpaired}
\leanfit{rademacherTraceProbability_expectationReal_sign_four_eq_one_same_pair}
\leanfit{rademacherTraceProbability_expectationReal_sign_four_eq_one_reversed_pair}
\leanfit{rademacherTraceProbability_expectationReal_sign_pair_mul_sign_pair_eq}
\leanfit{rademacherTraceProbability_expectationReal_sq_sum_distinctPair_mul_sign_pair_eq}
\leanfit{countSketchDistinctPairSwap}
\leanfit{countSketchDistinctPair_fourKernel_sum_le_two_sum_sq}
\leanfit{rowGram_preconditionRows_countSketchRows_sub_rowGram_eq_distinctPair_sum}
\leanfit{rademacherTraceProbability_expectationReal_countSketchRows_rowGram_entry_error_sq_eq}
\leanfit{rademacherTraceProbability_expectationReal_countSketchRows_rowGram_entry_error_sq_le}
\leanfit{countSketchHashProbability_expectationReal_collision_coeff_sq_eq_inv_mul}
\leanfit{countSketchProbability_expectationReal_rowGram_entry_error_sq_le}
\leanfit{countSketchProbability_expectationReal_rowGram_frob_error_sq_le}
\leanfit{countSketchProbability_eventProb_rowGram_frob_error_le_ge_one_sub}
\leanfit{countSketchRowGramFrobErrorEvent}
\leanfit{countSketchFlSparseGramDotRowGramFrobErrorEvent}
\leanfit{countSketchRowGramFrobErrorEvent_subset_flSparseGramDotRowGramFrobErrorEvent}
\leanfit{countSketchProbability_eventProb_flSparseGramDot_rowGram_frob_error_le_ge_one_sub}
\leanfit{countSketchDistinctPair_gramCoeffSq_sum_le_frobNormSqRect_sq}
\leanfit{countSketchProbability_eventProb_flSparseGramDot_rowGram_frob_error_le_ge_one_sub_frobNorm}
\leanfit{frobNormSqRect_eq_nat_of_hasOrthonormalColumns}
\leanfit{countSketchProbability_eventProb_flSparseGramDot_rowGram_frob_error_le_ge_one_sub_orthonormal}
\leanfit{countSketchFlSparseGramDotRowGramTwoSidedLoewnerEvent}
\leanfit{countSketchFlSparseGramDotRowGramFrobErrorEvent_subset_twoSidedLoewnerEvent}
\leanfit{countSketchProbability_eventProb_flSparseGramDot_rowGram_twoSidedLoewner_ge_one_sub}
\leanfit{countSketchProbability_eventProb_flSparseGramDot_rowGram_twoSidedLoewner_ge_one_sub_frobNorm}
\leanfit{countSketchProbability_eventProb_flSparseGramDot_rowGram_twoSidedLoewner_ge_one_sub_orthonormal}
\leanfit{countSketchProbability_eventProb_flSparseGramDot_rowGram_twoSidedLoewner_ge_one_sub_delta_of_coeff_budget}
\leanfit{countSketchProbability_eventProb_flSparseGramDot_rowGram_twoSidedLoewner_ge_one_sub_delta_of_frobNorm_budget}
\leanfit{countSketchProbability_eventProb_flSparseGramDot_rowGram_twoSidedLoewner_ge_one_sub_delta_of_orthonormal_budget}
\leanfit{countSketchProbability_eventProb_rowGram_twoSidedLoewner_productGrid_ge_one_sub_sum_gridNorm_add_nat_orthonormal}
\leanfit{countSketchProbability_eventProb_rowGram_twoSidedLoewner_productGrid_ge_one_sub_delta_of_sum_gridNorm_add_nat_orthonormal_budget}
\leanfit{countSketchProbability_eventProb_flSparseGramDot_rowGram_twoSidedLoewner_productGrid_ge_one_sub_sum_gridNorm_add_nat_orthonormal}
\leanfit{countSketchProbability_eventProb_flSparseGramDot_rowGram_twoSidedLoewner_productGrid_ge_one_sub_delta_of_sum_gridNorm_add_nat_orthonormal_budget}
\leanfit{countSketchProbability_eventProb_flSparseGramDotWithStoredSign_rowGram_twoSidedLoewner_productGrid_ge_one_sub_sum_gridNorm_add_nat_orthonormal}
\leanfit{countSketchProbability_eventProb_flSparseGramDotWithStoredSign_rowGram_twoSidedLoewner_productGrid_ge_one_sub_delta_of_sum_gridNorm_add_nat_orthonormal_budget}
\leanfit{countSketchProbability_eventProb_flSparseGramDotWithFlStoredSign_rowGram_twoSidedLoewner_productGrid_ge_one_sub_delta_of_sum_gridNorm_add_nat_orthonormal_budget}
\leanfit{countSketchProbability_eventProb_flSparseGramDotWithFlStoredSignAddZeroRight_rowGram_twoSidedLoewner_productGrid_ge_one_sub_delta_of_sum_gridNorm_add_nat_orthonormal_budget}
\leanfit{countSketchProbability_eventProb_flSparseGramDotWithFlStoredSignSubZeroRight_rowGram_twoSidedLoewner_productGrid_ge_one_sub_delta_of_sum_gridNorm_add_nat_orthonormal_budget}
\leanfit{countSketchUniformRowTraceProbability_eventProb_sparseStoredSignComputedPreconditioned_fl_uniformRowSampleGramDotWithComputedDen_rowGram_twoSidedLoewner_productGrid_ge_one_sub_delta_of_sum_gridNorm_add_nat_add_row_orthonormal_budget}
\leanfit{countSketchUniformRowTraceProbability_eventProb_sparseStoredSignComputedPreconditioned_fl_uniformRowSampleGramDotWithFlSqrtMulInvSqrtDen_rowGram_twoSidedLoewner_productGrid_ge_one_sub_delta_of_sum_gridNorm_add_nat_add_row_orthonormal_budget}
\leanfit{countSketchUniformRowTraceProbability_eventProb_sparseFlStoredSignComputedPreconditioned_fl_uniformRowSampleGramDotWithFlSqrtMulInvSqrtDen_rowGram_twoSidedLoewner_productGrid_ge_one_sub_delta_of_sum_gridNorm_add_nat_add_row_orthonormal_budget}
\leanfit{countSketchUniformRowTraceProbability_eventProb_sparseFlStoredSignAddZeroRightComputedPreconditioned_fl_uniformRowSampleGramDotWithFlSqrtMulInvSqrtDen_rowGram_twoSidedLoewner_productGrid_ge_one_sub_delta_of_sum_gridNorm_add_nat_add_row_orthonormal_budget}
\leanfit{countSketchUniformRowTraceProbability_eventProb_sparseFlStoredSignSubZeroRightComputedPreconditioned_fl_uniformRowSampleGramDotWithFlSqrtMulInvSqrtDen_rowGram_twoSidedLoewner_productGrid_ge_one_sub_delta_of_sum_gridNorm_add_nat_add_row_orthonormal_budget}
\leanfit{countSketchUniformRow_productGrid_orthonormal_coeff_add_frob_add_row_budget}
\leanfit{countSketchUniformRowTraceProbability_eventProb_sparseStoredSignPermutedComputedPreconditioned_fl_uniformRowSampleGramDotWithComputedDen_rowGram_twoSidedLoewner_productGrid_ge_one_sub_delta_of_sum_gridNorm_add_nat_add_row_orthonormal_budget}
\leanfit{countSketchUniformRowTraceProbability_eventProb_sparseStoredSignPermutedComputedPreconditioned_fl_uniformRowSampleGramDotWithFlSqrtMulInvSqrtDen_rowGram_twoSidedLoewner_productGrid_ge_one_sub_delta_of_sum_gridNorm_add_nat_add_row_orthonormal_budget}
\leanfit{countSketchUniformRowTraceProbability_eventProb_sparseFlStoredSignPermutedComputedPreconditioned_fl_uniformRowSampleGramDotWithFlSqrtMulInvSqrtDen_rowGram_twoSidedLoewner_productGrid_ge_one_sub_delta_of_sum_gridNorm_add_nat_add_row_orthonormal_budget}
\leanfit{countSketchUniformRowTraceProbability_eventProb_sparseFlStoredSignAddZeroRightPermutedComputedPreconditioned_fl_uniformRowSampleGramDotWithFlSqrtMulInvSqrtDen_rowGram_twoSidedLoewner_productGrid_ge_one_sub_delta_of_sum_gridNorm_add_nat_add_row_orthonormal_budget}
\leanfit{countSketchUniformRowTraceProbability_eventProb_sparseFlStoredSignSubZeroRightPermutedComputedPreconditioned_fl_uniformRowSampleGramDotWithFlSqrtMulInvSqrtDen_rowGram_twoSidedLoewner_productGrid_ge_one_sub_delta_of_sum_gridNorm_add_nat_add_row_orthonormal_budget}
\leanfit{countSketchUniformRowTraceProbability_eventProb_sparseStoredSignTreeComputedPreconditioned_fl_uniformRowSampleGramDotWithComputedDen_rowGram_twoSidedLoewner_productGrid_ge_one_sub_delta_of_sum_gridNorm_add_nat_add_row_orthonormal_budget}
\leanfit{countSketchUniformRowTraceProbability_eventProb_sparseStoredSignTreeComputedPreconditioned_fl_uniformRowSampleGramDotWithFlSqrtMulInvSqrtDen_rowGram_twoSidedLoewner_productGrid_ge_one_sub_delta_of_sum_gridNorm_add_nat_add_row_orthonormal_budget}
\leanfit{countSketchUniformRowTraceProbability_eventProb_sparseFlStoredSignTreeComputedPreconditioned_fl_uniformRowSampleGramDotWithFlSqrtMulInvSqrtDen_rowGram_twoSidedLoewner_productGrid_ge_one_sub_delta_of_sum_gridNorm_add_nat_add_row_orthonormal_budget}
\leanfit{countSketchUniformRowTraceProbability_eventProb_sparseFlStoredSignAddZeroRightTreeComputedPreconditioned_fl_uniformRowSampleGramDotWithFlSqrtMulInvSqrtDen_rowGram_twoSidedLoewner_productGrid_ge_one_sub_delta_of_sum_gridNorm_add_nat_add_row_orthonormal_budget}
\leanfit{countSketchUniformRowTraceProbability_eventProb_sparseFlStoredSignSubZeroRightTreeComputedPreconditioned_fl_uniformRowSampleGramDotWithFlSqrtMulInvSqrtDen_rowGram_twoSidedLoewner_productGrid_ge_one_sub_delta_of_sum_gridNorm_add_nat_add_row_orthonormal_budget}
\leanfit{countSketchHashInjectiveEvent}
\leanfit{countSketchHash_forall_pairNoCollision_iff_injective}
\leanfit{countSketchHashProbability_eventProb_injective_ge_one_sub_pair_sum}
\leanfit{countSketchDistinctPairBudget_le_square_inv}
\leanfit{countSketchHashProbability_eventProb_injective_ge_one_sub_square_inv}
\leanfit{countSketchRows_preconditionRows_entry}
\leanfit{CountSketchBucket}
\leanfit{countSketchBucketSize}
\leanfit{countSketchBucketIndex}
\leanfit{countSketchBucketExactTerm}
\leanfit{fl_countSketchBucketProduct}
\leanfit{fl_countSketchSparseApplyEntry}
\leanfit{countSketchRows_preconditionRows_bucket_sum_eq}
\leanfit{fl_countSketchSparseApplyEntry_error_bound}
\leanfit{countSketchBucketAbsSum_le_column_abs_sum}
\leanfit{fl_countSketchSparseApplyEntry_error_bound_of_abs_sign_le_one}
\leanfit{fl_countSketchSparseApply}
\leanfit{fl_countSketchSparseApply_entry}
\leanfit{fl_countSketchSparseApply_entry_error_bound}
\leanfit{fl_countSketchSparseApply_entry_error_bound_of_abs_sign_le_one}
\leanfit{rowSketch_abs_perturbed_le_of_abs_error}
\leanfit{rowSketch_abs_perturbed_mul_sum_le_of_abs_error}
\leanfit{rowSketchGram_entry_abs_error_bound_exact_of_entrywise}
\leanfit{rowSketchGram_frob_abs_error_bound_exact_of_entrywise}
\leanfit{fl_rowSketchGramDot}
\leanfit{rowSketchGramDotRoundoffExactBudget}
\leanfit{rowSketchGramAbsPerturbExactBudget}
\leanfit{rowSketchGramFullAbsFpExactBudget}
\leanfit{fl_rowSketchGramDot_roundoff_bound_of_abs_error}
\leanfit{fl_rowSketchGramDot_abs_perturb_bound_exact}
\leanfit{countSketchSparseApplyEntryFpAbsBudget}
\leanfit{countSketchSparseApplyFpAbsBudget}
\leanfit{fl_countSketchSparseGramDot}
\leanfit{countSketchSparseGramDotRoundoffBudget}
\leanfit{countSketchSparseGramApplyPerturbBudget}
\leanfit{countSketchSparseGramFullFpPerturbBudget}
\leanfit{fl_countSketchSparseGramDot_perturb_bound}
\leanfit{countSketchBucketSize_le}
\leanfit{fl_countSketchSparseGramDot_rowGram_perturb_bound_of_hash_injective}
\leanfit{countSketchHashFlSparseGramDotRowGramPerturbEvent}
\leanfit{countSketchHashInjectiveEvent_subset_flSparseGramDot_rowGram_perturbEvent}
\leanfit{countSketchHashProbability_eventProb_flSparseGramDot_rowGram_perturb_ge_one_sub_square_inv}
\leanfit{countSketchFlSparseGramDotRowGramPerturbEvent}
\leanfit{countSketchProbability_injectiveFst_subset_flSparseGramDot_rowGram_perturbEvent}
\leanfit{countSketchProbability_eventProb_flSparseGramDot_rowGram_perturb_ge_one_sub_square_inv}
\leanfit{countSketchRows_hasOrthonormalColumns_of_hash_injective}
\leanfit{rowGram_preconditionRows_countSketchRows_eq_of_hash_injective}
\leanfit{countSketchRows_preconditionRows_hasOrthonormalColumns_of_hash_injective}
\leanfit{rowNormSq_le_one_of_hasOrthonormalColumns}
\leanfit{countSketchUniformRowTraceProbability}
\leanfit{countSketchUniformRowSampleGramTwoSidedEvent}
\leanfit{countSketchUniformRowTraceProbability_eventProb_uniformRowSampleGram_two_sided_finiteLoewnerLe_ge_one_sub_square_inv_add_delta}
\leanfit{countSketchComputedPreconditionedFlUniformRowSampleGramWithComputedDenTwoSidedEvent}
\leanfit{countSketchComputedPreconditionedFlUniformRowPerturbWithComputedDenEvent}
\leanfit{countSketchSparseComputedPreconditionedBasis}
\leanfit{countSketchSparseUniformRowComputedDenPerturbBudget}
\leanfit{countSketchUniformRowTraceProbability_eventProb_sparseComputedPreconditioned_fl_uniformRowComputedDenPerturbEvent_eq_one}
\leanfit{countSketchUniformRowTraceProbability_eventProb_computedPreconditioned_fl_uniformRowSampleGramDotWithComputedDen_two_sided_finiteLoewnerLe_ge_one_sub_of_exact_event}
\leanfit{countSketchUniformRowTraceProbability_eventProb_sparseComputedPreconditioned_fl_uniformRowSampleGramDotWithComputedDen_two_sided_finiteLoewnerLe_ge_one_sub_square_inv_add_delta}
\noindent{\scriptsize\ttfamily\detokenize{countSketchUniformRowTraceProbability_eventProb_sparseComputedPreconditioned_fl}}\par
\noindent{\scriptsize\ttfamily\detokenize{_uniformRowSampleGramDotWithFlSqrtMulInvSqrtDen_two_sided_finiteLoewnerLe_ge_one_sub}}\par
\noindent{\scriptsize\ttfamily\detokenize{_square_inv_add_delta}}\par
\leanfit{uniformRowSampleGram_countSketchRectFactoredInput_error_eq_rightGramCongruence_error}
\leanfit{countSketchUniformRowFactoredInputSampleGramTwoSidedEvent}
\leanfit{countSketchUniformRowSampleGramTwoSidedEvent_subset_factoredInput}
\leanfit{countSketchUniformRowTraceProbability_eventProb_uniformRowFactoredInputSampleGram_two_sided_finiteLoewnerLe_ge_one_sub_square_inv_add_delta}
\leanfit{countSketchComputedPreconditionedFactoredInputFlUniformRowSampleGramWithComputedDenTwoSidedEvent}
\leanfit{countSketchUniformRowTraceProbability_eventProb_computedPreconditioned_factoredInput_fl_uniformRowSampleGramDotWithComputedDen_two_sided_finiteLoewnerLe_ge_one_sub_of_exact_event}
\leanfit{countSketchUniformRowTraceProbability_eventProb_sparseComputedPreconditioned_factoredInput_fl_uniformRowSampleGramDotWithComputedDen_two_sided_finiteLoewnerLe_ge_one_sub_square_inv_add_delta}
\leanfit{countSketchUniformRowTraceProbability_eventProb_sparseComputedPreconditioned_factoredInput_fl_uniformRowSampleGramDotWithComputedDen_two_sided_finiteLoewnerLe_ge_one_sub_delta_of_square_inv_budget}
\noindent{\scriptsize\ttfamily\detokenize{countSketchUniformRowTraceProbability_eventProb_sparseComputedPreconditioned_factoredInput_fl}}\par
\noindent{\scriptsize\ttfamily\detokenize{_uniformRowSampleGramDotWithFlSqrtMulInvSqrtDen_two_sided_finiteLoewnerLe_ge_one_sub}}\par
\noindent{\scriptsize\ttfamily\detokenize{_square_inv_add_delta}}\par
\noindent{\scriptsize\ttfamily\detokenize{countSketchUniformRowTraceProbability_eventProb_sparseComputedPreconditioned_factoredInput_fl}}\par
\noindent{\scriptsize\ttfamily\detokenize{_uniformRowSampleGramDotWithFlSqrtMulInvSqrtDen_two_sided_finiteLoewnerLe_ge_one_sub}}\par
\noindent{\scriptsize\ttfamily\detokenize{_delta_of_square_inv_budget}}\par
\leanfit{rademacherTraceProbability_expectationReal_countSketchRows_entry_mul_eq}
\leanfit{rademacherTraceProbability_expectationReal_countSketchRows_rowGram_eq}
\leanfit{countSketchProbability_expectationReal_rowGram_eq}
\leanfit{countSketchProbability_expectationReal_rowGram_eq_id_of_hasOrthonormalColumns}
\end{formalized}

\begin{corollary}[Exact signed-mixing sample-Gram composition]
Let \(G\in\R^{r\times m}\) and \(U\in\R^{m\times n}\) be exact matrices with
\(r>0\), \(n>0\), \(G^TG=I_m\), and \(U^TU=I_n\).  Let
\(\omega\in\{\pm1\}^m\) be exact Rademacher signs, define
\(\Pi_\omega(i,k)=G_{ik}\omega_k\), and after this preprocessing draw
\(s\) exact iid uniform row samples from the \(r\) rows of
\(\Pi_\omega U\).  Assume the exact entry-square cap
\[
  G_{ik}^2\le \alpha^2,\qquad \alpha>0,
\]
and fix \(B>0\), \(\theta>0\), \(\varepsilon\in\R\), and \(s>0\).  If
\[
  L = r n B^2,
\]
and
\[
\begin{aligned}
T_{\rm sm}(B,\alpha,\theta,\varepsilon)
  &=
  2rn\exp\!\left(-\frac{B^2}{2\alpha^2}\right)\\
  &\quad
  + n\exp\!\left(
      -\theta s\varepsilon
      + s\,
        \frac{\exp(\theta r n B^2)-\theta r n B^2-1}
             {r^2 n^2 B^4}
        \left(r^2n^2B^4+1\right)
    \right)\\
  &\quad
  + n\exp\!\left(
      -\theta s\varepsilon
      + s\,(\exp(\theta)-\theta-1)
        \left(r^2n^2B^4+1\right)
    \right),
\end{aligned}
\]
then the exact product law over signs and uniform row samples satisfies
\[
  \mathbb P\!\left\{
    -\varepsilon I_n
    \preceq
    {\rm Gram}_{s}(\Pi_\omega U)-I_n
    \preceq
    \varepsilon I_n
  \right\}
  \ge
  1-T_{\rm sm}(B,\alpha,\theta,\varepsilon).
\]
Equivalently, the Lean theorem permits separate budgets
\(\delta_{\rm pre}\) and \(\delta_{\rm samp}\) that dominate the first line
and the two matrix-Chernoff lines, respectively, and proves the bound
\(1-(\delta_{\rm pre}+\delta_{\rm samp})\).

In the asymptotic regime where \(r,n,s\) may vary, \(\alpha,B,\theta\) are
positive parameters, and the displayed exponents are kept in a non-vacuous
negative range, the failure expression has order
\[
\Theta\!\left(
rn e^{-B^2/(2\alpha^2)}
+ n e^{-\theta s\varepsilon
      + s\frac{e^{\theta rnB^2}-\theta rnB^2-1}{r^2n^2B^4}
        (r^2n^2B^4+1)}
+ n e^{-\theta s\varepsilon
      + s(e^\theta-\theta-1)(r^2n^2B^4+1)}
\right).
\]
The non-vacuity check is explicit: the probability statement is meaningful
when this total failure term is below \(1\).  This result is exact
probability and exact arithmetic.  The sampling laws are exact by project
convention, while applying/storing \(G\), forming \(\Pi_\omega U\), and
forming the sample Gram in floating point remain separate computed-quantity
obligations.
\end{corollary}

\begin{formalized}
\leanfit{signedMixingRows_preconditionRows_hasOrthonormalColumns}
\leanfit{signedMixingUniformRowTraceProbability}
\leanfit{signedMixingUniformRowSampleGramTwoSidedEvent}
\leanfit{signedMixingUniformRowTraceProbability_eventProb_uniformRowSampleGram_two_sided_finiteLoewnerLe_ge_one_sub_delta_of_entry_sq_le_uniform}
\end{formalized}

\begin{corollary}[Computed finite signed-mixing preconditioner formation]
Let \(G\in\R^{r\times m}\) be an exact deterministic mixing table and let
\(\omega\in\{\pm1\}^m\) be the exact realized sign vector.  The exact
analysis object is
\[
  \Pi_\omega(i,j)=G_{ij}\omega_j,
  \qquad\text{equivalently}\qquad
  \Pi_\omega=G\operatorname{diag}(\omega).
\]
The Lean theorem first proves this equivalence exactly:
\[
  \operatorname{preconditionRows}(G,\operatorname{diag}(\omega))
  =
  \operatorname{signedMixingRows}(G,\omega).
\]
The computable floating-point quantity in this proposition is the stored
matrix
\[
  \widehat\Pi_\omega
  =
  \operatorname{fl\_matMul}\bigl(G,\operatorname{diag}(\omega)\bigr),
\]
formed by rounded matrix multiplication.  The probabilities and the
Rademacher law are exact by project convention; the displayed bound charges
only this non-probability preconditioner-formation arithmetic.  If
\(\gamma_m\) is valid for the \(m\)-term product, then every entry satisfies
\[
  |\widehat\Pi_{\omega,ij}-G_{ij}\omega_j|
  \le
  \gamma_m\sum_{k=1}^m |G_{ik}|\,
    |\operatorname{diag}(\omega)_{kj}|.
\]
Since the diagonal has only one nonzero entry and \(|\omega_j|=1\), this
displayed sum has the readable exact-arithmetic simplification
\[
  \gamma_m |G_{ij}|.
\]
Thus, for fixed exact \(G\) and exact signs, the entrywise formation error has
order
\[
  \Theta(\gamma_m |G_{ij}|),
\]
and the non-vacuity check is immediate: it tends to zero as
\(\gamma_m\to0\).  The more general Lean constructor
\(\texttt{ComputedPreconditioner.flSignedMixing}\) also accepts computed
certificates for \(G\) and for the sign vector, so their storage/generation
errors are included before the rounded product is formed.
\end{corollary}

\begin{formalized}
\leanfit{preconditionRows_diagMatrix_eq_signedMixingRows}
\leanfit{ComputedPreconditioner.flSignedMixing}
\leanfit{ComputedPreconditioner.flSignedMixingExactFactors}
\leanfit{ComputedPreconditioner.flSignedMixingExactFactors_entry_error_bound}
\end{formalized}

\begin{corollary}[Computed finite signed-mixing sample Gram]
In the exact signed-mixing setting above, assume additionally that the
floating-point model has valid \(\gamma_m\) for length-\(m\) products and
valid \(\gamma_s\) for length-\(s\) Gram dot products.  The probability laws
for \(\omega\) and for the uniform row trace \(S=(S_1,\ldots,S_s)\) are exact.
All non-probability quantities below are computed.

For each realized sign vector define the exact analysis preconditioner
\(\Pi_\omega=G\operatorname{diag}(\omega)\).  The implementation first forms
\[
  \widehat\Pi_\omega=\operatorname{fl}(G\operatorname{diag}(\omega)),
\]
with exact-factor componentwise certificate
\[
  e^\Pi_{ia}
  =
  \gamma_m\sum_{b=1}^m |G_{ib}|\,
    |\operatorname{diag}(\omega)_{ba}|
  =
  \gamma_m |G_{ia}|.
\]
It then computes
\[
  \widehat V_\omega
  =
  \operatorname{fl}(\widehat\Pi_\omega U).
\]
The fully expanded entrywise budget used by Lean for this second product is
\[
  E_{ij}^{V}
  =
  \gamma_m\sum_{a=1}^m |\widehat\Pi_{\omega,ia}|\,|U_{aj}|
  +
  \sum_{a=1}^m e^\Pi_{ia}|U_{aj}|.
\]
Thus \( |\widehat V_{\omega,ij}-(\Pi_\omega U)_{ij}|\le E^V_{ij}\), and no
existence assumption for this perturbation is used.

Let \(d=\sqrt{s/r}\).  For a sampled row trace \(S\), define the exact
rescaled sketches
\[
  B_{tj}=\frac{(\Pi_\omega U)_{S_tj}}{d},
  \qquad
  \widetilde B_{tj}=\frac{\widehat V_{\omega,S_tj}}{d}.
\]
For the exact-denominator row-scaling path, the proved sample-dependent FP
radius is
\[
\begin{aligned}
  \tau_{\rm sm}^{\rm fp}(\omega,S)
  &=
  \Bigl((2u+u^2)+\gamma_s(1+u)^2\Bigr)
  \left\|
    \left[
      \sum_{t=1}^s |\widetilde B_{tj}|\,|\widetilde B_{tk}|
    \right]_{j,k=1}^n
  \right\|_F \\
  &\quad+
  \left\|
    \left[
      \sum_{t=1}^s
      \left(
        \frac{E^V_{S_tj}}{d}\,|\widetilde B_{tk}|
        +
        |B_{tj}|\,\frac{E^V_{S_tk}}{d}
      \right)
    \right]_{j,k=1}^n
  \right\|_F .
\end{aligned}
\]
The first line charges rounded row divisions and rounded length-\(s\) Gram dot
products after \(\widehat V_\omega\) has been formed.  The second line charges
the exact sampled-Gram drift from \(\widehat V_\omega\) back to
\(\Pi_\omega U\).

If the denominator is itself computed, let \(\widehat d\ne0\) satisfy
\(|\widehat d-d|\le\eta_d\).  Define
\[
  \widehat B_{tj}
  =
  \operatorname{fl}\!\left(\frac{\widehat V_{\omega,S_tj}}{\widehat d}\right),
  \qquad
  E^d_{tj}
  =
  \left|\frac{\widehat V_{\omega,S_tj}}{\widehat d}\right|u
  +
  \frac{|\widehat V_{\omega,S_tj}|\eta_d}{|\widehat d|\,d}.
\]
The computed-denominator radius is
\[
\begin{aligned}
  \tau_{\rm sm,den}^{\rm fp}(\omega,S)
  &=
  \left\|
    \left[
      \sum_{t=1}^s
      \left(E^d_{tj}|\widehat B_{tk}|+
            |\widetilde B_{tj}|E^d_{tk}\right)
    \right]_{j,k=1}^n
  \right\|_F \\
  &\quad+
  \left\|
    \left[
      \gamma_s\sum_{t=1}^s |\widehat B_{tj}|\,|\widehat B_{tk}|
    \right]_{j,k=1}^n
  \right\|_F \\
  &\quad+
  \left\|
    \left[
      \sum_{t=1}^s
      \left(
        \frac{E^V_{S_tj}}{d}\,|\widetilde B_{tk}|
        +
        |B_{tj}|\,\frac{E^V_{S_tk}}{d}
      \right)
    \right]_{j,k=1}^n
  \right\|_F .
\end{aligned}
\]
Consequently, with the same exact preprocessing and sampling budgets
\(\delta_{\rm pre}\) and \(\delta_{\rm samp}\) as the exact signed-mixing
sample-Gram theorem,
\[
\begin{aligned}
  \mathbb P\!\left\{
    -(\varepsilon+\tau_{\rm sm,den}^{\rm fp}(\omega,S))I_n
    \preceq
    \widehat G_{\omega,S}-I_n
    \preceq
    (\varepsilon+\tau_{\rm sm,den}^{\rm fp}(\omega,S))I_n
  \right\}
  \ge
  1-(\delta_{\rm pre}+\delta_{\rm samp}),
\end{aligned}
\]
where \(\widehat G_{\omega,S}\) is the Gram matrix computed from
\(\widehat V_\omega\) using the computed denominator \(\widehat d\), rounded
row divisions, and rounded length-\(s\) dot products.  There is no additional
\(\delta_{\rm comp}\): the computed perturbation event is proved with
probability one from the displayed concrete budget.

For the final concrete denominator path, Lean instantiates
\[
  \widehat d
  =
  \fl_{\rm mul}\!\left(
    \fl_{\rm sqrt}(s),
    \fl_{\rm div}(1,\fl_{\rm sqrt}(r))
  \right),
  \qquad
  d=\sqrt{s/r},
\]
with
\[
  \eta_d
  =
  \sqrt{s/r}\,
  \frac{4u+3u^2+u^3}{1-u}.
\]
Thus the displayed \(\tau_{\rm sm,den}^{\rm fp}\) has no hidden
denominator variable in the concrete theorem surface; the older
arbitrary-\(\widehat d\) statement remains only reusable infrastructure for
other denominator routines.

For fixed \(G,U,r,m,n,s\) and sampled rows, in the regime
\(u,\gamma_m,\gamma_s,\eta_d/d\to0\) with \(|\widehat d|\) bounded away from
zero,
\[
  \tau_{\rm sm,den}^{\rm fp}
  =
  \Theta\!\left(
    \left(u+\gamma_s+\frac{\eta_d}{d}\right)
    \left\|
      \left[\sum_t |\widehat B_{tj}|\,|\widehat B_{tk}|\right]_{jk}
    \right\|_F
    +
    \gamma_m
    \left\|
      \left[
        \sum_t
        \left(
          \frac{C_{S_tj}}{d}|\widetilde B_{tk}|
          + |B_{tj}|\frac{C_{S_tk}}{d}
        \right)
      \right]_{jk}
    \right\|_F
  \right),
\]
where the primitive coefficient
\[
  C_{ij}
  =
  \sum_{a=1}^m
  \bigl(|\Pi_{\omega,ia}|+|G_{ia}|\bigr)|U_{aj}|
\]
is the first-order expansion of \(E^V_{ij}/\gamma_m\).  The non-vacuity check
is explicit: for fixed exact data, the FP radius tends to \(0\) as
\(u,\gamma_m,\gamma_s,\eta_d\to0\), and the probability lower bound remains
the exact signed-mixing lower bound.
\end{corollary}

\begin{formalized}
\leanfit{signedMixingComputedLeftPreconditionedBasis}
\leanfit{signedMixingExactFactorPreconditioner}
\leanfit{signedMixingComputedLeftUniformRowPerturbBudget}
\leanfit{signedMixingComputedLeftUniformRowComputedDenPerturbBudget}
\leanfit{signedMixingUniformRowTraceProbability_eventProb_computedLeftPreconditioned_fl_uniformRowPerturbEvent_eq_one}
\leanfit{signedMixingUniformRowTraceProbability_eventProb_computedLeftPreconditioned_fl_uniformRowComputedDenPerturbEvent_eq_one}
\leanfit{signedMixingUniformRowTraceProbability_eventProb_exactFactorComputedLeftPreconditioned_fl_uniformRowSampleGramDot_two_sided_finiteLoewnerLe_ge_one_sub_delta_of_entry_sq_le_uniform}
\leanfit{signedMixingUniformRowTraceProbability_eventProb_exactFactorComputedLeftPreconditioned_fl_uniformRowSampleGramDotWithComputedDen_two_sided_finiteLoewnerLe_ge_one_sub_delta_of_entry_sq_le_uniform}
\leanfit{signedMixingUniformRowTraceProbability_eventProb_exactFactorComputedLeftPreconditioned_fl_uniformRowSampleGramDotWithFlSqrtMulInvSqrtDen_two_sided_finiteLoewnerLe_ge_one_sub_delta_of_entry_sq_le_uniform}
\end{formalized}

\begin{corollary}[Computed actual-input finite signed-mixing sample Gram]
Let \(r,m,q,n,s\) be positive dimensions.  Let
\(G\in\R^{r\times m}\) be an exact deterministic rectangular isometry,
\(G^TG=I_m\), satisfying the exact entry cap \(G_{ia}^2\le\alpha^2\).
Let the exact input matrix \(A\in\R^{m\times n}\) admit the exact analysis
factorization
\[
  A=UC,\qquad U\in\R^{m\times q},\quad U^TU=I_q,\quad C\in\R^{q\times n}.
\]
The matrices \(U\) and \(C\) are exact analysis witnesses in this theorem;
the algorithmic computation below is performed with the actual input \(A\).
The Rademacher signs \(\omega\) and the uniform row trace
\(S=(S_1,\ldots,S_s)\) are sampled from exact probability laws.  Every
non-probability quantity displayed next is computed in floating point.

Define the exact signed-mixing preconditioner
\[
  \Pi_\omega=G\operatorname{diag}(\omega),
\]
and compute
\[
  \widehat\Pi_{\omega,ia}
  =
  \operatorname{fl}\!\left(
    \sum_{b=1}^m G_{ib}\operatorname{diag}(\omega)_{ba}
  \right).
\]
The exact-factor preconditioner-formation radius is
\[
  e^\Pi_{ia}
  =
  \gamma_m\sum_{b=1}^m |G_{ib}|\,
    |\operatorname{diag}(\omega)_{ba}|
  =
  \gamma_m |G_{ia}|.
\]
The implementation then computes the preconditioned actual input
\[
  \widehat Y_{\omega}
  =
  \operatorname{fl}(\widehat\Pi_\omega A).
\]
Since \(A_{aj}=\sum_{\ell=1}^q U_{a\ell}C_{\ell j}\), the fully expanded
entrywise product budget is
\[
\begin{aligned}
  E^A_{ij}
  &=
  \gamma_m\sum_{a=1}^m |\widehat\Pi_{\omega,ia}|\,|A_{aj}|
  +
  \sum_{a=1}^m e^\Pi_{ia}|A_{aj}| \\
  &=
  \gamma_m
  \sum_{a=1}^m
  \left(
    \left|
      \operatorname{fl}\!\left(
        \sum_{b=1}^m G_{ib}\operatorname{diag}(\omega)_{ba}
      \right)
    \right|
    +|G_{ia}|
  \right)
  \left|\sum_{\ell=1}^q U_{a\ell}C_{\ell j}\right|.
\end{aligned}
\]
Lean proves
\[
  |\widehat Y_{\omega,ij}-(\Pi_\omega A)_{ij}|
  \le E^A_{ij};
\]
this is an instantiated arithmetic bound, not a perturbation-existence
assumption.

Set \(d=\sqrt{s/r}\).  For the sampled rows define
\[
  Y_{tj}=\frac{(\Pi_\omega A)_{S_tj}}{d},
  \qquad
  \widetilde Y_{tj}=\frac{\widehat Y_{\omega,S_tj}}{d}.
\]
Let \(\widehat d\ne0\) be the computed row-scale denominator with
\(|\widehat d-d|\le\eta_d\), and define
\[
  \widehat B_{tj}
  =
  \operatorname{fl}\!\left(
    \frac{\widehat Y_{\omega,S_tj}}{\widehat d}
  \right),
  \qquad
  E^d_{tj}
  =
  \left|\frac{\widehat Y_{\omega,S_tj}}{\widehat d}\right|u
  +
  \frac{|\widehat Y_{\omega,S_tj}|\eta_d}{|\widehat d|\,d}.
\]
In the final concrete denominator theorem this \(\widehat d\) is not an
abstract input.  It is
\[
  \widehat d
  =
  \fl_{\rm mul}\!\left(
    \fl_{\rm sqrt}(s),
    \fl_{\rm div}(1,\fl_{\rm sqrt}(r))
  \right),
  \qquad
  d=\sqrt{s/r},
  \qquad
  \eta_d
  =
  \sqrt{s/r}\,
  \frac{4u+3u^2+u^3}{1-u}.
\]
The formula below is therefore already irreducible after this substitution:
the only remaining floating-point parameters are \(u,\gamma_m,\gamma_s\) and
the explicitly displayed dimensions and exact data.
The computed-denominator perturbation radius used by the theorem is
\[
\begin{aligned}
  \tau_A^{\rm fp}(\omega,S)
  &=
  \left\|
    \left[
      \sum_{t=1}^s
      \left(E^d_{tj}|\widehat B_{tk}|+
            |\widetilde Y_{tj}|E^d_{tk}\right)
    \right]_{j,k=1}^n
  \right\|_F \\
  &\quad+
  \left\|
    \left[
      \gamma_s\sum_{t=1}^s |\widehat B_{tj}|\,|\widehat B_{tk}|
    \right]_{j,k=1}^n
  \right\|_F \\
  &\quad+
  \left\|
    \left[
      \sum_{t=1}^s
      \left(
        \frac{E^A_{S_tj}}{d}\,|\widetilde Y_{tk}|
        +
        |Y_{tj}|\,\frac{E^A_{S_tk}}{d}
      \right)
    \right]_{j,k=1}^n
  \right\|_F .
\end{aligned}
\]
The first line charges the computed denominator and rounded row divisions,
the second line charges rounded length-\(s\) Gram dot products, and the third
line charges the preconditioner/input matrix-product error.

Let \(\widehat{\mathcal G}_{\omega,S}\) be the matrix obtained by rounded
Gram dot products from the sampled rows \(\widehat B_{t,*}\).  With the same
exact preprocessing and sampling failures
\(\delta_{\rm pre}\) and \(\delta_{\rm samp}\) as the exact signed-mixing
sample-Gram theorem,
\[
  \mathbb P\!\left\{
    -\varepsilon A^TA-\tau_A^{\rm fp}(\omega,S)I_n
    \preceq
    \widehat{\mathcal G}_{\omega,S}-A^TA
    \preceq
    \varepsilon A^TA+\tau_A^{\rm fp}(\omega,S)I_n
  \right\}
  \ge
  1-(\delta_{\rm pre}+\delta_{\rm samp}).
\]
There is no additional \(\delta_{\rm comp}\): the computed perturbation event
is proved with probability one from the displayed radius.
Equivalently, if
\(\delta_{\rm pre}+\delta_{\rm samp}\le\delta\), the same fully computed
event has probability at least \(1-\delta\); this total-budget form is now a
separate Lean theorem, so later theorem surfaces do not need to recombine the
two exact-probability losses manually.

For fixed exact \(G,U,C,r,m,q,n,s\) and fixed sampled rows, assume
\(|\widehat d|\) is bounded below by a positive constant multiple of \(d\).
In the regime \(u,\gamma_m,\gamma_s,\eta_d/d\to0\),
\[
  \tau_A^{\rm fp}
  =
  \Theta\!\left(
    \left(u+\gamma_s+\frac{\eta_d}{d}\right)
    \left\|
      \left[\sum_t |\widehat B_{tj}|\,|\widehat B_{tk}|\right]_{jk}
    \right\|_F
    +
    \gamma_m
    \left\|
      \left[
        \sum_t
        \left(
          \frac{K_{S_tj}}{d}|\widetilde Y_{tk}|
          + |Y_{tj}|\frac{K_{S_tk}}{d}
        \right)
      \right]_{jk}
    \right\|_F
  \right),
\]
where the irreducible coefficient is
\[
  K_{ij}
  =
  \sum_{a=1}^m
  \left(
    \left|
      \operatorname{fl}\!\left(
        \sum_{b=1}^m G_{ib}\operatorname{diag}(\omega)_{ba}
      \right)
    \right|
    +|G_{ia}|
  \right)
  \left|\sum_{\ell=1}^q U_{a\ell}C_{\ell j}\right|.
\]
Equivalently, \(E^A_{ij}=\gamma_m K_{ij}\).  Hence the FP radius tends to
\(0\) as \(u,\gamma_m,\gamma_s,\eta_d\to0\) for fixed exact data and a
nonzero denominator, while the probability lower bound remains exactly the
signed-mixing lower bound.
\end{corollary}

\begin{formalized}
\leanfit{preconditionRows_preconditionColumns_assoc_rect}
\leanfit{uniformRowSampleGram_rectFactoredInput_error_eq_rightGramCongruence_error}
\leanfit{signedMixingUniformRowFactoredInputSampleGramTwoSidedEvent}
\leanfit{signedMixingUniformRowSampleGramTwoSidedEvent_subset_factoredInput}
\leanfit{signedMixingUniformRowTraceProbability_eventProb_uniformRowFactoredInputSampleGram_two_sided_finiteLoewnerLe_ge_one_sub_delta_of_entry_sq_le_uniform}
\leanfit{signedMixingComputedPreconditionedFactoredInputFlUniformRowSampleGramWithComputedDenTwoSidedEvent}
\leanfit{signedMixingUniformRowTraceProbability_eventProb_computedPreconditioned_factoredInput_fl_uniformRowSampleGramDotWithComputedDen_two_sided_finiteLoewnerLe_ge_one_sub_of_exact_event}
\leanfit{signedMixingUniformRowTraceProbability_eventProb_exactFactorComputedLeftPreconditioned_factoredInput_fl_uniformRowSampleGramDotWithComputedDen_two_sided_finiteLoewnerLe_ge_one_sub_delta_of_entry_sq_le_uniform}
\leanfit{signedMixingUniformRowTraceProbability_eventProb_exactFactorComputedLeftPreconditioned_factoredInput_fl_uniformRowSampleGramDotWithComputedDen_two_sided_finiteLoewnerLe_ge_one_sub_delta_of_total_budget}
\leanfit{signedMixingUniformRowTraceProbability_eventProb_exactFactorComputedLeftPreconditioned_factoredInput_fl_uniformRowSampleGramDotWithFlSqrtMulInvSqrtDen_two_sided_finiteLoewnerLe_ge_one_sub_delta_of_entry_sq_le_uniform}
\leanfit{signedMixingUniformRowTraceProbability_eventProb_exactFactorComputedLeftPreconditioned_factoredInput_fl_uniformRowSampleGramDotWithFlSqrtMulInvSqrtDen_two_sided_finiteLoewnerLe_ge_one_sub_delta_of_total_budget}
\end{formalized}

\begin{corollary}[Coordinate-Hoeffding row-norm bound]
Assume \(U\in\R^{m\times n}\) has orthonormal columns and
\[
  H_{ik}^2 = \frac1m
  \qquad\text{for all }i,k .
\]
For \(D_\omega=\operatorname{diag}(\omega)\), \(B\ge0\), and
\(\lambda>0\), define
\[
  \eta(B,\lambda,m)
  =
  2\exp(-\lambda B)
  \exp\left(\frac{\lambda^2}{2m}\right).
\]
Then
\[
  \mathbb P\left\{
    \forall i,\quad
    \Norm{(HD_\omega U)_{i,*}}^2\le nB^2
  \right\}
  \ge
  1-\sum_{i=1}^m\sum_{j=1}^n \eta(B,\lambda,m).
\]
\end{corollary}


\begin{formalized}
\leanfit{signedHadamard_entry_coeff_sum_sq_eq_inv}
\leanfit{rademacherTraceProbability_eventProb_entry_abs_signedHadamard_le_ge_one_sub_exp_sq_bound}
\leanfit{rademacherTraceProbability_eventProb_forall_entry_abs_signedHadamard_le_ge_one_sub_sum_exp_sq_bound}
\leanfit{rademacherTraceProbability_eventProb_forall_rowNormSq_signedHadamard_le_ge_one_sub_sum_exp_sq_bound}
\end{formalized}

\begin{corollary}[Coordinate-Hoeffding leverage probabilities]
Assume, in addition, that \(H\) is orthogonal and \(n>0\).  Let
\[
  p_i(HD_\omega U)
  =
  \frac{\Norm{(HD_\omega U)_{i,*}}^2}{n}
\]
be the equation~(6) leverage probability of the preconditioned basis.  With
the same \(\eta(B,\lambda,m)\) as above,
\[
  \mathbb P\left\{
    \forall i,\quad p_i(HD_\omega U)\le B^2
  \right\}
  \ge
  1-\sum_{i=1}^m\sum_{j=1}^n \eta(B,\lambda,m).
\]
In delta-budget form, if
\[
  \sum_{i=1}^m\sum_{j=1}^n \eta(B,\lambda,m)\le \delta,
\]
then the probability is at least \(1-\delta\).
\end{corollary}


\begin{formalized}
\leanfit{rademacherTraceProbability_eventProb_forall_leverageScoreProb_signedHadamard_le_ge_one_sub_sum_exp_sq_bound}
\leanfit{rademacherTraceProbability_eventProb_forall_leverageScoreProb_signedHadamard_le_ge_one_sub_delta}
\end{formalized}

\begin{lemma}[Uniform row one-step Loewner bound after preconditioning]
Let \(V\in\R^{m\times n}\) have orthonormal columns.  For uniform sampling of
one row, define the rank-one estimator
\[
  X_i(V)=m\,v_i v_i^T .
\]
Then \(X_i(V)\succeq0\) and
\[
  \frac1m\sum_{i=1}^m X_i(V)=I_n .
\]
Moreover, if the leverage probability
\[
  p_i(V)=\frac{\Norm{v_i}^2}{n}
\]
satisfies \(p_i(V)\le B^2\), then
\[
  X_i(V)\preceq mnB^2\,I_n .
\]
Consequently, under the signed-Hadamard hypotheses of the previous corollary,
if the coordinate-Hoeffding failure budget is at most \(\delta\), then
\[
  \mathbb P\left\{
    \forall i,\quad X_i(HD_\omega U)\preceq mnB^2 I_n
  \right\}
  \ge 1-\delta .
\]
\end{lemma}


\begin{formalized}
\lean{Algorithms/RandNLA/UniformRowSampling.lean}
\leanfit{uniformRowOuterGramSample}
\leanfit{finiteQuadraticForm_uniformRowOuterGramSample_eq}
\leanfit{finitePSD_uniformRowOuterGramSample}
\leanfit{uniform_rowOuterGramSample_mean_eq_id}
\leanfit{uniformRowOuterGramSample_finiteLoewnerLe_of_leverageScoreProb_le}
\lean{Algorithms/RandNLA/Preconditioning.lean}
\leanfit{rademacherTraceProbability_eventProb_forall_uniformRowOuterGramSample_signedHadamard_finiteLoewnerLe_ge_one_sub_delta}
\end{formalized}

\begin{theorem}[Uniform row sample-average concentration]
Let \(V\in\R^{m\times n}\) have orthonormal columns, and let
\[
  X_i(V)=m\,v_i v_i^T .
\]
Draw iid uniform rows \(i_1,\ldots,i_s\) and define the sampled Gram average
\[
  G_s(V)=\frac1s\sum_{t=1}^s X_{i_t}(V).
\]
Assume that for every row \(i\),
\[
  X_i(V)\preceq L I_n,
  \qquad
  m\Norm{v_i}^2\le L .
\]
Let \(\theta>0\), \(\varepsilon>0\), and \(s>0\).  Define
\[
  \beta_+(\theta,L)
  =
  \frac{\exp(\theta L)-\theta L-1}{L^2},
  \qquad
  \beta_-(\theta)
  =
  \exp(\theta)-\theta-1,
  \qquad
  W=L^2+1 .
\]
If
\[
  e^{-\theta s\varepsilon}\, n\,e^{s\beta_+(\theta,L)W}
  +
  e^{-\theta s\varepsilon}\, n\,e^{s\beta_-(\theta)W}
  \le \delta ,
\]
then
\[
  \mathbb P\left\{
    G_s(V)-I_n\preceq \varepsilon I_n
    \ \text{and}\
    -(G_s(V)-I_n)\preceq \varepsilon I_n
  \right\}
  \ge 1-\delta .
\]
\end{theorem}


\begin{formalized}
\lean{Algorithms/RandNLA/UniformRowSamplingMGF.lean}
\leanfit{uniformRowProb}
\leanfit{uniformRowTraceProbability}
\leanfit{uniformRowSampleGram}
\leanfit{uniformRowSampleGram_sub_finiteIdMatrix_eq_centered_average}
\leanfit{uniformRowOuterGramSample_centered_square_expectationCStarMatrix_le}
\leanfit{uniformRowTraceProbabilityFiniteRealTraceMGFLogBound_centered_le}
\leanfit{uniformRowTraceProbabilityFiniteRealTraceMGFLogBound_neg_centered_le}
\leanfit{uniformRowTraceProbability_eventProb_uniformRowSampleGram_two_sided_finiteLoewnerLe_ge_one_sub_delta_of_tail_budget}
\end{formalized}

\begin{corollary}[Signed-Hadamard preprocessing plus uniform row sampling]
Let \(U\in\R^{m\times n}\) have orthonormal columns.  Let \(H\in\R^{m\times m}\)
be orthogonal and flat, \(H_{ik}^2=1/m\).  Draw a Rademacher sign vector
\(\omega\), set \(V_\omega=HD_\omega U\), and then draw iid uniform rows
\(i_1,\ldots,i_s\).  Let
\[
  G_s(V_\omega)=\frac1s\sum_{t=1}^s m\,v_{\omega,i_t}v_{\omega,i_t}^T.
\]
Assume the coordinate-Hoeffding preprocessing budget
\[
  \sum_{i,j}2\exp(-\lambda B)\exp\!\left(\frac{\lambda^2}{2m}\right)
  \le \delta_{\mathrm{pre}}
\]
and define
\[
  L=mnB^2,\qquad
  \beta_+(\theta,L)=\frac{\exp(\theta L)-\theta L-1}{L^2},
  \qquad
  \beta_-(\theta)=\exp(\theta)-\theta-1 .
\]
If the uniform-sampling tail budget
\[
  e^{-\theta s\varepsilon}\,n\,e^{s\beta_+(\theta,L)(L^2+1)}
  +
  e^{-\theta s\varepsilon}\,n\,e^{s\beta_-(\theta)(L^2+1)}
  \le \delta_{\mathrm{sample}},
\]
then
\[
  \mathbb P\left\{
    G_s(V_\omega)-I_n\preceq \varepsilon I_n
    \ \text{and}\
    -(G_s(V_\omega)-I_n)\preceq \varepsilon I_n
  \right\}
  \ge 1-\delta_{\mathrm{pre}}-\delta_{\mathrm{sample}} .
\]
\end{corollary}


\begin{formalized}
\lean{Analysis/FiniteProbability.lean}
\leanfit{FiniteProbability.prod}
\leanfit{FiniteProbability.prod_eventProb_inter_dependent_ge_one_sub_add}
\lean{Algorithms/RandNLA/UniformRowSamplingComposition.lean}
\leanfit{signedHadamardUniformRowTraceProbability}
\leanfit{signedHadamardUniformRowSampleGramTwoSidedEvent}
\leanfit{signedHadamardUniformRowTraceProbability_eventProb_uniformRowSampleGram_two_sided_finiteLoewnerLe_ge_one_sub_delta}
\end{formalized}

\begin{corollary}[Source-sharp SRHT preprocessing plus uniform row sampling]
Under the same signed-Hadamard/SRHT setup, let \(t>0\) and define
\[
  S(t)=\sqrt{n/m}+t,\qquad
  L(t)=mS(t)^2 .
\]
Assume the source-sharp preprocessing budget
\[
  m\exp\!\left(-\frac{mt^2}{8}\right)
  \le \delta_{\mathrm{pre}}
\]
and define
\[
  \beta_+(\theta,L)=\frac{\exp(\theta L)-\theta L-1}{L^2},
  \qquad
  \beta_-(\theta)=\exp(\theta)-\theta-1 .
\]
If the uniform-sampling tail budget with \(L=L(t)\),
\[
  e^{-\theta s\varepsilon}\,n\,
    e^{s\beta_+(\theta,L)(L^2+1)}
  +
  e^{-\theta s\varepsilon}\,n\,
    e^{s\beta_-(\theta)(L^2+1)}
  \le \delta_{\mathrm{sample}},
\]
then
\[
  \mathbb P\left\{
    G_s(V_\omega)-I_n\preceq \varepsilon I_n
    \ \text{and}\
    -(G_s(V_\omega)-I_n)\preceq \varepsilon I_n
  \right\}
  \ge 1-\delta_{\mathrm{pre}}-\delta_{\mathrm{sample}} .
\]
\end{corollary}


\begin{formalized}
\lean{Algorithms/RandNLA/UniformRowSamplingComposition.lean}
\leanfit{signedHadamardUniformRowTraceProbability_eventProb_uniformRowSampleGram_two_sided_finiteLoewnerLe_ge_one_sub_delta_srht}
\leanfit{signedHadamardUniformRowTraceProbability_eventProb_uniformRowSampleGram_two_sided_finiteLoewnerLe_ge_one_sub_delta_srht_log_preprocess}
\end{formalized}

\begin{corollary}[Floating-point signed-Hadamard preprocessing plus uniform row sampling]
Under the hypotheses of the coordinate-Hoeffding signed-Hadamard preprocessing
plus uniform row-sampling corollary above, assume also a floating-point model
\(\mathrm{fp}\) with \(\gamma_s\) valid.  Let
\[
  \widehat G_{\omega,S}
  =
  \operatorname{fl}\!\left(\widehat V_{\omega,S}^{T}
    \widehat V_{\omega,S}\right)
\]
be the Gram matrix obtained by forming the uniformly rescaled sampled rows with
rounded division and computing each Gram entry with the local rounded dot
product.  Then the same event probability lower bound holds:
\[
  \mathbb P\left\{
    -(\varepsilon+\tau_{\omega,S})I_n
    \preceq
    \widehat G_{\omega,S}-I_n
    \preceq
    (\varepsilon+\tau_{\omega,S})I_n
  \right\}
  \ge
  1-\delta_{\mathrm{pre}}-\delta_{\mathrm{sample}},
\]
where
\[
  \tau_{\omega,S}
  =
  \texttt{uniformRowSampleGramFullFpPerturbBudget}
  (\mathrm{fp},s,V_\omega,S)
\]
is the explicit sample-dependent budget for rounded row scaling and rounded
Gram dot products.
\end{corollary}


\begin{formalized}
\lean{Analysis/MatrixAlgebra.lean}
\leanfit{finiteLoewnerLe_two_sided_add_of_frobNorm_le}
\lean{Algorithms/RandNLA/UniformRowSamplingFP.lean}
\leanfit{uniformRowSampleScaleDen}
\leanfit{uniformRowSampleSketch}
\leanfit{fl_uniformRowSampleSketch}
\leanfit{rowSketchGram_uniformRowSampleSketch_eq_uniformRowSampleGram}
\leanfit{fl_uniformRowSampleGramDot}
\leanfit{uniformRowSampleGramFullFpPerturbBudget}
\leanfit{fl_uniformRowSampleGramDot_perturb_bound}
\leanfit{signedHadamardFlUniformRowSampleGramTwoSidedEvent}
\leanfit{signedHadamardUniformRowTraceProbability_eventProb_fl_uniformRowSampleGramDot_two_sided_finiteLoewnerLe_ge_one_sub_delta}
\end{formalized}

\begin{corollary}[Computed-left \(\widehat V\) replacement for Corollary 171]
The floating-point corollary above is an exact-\(V_\omega=HD_\omega U\)
sampling-law result: the row probabilities are uniform and no leverage table is
computed, but the displayed \(\tau_{\omega,S}\) only charges rounded row scaling
and rounded Gram dot products after the exact preprocessed basis has been
supplied.  For an implementation that forms the basis, let
\[
  P_\omega=HD_\omega,\qquad
  V_\omega=P_\omega U,
\]
and let \(\widehat P_\omega\) be a computed transform with a
\(\mathrm{ComputedPreconditioner}\) certificate
\[
  |\widehat P_{\omega,ia}-P_{\omega,ia}|
  \le E^P_{\omega,ia},\qquad E^P_{\omega,ia}\ge0 .
\]
Define the actually formed basis and its entrywise error budget by
\[
  \widehat V_{\omega}
  =
  \operatorname{fl}(\widehat P_\omega U),
\]
\[
  E^V_{\omega,ij}
  =
  \gamma_m\sum_a |\widehat P_{\omega,ia}|\,|U_{aj}|
  +
  \sum_a E^P_{\omega,ia}|U_{aj}|.
\]
If \(\gamma_m\) is valid, then
\[
  |\widehat V_{\omega,ij}-V_{\omega,ij}|
  \le E^V_{\omega,ij}
  \qquad\text{for all } \omega,i,j .
\]
For a sampled row trace \(S\), set
\[
  \tau^{\mathrm{left}}_{\omega,S}
  =
  \texttt{uniformRowSampleGramFullFpPerturbBudget}
  (\mathrm{fp},s,\widehat V_\omega,S)
  +
  \texttt{uniformRowSampleGramBasisPerturbBudget}
  (V_\omega,\widehat V_\omega,E^V_\omega,S).
\]
Then the computed-left perturbation event has probability one:
\[
  \mathbb P\!\left[
    \Frob{
      \operatorname{fl}(\widehat V_{\omega,S}^T\widehat V_{\omega,S})
      -G_s(V_\omega)
    }
    \le \tau^{\mathrm{left}}_{\omega,S}
  \right]=1 .
\]
Consequently Corollary~171 has an implementation-facing form with
\[
  \mathbb P\left\{
    -(\varepsilon+\tau^{\mathrm{left}}_{\omega,S})I_n
    \preceq
    \operatorname{fl}(\widehat V_{\omega,S}^T\widehat V_{\omega,S})-I_n
    \preceq
    (\varepsilon+\tau^{\mathrm{left}}_{\omega,S})I_n
  \right\}
  \ge
  1-\delta_{\mathrm{pre}}-\delta_{\mathrm{sample}},
\]
once the exact signed-Hadamard/uniform-row event from Corollary~171 is
available.  The remaining implementation obligation is not hidden in this
statement: a concrete SRHT/FHT/sign generator or stored transform must
instantiate \(\widehat P_\omega\) and its
\(\mathrm{ComputedPreconditioner}\) certificate.
\end{corollary}

\begin{formalized}
\leanfit{flPreconditionRowsWithComputedLeftEntryErrorBudget}
\leanfit{fl_preconditionRowsWithComputedLeft_entry_error_budget_bound}
\leanfit{uniformRowSampleGram_frob_error_bound_of_basis_entrywise_abs}
\leanfit{signedHadamardComputedLeftPreconditionedBasis}
\leanfit{signedHadamardComputedLeftUniformRowPerturbBudget}
\leanfit{signedHadamardUniformRowTraceProbability_eventProb_computedLeftPreconditioned_fl_uniformRowPerturbEvent_eq_one}
\leanfit{signedHadamardUniformRowTraceProbability_eventProb_computedPreconditioned_fl_uniformRowSampleGramDot_two_sided_finiteLoewnerLe_ge_one_sub_delta}
\end{formalized}

\begin{corollary}[Computed-input \(\widehat U\) replacement for Corollary 171]
In the preceding computed-left setting, suppose the basis or singular-vector
table used by the implementation is itself computed.  Let \(\widehat U\) carry
a \(\mathrm{ComputedMatrix}\) certificate
\[
  |\widehat U_{aj}-U_{aj}|\le E^U_{aj},\qquad E^U_{aj}\ge0 .
\]
The implemented basis is now
\[
  \widehat V_\omega
  =
  \operatorname{fl}(\widehat P_\omega\widehat U).
\]
If \(\gamma_m\) is valid, then for every \(\omega,i,j\),
\[
  |\widehat V_{\omega,ij}-(HD_\omega U)_{ij}|
  \le
  \gamma_m\sum_a|\widehat P_{\omega,ia}|\,|\widehat U_{aj}|
  +
  \sum_aE^P_{\omega,ia}|\widehat U_{aj}|
  +
  \sum_a|P_{\omega,ia}|\,E^U_{aj}.
\]
Define \(E^{V,\mathrm{in}}_\omega\) by the right-hand side.  Let
\[
\begin{aligned}
  B^{\mathrm{fp}}_{\omega,S}
  &:=
  \mathcal B_{\mathrm{fp}}(\mathrm{fp},s,\widehat V_\omega,S),\\
  B^{\mathrm{basis}}_{\omega,S}
  &:=
  \mathcal B_{\mathrm{basis}}
  (HD_\omega U,\widehat V_\omega,E^{V,\mathrm{in}}_\omega,S),\\
  \tau^{\mathrm{in}}_{\omega,S}
  &:=
  B^{\mathrm{fp}}_{\omega,S}+B^{\mathrm{basis}}_{\omega,S}.
\end{aligned}
\]
Here \(\mathcal B_{\mathrm{fp}}\) and \(\mathcal B_{\mathrm{basis}}\) are the
Lean budgets
\texttt{uniformRowSampleGramFullFpPerturbBudget} and
\texttt{uniformRowSampleGramBasisPerturbBudget}, respectively.
Then the computed-input perturbation event has probability one under the same
exact Rademacher/uniform-row law:
\[
  \mathbb P\!\left[
    \Frob{
      \operatorname{fl}(\widehat V_{\omega,S}^T\widehat V_{\omega,S})
      -G_s(HD_\omega U)
    }
    \le \tau^{\mathrm{in}}_{\omega,S}
  \right]=1 .
\]
Thus a computed QR/SVD/singular-vector routine enters Algorithm 3 through the
matrix certificate \(E^U\), not through any probability-law perturbation.
\end{corollary}

\begin{formalized}
\leanfit{ComputedMatrix}
\leanfit{fl_preconditionRowsWithComputedLeftAndInput_total_error_bound}
\leanfit{flPreconditionRowsWithComputedLeftInputEntryErrorBudget}
\leanfit{fl_preconditionRowsWithComputedLeftInput_entry_error_budget_bound}
\leanfit{signedHadamardComputedLeftInputPreconditionedBasis}
\leanfit{signedHadamardComputedLeftInputUniformRowPerturbBudget}
\leanfit{signedHadamardUniformRowTraceProbability_eventProb_computedLeftInputPreconditioned_fl_uniformRowPerturbEvent_eq_one}
\end{formalized}

\begin{corollary}[Deterministic floating-point radius for the uniform row sketch]
In the setting of the preceding corollary, suppose a deterministic scalar
\(\tau\) dominates the sample-dependent floating-point budget for every joint
preprocessing and sampling outcome:
\[
  \texttt{uniformRowSampleGramFullFpPerturbBudget}
  (\mathrm{fp},s,V_\omega,S)
  \le \tau .
\]
Then the same probability lower bound holds with the fixed radius
\(\varepsilon+\tau\):
\[
  \mathbb P\left\{
    -(\varepsilon+\tau)I_n
    \preceq
    \widehat G_{\omega,S}-I_n
    \preceq
    (\varepsilon+\tau)I_n
  \right\}
  \ge
  1-\delta_{\mathrm{pre}}-\delta_{\mathrm{sample}}.
\]
Moreover, if \(R\ge 0\) and every sampled row satisfies
\[
  \Norm{(V_\omega)_{S_t,*}}^2\le R,
\]
then the sample-dependent budget is bounded deterministically by
\[
  \texttt{uniformRowSampleGramFullFpConstBudget}
  (\mathrm{fp},s,mR).
\]
The two parts of this fixed budget have the closed forms
\[
  n(2u+u^2)mR
  \qquad\text{and}\qquad
  n\,\gamma_s(1+u)^2mR,
\]
for row scaling and dot products, respectively.
\end{corollary}


\begin{formalized}
\lean{Analysis/MatrixAlgebra.lean}
\leanfit{finiteLoewnerLe_smul_finiteIdMatrix_mono}
\lean{Algorithms/RandNLA/UniformRowSamplingFP.lean}
\leanfit{uniformRowSampleGramFpPerturbConstBudget}
\leanfit{uniformRowSampleGramDotProductConstBudget}
\leanfit{uniformRowSampleGramFullFpConstBudget}
\leanfit{uniformRowSampleGramFpPerturbConstBudget_eq_nat_mul}
\leanfit{uniformRowSampleGramDotProductConstBudget_eq_nat_mul}
\leanfit{abs_mul_entry_le_rowNormSq}
\leanfit{uniformRowSampleSketch_abs_mul_sum_le_of_rowNormSq_le}
\leanfit{uniformRowSampleGramFullFpPerturbBudget_le_const_of_sample_rowNormSq_le}
\leanfit{signedHadamardFlUniformRowSampleGramTwoSidedConstBudgetEvent}
\leanfit{signedHadamardUniformRowTraceProbability_eventProb_fl_uniformRowSampleGramDot_two_sided_finiteLoewnerLe_ge_one_sub_delta_constBudget}
\end{formalized}

\begin{corollary}[Computed-\(\widehat V\) deterministic-radius replacement for
Corollary 172]
Let \(\widehat V_\omega\) be the implemented preprocessed basis.  Suppose the
exact signed-Hadamard/uniform-row event has probability at least
\[
  1-\delta_{\mathrm{pre}}-\delta_{\mathrm{sample}},
\]
and suppose the constant-radius computed-basis perturbation event satisfies
\[
  \mathbb P\!\left[
    \Frob{
      \operatorname{fl}(\widehat V_{\omega,S}^T\widehat V_{\omega,S})
      -G_s(V_\omega)
    }
    \le \tau_{\mathrm{comp}}
  \right]\ge 1-\delta_{\mathrm{comp}} .
\]
Then
\[
  \mathbb P\left\{
    -(\varepsilon+\tau_{\mathrm{comp}})I_n
    \preceq
    \operatorname{fl}(\widehat V_{\omega,S}^T\widehat V_{\omega,S})-I_n
    \preceq
    (\varepsilon+\tau_{\mathrm{comp}})I_n
  \right\}
  \ge
  1-\delta_{\mathrm{pre}}-\delta_{\mathrm{sample}}
    -\delta_{\mathrm{comp}} .
\]
The radius \(\tau_{\mathrm{comp}}\) must dominate every computed quantity used
to produce and use \(\widehat V_\omega\): transform generation or storage,
Hadamard/sign arithmetic, any computed input basis, the computed
\(\sqrt{s/m}\) denominator when one is used, rounded row scaling, and the final
Gram dot products.
\end{corollary}

\begin{formalized}
\leanfit{signedHadamardComputedPreconditionedFlUniformRowPerturbConstBudgetEvent}
\leanfit{signedHadamardUniformRowTraceProbability_eventProb_computedPreconditioned_fl_uniformRowSampleGramDot_two_sided_finiteLoewnerLe_ge_one_sub_delta_constBudget}
\end{formalized}

\begin{corollary}[Source-sharp SRHT floating-point uniform row sampling]
In the source-sharp SRHT setting, let
\[
  S(t)=\sqrt{n/m}+t,\qquad
  \tau_{\mathrm{SRHT}}(t)
  =
  \texttt{uniformRowSampleGramFullFpConstBudget}
  (\mathrm{fp},s,mS(t)^2).
\]
Under the same preprocessing and sample trace-MGF budgets as the source-sharp
exact theorem, the rounded sample Gram matrix satisfies
\[
  \mathbb P\left\{
    -(\varepsilon+\tau_{\mathrm{SRHT}}(t))I_n
    \preceq
    \widehat G_{\omega,S}-I_n
    \preceq
    (\varepsilon+\tau_{\mathrm{SRHT}}(t))I_n
  \right\}
  \ge
  1-\delta_{\mathrm{pre}}-\delta_{\mathrm{sample}} .
\]
The logarithmic version substitutes
\[
  t=\sqrt{8\log(m/\delta_{\mathrm{pre}})/m}.
\]
\end{corollary}


\begin{formalized}
\lean{Algorithms/RandNLA/UniformRowSamplingFP.lean}
\leanfit{signedHadamardUniformRowTraceProbability_eventProb_fl_uniformRowSampleGramDot_two_sided_finiteLoewnerLe_ge_one_sub_delta_constBudget_srht}
\leanfit{signedHadamardUniformRowTraceProbability_eventProb_fl_uniformRowSampleGramDot_two_sided_finiteLoewnerLe_ge_one_sub_delta_constBudget_srht_log_preprocess}
\end{formalized}

\begin{corollary}[Intermediate computed-\(\widehat V\) source-sharp SRHT
transfer]
In the source-sharp SRHT setting, the exact theorem controls
\(V_\omega=HD_\omega U\).  Let \(\widehat V_\omega\) be the basis produced by
an implementation.  This is reusable infrastructure for arbitrary
\(\widehat V_\omega\); the final concrete actual-input SRHT theorem below
instantiates the transform, the input multiplication, the denominator routine,
the row scaling, and the Gram dot products.  If the exact SRHT preprocessing
and sample trace-MGF
budgets give probability at least
\[
  1-\delta_{\mathrm{pre}}-\delta_{\mathrm{sample}},
\]
and the implemented-basis perturbation event satisfies
\[
  \mathbb P\!\left[
    \Frob{
      \operatorname{fl}(\widehat V_{\omega,S}^T\widehat V_{\omega,S})
      -G_s(V_\omega)
    }
    \le \tau_{\mathrm{SRHT,comp}}(t)
  \right]\ge 1-\delta_{\mathrm{comp}},
\]
then
\[
  \mathbb P\left\{
    -(\varepsilon+\tau_{\mathrm{SRHT,comp}}(t))I_n
    \preceq
    \operatorname{fl}(\widehat V_{\omega,S}^T\widehat V_{\omega,S})-I_n
    \preceq
    (\varepsilon+\tau_{\mathrm{SRHT,comp}}(t))I_n
  \right\}
  \ge
  1-\delta_{\mathrm{pre}}-\delta_{\mathrm{sample}}
    -\delta_{\mathrm{comp}} .
\]
The logarithmic source wrapper substitutes
\(t=\sqrt{8\log(m/\delta_{\mathrm{pre}})/m}\).  The exact SRHT concentration
result is unchanged; the new term
\(\tau_{\mathrm{SRHT,comp}}(t)\) is the place where the concrete
FFT/Hadamard routine, sign storage, scale denominators, computed basis or
singular-vector construction, row scaling, and Gram dot products must be
charged.
\end{corollary}

\begin{formalized}
\leanfit{signedHadamardUniformRowTraceProbability_eventProb_computedPreconditioned_fl_uniformRowSampleGramDot_two_sided_finiteLoewnerLe_ge_one_sub_delta_srht}
\leanfit{signedHadamardUniformRowTraceProbability_eventProb_computedPreconditioned_fl_uniformRowSampleGramDot_two_sided_finiteLoewnerLe_ge_one_sub_delta_constBudget_srht}
\leanfit{signedHadamardUniformRowTraceProbability_eventProb_computedPreconditioned_fl_uniformRowSampleGramDot_two_sided_finiteLoewnerLe_ge_one_sub_delta_constBudget_srht_log_preprocess}
\end{formalized}

\begin{corollary}[Computed actual-input SRHT with computed uniform denominator]
Let \(A=U_A C_A\) be the exact Algorithm 3 input factorization used only in the
analysis, where \(U_A^TU_A=I_r\).  The algorithm does not compute \(U_A\) or
\(C_A\); it computes with \(A\).  On the concrete scaled sign-pattern,
stored-sign SRHT path, form
\[
  \widehat P_\omega
  =
  \operatorname{fl}(\widehat H\widehat D_\omega),
  \qquad
  \widehat Y_\omega
  =
  \operatorname{fl}(\widehat P_\omega A),
  \qquad
  \widehat G_A(\omega,S)
  =
  \operatorname{fl}(\widehat B_{\omega,S}^T\widehat B_{\omega,S}),
\]
where \(\widehat B_{\omega,S}\) uses the computed nonzero denominator
\[
  \widehat d_\star
  =
  \operatorname{fl}\!\left(
    \operatorname{fl}(\sqrt{s})\,
    \operatorname{fl}(1/\operatorname{fl}(\sqrt m))
  \right),
\]
with certified relative denominator radius
\((4u+3u^2+u^3)/(1-u)\) for \(\sqrt{s/m}\).  The exact Rademacher and
uniform-row laws remain exact.  Every non-probability computation is charged
by the explicit
sample-dependent budget
\[
  \tau_A^{\mathrm{cd}}(\omega,S)
  =
  \texttt{uniformRowSampleGramComputedDenFullFpPerturbBudget}
  (\mathrm{fp},\widehat Y_\omega,\widehat d,S)
  +
  \texttt{uniformRowSampleGramBasisPerturbBudget}
  (P_\omega A,\widehat Y_\omega,E_\omega,S),
\]
where \(E_\omega\) is the entrywise product budget for
\(\operatorname{fl}(\widehat P_\omega A)\).  With
\[
  t=\sqrt{8\log(m/\delta_{\mathrm{pre}})/m},
  \qquad
  L=m\left(\sqrt{r/m}+t\right)^2,
\]
and the usual exact uniform-sampling trace-MGF budget at dimension \(r\), the
fully computed input Gram satisfies
\[
  \mathbb P\left[
    (1-\varepsilon)A^TA-\tau_A^{\mathrm{cd}}(\omega,S)I_n
    \preceq
    \widehat G_A(\omega,S)
    \preceq
    (1+\varepsilon)A^TA+\tau_A^{\mathrm{cd}}(\omega,S)I_n
  \right]
  \ge
  1-\delta_{\mathrm{pre}}-\delta_{\mathrm{sample}}.
\]
There is no perturbation-event hypothesis and no \(\delta_{\mathrm{comp}}\):
the computed-denominator perturbation event is proved with probability one for
this concrete path.
\end{corollary}

\begin{formalized}
\leanfit{rightGramCongruence}
\leanfit{uniformRowSampleGram_factoredInput_error_eq_rightGramCongruence_error}
\leanfit{signedHadamardUniformRowFactoredInputSampleGramTwoSidedEvent}
\leanfit{signedHadamardUniformRowSampleGramTwoSidedEvent_subset_factoredInput}
\leanfit{signedHadamardUniformRowTraceProbability_eventProb_uniformRowFactoredInputSampleGram_two_sided_finiteLoewnerLe_ge_one_sub_delta_srht_log_preprocess}
\leanfit{signedHadamardComputedPreconditionedFactoredInputFlUniformRowSampleGramWithComputedDenTwoSidedEvent}
\leanfit{signedHadamardUniformRowTraceProbability_eventProb_computedPreconditioned_factoredInput_fl_uniformRowSampleGramDotWithComputedDen_two_sided_finiteLoewnerLe_ge_one_sub_of_exact_event}
\leanfit{signedHadamardUniformRowTraceProbability_eventProb_scaledPatternStoredSignComputedLeftPreconditioned_factoredInput_fl_uniformRowSampleGramDotWithComputedDen_two_sided_finiteLoewnerLe_ge_one_sub_delta_srht_log_preprocess}
\leanfit{uniformRowFlSqrtMulInvSqrtScaleDen}
\leanfit{uniformRowFlSqrtMulInvSqrtScaleDen_den_abs_error}
\leanfit{signedHadamardUniformRowTraceProbability_eventProb_scaledPatternStoredSignComputedLeftPreconditioned_fl_uniformRowSampleGramDotWithFlSqrtMulInvSqrtDen_two_sided_finiteLoewnerLe_ge_one_sub_delta_srht_log_preprocess}
\leanfit{signedHadamardUniformRowTraceProbability_eventProb_scaledPatternStoredSignComputedLeftPreconditioned_factoredInput_fl_uniformRowSampleGramDotWithFlSqrtMulInvSqrtDen_two_sided_finiteLoewnerLe_ge_one_sub_delta_srht_log_preprocess}
\end{formalized}

\begin{corollary}[Exact-stored computed-left SRHT certificate]
Let \(P_\omega=HD_\omega\), and suppose the implementation supplies this
realized preconditioner exactly as its stored transform:
\[
  \widehat P_\omega=P_\omega,\qquad
  E^P_{\omega,ik}=0 .
\]
Define the implemented preconditioned basis by
\[
  \widehat V_\omega=\operatorname{fl}(\widehat P_\omega U).
\]
For every entry,
\[
  |\widehat V_{\omega,ij}-(P_\omega U)_{ij}|
  \le
  \gamma_m
  \sum_{k=1}^m |(P_\omega)_{ik}|\,|U_{kj}|.
\]
Consequently, under the same positive-dimension and \(\gamma\)-validity
hypotheses as the computed-left transfer, the joint exact
Rademacher/uniform-row probability law satisfies
\[
  \mathbb P\!\left[
    \Frob{
      \operatorname{fl}(\widehat V_{\omega,S}^T\widehat V_{\omega,S})
      -G_s(P_\omega U)
    }
    \le
    \tau_{\mathrm{row/Gram}}(\widehat V_\omega,S)
      +\tau_{\mathrm{basis}}(P_\omega U,\widehat V_\omega,E^\mathrm{exact}_\omega,S)
  \right]=1,
\]
where \(E^\mathrm{exact}_{\omega,ij}\) is the displayed matrix-product
rounding budget above.  Thus this specialization charges rounded formation of
\(\widehat V_\omega\), rounded row scaling, and rounded Gram dot products, but
it charges no preprocessing-matrix storage/generation error.  Non-exact FHT
routines, or sign-storage formats other than the modeled sign-copy
certificates below, must still instantiate a nonzero
\(\widehat P_\omega-P_\omega\) certificate.
\end{corollary}

\begin{formalized}
\leanfit{ComputedPreconditioner.exact}
\leanfit{flPreconditionRowsWithComputedLeftEntryErrorBudget_exact}
\leanfit{fl_preconditionRowsWithExactLeft_entry_error_bound}
\leanfit{signedHadamardExactStoredPreconditioner}
\leanfit{signedHadamardComputedLeftPreconditionedBasisEntryErrorBudget_exactStored}
\leanfit{signedHadamardUniformRowTraceProbability_eventProb_exactStoredComputedLeftPreconditioned_fl_uniformRowPerturbEvent_eq_one}
\end{formalized}

\begin{corollary}[Expected row norm after flat signed-Hadamard preprocessing]
Assume \(U\in\R^{m\times n}\) has orthonormal columns and
\[
  H_{ik}^2 = \frac1m
  \qquad\text{for all }i,k .
\]
For \(D_\omega=\operatorname{diag}(\omega)\), each row of \(HD_\omega U\)
has expected squared Euclidean norm
\[
  \mathbb E\Norm{(HD_\omega U)_{i,*}}^2
  =
  \frac{n}{m}.
\]
Consequently,
\[
  \mathbb E\Norm{(HD_\omega U)_{i,*}}
  \le
  \sqrt{\frac{n}{m}} .
\]
\end{corollary}


\begin{formalized}
\leanfit{rademacherTraceProbability_expectationReal_sign_eq_zero}
\leanfit{rademacherTraceProbability_expectationReal_sign_mul_eq_ite}
\leanfit{rademacherTraceProbability_expectationReal_sq_sum_mul_sign_eq_sum_sq}
\leanfit{rademacherTraceProbability_expectationReal_exp_sum_mul_sign_eq_prod}
\leanfit{rademacherTraceProbability_eventProb_sum_mul_sign_le_ge_one_sub_exp_mul_prod}
\leanfit{HadamardFlat}
\leanfit{signedHadamardPreconditionRows_entry}
\leanfit{rademacherTraceProbability_expectationReal_rowNormSq_signedHadamard_eq}
\leanfit{rademacherTraceProbability_expectationReal_sqrt_rowNormSq_signedHadamard_le}
\end{formalized}

\begin{lemma}[Signed-Hadamard row-norm convexity input]
For a fixed row \(i\), define
\[
  f_i(x)=\Norm{(H\operatorname{diag}(x)U)_{i,*}},
  \qquad x\in\R^m .
\]
Then \(f_i\) is convex:
\[
  f_i(\theta x+(1-\theta)y)
  \le
  \theta f_i(x)+(1-\theta)f_i(y)
  \qquad (0\le\theta\le1).
\]
\end{lemma}


\begin{formalized}
\lean{Analysis/MatrixAlgebra.lean}
\leanfit{FiniteVecConvex}
\leanfit{vecNorm2_linear_combination_convex}
\lean{Algorithms/RandNLA/Preconditioning.lean}
\leanfit{signedHadamard_row_vecNorm2_convex}
\end{formalized}

\begin{lemma}[Unit-cube to Rademacher scaling]
Let \(F:\R^m\to\R\) be convex and \(L\)-Lipschitz with \(L>0\).  Define
\[
  G(x)=\frac{1}{2L}\,F(2x-\mathbf 1).
\]
Then \(G\) is convex and 1-Lipschitz.  Consequently, Ledoux's \([0,1]^m\)
upper-tail statement with exponent \(\exp(-t^2/2)\) converts to Tropp's
Rademacher-sign statement with exponent \(\exp(-t^2/8)\).
\end{lemma}


\begin{formalized}
\lean{Analysis/MatrixAlgebra.lean}
\leanfit{FiniteVecLipschitzWith}
\leanfit{unitCubeToRademacherVec}
\leanfit{finiteVecConvex_scaled_unitCubeToRademacher}
\leanfit{finiteVecLipschitzWith_scaled_unitCubeToRademacher}
\end{formalized}

\begin{lemma}[Signed-Hadamard row-norm Lipschitz input]
Assume \(U\in\R^{m\times n}\) has orthonormal columns and
\[
  H_{ik}^2=\frac1m
  \qquad\text{for all }i,k .
\]
For deterministic sign-like vectors \(x,y\in\R^m\), define
\[
  f_i(x)=\Norm{(H\operatorname{diag}(x)U)_{i,*}} .
\]
Then
\[
  |f_i(x)-f_i(y)|
  \le
  \sqrt{\frac1m}\,\Norm{x-y} .
\]
\end{lemma}


\begin{formalized}
\lean{Algorithms/RandNLA/RowSamplingLeverage.lean}
\leanfit{hasOrthonormalColumns_vecNorm2Sq_mul_vec_eq}
\leanfit{hasOrthonormalColumns_transpose_mul_vecNorm2Sq_le}
\lean{Analysis/MatrixAlgebra.lean}
\leanfit{abs_vecNorm2_sub_le_vecNorm2_sub}
\lean{Algorithms/RandNLA/Preconditioning.lean}
\leanfit{signedHadamard_row_vecNorm2_lipschitz}
\end{formalized}

\begin{lemma}[Signed-Hadamard row-norm self-bounding]
Assume \(U\in\R^{m\times n}\) has orthonormal columns and
\[
  H_{ik}^2=\frac1m
  \qquad\text{for all }i,k .
\]
For a fixed row \(i\) and Rademacher sign trace \(\omega\), define
\[
  F_i(\omega)
  =
  \Norm{(H\operatorname{diag}(\varepsilon(\omega))U)_{i,*}},
\]
and let \(\omega^k\) denote \(\omega\) with only coordinate \(k\) flipped.
Then the positive coordinate-flip drops satisfy
\[
  \sum_{k=1}^m
  \max\{F_i(\omega)-F_i(\omega^k),0\}^2
  \le
  \frac4m .
\]
\end{lemma}


\begin{formalized}
\lean{Analysis/MatrixAlgebra.lean}
\leanfit{vecNorm2_inv_smul_self_of_pos}
\leanfit{vecInnerProduct_inv_smul_self_eq_norm}
\leanfit{vecNorm2_sub_le_inner_unit_diff}
\lean{Algorithms/RandNLA/Preconditioning.lean}
\leanfit{rademacherSignVector_flip_self}
\leanfit{rademacherSignVector_flip_of_ne}
\leanfit{rademacherSignVector_sub_flip}
\leanfit{signedHadamard_row_inner_sq_sum_eq_inv_mul}
\leanfit{signedHadamard_row_vec_sub_flip}
\leanfit{signedHadamard_row_vecNorm2_positive_flip_sq_sum_le}
\end{formalized}

\begin{lemma}[Positive-drop exponential-tilt bridge]
Let \(X\) be a real-valued function on the finite Rademacher cube and let
\(\omega^i\) denote coordinate \(i\) flipped.  If, pointwise in \(\omega\),
\[
  \sum_i \max\{X(\omega)-X(\omega^i),0\}^2 \le B ,
\]
then for every \(\lambda\ge0\),
\[
\begin{aligned}
&\sum_i
  \mathbb E\left[
  \left(
    e^{\lambda X(\omega)/2}
    -
    e^{\lambda X(\omega^i)/2}
  \right)^2
  \right] \\
&\qquad\le
  \frac{B}{2}\lambda^2\,
  \mathbb E e^{\lambda X(\omega)} .
\end{aligned}
\]
\end{lemma}


\begin{formalized}
\lean{Analysis/FiniteProbability.lean}
\leanfit{real_exp_half_sub_sq_le_quarter_mul_sq_mul_exp_of_le}
\leanfit{real_exp_half_sub_sq_le_lam_sq_quarter_pair_pos}
\lean{Algorithms/RandNLA/Preconditioning.lean}
\leanfit{rademacherTraceProbability_flip_tilt_sq_sum_bound_of_pointwise_posdiff_sq_sum_le}
\end{formalized}

\begin{theorem}[Source-sharp signed-Hadamard row-norm concentration]
Assume \(m>0\), \(U^TU=I_n\), and \(H_{ik}^2=1/m\).  For a fixed row
\[
  F_i(\omega)=\Norm{(H\operatorname{diag}(\varepsilon(\omega))U)_{i,*}},
\]
and every \(t>0\),
\[
  \mathbb P\{F_i(\omega)\le \mathbb E F_i + t\}
  \ge
  1-\exp\!\left(-\frac{m t^2}{8}\right).
\]
Consequently,
\[
  \mathbb P\left\{
    \forall i,\quad
    F_i(\omega)\le \sqrt{\frac nm}+t
  \right\}
  \ge
  1-m\exp\!\left(-\frac{m t^2}{8}\right),
\]
and
\[
  \mathbb P\left\{
    \forall i,\quad
    \ell_i(HD_\omega U)
      \le
      \frac{(\sqrt{n/m}+t)^2}{n}
  \right\}
  \ge
  1-m\exp\!\left(-\frac{m t^2}{8}\right),
\]
where \(\ell_i(V)=\Norm{V_{i,*}}^2/n\) is the equation~(6) leverage
probability for an orthonormal-column matrix.
\end{theorem}


\begin{formalized}
\lean{Algorithms/RandNLA/Preconditioning.lean}
\leanfit{rademacherTraceProbability_flip_tilt_sq_sum_bound_signedHadamard_row_vecNorm2}
\leanfit{rademacherTraceProbability_eventProb_vecNorm2_signedHadamard_le_mean_add_ge_one_sub_exp_m_t_sq_div_eight}
\leanfit{rademacherTraceProbability_eventProb_forall_vecNorm2_signedHadamard_le_sqrt_add_ge_one_sub_m_exp_m_t_sq_div_eight}
\leanfit{rademacherTraceProbability_eventProb_forall_rowNormSq_signedHadamard_le_sqrt_add_sq_ge_one_sub_m_exp_m_t_sq_div_eight}
\leanfit{rademacherTraceProbability_eventProb_forall_leverageScoreProb_signedHadamard_le_sqrt_add_sq_div_nat_ge_one_sub_m_exp_m_t_sq_div_eight}
\leanfit{rademacherTraceProbability_eventProb_forall_leverageScoreProb_signedHadamard_le_sqrt_add_sq_div_nat_ge_one_sub_delta}
\end{formalized}

\begin{corollary}[Logarithmic SRHT preprocessing radius]
Let \(0<\delta<m\) and choose
\[
  t_\delta=\sqrt{\frac{8\log(m/\delta)}{m}} .
\]
Then the preceding theorem gives
\[
  \mathbb P\left\{
    \forall i,\quad
    \Norm{(HD_\omega U)_{i,*}}^2
      \le
      \left(\sqrt{\frac nm}+t_\delta\right)^2
  \right\}
  \ge 1-\delta
\]
and
\[
  \mathbb P\left\{
    \forall i,\quad
    \ell_i(HD_\omega U)
      \le
      \frac{(\sqrt{n/m}+t_\delta)^2}{n}
  \right\}
  \ge 1-\delta .
\]
\end{corollary}


\begin{formalized}
\lean{Analysis/MatrixConcentration.lean}
\leanfit{real_sqrt_eight_log_div_pos_of_pos_lt}
\leanfit{real_mul_exp_neg_mul_sqrt_eight_log_div_sq_div_eight_eq}
\lean{Algorithms/RandNLA/Preconditioning.lean}
\leanfit{rademacherTraceProbability_eventProb_forall_rowNormSq_signedHadamard_le_log_delta_ge_one_sub_delta}
\leanfit{rademacherTraceProbability_eventProb_forall_leverageScoreProb_signedHadamard_le_log_delta_ge_one_sub_delta}
\end{formalized}

\begin{corollary}[Weak all-row Markov bound for flat signed-Hadamard preprocessing]
Assume the hypotheses of the preceding corollary, and let \(T>0\).  Then
\[
  \mathbb P\left\{
    \forall i,\quad
    \Norm{(HD_\omega U)_{i,*}}^2 \le T
  \right\}
  \ge
  1 - m\,\frac{n/m}{T}.
\]
Equivalently, if \(m((n/m)/T)\le\delta\), then
\[
  \mathbb P\left\{
    \forall i,\quad
    \Norm{(HD_\omega U)_{i,*}}^2 \le T
  \right\}
  \ge
  1-\delta .
\]
\end{corollary}


\begin{formalized}
\leanfit{rademacherTraceProbability_eventProb_rowNormSq_signedHadamard_le_ge_one_sub}
\leanfit{rademacherTraceProbability_eventProb_forall_rowNormSq_signedHadamard_le_ge_one_sub_sum}
\leanfit{rademacherTraceProbability_eventProb_forall_rowNormSq_signedHadamard_le_ge_one_sub}
\leanfit{rademacherTraceProbability_eventProb_forall_rowNormSq_signedHadamard_le_ge_one_sub_delta}
\end{formalized}

The Markov/union theorem and the coordinate-Hoeffding theorem remain in the
library as weaker alternatives.  The source-sharp SRHT row-norm and leverage
route is now also formalized in both explicit \(t\)-parameter and logarithmic
target-failure forms, composed with the exact iid uniform-row sampling theorem,
and transferred through the rounded sample-Gram computation with the fixed
row-norm-derived floating-point budget.

\begin{theorem}[One-sided floating-point preprocessing]
For \(\Pi_L\in\R^{r\times m}\), \(A\in\R^{m\times n}\), and valid
\(\gamma_m\),
\[
  \left|
    \fl(\Pi_LA)_{ij}-(\Pi_LA)_{ij}
  \right|
  \le
  \gamma_m\sum_k |(\Pi_L)_{ik}|\,|A_{kj}|.
\]
For \(A\in\R^{m\times n}\), \(\Pi_R\in\R^{n\times q}\), and valid
\(\gamma_n\),
\[
  \left|
    \fl(A\Pi_R)_{ij}-(A\Pi_R)_{ij}
  \right|
  \le
  \gamma_n\sum_k |A_{ik}|\,|(\Pi_R)_{kj}|.
\]
\end{theorem}

\begin{formalized}
\leanfit{fl_preconditionRows_error_bound}
\leanfit{fl_preconditionColumns_error_bound}
\leanfit{ComputedPreconditioner}
\leanfit{fl_preconditionRowsWithComputedLeft_total_error_bound}
\leanfit{fl_preconditionColumnsWithComputedRight_total_error_bound}
\end{formalized}

\begin{corollary}[Computed signed-Hadamard preconditioner]
Let \(H\in\mathbb R^{m\times m}\) be the exact scaled Hadamard/FHT factor and
let \(\sigma\in\mathbb R^m\) be the realized sign vector drawn from the exact
Rademacher law.  Put \(D=\operatorname{diag}(\sigma)\).  Suppose the implementation stores or
computes \(\widehat H\) and \(\widehat\sigma\) with certificates
\[
  |\widehat H_{ik}-H_{ik}|\le E^H_{ik},\qquad
  |\widehat\sigma_k-\sigma_k|\le E^\sigma_k,
  \qquad E^H_{ik},E^\sigma_k\ge0.
\]
Let \(\widehat D=\operatorname{diag}(\widehat\sigma)\), and form the realized computed
preconditioner by the rounded product
\[
  \widehat\Pi=\fl(\widehat H\,\widehat D).
\]
If \(\gamma_m\) is valid, then for every \(i,j\),
\[
\begin{aligned}
  |\widehat\Pi_{ij}-(HD)_{ij}|
  &\le
  \gamma_m\sum_k |\widehat H_{ik}|\,|\widehat D_{kj}|
  +\sum_k E^H_{ik}|\widehat D_{kj}|
  +\sum_k |H_{ik}|\,E^D_{kj}
  \;=:\;E^\Pi_{ij},
\end{aligned}
\]
where \(E^D_{kj}=E^\sigma_k\) when \(k=j\) and \(E^D_{kj}=0\) otherwise.
Consequently, for any \(A\in\mathbb R^{m\times n}\),
\[
  |\fl(\widehat\Pi A)_{ij}-(HDA)_{ij}|
  \le
  \gamma_m\sum_k |\widehat\Pi_{ik}|\,|A_{kj}|
  +\sum_k E^\Pi_{ik}|A_{kj}|.
\]
The probability law of \(\sigma\) is not perturbed: this result charges only
storage and arithmetic in the realized transform object.
\end{corollary}

\begin{formalized}
\leanfit{ComputedVector}
\leanfit{ComputedMatrix.diag}
\leanfit{fl_computedMatrixProduct_total_error_bound}
\leanfit{flComputedMatrixProductEntryErrorBudget}
\leanfit{ComputedMatrix.flProduct}
\leanfit{ComputedPreconditioner.flSignedHadamard}
\leanfit{fl_preconditionRowsWithComputedLeft_total_error_bound}
\end{formalized}

\begin{corollary}[Exact-factor signed-Hadamard preconditioner]
In the setting of the preceding corollary, suppose instead that \(H\) and the
realized sign vector \(\sigma\) are supplied exactly as mathematical factors.
The probability law of \(\sigma\) is still the exact Rademacher law.  The
implemented preconditioner is nevertheless formed by rounded arithmetic:
\[
  \widehat\Pi=\fl(HD),\qquad D=\operatorname{diag}(\sigma).
\]
If \(\gamma_m\) is valid, then for every \(i,j\),
\[
  |\widehat\Pi_{ij}-(HD)_{ij}|
  \le
  \gamma_m\sum_k |H_{ik}|\,|D_{kj}|
  =:E^{\Pi,\mathrm{fac}}_{ij}.
\]
Consequently, for any \(A\in\mathbb R^{m\times n}\),
\[
  |\fl(\widehat\Pi A)_{ij}-(HDA)_{ij}|
  \le
  \gamma_m\sum_k |\widehat\Pi_{ik}|\,|A_{kj}|
  +\sum_k E^{\Pi,\mathrm{fac}}_{ik}|A_{kj}|.
\]
Moreover, for the joint exact Rademacher/uniform-row law,
\[
  \mathbb P\!\left[
    \Frob{
      \operatorname{fl}(\widehat V_{\omega,S}^T\widehat V_{\omega,S})
      -G_s(HD_\omega U)
    }
    \le \tau^{\mathrm{fac}}_{\omega,S}
  \right]=1,
\]
where \(\widehat V_\omega=\fl(\widehat\Pi_\omega U)\) and
\(\tau^{\mathrm{fac}}_{\omega,S}\) is the existing computed-left budget with
\(E^{\Pi,\mathrm{fac}}\).  This closes the exact supplied-factor path; a
fast FHT recurrence/apply routine, or a sign-storage format other than the
modeled sign-copy certificates below, must still instantiate the nonzero-error
certificates above.
\end{corollary}

\begin{formalized}
\leanfit{ComputedPreconditioner.flSignedHadamardExactFactors}
\leanfit{ComputedPreconditioner.flSignedHadamardExactFactors_matrix}
\leanfit{ComputedPreconditioner.flSignedHadamardExactFactors_abs_error}
\leanfit{ComputedPreconditioner.flSignedHadamardExactFactors_entry_error_bound}
\leanfit{signedHadamardExactFactorPreconditioner}
\leanfit{signedHadamardUniformRowTraceProbability_eventProb_exactFactorComputedLeftPreconditioned_fl_uniformRowPerturbEvent_eq_one}
\end{formalized}

\begin{corollary}[Rounded scaled-Hadamard sign-pattern table]
Let \(S\in\mathbb R^{m\times m}\) be a supplied sign-pattern table satisfying
\[
  S_{ij}^2=1 \qquad\text{for all } i,j.
\]
Define the ideal scaled Hadamard-style table and the computed table by
\[
  H_{ij}=\sqrt{m^{-1}}\,S_{ij},
  \qquad
  \widehat H_{ij}=\operatorname{flsqrt}(m^{-1})\,S_{ij}.
\]
Here the sign pattern \(S\), the Rademacher law used later, and all sampling
laws are exact mathematical inputs; the only charged computation in this
corollary is the rounded square-root scale.  For every \(i,j\),
\[
  |\widehat H_{ij}-H_{ij}|
  \le
  \sqrt{m^{-1}}\,u\,|S_{ij}|
  =:E^H_{ij}.
\]
Moreover the ideal table \(H\) satisfies the deterministic SRHT flatness
hypothesis
\[
  H_{ij}^2=m^{-1}\qquad\text{for all }i,j,
\]
so it can be used as the exact analysis object in the SRHT row-norm
concentration theorem.  The certificate \(E^H\) can then be supplied to the
computed signed-Hadamard preconditioner corollary above.  This does not
formalize a concrete fast FHT recurrence/apply routine.  The next results handle the
modeled rounded-copy sign-storage paths \( \fl(\sigma_i\cdot 1) \),
\( \fl(\sigma_i+0) \), and \( \fl(\sigma_i-0) \); other sign-storage
representations remain ledger obligations if an implementation does not use
one of those modeled paths.
\end{corollary}

\begin{formalized}
\leanfit{ComputedMatrix.flSqrtInvNatScaledPattern}
\leanfit{ComputedMatrix.flSqrtInvNatScaledPattern_matrix}
\leanfit{ComputedMatrix.flSqrtInvNatScaledPattern_abs_error}
\leanfit{ComputedMatrix.flSqrtInvNatScaledPattern_entry_error_bound}
\leanfit{HadamardSignPattern}
\leanfit{hadamardFlat_sqrt_inv_nat_mul_signPattern}
\end{formalized}

\begin{corollary}[Scaled-pattern computed-left preconditioner event]
In the setting of the preceding corollary, let
\[
  H_{ij}=\sqrt{m^{-1}}\,S_{ij},\qquad
  \widehat H_{ij}=\operatorname{flsqrt}(m^{-1})\,S_{ij}.
\]
For a realized Rademacher sign vector \(\sigma_\omega\), define
\[
  D_\omega=\operatorname{diag}(\sigma_\omega),
  \qquad
  \widehat\Pi_\omega=\fl(\widehat H D_\omega),
  \qquad
  \widehat V_\omega=\fl(\widehat\Pi_\omega U).
\]
The laws of \(\omega\) and of the subsequent uniform row sample are exact
mathematical laws.  If the matrix-product \(\gamma_m\) and row-scaling
\(\gamma_s\) hypotheses are valid, then the computed-left perturbation event
holds with probability one:
\[
  \mathbb P\!\left[
    \Frob{
      \operatorname{fl}(\widehat V_{\omega,S}^T\widehat V_{\omega,S})
      -G_s(H D_\omega U)
    }
    \le \tau^{\mathrm{scale}}_{\omega,S}
  \right]=1.
\]
The radius \(\tau^{\mathrm{scale}}_{\omega,S}\) is the existing
computed-left budget instantiated with the rounded scale-table certificate for
\(\widehat H-H\), the rounded product certificate for
\(\widehat\Pi_\omega-HD_\omega\), the rounded formation of \(\widehat V\), the
uniform row-scaling arithmetic, and the final Gram dot products.  No
probability-construction or sampling-law error is introduced.
\end{corollary}

\begin{formalized}
\leanfit{ComputedPreconditioner.flSignedHadamardScaledPattern}
\leanfit{ComputedPreconditioner.flSignedHadamardScaledPattern_matrix}
\leanfit{ComputedPreconditioner.flSignedHadamardScaledPattern_abs_error}
\leanfit{ComputedPreconditioner.flSignedHadamardScaledPattern_entry_error_bound}
\leanfit{signedHadamardScaledPatternPreconditioner}
\leanfit{signedHadamardUniformRowTraceProbability_eventProb_scaledPatternComputedLeftPreconditioned_fl_uniformRowPerturbEvent_eq_one}
\end{formalized}

\begin{corollary}[Scaled-pattern event with rounded stored Rademacher signs]
In the setting of the preceding corollary, suppose the realized Rademacher
sign vector is stored by the rounded copy routine
\[
  \widehat\sigma_{\omega,i}=\fl(\sigma_{\omega,i}\cdot 1).
\]
Since the exact Rademacher law gives \(|\sigma_{\omega,i}|=1\), the stored
sign certificate satisfies
\[
  |\widehat\sigma_{\omega,i}-\sigma_{\omega,i}|
  \le u
  =:E^\sigma_{\omega,i}.
\]
Let
\[
  \widehat D_\omega=\operatorname{diag}(\widehat\sigma_\omega),
  \qquad
  \widehat\Pi_\omega=\fl(\widehat H\widehat D_\omega),
  \qquad
  \widehat V_\omega=\fl(\widehat\Pi_\omega U),
\]
where \(\widehat H_{ij}=\operatorname{flsqrt}(m^{-1})S_{ij}\).  If the
matrix-product \(\gamma_m\) and row-scaling \(\gamma_s\) hypotheses are valid,
then, under the exact Rademacher/uniform-row product law,
\[
  \mathbb P\!\left[
    \Frob{
      \operatorname{fl}(\widehat V_{\omega,S}^T\widehat V_{\omega,S})
      -G_s(H D_\omega U)
    }
    \le \tau^{\mathrm{store}}_{\omega,S}
  \right]=1.
\]
The radius \(\tau^{\mathrm{store}}_{\omega,S}\) is the existing computed-left
budget instantiated with the rounded scale-table certificate
\(\widehat H-H\), the stored-sign certificate
\(\widehat\sigma_\omega-\sigma_\omega\), the diagonal lift
\(\widehat D_\omega-D_\omega\), the rounded product certificate for
\(\widehat\Pi_\omega-H D_\omega\), the rounded formation of \(\widehat V\),
the uniform row-scaling arithmetic, and the final Gram dot products.  No
probability-construction or sampling-law error is introduced.
\end{corollary}

\begin{formalized}
\leanfit{ComputedVector.flMulOne}
\leanfit{ComputedVector.flStoredSign}
\leanfit{rademacherSignVector_abs}
\leanfit{ComputedPreconditioner.flSignedHadamardScaledPatternStoredSign}
\leanfit{ComputedPreconditioner.flSignedHadamardScaledPatternStoredSign_matrix}
\leanfit{ComputedPreconditioner.flSignedHadamardScaledPatternStoredSign_abs_error}
\leanfit{ComputedPreconditioner.flSignedHadamardScaledPatternStoredSign_entry_error_bound}
\leanfit{signedHadamardScaledPatternStoredSignPreconditioner}
\leanfit{signedHadamardUniformRowTraceProbability_eventProb_scaledPatternStoredSignComputedLeftPreconditioned_fl_uniformRowPerturbEvent_eq_one}
\end{formalized}

\begin{corollary}[Scaled-pattern event with add-zero stored Rademacher signs]
In the setting of the preceding corollary, suppose the realized Rademacher
sign vector is stored by the rounded add-zero copy routine
\[
  \widehat\sigma^{+}_{\omega,i}=\fl(\sigma_{\omega,i}+0).
\]
Since the exact Rademacher law gives \(|\sigma_{\omega,i}|=1\), the stored
sign certificate satisfies
\[
  |\widehat\sigma^{+}_{\omega,i}-\sigma_{\omega,i}|
  \le u
  =:E^{\sigma,+}_{\omega,i}.
\]
Let
\[
  \widehat D^{+}_\omega=\operatorname{diag}(\widehat\sigma^{+}_\omega),
  \qquad
  \widehat\Pi^{+}_\omega=\fl(\widehat H\widehat D^{+}_\omega),
  \qquad
  \widehat V^{+}_\omega=\fl(\widehat\Pi^{+}_\omega U),
\]
where \(\widehat H_{ij}=\operatorname{flsqrt}(m^{-1})S_{ij}\).  If the
matrix-product \(\gamma_m\) and row-scaling \(\gamma_s\) hypotheses are valid,
then, under the exact Rademacher/uniform-row product law,
\[
  \mathbb P\!\left[
    \Frob{
      \operatorname{fl}((\widehat V^{+}_{\omega,S})^T
        \widehat V^{+}_{\omega,S})
      -G_s(H D_\omega U)
    }
    \le \tau^{\mathrm{store},+}_{\omega,S}
  \right]=1.
\]
The radius \(\tau^{\mathrm{store},+}_{\omega,S}\) is the existing
computed-left budget instantiated with the rounded scale-table certificate
\(\widehat H-H\), the add-zero stored-sign certificate
\(\widehat\sigma^{+}_\omega-\sigma_\omega\), the diagonal lift
\(\widehat D^{+}_\omega-D_\omega\), the rounded product certificate for
\(\widehat\Pi^{+}_\omega-H D_\omega\), the rounded formation of
\(\widehat V^{+}\), the uniform row-scaling arithmetic, and the final Gram
dot products.  No probability-construction or sampling-law error is introduced.
\end{corollary}

\begin{formalized}
\leanfit{ComputedVector.flAddZeroRight}
\leanfit{ComputedVector.flStoredSignAddZeroRight}
\leanfit{rademacherSignVector_abs}
\leanfit{ComputedPreconditioner.flSignedHadamardScaledPatternStoredSignAddZeroRight}
\leanfit{ComputedPreconditioner.flSignedHadamardScaledPatternStoredSignAddZeroRight_matrix}
\leanfit{ComputedPreconditioner.flSignedHadamardScaledPatternStoredSignAddZeroRight_abs_error}
\leanfit{ComputedPreconditioner.flSignedHadamardScaledPatternStoredSignAddZeroRight_entry_error_bound}
\leanfit{signedHadamardScaledPatternStoredSignAddZeroRightPreconditioner}
\leanfit{signedHadamardUniformRowTraceProbability_eventProb_scaledPatternStoredSignAddZeroRightComputedLeftPreconditioned_fl_uniformRowPerturbEvent_eq_one}
\end{formalized}

\begin{corollary}[Scaled-pattern event with subtract-zero stored Rademacher signs]
In the setting of the scaled-pattern computed-left corollary, suppose the
realized Rademacher sign vector is stored by the rounded subtract-zero copy
routine
\[
  \widehat\sigma^{-}_{\omega,i}=\fl(\sigma_{\omega,i}-0).
\]
Since the exact Rademacher law gives \(|\sigma_{\omega,i}|=1\), the stored
sign certificate satisfies
\[
  |\widehat\sigma^{-}_{\omega,i}-\sigma_{\omega,i}|
  \le u
  =:E^{\sigma,-}_{\omega,i}.
\]
Let
\[
  \widehat D^{-}_\omega=\operatorname{diag}(\widehat\sigma^{-}_\omega),
  \qquad
  \widehat\Pi^{-}_\omega=\fl(\widehat H\widehat D^{-}_\omega),
  \qquad
  \widehat V^{-}_\omega=\fl(\widehat\Pi^{-}_\omega U),
\]
where \(\widehat H_{ij}=\operatorname{flsqrt}(m^{-1})S_{ij}\).  If the
matrix-product \(\gamma_m\) and row-scaling \(\gamma_s\) hypotheses are valid,
then, under the exact Rademacher/uniform-row product law,
\[
  \mathbb P\!\left[
    \Frob{
      \operatorname{fl}((\widehat V^{-}_{\omega,S})^T
        \widehat V^{-}_{\omega,S})
      -G_s(H D_\omega U)
    }
    \le \tau^{\mathrm{store},-}_{\omega,S}
  \right]=1.
\]
The radius \(\tau^{\mathrm{store},-}_{\omega,S}\) is the existing
computed-left budget instantiated with the rounded scale-table certificate
\(\widehat H-H\), the subtract-zero stored-sign certificate
\(\widehat\sigma^{-}_\omega-\sigma_\omega\), the diagonal lift
\(\widehat D^{-}_\omega-D_\omega\), the rounded product certificate for
\(\widehat\Pi^{-}_\omega-H D_\omega\), the rounded formation of
\(\widehat V^{-}\), the uniform row-scaling arithmetic, and the final Gram
dot products.  No probability-construction or sampling-law error is introduced.
\end{corollary}

\begin{formalized}
\leanfit{ComputedVector.flSubZeroRight}
\leanfit{ComputedVector.flStoredSignSubZeroRight}
\leanfit{rademacherSignVector_abs}
\leanfit{ComputedPreconditioner.flSignedHadamardScaledPatternStoredSignSubZeroRight}
\leanfit{ComputedPreconditioner.flSignedHadamardScaledPatternStoredSignSubZeroRight_matrix}
\leanfit{ComputedPreconditioner.flSignedHadamardScaledPatternStoredSignSubZeroRight_abs_error}
\leanfit{ComputedPreconditioner.flSignedHadamardScaledPatternStoredSignSubZeroRight_entry_error_bound}
\leanfit{signedHadamardScaledPatternStoredSignSubZeroRightPreconditioner}
\leanfit{signedHadamardUniformRowTraceProbability_eventProb_scaledPatternStoredSignSubZeroRightComputedLeftPreconditioned_fl_uniformRowPerturbEvent_eq_one}
\end{formalized}

\begin{corollary}[Generated Sylvester/Walsh scaled-pattern event]
Let \(p\in\mathbb N\) and \(m=2^p\).  For \(i,j\in\{0,\ldots,m-1\}\), define
the exact bit-parity predicate
\[
  b(i,j)=
  \operatorname{odd}\!\left(
    \sum_{\ell=0}^{p-1}
      \mathbf 1\{\operatorname{bit}_\ell(i)=1
      \text{ and } \operatorname{bit}_\ell(j)=1\}
  \right)
\]
and the generated Sylvester/Walsh sign pattern
\[
  S^{\mathrm{syl}}_{ij}=
  \begin{cases}
    -1, & b(i,j),\\
    1, & \neg b(i,j).
  \end{cases}
\]
The generation of \(S^{\mathrm{syl}}\) is exact integer/Boolean arithmetic, not
a sampling-probability computation.  Hence
\[
  |S^{\mathrm{syl}}_{ij}|=1,\qquad
  (S^{\mathrm{syl}}_{ij})^2=1,
\]
and the ideal scaled table
\[
  H^{\mathrm{syl}}_{ij}=\sqrt{m^{-1}}\,S^{\mathrm{syl}}_{ij}
\]
satisfies the flatness hypothesis \((H^{\mathrm{syl}}_{ij})^2=m^{-1}\).
The computed table is
\[
  \widehat H^{\mathrm{syl}}_{ij}
    =\operatorname{flsqrt}(m^{-1})\,S^{\mathrm{syl}}_{ij},
\]
and the rounded scale certificate is
\[
  |\widehat H^{\mathrm{syl}}_{ij}-H^{\mathrm{syl}}_{ij}|
  \le \sqrt{m^{-1}}\,u
  =:E^{H,\mathrm{syl}}_{ij}.
\]
For an exact Rademacher sign vector \(\sigma_\omega\), let
\[
  D_\omega=\operatorname{diag}(\sigma_\omega),\qquad
  \widehat\Pi^{\mathrm{syl}}_\omega
    =\fl(\widehat H^{\mathrm{syl}}D_\omega),\qquad
  \widehat V^{\mathrm{syl}}_\omega
    =\fl(\widehat\Pi^{\mathrm{syl}}_\omega U).
\]
If the matrix-product \(\gamma_m\) and row-scaling \(\gamma_s\) hypotheses are
valid, then under the exact Rademacher/uniform-row product law,
\[
  \mathbb P\!\left[
    \Frob{
      \operatorname{fl}((\widehat V^{\mathrm{syl}}_{\omega,S})^T
        \widehat V^{\mathrm{syl}}_{\omega,S})
      -G_s(H^{\mathrm{syl}}D_\omega U)
    }
    \le \tau^{\mathrm{syl}}_{\omega,S}
  \right]=1.
\]
The same probability-one statement holds for the three rounded sign-storage
variants \(\widehat\sigma_{\omega,i}=\fl(\sigma_{\omega,i}\cdot1)\),
\(\widehat\sigma^+_{\omega,i}=\fl(\sigma_{\omega,i}+0)\), and
\(\widehat\sigma^-_{\omega,i}=\fl(\sigma_{\omega,i}-0)\), with radii
\(\tau^{\mathrm{syl,store}}_{\omega,S}\),
\(\tau^{\mathrm{syl,store},+}_{\omega,S}\), and
\(\tau^{\mathrm{syl,store},-}_{\omega,S}\), respectively.  These radii charge
the rounded scale table, optional rounded sign storage, the diagonal lift, the
rounded \(H D_\omega\) product, rounded \(\widehat V\) formation, uniform
row-scaling arithmetic, and final Gram dot products.  They do not introduce
probability-construction, probability-normalization, or sampling-law transfer
terms.  A fast in-place FHT recurrence/apply implementation remains a separate
not-proved ledger obligation unless it materializes this generated table.
\end{corollary}

\begin{formalized}
\leanfit{sylvesterHadamardParity}
\leanfit{sylvesterHadamardSignPattern}
\leanfit{sylvesterHadamardSignPattern_abs}
\leanfit{sylvesterHadamardSignPattern_isSignPattern}
\leanfit{hadamardFlat_sqrt_inv_nat_mul_sylvesterSignPattern}
\leanfit{ComputedMatrix.flSqrtInvNatScaledSylvesterPattern}
\leanfit{ComputedMatrix.flSqrtInvNatScaledSylvesterPattern_entry_error_bound}
\leanfit{ComputedPreconditioner.flSignedHadamardSylvesterPattern}
\leanfit{ComputedPreconditioner.flSignedHadamardSylvesterPatternStoredSign}
\leanfit{ComputedPreconditioner.flSignedHadamardSylvesterPatternStoredSignAddZeroRight}
\leanfit{ComputedPreconditioner.flSignedHadamardSylvesterPatternStoredSignSubZeroRight}
\leanfit{signedHadamardSylvesterPatternPreconditioner}
\leanfit{signedHadamardSylvesterPatternStoredSignPreconditioner}
\leanfit{signedHadamardSylvesterPatternStoredSignAddZeroRightPreconditioner}
\leanfit{signedHadamardSylvesterPatternStoredSignSubZeroRightPreconditioner}
\leanfit{signedHadamardUniformRowTraceProbability_eventProb_sylvesterPatternComputedLeftPreconditioned_fl_uniformRowPerturbEvent_eq_one}
\leanfit{signedHadamardUniformRowTraceProbability_eventProb_sylvesterPatternStoredSignComputedLeftPreconditioned_fl_uniformRowPerturbEvent_eq_one}
\leanfit{signedHadamardUniformRowTraceProbability_eventProb_sylvesterPatternStoredSignAddZeroRightComputedLeftPreconditioned_fl_uniformRowPerturbEvent_eq_one}
\leanfit{signedHadamardUniformRowTraceProbability_eventProb_sylvesterPatternStoredSignSubZeroRightComputedLeftPreconditioned_fl_uniformRowPerturbEvent_eq_one}
\end{formalized}

\begin{corollary}[Sylvester/Walsh parity-count sign recurrence]
Let \(p\in\mathbb N\), and let \(S^{\mathrm{syl}}\) be the exact
Sylvester/Walsh sign table defined by the bit-parity count above.  For
\(i,i',j\in\{0,\ldots,2^p-1\}\) and a Boolean \(q\), suppose
\[
  \sum_{\ell=0}^{p-1}
    \mathbf 1\{\operatorname{bit}_\ell(i')=1
      \text{ and }\operatorname{bit}_\ell(j)=1\}
  =
  \sum_{\ell=0}^{p-1}
    \mathbf 1\{\operatorname{bit}_\ell(i)=1
      \text{ and }\operatorname{bit}_\ell(j)=1\}
  + \mathbf 1\{q\}.
\]
Then
\[
  S^{\mathrm{syl}}_{i'j}
  =
  \begin{cases}
    -S^{\mathrm{syl}}_{ij}, & q,\\
    \phantom{-}S^{\mathrm{syl}}_{ij}, & \neg q.
  \end{cases}
\]
In particular, equal parity counts give equal signs, while increasing the
parity count by one gives the negated sign.  This is the exact algebraic
bridge from the one-stage partner bit facts to the Sylvester/Walsh recurrence:
once the bit-toggle proof supplies the displayed count identity, the sign
recurrence follows immediately.  It is deterministic integer/Boolean
arithmetic, not a probability-law theorem and not a floating-point theorem.
The Rademacher and uniform-row sampling laws remain exact mathematical inputs.
\end{corollary}

\begin{formalized}
\leanfit{sylvesterHadamardSignPattern_eq_or_neg_of_parityWeight_eq_add_bool}
\leanfit{sylvesterHadamardSignPattern_eq_of_parityWeight_eq}
\leanfit{sylvesterHadamardSignPattern_neg_of_parityWeight_eq_add_one}
\end{formalized}

\begin{corollary}[Abstract stage-bit parity-count recurrence]
Let \(p\in\mathbb N\), let \(0\le r<p\), and let
\(S^{\mathrm{syl}}\) be the exact Sylvester/Walsh sign table.  Let
\(i,i',j\in\{0,\ldots,2^p-1\}\).  Suppose that
\[
  \operatorname{bit}_r(i)=0,\qquad
  \operatorname{bit}_r(i')=1,
\]
and that the two row indices agree in every non-stage bit:
\[
  \operatorname{bit}_\ell(i')=\operatorname{bit}_\ell(i)
  \qquad\text{for every }0\le \ell<p,\ \ell\ne r.
\]
Then the parity-count identity needed by the FHT recurrence is
\[
  \sum_{\ell=0}^{p-1}
    \mathbf 1\{\operatorname{bit}_\ell(i')=1
      \text{ and }\operatorname{bit}_\ell(j)=1\}
  =
  \sum_{\ell=0}^{p-1}
    \mathbf 1\{\operatorname{bit}_\ell(i)=1
      \text{ and }\operatorname{bit}_\ell(j)=1\}
  + \mathbf 1\{\operatorname{bit}_r(j)=1\}.
\]
Consequently,
\[
  S^{\mathrm{syl}}_{i'j}
  =
  \begin{cases}
    -S^{\mathrm{syl}}_{ij}, & \operatorname{bit}_r(j)=1,\\
    \phantom{-}S^{\mathrm{syl}}_{ij}, & \operatorname{bit}_r(j)=0.
  \end{cases}
\]
This closes the abstract finite-sum stage-bit part of the generated-FHT
realization proof.  The concrete generated partners \(i+2^r\) and \(i-2^r\)
are instantiated in the next corollary.  The result is exact deterministic
integer/Boolean arithmetic; it does not add floating-point error or
probability-law transfer.  The Rademacher and uniform row sampling laws remain
exact mathematical inputs.
\end{corollary}

\begin{formalized}
\leanfit{sylvesterHadamardParityWeight_partner_eq_add_stage_of_bits}
\leanfit{sylvesterHadamardSignPattern_partner_eq_or_neg_of_stage_bit}
\end{formalized}

\begin{corollary}[Generated stage bit-toggle parity recurrences]
Let \(p\in\mathbb N\), \(0\le r<p\), and let
\(i,j\in\{0,\ldots,2^p-1\}\).  The generic bit lemmas prove:
if \(\operatorname{bit}_r(a)=0\), then
\[
  \operatorname{bit}_\ell(a+2^r)=\operatorname{bit}_\ell(a)
  \qquad(\ell\ne r),
\]
and if \(\operatorname{bit}_r(a)=1\), then
\[
  \operatorname{bit}_\ell(a-2^r)=\operatorname{bit}_\ell(a)
  \qquad(\ell\ne r).
\]
Consequently the generated lower-half block test
\[
  i\bmod (2\cdot 2^r) < 2^r
\]
implies the concrete parity-count recurrence
\[
  \sum_{\ell=0}^{p-1}
    \mathbf 1\{\operatorname{bit}_\ell(i+2^r)=1
      \text{ and }\operatorname{bit}_\ell(j)=1\}
  =
  \sum_{\ell=0}^{p-1}
    \mathbf 1\{\operatorname{bit}_\ell(i)=1
      \text{ and }\operatorname{bit}_\ell(j)=1\}
  + \mathbf 1\{\operatorname{bit}_r(j)=1\},
\]
and therefore
\[
  S^{\mathrm{syl}}_{i+2^r,j}
  =
  \begin{cases}
    -S^{\mathrm{syl}}_{ij}, & \operatorname{bit}_r(j)=1,\\
    \phantom{-}S^{\mathrm{syl}}_{ij}, & \operatorname{bit}_r(j)=0.
  \end{cases}
\]
Similarly, the generated upper-half block test
\[
  2^r \le i\bmod (2\cdot 2^r)
\]
implies
\[
  \sum_{\ell=0}^{p-1}
    \mathbf 1\{\operatorname{bit}_\ell(i)=1
      \text{ and }\operatorname{bit}_\ell(j)=1\}
  =
  \sum_{\ell=0}^{p-1}
    \mathbf 1\{\operatorname{bit}_\ell(i-2^r)=1
      \text{ and }\operatorname{bit}_\ell(j)=1\}
  + \mathbf 1\{\operatorname{bit}_r(j)=1\},
\]
and therefore
\[
  S^{\mathrm{syl}}_{ij}
  =
  \begin{cases}
    -S^{\mathrm{syl}}_{i-2^r,j}, & \operatorname{bit}_r(j)=1,\\
    \phantom{-}S^{\mathrm{syl}}_{i-2^r,j}, & \operatorname{bit}_r(j)=0.
  \end{cases}
\]
This closes the concrete generated-partner bit arithmetic feeding the exact
FHT realization proof.  The remaining exact theorem is the one-stage
parity-table recurrence that combines these sign recurrences with the
lower/upper butterfly coordinate formulas, followed by the induction over
stages.  These are exact deterministic integer/FHT facts; no probability
construction, probability normalization, sampling-law transfer, or
floating-point arithmetic is introduced here.
\end{corollary}

\begin{formalized}
\leanfit{nat_testBit_add_two_pow_eq_of_testBit_eq_false_of_ne}
\leanfit{two_pow_le_of_testBit_eq_true}
\leanfit{nat_testBit_sub_two_pow_eq_of_testBit_eq_true_of_ne}
\leanfit{fhtSylvesterStage_upper_partner_testBit_eq_of_ne_of_mod_lt}
\leanfit{fhtSylvesterStage_lower_partner_testBit_eq_of_ne_of_mod_ge}
\leanfit{sylvesterHadamardParityWeight_upper_partner_eq_add_stage_of_mod_lt}
\leanfit{sylvesterHadamardSignPattern_upper_partner_eq_or_neg_of_mod_lt}
\leanfit{sylvesterHadamardParityWeight_upper_eq_lower_partner_add_stage_of_mod_ge}
\leanfit{sylvesterHadamardSignPattern_upper_eq_or_neg_lower_partner_of_mod_ge}
\end{formalized}

\begin{corollary}[One-butterfly FHT arithmetic certificate]
Consider one scalar pair update in a fast Hadamard transform stage.  The exact
butterfly sends
\[
  (a,b)\longmapsto
  (y_+,y_-)=(a+b,a-b),
\]
while the implemented floating-point butterfly is
\[
  (\widehat y_+,\widehat y_-)
  =
  \bigl(\operatorname{fladd}(a,b),\operatorname{flsub}(a,b)\bigr).
\]
Under the floating-point model with unit roundoff \(u\), the two computed
outputs satisfy
\[
  |\widehat y_+-y_+|
  \le
  u\,|y_+|,
  \qquad
  |\widehat y_--y_-|
  \le
  u\,|y_-|.
\]
This is an implementation-facing certificate for the non-probability
arithmetic inside one FHT butterfly: it charges the rounded addition and
rounded subtraction.  It does not introduce probability-construction,
probability-normalization, or sampling-law transfer terms; the Rademacher and
uniform row laws remain exact mathematical laws.  It is also not yet a full
fast-Hadamard apply theorem.  A complete FHT implementation must still compose
this primitive across a concrete stage schedule and account for scaling,
storage/overwrite effects, vector or matrix application, and the resulting
\texttt{ComputedMatrix} or \texttt{ComputedPreconditioner} certificate.
\end{corollary}

\begin{formalized}
\leanfit{fhtButterflyExact}
\leanfit{flFhtButterfly}
\leanfit{flFhtButterfly_fst}
\leanfit{flFhtButterfly_snd}
\leanfit{flFhtButterfly_add_error_bound}
\leanfit{flFhtButterfly_sub_error_bound}
\end{formalized}

\begin{corollary}[Vector-pair FHT update certificate]
Let \(x\in\mathbb R^n\), and let \(p,q\in\{0,\ldots,n-1\}\) be one pair of
coordinates selected by an FHT stage.  Define the exact pair update \(x'\) by
\[
  x'_i =
  \begin{cases}
    x_p+x_q, & i=p,\\
    x_p-x_q, & i=q\text{ and }i\ne p,\\
    x_i, & \text{otherwise},
  \end{cases}
\]
and define the computed pair update \(\widehat x'\) by
\[
  \widehat x'_i =
  \begin{cases}
    \operatorname{fladd}(x_p,x_q), & i=p,\\
    \operatorname{flsub}(x_p,x_q), & i=q\text{ and }i\ne p,\\
    x_i, & \text{otherwise}.
  \end{cases}
\]
Set the entrywise butterfly budget
\[
  B_i =
  \begin{cases}
    u\,|x_p+x_q|, & i=p,\\
    u\,|x_p-x_q|, & i=q\text{ and }i\ne p,\\
    0, & \text{otherwise}.
  \end{cases}
\]
Then \(B_i\ge0\) for every \(i\), and
\[
  |\widehat x'_i-x'_i|\le B_i
  \qquad\text{for every }i.
\]
This lifts the one-butterfly scalar arithmetic certificate to the vector shape
used by an in-place FHT stage.  The unchanged coordinates are treated as exact
reference entries; if an implementation copies or overwrites them through
rounded storage, that storage path needs its own certificate.  This result
still does not compose stages or scaling, and it introduces no probability-law
or probability-construction error term.
\end{corollary}

\begin{formalized}
\leanfit{fhtPairUpdateExact}
\leanfit{flFhtPairUpdate}
\leanfit{fhtPairUpdateErrorBudget}
\leanfit{fhtPairUpdateErrorBudget_nonneg}
\leanfit{flFhtPairUpdate_error_bound}
\end{formalized}

\begin{corollary}[Propagated vector-pair FHT update certificate]
Let \(x,\widehat x\in\mathbb R^n\), let \(p,q\in\{0,\ldots,n-1\}\), and
suppose the input to one FHT pair update already satisfies an entrywise
certificate
\[
  |\widehat x_i-x_i|\le E_i,\qquad E_i\ge0
  \qquad\text{for every }i .
\]
Let \(x'\) be the exact pair update applied to \(x\),
\[
  x'_i =
  \begin{cases}
    x_p+x_q, & i=p,\\
    x_p-x_q, & i=q\text{ and }i\ne p,\\
    x_i, & \text{otherwise},
  \end{cases}
\]
and let \(\widehat x'\) be the rounded pair update applied to the already
computed vector \(\widehat x\),
\[
  \widehat x'_i =
  \begin{cases}
    \operatorname{fladd}(\widehat x_p,\widehat x_q), & i=p,\\
    \operatorname{flsub}(\widehat x_p,\widehat x_q), & i=q\text{ and }i\ne p,\\
    \widehat x_i, & \text{otherwise}.
  \end{cases}
\]
Define the propagated one-pair budget
\[
  B^{\rm prop}_i =
  \begin{cases}
    u\,|\widehat x_p+\widehat x_q|+E_p+E_q, & i=p,\\
    u\,|\widehat x_p-\widehat x_q|+E_p+E_q, & i=q\text{ and }i\ne p,\\
    E_i, & \text{otherwise}.
  \end{cases}
\]
Then \(B^{\rm prop}_i\ge0\) and
\[
  |\widehat x'_i-x'_i|\le B^{\rm prop}_i
  \qquad\text{for every }i.
\]
The terms \(E_p+E_q\) propagate the already-present non-probability input
error from earlier FHT stages, while the \(u|\widehat x_p\pm\widehat x_q|\)
terms charge the new rounded addition/subtraction in this pair update.  The
unchanged coordinates retain their previous entrywise certificate.  This is
still a one-pair certificate, not a full fast-Hadamard recurrence: a complete
implementation-facing FHT theorem must compose these budgets over the concrete
stage schedule and add scale, copy, overwrite, vector-apply, and matrix-apply
certificates as needed.  Sampling probabilities, Rademacher signs as random
laws, and uniform row laws remain exact mathematical inputs.
\end{corollary}

\begin{formalized}
\leanfit{abs_add_pair_sub_add_pair_le}
\leanfit{abs_sub_pair_sub_sub_pair_le}
\leanfit{flFhtButterfly_add_error_bound_of_input_error}
\leanfit{flFhtButterfly_sub_error_bound_of_input_error}
\leanfit{fhtPairUpdatePropagatedErrorBudget}
\leanfit{fhtPairUpdatePropagatedErrorBudget_nonneg}
\leanfit{flFhtPairUpdate_propagated_error_bound}
\end{formalized}

\begin{corollary}[Ordered FHT pair-schedule propagation certificate]
Let
\[
  \mathcal P=((p_1,q_1),\ldots,(p_T,q_T))
\]
be any finite ordered list of FHT pair updates on \(\mathbb R^n\).  Starting
from exact and computed inputs \(x^{(0)}=x\), \(\widehat x^{(0)}=\widehat x\),
and an entrywise initial budget \(E^{(0)}=E\), define recursively
\[
  x^{(r+1)}=\operatorname{PairExact}_{p_{r+1},q_{r+1}}(x^{(r)}),
  \qquad
  \widehat x^{(r+1)}
    =\operatorname{PairFl}_{p_{r+1},q_{r+1}}(\widehat x^{(r)}),
\]
and set \(E^{(r+1)}=B^{\rm prop}(p_{r+1},q_{r+1},\widehat x^{(r)},E^{(r)})\),
where \(B^{\rm prop}\) is the one-pair budget from Corollary 203:
\[
  E^{(r+1)}_i =
  \begin{cases}
    u\,|\widehat x^{(r)}_{p_{r+1}}+\widehat x^{(r)}_{q_{r+1}}|
      +E^{(r)}_{p_{r+1}}+E^{(r)}_{q_{r+1}}, & i=p_{r+1},\\
    u\,|\widehat x^{(r)}_{p_{r+1}}-\widehat x^{(r)}_{q_{r+1}}|
      +E^{(r)}_{p_{r+1}}+E^{(r)}_{q_{r+1}},
      & i=q_{r+1}\text{ and }i\ne p_{r+1},\\
    E^{(r)}_i, & \text{otherwise}.
  \end{cases}
\]
If the initial budget satisfies \(E_i\ge0\) and
\[
  |\widehat x_i-x_i|\le E_i
  \qquad\text{for every }i,
\]
then every recursively propagated budget is nonnegative and the final exact
and rounded schedule outputs satisfy
\[
  |\widehat x^{(T)}_i-x^{(T)}_i|\le E^{(T)}_i
  \qquad\text{for every }i.
\]
This composes the visible rounded-addition/subtraction and previously
propagated non-probability error across any supplied ordered pair schedule.
It still does not add the FHT normalization scale, prove a concrete Hadamard
stage-list generator, model storage/copy paths beyond the functional update,
or lift the vector schedule to a matrix/vector apply certificate.  The
sampling laws remain exact mathematical inputs.
\end{corollary}

\begin{formalized}
\leanfit{fhtPairScheduleExact}
\leanfit{flFhtPairSchedule}
\leanfit{fhtPairSchedulePropagatedErrorBudget}
\leanfit{fhtPairSchedulePropagatedErrorBudget_nonneg}
\leanfit{flFhtPairSchedule_propagated_error_bound}
\end{formalized}

\begin{corollary}[Scaled ordered FHT pair-schedule certificate]
Let \(\mathcal P\) be any finite ordered list of FHT pair updates on
\(\mathbb R^n\).  Let
\[
  y_i=\operatorname{FHTSchedule}_{\mathcal P}(x)_i,\qquad
  \widehat y_i=\operatorname{flFHTSchedule}_{\mathcal P}(\widehat x)_i,
\]
and let \(E^{\mathcal P}_i\) be the propagated schedule budget from
Corollary 204.  Assume
\[
  E_i\ge0,\qquad |\widehat x_i-x_i|\le E_i
  \quad\text{for every }i,
\]
and suppose the FHT normalization scale is represented by an exact analysis
value \(c\), a computed value \(\widehat c\), and a certificate
\[
  |\widehat c-c|\le\eta,\qquad \eta\ge0.
\]
Define the exact scaled schedule output and the rounded implementation output
by
\[
  z_i=c\,y_i,\qquad
  \widehat z_i=\operatorname{flmul}(\widehat c,\widehat y_i).
\]
For each coordinate set
\[
  B^{\rm scale}_i
   =
   u\,|\widehat c\,\widehat y_i|
   +|\widehat c|\,E^{\mathcal P}_i
   +\eta\bigl(|\widehat y_i|+E^{\mathcal P}_i\bigr).
\]
Then \(B^{\rm scale}_i\ge0\) and
\[
  |\widehat z_i-z_i|\le B^{\rm scale}_i
  \qquad\text{for every }i.
\]
Thus the generic scaled FHT schedule charges the final rounded multiplication,
the already-propagated schedule error, and the computed-scale radius.  This is
an implementation-facing non-probability arithmetic certificate: Rademacher
and uniform row sampling laws remain exact mathematical inputs, and no
probability-construction or law-transfer term is introduced.  A full fast FHT
implementation still needs a concrete stage-list generator, a concrete
normalization-scale routine when \(\widehat c\) is computed by code,
storage/copy/overwrite certificates for the chosen layout, vector/matrix apply
wrappers, and a final \texttt{ComputedMatrix}/\texttt{ComputedPreconditioner}
instantiation.
\end{corollary}

\begin{formalized}
\leanfit{fhtScaledPairScheduleExact}
\leanfit{flFhtScaledPairSchedule}
\leanfit{fhtScaledPairScheduleErrorBudget}
\leanfit{fhtScaledPairScheduleErrorBudget_nonneg}
\leanfit{flFhtScaledPairSchedule_error_bound}
\end{formalized}

\begin{corollary}[Columnwise computed-matrix FHT schedule certificate]
Let \(A\in\mathbb R^{m\times n}\) be the exact analysis input matrix and let
\(\widehat A\) be a computed input matrix with entrywise certificate
\[
  |\widehat A_{ij}-A_{ij}|\le E^A_{ij},\qquad E^A_{ij}\ge0.
\]
Let \(\mathcal P\) be a supplied ordered list of pair updates on the row index
\(\{1,\ldots,m\}\), and let the scale certificate satisfy
\[
  |\widehat c-c|\le\eta,\qquad \eta\ge0.
\]
For each column \(j\), apply the exact scaled schedule to \(A_{\ast j}\) and
the rounded scaled schedule to \(\widehat A_{\ast j}\):
\[
  Y_{ij}
    =\operatorname{FHTScheduleMatrix}_{\mathcal P,c}(A)_{ij},
  \qquad
  \widehat Y_{ij}
    =\operatorname{flFHTScheduleMatrix}_{\mathcal P,\widehat c}(\widehat A)_{ij}.
\]
Let \(B^{\rm mat}_{ij}\) be the columnwise budget obtained from Corollary 205,
namely
\[
  B^{\rm mat}_{ij}
   =
   u\,|\widehat c\,\widehat y_{ij}|
   +|\widehat c|\,E^{\mathcal P}_{ij}
   +\eta\bigl(|\widehat y_{ij}|+E^{\mathcal P}_{ij}\bigr),
\]
where \(\widehat y_{ij}\) is the rounded unscaled schedule output in column
\(j\) and \(E^{\mathcal P}_{ij}\) is the propagated unscaled schedule budget
for that column.  Then \(B^{\rm mat}_{ij}\ge0\) and
\[
  |\widehat Y_{ij}-Y_{ij}|\le B^{\rm mat}_{ij}
  \qquad\text{for every }i,j.
\]
Consequently the rounded columnwise scheduled FHT output is a
\texttt{ComputedMatrix} certificate for the exact supplied-schedule transform.
This closes the vector-to-matrix lift for a supplied pair list.  It still does
not prove that the pair list is the Sylvester/Hadamard stage generator, nor
does the generic adapter itself choose a normalization-scale routine,
layout-specific storage/copy/overwrite model, or final chosen-transform
\texttt{ComputedPreconditioner}.  The following corollary instantiates the
rounded square-root normalization scale.
\end{corollary}

\begin{formalized}
\leanfit{fhtScaledPairScheduleMatrixExact}
\leanfit{flFhtScaledPairScheduleMatrix}
\leanfit{fhtScaledPairScheduleMatrixErrorBudget}
\leanfit{fhtScaledPairScheduleMatrixErrorBudget_nonneg}
\leanfit{flFhtScaledPairScheduleMatrix_error_bound}
\leanfit{ComputedMatrix.flScaledFhtPairScheduleColumns}
\leanfit{ComputedMatrix.flScaledFhtPairScheduleColumns_entry_error_bound}
\end{formalized}

\begin{corollary}[Rounded square-root FHT scale certificate]
Let \(m\) be the row dimension and let the exact FHT normalization scale,
computed scale, and scale-error radius be
\[
  c_m=\sqrt{m^{-1}},\qquad
  \widehat c_m=\operatorname{flsqrt}(m^{-1}),\qquad
  \eta_m=c_m u .
\]
Then
\[
  \eta_m\ge0,
  \qquad
  |\widehat c_m-c_m|\le\eta_m .
\]
Consequently the generic scaled-schedule certificates above may be instantiated
with \(c=c_m\), \(\widehat c=\widehat c_m\), and \(\eta=\eta_m\).

For a vector input, suppose
\[
  |\widehat x_i-x_i|\le E_i,\qquad E_i\ge0,
\]
and let \(z_i=\operatorname{FHTSchedule}_{\mathcal P,c_m}(x)_i\) and
\(\widehat z_i=
\operatorname{flFHTSchedule}_{\mathcal P,\widehat c_m}(\widehat x)_i\).
If \(\widehat y_i\) is the rounded unscaled schedule output and
\(E_i^{\mathcal P}\) is the propagated unscaled schedule budget, then
\[
  |\widehat z_i-z_i|
  \le
  u\,|\widehat c_m\widehat y_i|
  +|\widehat c_m|\,E_i^{\mathcal P}
  +\eta_m\bigl(|\widehat y_i|+E_i^{\mathcal P}\bigr).
\]
For a matrix input, suppose
\[
  |\widehat A_{ij}-A_{ij}|\le E^A_{ij},\qquad E^A_{ij}\ge0,
\]
and apply the same supplied pair schedule columnwise.  Then
\[
  \left|
    \operatorname{flFHTScheduleMatrix}_{\mathcal P,\widehat c_m}
      (\widehat A)_{ij}
    -
    \operatorname{FHTScheduleMatrix}_{\mathcal P,c_m}(A)_{ij}
  \right|
  \le
  u\,|\widehat c_m\widehat y_{ij}|
  +|\widehat c_m|\,E_{ij}^{\mathcal P}
  +\eta_m\bigl(|\widehat y_{ij}|+E_{ij}^{\mathcal P}\bigr).
\]
The columnwise constructor packages this as a \texttt{ComputedMatrix}
certificate.  This is a fully computed-object result for the square-root FHT
normalization scale in a supplied-schedule implementation: the exact pair
schedule and exact input are the analysis references, while the rounded pair
arithmetic, rounded final scale multiplication, input perturbation, and rounded
\(\operatorname{flsqrt}(m^{-1})\) scale are charged.  Rademacher signs,
uniform row laws, and all sampling probabilities remain exact mathematical
inputs.
\end{corollary}

\begin{formalized}
\leanfit{fhtSqrtInvNatScale}
\leanfit{flFhtSqrtInvNatScale}
\leanfit{fhtSqrtInvNatScaleErrorRadius}
\leanfit{fhtSqrtInvNatScaleErrorRadius_nonneg}
\leanfit{flFhtSqrtInvNatScale_error_bound}
\leanfit{flFhtScaledPairSchedule_sqrtInvNatScale_error_bound}
\leanfit{flFhtScaledPairScheduleMatrix_sqrtInvNatScale_error_bound}
\leanfit{ComputedMatrix.flScaledFhtPairScheduleColumnsSqrtInvNat}
\leanfit{ComputedMatrix.flScaledFhtPairScheduleColumnsSqrtInvNat_entry_error_bound}
\end{formalized}

\begin{corollary}[Concrete one-stage FHT pair-list generator]
Let \(n\) be a finite row dimension and let \(s\) be a stage stride.  For
indices \(a,b\in\{0,\ldots,n-1\}\), define the one-stage FHT pair predicate
\[
  G_s(a,b)
  \quad\Longleftrightarrow\quad
  0<s,\qquad b=a+s,\qquad a\bmod (2s)<s .
\]
Let \(\operatorname{Pairs}(n,s)\) be the finite list of all pairs satisfying
this predicate.  Then, for every pair \((a,b)\),
\[
  (a,b)\in\operatorname{Pairs}(n,s)
  \quad\Longleftrightarrow\quad
  G_s(a,b).
\]
Consequently, any generated pair satisfies
\[
  0<s,\qquad b=a+s,\qquad a\bmod(2s)<s,\qquad a<b,
  \qquad a\ne b .
\]
The one-stage list also has the first partner-uniqueness and lower/upper-half
separation facts needed for the later exact stage formula.  If
\((a,b),(a',b')\in\operatorname{Pairs}(n,s)\), then
\[
  a=a'\Longrightarrow b=b',
  \qquad
  b=b'\Longrightarrow a=a',
  \qquad
  a\ne b',
  \qquad
  b\ne a' .
\]
Equivalently, equal first or equal second coordinates identify the generated
pair itself:
\[
  a=a'\Longrightarrow (a,b)=(a',b'),
  \qquad
  b=b'\Longrightarrow (a,b)=(a',b') .
\]
The modular arithmetic used for the separation is explicit: whenever
\[
  0<s,\qquad a\bmod(2s)<s,
\]
one has
\[
  s\le (a+s)\bmod(2s).
\]
For the Sylvester/Walsh dimension \(n=2^p\) and stage \(r\), the generated
stage list uses stride \(2^r\).  In this case
\[
  (a,b)\in\operatorname{SylvesterPairs}(p,r)
  \quad\Longleftrightarrow\quad
  b=a+2^r,\qquad a\bmod(2\cdot2^r)<2^r .
\]
This closes the exact finite-index list-generation layer for one FHT stage.
It is not a probability law and it does not charge probability construction.
The later partial-transform recurrence and range induction close the
transform-correctness proof that the composed generated stages realize the
desired Sylvester/Hadamard transform.  What remains after that exact result is
to combine the rounded arithmetic with storage/layout and final
\texttt{ComputedPreconditioner} certificates for a concrete implementation.
\end{corollary}

\begin{formalized}
\leanfit{fhtStagePairPredicate}
\leanfit{fhtStagePairs}
\leanfit{mem_fhtStagePairs_iff}
\leanfit{fhtStagePairs_stride_pos}
\leanfit{fhtStagePairs_second_val_eq}
\leanfit{fhtStagePairs_first_val_mod_lt}
\leanfit{fhtStagePairs_first_val_lt_second_val}
\leanfit{fhtStagePairs_fst_ne_snd}
\leanfit{fhtStagePairs_snd_eq_of_fst_eq}
\leanfit{fhtStagePairs_fst_eq_of_snd_eq}
\leanfit{nat_stride_add_mod_two_stride_ge_of_mod_lt}
\leanfit{fhtStagePairs_fst_ne_snd_of_mem_mem}
\leanfit{fhtStagePairs_snd_ne_fst_of_mem_mem}
\leanfit{fhtStagePairs_eq_of_fst_eq}
\leanfit{fhtStagePairs_eq_of_snd_eq}
\leanfit{fhtSylvesterStagePairs}
\leanfit{mem_fhtSylvesterStagePairs_iff}
\end{formalized}

\begin{corollary}[Generated one-stage FHT pair updates commute]
Let \(T_{a,b}\) denote the exact vector butterfly update that replaces
\((x_a,x_b)\) by \((x_a+x_b,x_a-x_b)\) and leaves all other coordinates
unchanged.  If
\[
  (a,b),(a',b')\in\operatorname{Pairs}(n,s),
  \qquad
  (a,b)\ne(a',b'),
\]
then the generated pairs have disjoint coordinates:
\[
  a\ne a',
  \qquad a\ne b',
  \qquad b\ne a',
  \qquad b\ne b' .
\]
Consequently, for every vector \(x\),
\[
  T_{a,b}\bigl(T_{a',b'}x\bigr)
  =
  T_{a',b'}\bigl(T_{a,b}x\bigr).
\]
The same disjoint-coordinate conclusion is now formalized for the rounded
functional update \(\widehat T_{a,b}\) and for the three modified-coordinate
writeback updates that store only the two butterfly outputs through
\(\operatorname{fl\_add}(y,0)\), \(\operatorname{fl\_mul}(y,1)\), or
\(\operatorname{fl\_sub}(y,0)\).  Thus, for the modeled
functional schedules, same-stage generated-pair order does not introduce an
extra aliasing or memory-order hypothesis beyond the already charged rounded
butterfly and writeback arithmetic.  Concrete Sylvester/Walsh wrappers
specialize the same no-alias/list-order facts to
\(\operatorname{Pairs}_{\rm syl}(p,r)\), so downstream generated-FHT proofs can
use the \(2^p\)-dimensional stage interface directly without unfolding it to a
generic stride schedule.
Lean also proves the companion no-touch property for the rounded base schedule
and all three modified-coordinate writeback schedules: if a coordinate is not
the first or second entry of any pair in the supplied schedule list, then that
coordinate is returned exactly unchanged.  The generated-stage wrappers apply
the same statement to \(\operatorname{Pairs}(n,s)\).  This preservation result
is not asserted for all-coordinate copy/writeback routines, which deliberately
rewrite every entry.
The same no-touch statement is also formalized for the propagated error
budgets: under the same hypothesis, the base rounded budget and the three
modified-coordinate budgets return exactly the incoming radius \(E_i\).  Thus
the displayed modified-coordinate recurrences do not hide any copy contribution
on coordinates outside the active pair list.
The concrete Sylvester/Walsh wrappers specialize these value and budget facts
to the full generated list \(\operatorname{Pairs}_{\rm syl}(p)\) and to a
single generated stage \(\operatorname{Pairs}_{\rm syl}(p,r)\).  Therefore a
downstream proof can use the concrete generated-FHT interfaces directly: if
coordinate \(i\) is outside every pair in the relevant full schedule or stage,
then the rounded base schedule and the three modified-coordinate schedules
return \(\widehat x_i\), and the corresponding propagated budgets return
exactly \(E_i\).  These wrappers add no probability-law change and no extra
floating-point term; all-coordinate copy routines remain outside this result.
More generally, let \(L\) be an exact schedule list such that every
\((c,d)\in L\) belongs to \(\operatorname{Pairs}(n,s)\) and
\((a,b)\ne(c,d)\).  If \(S_L\) denotes the exact schedule obtained by applying
the updates in \(L\), then
\[
  S_L\bigl(T_{a,b}x\bigr)
  =
  T_{a,b}\bigl(S_Lx\bigr).
\]
This is the exact list-order bridge needed before proving that one generated
stage realizes the simultaneous lower/upper-half FHT butterfly map.  It does
not introduce any probability or sampling-law approximation.
\end{corollary}

\begin{formalized}
\leanfit{fhtPairUpdateExact_commute_of_disjoint}
\leanfit{flFhtPairUpdate_commute_of_disjoint}
\leanfit{fhtStagePairs_disjoint_of_ne}
\leanfit{fhtStagePairs_pairUpdateExact_commute_of_ne}
\leanfit{fhtPairScheduleExact_commute_update_of_forall}
\leanfit{fhtStagePairs_pairUpdateExact_commute_schedule_of_ne}
\leanfit{flFhtPairSchedule_commute_update_of_forall}
\leanfit{fhtStagePairs_flFhtPairUpdate_commute_of_ne}
\leanfit{fhtStagePairs_flFhtPairUpdate_commute_schedule_of_ne}
\leanfit{flFhtPairUpdateModifiedStoredAddZeroRight_commute_of_disjoint}
\leanfit{flFhtPairUpdateModifiedStoredMulOne_commute_of_disjoint}
\leanfit{flFhtPairUpdateModifiedStoredSubZeroRight_commute_of_disjoint}
\leanfit{flFhtPairScheduleModifiedStoredAddZeroRight_commute_update_of_forall}
\leanfit{flFhtPairScheduleModifiedStoredMulOne_commute_update_of_forall}
\leanfit{flFhtPairScheduleModifiedStoredSubZeroRight_commute_update_of_forall}
\leanfit{fhtStagePairs_flFhtPairUpdateModifiedStoredAddZeroRight_commute_of_ne}
\leanfit{fhtStagePairs_flFhtPairUpdateModifiedStoredAddZeroRight_commute_schedule_of_ne}
\leanfit{fhtStagePairs_flFhtPairUpdateModifiedStoredMulOne_commute_of_ne}
\leanfit{fhtStagePairs_flFhtPairUpdateModifiedStoredMulOne_commute_schedule_of_ne}
\leanfit{fhtStagePairs_flFhtPairUpdateModifiedStoredSubZeroRight_commute_of_ne}
\leanfit{fhtStagePairs_flFhtPairUpdateModifiedStoredSubZeroRight_commute_schedule_of_ne}
\leanfit{fhtSylvesterStagePairs_disjoint_of_ne}
\leanfit{fhtSylvesterStagePairs_pairUpdateExact_commute_of_ne}
\leanfit{fhtSylvesterStagePairs_pairUpdateExact_commute_schedule_of_ne}
\leanfit{fhtSylvesterStagePairs_flFhtPairUpdate_commute_of_ne}
\leanfit{fhtSylvesterStagePairs_flFhtPairUpdate_commute_schedule_of_ne}
\leanfit{fhtSylvesterStagePairs_flFhtPairUpdateModifiedStoredAddZeroRight_commute_of_ne}
\leanfit{fhtSylvesterStagePairs_flFhtPairUpdateModifiedStoredAddZeroRight_commute_schedule_of_ne}
\leanfit{fhtSylvesterStagePairs_flFhtPairUpdateModifiedStoredMulOne_commute_of_ne}
\leanfit{fhtSylvesterStagePairs_flFhtPairUpdateModifiedStoredMulOne_commute_schedule_of_ne}
\leanfit{fhtSylvesterStagePairs_flFhtPairUpdateModifiedStoredSubZeroRight_commute_of_ne}
\leanfit{fhtSylvesterStagePairs_flFhtPairUpdateModifiedStoredSubZeroRight_commute_schedule_of_ne}
\leanfit{flFhtPairSchedule_apply_of_forall_not_mem}
\leanfit{flFhtPairScheduleModifiedStoredAddZeroRight_apply_of_forall_not_mem}
\leanfit{flFhtPairScheduleModifiedStoredMulOne_apply_of_forall_not_mem}
\leanfit{flFhtPairScheduleModifiedStoredSubZeroRight_apply_of_forall_not_mem}
\leanfit{fhtStagePairs_flFhtPairSchedule_apply_of_forall_not_mem}
\leanfit{fhtStagePairs_flFhtPairScheduleModifiedStoredAddZeroRight_apply_of_forall_not_mem}
\leanfit{fhtStagePairs_flFhtPairScheduleModifiedStoredMulOne_apply_of_forall_not_mem}
\leanfit{fhtStagePairs_flFhtPairScheduleModifiedStoredSubZeroRight_apply_of_forall_not_mem}
\leanfit{fhtPairSchedulePropagatedErrorBudget_apply_of_forall_not_mem}
\leanfit{fhtPairScheduleModifiedStoredAddZeroRightPropagatedErrorBudget_apply_of_forall_not_mem}
\leanfit{fhtPairScheduleModifiedStoredMulOnePropagatedErrorBudget_apply_of_forall_not_mem}
\leanfit{fhtPairScheduleModifiedStoredSubZeroRightPropagatedErrorBudget_apply_of_forall_not_mem}
\leanfit{fhtStagePairs_fhtPairSchedulePropagatedErrorBudget_apply_of_forall_not_mem}
\leanfit{fhtStagePairs_fhtPairScheduleModifiedStoredAddZeroRightPropagatedErrorBudget_apply_of_forall_not_mem}
\leanfit{fhtStagePairs_fhtPairScheduleModifiedStoredMulOnePropagatedErrorBudget_apply_of_forall_not_mem}
\leanfit{fhtStagePairs_fhtPairScheduleModifiedStoredSubZeroRightPropagatedErrorBudget_apply_of_forall_not_mem}
\leanfit{flFhtSylvesterSchedule_apply_of_forall_not_mem}
\leanfit{flFhtSylvesterScheduleModifiedStoredAddZeroRight_apply_of_forall_not_mem}
\leanfit{flFhtSylvesterScheduleModifiedStoredMulOne_apply_of_forall_not_mem}
\leanfit{flFhtSylvesterScheduleModifiedStoredSubZeroRight_apply_of_forall_not_mem}
\leanfit{fhtSylvesterSchedulePropagatedErrorBudget_apply_of_forall_not_mem}
\leanfit{fhtSylvesterScheduleModifiedStoredAddZeroRightPropagatedErrorBudget_apply_of_forall_not_mem}
\leanfit{fhtSylvesterScheduleModifiedStoredMulOnePropagatedErrorBudget_apply_of_forall_not_mem}
\leanfit{fhtSylvesterScheduleModifiedStoredSubZeroRightPropagatedErrorBudget_apply_of_forall_not_mem}
\leanfit{flFhtSylvesterStageScheduleModifiedStoredAddZeroRight}
\leanfit{flFhtSylvesterStageScheduleModifiedStoredMulOne}
\leanfit{flFhtSylvesterStageScheduleModifiedStoredSubZeroRight}
\leanfit{fhtSylvesterStageScheduleModifiedStoredAddZeroRightPropagatedErrorBudget}
\leanfit{fhtSylvesterStageScheduleModifiedStoredMulOnePropagatedErrorBudget}
\leanfit{fhtSylvesterStageScheduleModifiedStoredSubZeroRightPropagatedErrorBudget}
\leanfit{flFhtSylvesterStageSchedule_apply_of_forall_not_mem}
\leanfit{flFhtSylvesterStageScheduleModifiedStoredAddZeroRight_apply_of_forall_not_mem}
\leanfit{flFhtSylvesterStageScheduleModifiedStoredMulOne_apply_of_forall_not_mem}
\leanfit{flFhtSylvesterStageScheduleModifiedStoredSubZeroRight_apply_of_forall_not_mem}
\leanfit{fhtSylvesterStageSchedulePropagatedErrorBudget_apply_of_forall_not_mem}
\leanfit{fhtSylvesterStageScheduleModifiedStoredAddZeroRightPropagatedErrorBudget_apply_of_forall_not_mem}
\leanfit{fhtSylvesterStageScheduleModifiedStoredMulOnePropagatedErrorBudget_apply_of_forall_not_mem}
\leanfit{fhtSylvesterStageScheduleModifiedStoredSubZeroRightPropagatedErrorBudget_apply_of_forall_not_mem}
\end{formalized}

\begin{corollary}[Exact output of one generated FHT stage]
Let \(\operatorname{Pairs}(n,s)\) be the generated one-stage pair list above,
and let \(S_{n,s}\) be the exact ordered schedule obtained by applying all
butterfly updates in that list.  For every generated pair
\((a,b)\in\operatorname{Pairs}(n,s)\) and every vector \(x\),
\[
  S_{n,s}(x)_a=x_a+x_b,\qquad
  S_{n,s}(x)_b=x_a-x_b .
\]
Equivalently, in the Sylvester/Walsh specialization, if
\((a,b)\in\operatorname{SylvesterPairs}(p,r)\), then the one-stage schedule
\(S_{p,r}\) satisfies
\[
  S_{p,r}(x)_a=x_a+x_b,\qquad
  S_{p,r}(x)_b=x_a-x_b .
\]
The proof uses only exact finite-list arithmetic: the generated stage list has
no duplicate pairs, and every other pair in the same stage is disjoint from
\((a,b)\), hence later updates leave the target coordinates unchanged.  This is
an exact-object theorem, not a probability-law theorem, and it introduces no
probability-construction or sampling-law approximation.
\end{corollary}

\begin{formalized}
\leanfit{fhtStagePairs_nodup}
\leanfit{fhtPairScheduleExact_apply_of_forall_not_mem}
\leanfit{fhtPairScheduleExact_apply_pair_of_mem_stage_list}
\leanfit{fhtStagePairs_pairScheduleExact_apply_pair_of_mem}
\leanfit{fhtStagePairs_pairScheduleExact_apply_fst_of_mem}
\leanfit{fhtStagePairs_pairScheduleExact_apply_snd_of_mem}
\leanfit{fhtSylvesterStageScheduleExact_apply_pair_of_mem}
\leanfit{fhtSylvesterStageScheduleExact_apply_fst_of_mem}
\leanfit{fhtSylvesterStageScheduleExact_apply_snd_of_mem}
\end{formalized}

\begin{corollary}[Coordinate formulas for one generated Sylvester FHT stage]
Let \(p,r\in\mathbb N\), let \(S_{p,r}\) be the exact generated stage with
stride \(2^r\) in dimension \(2^p\), and let \(x\in\mathbb R^{2^p}\).
If \(i\) is a lower-half coordinate in its stride block, meaning
\[
  i+2^r<2^p,\qquad
  i\bmod(2\cdot 2^r)<2^r,
\]
then the generated stage satisfies the explicit coordinate formula
\[
  S_{p,r}(x)_i=x_i+x_{i+2^r}.
\]
If \(i\) is treated as an upper-half coordinate with lower partner
\(\ell=i-2^r\), meaning
\[
  2^r\le i,\qquad
  (i-2^r)\bmod(2\cdot 2^r)<2^r,
\]
then
\[
  S_{p,r}(x)_i=x_{i-2^r}-x_i .
\]
Equivalently, the exact finite-index generator contains the constructed pairs
\[
  (i,i+2^r)\in\operatorname{SylvesterPairs}(p,r),
  \qquad
  (i-2^r,i)\in\operatorname{SylvesterPairs}(p,r)
\]
under the displayed hypotheses.  The upper-bound and modulus hypotheses are
kept explicit in this statement; the next corollary discharges them from the
power-of-two block tests used by the generated stage.  This is an exact
arithmetic result and does not introduce probability-construction or
sampling-law error.
\end{corollary}

\begin{formalized}
\leanfit{fhtStagePairs_mem_lower_mk}
\leanfit{fhtSylvesterStagePairs_mem_lower_mk}
\leanfit{fhtSylvesterStagePairs_mem_upper_mk}
\leanfit{fhtSylvesterStageScheduleExact_apply_lower_mk}
\leanfit{fhtSylvesterStageScheduleExact_apply_upper_mk}
\end{formalized}

\begin{corollary}[Power-of-two block-test formulas for one generated Sylvester FHT stage]
Let \(p,r\in\mathbb N\) with \(r<p\), let \(S_{p,r}\) be the exact generated
stage with stride \(2^r\) in dimension \(2^p\), and let
\(x\in\mathbb R^{2^p}\).  If a coordinate \(i\) satisfies the lower-half block
test
\[
  i\bmod(2\cdot 2^r)<2^r,
\]
then the upper partner is automatically in bounds,
\[
  i+2^r<2^p,
\]
and the exact generated stage satisfies
\[
  S_{p,r}(x)_i=x_i+x_{i+2^r}.
\]
If \(i\) satisfies the upper-half block test
\[
  2^r\le i\bmod(2\cdot 2^r),
\]
then \(2^r\le i\), the lower partner \(i-2^r\) satisfies
\[
  (i-2^r)\bmod(2\cdot 2^r)<2^r,
\]
and
\[
  S_{p,r}(x)_i=x_{i-2^r}-x_i.
\]
The proof uses only exact integer arithmetic: \(2\cdot 2^r\mid 2^p\), the
lower-half in-bounds lemma, and the upper-half subtract-stride modulus lemma.
It is therefore a deterministic exact-object FHT result.  The Rademacher and
uniform-row sampling laws remain exact mathematical inputs, and no
probability-construction, probability-normalization, or law-transfer error is
introduced.
\end{corollary}

\begin{formalized}
\leanfit{nat_add_stride_lt_of_mod_lt_of_dvd}
\leanfit{nat_sub_stride_mod_two_stride_lt_of_mod_ge}
\leanfit{two_mul_two_pow_dvd_two_pow_of_lt}
\leanfit{fhtSylvesterStage_lower_partner_lt_of_mod_lt}
\leanfit{fhtSylvesterStage_upper_value_le_of_mod_ge}
\leanfit{fhtSylvesterStage_upper_partner_mod_lt_of_mod_ge}
\leanfit{fhtSylvesterStageScheduleExact_apply_lower_of_mod_lt}
\leanfit{fhtSylvesterStageScheduleExact_apply_upper_of_mod_ge}
\end{formalized}

\begin{corollary}[Generated Sylvester stage block tests identify the stage bit]
Let \(p,r\in\mathbb N\), let \(i\in\{0,\ldots,2^p-1\}\), and consider the
same generated stage with stride \(2^r\).  If \(i\) is in the lower half of its
stride block,
\[
  i\bmod(2\cdot 2^r)<2^r,
\]
then the \(r\)-th bit of \(i\) is cleared:
\[
  \operatorname{testBit}(i,r)=\mathsf{false}.
\]
If \(i\) is in the upper half of its stride block,
\[
  2^r\le i\bmod(2\cdot 2^r),
\]
then the \(r\)-th bit is set:
\[
  \operatorname{testBit}(i,r)=\mathsf{true}.
\]
Thus the lower/upper cases in the generated FHT stage match the exact bit
predicate appearing in \(\operatorname{sylvesterHadamardParity}\).  The proof is
pure integer arithmetic using reduction modulo \(2^{r+1}\) and quotienting by
\(2^r\).  It is deterministic and exact; no sampling-law or probability
construction error is introduced.
\end{corollary}

\begin{formalized}
\leanfit{nat_testBit_eq_false_of_mod_two_mul_two_pow_lt}
\leanfit{nat_testBit_eq_true_of_two_pow_le_mod_two_mul_two_pow}
\leanfit{fhtSylvesterStage_testBit_eq_false_of_mod_lt}
\leanfit{fhtSylvesterStage_testBit_eq_true_of_mod_ge}
\end{formalized}

\begin{corollary}[Generated Sylvester stage partner bits]
Let \(p,r\in\mathbb N\), let \(i\in\{0,\ldots,2^p-1\}\), and consider the
generated Sylvester/Walsh stage with stride \(2^r\).  If \(i\) is a lower-half
coordinate,
\[
  i\bmod(2\cdot 2^r)<2^r,
\]
then the generated upper partner \(i+2^r\) has its \(r\)-th bit set:
\[
  \operatorname{testBit}(i+2^r,r)=\mathsf{true}.
\]
If \(i\) is an upper-half coordinate,
\[
  2^r\le i\bmod(2\cdot 2^r),
\]
then the generated lower partner \(i-2^r\) has its \(r\)-th bit cleared:
\[
  \operatorname{testBit}(i-2^r,r)=\mathsf{false}.
\]
These two partner-bit facts are the exact integer bridge needed when the
one-stage coordinate formulas \(x_i+x_{i+2^r}\) and \(x_{i-2^r}-x_i\) are
rewritten as the Sylvester/Walsh parity recurrence.  The theorem is
deterministic exact arithmetic for the generated FHT object.  The Rademacher
and uniform-row sampling laws remain exact mathematical inputs; no
probability-construction, normalization, repair, or law-transfer error is
introduced.
\end{corollary}

\begin{formalized}
\leanfit{fhtSylvesterStage_upper_partner_testBit_eq_true_of_mod_lt}
\leanfit{fhtSylvesterStage_lower_partner_testBit_eq_false_of_mod_ge}
\end{formalized}

\begin{corollary}[Generated Sylvester FHT schedule with explicit FP radius]
Let \(p\in\mathbb N\), and let the full generated Sylvester/Walsh FHT schedule be
\[
  \mathcal F_p
  =\bigcup_{\text{ordered }r=0}^{p-1}\operatorname{SylvesterPairs}(p,r),
\]
implemented as the concatenated finite list of stages with strides \(2^r\).
Then for every generated pair \((a,b)\),
\[
  (a,b)\in\mathcal F_p
  \quad\Longleftrightarrow\quad
  \exists r<p:\quad b=a+2^r,\qquad
  a\bmod(2\cdot 2^r)<2^r .
\]
In particular every generated pair is ordered and non-self:
\[
  a<b,\qquad a\ne b .
\]

Let \(x,\widehat x:\{0,\ldots,2^p-1\}\to\mathbb R\) and suppose
\[
  |\widehat x_i-x_i|\le E_i\qquad\text{for all }i .
\]
Let \(Y_p(x)\) be the exact vector obtained by applying the generated schedule
\(\mathcal F_p\), and let \(\widehat Y_p(\widehat x)\) be the result of applying
the rounded butterfly schedule, with each butterfly using
\((\operatorname{fladd},\operatorname{flsub})\).  If
\[
  B_i =
  \operatorname{FHTScheduleBudget}_{p,i}(\widehat x,E),
\]
then
\[
  |\widehat Y_p(\widehat x)_i-Y_p(x)_i|\le B_i,
  \qquad
  E_i\ge0\ \forall i \Longrightarrow B_i\ge0 .
\]

For a matrix input \(A,\widehat A\in\mathbb R^{2^p\times n}\) with
\[
  |\widehat A_{ij}-A_{ij}|\le E^A_{ij},
\]
let the exact scaled output use \(c=\sqrt{(2^p)^{-1}}\), and let the computed
output use
\[
  \widehat c=\operatorname{flsqrt}((2^p)^{-1}).
\]
Then the full generated columnwise implementation satisfies
\[
\begin{aligned}
&\left|
  \operatorname{flFHTSylvesterMatrix}_{\widehat c}(\widehat A)_{ij}
  - \operatorname{FHTSylvesterMatrix}_{c}(A)_{ij}
\right| \\
&\qquad\le
  \operatorname{FHTSylvesterMatrixBudget}_{ij}
  \left(
    \widehat c,\,
    c\,u,\,
    \widehat A,\,
    E^A
  \right).
\end{aligned}
\]
The corresponding \texttt{ComputedMatrix} constructor packages this entrywise
radius for the full generated schedule.

Moreover, the generated flat schedule is exactly the stage-by-stage recurrence
over the ordered stage list \((0,1,\ldots,p-1)\).  If
\(\operatorname{StageExact}_p\),
\(\operatorname{StageFP}_p\), and \(\operatorname{StageBudget}_p\) denote the
exact, rounded, and propagated-budget recurrences that apply one generated
stage at a time, then
\[
  Y_p(x)=\operatorname{StageExact}_p(0,1,\ldots,p-1;x),
\]
\[
  \widehat Y_p(\widehat x)
  =\operatorname{StageFP}_p(0,1,\ldots,p-1;\widehat x),
\]
and
\[
  B
  =\operatorname{StageBudget}_p(0,1,\ldots,p-1;\widehat x,E).
\]
Thus the transform-correctness theorem may be stated purely as a
stage-recurrence identity against the Sylvester/Hadamard sign table; the later
partial-transform recurrence and range induction below close that identity.  No
probability-law or sampling-distribution approximation is introduced in this
handoff.

This is an implementation-facing
non-probability arithmetic result: the stage indices are exact integer
computations, and the Rademacher signs, uniform row laws, and all sampling
probabilities remain exact mathematical inputs.  The exact transform-correctness
theorem is closed below; what remains open is layout/storage/overwrite effects
and final \texttt{ComputedPreconditioner} wiring for a chosen implementation
path.
\end{corollary}

\begin{formalized}
\leanfit{fhtSylvesterSchedulePairs}
\leanfit{mem_fhtSylvesterSchedulePairs_iff}
\leanfit{mem_fhtSylvesterSchedulePairs_iff_stage_rule}
\leanfit{fhtSylvesterSchedulePairs_stage_exists}
\leanfit{fhtSylvesterSchedulePairs_first_val_lt_second_val}
\leanfit{fhtSylvesterSchedulePairs_fst_ne_snd}
\leanfit{fhtPairScheduleExact_append}
\leanfit{flFhtPairSchedule_append}
\leanfit{fhtPairSchedulePropagatedErrorBudget_append}
\leanfit{fhtSylvesterScheduleExact}
\leanfit{flFhtSylvesterSchedule}
\leanfit{fhtSylvesterSchedulePropagatedErrorBudget}
\leanfit{fhtSylvesterSchedulePropagatedErrorBudget_nonneg}
\leanfit{flFhtSylvesterSchedule_propagated_error_bound}
\leanfit{fhtSylvesterStageScheduleExact}
\leanfit{flFhtSylvesterStageSchedule}
\leanfit{fhtSylvesterStageSchedulePropagatedErrorBudget}
\leanfit{fhtSylvesterStageScheduleListExact}
\leanfit{flFhtSylvesterStageScheduleList}
\leanfit{fhtSylvesterStageScheduleListPropagatedErrorBudget}
\leanfit{fhtSylvesterStageScheduleListExact_append}
\leanfit{flFhtSylvesterStageScheduleList_append}
\leanfit{fhtSylvesterStageScheduleListPropagatedErrorBudget_append}
\leanfit{fhtSylvesterStageScheduleListExact_range_succ}
\leanfit{flFhtSylvesterStageScheduleList_range_succ}
\leanfit{fhtSylvesterStageScheduleListPropagatedErrorBudget_range_succ}
\leanfit{fhtSylvesterStageScheduleListPropagatedErrorBudget_nonneg}
\leanfit{fhtPairScheduleExact_flatMap_sylvesterStagePairs}
\leanfit{flFhtPairSchedule_flatMap_sylvesterStagePairs}
\leanfit{fhtPairSchedulePropagatedErrorBudget_flatMap_sylvesterStagePairs}
\leanfit{fhtSylvesterScheduleExact_eq_stageScheduleListExact}
\leanfit{flFhtSylvesterSchedule_eq_stageScheduleList}
\leanfit{fhtSylvesterSchedulePropagatedErrorBudget_eq_stageScheduleList}
\leanfit{fhtScaledSylvesterScheduleMatrixExact}
\leanfit{flFhtScaledSylvesterScheduleMatrix}
\leanfit{fhtScaledSylvesterScheduleMatrixErrorBudget}
\leanfit{fhtScaledSylvesterScheduleMatrixErrorBudget_nonneg}
\leanfit{flFhtScaledSylvesterScheduleMatrix_error_bound}
\leanfit{flFhtScaledSylvesterScheduleMatrix_sqrtInvNatScale_error_bound}
\leanfit{ComputedMatrix.flScaledFhtSylvesterScheduleColumnsSqrtInvNat}
\leanfit{ComputedMatrix.flScaledFhtSylvesterScheduleColumnsSqrtInvNat_matrix}
\leanfit{ComputedMatrix.flScaledFhtSylvesterScheduleColumnsSqrtInvNat_abs_error}
\leanfit{ComputedMatrix.flScaledFhtSylvesterScheduleColumnsSqrtInvNat_entry_error_bound}
\end{formalized}

\begin{corollary}[Stage-list append and range-succ recurrence for generated FHT]
Let \(p\in\mathbb N\), let \(L_1,L_2\) be finite lists of generated stage
indices, and let \(x,\widehat x,E\in\mathbb R^{2^p}\).  Write
\[
  \operatorname{StageExact}_p(L;x),\qquad
  \operatorname{StageFP}_p(L;\widehat x),\qquad
  \operatorname{StageBudget}_p(L;\widehat x,E)
\]
for the exact stage recurrence, rounded stage recurrence, and propagated
entrywise error budget.  Then stage-list append composes exactly:
\[
  \operatorname{StageExact}_p(L_1\mathbin{+\!\!+}L_2;x)
  =
  \operatorname{StageExact}_p
  \bigl(L_2;\operatorname{StageExact}_p(L_1;x)\bigr),
\]
\[
  \operatorname{StageFP}_p(L_1\mathbin{+\!\!+}L_2;\widehat x)
  =
  \operatorname{StageFP}_p
  \bigl(L_2;\operatorname{StageFP}_p(L_1;\widehat x)\bigr),
\]
and
\[
\begin{aligned}
&\operatorname{StageBudget}_p(L_1\mathbin{+\!\!+}L_2;\widehat x,E)\\
&\quad =
  \operatorname{StageBudget}_p
  \bigl(
    L_2;\operatorname{StageFP}_p(L_1;\widehat x),
    \operatorname{StageBudget}_p(L_1;\widehat x,E)
  \bigr).
\end{aligned}
\]
In particular, with \(R_r=(0,1,\ldots,r-1)\), the successor recurrence is
\[
  \operatorname{StageExact}_p(R_{r+1};x)
  =
  \operatorname{StageExact}_p
  \bigl([r];\operatorname{StageExact}_p(R_r;x)\bigr),
\]
\[
  \operatorname{StageFP}_p(R_{r+1};\widehat x)
  =
  \operatorname{StageFP}_p
  \bigl([r];\operatorname{StageFP}_p(R_r;\widehat x)\bigr),
\]
and the propagated budget obeys the analogous one-stage update.  Equivalently,
the exact induction step is all previous generated stages followed by the
single current stage \(r\).

This closes the deterministic list-recursion spine needed for the generated
FHT realization proof.  The next recurrence corollary identifies the one-stage
step with the partial Sylvester/Walsh table, using the lower/upper coordinate
formulas and the concrete sign recurrences from the preceding corollaries.  No
probability construction, probability normalization, sampling-law transfer, or
floating-point probability error is introduced.
\end{corollary}

\begin{formalized}
\leanfit{fhtSylvesterStageScheduleListExact_append}
\leanfit{flFhtSylvesterStageScheduleList_append}
\leanfit{fhtSylvesterStageScheduleListPropagatedErrorBudget_append}
\leanfit{fhtSylvesterStageScheduleListExact_range_succ}
\leanfit{flFhtSylvesterStageScheduleList_range_succ}
\leanfit{fhtSylvesterStageScheduleListPropagatedErrorBudget_range_succ}
\end{formalized}

\begin{corollary}[Partial Sylvester/Walsh table anchors for generated FHT]
For \(p,t\in\mathbb N\), define the partial sign table
\[
  S_{p,t}(i,j)
  =
  (-1)^{
    \sum_{\ell=0}^{t-1}
      \mathbf 1\{\operatorname{bit}_\ell(i)=1
      \text{ and }\operatorname{bit}_\ell(j)=1\}
  }.
\]
The corresponding partial transform is
\[
  T_{p,t}(x)_i
  =
  \sum_{j:\ \lfloor j/2^t\rfloor=\lfloor i/2^t\rfloor}
    S_{p,t}(i,j)x_j .
\]
Then the endpoints of the generated-stage induction are exact:
\[
  T_{p,0}(x)=x,
\]
because the quotient support with \(2^0=1\) forces \(j=i\) and the empty
parity sum has sign \(1\), and
\[
  T_{p,p}(x)_i
  =
  \sum_j S_p(i,j)x_j,
\]
because every \(i,j<2^p\) satisfies
\(\lfloor i/2^p\rfloor=\lfloor j/2^p\rfloor=0\), while the partial parity sum
uses exactly the \(p\) bit coordinates of the concrete Sylvester/Walsh table.

Thus the generated-FHT realization proof has a precise induction
target: one stage maps \(T_{p,r}\) to \(T_{p,r+1}\), and the range-succ
stage-list recurrence plus the full-depth identity complete the proof.  The
next corollary closes this one-stage step.
This is exact deterministic integer/Boolean arithmetic; no probability
construction, sampling-law transfer, or floating-point probability error is
introduced.
\end{corollary}

\begin{formalized}
\leanfit{sylvesterHadamardPartialParity}
\leanfit{sylvesterHadamardPartialSignPattern}
\leanfit{sylvesterHadamardPartialUnscaledApply}
\leanfit{sylvesterHadamardPartialSignPattern_zero}
\leanfit{sylvesterHadamardPartialUnscaledApply_zero}
\leanfit{sylvesterHadamardPartialSignPattern_full}
\leanfit{sylvesterHadamardPartialUnscaledApply_full}
\end{formalized}

\begin{corollary}[One-stage partial Sylvester/Walsh recurrence for generated FHT]
Let \(r<p\).  For each \(i\in\{0,\ldots,2^p-1\}\), write
\[
  \operatorname{low}_r(i)
  \quad\Longleftrightarrow\quad
  i\bmod (2\cdot 2^r)<2^r .
\]
Then the exact generated FHT stage advances the partial table by one bit:
\[
  \operatorname{Stage}_{p,r}\bigl(T_{p,r}(x)\bigr)
  =
  T_{p,r+1}(x).
\]
More explicitly, in the lower half,
\[
  T_{p,r+1}(x)_i
  =
  T_{p,r}(x)_i+T_{p,r}(x)_{i+2^r},
  \qquad \operatorname{low}_r(i),
\]
and in the upper half,
\[
  T_{p,r+1}(x)_i
  =
  T_{p,r}(x)_{i-2^r}-T_{p,r}(x)_i,
  \qquad \neg\operatorname{low}_r(i).
\]
The proof splits each quotient block
\(\lfloor j/2^{r+1}\rfloor=\lfloor i/2^{r+1}\rfloor\) into the two
\(2^r\)-blocks used by the butterfly, proves the partial sign recurrence
\[
  S_{p,r+1}(i,j)
  =
  \begin{cases}
    -S_{p,r}(i,j),&
      \operatorname{bit}_r(i)=\operatorname{bit}_r(j)=1,\\
     S_{p,r}(i,j),&\text{otherwise,}
  \end{cases}
\]
and uses the generated-partner bit-agreement facts for \(i\pm 2^r\).
This is exact deterministic integer/FHT arithmetic and does not introduce any
probability-construction, sampling-law, or FP-probability term.
\end{corollary}

\begin{formalized}
\leanfit{nat_div_stride_eq_two_mul_block_div_add_mod_div}
\leanfit{nat_div_stride_eq_two_mul_block_div_of_mod_lt}
\leanfit{nat_div_stride_eq_two_mul_block_div_add_one_of_mod_ge}
\leanfit{nat_add_stride_div_stride_eq_two_mul_block_div_add_one_of_mod_lt}
\leanfit{nat_add_stride_div_two_stride_eq_of_mod_lt}
\leanfit{nat_sub_stride_div_stride_eq_two_mul_block_div_of_mod_ge}
\leanfit{nat_sub_stride_div_two_stride_eq_of_mod_ge}
\leanfit{nat_div_two_mul_stride_eq_div_stride_div_two}
\leanfit{nat_div_two_stride_eq_of_div_stride_eq}
\leanfit{sylvesterHadamardPartialParityWeight_succ}
\leanfit{sylvesterHadamardPartialSignPattern_succ_eq_or_neg}
\leanfit{sylvesterHadamardPartialSignPattern_succ_eq_of_stage_bit_false}
\leanfit{sylvesterHadamardPartialSignPattern_succ_eq_or_neg_of_stage_bit_true}
\leanfit{sylvesterHadamardPartialParityWeight_eq_of_bits}
\leanfit{sylvesterHadamardPartialSignPattern_eq_of_bits}
\leanfit{sylvesterHadamardPartialSignPattern_upper_partner_eq_of_mod_lt}
\leanfit{sylvesterHadamardPartialSignPattern_upper_eq_lower_partner_of_mod_ge}
\leanfit{sylvesterHadamardPartialUnscaledApply_succ_lower}
\leanfit{sylvesterHadamardPartialUnscaledApply_succ_upper}
\leanfit{fhtSylvesterStageScheduleExact_partialUnscaledApply_eq_succ}
\end{formalized}

\begin{corollary}[Generated FHT realization and signed-preconditioner bridge]
For \(p\in\mathbb N\), let \(S_p\) be the concrete Sylvester/Walsh sign table
\[
  S_p(i,j)=
  \begin{cases}
    -1, & \text{if }\operatorname{parity}(i\wedge j)=1,\\
     1, & \text{otherwise.}
  \end{cases}
\]
The exact generated FHT schedule realizes this table:
\[
  Y_p(x)_i=\sum_j S_p(i,j)x_j
  \quad\text{for every }x\text{ and }i,
\]
where \(Y_p(x)\) is the exact generated FHT schedule from the preceding
corollary.  The proof first shows, for every \(t\le p\),
\[
  \operatorname{StageExact}_p(0,\ldots,t-1;x)=T_{p,t}(x),
\]
and then sets \(t=p\), where \(T_{p,p}\) is the full bit-parity transform.
Consequently, for every matrix
\(A\in\mathbb R^{2^p\times n}\) and every exact scale \(c\),
\[
  \operatorname{FHTSylvesterMatrix}_{c}(A)_{ij}
  =
  c\sum_k S_p(i,k)A_{kj}.
\]
In particular, for a realized sign vector \(\sigma\) and an input basis \(U\),
\[
  \operatorname{FHTSylvesterMatrix}_{c}
    \bigl((k,j)\mapsto \sigma_k U_{kj}\bigr)
  =
  \operatorname{preconditionRows}
    \bigl((cS_p)\operatorname{diag}(\sigma)\bigr)U .
\]
Thus the generated FHT implementation is tied to the exact \(HDU\)
preconditioner used by the existing SRHT exact-law and floating-point transfer
theorems without a remaining realization hypothesis.  This result is exact and
deterministic; it does not introduce any probability-construction, sampling-law
normalization, or law-transfer error.  Layout-specific rounded storage and
overwrite effects remain separate implementation obligations.
\end{corollary}

\begin{formalized}
\leanfit{sylvesterHadamardUnscaledApply}
\leanfit{fhtSylvesterScheduleRealizesSignPattern}
\leanfit{fhtSylvesterScheduleRealizesSignPattern_zero}
\leanfit{fhtSylvesterStageScheduleListExact_range_eq_partialUnscaledApply}
\leanfit{fhtSylvesterScheduleRealizesSignPattern_generated}
\leanfit{sylvesterHadamardScaledMatrixApply}
\leanfit{fhtScaledSylvesterScheduleMatrixExact_eq_sylvesterHadamardScaledMatrixApply_of_realizes}
\leanfit{fhtScaledSylvesterScheduleMatrixExact_eq_sylvesterHadamardScaledMatrixApply}
\leanfit{fhtScaledSylvesterScheduleMatrixExact_signed_eq_preconditionRows_of_realizes}
\leanfit{fhtScaledSylvesterScheduleMatrixExact_signed_eq_preconditionRows}
\end{formalized}

\begin{corollary}[Fast generated-FHT computed preconditioner]
Let \(S_p\) be the concrete Sylvester/Walsh sign table in dimension \(2^p\),
let \(c=\sqrt{((2^p)^{-1})}\), and let \(\sigma\in\{\pm1\}^{2^p}\) be a
realized Rademacher sign vector.  The sampling law for \(\sigma\) is exact.
Suppose the implementation stores or computes the sign vector as
\(\widehat\sigma\) with entrywise certificate
\[
  |\widehat\sigma_i-\sigma_i|\le E^\sigma_i,\qquad E^\sigma_i\ge0,
\]
and forms the computed diagonal input
\[
  \widehat D_{ij}=
  \begin{cases}
    \widehat\sigma_i, & i=j,\\
    0, & i\ne j,
  \end{cases}
  \qquad
  E^D_{ij}=
  \begin{cases}
    E^\sigma_i, & i=j,\\
    0, & i\ne j.
  \end{cases}
\]
Let \(\widehat c=\operatorname{flsqrt}((2^p)^{-1})\) and
\(\eta=c\,u\).  The fast generated-FHT preconditioner is the rounded
columnwise schedule
\[
  \widehat P
  =
  \operatorname{flFHTSylvesterMatrix}_{\widehat c}(\widehat D),
\]
and the exact analysis preconditioner is
\[
  P=(cS_p)\operatorname{diag}(\sigma).
\]
Then for every \(i,j\),
\[
  |\widehat P_{ij}-P_{ij}|
  \le
  B^{\rm FHT}_{ij},
\]
where \(B^{\rm FHT}_{ij}\) is the generated-FHT propagated budget obtained from
\[
  \operatorname{FHTBudget}_{p,\widehat c,\eta}(\widehat D,E^D)_{ij}.
\]
Equivalently, the implementation supplies a \texttt{ComputedPreconditioner}
certificate for \(P\).  The exact bridge used in the certificate is
\[
  \operatorname{FHTSylvesterMatrix}_{c}(\operatorname{diag}(\sigma))
  =
  (cS_p)\operatorname{diag}(\sigma).
\]
Thus the implementation no longer silently uses an exact materialized \(H D\)
matrix: sign storage, diagonal input error, butterfly arithmetic, and rounded
normalization are all charged.  Under the exact Rademacher/uniform-row product
law, the resulting computed-left perturbation event for
\(\widehat V=\fl(\widehat P U)\) has probability one; the downstream budget
adds \(\widehat V\)-formation, row-scaling, and Gram-dot-product errors.

If an implementation forms the signed input \(D_\sigma U\) directly before the
FHT stages, a separate rowwise multiplication certificate is available:
\[
  |\fl(\widehat\sigma_i\widehat U_{ij})-\sigma_iU_{ij}|
  \le
  u|\widehat\sigma_i\widehat U_{ij}|
  +|\widehat\sigma_i|\,E^U_{ij}
  +E^\sigma_i|U_{ij}|.
\]
For arbitrary array implementations, the remaining FHT obligation is
layout-specific in-place storage, copy, aliasing, or overwrite behavior if the
chosen routine differs from the functional pair-update schedule formalized
here.  The next result closes one concrete writeback/copy model.
\end{corollary}

\begin{formalized}
\leanfit{ComputedMatrix.flRowSignMul}
\leanfit{ComputedMatrix.flRowSignMul_entry_error_bound}
\leanfit{fhtScaledSylvesterScheduleMatrixExact_diag_eq_matMul_diag}
\leanfit{ComputedPreconditioner.flSignedHadamardSylvesterFhtSchedule}
\leanfit{ComputedPreconditioner.flSignedHadamardSylvesterFhtSchedule_entry_error_bound}
\leanfit{signedHadamardSylvesterFhtSchedulePreconditioner}
\leanfit{signedHadamardSylvesterFhtScheduleStoredSignPreconditioner}
\leanfit{signedHadamardUniformRowTraceProbability_eventProb_sylvesterFhtScheduleComputedLeftPreconditioned_fl_uniformRowPerturbEvent_eq_one}
\leanfit{signedHadamardUniformRowTraceProbability_eventProb_sylvesterFhtScheduleStoredSignComputedLeftPreconditioned_fl_uniformRowPerturbEvent_eq_one}
\end{formalized}

\begin{theorem}[Generated-FHT SRHT endpoint without supplied orthogonality]
Fix \(p,r,n,s\in\mathbb N\), let \(m=2^p\), and let
\[
  S_p(i,k)=\operatorname{sylvesterHadamardSignPattern}(p,i,k),
  \qquad
  H_p(i,k)=\sqrt{m^{-1}}\,S_p(i,k).
\]
Let \(U\in\mathbb R^{m\times r}\) have orthonormal columns, let
\(C\in\mathbb R^{r\times n}\), and let the exact analysis input be
\[
  A=UC=\operatorname{preconditionColumns}(U,C).
\]
The Rademacher sign law and the downstream uniform row law are exact
mathematical laws.  For each realized sign vector, the computed
non-probability objects are:
\[
  \widehat P_\omega
    =\texttt{signedHadamardSylvesterFhtScheduleStoredSignPreconditioner}
      (fp,\omega),
  \qquad
  \widehat d
    =\texttt{uniformRowFlSqrtMulInvSqrtScaleDen}(fp,m,s),
\]
followed by the rounded product
\(\widehat Y_\omega=\fl(\widehat P_\omega A)\), rounded sampled-row
divisions by \(\widehat d\), and rounded sampled-Gram dot products.  Assume
the same positivity, \(\gamma_m\), \(\gamma_s\), \(\theta\), and SRHT
sample-budget hypotheses as in the source-sharp SRHT logarithmic theorem, with
\[
  t=\sqrt{\frac{8\log(m/\delta_{\rm pre})}{m}},\quad
  S_{\rm rad}=\sqrt{r/m}+t,\quad
  L=mS_{\rm rad}^2,
\]
\[
  \beta_+=\frac{\exp(\theta L)-\theta L-1}{L^2},
  \qquad
  \beta_-=\exp(\theta)-\theta-1,
\]
and
\[
  e^{-\theta s\varepsilon}r
  \left(e^{s\beta_+(L^2+1)}+e^{s\beta_-(L^2+1)}\right)
  \le \delta_{\rm sample}.
\]
Then the computed two-sided sampled-Gram event centered at \(A^TA\) satisfies
\[
  \mathbb P\!\left[
    -\varepsilon I \preceq
      \widehat G_{\rm SRHT}^{\rm fp}-A^TA
    \preceq \varepsilon I
    \ \text{with radius enlarged by }\
    \texttt{signedHadamardComputedLeftUniformRowComputedDenPerturbBudget}
  \right]
  \ge 1-(\delta_{\rm pre}+\delta_{\rm sample}).
\]
The exact objects in the theorem are \(U\), \(C\), \(A=UC\), \(S_p\), \(H_p\),
and the two sampling laws.  The computed objects are \(\widehat P_\omega\),
\(\widehat Y_\omega\), \(\widehat d\), the sampled scaled rows, and the Gram
dot products.  The normalized Sylvester/Walsh table is proved orthogonal
inside Lean, so no external \(H\)-orthogonality or flatness hypothesis appears.
The stored-input siblings close the implementation paths in which the actual
input matrix is first realized in floating point before preconditioning:
\[
  \widehat A^{\rm in}_{ij}\in
  \{\operatorname{flmul}(A_{ij},1),
    \operatorname{fladd}(A_{ij},0),
    \operatorname{flsub}(A_{ij},0)\},
  \qquad
  A=UC.
\]
Here \(U,C,A,S_p,H_p\) and the sampling laws are exact analysis references,
while \(\widehat A^{\rm in}\), \(\widehat P_\omega\),
\(\widehat Y^{\rm in}_\omega=\fl(\widehat P_\omega\widehat A^{\rm in})\),
the denominator, row divisions, and Gram dot products are computed
non-probability objects.  The new product budget adds the irreducible input
storage term
\[
  \sum_a |(P_\omega)_{ia}|\,u|A_{aj}|
\]
inside the entrywise bound for
\(\fl(\widehat P_\omega\widehat A^{\rm in})-P_\omega A\).  Thus these
endpoints have no perturbation-event or certificate-existence assumption for
the stored input matrix.
The same final theorem is also instantiated for modeled writeback/copy
routines.  In the all-coordinate add-zero routine, after every generated FHT
pair update, each output coordinate is stored as
\(\operatorname{fladd}(y_i,0)\).  In the all-coordinate multiply-one and
subtract-zero routines, each output coordinate is stored as
\(\operatorname{flmul}(y_i,1)\) or \(\operatorname{flsub}(y_i,0)\).  In the
modified-coordinate add-zero, multiply-one, and subtract-zero routines, only
the two coordinates changed by the butterfly receive the corresponding
writeback.  In all cases, the writeback radius is included in the same
computed-preconditioner and downstream sampled-Gram budget.
\end{theorem}

\begin{formalized}
\leanfit{fhtButterflyExact_sq_sum}
\leanfit{fhtButterflyExact_inner_sum}
\leanfit{fhtPairUpdateExact_inner_sum}
\leanfit{vecNorm2Sq_fhtPairUpdateExact}
\leanfit{sylvesterHadamardSignPattern_symm}
\leanfit{exists_testBit_ne_of_fin_two_pow_ne}
\leanfit{sylvesterStageBitFlip}
\leanfit{sylvesterStageBitFlip_involutive}
\leanfit{sylvesterStageBitFlip_bijective}
\leanfit{sylvesterHadamardSignPattern_mul_stageBitFlip_eq_neg}
\leanfit{sylvesterHadamardSignPattern_col_inner_eq_zero_of_ne}
\leanfit{sylvesterHadamardSignPattern_col_inner}
\leanfit{sylvesterHadamardSignPattern_row_inner}
\leanfit{sylvesterHadamardScaled_col_inner}
\leanfit{sylvesterHadamardScaled_row_inner}
\leanfit{isOrthogonal_sqrt_inv_nat_mul_sylvesterSignPattern}
\leanfit{ComputedMatrix.flMulOne}
\leanfit{ComputedMatrix.flMulOne_entry_error_bound}
\leanfit{ComputedMatrix.flAddZeroRight}
\leanfit{ComputedMatrix.flAddZeroRight_entry_error_bound}
\leanfit{ComputedMatrix.flSubZeroRight}
\leanfit{ComputedMatrix.flSubZeroRight_entry_error_bound}
\leanfit{signedHadamardComputedLeftInputUniformRowComputedDenPerturbBudget}
\leanfit{signedHadamardUniformRowTraceProbability_eventProb_computedLeftInputPreconditioned_fl_uniformRowComputedDenPerturbEvent_eq_one}
\noindent{\scriptsize\nolinkurl{signedHadamardUniformRowTraceProbability_eventProb_}}\par
\noindent{\scriptsize\quad\nolinkurl{sylvesterFhtScheduleStoredSignComputedLeftPreconditioned_factoredInput_}}\par
\noindent{\scriptsize\quad\nolinkurl{fl_uniformRowSampleGramDotWithFlSqrtMulInvSqrtDen_two_sided_finiteLoewnerLe_}}\par
\noindent{\scriptsize\quad\nolinkurl{ge_one_sub_delta_srht_log_preprocess}}\par
\noindent{\scriptsize\nolinkurl{signedHadamardUniformRowTraceProbability_eventProb_}}\par
\noindent{\scriptsize\quad\nolinkurl{sylvesterFhtScheduleStoredSignStoredInputMulOneComputedLeftInputPreconditioned_}}\par
\noindent{\scriptsize\quad\nolinkurl{factoredInput_fl_uniformRowSampleGramDotWithFlSqrtMulInvSqrtDen_}}\par
\noindent{\scriptsize\quad\nolinkurl{two_sided_finiteLoewnerLe_ge_one_sub_delta_srht_log_preprocess}}\par
\noindent{\scriptsize\nolinkurl{signedHadamardUniformRowTraceProbability_eventProb_}}\par
\noindent{\scriptsize\quad\nolinkurl{sylvesterFhtScheduleStoredSignStoredInputAddZeroRightComputedLeftInputPreconditioned_}}\par
\noindent{\scriptsize\quad\nolinkurl{factoredInput_fl_uniformRowSampleGramDotWithFlSqrtMulInvSqrtDen_}}\par
\noindent{\scriptsize\quad\nolinkurl{two_sided_finiteLoewnerLe_ge_one_sub_delta_srht_log_preprocess}}\par
\noindent{\scriptsize\nolinkurl{signedHadamardUniformRowTraceProbability_eventProb_}}\par
\noindent{\scriptsize\quad\nolinkurl{sylvesterFhtScheduleStoredSignStoredInputSubZeroRightComputedLeftInputPreconditioned_}}\par
\noindent{\scriptsize\quad\nolinkurl{factoredInput_fl_uniformRowSampleGramDotWithFlSqrtMulInvSqrtDen_}}\par
\noindent{\scriptsize\quad\nolinkurl{two_sided_finiteLoewnerLe_ge_one_sub_delta_srht_log_preprocess}}\par
\noindent{\scriptsize\nolinkurl{signedHadamardUniformRowTraceProbability_eventProb_}}\par
\noindent{\scriptsize\quad\nolinkurl{sylvesterFhtScheduleStoredSignStoredAddZeroRightComputedLeftPreconditioned_}}\par
\noindent{\scriptsize\quad\nolinkurl{factoredInput_fl_uniformRowSampleGramDotWithFlSqrtMulInvSqrtDen_}}\par
\noindent{\scriptsize\quad\nolinkurl{two_sided_finiteLoewnerLe_ge_one_sub_delta_srht_log_preprocess}}\par
\noindent{\scriptsize\nolinkurl{signedHadamardUniformRowTraceProbability_eventProb_}}\par
\noindent{\scriptsize\quad\nolinkurl{sylvesterFhtScheduleStoredSignModifiedStoredAddZeroRightComputedLeftPreconditioned_}}\par
\noindent{\scriptsize\quad\nolinkurl{factoredInput_fl_uniformRowSampleGramDotWithFlSqrtMulInvSqrtDen_}}\par
\noindent{\scriptsize\quad\nolinkurl{two_sided_finiteLoewnerLe_ge_one_sub_delta_srht_log_preprocess}}\par
\end{formalized}

\begin{corollary}[Stored add-zero writeback for fast generated-FHT]
Use the notation and exact probability laws of the preceding corollary.  Let
\(\mathcal P_p\) be the generated Sylvester/Walsh pair schedule in dimension
\(2^p\).  For one pair \((r,s)\), let
\[
  y^{\rm pair}=\operatorname{FHTPairUpdate}_{r,s}(x),\qquad
  \widehat y^{\rm pair}
    =\operatorname{flFHTPairUpdate}_{r,s}(\widehat x),
\]
and define the stored output by
\[
  \widehat y^{\rm store}_i
  =\fl_{\rm add}(\widehat y^{\rm pair}_i,0).
\]
If the input certificate satisfies
\[
  |\widehat x_i-x_i|\le E_i,\qquad E_i\ge0,
\]
then the stored pair update has the entrywise bound
\[
  |\widehat y^{\rm store}_i-y^{\rm pair}_i|
  \le
  B^{\rm pair}_{r,s,i}(\widehat x,E)
  +u\,|\widehat y^{\rm pair}_i|,
\]
where \(B^{\rm pair}_{r,s,i}\) is the usual propagated rounded-butterfly
budget.  Equivalently,
\[
  E^{\rm store,+}_{r,s,i}
  =
  B^{\rm pair}_{r,s,i}(\widehat x,E)
  +u\,|\operatorname{flFHTPairUpdate}_{r,s}(\widehat x)_i|
\]
is a valid propagated radius for one rounded pair update followed by an
add-zero output writeback/copy.

Composing this recurrence over \(\mathcal P_p\), let
\[
  \widehat z_i
  =
  \operatorname{flFHTSchedule}^{\rm store,+}_{\mathcal P_p}
  (\widehat x)_i,\qquad
  E^{\rm store,+}_{p,i}
  =
  \operatorname{FHTScheduleBudget}^{\rm store,+}_{p,i}(\widehat x,E).
\]
For the scaled transform, let \(c=\sqrt{((2^p)^{-1})}\),
\(\widehat c=\operatorname{flsqrt}((2^p)^{-1})\), and \(\eta=c\,u\).  Then
\[
  \left|
    \fl_{\rm mul}(\widehat c,\widehat z_i)
    -
    c\,\operatorname{FHTSchedule}_{\mathcal P_p}(x)_i
  \right|
  \le
  u|\widehat c\,\widehat z_i|
  +|\widehat c|\,E^{\rm store,+}_{p,i}
  +\eta\bigl(|\widehat z_i|+E^{\rm store,+}_{p,i}\bigr).
\]
Columnwise, this gives a computed matrix certificate
\[
  \left|
    \operatorname{flFHTSylvesterMatrix}^{\rm store,+}_{\widehat c}(\widehat A)_{ij}
    -
    \operatorname{FHTSylvesterMatrix}_{c}(A)_{ij}
  \right|
  \le
  B^{\rm store,+}_{ij},
\]
where \(B^{\rm store,+}\) is
\[
  \operatorname{FHTSylvesterMatrixBudget}^{\rm store,+}_{p,\widehat c,\eta}
  (\widehat A,E).
\]

In particular, with the computed diagonal sign input \(\widehat D\) and its
entrywise sign-storage certificate \(E^D\), the implemented preconditioner
\[
  \widehat P^{\rm store,+}
  =
  \operatorname{flFHTSylvesterMatrix}^{\rm store,+}_{\widehat c}(\widehat D)
\]
satisfies, for every \(i,j\),
\[
  |\widehat P^{\rm store,+}_{ij}
    -((cS_p)\operatorname{diag}(\sigma))_{ij}|
  \le
  B^{\rm store,+}_{ij}.
\]
Thus this is a fully computed-object implementation-facing certificate for
the modeled routine ``rounded butterfly, then \(\fl_{\rm add}(\cdot,0)\) output
writeback/copy.''  The Rademacher and uniform-row sampling laws remain exact.
Under their exact product law, the computed-left perturbation event for
\[
  \widehat V=\fl(\widehat P^{\rm store,+}U)
\]
has probability one, with the downstream radius additionally charging
\(\widehat V\)-formation, row scaling, and Gram-dot-product errors.  If the
realized signs are also stored by the existing rounded sign-copy certificate,
the same probability-one event holds with that sign-storage radius included in
\(E^D\).  Memory order, aliasing, vectorized layout, or in-place overwrite
semantics not represented by this pair-update plus add-zero writeback model
remain separate implementation obligations.
\end{corollary}

\begin{formalized}
\leanfit{flFhtPairUpdateStoredAddZeroRight}
\leanfit{fhtPairUpdateStoredAddZeroRightPropagatedErrorBudget}
\leanfit{flFhtPairUpdateStoredAddZeroRight_propagated_error_bound}
\leanfit{flFhtPairScheduleStoredAddZeroRight}
\leanfit{fhtPairScheduleStoredAddZeroRightPropagatedErrorBudget}
\leanfit{flFhtPairScheduleStoredAddZeroRight_propagated_error_bound}
\leanfit{flFhtSylvesterScheduleStoredAddZeroRight}
\leanfit{flFhtSylvesterScheduleStoredAddZeroRight_propagated_error_bound}
\leanfit{flFhtScaledPairScheduleStoredAddZeroRight}
\leanfit{flFhtScaledPairScheduleStoredAddZeroRight_sqrtInvNatScale_error_bound}
\leanfit{flFhtScaledPairScheduleMatrixStoredAddZeroRight}
\leanfit{flFhtScaledPairScheduleMatrixStoredAddZeroRight_sqrtInvNatScale_error_bound}
\leanfit{flFhtScaledSylvesterScheduleMatrixStoredAddZeroRight}
\leanfit{flFhtScaledSylvesterScheduleMatrixStoredAddZeroRight_sqrtInvNatScale_error_bound}
\leanfit{ComputedMatrix.flScaledFhtSylvesterScheduleColumnsSqrtInvNatStoredAddZeroRight}
\leanfit{ComputedMatrix.flScaledFhtSylvesterScheduleColumnsSqrtInvNatStoredAddZeroRight_entry_error_bound}
\leanfit{ComputedPreconditioner.flSignedHadamardSylvesterFhtScheduleStoredAddZeroRight}
\leanfit{ComputedPreconditioner.flSignedHadamardSylvesterFhtScheduleStoredAddZeroRight_entry_error_bound}
\leanfit{signedHadamardSylvesterFhtScheduleStoredAddZeroRightPreconditioner}
\leanfit{signedHadamardSylvesterFhtScheduleStoredSignStoredAddZeroRightPreconditioner}
\noindent{\scriptsize\nolinkurl{signedHadamardUniformRowTraceProbability_eventProb_}}\par
\noindent{\scriptsize\quad\nolinkurl{sylvesterFhtScheduleStoredAddZeroRightComputedLeftPreconditioned_}}\par
\noindent{\scriptsize\quad\nolinkurl{fl_uniformRowPerturbEvent_eq_one}}\par
\noindent{\scriptsize\nolinkurl{signedHadamardUniformRowTraceProbability_eventProb_}}\par
\noindent{\scriptsize\quad\nolinkurl{sylvesterFhtScheduleStoredSignStoredAddZeroRightComputedLeftPreconditioned_}}\par
\noindent{\scriptsize\quad\nolinkurl{fl_uniformRowPerturbEvent_eq_one}}\par
\end{formalized}

\begin{corollary}[Generated-SRHT multiply-one/subtract-zero FHT writeback]
Fix one FHT butterfly pair \((r,s)\) and let
\[
  y^{\rm pair}
    =\operatorname{FHTPairUpdate}_{r,s}(x),\qquad
  \widehat y^{\rm pair}
    =\operatorname{flFHTPairUpdate}_{r,s}(\widehat x).
\]
Assume the input certificate
\[
  |\widehat x_i-x_i|\le E_i,\qquad E_i\ge0
\]
for every coordinate \(i\).  Let
\[
  B^{\rm pair}_{r,s,i}(\widehat x,E)
\]
be the existing propagated radius for one rounded pair update.  Define two
stored outputs
\[
  \widehat y^{\times 1}_i
    =\fl_{\rm mul}(\widehat y^{\rm pair}_i,1),
  \qquad
  \widehat y^{-0}_i
    =\fl_{\rm sub}(\widehat y^{\rm pair}_i,0).
\]
Then both all-coordinate writeback/copy variants satisfy the same explicit
one-pair radius:
\[
  |\widehat y^{\times 1}_i-y^{\rm pair}_i|
  \le
  B^{\rm pair}_{r,s,i}(\widehat x,E)
  +u\,|\widehat y^{\rm pair}_i|
\]
and
\[
  |\widehat y^{-0}_i-y^{\rm pair}_i|
  \le
  B^{\rm pair}_{r,s,i}(\widehat x,E)
  +u\,|\widehat y^{\rm pair}_i|.
\]
Equivalently,
\[
  E^{\times 1}_{r,s,i}
  =
  B^{\rm pair}_{r,s,i}(\widehat x,E)
  +u\,|\operatorname{flFHTPairUpdate}_{r,s}(\widehat x)_i|
\]
and
\[
  E^{-0}_{r,s,i}
  =
  B^{\rm pair}_{r,s,i}(\widehat x,E)
  +u\,|\operatorname{flFHTPairUpdate}_{r,s}(\widehat x)_i|
\]
are valid nonnegative propagated radii for the two one-pair routines.

Now let \(\mathcal P\) be any ordered list of FHT pairs.  Define
\[
  \widehat z^{\times 1}
  =
  \operatorname{flFHTSchedule}^{\times 1}_{\mathcal P}(\widehat x),
  \qquad
  \widehat z^{-0}
  =
  \operatorname{flFHTSchedule}^{-0}_{\mathcal P}(\widehat x)
\]
by applying the corresponding one-pair routine at each pair in the list.  Let
\[
  B^{\times 1}_{\mathcal P,i}(\widehat x,E),
  \qquad
  B^{-0}_{\mathcal P,i}(\widehat x,E)
\]
be the recursively propagated schedules of the one-pair radii above.  Then,
under the same input certificate,
\[
  |\widehat z^{\times 1}_i
    -\operatorname{FHTSchedule}_{\mathcal P}(x)_i|
  \le
  B^{\times 1}_{\mathcal P,i}(\widehat x,E)
\]
and
\[
  |\widehat z^{-0}_i
    -\operatorname{FHTSchedule}_{\mathcal P}(x)_i|
  \le
  B^{-0}_{\mathcal P,i}(\widehat x,E).
\]
The two schedule budgets are nonnegative whenever the input budget \(E\) is
coordinatewise nonnegative.

Specialize now to the generated Sylvester/Walsh schedule
\(\mathcal P_p\) in dimension \(2^p\).  Let
\[
  c=\sqrt{((2^p)^{-1})},\qquad
  \widehat c=\fl_{\rm sqrt}((2^p)^{-1}),\qquad
  \eta=c\,u .
\]
For a computed matrix \(\widehat A\) with certificate
\[
  |\widehat A_{ij}-A_{ij}|\le E_{ij},\qquad E_{ij}\ge0,
\]
define the two scaled generated-FHT outputs
\[
  \widehat Z^{\times 1}
   =
   \operatorname{flScaledFHT}^{\times1}_{p}(\widehat c,\widehat A),
  \qquad
  \widehat Z^{-0}
   =
   \operatorname{flScaledFHT}^{-0}_{p}(\widehat c,\widehat A).
\]
Then there are explicit entrywise budgets
\[
  C^{\times1}_{p,ij}(\widehat A,E),\qquad
  C^{-0}_{p,ij}(\widehat A,E)
\]
obtained by the recursive schedule budgets above plus the rounded final scale
and scale-radius terms, and they satisfy
\[
  |\widehat Z^{\times1}_{ij}
    -\operatorname{ScaledFHT}_{p,c}(A)_{ij}|
  \le C^{\times1}_{p,ij}(\widehat A,E),
\]
\[
  |\widehat Z^{-0}_{ij}
    -\operatorname{ScaledFHT}_{p,c}(A)_{ij}|
  \le C^{-0}_{p,ij}(\widehat A,E).
\]
The same inequalities package as computed-matrix certificates.  Applying them
to the diagonal sign matrix gives computed preconditioners
\[
  \widehat\Pi^{\times1}_{\omega}
  \quad\hbox{and}\quad
  \widehat\Pi^{-0}_{\omega}
\]
for the ideal signed-Hadamard matrix
\[
  \Pi_\omega
  =
  H_p\,D_{\omega},
  \qquad
  (H_p)_{ik}=c\,S_p(i,k).
\]
The sign vector may be supplied exactly or through the existing rounded
\(\fl_{\rm mul}({\rm sign}_i,1)\) sign-storage certificate.

Finally, with the exact Rademacher law for \(\omega\) and the exact uniform-row
law for the row trace, the computed-left perturbation event for the resulting
preconditioned basis has probability one:
\[
  \Pr_{\omega,\rho}
  \bigl[
    \mathcal E_{\rm fp}(
      \widehat\Pi^{\times1}_{\omega},U,\rho)
  \bigr]=1,\qquad
  \Pr_{\omega,\rho}
  \bigl[
    \mathcal E_{\rm fp}(
      \widehat\Pi^{-0}_{\omega},U,\rho)
  \bigr]=1,
\]
and the same probability-one statements hold for the stored-sign variants.
These are implementation-facing non-probability arithmetic certificates:
butterfly arithmetic, writeback/copy, rounded normalization, computed
preconditioner formation, \(\widehat V\) formation, row scaling, and Gram dot
products are charged.  The Rademacher sign law and the uniform row-sampling law
remain exact mathematical inputs.  Array memory order, aliasing, in-place
overwrite, or vectorized routines with different semantics still require
separate routine certificates.
\end{corollary}

\begin{formalized}
\leanfit{flFhtPairUpdateStoredMulOne}
\leanfit{fhtPairUpdateStoredMulOnePropagatedErrorBudget}
\leanfit{fhtPairUpdateStoredMulOnePropagatedErrorBudget_nonneg}
\leanfit{flFhtPairUpdateStoredMulOne_propagated_error_bound}
\leanfit{flFhtPairUpdateStoredSubZeroRight}
\leanfit{fhtPairUpdateStoredSubZeroRightPropagatedErrorBudget}
\leanfit{fhtPairUpdateStoredSubZeroRightPropagatedErrorBudget_nonneg}
\leanfit{flFhtPairUpdateStoredSubZeroRight_propagated_error_bound}
\leanfit{flFhtPairScheduleStoredMulOne}
\leanfit{fhtPairScheduleStoredMulOnePropagatedErrorBudget}
\leanfit{fhtPairScheduleStoredMulOnePropagatedErrorBudget_nonneg}
\leanfit{flFhtPairScheduleStoredMulOne_propagated_error_bound}
\leanfit{flFhtPairScheduleStoredSubZeroRight}
\leanfit{fhtPairScheduleStoredSubZeroRightPropagatedErrorBudget}
\leanfit{fhtPairScheduleStoredSubZeroRightPropagatedErrorBudget_nonneg}
\leanfit{flFhtPairScheduleStoredSubZeroRight_propagated_error_bound}
\leanfit{flFhtScaledPairScheduleStoredMulOne_sqrtInvNatScale_error_bound}
\leanfit{flFhtScaledPairScheduleStoredSubZeroRight_sqrtInvNatScale_error_bound}
\leanfit{flFhtScaledPairScheduleMatrixStoredMulOne_sqrtInvNatScale_error_bound}
\leanfit{flFhtScaledPairScheduleMatrixStoredSubZeroRight_sqrtInvNatScale_error_bound}
\leanfit{flFhtScaledSylvesterScheduleMatrixStoredMulOne_sqrtInvNatScale_error_bound}
\leanfit{flFhtScaledSylvesterScheduleMatrixStoredSubZeroRight_sqrtInvNatScale_error_bound}
\leanfit{ComputedMatrix.flScaledFhtSylvesterScheduleColumnsSqrtInvNatStoredMulOne_entry_error_bound}
\leanfit{ComputedMatrix.flScaledFhtSylvesterScheduleColumnsSqrtInvNatStoredSubZeroRight_entry_error_bound}
\leanfit{ComputedPreconditioner.flSignedHadamardSylvesterFhtScheduleStoredMulOne_entry_error_bound}
\leanfit{ComputedPreconditioner.flSignedHadamardSylvesterFhtScheduleStoredSubZeroRight_entry_error_bound}
\leanfit{signedHadamardSylvesterFhtScheduleStoredMulOnePreconditioner}
\leanfit{signedHadamardSylvesterFhtScheduleStoredSubZeroRightPreconditioner}
\leanfit{signedHadamardSylvesterFhtScheduleStoredSignStoredMulOnePreconditioner}
\leanfit{signedHadamardSylvesterFhtScheduleStoredSignStoredSubZeroRightPreconditioner}
\noindent{\scriptsize\nolinkurl{signedHadamardUniformRowTraceProbability_eventProb_}}\par
\noindent{\scriptsize\quad\nolinkurl{sylvesterFhtScheduleStoredMulOneComputedLeftPreconditioned_}}\par
\noindent{\scriptsize\quad\nolinkurl{fl_uniformRowPerturbEvent_eq_one}}\par
\noindent{\scriptsize\nolinkurl{signedHadamardUniformRowTraceProbability_eventProb_}}\par
\noindent{\scriptsize\quad\nolinkurl{sylvesterFhtScheduleStoredSubZeroRightComputedLeftPreconditioned_}}\par
\noindent{\scriptsize\quad\nolinkurl{fl_uniformRowPerturbEvent_eq_one}}\par
\noindent{\scriptsize\nolinkurl{signedHadamardUniformRowTraceProbability_eventProb_}}\par
\noindent{\scriptsize\quad\nolinkurl{sylvesterFhtScheduleStoredSignStoredMulOneComputedLeftPreconditioned_}}\par
\noindent{\scriptsize\quad\nolinkurl{fl_uniformRowPerturbEvent_eq_one}}\par
\noindent{\scriptsize\nolinkurl{signedHadamardUniformRowTraceProbability_eventProb_}}\par
\noindent{\scriptsize\quad\nolinkurl{sylvesterFhtScheduleStoredSignStoredSubZeroRightComputedLeftPreconditioned_}}\par
\noindent{\scriptsize\quad\nolinkurl{fl_uniformRowPerturbEvent_eq_one}}\par
\noindent{\scriptsize\nolinkurl{signedHadamardUniformRowTraceProbability_eventProb_}}\par
\noindent{\scriptsize\quad\nolinkurl{sylvesterFhtScheduleStoredSignStoredMulOneComputedLeftPreconditioned_}}\par
\noindent{\scriptsize\quad\nolinkurl{factoredInput_fl_uniformRowSampleGramDotWithFlSqrtMulInvSqrtDen_}}\par
\noindent{\scriptsize\quad\nolinkurl{two_sided_finiteLoewnerLe_ge_one_sub_delta_srht_log_preprocess}}\par
\noindent{\scriptsize\nolinkurl{signedHadamardUniformRowTraceProbability_eventProb_}}\par
\noindent{\scriptsize\quad\nolinkurl{sylvesterFhtScheduleStoredSignStoredSubZeroRightComputedLeftPreconditioned_}}\par
\noindent{\scriptsize\quad\nolinkurl{factoredInput_fl_uniformRowSampleGramDotWithFlSqrtMulInvSqrtDen_}}\par
\noindent{\scriptsize\quad\nolinkurl{two_sided_finiteLoewnerLe_ge_one_sub_delta_srht_log_preprocess}}\par
\end{formalized}

\begin{corollary}[Modified-coordinate add/multiply/subtract writeback for fast generated-FHT]
Use the one-pair notation above and the exact probability laws already stated
for Algorithm 3.  Let
\(\mathcal P_p\) be the generated Sylvester/Walsh pair schedule in dimension
\(2^p\).  For one pair \((r,s)\), let
\[
  y^{\rm pair}=\operatorname{FHTPairUpdate}_{r,s}(x),\qquad
  \widehat y^{\rm pair}
    =\operatorname{flFHTPairUpdate}_{r,s}(\widehat x).
\]
For \(\tau\in\{+,\times 1,-0\}\), define the rounded copy map
\[
  \psi_+(a)=\fl_{\rm add}(a,0),\qquad
  \psi_{\times 1}(a)=\fl_{\rm mul}(a,1),\qquad
  \psi_{-0}(a)=\fl_{\rm sub}(a,0),
\]
and define the modified-coordinate stored output by
\[
  \widehat y^{\rm mod,\tau}_i
  =
  \begin{cases}
    \psi_\tau(\widehat y^{\rm pair}_i),
      & i=r\ {\rm or}\ i=s,\\
    \widehat y^{\rm pair}_i,
      & i\ne r\ {\rm and}\ i\ne s.
  \end{cases}
\]
If the input certificate satisfies
\[
  |\widehat x_i-x_i|\le E_i,\qquad E_i\ge0,
\]
then, for each \(\tau\in\{+,\times 1,-0\}\), the modified-coordinate stored
pair update satisfies
\[
  |\widehat y^{\rm mod,\tau}_i-y^{\rm pair}_i|
  \le
  B^{\rm pair}_{r,s,i}(\widehat x,E)
  +{\bf 1}_{\{i=r\ {\rm or}\ i=s\}}\,
    u\,|\widehat y^{\rm pair}_i|.
\]
Equivalently,
\[
  E^{\rm mod,\tau}_{r,s,i}
  =
  B^{\rm pair}_{r,s,i}(\widehat x,E)
  +
  \begin{cases}
    u\,|\operatorname{flFHTPairUpdate}_{r,s}(\widehat x)_i|,
      & i=r\ {\rm or}\ i=s,\\
    0,& \text{otherwise}
  \end{cases}
\]
is a valid propagated radius for one rounded pair update followed by a
writeback only of the modified butterfly outputs.  The copy map \(\psi_\tau\)
is the only part changed across the three variants.

Composing this recurrence over \(\mathcal P_p\), let
\[
  \widehat z^\tau_i
  =
  \operatorname{flFHTSchedule}^{\rm mod,\tau}_{\mathcal P_p}
  (\widehat x)_i,\qquad
  E^{\rm mod,\tau}_{p,i}
  =
  \operatorname{FHTScheduleBudget}^{\rm mod,\tau}_{p,i}(\widehat x,E).
\]
For the scaled transform, let \(c=\sqrt{((2^p)^{-1})}\),
\(\widehat c=\operatorname{flsqrt}((2^p)^{-1})\), and \(\eta=c\,u\).  Then
for each \(\tau\),
\[
  \left|
    \fl_{\rm mul}(\widehat c,\widehat z^\tau_i)
    -
    c\,\operatorname{FHTSchedule}_{\mathcal P_p}(x)_i
  \right|
  \le
  u|\widehat c\,\widehat z^\tau_i|
  +|\widehat c|\,E^{\rm mod,\tau}_{p,i}
  +\eta\bigl(|\widehat z^\tau_i|+E^{\rm mod,\tau}_{p,i}\bigr).
\]
Columnwise, this gives computed matrix certificates
\[
  \left|
    \operatorname{flFHTSylvesterMatrix}^{\rm mod,\tau}_{\widehat c}(\widehat A)_{ij}
    -
    \operatorname{FHTSylvesterMatrix}_{c}(A)_{ij}
  \right|
  \le
  B^{\rm mod,\tau}_{ij},
\]
where
\[
  B^{\rm mod,\tau}
  =
  \operatorname{FHTSylvesterMatrixBudget}^{\rm mod,\tau}_{p,\widehat c,\eta}
  (\widehat A,E).
\]

With computed diagonal sign input \(\widehat D\) and sign-storage certificate
\(E^D\), the implemented preconditioners
\[
  \widehat P^{\rm mod,\tau}
  =
  \operatorname{flFHTSylvesterMatrix}^{\rm mod,\tau}_{\widehat c}(\widehat D)
\]
satisfy, for every \(i,j\),
\[
  |\widehat P^{\rm mod,\tau}_{ij}
    -((cS_p)\operatorname{diag}(\sigma))_{ij}|
  \le
  B^{\rm mod,\tau}_{ij}.
\]
Thus these are fully computed-object implementation-facing certificates for
the modeled routines ``rounded butterfly, then
\(\fl_{\rm add}(\cdot,0)\), \(\fl_{\rm mul}(\cdot,1)\), or
\(\fl_{\rm sub}(\cdot,0)\) writeback only on the two modified pair outputs.''
The Rademacher and uniform-row sampling laws remain exact mathematical inputs.
Under their exact product law, for each \(\tau\), the computed-left perturbation
event for
\[
  \widehat V^\tau=\fl(\widehat P^{\rm mod,\tau}U)
\]
has probability one, with the downstream radius additionally charging
\(\widehat V^\tau\)-formation, row scaling, and Gram-dot-product errors.  If
the realized signs are also stored by the rounded sign-copy certificate, the
same probability-one event holds with that sign-storage radius included in
\(E^D\).  Memory order, aliasing, vectorized layout, or in-place overwrite
semantics not represented by these modified-coordinate writeback models remain
separate implementation obligations.
\end{corollary}

\begin{formalized}
\leanfit{flFhtPairUpdateModifiedStoredAddZeroRight}
\leanfit{fhtPairUpdateModifiedStoredAddZeroRightPropagatedErrorBudget}
\leanfit{flFhtPairUpdateModifiedStoredAddZeroRight_propagated_error_bound}
\leanfit{flFhtPairUpdateModifiedStoredMulOne}
\leanfit{fhtPairUpdateModifiedStoredMulOnePropagatedErrorBudget}
\leanfit{flFhtPairUpdateModifiedStoredMulOne_propagated_error_bound}
\leanfit{flFhtPairUpdateModifiedStoredSubZeroRight}
\leanfit{fhtPairUpdateModifiedStoredSubZeroRightPropagatedErrorBudget}
\leanfit{flFhtPairUpdateModifiedStoredSubZeroRight_propagated_error_bound}
\leanfit{flFhtPairScheduleModifiedStoredAddZeroRight}
\leanfit{fhtPairScheduleModifiedStoredAddZeroRightPropagatedErrorBudget}
\leanfit{flFhtPairScheduleModifiedStoredAddZeroRight_propagated_error_bound}
\leanfit{flFhtPairScheduleModifiedStoredMulOne}
\leanfit{fhtPairScheduleModifiedStoredMulOnePropagatedErrorBudget}
\leanfit{flFhtPairScheduleModifiedStoredMulOne_propagated_error_bound}
\leanfit{flFhtPairScheduleModifiedStoredSubZeroRight}
\leanfit{fhtPairScheduleModifiedStoredSubZeroRightPropagatedErrorBudget}
\leanfit{flFhtPairScheduleModifiedStoredSubZeroRight_propagated_error_bound}
\leanfit{flFhtSylvesterScheduleModifiedStoredAddZeroRight}
\leanfit{flFhtSylvesterScheduleModifiedStoredAddZeroRight_propagated_error_bound}
\leanfit{flFhtSylvesterScheduleModifiedStoredMulOne}
\leanfit{flFhtSylvesterScheduleModifiedStoredMulOne_propagated_error_bound}
\leanfit{flFhtSylvesterScheduleModifiedStoredSubZeroRight}
\leanfit{flFhtSylvesterScheduleModifiedStoredSubZeroRight_propagated_error_bound}
\leanfit{flFhtScaledPairScheduleModifiedStoredAddZeroRight}
\leanfit{flFhtScaledPairScheduleModifiedStoredAddZeroRight_sqrtInvNatScale_error_bound}
\leanfit{flFhtScaledPairScheduleModifiedStoredMulOne}
\leanfit{flFhtScaledPairScheduleModifiedStoredMulOne_sqrtInvNatScale_error_bound}
\leanfit{flFhtScaledPairScheduleModifiedStoredSubZeroRight}
\leanfit{flFhtScaledPairScheduleModifiedStoredSubZeroRight_sqrtInvNatScale_error_bound}
\leanfit{flFhtScaledPairScheduleMatrixModifiedStoredAddZeroRight}
\leanfit{flFhtScaledPairScheduleMatrixModifiedStoredAddZeroRight_sqrtInvNatScale_error_bound}
\leanfit{flFhtScaledPairScheduleMatrixModifiedStoredMulOne_sqrtInvNatScale_error_bound}
\leanfit{flFhtScaledPairScheduleMatrixModifiedStoredSubZeroRight_sqrtInvNatScale_error_bound}
\leanfit{flFhtScaledSylvesterScheduleMatrixModifiedStoredAddZeroRight}
\leanfit{flFhtScaledSylvesterScheduleMatrixModifiedStoredAddZeroRight_sqrtInvNatScale_error_bound}
\leanfit{flFhtScaledSylvesterScheduleMatrixModifiedStoredMulOne_sqrtInvNatScale_error_bound}
\leanfit{flFhtScaledSylvesterScheduleMatrixModifiedStoredSubZeroRight_sqrtInvNatScale_error_bound}
\leanfit{ComputedMatrix.flScaledFhtSylvesterScheduleColumnsSqrtInvNatModifiedStoredAddZeroRight}
\leanfit{ComputedMatrix.flScaledFhtSylvesterScheduleColumnsSqrtInvNatModifiedStoredAddZeroRight_entry_error_bound}
\leanfit{ComputedMatrix.flScaledFhtSylvesterScheduleColumnsSqrtInvNatModifiedStoredMulOne_entry_error_bound}
\leanfit{ComputedMatrix.flScaledFhtSylvesterScheduleColumnsSqrtInvNatModifiedStoredSubZeroRight_entry_error_bound}
\leanfit{ComputedPreconditioner.flSignedHadamardSylvesterFhtScheduleModifiedStoredAddZeroRight}
\leanfit{ComputedPreconditioner.flSignedHadamardSylvesterFhtScheduleModifiedStoredAddZeroRight_entry_error_bound}
\leanfit{ComputedPreconditioner.flSignedHadamardSylvesterFhtScheduleModifiedStoredMulOne_entry_error_bound}
\leanfit{ComputedPreconditioner.flSignedHadamardSylvesterFhtScheduleModifiedStoredSubZeroRight_entry_error_bound}
\leanfit{signedHadamardSylvesterFhtScheduleModifiedStoredAddZeroRightPreconditioner}
\leanfit{signedHadamardSylvesterFhtScheduleModifiedStoredMulOnePreconditioner}
\leanfit{signedHadamardSylvesterFhtScheduleModifiedStoredSubZeroRightPreconditioner}
\leanfit{signedHadamardSylvesterFhtScheduleStoredSignModifiedStoredAddZeroRightPreconditioner}
\leanfit{signedHadamardSylvesterFhtScheduleStoredSignModifiedStoredMulOnePreconditioner}
\leanfit{signedHadamardSylvesterFhtScheduleStoredSignModifiedStoredSubZeroRightPreconditioner}
\noindent{\scriptsize\nolinkurl{signedHadamardUniformRowTraceProbability_eventProb_}}\par
\noindent{\scriptsize\quad\nolinkurl{sylvesterFhtScheduleModifiedStoredAddZeroRightComputedLeftPreconditioned_}}\par
\noindent{\scriptsize\quad\nolinkurl{fl_uniformRowPerturbEvent_eq_one}}\par
\noindent{\scriptsize\nolinkurl{signedHadamardUniformRowTraceProbability_eventProb_}}\par
\noindent{\scriptsize\quad\nolinkurl{sylvesterFhtScheduleModifiedStoredMulOneComputedLeftPreconditioned_}}\par
\noindent{\scriptsize\quad\nolinkurl{fl_uniformRowPerturbEvent_eq_one}}\par
\noindent{\scriptsize\nolinkurl{signedHadamardUniformRowTraceProbability_eventProb_}}\par
\noindent{\scriptsize\quad\nolinkurl{sylvesterFhtScheduleModifiedStoredSubZeroRightComputedLeftPreconditioned_}}\par
\noindent{\scriptsize\quad\nolinkurl{fl_uniformRowPerturbEvent_eq_one}}\par
\noindent{\scriptsize\nolinkurl{signedHadamardUniformRowTraceProbability_eventProb_}}\par
\noindent{\scriptsize\quad\nolinkurl{sylvesterFhtScheduleStoredSignModifiedStoredAddZeroRightComputedLeftPreconditioned_}}\par
\noindent{\scriptsize\quad\nolinkurl{fl_uniformRowPerturbEvent_eq_one}}\par
\noindent{\scriptsize\nolinkurl{signedHadamardUniformRowTraceProbability_eventProb_}}\par
\noindent{\scriptsize\quad\nolinkurl{sylvesterFhtScheduleStoredSignModifiedStoredMulOneComputedLeftPreconditioned_}}\par
\noindent{\scriptsize\quad\nolinkurl{fl_uniformRowPerturbEvent_eq_one}}\par
\noindent{\scriptsize\nolinkurl{signedHadamardUniformRowTraceProbability_eventProb_}}\par
\noindent{\scriptsize\quad\nolinkurl{sylvesterFhtScheduleStoredSignModifiedStoredSubZeroRightComputedLeftPreconditioned_}}\par
\noindent{\scriptsize\quad\nolinkurl{fl_uniformRowPerturbEvent_eq_one}}\par
\noindent{\scriptsize\nolinkurl{signedHadamardUniformRowTraceProbability_eventProb_}}\par
\noindent{\scriptsize\quad\nolinkurl{sylvesterFhtScheduleStoredSignModifiedStoredMulOneComputedLeftPreconditioned_}}\par
\noindent{\scriptsize\quad\nolinkurl{factoredInput_fl_uniformRowSampleGramDotWithFlSqrtMulInvSqrtDen_}}\par
\noindent{\scriptsize\quad\nolinkurl{two_sided_finiteLoewnerLe_ge_one_sub_delta_srht_log_preprocess}}\par
\noindent{\scriptsize\nolinkurl{signedHadamardUniformRowTraceProbability_eventProb_}}\par
\noindent{\scriptsize\quad\nolinkurl{sylvesterFhtScheduleStoredSignModifiedStoredSubZeroRightComputedLeftPreconditioned_}}\par
\noindent{\scriptsize\quad\nolinkurl{factoredInput_fl_uniformRowSampleGramDotWithFlSqrtMulInvSqrtDen_}}\par
\noindent{\scriptsize\quad\nolinkurl{two_sided_finiteLoewnerLe_ge_one_sub_delta_srht_log_preprocess}}\par
\end{formalized}

\begin{corollary}[Computed-projector replacement for Theorem 183]
Let \(Q\in\mathbb R^{m\times d}\) be the exact analysis basis or singular-vector
table used to define the projection matrix
\[
  \Pi=QQ^T,\qquad
  \Pi_{ij}=\sum_a Q_{ia}Q_{ja}.
\]
Suppose the implementation computes \(\widehat Q\) with certificate
\[
  |\widehat Q_{ia}-Q_{ia}|\le E^Q_{ia},\qquad E^Q_{ia}\ge0,
\]
and forms the projection matrix by the rounded product
\[
  \widehat\Pi=\fl(\widehat Q\widehat Q^T).
\]
If \(\gamma_d\) is valid, then for every \(i,j\),
\[
\begin{aligned}
  |\widehat\Pi_{ij}-\Pi_{ij}|
  \le\;&
  \gamma_d\sum_a |\widehat Q_{ia}|\,|\widehat Q_{ja}|\\
  &+\sum_a E^Q_{ia}|\widehat Q_{ja}|
  +\sum_a |Q_{ia}|E^Q_{ja}
  =: E^\Pi_{ij}.
\end{aligned}
\]
Consequently the one-sided Algorithm~3 preprocessing bounds become
\[
  |\fl(\widehat\Pi A)_{ij}-(\Pi A)_{ij}|
  \le
  \gamma_m\sum_k|\widehat\Pi_{ik}|\,|A_{kj}|
  +\sum_k E^\Pi_{ik}|A_{kj}|,
\]
and analogously on the right,
\[
  |\fl(A\widehat\Pi)_{ij}-(A\Pi)_{ij}|
  \le
  \gamma_m\sum_k|A_{ik}|\,|\widehat\Pi_{kj}|
  +\sum_k |A_{ik}|E^\Pi_{kj}.
\]
Thus a computed SVD/QR/singular-vector routine enters through the
\(E^Q\) certificate and the rounded formation of \(QQ^T\), not through a hidden
exact projector.
\end{corollary}

\begin{formalized}
\leanfit{basisColumnProjector}
\leanfit{fl_basisColumnProjector}
\leanfit{fl_basisColumnProjector_total_error_bound}
\leanfit{flBasisColumnProjectorEntryErrorBudget}
\leanfit{fl_basisColumnProjector_entry_error_budget_bound}
\leanfit{ComputedPreconditioner.flBasisColumnProjector}
\leanfit{fl_preconditionRowsWithComputedLeft_total_error_bound}
\leanfit{fl_preconditionColumnsWithComputedRight_total_error_bound}
\end{formalized}

\begin{corollary}[Certified QR/SVD basis routine handoff for Theorem 183]
Let a concrete QR, SVD, singular-vector, or orthonormal-basis routine return
\(\widehat Q\in\mathbb R^{m\times d}\) together with an entrywise certificate
against the exact analysis basis \(Q\),
\[
  |\widehat Q_{ia}-Q_{ia}|\le E^Q_{ia},
  \qquad E^Q_{ia}\ge0 .
\]
Define
\[
  \widehat\Pi=\fl(\widehat Q\widehat Q^T),
  \qquad
  \Pi=QQ^T,
\]
and
\[
\begin{aligned}
  E^\Pi_{ij}
  =
  &\ \gamma_d\sum_a |\widehat Q_{ia}|\,|\widehat Q_{ja}|\\
  &+\sum_a E^Q_{ia}|\widehat Q_{ja}|
  +\sum_a |Q_{ia}|E^Q_{ja}.
\end{aligned}
\]
If the length-\(d\) dot products used to form
\(\widehat Q\widehat Q^T\) are \(\gamma\)-valid, then
\[
  |\widehat\Pi_{ij}-\Pi_{ij}|\le E^\Pi_{ij}
  \qquad\text{for all }i,j .
\]
Thus a routine-produced basis table feeds Algorithm~3 through the explicit
certificate \(E^Q\), not by replacing \(\widehat Q\) with the exact analysis
basis.

For the two-sided projector branch, suppose two routines return
\(\widehat Q_L,\widehat Q_R\) with entrywise radii
\(E^{Q_L},E^{Q_R}\), and define the corresponding projector budgets
\(E^{\Pi_L},E^{\Pi_R}\) by the same formula.  If the matrix-product
dimensions are \(\gamma\)-valid, then for every \(i,j\),
\[
\begin{aligned}
&|\fl(\widehat\Pi_L A\widehat\Pi_R)_{ij}
  -(\Pi_LA\Pi_R)_{ij}|\\
&\quad\le
\gamma_n\sum_t
  |\fl(\widehat\Pi_LA)_{it}|\,|\widehat\Pi_{R,tj}|\\
&\qquad+
\sum_t
  \left(\gamma_m\sum_a|\widehat\Pi_{L,ia}|\,|A_{at}|\right)
  |\widehat\Pi_{R,tj}|\\
&\qquad+
\sum_t
  \left(\sum_a E^{\Pi_L}_{ia}|A_{at}|\right)
  |\widehat\Pi_{R,tj}|
  +
  \sum_t |(\Pi_LA)_{it}|E^{\Pi_R}_{tj}.
\end{aligned}
\]
The displayed bound charges basis-routine error, projector formation, and the
two rounded preprocessing products.  Sampling probabilities and sampling laws
remain exact mathematical inputs.
\end{corollary}

\begin{formalized}
\leanfit{ComputedMatrix.ofEntrywiseBound}
\leanfit{certifiedBasisProjectorEntryErrorBudget}
\leanfit{certifiedBasisProjectorEntryErrorBudget_nonneg}
\leanfit{fl_basisColumnProjector_of_certifiedBasis_entry_error_bound}
\leanfit{ComputedPreconditioner.flBasisColumnProjectorOfCertifiedBasis}
\leanfit{ComputedPreconditioner.flBasisColumnProjectorOfCertifiedBasis_entry_error_bound}
\leanfit{fl_preconditionElementsWithCertifiedBasisProjectors_total_error_bound}
\end{formalized}

\begin{corollary}[Right-orthogonal QR/SVD reference handoff for Theorem 183]
Let \(Q\in\mathbb R^{m\times d}\) be the exact analysis basis and let
\(O\in\mathbb R^{d\times d}\) be exact and orthogonal.  Define the rotated
reference table
\[
  Q^O=QO .
\]
Then
\[
  Q^O(Q^O)^T=QQ^T .
\]
Consequently, if a concrete QR, SVD, singular-vector, or orthonormal-basis
routine returns \(\widehat Q\) with a certificate against the rotated
reference,
\[
  |\widehat Q_{ia}-Q^O_{ia}|\le E^Q_{ia},
  \qquad E^Q_{ia}\ge0 ,
\]
and
\[
  \widehat\Pi=\fl(\widehat Q\widehat Q^T),
  \qquad
  \Pi=QQ^T,
\]
then the budget
\[
\begin{aligned}
  E^{\Pi,O}_{ij}
  =
  &\ \gamma_d\sum_a |\widehat Q_{ia}|\,|\widehat Q_{ja}|\\
  &+\sum_a E^Q_{ia}|\widehat Q_{ja}|
  +\sum_a |Q^O_{ia}|E^Q_{ja}
\end{aligned}
\]
satisfies
\[
  |\widehat\Pi_{ij}-\Pi_{ij}|\le E^{\Pi,O}_{ij}
  \qquad\text{for all }i,j ,
\]
provided the length-\(d\) projector dot products are \(\gamma\)-valid.

The same correction applies to generated-then-stored basis tables.  Suppose
the routine returns \(Q^{\rm raw}\) with
\[
  |Q^{\rm raw}_{ia}-Q^O_{ia}|\le E^Q_{ia},
\]
and the algorithm stores the table actually used in projector formation as
\[
  |Q^{\rm store}_{ia}-Q^{\rm raw}_{ia}|\le C^Q_{ia},
  \qquad C^Q_{ia}\ge0 .
\]
With
\[
  \widehat\Pi=\fl(Q^{\rm store}(Q^{\rm store})^T),
  \qquad
  \Pi=QQ^T,
\]
the stored-table budget
\[
\begin{aligned}
  E^{\Pi,O/{\rm store}}_{ij}
  =
  &\ \gamma_d\sum_a |Q^{\rm store}_{ia}|\,|Q^{\rm store}_{ja}|\\
  &+\sum_a (C^Q_{ia}+E^Q_{ia})|Q^{\rm store}_{ja}|\\
  &+\sum_a |Q^O_{ia}|(C^Q_{ja}+E^Q_{ja})
\end{aligned}
\]
satisfies
\[
  |\widehat\Pi_{ij}-\Pi_{ij}|
  \le E^{\Pi,O/{\rm store}}_{ij}
  \qquad\text{for all }i,j .
\]
Thus QR/SVD sign choices or right-orthogonal basis rotations are handled as an
exact projector equivalence, while all computed non-probability quantities
\(\widehat Q,Q^{\rm raw},Q^{\rm store}\) remain explicitly certified.
Sampling probabilities and sampling laws remain exact mathematical inputs.
\end{corollary}

\begin{formalized}
\leanfit{basisColumnProjector_matMulRectRight_orthogonal}
\leanfit{fl_basisColumnProjector_entry_error_budget_bound_rightOrthogonalReference}
\leanfit{fl_basisColumnProjector_of_rightOrthogonalCertifiedBasis_entry_error_bound}
\leanfit{fl_basisColumnProjector_of_rightOrthogonalCertifiedStoredBasis_entry_error_bound}
\end{formalized}

\begin{corollary}[Generated-then-stored QR/SVD basis handoff for Theorem 183]
Let a concrete QR, SVD, singular-vector, or orthonormal-basis routine first
return a raw table \(Q^{\rm raw}\in\mathbb R^{m\times d}\) together with an
entrywise generation certificate against the exact analysis basis \(Q\),
\[
  |Q^{\rm raw}_{ia}-Q_{ia}|\le E^Q_{ia},
  \qquad E^Q_{ia}\ge0 .
\]
Suppose the implemented algorithm does not use \(Q^{\rm raw}\) directly, but
stores or copies it as the table \(Q^{\rm store}\) used in the projector
formation, with a separate storage certificate
\[
  |Q^{\rm store}_{ia}-Q^{\rm raw}_{ia}|\le C^Q_{ia},
  \qquad C^Q_{ia}\ge0 .
\]
Define
\[
  \widehat\Pi=\fl(Q^{\rm store}(Q^{\rm store})^T),
  \qquad
  \Pi=QQ^T,
\]
and
\[
\begin{aligned}
  E^{\Pi,\rm store}_{ij}
  =
  &\ \gamma_d\sum_a
      |Q^{\rm store}_{ia}|\,|Q^{\rm store}_{ja}|\\
  &+\sum_a (C^Q_{ia}+E^Q_{ia})|Q^{\rm store}_{ja}|\\
  &+\sum_a |Q_{ia}|(C^Q_{ja}+E^Q_{ja}) .
\end{aligned}
\]
If the length-\(d\) dot products used to form
\(Q^{\rm store}(Q^{\rm store})^T\) are \(\gamma\)-valid, then
\[
  |\widehat\Pi_{ij}-\Pi_{ij}|
  \le E^{\Pi,\rm store}_{ij}
  \qquad\text{for all }i,j .
\]
For the concrete rounded-copy routines
\[
  Q^{\rm store}_{ia}=\operatorname{fl\_mul}(Q^{\rm raw}_{ia},1),\qquad
  Q^{\rm store}_{ia}=\operatorname{fl\_add}(Q^{\rm raw}_{ia},0),\qquad
  Q^{\rm store}_{ia}=\operatorname{fl\_sub}(Q^{\rm raw}_{ia},0),
\]
the formalized storage certificate instantiates
\[
  C^Q_{ia}=u\,|Q^{\rm raw}_{ia}|.
\]

For the two-sided projector branch, suppose left and right routines provide
\((Q_L^{\rm raw},Q_L^{\rm store},E^{Q_L},C^{Q_L})\) and
\((Q_R^{\rm raw},Q_R^{\rm store},E^{Q_R},C^{Q_R})\) satisfying the same
generation and storage certificates, and define
\(E^{\Pi_L,\rm store}\) and \(E^{\Pi_R,\rm store}\) by the displayed formula.
If the matrix-product dimensions are \(\gamma\)-valid, then for every \(i,j\),
\[
\begin{aligned}
&|\fl(\widehat\Pi_L A\widehat\Pi_R)_{ij}
  -(\Pi_LA\Pi_R)_{ij}|\\
&\quad\le
\gamma_n\sum_t
  |\fl(\widehat\Pi_LA)_{it}|\,|\widehat\Pi_{R,tj}|\\
&\qquad+
\sum_t
  \left(\gamma_m\sum_a|\widehat\Pi_{L,ia}|\,|A_{at}|\right)
  |\widehat\Pi_{R,tj}|\\
&\qquad+
\sum_t
  \left(\sum_a E^{\Pi_L,\rm store}_{ia}|A_{at}|\right)
  |\widehat\Pi_{R,tj}|
  +
  \sum_t |(\Pi_LA)_{it}|E^{\Pi_R,\rm store}_{tj}.
\end{aligned}
\]
The result charges the raw basis-generation error, the actual table-storage
error, floating-point projector formation, and the two rounded preprocessing
products.  Sampling probabilities and sampling laws remain exact mathematical
inputs.
\end{corollary}

\begin{formalized}
\leanfit{ComputedMatrix.ofEntrywiseBoundThenStorage}
\leanfit{ComputedMatrix.ofEntrywiseBoundStoredMulOne}
\leanfit{ComputedMatrix.ofEntrywiseBoundStoredAddZeroRight}
\leanfit{ComputedMatrix.ofEntrywiseBoundStoredSubZeroRight}
\leanfit{certifiedStoredBasisProjectorEntryErrorBudget}
\leanfit{certifiedStoredBasisProjectorEntryErrorBudget_nonneg}
\leanfit{fl_basisColumnProjector_of_certifiedStoredBasis_entry_error_bound}
\leanfit{ComputedPreconditioner.flBasisColumnProjectorOfCertifiedStoredBasis}
\leanfit{ComputedPreconditioner.flBasisColumnProjectorOfCertifiedStoredBasis_entry_error_bound}
\leanfit{fl_preconditionElementsWithCertifiedStoredBasisProjectors_total_error_bound}
\end{formalized}

\begin{corollary}[Normwise generated-then-stored QR/SVD basis handoff for Theorem 183]
Let a concrete QR, SVD, singular-vector, or orthonormal-basis routine first
return a raw table \(Q^{\rm raw}\in\mathbb R^{m\times d}\).  Suppose the
algorithm then stores or copies this table as \(Q^{\rm store}\), and the
projector is formed from \(Q^{\rm store}\), with storage certificate
\[
  |Q^{\rm store}_{ia}-Q^{\rm raw}_{ia}|\le C^Q_{ia},
  \qquad C^Q_{ia}\ge0 .
\]
Assume that the raw routine supplies one of the following generation
certificates against the exact analysis basis \(Q\):
\[
\begin{array}{ll}
\text{\rm Frobenius:}
  & \|Q^{\rm raw}-Q\|_F\le \eta,\qquad \eta\ge0,\\[2mm]
\text{\rm columnwise:}
  & \|Q^{\rm raw}_{:a}-Q_{:a}\|_2\le \eta_a,\qquad \eta_a\ge0,\\[2mm]
\text{\rm operator:}
  & \|(Q^{\rm raw}-Q)x\|_2\le \eta\|x\|_2
     \quad\hbox{for all }x,\qquad \eta\ge0 .
\end{array}
\]
Define the raw entry radius
\[
  \rho_{ia} =
  \begin{cases}
    \eta, & \text{in the Frobenius case},\\
    \eta_a, & \text{in the columnwise case},\\
    \eta, & \text{in the operator case}.
  \end{cases}
\]
The formalized constructors prove
\[
  |Q^{\rm store}_{ia}-Q_{ia}|\le C^Q_{ia}+\rho_{ia}.
\]
Consequently, with
\[
  \widehat\Pi=\fl(Q^{\rm store}(Q^{\rm store})^T),
  \qquad
  \Pi=QQ^T,
\]
and
\[
\begin{aligned}
  E^{\Pi,\rm norm/store}_{ij}
  =
  &\ \gamma_d\sum_a
      |Q^{\rm store}_{ia}|\,|Q^{\rm store}_{ja}|\\
  &+\sum_a (C^Q_{ia}+\rho_{ia})|Q^{\rm store}_{ja}|\\
  &+\sum_a |Q_{ia}|(C^Q_{ja}+\rho_{ja}),
\end{aligned}
\]
if the length-\(d\) dot products used to form
\(Q^{\rm store}(Q^{\rm store})^T\) are \(\gamma\)-valid, then
\[
  |\widehat\Pi_{ij}-\Pi_{ij}|
  \le E^{\Pi,\rm norm/store}_{ij}
  \qquad\text{for all }i,j .
\]
The concrete rounded-copy paths from the previous corollary instantiate
\(C^Q_{ia}=u|Q^{\rm raw}_{ia}|\) for
\(\operatorname{fl\_mul}(Q^{\rm raw}_{ia},1)\),
\(\operatorname{fl\_add}(Q^{\rm raw}_{ia},0)\), or
\(\operatorname{fl\_sub}(Q^{\rm raw}_{ia},0)\).

For the two-sided projector branch, suppose left and right routines provide
raw/stored tables with storage radii \(C^{Q_L},C^{Q_R}\) and raw normwise
radii \(\rho^L,\rho^R\), where each \(\rho\) is obtained from a Frobenius,
columnwise, or operator certificate as above.  Let
\(E^{\Pi_L,\rm norm/store}\) and \(E^{\Pi_R,\rm norm/store}\) be the
corresponding displayed projector budgets.  If the matrix-product dimensions
are \(\gamma\)-valid, then for every \(i,j\),
\[
\begin{aligned}
&|\fl(\widehat\Pi_L A\widehat\Pi_R)_{ij}
  -(\Pi_LA\Pi_R)_{ij}|\\
&\quad\le
\gamma_n\sum_t
  |\fl(\widehat\Pi_LA)_{it}|\,|\widehat\Pi_{R,tj}|\\
&\qquad+
\sum_t
  \left(\gamma_m\sum_a|\widehat\Pi_{L,ia}|\,|A_{at}|\right)
  |\widehat\Pi_{R,tj}|\\
&\qquad+
\sum_t
  \left(\sum_a E^{\Pi_L,\rm norm/store}_{ia}|A_{at}|\right)
  |\widehat\Pi_{R,tj}|
  +
  \sum_t |(\Pi_LA)_{it}|E^{\Pi_R,\rm norm/store}_{tj}.
\end{aligned}
\]
Thus the two-sided implementation-facing theorem charges the raw normwise
basis-generation certificates, the actual table-storage errors, floating-point
projector formation, and the two rounded preprocessing products.  Sampling
probabilities and sampling laws remain exact mathematical inputs.
\end{corollary}

\begin{formalized}
\leanfit{ComputedMatrix.ofFrobeniusBoundThenStorage}
\leanfit{ComputedMatrix.ofColumnVecNorm2BoundThenStorage}
\leanfit{ComputedMatrix.ofRectOpNorm2BoundThenStorage}
\leanfit{fl_basisColumnProjector_of_frobeniusCertifiedStoredBasis_entry_error_bound}
\leanfit{fl_basisColumnProjector_of_columnwiseCertifiedStoredBasis_entry_error_bound}
\leanfit{fl_basisColumnProjector_of_opNormCertifiedStoredBasis_entry_error_bound}
\leanfit{ComputedPreconditioner.flBasisColumnProjectorOfFrobeniusCertifiedStoredBasis}
\leanfit{ComputedPreconditioner.flBasisColumnProjectorOfColumnwiseCertifiedStoredBasis}
\leanfit{ComputedPreconditioner.flBasisColumnProjectorOfOpNormCertifiedStoredBasis}
\leanfit{fl_preconditionElementsWithFrobeniusCertifiedStoredBasisProjectors_total_error_bound}
\leanfit{fl_preconditionElementsWithColumnwiseCertifiedStoredBasisProjectors_total_error_bound}
\leanfit{fl_preconditionElementsWithOpNormCertifiedStoredBasisProjectors_total_error_bound}
\end{formalized}

\begin{corollary}[Normwise storage certificates for generated QR/SVD basis handoffs]
Let \(Q^{\rm raw}\) be the table produced by a QR, SVD, singular-vector, or
orthonormal-basis routine, and let \(Q^{\rm store}\) be the table actually
stored and used to form the projector.  Suppose the routine and storage path
provide one of the following certificate pairs against the exact analysis
basis \(Q\):
\[
\begin{array}{ll}
\text{\rm entrywise/Frobenius storage:}
  & |Q^{\rm raw}_{ia}-Q_{ia}|\le E_{ia},\quad
    \|Q^{\rm store}-Q^{\rm raw}\|_F\le\sigma,\\[1mm]
\text{\rm Frobenius/Frobenius:}
  & \|Q^{\rm raw}-Q\|_F\le\eta,\quad
    \|Q^{\rm store}-Q^{\rm raw}\|_F\le\sigma,\\[1mm]
\text{\rm columnwise/columnwise:}
  & \|Q^{\rm raw}_{:a}-Q_{:a}\|_2\le\eta_a,\quad
    \|Q^{\rm store}_{:a}-Q^{\rm raw}_{:a}\|_2\le\sigma_a,\\[1mm]
\text{\rm operator/operator:}
  & \|(Q^{\rm raw}-Q)x\|_2\le\eta\|x\|_2,\quad
    \|(Q^{\rm store}-Q^{\rm raw})x\|_2\le\sigma\|x\|_2 .
\end{array}
\]
Assume all displayed radii are nonnegative.  Define entry radii
\[
(\rho_{ia},\tau_{ia}) =
\begin{cases}
  (E_{ia},\sigma), & \text{entrywise/Frobenius storage},\\
  (\eta,\sigma), & \text{Frobenius/Frobenius},\\
  (\eta_a,\sigma_a), & \text{columnwise/columnwise},\\
  (\eta,\sigma), & \text{operator/operator}.
\end{cases}
\]
The formalized constructors prove
\[
  |Q^{\rm store}_{ia}-Q_{ia}|\le \tau_{ia}+\rho_{ia}.
\]
Consequently, with
\[
  \widehat\Pi=\fl(Q^{\rm store}(Q^{\rm store})^T),
  \qquad \Pi=QQ^T,
\]
and
\[
\begin{aligned}
  E^{\Pi,\rm norm/store2}_{ij}
  =
  &\ \gamma_d\sum_a
    |Q^{\rm store}_{ia}|\,|Q^{\rm store}_{ja}|\\
  &+\sum_a(\tau_{ia}+\rho_{ia})|Q^{\rm store}_{ja}|
   +\sum_a |Q_{ia}|(\tau_{ja}+\rho_{ja}),
\end{aligned}
\]
if the length-\(d\) dot products used to form the projector are
\(\gamma\)-valid, then
\[
  |\widehat\Pi_{ij}-\Pi_{ij}|
  \le E^{\Pi,\rm norm/store2}_{ij}
  \qquad\text{for all }i,j .
\]
For two-sided Algorithm~3 preprocessing, the new `ComputedMatrix`
constructors supply the left and right computed-basis certificates consumed by
the existing computed-projector theorem
\(\fl(\widehat\Pi_L A\widehat\Pi_R)\).  Thus storage layouts with only
normwise certificates are charged explicitly; no floating-point or
law-transfer error is added for sampling probabilities or sampling laws.
\end{corollary}

\begin{formalized}
\leanfit{ComputedMatrix.ofEntrywiseBoundThenFrobeniusStorage}
\leanfit{ComputedMatrix.ofFrobeniusBoundThenFrobeniusStorage}
\leanfit{ComputedMatrix.ofColumnVecNorm2BoundThenColumnVecNorm2Storage}
\leanfit{ComputedMatrix.ofRectOpNorm2BoundThenRectOpNorm2Storage}
\leanfit{fl_basisColumnProjector_of_frobeniusCertifiedFrobeniusStoredBasis_entry_error_bound}
\leanfit{fl_basisColumnProjector_of_columnwiseCertifiedColumnwiseStoredBasis_entry_error_bound}
\leanfit{fl_basisColumnProjector_of_opNormCertifiedOpNormStoredBasis_entry_error_bound}
\leanfit{ComputedPreconditioner.flBasisColumnProjector}
\leanfit{fl_preconditionElementsWithComputedBasisProjectors_total_error_bound}
\end{formalized}

\begin{corollary}[Frobenius-certified QR/SVD basis routine handoff for Theorem 183]
Let a concrete QR, SVD, singular-vector, or orthonormal-basis routine return
\(\widehat Q\in\mathbb R^{m\times d}\) together with a normwise certificate
against the exact analysis basis \(Q\),
\[
  \|\widehat Q-Q\|_F\le\eta,
  \qquad \eta\ge0 .
\]
Define
\[
  \widehat\Pi=\fl(\widehat Q\widehat Q^T),
  \qquad
  \Pi=QQ^T,
\]
and
\[
\begin{aligned}
  E^{\Pi,F}_{ij}
  =
  &\ \gamma_d\sum_a |\widehat Q_{ia}|\,|\widehat Q_{ja}|\\
  &+\sum_a \eta\,|\widehat Q_{ja}|
  +\sum_a |Q_{ia}|\,\eta .
\end{aligned}
\]
If the length-\(d\) dot products used to form
\(\widehat Q\widehat Q^T\) are \(\gamma\)-valid, then
\[
  |\widehat\Pi_{ij}-\Pi_{ij}|\le E^{\Pi,F}_{ij}
  \qquad\text{for all }i,j .
\]
The proof first uses the Frobenius entry bound
\(|\widehat Q_{ia}-Q_{ia}|\le\|\widehat Q-Q\|_F\le\eta\), then applies the
entrywise certified-projector theorem.  Thus a normwise QR/SVD certificate is
not silently treated as exact singular-vector data; it is converted into the
displayed implementation-facing projector radius.

For the two-sided projector branch, suppose two routines return
\(\widehat Q_L,\widehat Q_R\) with Frobenius radii \(\eta_L,\eta_R\), and
define \(E^{\Pi_L,F}\) and \(E^{\Pi_R,F}\) by the same formula.  If the
matrix-product dimensions are \(\gamma\)-valid, then for every \(i,j\),
\[
\begin{aligned}
&|\fl(\widehat\Pi_L A\widehat\Pi_R)_{ij}
  -(\Pi_LA\Pi_R)_{ij}|\\
&\quad\le
\gamma_n\sum_t
  |\fl(\widehat\Pi_LA)_{it}|\,|\widehat\Pi_{R,tj}|\\
&\qquad+
\sum_t
  \left(\gamma_m\sum_a|\widehat\Pi_{L,ia}|\,|A_{at}|\right)
  |\widehat\Pi_{R,tj}|\\
&\qquad+
\sum_t
  \left(\sum_a E^{\Pi_L,F}_{ia}|A_{at}|\right)
  |\widehat\Pi_{R,tj}|
  +
  \sum_t |(\Pi_LA)_{it}|E^{\Pi_R,F}_{tj}.
\end{aligned}
\]
The displayed bound charges the normwise basis-routine certificates,
projector formation, and the two rounded preprocessing products.  Sampling
probabilities and sampling laws remain exact mathematical inputs.
\end{corollary}

\begin{formalized}
\leanfit{ComputedMatrix.ofFrobeniusBound}
\leanfit{frobeniusCertifiedBasisProjectorEntryErrorBudget}
\leanfit{frobeniusCertifiedBasisProjectorEntryErrorBudget_nonneg}
\leanfit{fl_basisColumnProjector_of_frobeniusCertifiedBasis_entry_error_bound}
\leanfit{ComputedPreconditioner.flBasisColumnProjectorOfFrobeniusCertifiedBasis}
\leanfit{ComputedPreconditioner.flBasisColumnProjectorOfFrobeniusCertifiedBasis_entry_error_bound}
\leanfit{fl_preconditionElementsWithFrobeniusCertifiedBasisProjectors_total_error_bound}
\end{formalized}

\begin{corollary}[Columnwise-certified QR/SVD basis routine handoff for Theorem 183]
Let a concrete QR, SVD, singular-vector, or orthonormal-basis routine return
\(\widehat Q\in\mathbb R^{m\times d}\) together with per-column Euclidean
certificates against the exact analysis basis \(Q\),
\[
  \|\widehat Q_{:a}-Q_{:a}\|_2\le\eta_a,
  \qquad \eta_a\ge0,
  \qquad a=1,\ldots,d .
\]
Define
\[
  \widehat\Pi=\fl(\widehat Q\widehat Q^T),
  \qquad
  \Pi=QQ^T,
\]
and
\[
\begin{aligned}
  E^{\Pi,\mathrm{col}}_{ij}
  =
  &\ \gamma_d\sum_a |\widehat Q_{ia}|\,|\widehat Q_{ja}|\\
  &+\sum_a \eta_a\,|\widehat Q_{ja}|
  +\sum_a |Q_{ia}|\,\eta_a .
\end{aligned}
\]
If the length-\(d\) dot products used to form
\(\widehat Q\widehat Q^T\) are \(\gamma\)-valid, then
\[
  |\widehat\Pi_{ij}-\Pi_{ij}|
  \le E^{\Pi,\mathrm{col}}_{ij}
  \qquad\text{for all }i,j .
\]
Indeed, coordinate domination gives
\(|\widehat Q_{ia}-Q_{ia}|\le
\|\widehat Q_{:a}-Q_{:a}\|_2\le\eta_a\) for each column \(a\), and the
entrywise certified-projector theorem applies with this column-dependent
radius.

For the two-sided projector branch, suppose two routines return
\(\widehat Q_L,\widehat Q_R\) with columnwise radii
\(\eta^L_a,\eta^R_a\), and define \(E^{\Pi_L,\mathrm{col}}\) and
\(E^{\Pi_R,\mathrm{col}}\) by the same formula.  If the matrix-product
dimensions are \(\gamma\)-valid, then for every \(i,j\),
\[
\begin{aligned}
&|\fl(\widehat\Pi_L A\widehat\Pi_R)_{ij}
  -(\Pi_LA\Pi_R)_{ij}|\\
&\quad\le
\gamma_n\sum_t
  |\fl(\widehat\Pi_LA)_{it}|\,|\widehat\Pi_{R,tj}|\\
&\qquad+
\sum_t
  \left(\gamma_m\sum_a|\widehat\Pi_{L,ia}|\,|A_{at}|\right)
  |\widehat\Pi_{R,tj}|\\
&\qquad+
\sum_t
  \left(\sum_a E^{\Pi_L,\mathrm{col}}_{ia}|A_{at}|\right)
  |\widehat\Pi_{R,tj}|
  +
  \sum_t |(\Pi_LA)_{it}|E^{\Pi_R,\mathrm{col}}_{tj}.
\end{aligned}
\]
The displayed bound charges the per-column basis-routine certificates,
projector formation, and the two rounded preprocessing products.  Sampling
probabilities and sampling laws remain exact mathematical inputs.
\end{corollary}

\begin{formalized}
\leanfit{ComputedMatrix.ofColumnVecNorm2Bound}
\leanfit{columnwiseCertifiedBasisProjectorEntryErrorBudget}
\leanfit{columnwiseCertifiedBasisProjectorEntryErrorBudget_nonneg}
\leanfit{fl_basisColumnProjector_of_columnwiseCertifiedBasis_entry_error_bound}
\leanfit{ComputedPreconditioner.flBasisColumnProjectorOfColumnwiseCertifiedBasis}
\leanfit{ComputedPreconditioner.flBasisColumnProjectorOfColumnwiseCertifiedBasis_entry_error_bound}
\leanfit{fl_preconditionElementsWithColumnwiseCertifiedBasisProjectors_total_error_bound}
\end{formalized}

\begin{corollary}[Operator-certified QR/SVD basis routine handoff for Theorem 183]
Let a concrete QR, SVD, singular-vector, orthonormal-basis, or subspace
perturbation routine return
\(\widehat Q\in\mathbb R^{m\times d}\).  Suppose it supplies the rectangular
operator certificate
\[
  \|(\widehat Q-Q)x\|_2\le \eta\|x\|_2
  \qquad\text{for all }x\in\mathbb R^d,
  \qquad \eta\ge0,
\]
against the exact analysis basis \(Q\).  Define
\[
  \widehat\Pi=\fl(\widehat Q\widehat Q^T),
  \qquad
  \Pi=QQ^T,
\]
and
\[
\begin{aligned}
  E^{\Pi,2}_{ij}
  =
  &\ \gamma_d\sum_a |\widehat Q_{ia}|\,|\widehat Q_{ja}|\\
  &+\sum_a \eta\,|\widehat Q_{ja}|
  +\sum_a |Q_{ia}|\,\eta .
\end{aligned}
\]
If the length-\(d\) dot products used to form
\(\widehat Q\widehat Q^T\) are \(\gamma\)-valid, then
\[
  |\widehat\Pi_{ij}-\Pi_{ij}|\le E^{\Pi,2}_{ij}
  \qquad\text{for all }i,j .
\]
Indeed, testing the operator certificate on the standard basis vector \(e_a\)
gives \(\|(\widehat Q-Q)e_a\|_2\le\eta\); coordinate domination then gives
\(|\widehat Q_{ia}-Q_{ia}|\le\eta\).  The entrywise certified-projector
theorem applies with this uniform radius.

For the two-sided projector branch, suppose two routines return
\(\widehat Q_L,\widehat Q_R\) with operator radii \(\eta_L,\eta_R\), and
define \(E^{\Pi_L,2}\) and \(E^{\Pi_R,2}\) by the same formula.  If the
matrix-product dimensions are \(\gamma\)-valid, then for every \(i,j\),
\[
\begin{aligned}
&|\fl(\widehat\Pi_L A\widehat\Pi_R)_{ij}
  -(\Pi_LA\Pi_R)_{ij}|\\
&\quad\le
\gamma_n\sum_t
  |\fl(\widehat\Pi_LA)_{it}|\,|\widehat\Pi_{R,tj}|\\
&\qquad+
\sum_t
  \left(\gamma_m\sum_a|\widehat\Pi_{L,ia}|\,|A_{at}|\right)
  |\widehat\Pi_{R,tj}|\\
&\qquad+
\sum_t
  \left(\sum_a E^{\Pi_L,2}_{ia}|A_{at}|\right)
  |\widehat\Pi_{R,tj}|
  +
  \sum_t |(\Pi_LA)_{it}|E^{\Pi_R,2}_{tj}.
\end{aligned}
\]
The displayed bound charges the operator-norm routine certificates,
projector formation, and the two rounded preprocessing products.  Sampling
probabilities and sampling laws remain exact mathematical inputs.
\end{corollary}

\begin{formalized}
\leanfit{ComputedMatrix.ofRectOpNorm2Bound}
\leanfit{opNormCertifiedBasisProjectorEntryErrorBudget}
\leanfit{opNormCertifiedBasisProjectorEntryErrorBudget_nonneg}
\leanfit{fl_basisColumnProjector_of_opNormCertifiedBasis_entry_error_bound}
\leanfit{ComputedPreconditioner.flBasisColumnProjectorOfOpNormCertifiedBasis}
\leanfit{ComputedPreconditioner.flBasisColumnProjectorOfOpNormCertifiedBasis_entry_error_bound}
\leanfit{fl_preconditionElementsWithOpNormCertifiedBasisProjectors_total_error_bound}
\end{formalized}

\begin{corollary}[Add-zero stored basis before projector formation]
Let \(Q\in\mathbb R^{m\times k}\) be the exact basis or singular-vector table
used in the analysis, and suppose the implemented table is stored entrywise by
the rounded arithmetic copy
\[
  \widetilde Q_{ia}=\fl_{\rm add}(Q_{ia},0),
  \qquad
  |\widetilde Q_{ia}-Q_{ia}|\le u\,|Q_{ia}|.
\]
Let the implemented projector be formed by rounded dot products,
\[
  \widehat\Pi_{ij}
  = \fl\!\left(\sum_{a=1}^k \widetilde Q_{ia}\widetilde Q_{ja}\right),
  \qquad
  \Pi_{ij}=(Q Q^T)_{ij}.
\]
If the length-\(k\) dot products satisfy the usual \(\gamma_k\) validity
condition, then for every \(i,j\)
\[
\begin{aligned}
|\widehat\Pi_{ij}-\Pi_{ij}|
&\le
\gamma_k \sum_{a=1}^k |\widetilde Q_{ia}|\,|\widetilde Q_{ja}| \\
&\quad
+ u\sum_{a=1}^k |Q_{ia}|\,|\widetilde Q_{ja}|
+ u\sum_{a=1}^k |Q_{ia}|\,|Q_{ja}|.
\end{aligned}
\]
Thus the implemented projector path charges the arithmetic storage of the
basis/singular-vector table and the rounded formation of
\(\widetilde Q\widetilde Q^T\).  The sampling probabilities and sampling laws
are unchanged exact mathematical inputs.
\end{corollary}

\begin{formalized}
\leanfit{ComputedMatrix.flAddZeroRight}
\leanfit{ComputedMatrix.flAddZeroRight_entry_error_bound}
\leanfit{ComputedPreconditioner.flBasisColumnProjectorStoredBasisAddZeroRight}
\leanfit{ComputedPreconditioner.flBasisColumnProjectorStoredBasisAddZeroRight_entry_error_bound}
\end{formalized}

\begin{corollary}[Multiply-one and subtract-zero stored bases before projector formation]
Let \(Q\in\mathbb R^{m\times k}\) be the exact basis or singular-vector table
used in the analysis.  Define the two additional implemented storage paths
\[
  \widetilde Q^{(\times)}_{ia}=\fl_{\rm mul}(Q_{ia},1),
  \qquad
  \widetilde Q^{(-)}_{ia}=\fl_{\rm sub}(Q_{ia},0).
\]
For each \(\star\in\{(\times),(-)\}\), suppose
\[
  |\widetilde Q^{\star}_{ia}-Q_{ia}|\le u\,|Q_{ia}|
  \qquad\text{for all }i,a,
\]
and form
\[
  \widehat\Pi^{\star}_{ij}
  =\fl\!\left(\sum_{a=1}^k
      \widetilde Q^{\star}_{ia}\widetilde Q^{\star}_{ja}\right),
  \qquad
  \Pi_{ij}=(Q Q^T)_{ij}.
\]
If the length-\(k\) projector dot products satisfy the usual \(\gamma_k\)
validity condition, then for every \(i,j\) and every
\(\star\in\{(\times),(-)\}\),
\[
\begin{aligned}
|\widehat\Pi^{\star}_{ij}-\Pi_{ij}|
&\le
\gamma_k \sum_{a=1}^k
  |\widetilde Q^{\star}_{ia}|\,|\widetilde Q^{\star}_{ja}| \\
&\quad
+ u\sum_{a=1}^k |Q_{ia}|\,|\widetilde Q^{\star}_{ja}|
+ u\sum_{a=1}^k |Q_{ia}|\,|Q_{ja}|.
\end{aligned}
\]
Thus the implementation-facing projector theorem now covers the three rounded
copy paths \( \fl_{\rm mul}(Q_{ia},1)\), \(\fl_{\rm add}(Q_{ia},0)\), and
\(\fl_{\rm sub}(Q_{ia},0)\).  These are non-probability storage operations;
sampling probabilities and sampling laws remain exact mathematical inputs.
\end{corollary}

\begin{formalized}
\leanfit{ComputedMatrix.flMulOne}
\leanfit{ComputedMatrix.flMulOne_entry_error_bound}
\leanfit{ComputedMatrix.flSubZeroRight}
\leanfit{ComputedMatrix.flSubZeroRight_entry_error_bound}
\leanfit{ComputedPreconditioner.flBasisColumnProjectorStoredBasisMulOne}
\leanfit{ComputedPreconditioner.flBasisColumnProjectorStoredBasisMulOne_entry_error_bound}
\leanfit{ComputedPreconditioner.flBasisColumnProjectorStoredBasisSubZeroRight}
\leanfit{ComputedPreconditioner.flBasisColumnProjectorStoredBasisSubZeroRight_entry_error_bound}
\end{formalized}

\begin{corollary}[Two-sided floating-point preprocessing]
Let \(B=\Pi_LA\) and \(\widehat B=\fl(\Pi_LA)\).  If the computed two-sided
branch is \(\fl(\widehat B\Pi_R)\), then
\[
\begin{aligned}
&\left|
  \fl(\widehat B\Pi_R)_{ij}-(\Pi_LA\Pi_R)_{ij}
\right| \\
&\quad\le
\gamma_n\sum_k |\widehat B_{ik}|\,|(\Pi_R)_{kj}|
+
\sum_k
\left(
  \gamma_m\sum_a |(\Pi_L)_{ia}|\,|A_{ak}|
\right)
|(\Pi_R)_{kj}|.
\end{aligned}
\]
\end{corollary}

\begin{formalized}
\leanfit{preconditionColumns_entry_error_bound_of_entrywise}
\leanfit{fl_preconditionElements_error_bound}
\leanfit{fl_preconditionElementsWithComputed_error_bound}
\leanfit{fl_preconditionElementsWithComputed_total_error_bound}
\end{formalized}

\begin{corollary}[Computed-projector replacement for Corollary 184]
Let
\[
  \Pi_L=Q_LQ_L^T,\qquad \Pi_R=Q_RQ_R^T
\]
be the exact analysis projectors, and let
\[
  \widehat\Pi_L=\fl(\widehat Q_L\widehat Q_L^T),
  \qquad
  \widehat\Pi_R=\fl(\widehat Q_R\widehat Q_R^T)
\]
be the projectors formed in the implementation from computed basis or
singular-vector tables.  Suppose the preceding computed-projector bound gives
entrywise radii \(E^{\Pi_L}\) and \(E^{\Pi_R}\).  If the product dimensions are
\(\gamma\)-valid, then for every \(i,j\),
\[
\begin{aligned}
&|\fl(\widehat\Pi_L A\widehat\Pi_R)_{ij}
  -(\Pi_LA\Pi_R)_{ij}|\\
&\quad\le
\gamma_n\sum_t
  |\fl(\widehat\Pi_LA)_{it}|\,|\widehat\Pi_{R,tj}|\\
&\qquad+
\sum_t
  \left(\gamma_m\sum_a|\widehat\Pi_{L,ia}|\,|A_{at}|\right)
  |\widehat\Pi_{R,tj}|\\
&\qquad+
\sum_t
  \left(\sum_a E^{\Pi_L}_{ia}|A_{at}|\right)
  |\widehat\Pi_{R,tj}|
  +
  \sum_t |(\Pi_LA)_{it}|E^{\Pi_R}_{tj}.
\end{aligned}
\]
The first two lines are the rounded second and first products; the last two
lines are the left and right computed-projector formation errors.  Sampling
probabilities and sampling laws remain exact under the current convention.
\end{corollary}

\begin{formalized}
\leanfit{fl_preconditionElementsWithComputedBasisProjectors_total_error_bound}
\end{formalized}

\section{Equation (8): Least-Squares Sketching}

Equation~(8) is the least-squares problem
\[
  x_{\rm opt}=\arg\min_{x\in\R^n}\Norm{Ax-b}.
\]
For a sketched problem \((SA,Sb)\), define
\[
  f(x)=\Norm{Ax-b}^2,
  \qquad
  \tilde f(x)=\Norm{SAx-Sb}^2.
\]

\begin{definition}[Residual-objective preserving sketch]
The sketched problem preserves the least-squares objective with accuracy
\(\varepsilon\) if, for every \(x\),
\[
  (1-\varepsilon)f(x)\le \tilde f(x)\le (1+\varepsilon)f(x).
\]
\end{definition}

\begin{theorem}[Quadratic subspace-embedding bridge]
Suppose every residual has coordinates \(y(x)\in\R^d\) in an orthonormal
residual coordinate system, so that
\[
  f(x)=\Norm{y(x)}^2.
\]
Suppose further that, for a fixed sketch outcome,
\[
  \tilde f(x)
  =
  \Norm{y(x)}^2+y(x)^T E\,y(x)
  \quad\text{for every }x,
\]
and that the coordinate Gram-error matrix satisfies the vector-action
operator bound
\[
  \Norm{Ez}\le\varepsilon\Norm{z}
  \quad\text{for every }z\in\R^d.
\]
Then the sketch preserves the least-squares residual objective:
\[
  (1-\varepsilon)f(x)\le \tilde f(x)\le (1+\varepsilon)f(x)
  \quad\text{for every }x.
\]
Moreover, on any finite probability space, a proved lower bound for the
operator event \(\Norm{Ez}\le\varepsilon\Norm{z}\) transfers directly to the
least-squares preservation event.
\end{theorem}


\begin{formalized}
\leanfit{abs_vecInnerProduct_le_vecNorm2_mul}
\leanfit{abs_vecInnerProduct_matMulVec_le_of_opNorm2Le}
\leanfit{preservesLSObjective_of_coordinate_quadratic_error}
\leanfit{eventProb_preservesLSObjective_of_coordinate_quadratic_error}
\leanfit{preservesLSObjective_of_coordinate_finiteLoewner_error}
\leanfit{eventProb_preservesLSObjective_of_coordinate_finiteLoewner_error}
\end{formalized}

\begin{theorem}[Finite-Loewner subspace-embedding bridge]
In the same coordinate model, suppose instead that the Gram-error matrix
satisfies the two-sided Loewner event
\[
  E\preceq \varepsilon I_d,
  \qquad
  -E\preceq \varepsilon I_d .
\]
Then the sketch preserves every residual objective with accuracy
\(\varepsilon\).  On a finite probability space, any proved lower bound for
this two-sided Loewner event transfers to the least-squares preservation event.
\end{theorem}


\begin{theorem}[Sketched minimizer residual guarantee]
Assume \(\varepsilon<1\), the sketch preserves every residual objective as
above, and \(\widehat x\) is an exact minimizer of the sketched problem.  Then,
for every comparison vector \(x_{\rm ref}\),
\[
  \Norm{A\widehat x-b}^2
  \le
  \frac{1+\varepsilon}{1-\varepsilon}
  \Norm{Ax_{\rm ref}-b}^2.
\]
In particular, this applies to an exact minimizer \(x_{\rm opt}\) of the
original problem.
\end{theorem}

\begin{formalized}
\leanfit{lsResidual}
\leanfit{lsObjective}
\leanfit{IsLeastSquaresMinimizer}
\leanfit{PreservesLSObjective}
\leanfit{lsObjective_le_of_sketch_preserves}
\end{formalized}

\begin{corollary}[\(1+\eta\) residual-objective form]
If additionally
\[
  \frac{1+\varepsilon}{1-\varepsilon}\le 1+\eta,
\]
then
\[
  \Norm{A\widehat x-b}^2
  \le
  (1+\eta)\Norm{Ax_{\rm ref}-b}^2.
\]
\end{corollary}

\begin{formalized}
\leanfit{lsObjective_le_one_add_eta_of_sketch_preserves}
\end{formalized}

\begin{corollary}[Randomized sketched-minimizer transfer]
On any finite probability space, if the sketch-preservation event has
probability at least \(\alpha\), and for every outcome \(\omega\) the vector
\(\widehat x(\omega)\) is an exact minimizer of the sketched problem
\((S_\omega A,S_\omega b)\), then the same probability lower bound applies to
the sketched-minimizer objective guarantee:
\[
  \mathbb P\!\left[
    \Norm{A\widehat x(\omega)-b}^2
    \le
    \frac{1+\varepsilon}{1-\varepsilon}
    \Norm{Ax_{\rm ref}-b}^2
  \right]\ge \alpha .
\]
The \(1+\eta\) version follows by the same scalar comparison as above.
\end{corollary}

\begin{formalized}
\leanfit{eventProb_lsObjective_le_of_preserves}
\leanfit{eventProb_lsObjective_le_one_add_eta_of_preserves}
\leanfit{eventProb_lsObjective_le_one_add_eta_of_coordinate_finiteLoewner_error}
\end{formalized}

\begin{corollary}[Leverage-score row sampling to least-squares objective]
Let \(U\in\R^{m\times d}\) have orthonormal columns, and let the Algorithm~2
row trace be sampled from the leverage-score law
\[
  p_i=\frac{\|U_{i,*}\|_2^2}{d}.
\]
Define the canonical residual coordinates
\[
  y(x)=U^T(Ax-b).
\]
Assume that every residual lies in the column span of \(U\), equivalently
\[
  (Ax-b)_i=\sum_{\ell=1}^d U_{i\ell}\,y_\ell(x),
  \qquad 1\le i\le m .
\]
The repository proves from this column-space condition and \(U^T U=I\) that
\[
  \Norm{Ax-b}^2=\Norm{y(x)}^2 .
\]
Let the concrete sampled least-squares instance use rows
\[
  (S_\omega A)_{t,*}
  =
  \frac{A_{i_t,*}}{\sqrt{s\,p_{i_t}}},
  \qquad
  (S_\omega b)_t
  =
  \frac{b_{i_t}}{\sqrt{s\,p_{i_t}}}.
\]
Then the repository proves the sampled objective identity
\[
  \tilde f_\omega(x)
  =
  \|y(x)\|_2^2+
  y(x)^T\!\left(\widetilde U(\omega)^T\widetilde U(\omega)-I\right)y(x).
\]
Then the leverage-score event for equation~(7) transfers to the least-squares
sketched-minimizer guarantee.  The older operator-event route gives
\[
  \mathbb P\!\left[
    \Norm{A\widehat x(\omega)-b}^2
    \le
    (1+\eta)\Norm{Ax_{\rm opt}-b}^2
  \right]
  \ge
  1-\frac{1}{s(\varepsilon/d)^2},
\]
provided \(\varepsilon<1\) and
\((1+\varepsilon)/(1-\varepsilon)\le1+\eta\).

The sharper finite-Loewner theorem proved for equation~(7) gives the
source-aligned Bennett sample-budget version.  For the augmented-span basis,
let \(d=\dim\operatorname{span}\{A_{*,j},b\}\).  If \(d>1\), \(\varepsilon>0\),
\(\delta>0\), and
\[
  \log\frac{2d}{\delta}
    \le
  s\frac{\varepsilon^2}{2(d-1)+(2/3)d\varepsilon},
  \qquad
  \log\frac{2d}{\delta}
    \le
  s\frac{\varepsilon^2}{2(d-1)+(2/3)\varepsilon},
\]
then the same sketched-minimizer event holds with probability at least
\(1-\delta\).

The fully floating-point version replaces
\(\widetilde U^T\widetilde U-I\) by the rounded Gram matrix and replaces the
preservation radius \(\varepsilon\) by
\[
  \varepsilon+\tau_{\rm Gram}^{\rm fp},\qquad
  \tau_{\rm Gram}^{\rm fp}
  =
  \operatorname{rowSampleGramFullFpPerturbBudget}(\mathrm{fp},s,U).
\]
The budget is proved nonnegative locally.
\end{corollary}

\begin{formalized}
\leanfit{rowSampleGramFullFpPerturbBudget_nonneg}
\leanfit{rowSampleLSMatrixWithBasisScale}
\leanfit{rowSampleLSVectorWithBasisScale}
\leanfit{fl_rowSampleLSMatrixWithBasisScale}
\leanfit{fl_rowSampleLSVectorWithBasisScale}
\leanfit{fl_rowSampleLSMatrixWithBasisScale_error_bound}
\leanfit{fl_rowSampleLSVectorWithBasisScale_error_bound}
\leanfit{fl_rowSampleLSResidualWithBasisScale_error_bound}
\leanfit{fl_rowSampleLSResidualWithBasisScale_error_bound_of_positiveProb}
\leanfit{lsObjective_residual_difference_bound}
\leanfit{lsObjective_residual_budget_bound}
\leanfit{rowSampleLSResidualFpBudget}
\leanfit{rowSampleLSResidualFpBudget_nonneg}
\leanfit{rowSampleLSObjectiveFpBudget}
\leanfit{rowSampleLSObjectiveFpBudget_nonneg}
\leanfit{fl_rowSampleLSObjectiveWithBasisScale_error_bound_of_positiveProb}
\leanfit{fl_rowSampleLSObjectiveWithBasisScale_error_bound_of_positiveProb_budget}
\leanfit{lsObjective_le_of_sketch_preserves_with_objective_error}
\leanfit{lsObjective_le_one_add_eta_of_sketch_preserves_with_objective_error}
\leanfit{eventProb_lsObjective_le_one_add_eta_of_preserves_with_objective_error}
\leanfit{eventProb_lsObjective_le_one_add_eta_of_preserves_with_pointwise_objective_error}
\leanfit{eventProb_lsObjective_le_one_add_eta_of_preserves_with_objective_error_on_event}
\leanfit{eventProb_lsObjective_le_one_add_eta_of_coordinate_finiteLoewner_error_with_objective_error_on_event}
\leanfit{residualCoordinates}
\leanfit{ResidualsInColumnSpace}
\leanfit{ColumnsAndRhsInColumnSpace}
\leanfit{residualCoordinatesFromColumns}
\leanfit{residualsInColumnSpace_of_residual_representation}
\leanfit{lsResidual_eq_basis_sum_of_columnsAndRhsInColumnSpace}
\leanfit{residualsInColumnSpace_of_columnsAndRhsInColumnSpace}
\leanfit{augmentedDataVector}
\leanfit{augmentedDataSpan}
\leanfit{augmentedSpanBasisMatrix}
\leanfit{augmentedSpanColumnCoords}
\leanfit{augmentedSpanRhsCoords}
\leanfit{hasOrthonormalColumns_augmentedSpanBasisMatrix}
\leanfit{columnsAndRhsInColumnSpace_augmentedSpanBasisMatrix}
\leanfit{residualsInColumnSpace_augmentedSpanBasisMatrix}
\leanfit{lsObjective_eq_vecNorm2Sq_residualCoordinates_of_residualsInColumnSpace}
\leanfit{rowSampleLSResidualWithBasisScale_eq_coord}
\leanfit{rowSampleLSObjectiveWithBasisScale_eq_coordinate_quadratic_error}
\leanfit{eventProb_lsObjective_le_one_add_eta_of_coordinate_quadratic_error}
\leanfit{eventProb_lsObjective_le_one_add_eta_of_coordinate_finiteLoewner_error}
\leanfit{leverageTraceProbability_eventProb_lsObjective_le_one_add_eta_of_coordinate_quadratic_error}
\leanfit{leverageTraceProbability_eventProb_rowSampleLSObjective_le_one_add_eta_of_residual_coordinates}
\leanfit{leverageTraceProbability_eventProb_rowSampleLSObjective_le_one_add_eta_of_residualsInColumnSpace}
\leanfit{leverageTraceProbability_eventProb_rowSampleLSObjective_le_one_add_eta_of_augmentedSpan}
\leanfit{leverageTraceProbability_eventProb_rowSampleLSObjective_le_one_add_eta_of_augmentedSpan_sample_budget}
\leanfit{leverageTraceProbability_eventProb_fl_lsObjective_le_one_add_eta_of_coordinate_quadratic_error}
\leanfit{leverageTraceProbability_eventProb_fl_lsObjective_le_one_add_eta_of_augmentedSpan}
\leanfit{leverageTraceProbability_eventProb_fl_lsObjective_le_one_add_eta_of_augmentedSpan_sample_budget}
\leanfit{leverageTraceProbability_eventProb_fl_rowSampleLSObjective_le_one_add_eta_of_augmentedSpan_sample_budget_of_objective_error}
\leanfit{leverageTraceProbability_eventProb_fl_lsObjective_le_one_add_eta_of_idBasis}
\end{formalized}

\begin{theorem}[Literal rounded sampled-row residual perturbation]
Let \(S_\omega A\) and \(S_\omega b\) denote the exact sampled and scaled
least-squares matrix and right-hand side obtained from Algorithm~2 equation~(6).
Let \(\fl(S_\omega A)\) and \(\fl(S_\omega b)\) be the concrete rounded
implementation that applies the repository's floating-point division primitive
to each sampled/scaled entry.  If the sampled scaling denominator at row \(t\)
is nonzero, then
\[
\begin{aligned}
&\left|
  \bigl(\fl(S_\omega A)x-\fl(S_\omega b)\bigr)_t
  -
  \bigl(S_\omega Ax-S_\omega b\bigr)_t
 \right| \\
&\qquad\le
 \sum_j
   u_{\rm fp}\,\left|(S_\omega A)_{tj}\right|\,|x_j|
   +
   u_{\rm fp}\,\left|(S_\omega b)_t\right|.
\end{aligned}
\]
Under the product sampling law, the positive-probability support event
discharges this nonzero-denominator premise with probability one for sampled
rows.
\end{theorem}


\begin{formalized}
\leanfit{fl_rowSampleLSMatrixWithBasisScale}
\leanfit{fl_rowSampleLSVectorWithBasisScale}
\leanfit{fl_rowSampleLSMatrixWithBasisScale_error_bound}
\leanfit{fl_rowSampleLSVectorWithBasisScale_error_bound}
\leanfit{fl_rowSampleLSResidualWithBasisScale_error_bound}
\leanfit{fl_rowSampleLSResidualWithBasisScale_error_bound_of_positiveProb}
\end{formalized}

\begin{corollary}[Literal rounded sampled-row objective perturbation]
For fixed sampled rows \(\omega\) and a fixed vector \(x\), define the rowwise
budget
\[
  B_t(x)=
  \sum_j u_{\rm fp}\,\left|(S_\omega A)_{tj}\right|\,|x_j|
  +
  u_{\rm fp}\,\left|(S_\omega b)_t\right|.
\]
On the positive-probability support event of the row sampler,
\[
\begin{aligned}
&\left|
  \Norm{\fl(S_\omega A)x-\fl(S_\omega b)}^2
  -
  \Norm{S_\omega Ax-S_\omega b}^2
 \right| \\
&\qquad\le
 2\,\Norm{S_\omega Ax-S_\omega b}\,\Norm{B(x)}
 + \Norm{B(x)}^2 .
\end{aligned}
\]
\end{corollary}


\begin{formalized}
\leanfit{vecNorm2Sq_le_of_abs_le}
\leanfit{vecNorm2_le_of_abs_le}
\leanfit{abs_vecNorm2Sq_add_sub_le}
\leanfit{lsObjective_residual_difference_bound}
\leanfit{lsObjective_residual_budget_bound}
\leanfit{rowSampleLSResidualFpBudget}
\leanfit{rowSampleLSResidualFpBudget_nonneg}
\leanfit{rowSampleLSObjectiveFpBudget}
\leanfit{rowSampleLSObjectiveFpBudget_nonneg}
\leanfit{fl_rowSampleLSObjectiveWithBasisScale_error_bound_of_positiveProb}
\leanfit{fl_rowSampleLSObjectiveWithBasisScale_error_bound_of_positiveProb_budget}
\end{formalized}

\begin{corollary}[High-probability literal rounded sampled-row minimizer]
Let \(U\) be the augmented-span leverage basis for
\(\operatorname{span}\{A_{*,1},\ldots,A_{*,n},b\}\), and let
\(d=\dim\operatorname{span}\{A_{*,1},\ldots,A_{*,n},b\}>1\).  Assume
\(\varepsilon>0\), \(\varepsilon<1\), \(\delta>0\), and the two Bennett
sample-budget inequalities
\[
  \log\frac{2d}{\delta}
    \le
  s\frac{\varepsilon^2}{2(d-1)+(2/3)d\varepsilon},
  \qquad
  \log\frac{2d}{\delta}
    \le
  s\frac{\varepsilon^2}{2(d-1)+(2/3)\varepsilon}.
\]
Let \(\widehat x(\omega)\) be an exact minimizer of the literal rounded
sampled/scaled problem
\[
  \min_x \Norm{\fl(S_\omega A)x-\fl(S_\omega b)}^2.
\]
For \(x\), define
\[
  T_\omega(x)
  =
  2\,\Norm{S_\omega Ax-S_\omega b}\,\Norm{B_\omega(x)}
  +\Norm{B_\omega(x)}^2,
\]
where \(B_\omega(x)\) is the rowwise residual budget from the previous
corollary.  If, on the positive-probability support event,
\[
  T_\omega(\widehat x(\omega))+T_\omega(x_{\rm opt})
  \le
  \bigl((1+\eta)(1-\varepsilon)-(1+\varepsilon)\bigr)
  \Norm{Ax_{\rm opt}-b}^2,
\]
then
\[
  \mathbb P\!\left[
    \Norm{A\widehat x(\omega)-b}^2
    \le
    (1+\eta)\Norm{Ax_{\rm opt}-b}^2
  \right]
  \ge 1-\delta .
\]
\end{corollary}


\begin{formalized}
\leanfit{rowSampleLSObjectiveFpBudget}
\leanfit{rowSampleLSObjectiveFpBudget_nonneg}
\leanfit{fl_rowSampleLSObjectiveWithBasisScale_error_bound_of_positiveProb_budget}
\leanfit{lsObjective_le_of_sketch_preserves_with_objective_error}
\leanfit{lsObjective_le_one_add_eta_of_sketch_preserves_with_objective_error}
\leanfit{eventProb_lsObjective_le_one_add_eta_of_preserves_with_objective_error}
\leanfit{eventProb_lsObjective_le_one_add_eta_of_preserves_with_pointwise_objective_error}
\leanfit{eventProb_lsObjective_le_one_add_eta_of_preserves_with_objective_error_on_event}
\leanfit{eventProb_lsObjective_le_one_add_eta_of_coordinate_finiteLoewner_error_with_objective_error_on_event}
\leanfit{leverageTraceProbability_eventProb_fl_rowSampleLSObjective_le_one_add_eta_of_augmentedSpan_sample_budget_of_objective_error}
\end{formalized}

\begin{corollary}[High-probability literal rounded sampled-row approximate solver]
Use the notation and Bennett sample-budget assumptions of the previous
corollary.  Let \(\widehat x(\omega)\) be an additive-gap approximate
minimizer of the literal rounded sampled/scaled problem:
\[
  \Norm{\fl(S_\omega A)\widehat x(\omega)-\fl(S_\omega b)}^2
  \le
  \Norm{\fl(S_\omega A)y-\fl(S_\omega b)}^2+\sigma_\omega
  \quad\text{for all }y .
\]
If, on the positive-probability support event,
\[
  T_\omega(\widehat x(\omega))+T_\omega(x_{\rm opt})+\sigma_\omega
  \le
  \bigl((1+\eta)(1-\varepsilon)-(1+\varepsilon)\bigr)
  \Norm{Ax_{\rm opt}-b}^2,
\]
then
\[
  \mathbb P\!\left[
    \Norm{A\widehat x(\omega)-b}^2
    \le
    (1+\eta)\Norm{Ax_{\rm opt}-b}^2
  \right]
  \ge 1-\delta .
\]
\end{corollary}


\begin{formalized}
\leanfit{IsLeastSquaresApproxMinimizer}
\leanfit{isLeastSquaresApproxMinimizer_of_minimizer}
\leanfit{lsObjective_le_of_sketch_preserves_with_objective_error_and_solver_gap}
\leanfit{lsObjective_le_one_add_eta_of_sketch_preserves_with_objective_error_and_solver_gap}
\leanfit{eventProb_lsObjective_le_one_add_eta_of_preserves_with_objective_error_and_solver_gap_on_event}
\leanfit{eventProb_lsObjective_le_one_add_eta_of_coordinate_finiteLoewner_error_with_objective_error_and_solver_gap_on_event}
\leanfit{leverageTraceProbability_eventProb_fl_rowSampleLSObjective_le_one_add_eta_of_augmentedSpan_sample_budget_of_objective_error_and_solver_gap}
\end{formalized}

\begin{corollary}[High-probability literal rounded sampled-row solver with a forward-error certificate]
Use the notation and Bennett sample-budget assumptions of the previous two
corollaries.  For each sample trace \(\omega\), let \(x_*(\omega)\) be an
exact minimizer of the literal rounded sampled/scaled problem, and let
\(\widehat x(\omega)\) be the vector returned by a solver.  Suppose the solver
also provides a componentwise certificate \(d_\omega\ge 0\) such that
\[
  |\widehat x_j(\omega)-x_{*,j}(\omega)|\le d_{\omega,j},
  \qquad j=1,\ldots,n .
\]
For the literal rounded sampled/scaled matrix and vector, define the induced
residual budget
\[
  R_\omega(i)=
  \sum_{j=1}^n
    |\fl(S_\omega A)_{ij}|\,d_{\omega,j},
  \qquad i=1,\ldots,s,
\]
and the induced objective gap
\[
  \Sigma_\omega
  =
  2\,
  \Norm{\fl(S_\omega A)x_*(\omega)-\fl(S_\omega b)}
  \Norm{R_\omega}
  +\Norm{R_\omega}^2 .
\]
If, on the positive-probability support event,
\[
  T_\omega(\widehat x(\omega))+T_\omega(x_{\rm opt})+\Sigma_\omega
  \le
  \bigl((1+\eta)(1-\varepsilon)-(1+\varepsilon)\bigr)
  \Norm{Ax_{\rm opt}-b}^2,
\]
then
\[
  \mathbb P\!\left[
    \Norm{A\widehat x(\omega)-b}^2
    \le
    (1+\eta)\Norm{Ax_{\rm opt}-b}^2
  \right]
  \ge 1-\delta .
\]
\end{corollary}


\begin{formalized}
\leanfit{lsSolutionForwardResidualBudget}
\leanfit{lsSolutionForwardObjectiveGap}
\leanfit{lsResidual_difference_bound_of_solution_abs_le}
\leanfit{lsObjective_solution_forward_error_bound}
\leanfit{isLeastSquaresApproxMinimizer_of_solution_abs_le}
\leanfit{leverageTraceProbability_eventProb_fl_rowSampleLSObjective_le_one_add_eta_of_augmentedSpan_sample_budget_of_objective_error_and_forward_error}
\end{formalized}

\begin{corollary}[High-probability literal rounded sampled-row solver with a perturbed Gram certificate]
Use the notation and Bennett sample-budget assumptions of the previous
corollary.  Let
\[
  G_\omega=\fl(S_\omega A)^T\fl(S_\omega A),
  \qquad
  g_\omega=\fl(S_\omega A)^T\fl(S_\omega b).
\]
Suppose \(G_\omega\) has inverse \(H_\omega\), the exact rounded minimizer
\(x_*(\omega)\) satisfies the normal equations
\[
  G_\omega x_*(\omega)=g_\omega ,
\]
and the computed solver output satisfies a perturbed Gram system
\[
  (G_\omega+\Delta G_\omega)\widehat x(\omega)
  =
  g_\omega+\Delta g_\omega
\]
with componentwise perturbation radii
\[
  |(\Delta G_\omega)_{jk}|\le \varepsilon^G_\omega,
  \qquad
  |(\Delta g_\omega)_j|\le \varepsilon^g_\omega .
\]
Define
\[
  d_{\omega,i}
  =
  \sum_{j=1}^n |(H_\omega)_{ij}|
  \left(\varepsilon^G_\omega\sum_{k=1}^n
    |\widehat x_k(\omega)|+\varepsilon^g_\omega\right).
\]
Then \(d_\omega\) is a componentwise forward-error certificate for the
previous corollary.  Consequently, if the same support-event budget condition
holds with \(d_\omega\) in the induced \(\Sigma_\omega\), then
\[
  \mathbb P\!\left[
    \Norm{A\widehat x(\omega)-b}^2
    \le
    (1+\eta)\Norm{Ax_{\rm opt}-b}^2
  \right]
  \ge 1-\delta .
\]
\end{corollary}


\begin{formalized}
\leanfit{lsNormalMatrix}
\leanfit{lsNormalRhs}
\leanfit{gramForwardSolverDx}
\leanfit{gramForwardSolverDx_nonneg}
\leanfit{gram_forward_error_certificate_of_perturbed_gram_system}
\leanfit{leverageTraceProbability_eventProb_fl_rowSampleLSObjective_le_one_add_eta_of_augmentedSpan_sample_budget_of_objective_error_and_perturbed_gram_solver}
\end{formalized}

\begin{corollary}[High-probability literal rounded sampled-row solver from the local QR backward-error spec]
Use the notation and Bennett sample-budget assumptions of the previous
corollary.  Suppose that, for each sample trace \(\omega\), the solver output
\(\widehat x(\omega)\) satisfies the repository's local QR least-squares
backward-error specification for the rounded normal equations.  Concretely,
there exist perturbations \(\Delta G_\omega,\Delta g_\omega\) such that
\[
  (G_\omega+\Delta G_\omega)\widehat x(\omega)
  =
  g_\omega+\Delta g_\omega,
  \qquad
  \|\Delta G_\omega\|_F\le c^G_\omega,
  \qquad
  |(\Delta g_\omega)_j|\le c^g_\omega .
\]
Define
\[
  d_{\omega,i}
  =
  \sum_{j=1}^n |(H_\omega)_{ij}|
  \left(c^G_\omega\sum_{k=1}^n
    |\widehat x_k(\omega)|+\max\{c^g_\omega,0\}\right).
\]
Then \(d_\omega\) is a componentwise forward-error certificate for the
forward-error corollary.  Hence, if the same support-event budget condition
holds with this \(d_\omega\), then
\[
  \mathbb P\!\left[
    \Norm{A\widehat x(\omega)-b}^2
    \le
    (1+\eta)\Norm{Ax_{\rm opt}-b}^2
  \right]
  \ge 1-\delta .
\]
\end{corollary}


\begin{formalized}
\leanfit{abs_entry_le_frobNorm}
\leanfit{lsQRSolveBackwardSolverDx}
\leanfit{lsQRSolveBackwardSolverDx_nonneg}
\leanfit{gram_forward_error_certificate_of_ls_qr_solve_backward_error}
\leanfit{leverageTraceProbability_eventProb_fl_rowSampleLSObjective_le_one_add_eta_of_augmentedSpan_sample_budget_of_objective_error_and_ls_qr_backward_error_solver}
\end{formalized}

\begin{corollary}[High-probability rounded sampled-row solver from stored QR off-diagonal control]
Use the same equation~(8) sampled least-squares setting and Bennett
sample-budget assumptions.  For each sample trace \(\omega\), suppose the
rounded sampled/scaled matrix and right-hand side are fed to a stored
Householder QR process whose final top block is
\[
  R_\omega=\widehat A_{\omega,n}(1{:}n,1{:}n),
  \qquad
  c_\omega=\widehat b_{\omega,n}(1{:}n),
\]
and whose reported coefficient vector is
\[
  \widehat x(\omega)=\operatorname{flBackSub}(R_\omega,c_\omega).
\]
Assume the stored process satisfies the visible packaged route-1 invariant
\(\lean{StoredQROffDiagonalControlInvariant}\) for every sample trace, together
with signed-alpha Householder steps and the final QR Gram/RHS radius equations.
Equivalently, one may supply the source-shaped off-diagonal-control fields and
derive the packaged invariant locally.  Then the local stored-QR theorem supplies
\[
  \lean{LSQRSolveBackwardError}\bigl(G_\omega,g_\omega,
    \widehat x(\omega),c^G_\omega,c^g_\omega\bigr).
\]
Consequently the preceding QR-spec corollary applies: if the same
support-event objective-budget condition holds with the induced solver
certificate, then
\[
  \mathbb P\!\left[
    \Norm{A\widehat x(\omega)-b}^2
    \le
    (1+\eta)\Norm{Ax_{\rm opt}-b}^2
  \right]\ge 1-\delta .
\]
\end{corollary}


\begin{formalized}
\leanfit{storedQRFinalR}
\leanfit{storedQRFinalTopRhs}
\leanfit{storedQRBackSubSolution}
\leanfit{leverageTraceProbability_eventProb_fl_rowSampleLSObjective_le_one_add_eta_of_augmentedSpan_sample_budget_of_objective_error_and_stored_qr_offDiagonalControl_solver}
\leanfit{leverageTraceProbability_eventProb_fl_rowSampleLSObjective_le_one_add_eta_of_augmentedSpan_sample_budget_of_objective_error_and_stored_qr_offDiagonalControl_solver_of_diagDominant_finiteMaxSmallness}
\leanfit{leverageTraceProbability_eventProb_fl_rowSampleLSObjective_le_one_add_eta_of_augmentedSpan_sample_budget_of_objective_error_and_stored_qr_sourceOffDiagonalControl_solver}
\leanfit{leverageTraceProbability_eventProb_fl_rowSampleLSObjective_le_one_add_eta_of_augmentedSpan_sample_budget_of_objective_error_and_stored_qr_leadingBlock_det_ne_zero_normSqBudget_offdiag_product_solver}
\leanfit{leverageTraceProbability_eventProb_fl_rowSampleLSObjective_le_one_add_eta_of_augmentedSpan_sample_budget_of_objective_error_and_stored_qr_rowBudgetControl_solver}
\leanfit{leverageTraceProbability_eventProb_fl_rowSampleLSObjective_le_one_add_eta_of_augmentedSpan_sample_budget_of_objective_error_and_stored_qr_rowBudgetControl_kappaInf_dualBudget_solver}
\leanfit{leverageTraceProbability_eventProb_fl_rowSampleLSObjective_le_one_add_eta_of_augmentedSpan_sample_budget_of_objective_error_and_stored_qr_rowBudgetControl_globalProduct_kappaInf_dualBudget_solver}
\leanfit{leverageTraceProbability_eventProb_fl_rowSampleLSObjective_le_one_add_eta_of_augmentedSpan_sample_budget_of_objective_error_and_stored_qr_rowMaxDiagDefect_globalProduct_solver}
\end{formalized}

\begin{corollary}[Source-shaped stored QR off-diagonal-control certificate]
The packaged invariant in the preceding corollary may be discharged from the
primitive source-shaped certificate
\(\lean{StoredQRSourceOffDiagonalControl}\): local upper-triangular leading
blocks, nonzero displayed leading diagonals, row-wise off-diagonal domination,
and the stored-sequence compact-product inequality.  Under these source-shaped
fields, the same high-probability rounded sampled-row objective guarantee holds.
\end{corollary}


\begin{formalized}
\leanfit{StoredQROffDiagonalControlInvariant.of_sourceOffDiagonalControl}
\leanfit{LSQRSolveBackwardError.of_stored_trailing_householder_sequence_topBlock_fl_backSub_gamma_bound_explicitCompactBudget_of_signed_alpha_offDiagonalControl}
\leanfit{LSQRSolveBackwardError.of_stored_trailing_householder_sequence_topBlock_fl_backSub_gamma_bound_explicitCompactBudget_of_signed_alpha_sourceOffDiagonalControl}
\leanfit{leverageTraceProbability_eventProb_fl_rowSampleLSObjective_le_one_add_eta_of_augmentedSpan_sample_budget_of_objective_error_and_stored_qr_sourceOffDiagonalControl_solver}
\end{formalized}

\begin{lemma}[Rectangular perturbed normal equations feed the local QR specification]
Let \(A\in\mathbb R^{m\times n}\), \(b\in\mathbb R^m\), and let
\(\Delta A,\Delta b\) be rectangular data perturbations.  Define
\[
  G=A^T A,\qquad g=A^T b,
\]
and the induced Gram and right-hand-side perturbations
\[
  \Delta G=(A+\Delta A)^T(A+\Delta A)-A^T A,
  \qquad
  \Delta g=(A+\Delta A)^T(b+\Delta b)-A^T b .
\]
If a computed vector \(\widehat x\) satisfies the perturbed rectangular normal
equations
\[
  (A+\Delta A)^T(A+\Delta A)\widehat x
  =
  (A+\Delta A)^T(b+\Delta b),
\]
then it satisfies the local perturbed Gram system
\[
  (G+\Delta G)\widehat x=g+\Delta g .
\]
Consequently, if \(\|\Delta G\|_F\le c_G\) and
\(|(\Delta g)_j|\le c_g\) for every \(j\), then
\(\widehat x\) satisfies the repository's
\(\lean{LSQRSolveBackwardError}\) specification with radii \(c_G,c_g\).
\end{lemma}


\begin{formalized}
\leanfit{rectLSGram}
\leanfit{rectLSRhs}
\leanfit{RectLSNormalEquations}
\leanfit{rectLSGramPerturbation}
\leanfit{rectLSRhsPerturbation}
\leanfit{rectLSNormalEquations_perturbed_to_gram_system}
\leanfit{LSQRSolveBackwardError.of_rectangular_perturbed_normal_equations}
\end{formalized}

\begin{lemma}[Explicit budgets for induced rectangular perturbations]
With the notation of the previous lemma, the induced perturbations satisfy the
entrywise identities
\[
  (\Delta G)_{jk}
  =
  \sum_{i=1}^m
  \bigl(A_{ij}(\Delta A)_{ik}
    +(\Delta A)_{ij}A_{ik}
    +(\Delta A)_{ij}(\Delta A)_{ik}\bigr),
\]
\[
  (\Delta g)_j
  =
  \sum_{i=1}^m
  \bigl(A_{ij}(\Delta b)_i
    +(\Delta A)_{ij}b_i
    +(\Delta A)_{ij}(\Delta b)_i\bigr).
\]
Consequently
\[
  |(\Delta G)_{jk}|
  \le
  \sum_{i=1}^m
  \bigl(|A_{ij}|\,|(\Delta A)_{ik}|
    +|(\Delta A)_{ij}|\,|A_{ik}|
    +|(\Delta A)_{ij}|\,|(\Delta A)_{ik}|\bigr),
\]
and the analogous bound holds for \(|(\Delta g)_j|\).  In particular, if
\(\|\Delta A\|_F\le c_A\) and \(\|\Delta b\|_2\le c_b\), then
\[
  |(\Delta G)_{jk}|
  \le
  \sum_{i=1}^m
  \bigl(|A_{ij}|c_A+c_A|A_{ik}|+c_A^2\bigr),
  \qquad
  |(\Delta g)_j|
  \le
  \sum_{i=1}^m
  \bigl(|A_{ij}|c_b+c_A|b_i|+c_Ac_b\bigr).
\]
The Frobenius norm of \(\Delta G\) is bounded by the Frobenius norm of either
displayed Gram-budget matrix.
\end{lemma}


\begin{formalized}
\leanfit{rectLSGramPerturbationEntryBudget}
\leanfit{rectLSRhsPerturbationEntryBudget}
\leanfit{rectLSGramPerturbationNormBudget}
\leanfit{rectLSRhsPerturbationNormBudget}
\leanfit{rectLSGramPerturbation_eq_sum}
\leanfit{rectLSRhsPerturbation_eq_sum}
\leanfit{rectLSGramPerturbation_abs_le_entryBudget}
\leanfit{rectLSRhsPerturbation_abs_le_entryBudget}
\leanfit{rectLSGramPerturbation_frobNorm_le_entryBudget}
\leanfit{rectLSGramPerturbation_abs_le_normBudget}
\leanfit{rectLSRhsPerturbation_abs_le_normBudget}
\leanfit{rectLSGramPerturbation_frobNorm_le_normBudget}
\leanfit{LSQRSolveBackwardError.of_rectangular_perturbed_normal_equations_normBudget}
\end{formalized}

\begin{lemma}[Rectangular one-step orthogonal-transformation accumulation]
Let \(A,\widehat A,\Delta A\in\mathbb R^{m\times n}\), let
\(Q,P\in\mathbb R^{m\times m}\) be orthogonal, and suppose
\[
  \widehat A=Q^T(A+\Delta A).
\]
If the next computed transformation is
\[
  A_{\mathrm{next}}=(P+\Delta P)\widehat A
  \qquad\text{with}\qquad
  \|\Delta P\|_F\le c_{\mathrm{step}},
\]
then there are an orthogonal \(Q'\in\mathbb R^{m\times m}\) and a perturbation
\(\Delta A'\in\mathbb R^{m\times n}\) such that
\[
  A_{\mathrm{next}}=(Q')^T(A+\Delta A')
\]
and
\[
  \|\Delta A'\|_F
  \le
  \|\Delta A\|_F
  +c_{\mathrm{step}}\|A+\Delta A\|_F .
\]
\end{lemma}


\begin{formalized}
\lean{Algorithms/QR/HouseholderQR.lean}
\leanfit{rect_orthogonal_sequence_one_step}
\end{formalized}

\begin{theorem}[Rectangular multi-step orthogonal-transformation accumulation]
Let \(A\in\mathbb R^{m\times n}\), let \(A_0=A\), and suppose that for
\(k=0,\ldots,r-1\),
\[
  A_{k+1}=(P_k+\Delta P_k)A_k,
  \qquad
  P_k^TP_k=I,
  \qquad
  \|\Delta P_k\|_F\le c,
\]
with \(c\ge0\).  Then there are an orthogonal
\(Q\in\mathbb R^{m\times m}\) and a perturbation
\(\Delta A\in\mathbb R^{m\times n}\) such that
\[
  A_r=Q^T(A+\Delta A)
\]
and
\[
  \|\Delta A\|_F
  \le
  \bigl((1+c)^r-1\bigr)\,\|A\|_F .
\]
\end{theorem}


\begin{formalized}
\lean{Algorithms/QR/HouseholderQR.lean}
\leanfit{rect_orthogonal_sequence_geometric}
\end{formalized}

\begin{theorem}[Right-hand-side orthogonal-transformation accumulation]
Let \(b\in\mathbb R^m\), let \(b_0=b\), and suppose that for
\(k=0,\ldots,r-1\),
\[
  b_{k+1}=(P_k+\Delta P_k)b_k,
  \qquad
  P_k^TP_k=I,
  \qquad
  \|\Delta P_k\|_F\le c,
\]
with \(c\ge0\).  Then there are an orthogonal
\(Q\in\mathbb R^{m\times m}\) and a perturbation
\(\Delta b\in\mathbb R^m\) such that
\[
  b_r=Q^T(b+\Delta b)
\]
and
\[
  \|\Delta b\|_2
  \le
  \bigl((1+c)^r-1\bigr)\,\|b\|_2 .
\]
\end{theorem}


\begin{formalized}
\lean{Algorithms/QR/HouseholderQR.lean}
\leanfit{orthogonal_vector_sequence_one_step}
\leanfit{orthogonal_vector_sequence_geometric}
\end{formalized}

\begin{theorem}[Common-\(Q\) matrix/right-hand-side accumulation]
Let \(A\in\mathbb R^{m\times n}\), \(b\in\mathbb R^m\), \(A_0=A\), and
\(b_0=b\).  Suppose that the same perturbed orthogonal transformations are
applied to both data blocks:
\[
  A_{k+1}=(P_k+\Delta P_k)A_k,\qquad
  b_{k+1}=(P_k+\Delta P_k)b_k,
\]
where \(P_k^TP_k=I\), \(\|\Delta P_k\|_F\le c\), and \(c\ge0\).  Then there
are one common orthogonal matrix \(Q\), a matrix perturbation \(\Delta A\), and
a vector perturbation \(\Delta b\) such that
\[
  A_r=Q^T(A+\Delta A),\qquad
  b_r=Q^T(b+\Delta b),
\]
with
\[
  \|\Delta A\|_F\le\bigl((1+c)^r-1\bigr)\|A\|_F,
  \qquad
  \|\Delta b\|_2\le\bigl((1+c)^r-1\bigr)\|b\|_2 .
\]
\end{theorem}


\begin{formalized}
\lean{Algorithms/QR/HouseholderQR.lean}
\leanfit{rect_orthogonal_matrix_vector_sequence_one_step}
\leanfit{rect_orthogonal_matrix_vector_sequence_geometric}
\end{formalized}

\begin{lemma}[Top-block QR solve gives transformed normal equations]
Let \(\widehat A\in\mathbb R^{m\times n}\) and
\(\widehat b\in\mathbb R^m\) be transformed QR data.  Suppose there are
\(R\in\mathbb R^{n\times n}\) and \(c\in\mathbb R^n\) such that
\[
  \widehat A_{i,*}=R_{i,*},\quad \widehat b_i=c_i
  \qquad\text{for every } i<n,
\]
and
\[
  \widehat A_{i,*}=0\qquad\text{for every }i\ge n.
\]
If \(\widehat x\) solves the top system
\[
  R\widehat x=c,
\]
then \(\widehat x\) satisfies the rectangular normal equations for
\((\widehat A,\widehat b)\):
\[
  \widehat A^T\widehat A\,\widehat x=\widehat A^T\widehat b .
\]
No condition is imposed on the lower entries of \(\widehat b\); they are the
residual block in the usual tall-QR least-squares solve and disappear because
the corresponding rows of \(\widehat A\) are zero.
\end{lemma}


\begin{formalized}
\lean{Algorithms/LeastSquares/LSQRSolve.lean}
\leanfit{RectLSNormalEquations.of_rowwise_normal}
\leanfit{RectLSNormalEquations.of_top_solve_zero_bottom}
\end{formalized}

\begin{lemma}[Rounded top-block solve gives perturbed transformed normal equations]
Assume the top QR block \(R\in\mathbb R^{n\times n}\) is nonsingular on the
diagonal and upper triangular in the sense required by the repository's
back-substitution theorem.  Let
\[
  \widehat x = \operatorname{flBackSub}(R,c).
\]
Then there is a perturbation \(\Delta R\) such that
\[
  |\Delta R_{ij}| \le \gamma_n |R_{ij}|
\]
and \(\widehat x\) satisfies the rectangular normal equations for the
perturbed transformed QR data
\[
  \widetilde A =
  \begin{bmatrix} R+\Delta R \\ 0 \end{bmatrix},
  \qquad
  \widetilde b_i=c_i\quad (i<n),
\]
with no condition on the lower entries of \(\widetilde b\).
\end{lemma}


\begin{formalized}
\lean{Algorithms/LeastSquares/LSQRSolve.lean}
\leanfit{rectTopBlock}
\leanfit{rectTopBlock_top}
\leanfit{rectTopBlock_bottom}
\leanfit{RectLSNormalEquations.exists_topBlock_of_fl_backSub}
\end{formalized}

\begin{lemma}[Orthogonal handoff for rectangular normal equations]
Let \(U\in\mathbb R^{m\times m}\) be orthogonal, let
\(\widehat A=UA\), and let \(\widehat b=Ub\).  Then
\[
  \widehat A^T\widehat A=A^TA,
  \qquad
  \widehat A^T\widehat b=A^Tb.
\]
Consequently, if \(\widehat x\) satisfies the exact normal equations for
\((\widehat A,\widehat b)\), then \(\widehat x\) satisfies the exact normal
equations for \((A,b)\).
\end{lemma}


\begin{formalized}
\lean{Algorithms/LeastSquares/LSQRSolve.lean}
\leanfit{rectLSGram_matMulRectLeft_orthogonal}
\leanfit{rectLSRhs_matMulRectLeft_orthogonal}
\leanfit{RectLSNormalEquations.of_orthogonal_left}
\end{formalized}

\begin{lemma}[Pulling the rounded top-block perturbation back to the original data]
Assume that transformed QR data are represented with one common orthogonal
factor,
\[
  \widehat A=Q^T(A+\Delta A),\qquad
  \widehat b=Q^T(b+\Delta b),
\]
and that \(\widehat A\) has top block \(R\) and zero lower block.  If
\(\widehat x=\operatorname{flBackSub}(R,c)\), then there are a triangular-solve
perturbation \(\Delta R\) and an original-data matrix perturbation
\(\Delta A_{\rm tot}\) such that
\[
  |\Delta R_{ij}|\le\gamma_n |R_{ij}|,
  \qquad
  \Delta A_{\rm tot}
    = \Delta A + Q\begin{bmatrix}\Delta R\\0\end{bmatrix},
\]
\[
  \|\Delta A_{\rm tot}\|_F
  \le
  \|\Delta A\|_F+
  \left\|\begin{bmatrix}\Delta R\\0\end{bmatrix}\right\|_F,
\]
and \(\widehat x\) satisfies the rectangular normal equations for the
original perturbed data:
\[
  (A+\Delta A_{\rm tot})^T(A+\Delta A_{\rm tot})\widehat x
  =
  (A+\Delta A_{\rm tot})^T(b+\Delta b).
\]
\end{lemma}


\begin{formalized}
\lean{Algorithms/LeastSquares/LSQRSolve.lean}
\leanfit{rectTopBlock_add}
\leanfit{RectLSNormalEquations.exists_original_perturbation_of_commonQ_topBlock_fl_backSub}
\end{formalized}

\begin{corollary}[Gamma-budget version of the pulled-back top-block theorem]
Under the hypotheses of the previous lemma, the abstract top-block radius can
be discharged using the triangular-solve entrywise budget:
\[
  \left\|\begin{bmatrix}\Delta R\\0\end{bmatrix}\right\|_F
  \le
  \gamma_n
  \left\|\begin{bmatrix}R\\0\end{bmatrix}\right\|_F .
\]
Consequently the original-data perturbation can be chosen so that
\[
  \|\Delta A_{\rm tot}\|_F
  \le
  \|\Delta A\|_F+
  \gamma_n
  \left\|\begin{bmatrix}R\\0\end{bmatrix}\right\|_F .
\]
\end{corollary}


\begin{formalized}
\lean{Analysis/MatrixAlgebra.lean}
\leanfit{frobNormSqRect_abs}
\leanfit{frobNormRect_abs}
\lean{Algorithms/LeastSquares/LSQRSolve.lean}
\leanfit{rectTopBlock_frobNorm_perturb_bound}
\leanfit{rectTopBlock_frobNorm_perturb_bound_of_gamma}
\leanfit{RectLSNormalEquations.exists_original_perturbation_of_commonQ_topBlock_fl_backSub_gamma_bound}
\end{formalized}

\begin{corollary}[Supplied-transform rectangular QR route into the local solver specification]
Let \(A\in\mathbb R^{m\times n}\), \(b\in\mathbb R^m\), and suppose a
floating-point orthogonal-transformation pipeline produces transformed data
\(\widehat A,\widehat b\) with one common exact orthogonal representation
\[
  \widehat A = Q^T(A+\Delta A),
  \qquad
  \widehat b = Q^T(b+\Delta b).
\]
Assume that the final transformed matrix has QR top-block form
\[
  \widehat A =
  \begin{bmatrix}
    R\\0
  \end{bmatrix},
\]
that the first \(n\) entries of \(\widehat b\) are \(c\), and that \(R\) is
upper triangular with nonzero diagonal.  If back-substitution is performed as
\[
  \widehat x=\operatorname{flBackSub}(R,c),
\]
then the local QR least-squares backward-error specification holds with
matrix perturbation radius
\[
  c_A
  =
  \|\Delta A\|_F+
  \gamma_n
  \left\|\begin{bmatrix}R\\0\end{bmatrix}\right\|_F
\]
and right-hand-side radius \(c_b=\|\Delta b\|_2\), provided the displayed
radii are large enough to satisfy the induced Gram and right-hand-side
norm-budget hypotheses of the local specification.

In particular, if \(\widehat A_r,\widehat b_r\) are produced by a supplied
sequence of \(r\) perturbed orthogonal transformations whose per-step
Frobenius radius is \(c\), then the same conclusion holds with
\[
  c_A =
  \bigl((1+c)^r-1\bigr)\|A\|_F
  +
  \gamma_n
  \left\|\begin{bmatrix}R\\0\end{bmatrix}\right\|_F,
  \qquad
  c_b =
  \bigl((1+c)^r-1\bigr)\|b\|_2 .
\]
\end{corollary}


\begin{formalized}
\lean{Algorithms/LeastSquares/LSQRSolve.lean}
\leanfit{LSQRSolveBackwardError.of_commonQ_topBlock_fl_backSub_gamma_bound_normBudget}
\leanfit{LSQRSolveBackwardError.of_rect_orthogonal_sequence_topBlock_fl_backSub_gamma_bound_normBudget}
\end{formalized}

\paragraph{Scope.}
This corollary is intentionally a supplied-transform theorem.  It does not
claim that a concrete implementation of Householder QR has been formalized.
The remaining red bottleneck is the algorithm-level theorem that constructs
the supplied reflectors/updates, proves the final \([R;0]\) shape and
transformed right-hand side, and discharges the per-step floating-point
recurrences from actual rounded operations.

\begin{lemma}[Exact embedded Householder reflector with a zero prefix]
Let
\[
  P = I-\beta vv^T
\]
be the full-size Householder reflector represented in the repository.  If
\[
  v_i=0\qquad\text{for every }i<k,
\]
then \(P\) agrees with the identity on the first \(k\) rows and columns:
\[
  P_{ij}=\delta_{ij}\quad (i<k),
  \qquad
  P_{ij}=\delta_{ij}\quad (j<k).
\]
Consequently, for every vector \(x\) and every matrix \(A\),
\[
  (Px)_i=x_i\quad (i<k),
  \qquad
  (PA)_{ij}=A_{ij}\quad (i<k).
\]
The matrix statement is formalized both for square matrices and for
rectangular matrices.
\end{lemma}


\begin{formalized}
\lean{Algorithms/QR/HouseholderSpec.lean}
\leanfit{householder_row_eq_id_of_zero_prefix}
\leanfit{householder_col_eq_id_of_zero_prefix}
\leanfit{matMulVec_householder_eq_self_of_zero_prefix}
\leanfit{matMul_householder_eq_self_row_of_zero_prefix}
\leanfit{matMulRectLeft_householder_eq_self_row_of_zero_prefix}
\end{formalized}

\paragraph{Role in the bottleneck.}
This is a genuine exact foundation for the concrete rectangular Householder
route: it proves that a trailing reflector does not disturb already-finished
rows/columns.  The next lemma closes the exact active-column zeroing algebra.

\begin{lemma}[Exact active-column zeroing for the constructed Householder reflector]
Let \(x\in\mathbb R^m\), choose a pivot index \(p\), and define
\[
  v = x-\alpha e_p,\qquad
  \beta = \frac{2}{v^T v}.
\]
Assume
\[
  \alpha^2=\|x\|_2^2,
  \qquad
  v^T v \ne 0.
\]
Then the exact Householder reflector \(P=I-\beta vv^T\) maps the active
column vector to a single pivot entry:
\[
  Px=\alpha e_p.
\]
Equivalently,
\[
  (Px)_p=\alpha,\qquad (Px)_i=0\quad(i\ne p).
\]
\end{lemma}


\begin{formalized}
\lean{Algorithms/QR/HouseholderSpec.lean}
\leanfit{householderActiveVector}
\leanfit{householderBeta}
\leanfit{householderActiveVector_inner_x}
\leanfit{householderActiveVector_inner_self}
\leanfit{householderActiveVector_inner_self_eq_two_inner_x}
\leanfit{householderBeta_mul_activeVector_inner_x}
\leanfit{matMulVec_householder_activeVector_eq_alpha_basis}
\leanfit{matMulVec_householder_activeVector_eq_zero_of_ne}
\end{formalized}

\paragraph{Role in the bottleneck.}
This closes the exact active-column algebra for route A.  It is still not a
floating-point construction of the reflector and not a full QR implementation:
the QR-specific construction also needs the reflector vector to have a zero
prefix.  The next theorem closes that exact trailing-shape point.

\begin{definition}[Trailing Householder vector for rectangular QR]
For a pivot index \(p\), split a vector \(x\in\mathbb R^m\) as
\[
  x=x_{<p}+x_{\ge p},
  \qquad
  (x_{<p})_i =
  \begin{cases}x_i,&i<p,\\0,&i\ge p,\end{cases}
  \qquad
  (x_{\ge p})_i =
  \begin{cases}0,&i<p,\\x_i,&i\ge p.\end{cases}
\]
The QR trailing active vector is
\[
  v=x_{\ge p}-\alpha e_p.
\]
Thus \(v_i=0\) for all \(i<p\).  This is the source-faithful rectangular QR
construction; it differs from the full-vector construction \(x-\alpha e_p\)
unless the prefix of \(x\) is already zero.
\end{definition}

\begin{lemma}[Exact trailing-column Householder zeroing]
Let \(v=x_{\ge p}-\alpha e_p\), let
\[
  \beta=\frac{2}{v^Tv},
\]
and assume
\[
  \alpha^2=\|x_{\ge p}\|_2^2,
  \qquad
  v^Tv\ne0.
\]
For \(P=I-\beta vv^T\),
\[
  (Px)_i =
  \begin{cases}
    x_i, & i<p,\\
    \alpha, & i=p,\\
    0, & i>p.
  \end{cases}
\]
\end{lemma}


\begin{theorem}[Exact trailing Householder loop shape]
Let \(m\ge n\).  Suppose an exact rectangular Householder recurrence applies,
at step \(k<n\), the trailing reflector built from the current pivot column
\(\widehat A_k(:,k)\).  Assume the corresponding normalization hypotheses
\[
  \alpha_k^2=\|(\widehat A_k(:,k))_{\ge k}\|_2^2,
  \qquad
  v_k^Tv_k\ne0
\]
for each step.  Then after \(n\) steps,
\[
  j<i \quad\Longrightarrow\quad \widehat A_n(i,j)=0.
\]
Equivalently, the final transformed matrix is lower-zero/trapezoidal in the
QR sense.
\end{theorem}

\begin{corollary}[Top-block shape facts from the lower-zero invariant]
Define
\[
  R_{ij}=\widehat A_n(i,j)\quad(1\le i,j\le n),
  \qquad
  (c_{\rm top})_i=(\widehat b_n)_i\quad(1\le i\le n).
\]
If \(\widehat A_n(i,j)=0\) whenever \(j<i\), then
\[
  \widehat A_n=\begin{bmatrix}R\\0\end{bmatrix},
  \qquad
  (\widehat b_n)_{1:n}=c_{\rm top},
  \qquad
  R_{ij}=0\quad(j<i).
\]
The nonzero diagonal condition for triangular back substitution is not a
consequence of this shape theorem; it remains a visible rank/nonbreakdown
assumption unless separately proved from a rank hypothesis.
\end{corollary}

\begin{definition}[Stored rounded Householder panel and RHS steps]
For a pivot step \(k\), let \(\mathrm{raw}_k(A)\) denote the compact rounded
Householder panel update computed by the dot--scale--subtract primitive.  The
stored panel step is
\[
  A_{k+1}(i,j)=
  \begin{cases}
    A_k(i,j), & j<k,\\
    0,        & j=k \text{ and } i>k,\\
    \mathrm{raw}_k(A_k)(i,j), & \text{otherwise}.
  \end{cases}
\]
For the right-hand side, the stored step preserves rows above the active pivot
and applies the compact update on the active tail:
\[
  b_{k+1}(i)=
  \begin{cases}
    b_k(i), & i<k,\\
    \mathrm{raw}_k(b_k)(i), & i\ge k.
  \end{cases}
\]
\end{definition}

\begin{theorem}[Stored rounded Householder loop shape]
Suppose \(A_{k+1}\) is obtained from \(A_k\) by the stored panel step above for
each \(k<n\).  Then the final stored matrix satisfies
\[
  j<i \quad\Longrightarrow\quad A_n(i,j)=0.
\]
\end{theorem}

\begin{corollary}[Stored top-block shape facts]
Let
\[
  R_{ij}=A_n(i,j)\quad(1\le i,j\le n),
  \qquad
  (c_{\rm top})_i=(b_n)_i\quad(1\le i\le n).
\]
For a stored rounded loop,
\[
  A_n=\begin{bmatrix}R\\0\end{bmatrix},
  \qquad
  (b_n)_{1:n}=c_{\rm top},
  \qquad
  R_{ij}=0\quad(j<i).
\]
This is a storage-shape theorem only.  The perturbation contract for one stored
step is stated next, under explicit preservation, zeroing, and budget
hypotheses.  Neither statement proves nonzero diagonal pivots.
\end{corollary}


\begin{formalized}
\lean{Algorithms/QR/HouseholderSpec.lean}
\leanfit{householderPrefixPart}
\leanfit{householderTrailingPart}
\leanfit{householderTrailingNorm2Sq}
\leanfit{householderTrailingActiveVector}
\leanfit{householderPrefixPart_add_trailingPart}
\leanfit{householderPrefixPart_support}
\leanfit{householderTrailingActiveVector_zero_prefix}
\leanfit{matMulVec_householder_eq_self_of_zero_prefix_support}
\leanfit{matMulVec_householder_trailingActiveVector_eq_prefix_alpha_zero}
\leanfit{matMulVec_householder_trailingActiveVector_eq_zero_of_pivot_lt}
\lean{Algorithms/QR/HouseholderQR.lean}
\leanfit{exact_trailing_householder_sequence_lower_zero}
\leanfit{rectangular_topBlock_shape_facts_of_lower_zero}
\lean{Algorithms/QR/HouseholderApply.lean}
\leanfit{fl_householderStoredPanelStep}
\leanfit{fl_householderStoredRhsStep}
\lean{Algorithms/QR/HouseholderQR.lean}
\leanfit{fl_householderStoredPanel_sequence_lower_zero}
\leanfit{fl_householderStoredPanel_sequence_topBlock_shape_facts}
\end{formalized}

\begin{theorem}[Stored Householder step satisfies the columnwise perturbation contract]
Fix one pivot step \(k\).  Let \(P=I-\beta vv^T\) be the exact Householder
reflector and let \(\widehat A,\widehat b\) be computed by the stored rounded
panel and right-hand-side steps above.  Assume:
\[
  P^TP=I,
\]
completed columns are exactly preserved by the reflector,
\[
  j<k \Longrightarrow (PA_{\cdot j})_i=A_{ij},
\]
the entries stored as zero in the pivot column are exact zeros of the reflected
column,
\[
  i>k \Longrightarrow (PA_{\cdot k})_i=0,
\]
the right-hand-side prefix is exactly preserved,
\[
  i<k \Longrightarrow (Pb)_i=b_i,
\]
and the explicit compact Householder budgets satisfy
\[
  \|B_j\|_2\le c\|A_{\cdot j}\|_2,\qquad
  \|B_b\|_2\le c\|b\|_2 .
\]
Here \(B_j\) and \(B_b\) are the deterministic componentwise budgets produced
by the rounded dot--scale--subtract primitive.  Then
\[
  \mathrm{HouseholderColumnwisePanelAppError}
     (P,A,\widehat A,b,\widehat b;c)
\]
holds.
\end{theorem}


\begin{formalized}
\lean{Algorithms/QR/HouseholderApply.lean}
\leanfit{householderCompactComponentBudget_nonneg}
\leanfit{fl_householderStoredRhsStep_componentwise_error_bound}
\leanfit{fl_householderStoredRhsStep_forward_error_bound}
\leanfit{fl_householderStoredPanelStep_column_componentwise_error_bound}
\leanfit{fl_householderStoredPanelStep_column_forward_error_bound}
\leanfit{fl_householderStoredPanelStep_HouseholderColumnwisePanelAppError_of_budget}
\end{formalized}

\begin{corollary}[One stored trailing Householder QR step]
Let \(p\) be the pivot row and let \(v\) be the trailing Householder vector
formed from the current pivot column.  Suppose the current stored matrix has
the lower-zero invariant for all completed columns:
\[
  j<k,\quad j<i \quad\Longrightarrow\quad A_{ij}=0.
\]
Assume the Householder denominator is nonzero and the compact component budgets
for this step satisfy the same domination inequalities as in the previous
theorem.  Then the stored rounded trailing Householder step satisfies
\[
  \mathrm{HouseholderColumnwisePanelAppError}
     (P,A,\widehat A,b,\widehat b;c).
\]
\end{corollary}


\begin{formalized}
\lean{Algorithms/QR/HouseholderQR.lean}
\leanfit{fl_householderStoredTrailingPanelStep_HouseholderColumnwisePanelAppError_of_budget}
\end{formalized}

\begin{theorem}[Stored trailing Householder QR sequence]
Let \(m\ge n\), \(A_0=A\), and \(b_0=b\).  For each pivot \(k<n\), form the
trailing Householder vector from the current stored pivot column, set
\(\beta_k=2/(v_k^Tv_k)\), and define the stored rounded panel and right-hand
side recurrences
\[
  A_{k+1}=\mathrm{StoredPanelStep}(A_k,k,v_k,\beta_k),
  \qquad
  b_{k+1}=\mathrm{StoredRhsStep}(b_k,k,v_k,\beta_k).
\]
Assume that the Householder denominator \(v_k^Tv_k\) is nonzero, that the
chosen sign parameter satisfies \(\alpha_k^2=\sum_{i\ge k} A_{k,i k}^2\), and
that the compact component budgets for every stored panel column and right-hand
side update are dominated by a common constant \(c\ge0\).  Then there exist an
orthogonal matrix \(Q\), a data perturbation \(\Delta A\), and a right-hand-side
perturbation \(\Delta b\) such that
\[
  A_n = Q^T(A+\Delta A),\qquad
  b_n = Q^T(b+\Delta b),
\]
and, for every column \(j\),
\[
  \|\Delta A_{:j}\|_2
    \le \bigl((1+c)^n-1\bigr)\|A_{:j}\|_2,
  \qquad
  \|\Delta b\|_2
    \le \bigl((1+c)^n-1\bigr)\|b\|_2 .
\]
\end{theorem}


\begin{formalized}
\lean{Algorithms/QR/HouseholderQR.lean}
\leanfit{fl_householderStoredTrailingPanel_rect_orthogonal_columnwise_vector_sequence_geometric}
\end{formalized}

\begin{theorem}[Higham columnwise stored QR factorization]
Under the hypotheses of the stored trailing Householder QR sequence theorem,
let
\[
  R_{ij}=(A_n)_{ij},\qquad
  c_i=(b_n)_i,\qquad 0\le i,j<n,
\]
denote the top block of the final stored matrix and the top part of the final
stored right-hand side.  Then there exist an orthogonal matrix \(Q\), a data
perturbation \(\Delta A\), and a right-hand-side perturbation \(\Delta b\) such
that
\[
  A_n=Q^T(A+\Delta A),\qquad b_n=Q^T(b+\Delta b),
\]
\[
  \|\Delta A_{:j}\|_2
    \le \bigl((1+c)^n-1\bigr)\|A_{:j}\|_2,\qquad
  \|\Delta b\|_2
    \le \bigl((1+c)^n-1\bigr)\|b\|_2 .
\]
Moreover the final stored matrix has the block form
\[
  A_n=\begin{bmatrix}R\\0\end{bmatrix},
\]
the vector \(c\) is exactly the first \(n\) entries of \(b_n\), and \(R\) is
upper triangular.
\end{theorem}


\begin{formalized}
\lean{Algorithms/QR/HouseholderQR.lean}
\leanfit{fl_householderStoredTrailingPanel_higham_columnwise_factorization}
\end{formalized}

\begin{corollary}[Stored trailing Householder QR solve certificate]
Let \(R\) be the \(n\times n\) top block of the final stored matrix
\(A_n\), and let \(c_{\mathrm{top}}\) be the first \(n\) entries of the final
stored right-hand side \(b_n\).  Under the hypotheses of the stored trailing
sequence theorem, assume in addition that \(R\) has nonzero diagonal, that the
top-block arithmetic is gamma-valid, and that the induced Gram/RHS perturbation
budgets are bounded by \(c_G,c_g\).  Then the computed triangular solve
\[
  \widehat x = \mathrm{fl\_backSub}(R,c_{\mathrm{top}})
\]
satisfies the repository's least-squares QR backward-error specification
\[
  \mathrm{LSQRSolveBackwardError}
    (A^TA,A^Tb,\widehat x;c_G,c_g).
\]
\end{corollary}


\begin{formalized}
\lean{Algorithms/LeastSquares/LSQRSolve.lean}
\leanfit{LSQRSolveBackwardError.of_stored_trailing_householder_sequence_topBlock_fl_backSub_gamma_bound_normBudget}
\end{formalized}

\begin{lemma}[Pivot-budget nonbreakdown for stored trailing QR]
At pivot \(k\), let \(v_k\) be the trailing Householder vector, let
\(\beta_k=2/(v_k^Tv_k)\), and let \(B_k\) be the componentwise compact-update
budget for the stored diagonal entry.  If
\[
  B_k < |\alpha_k|,
\]
then the stored rounded diagonal entry written at that pivot is nonzero.  If
this inequality holds for every pivot, then the final stored top block \(R\)
has nonzero diagonal.
\end{lemma}

\begin{corollary}[Stored QR solve certificate with pivot-budget nonbreakdown]
In the stored trailing Householder QR solve certificate above, the hypothesis
that \(R\) has nonzero diagonal may be replaced by the per-pivot inequalities
\[
  B_k < |\alpha_k|,\qquad 0\le k<n.
\]
Under those inequalities, the same computed vector
\[
  \widehat x=\mathrm{fl\_backSub}(R,c_{\mathrm{top}})
\]
satisfies
\[
  \mathrm{LSQRSolveBackwardError}
    (A^TA,A^Tb,\widehat x;c_G,c_g).
\]
\end{corollary}


\begin{formalized}
\lean{Algorithms/QR/HouseholderQR.lean}
\leanfit{fl_householderStoredTrailingPanelStep_diag_nonzero_of_budget_lt_abs_alpha}
\leanfit{fl_householderStoredPanel_sequence_diag_nonzero_of_step_diag_nonzero}
\leanfit{fl_householderStoredTrailingPanel_sequence_diag_nonzero_of_budget_lt_abs_alpha}
\lean{Algorithms/LeastSquares/LSQRSolve.lean}
\leanfit{LSQRSolveBackwardError.of_stored_trailing_householder_sequence_topBlock_fl_backSub_gamma_bound_normBudget_of_pivot_error_lt_abs_alpha}
\end{formalized}

\begin{lemma}[Pivot-value nonbreakdown for the Householder denominator]
Let \(x\) be the current pivot column and
\[
  v=x_{\mathrm{tail}}-\alpha_k e_k
\]
be the trailing Householder vector.  If the active pivot entry differs from
the chosen sign parameter,
\[
  x_k\ne \alpha_k,
\]
then \(v^Tv\ne 0\).  Thus the denominator needed to form
\(\beta_k=2/(v^Tv)\) follows from a scalar pivot nonbreakdown condition rather
than from an unexplained sum-of-squares hypothesis.
\end{lemma}

\begin{corollary}[Stored QR solve certificate with pivot-value nonbreakdown]
In the stored trailing Householder QR solve certificate with pivot-budget
nonbreakdown, the Householder denominator assumptions may be replaced by
\[
  A_k(k,k)\ne \alpha_k,\qquad 0\le k<n,
\]
where \(A_k\) is the stored panel before pivot \(k\).  Together with
\(B_k<|\alpha_k|\), the computed vector
\[
  \widehat x=\mathrm{fl\_backSub}(R,c_{\mathrm{top}})
\]
satisfies the same
\[
  \mathrm{LSQRSolveBackwardError}
    (A^TA,A^Tb,\widehat x;c_G,c_g).
\]
\end{corollary}


\begin{formalized}
\lean{Algorithms/QR/HouseholderSpec.lean}
\leanfit{householderTrailingActiveVector_inner_self_ne_zero_of_pivot_ne_alpha}
\lean{Algorithms/QR/HouseholderQR.lean}
\leanfit{fl_householderStoredTrailingPanel_sequence_diag_nonzero_of_pivot_ne_alpha}
\lean{Algorithms/LeastSquares/LSQRSolve.lean}
\leanfit{LSQRSolveBackwardError.of_stored_trailing_householder_sequence_topBlock_fl_backSub_gamma_bound_normBudget_of_pivot_ne_alpha_and_pivot_error_lt_abs_alpha}
\end{formalized}

\begin{lemma}[Sign-choice nonbreakdown for the Householder denominator]
Let \(x\) be the current pivot column and suppose the Householder parameter
\(\alpha_k\) satisfies
\[
  \alpha_k^2=\|x_{\mathrm{tail}}\|_2^2,\qquad
  \|x_{\mathrm{tail}}\|_2^2>0,\qquad
  \alpha_k x_k\le 0 .
\]
Then \(x_k\ne \alpha_k\), and hence the trailing Householder vector
\[
  v=x_{\mathrm{tail}}-\alpha_k e_k
\]
satisfies \(v^Tv\ne 0\).
\end{lemma}

\begin{corollary}[Stored QR solve certificate with sign-choice nonbreakdown]
In the stored trailing Householder QR solve certificate with pivot-budget
nonbreakdown, the Householder denominator assumptions may be replaced by the
source-style local conditions
\[
  \alpha_k^2=\|A_k(k{:}m,k)\|_2^2,\qquad
  \|A_k(k{:}m,k)\|_2^2>0,\qquad
  \alpha_k A_k(k,k)\le 0,\qquad 0\le k<n .
\]
Together with \(B_k<|\alpha_k|\), the computed vector
\[
  \widehat x=\mathrm{fl\_backSub}(R,c_{\mathrm{top}})
\]
satisfies the same
\[
  \mathrm{LSQRSolveBackwardError}
    (A^TA,A^Tb,\widehat x;c_G,c_g).
\]
\end{corollary}


\begin{formalized}
\lean{Algorithms/QR/HouseholderSpec.lean}
\leanfit{householder_pivot_ne_alpha_of_trailingNorm2Sq_pos_mul_nonpos}
\leanfit{householderTrailingActiveVector_inner_self_ne_zero_of_trailingNorm2Sq_pos_mul_nonpos}
\lean{Algorithms/QR/HouseholderQR.lean}
\leanfit{fl_householderStoredTrailingPanel_sequence_diag_nonzero_of_trailingNorm_pos_mul_nonpos}
\lean{Algorithms/LeastSquares/LSQRSolve.lean}
\leanfit{LSQRSolveBackwardError.of_stored_trailing_householder_sequence_topBlock_fl_backSub_gamma_bound_normBudget_of_trailingNorm_pos_mul_nonpos_and_pivot_error_lt_abs_alpha}
\end{formalized}

\begin{lemma}[Row sorting preserves rectangular least-squares data]
Let \(\sigma\) be a permutation of the \(m\) row indices.  Define
\[
  A^\sigma_{ij}=A_{\sigma(i),j},
  \qquad
  b^\sigma_i=b_{\sigma(i)} .
\]
Then row sorting does not change the least-squares residual objective:
\[
  \|A^\sigma x-b^\sigma\|_2^2
    =\|Ax-b\|_2^2 .
\]
It also preserves the normal-equation data:
\[
  (A^\sigma)^TA^\sigma=A^TA,
  \qquad
  (A^\sigma)^Tb^\sigma=A^Tb .
\]
The same permutation invariance is formalized for rectangular Frobenius norms
and for rectangular matrix-vector multiplication.
\end{lemma}


\begin{formalized}
\lean{Analysis/MatrixAlgebra.lean}
\leanfit{vecPermute}
\leanfit{rectPermuteRows}
\leanfit{rectPermuteCols}
\leanfit{vecNorm2Sq_permute}
\leanfit{frobNormSqRect_permuteRows}
\leanfit{frobNormSqRect_permuteCols}
\leanfit{frobNormRect_permuteRows}
\leanfit{frobNormRect_permuteCols}
\leanfit{rectMatMulVec_permuteRows}
\lean{Algorithms/LeastSquares/LSQRSolve.lean}
\leanfit{rectLSGram_permuteRows}
\leanfit{rectLSRhs_permuteRows}
\lean{Algorithms/RandNLA/LeastSquaresSketch.lean}
\leanfit{lsResidual_permuteRows}
\leanfit{lsObjective_permuteRows}
\end{formalized}

\begin{lemma}[Column pivoting preserves rectangular least-squares data]
Let \(\pi\) be a permutation of the \(n\) column indices and define
\[
  A^\pi_{ij}=A_{i,\pi(j)} .
\]
Then column pivoting is a relabeling of the coefficient vector:
\[
  A^\pi x = A(x\circ\pi^{-1}),
  \qquad
  \|A^\pi x-b\|_2^2=\|A(x\circ\pi^{-1})-b\|_2^2 .
\]
The normal-equation data are relabeled by the same column permutation,
\[
  (A^\pi)^T A^\pi_{jk}=(A^TA)_{\pi(j),\pi(k)},
  \qquad
  (A^\pi)^T b_j=(A^Tb)_{\pi(j)} .
\]
Consequently, a solution or exact minimizer of the column-pivoted
least-squares problem maps back to the original coordinates by
\(x\mapsto x\circ\pi^{-1}\).
\end{lemma}


\begin{formalized}
\lean{Analysis/MatrixAlgebra.lean}
\leanfit{vecPermute_symm_vecPermute}
\leanfit{vecPermute_vecPermute_symm}
\leanfit{rectMatMulVec_permuteCols}
\lean{Algorithms/LeastSquares/LSQRSolve.lean}
\leanfit{rectLSGram_permuteCols}
\leanfit{rectLSRhs_permuteCols}
\leanfit{RectLSNormalEquations.of_permuteCols}
\lean{Algorithms/RandNLA/LeastSquaresSketch.lean}
\leanfit{lsResidual_permuteCols}
\leanfit{lsObjective_permuteCols}
\leanfit{IsLeastSquaresMinimizer.of_permuteCols}
\end{formalized}

\begin{lemma}[Combined row sorting and column pivoting for least squares]
Let \(\sigma\) be a row permutation and \(\pi\) a column permutation.  If
\[
  \widetilde A_{ij}=A_{\sigma(i),\pi(j)},\qquad
  \widetilde b_i=b_{\sigma(i)},
\]
then the row permutation disappears from the normal-equation data and the
column permutation only relabels coefficient coordinates:
\[
  \widetilde A^T\widetilde A_{jk}=(A^TA)_{\pi(j),\pi(k)},\qquad
  \widetilde A^T\widetilde b_j=(A^Tb)_{\pi(j)} .
\]
Moreover,
\[
  \|\widetilde A x-\widetilde b\|_2^2
  =
  \|A(x\circ\pi^{-1})-b\|_2^2 .
\]
Consequently, any solution of the row-sorted, column-pivoted normal equations,
and any exact minimizer of that row-sorted, column-pivoted least-squares
problem, maps back to the original coordinates by
\(x\mapsto x\circ\pi^{-1}\).
\end{lemma}


\begin{formalized}
\lean{Algorithms/LeastSquares/LSQRSolve.lean}
\leanfit{rectLSGram_permuteRowsCols}
\leanfit{rectLSRhs_permuteRowsCols}
\leanfit{RectLSNormalEquations.of_permuteRowsCols}
\lean{Algorithms/RandNLA/LeastSquaresSketch.lean}
\leanfit{lsResidual_permuteRowsCols}
\leanfit{lsObjective_permuteRowsCols}
\leanfit{IsLeastSquaresMinimizer.of_permuteRowsCols}
\end{formalized}

\begin{lemma}[Cox--Higham signed Householder denominator bound]
Let \(x_{\mathrm{tail}}\) be the active trailing pivot column segment and let
\[
  v=x_{\mathrm{tail}}-\alpha e_k .
\]
If
\[
  \alpha^2=\|x_{\mathrm{tail}}\|_2^2
  \qquad\text{and}\qquad
  \alpha x_k\le 0,
\]
then
\[
  2\|x_{\mathrm{tail}}\|_2^2 \le v^Tv .
\]
For the concrete signed Householder choice
\[
  \alpha=\mathrm{signedHouseholderAlpha}
    \bigl(\sqrt{\|x_{\mathrm{tail}}\|_2^2},x_k\bigr),
\]
the two hypotheses above are supplied by the local signed-alpha lemmas.
\end{lemma}


\begin{formalized}
\lean{Algorithms/QR/HouseholderSpec.lean}
\leanfit{householderTrailingActiveVector_inner_self_ge_two_trailingNorm2Sq_of_mul_nonpos}
\leanfit{householderTrailingActiveVector_inner_self_ge_two_trailingNorm2Sq_signed}
\end{formalized}

\begin{lemma}[Cox--Higham active-column pivot comparison]
Fix a Householder step \(k\) and define
\[
  \tau_j=\|A_{k:m,j}\|_2^2
\]
for each remaining column \(j\ge k\).  There exists a remaining column \(j_*\)
such that
\[
  \tau_j\le \tau_{j_*}\qquad(j\ge k).
\]
If \(x=A_{:,j_*}\), \(y=A_{:,j}\), and
\[
  v=x_{\mathrm{tail}}-\alpha e_k,
\]
then the pivot-max property and Cauchy--Schwarz give
\[
  |v^T y_{\mathrm{tail}}|
    \le \|v\|_2\,\|y_{\mathrm{tail}}\|_2
    \le \|v\|_2\,\|x_{\mathrm{tail}}\|_2 .
\]
\end{lemma}


\begin{formalized}
\lean{Algorithms/QR/HouseholderSpec.lean}
\leanfit{householderTrailingColumnNorm2Sq}
\leanfit{exists_householderTrailingColumnNorm2Sq_active_max}
\leanfit{abs_inner_householderTrailingActiveVector_trailingPart_le_vecNorm2_mul_sqrt}
\leanfit{abs_inner_householderTrailingActiveVector_column_le_vecNorm2_mul_sqrt_of_pivot_max}
\end{formalized}

\begin{lemma}[Cox--Higham scalar update and one-step row growth]
With the same signed Householder vector \(v\), exact coefficient
\[
  \beta=\frac{2}{v^Tv},
\]
and a remaining column \(y\), the previous two lemmas imply
\[
  |\beta\,v^T y_{\mathrm{tail}}|\le \sqrt 2 .
\]
Consequently, if a row has
\[
  |a_{ij}|\le B,\qquad |a_{ik}|\le B,
\]
then the corresponding off-pivot Householder update with
\(\phi=\beta v^Ty_{\mathrm{tail}}\) satisfies
\[
  |a_{ij}-\phi a_{ik}|\le (1+\sqrt 2)B .
\]
This is the single-step entry growth estimate used in Cox--Higham's proof of
equation (4.3).  It is not yet the multi-stage row-sorting or row-pivoting
accumulation theorem.
\end{lemma}


\begin{formalized}
\lean{Algorithms/QR/HouseholderSpec.lean}
\leanfit{abs_two_div_mul_le_sqrt_two_of_abs_le_sqrt_mul_sqrt}
\leanfit{abs_householderBeta_mul_inner_trailingPart_le_sqrt_two_of_mul_nonpos}
\leanfit{abs_householderBeta_mul_inner_trailingPart_le_sqrt_two_signed}
\leanfit{abs_householderBeta_mul_inner_column_le_sqrt_two_of_signed_pivot_max}
\leanfit{abs_sub_mul_le_one_add_sqrt_two_mul_bound}
\leanfit{abs_householder_signed_pivot_update_entry_le_one_add_sqrt_two_mul_row_bound}
\end{formalized}

\begin{lemma}[Cox--Higham row-sorting stage accumulation]
Let \(c=1+\sqrt2\).  Suppose that the initial row-size bounds are sorted
among the active rows, so that for every \(r\ge k\),
\[
  B_r^{(0)}\le B_k^{(0)} .
\]
If a fixed active entry grows at each previous Householder stage by at most the
one-step factor \(c\), then after \(k\) stages
\[
  |A_k(r,j)|\le c^k B_k^{(0)} .
\]
Equivalently, row sorting turns the one-step estimate
\(|A_{t+1}(r,j)|\le (1+\sqrt2)|A_t(r,j)|\) into the accumulated
Cox--Higham row bound used in the proof of Lemma 4.2.
\end{lemma}


\begin{formalized}
\lean{Algorithms/QR/HouseholderSpec.lean}
\leanfit{coxHighamGrowthFactor_nonneg}
\leanfit{scalar_growth_iterate_bound}
\leanfit{coxHigham_rowSorting_active_entry_bound_of_prior_growth}
\leanfit{coxHigham_rowSorting_active_entry_bound_of_stage_growth}
\end{formalized}

\begin{lemma}[Cox--Higham pivot-row active-tail norm step]
Let \(p=k\) be the current pivot row and suppose the next pivot-row entry in
column \(j\) is bounded by the Euclidean norm of the active tail of the
previous column,
\[
  |A_{k+1}(k,j)|\le \|A_k(k:m,j)\|_2 .
\]
If every active-tail entry of that column is bounded by \(B\), then
\[
  |A_{k+1}(k,j)|\le \sqrt m\,B .
\]
Consequently, under the same row-sorting scale as above,
\[
  |A_{k+1}(k,j)|
  \le \sqrt m\,(1+\sqrt2)^k B_k^{(0)} .
\]
This is the repository's ambient-\(\mathbb R^m\) form of the norm step in
Cox--Higham equations (4.4)--(4.5).  The sharper source factor
\(\sqrt{m-k+1}\) is not hidden; it remains a possible refinement once the
library has a dedicated active-tail dimension lemma.
\end{lemma}


\begin{formalized}
\lean{Analysis/MatrixAlgebra.lean}
\leanfit{vecNorm2_le_sqrt_card_mul_of_abs_le}
\lean{Algorithms/QR/HouseholderSpec.lean}
\leanfit{coxHigham_pivot_row_entry_bound_of_active_tail_entry_bound}
\leanfit{coxHigham_pivot_row_entry_bound_of_stage_entry_bound}
\end{formalized}

\begin{lemma}[Exact signed pivot-row same-reflector bridge]
Let \(p\) be the current pivot row and let the signed Householder reflector
\[
  P=I-\beta vv^T,\qquad
  v=A(:,k)_{\mathrm{tail}}-\alpha e_p,\qquad
  \alpha=-\operatorname{sign}(A_{pk})\|A(p:m,k)\|_2
\]
be formed from a pivot column \(k\) with positive active trailing norm.  Then
for any updated column \(j\),
\[
  |(PA(:,j))_p|\le \|A(p:m,j)\|_2 .
\]
Consequently, if the active entries of column \(j\) at stage \(p\) satisfy
\[
  |A_{ij}|\le (1+\sqrt2)^p B_p^{(0)},\qquad i\ge p,
\]
then
\[
  |(PA(:,j))_p|
  \le
  \sqrt m\,(1+\sqrt2)^p B_p^{(0)} .
\]
This supplies the exact same-reflector pivot-row field corresponding to
Cox--Higham equations (4.4)--(4.5), in the repository's ambient-dimension
form.
\end{lemma}


\begin{formalized}
\lean{Algorithms/QR/HouseholderSpec.lean}
\leanfit{householderBeta_mul_inner_self_eq_two}
\leanfit{abs_matMulVec_householder_pivot_le_trailing_vecNorm2_of_zero_prefix_orthogonal}
\leanfit{coxHigham_signed_pivot_row_update_abs_le_trailing_vecNorm2}
\leanfit{coxHigham_exact_signed_pivot_row_entry_bound_of_stage_entry_bound}
\end{formalized}

\begin{lemma}[Exact same-reflector row-growth bridge]
Let \(p\) be the current pivot row and \(k\) the current pivot column.  Form
the signed Householder vector
\[
  v=A(:,k)_{\mathrm{tail}}-\alpha e_p,\qquad
  \alpha=-\operatorname{sign}(A_{pk})\|A(p:m,k)\|_2,\qquad
  \beta=\frac{2}{v^Tv}.
\]
For any non-pivot active row \(r>p\), the exact matrix-vector update of column
\(j\) is the scalar Cox--Higham row update
\[
  \bigl((I-\beta vv^T)A(:,j)\bigr)_r
  =
  A_{rj}
  -
  \left(\beta\,v^T A(p:m,j)\right)A_{rk}.
\]
Consequently, if column \(k\) maximizes the active trailing column norm among
remaining columns, \(j\ge k\), and the active row entries obey
\[
  |A_{r\ell}|\le B,\qquad \ell\ge k,
\]
then
\[
  \left|\bigl((I-\beta vv^T)A(:,j)\bigr)_r\right|
  \le (1+\sqrt2)B .
\]
This is the one-step exact same-reflector field needed by the direct stored
rounded-panel recurrence for non-pivot active rows.
\end{lemma}


\begin{formalized}
\lean{Algorithms/QR/HouseholderApply.lean}
\leanfit{matMulVec_householder_signed_pivot_update_entry_eq}
\leanfit{coxHigham_exact_same_reflector_row_growth_of_signed_pivot_row_bound}
\end{formalized}

\begin{corollary}[One-step exact active-row signed-pivot bound]
In one signed pivoted Householder step, suppose the current row is active,
\(r\ge p\), the pivot column has positive active trailing norm, column \(k\)
maximizes the active trailing column norm among the remaining columns, and the
needed row and column active-entry bounds are bounded by \(B\ge0\).  Then
\[
  |(PA(:,j))_r|
  \le
  \max\{1+\sqrt2,\sqrt m\}\,B .
\]
For \(r>p\) this is the non-pivot Cox--Higham scalar row-growth bridge; for
\(r=p\) it is the signed pivot-row active-tail bridge.  The statement is a
one-step case split, not yet the concrete multi-stage sorted-loop theorem.
\end{corollary}


\begin{formalized}
\lean{Algorithms/QR/HouseholderApply.lean}
\leanfit{coxHigham_exact_signed_pivot_active_row_entry_bound_of_stage_bounds}
\end{formalized}

\begin{corollary}[Active-block budget supplies the one-step signed-pivot fields]
In the preceding one-step signed-pivot setting, suppose a single active-block
budget is available:
\[
  |A_{i\ell}|\le B
  \qquad
  (i\ge p,\ \ell\ge k),
  \qquad B\ge0 .
\]
If the pivot column has positive active trailing norm and maximizes the active
trailing column norm among the remaining columns, then for every active row
\(r\ge p\) and remaining column \(j\ge k\),
\[
  |(PA(:,j))_r|
  \le
  \max\{1+\sqrt2,\sqrt m\}\,B .
\]
Thus a block invariant of the form used in a sorted or pivoted loop supplies
the separate row and column hypotheses in the one-step signed-pivot theorem.
This is still a one-step field adapter: it does not prove that a particular
sorting policy propagates the active-block budget over all stages.
\end{corollary}


\begin{formalized}
\lean{Algorithms/QR/HouseholderApply.lean}
\leanfit{coxHigham_exactSignedPivotPanelStep_active_block_bound_of_stage_bound}
\end{formalized}

\begin{corollary}[Exact signed-pivot panel sequence with visible budgets]
Let \(A_t\) be the exact sequence generated by applying, at each stage \(t\),
the signed-pivot Householder panel step to the current matrix \(A_t\).  Let
\[
  c_m=\max\{1+\sqrt2,\sqrt m\}.
\]
Suppose that, for every active stage \(t<q\), the pivot column has positive
active trailing norm, maximizes the active trailing norm among the remaining
columns, the target column is remaining, and the row and column stage budgets
are bounded by \(B_t\ge0\).  If
\[
  c_m B_t\le B_{t+1}\qquad(0\le t<q),
  \qquad |A_0(r,j)|\le B_0,
\]
then
\[
  |A_q(r,j)|\le B_q .
\]
In particular, if the source sorting/pivoting invariant supplies
\(B_t=c_m^tB_0\), then
\[
  |A_q(r,j)|\le c_m^qB_0 .
\]
This is the concrete exact-loop budget propagation theorem needed by the
Cox--Higham route.  It still keeps the stage-budget/sorting fields visible; it
does not yet prove that a particular sorting policy supplies those fields, and
it is not the final QR/preconditioner theorem.
\end{corollary}


\begin{formalized}
\lean{Algorithms/QR/HouseholderSpec.lean}
\leanfit{coxHighamActiveRowGrowthFactor}
\leanfit{coxHighamActiveRowGrowthFactor_nonneg}
\lean{Algorithms/QR/HouseholderApply.lean}
\leanfit{exactSignedPivotHouseholderPanelStep}
\leanfit{coxHigham_exactSignedPivotPanelStep_active_entry_bound_of_stage_bounds}
\leanfit{coxHigham_exactSignedPivotPanel_sequence_active_entry_bound_of_stage_budgets}
\leanfit{coxHigham_exactSignedPivotPanel_sequence_active_entry_bound_of_geometric_stage_budgets}
\end{formalized}

\begin{corollary}[Exact signed-pivot sequence with active-block budgets]
The preceding sequence theorem can be stated using one active-block budget at
each stage.  Suppose that, for every stage \(t<q\),
\[
  |A_t(i,\ell)|\le B_t
  \qquad
  (i\ge p_t,\ \ell\ge k_t),
\]
and suppose the same positive active trailing norm, pivot maximality, stage
recurrence \(c_mB_t\le B_{t+1}\), and initial bound assumptions hold.  Then
\[
  |A_q(r,j)|\le B_q .
\]
In the geometric case, if the active block satisfies
\[
  |A_t(i,\ell)|\le c_m^tB_0
  \qquad
  (i\ge p_t,\ \ell\ge k_t),
\]
then
\[
  |A_q(r,j)|\le c_m^qB_0 .
\]
This is the exact sequence theorem in the shape used by a sorting or pivoting
invariant: one block bound supplies all row and column fields at every stage.
It still does not prove that a particular sorting policy supplies the displayed
active-block bounds.
\end{corollary}


\begin{formalized}
\lean{Algorithms/QR/HouseholderApply.lean}
\leanfit{coxHigham_exactSignedPivotPanel_sequence_active_block_bound_of_stage_budgets}
\leanfit{coxHigham_exactSignedPivotPanel_sequence_active_block_bound_of_geometric_stage_budgets}
\end{formalized}

\begin{theorem}[Active-block budgets from monotone active windows]
Let \(A_t\) be the exact signed-pivot panel sequence above and let
\(c_m=\max\{1+\sqrt2,\sqrt m\}\).  Suppose the initial matrix has the
entrywise bound
\[
  |A_0(i,\ell)|\le B_0
  \qquad (1\le i\le m,\ 1\le \ell\le n),
\]
with \(B_0\ge0\).  Suppose also that the active row starts \(p_t\) and
remaining-column starts \(k_t\) are monotone in \(t\), and that every stage
satisfies the visible positive active trailing-norm and column-pivot maximality
conditions.  Then, for every \(t\le q\),
\[
  |A_t(i,\ell)|\le c_m^tB_0
  \qquad (i\ge p_t,\ \ell\ge k_t).
\]
\end{theorem}


\begin{formalized}
\lean{Algorithms/QR/HouseholderApply.lean}
\leanfit{coxHigham_exactSignedPivotPanel_sequence_active_block_bound_of_initial_block_bound}
\end{formalized}

\begin{corollary}[Active-block budgets from nonzero active mass]
In the setting of the previous theorem, the positive active trailing-norm
hypothesis can be replaced by the following source-shaped nonbreakdown
condition: at each stage \(t<q\), some entry in the remaining active block is
nonzero,
\[
  \exists \ell\ge k_t,\ \exists i\ge p_t:\quad A_t(i,\ell)\ne0.
\]
Together with column-pivot maximality, this implies that the selected pivot
column has positive active trailing norm.  Hence the same conclusion holds:
\[
  |A_t(i,\ell)|\le c_m^tB_0
  \qquad (t\le q,\ i\ge p_t,\ \ell\ge k_t).
\]
This is still not a proof that a concrete implementation supplies the
nonzero-active-block and pivot-max fields; it removes only the direct
positive-norm assumption from the exact sequence theorem.
\end{corollary}


\begin{formalized}
\lean{Algorithms/QR/HouseholderSpec.lean}
\leanfit{householderTrailingColumnNorm2Sq_pos_of_pivot_max_exists_pos}
\leanfit{householderTrailingColumnNorm2Sq_pos_of_pivot_max_exists_active_entry_ne}
\lean{Algorithms/QR/HouseholderApply.lean}
\leanfit{coxHigham_exactSignedPivotPanel_sequence_active_block_bound_of_initial_block_bound_of_active_block_nonzero}
\end{formalized}

\begin{corollary}[Active-block budgets from an active max-pivot policy]
In the setting of the preceding corollary, the raw column-pivot maximality
field can be replaced by an explicit finite max-pivot policy.  In the current
exact-sequence formalism, the remaining-column start \(k_t\) denotes the
current column after pivoting/sorting has placed the selected column in active
position.  The policy hypothesis says that this current column is exactly the
output of \(\lean{householderActiveMaxPivotColumn}\) for the displayed active
block.  Under this policy and the same nonzero-active-block condition,
\[
  |A_t(i,\ell)|\le c_m^tB_0
  \qquad (t\le q,\ i\ge p_t,\ \ell\ge k_t).
\]
This is still an exact-sequence dependency.  It does not yet prove that the
rounded stored implementation performs the corresponding pivoting/sorting, nor
does it connect the result to the QR/preconditioner solve theorem.
\end{corollary}


\begin{formalized}
\lean{Algorithms/QR/HouseholderSpec.lean}
\leanfit{householderActiveMaxPivotColumn}
\leanfit{householderActiveMaxPivotColumn_ge}
\leanfit{householderActiveMaxPivotColumn_pivot_max}
\lean{Algorithms/QR/HouseholderApply.lean}
\leanfit{coxHigham_exactSignedPivotPanel_sequence_active_block_bound_of_initial_block_bound_of_active_max_pivot}
\end{formalized}

\begin{corollary}[Active-block budgets from positive active-block mass]
In the setting of the active max-pivot policy corollary, the remaining
nonzero-active-block hypothesis can be replaced by a scalar active-block mass
condition.  Define
\[
  \mathcal M_t
    = \sum_{i,\ell}
        \mathbf 1_{\{i\ge p_t,\ \ell\ge k_t\}}\; A_t(i,\ell)^2 .
\]
If \(\mathcal M_t>0\) at every active stage and the current column is selected
by \(\lean{householderActiveMaxPivotColumn}\), then
\[
  |A_t(i,\ell)|\le c_m^tB_0
  \qquad (t\le q,\ i\ge p_t,\ \ell\ge k_t).
\]
This is still an exact-sequence dependency: a concrete pivoted/sorted loop
must prove the positive active-block mass and policy fields.
\end{corollary}


\begin{formalized}
\lean{Algorithms/QR/HouseholderSpec.lean}
\leanfit{householderActiveBlockNorm2Sq}
\leanfit{exists_active_entry_ne_of_householderActiveBlockNorm2Sq_pos}
\lean{Algorithms/QR/HouseholderApply.lean}
\leanfit{coxHigham_exactSignedPivotPanel_sequence_active_block_bound_of_initial_block_bound_of_active_block_norm_pos}
\end{formalized}

\begin{corollary}[Active-block budgets after active max-column swapping]
Let \(q_t\) be the active column with maximal trailing column norm in a raw
stage \(A_t^{\rm raw}\), and let the displayed stage be obtained by swapping
columns \(k_t\) and \(q_t\):
\[
  A_t
    =
    \operatorname{swapcols}(A_t^{\rm raw},k_t,q_t).
\]
Then the displayed column \(k_t\) is pivot-maximal on the active suffix.  Hence,
under the same monotone active-window assumptions as above and the positive
active-block mass condition for the displayed stages, the exact signed-pivot
sequence satisfies
\[
  |A_t(i,\ell)|\le c_m^tB_0
  \qquad (t\le q,\ i\ge p_t,\ \ell\ge k_t).
\]
This closes the concrete column-swap policy field for the exact Cox--Higham
route.  It still leaves the positive active-block mass/nonbreakdown invariant
and the rounded QR/preconditioner solve connection open.
\end{corollary}


\begin{formalized}
\lean{Algorithms/QR/HouseholderSpec.lean}
\leanfit{householderSwapColumns}
\leanfit{householderTrailingColumnNorm2Sq_swapColumns_left}
\leanfit{householderTrailingColumnNorm2Sq_swapColumns_of_ne}
\leanfit{householderSwapColumns_activeMaxPivotColumn_pivot_max}
\lean{Algorithms/QR/HouseholderApply.lean}
\leanfit{coxHigham_exactSignedPivotPanel_sequence_active_block_bound_of_initial_block_bound_of_swapped_active_max_pivot}
\end{formalized}

\begin{corollary}[Raw active-block mass survives the active max-column swap]
In the previous corollary, the positive active-block mass condition may be
placed before the column swap.  If
\[
  \mathcal M_t^{\rm raw}
    =
    \sum_{i,\ell}
      \mathbf 1_{\{i\ge p_t,\ \ell\ge k_t\}}\,
      A_t^{\rm raw}(i,\ell)^2
    > 0
\]
and \(q_t\) lies in the active suffix, then
\[
  \mathcal M_t(A_t)>0
  \qquad
  \text{for }
  A_t=\operatorname{swapcols}(A_t^{\rm raw},k_t,q_t).
\]
Consequently the exact swapped-stage Cox--Higham sequence again satisfies
\[
  |A_t(i,\ell)|\le c_m^tB_0
  \qquad (t\le q,\ i\ge p_t,\ \ell\ge k_t).
\]
Thus the remaining nonbreakdown obligation is raw-stage active-block mass,
not an artifact introduced after sorting.
\end{corollary}


\begin{formalized}
\lean{Algorithms/QR/HouseholderSpec.lean}
\leanfit{householderActiveBlockNorm2Sq_pos_of_exists_active_entry_ne}
\leanfit{householderActiveBlockNorm2Sq_swapColumns_pos_of_pos}
\lean{Algorithms/QR/HouseholderApply.lean}
\leanfit{coxHigham_exactSignedPivotPanel_sequence_active_block_bound_of_initial_block_bound_of_swapped_active_max_pivot_of_raw_active_block_norm_pos}
\end{formalized}

\begin{corollary}[Raw active-block mass from QR nonbreakdown]
Fix a QR stage \(k\).  Suppose the current column is not in the span of the
previous columns, and suppose the previous columns span every vector supported
above the active row.  Then the raw active block has positive mass:
\[
  \mathcal M_k(A)>0 .
\]
Equivalently, the same conclusion follows from the concrete QR witnesses used
elsewhere in the library: a previous-block coefficient matrix plus a leading
column left inverse, local left inverses for the previous and current leading
blocks, or nonzero determinants of the previous and current leading blocks.

In stagewise determinant form, if for every stored QR stage \(k<n\) the
previous leading block and current leading block are nonsingular, and the
already completed columns have the stored lower-zero shape, then
\[
  \mathcal M_k(\widehat A_k)>0
  \qquad (k<n).
\]
Thus the raw-stage nonbreakdown field in the Cox--Higham active-block theorem
can be supplied by the existing QR rank/determinant infrastructure.
\end{corollary}


\begin{formalized}
\lean{Algorithms/QR/HouseholderQR.lean}
\leanfit{householderActiveBlockNorm2Sq_pos_of_column_notInPreviousSpan}
\leanfit{householderActiveBlockNorm2Sq_pos_of_leading_witnesses}
\leanfit{householderActiveBlockNorm2Sq_pos_of_leading_block_leftInverses}
\leanfit{householderActiveBlockNorm2Sq_pos_of_leading_block_det_ne_zero}
\leanfit{householderActiveBlockNorm2Sq_pos_sequence_of_leading_block_det_ne_zero}
\end{formalized}

\begin{corollary}[Rounded active-block budgets for signed-pivot stored panels]
Let
\[
  \widehat A_{t+1}
  =
  \operatorname{fl}\bigl((I-\beta_t v_t v_t^T)\widehat A_t\bigr)
\]
be the stored compact Householder panel update, where \(v_t\) and \(\beta_t\)
are the signed-pivot Householder vector and coefficient built from
\(\widehat A_t\), the active row \(p_t\), and the pivot column \(k_t\).  Assume
the active row and column windows are monotone, the pivot column has positive
active trailing norm, the pivot column is maximal among the remaining active
columns, completed columns are stored exactly, and the pivot-column storage
equation has the required zero entries.  Let \(B_t\ge 0\) be an active-block
budget for the stored matrix.  If
\[
  c_m B_t
  +
  \eta_t(i,\ell)
  \le B_{t+1}
\]
for every next active entry \((i,\ell)\), where
\[
  c_m=\max\{1+\sqrt2,\sqrt m\}
\]
and \(\eta_t(i,\ell)\) is the compact Householder componentwise FP budget, then
\[
  |\widehat A_t(i,\ell)|\le B_t
  \qquad
  (t\le q,\ i\ge p_t,\ \ell\ge k_t).
\]
\end{corollary}


\begin{formalized}
\lean{Algorithms/QR/HouseholderApply.lean}
\leanfit{signedPivotHouseholderVector}
\leanfit{signedPivotHouseholderBeta}
\leanfit{coxHigham_storedPanel_sequence_active_block_bound_of_signed_pivot_stage_bounds}
\end{formalized}

\begin{lemma}[Cox--Higham row-wise accumulated perturbation]
Let \(c=1+\sqrt2\).  Suppose a computed row-entry sequence
\(\widehat A_t(r,j)\) is compared against an exact sequence \(A_t(r,j)\), and
the entrywise discrepancy
\[
  E_t=|\widehat A_t(r,j)-A_t(r,j)|
\]
satisfies the affine recurrence
\[
  E_{t+1}\le c E_t+\eta_t,\qquad 0\le t<k .
\]
Define the accumulated additive budget recursively by
\[
  \mathcal E_0=0,\qquad
  \mathcal E_{t+1}=c\mathcal E_t+\eta_t .
\]
Then
\[
  E_k\le c^k E_0+\mathcal E_k .
\]
Consequently, if the exact sequence satisfies the sorted Cox--Higham row bound
\(|A_k(r,j)|\le c^k B_k^{(0)}\), then the computed sequence satisfies
\[
  |\widehat A_k(r,j)|
  \le c^k B_k^{(0)} + c^k E_0+\mathcal E_k .
\]
\end{lemma}


\begin{formalized}
\lean{Algorithms/QR/HouseholderSpec.lean}
\leanfit{scalarAffineGrowthBudget}
\leanfit{scalar_affine_growth_iterate_bound}
\leanfit{coxHigham_rowSorting_active_entry_bound_of_stage_growth_with_additive}
\leanfit{coxHigham_rowwise_error_accumulation_bound}
\leanfit{coxHigham_abs_entry_le_exact_bound_add_error}
\leanfit{coxHigham_rowSorting_active_entry_bound_with_accumulated_error}
\end{formalized}

\begin{lemma}[Concrete stored-panel FP budget for Cox--Higham row errors]
Consider one stored rounded compact Householder panel step
\[
  \widehat A_{t+1}
  =
  \operatorname{fl}\bigl((I-\beta_t v_t v_t^T)\widehat A_t\bigr)
\]
with the repository's stored-QR convention: completed columns are preserved and
the current pivot column is stored with zeros below the pivot.  If column \(j\)
is active at stage \(t\), then for each row \(r\)
\[
  \left|
    \widehat A_{t+1}(r,j)
    -
    (I-\beta_t v_t v_t^T)\widehat A_t(:,j)_r
  \right|
  \le
  B^{\mathrm{fp}}_t(r,j),
\]
where \(B^{\mathrm{fp}}_t(r,j)\) is exactly the compact Householder component
budget formalized in
\(\mathrm{householderCompactComponentBudget}\).  Consequently, if the exact
Cox--Higham part supplies the visible Lipschitz field
\[
  \left|
    (I-\beta_t v_t v_t^T)\widehat A_t(:,j)_r
      - A_{t+1}(r,j)
  \right|
  \le
  (1+\sqrt2)\,|\widehat A_t(r,j)-A_t(r,j)|,
\]
then
\[
  |\widehat A_{t+1}(r,j)-A_{t+1}(r,j)|
  \le
  (1+\sqrt2)|\widehat A_t(r,j)-A_t(r,j)|
  + B^{\mathrm{fp}}_t(r,j).
\]
Iterating this recurrence gives the same affine accumulated budget
\[
  |\widehat A_k(r,j)-A_k(r,j)|
  \le
  (1+\sqrt2)^k|\widehat A_0(r,j)-A_0(r,j)|
  +\mathcal E_k,
\]
with \(\mathcal E_{t+1}=(1+\sqrt2)\mathcal E_t+B^{\mathrm{fp}}_t(r,j)\).
\end{lemma}


\begin{formalized}
\lean{Algorithms/QR/HouseholderApply.lean}
\leanfit{fl_householderStoredPanelStep_active_entry_componentwise_error_bound}
\leanfit{coxHigham_storedPanelStep_row_error_recurrence_of_exact_lipschitz}
\leanfit{coxHigham_storedPanel_sequence_rowwise_error_accumulation_bound_of_exact_lipschitz}
\end{formalized}

\begin{lemma}[Direct rounded row-magnitude recurrence for stored panels]
Under the same stored-panel hypotheses, suppose the exact same-reflector update
of the current stored panel satisfies the Cox--Higham row-growth estimate
\[
  \left|
    (I-\beta_t v_t v_t^T)\widehat A_t(:,j)_r
  \right|
  \le
  (1+\sqrt2)|\widehat A_t(r,j)|.
\]
Then the rounded stored entry satisfies
\[
  |\widehat A_{t+1}(r,j)|
  \le
  (1+\sqrt2)|\widehat A_t(r,j)|
  + B^{\mathrm{fp}}_t(r,j).
\]
Consequently, after \(k\) stored panel steps,
\[
  |\widehat A_k(r,j)|
  \le
  (1+\sqrt2)^k|\widehat A_0(r,j)|+\mathcal E_k,
  \qquad
  \mathcal E_{t+1}=(1+\sqrt2)\mathcal E_t+B^{\mathrm{fp}}_t(r,j).
\]
This is the source-shaped floating-point dependency for the Cox--Higham route:
the remaining exact growth/sorting field is explicit, and the local FP
contribution is exactly the repository's compact Householder component budget.
\end{lemma}


\begin{formalized}
\lean{Algorithms/QR/HouseholderApply.lean}
\leanfit{coxHigham_storedPanelStep_active_entry_bound_of_exact_growth}
\leanfit{coxHigham_storedPanel_sequence_active_entry_bound_of_exact_growth}
\end{formalized}

\begin{corollary}[Stored-panel handoff for the active-row Cox--Higham factor]
Let
\[
  c_m=\max\{1+\sqrt2,\sqrt m\}.
\]
Suppose the stored rounded panel loop satisfies the same concrete compact
Householder update hypotheses as above, and suppose the exact same-reflector
part now supplies the active-row growth field
\[
  \left|
    (I-\beta_t v_t v_t^T)\widehat A_t(:,j)_r
  \right|
  \le
  c_m|\widehat A_t(r,j)|.
\]
Then, with the same compact floating-point component budget
\(B^{\mathrm{fp}}_t(r,j)\),
\[
  |\widehat A_k(r,j)|
  \le
  c_m^k|\widehat A_0(r,j)|
  +\mathcal E_k,
  \qquad
  \mathcal E_{t+1}=c_m\mathcal E_t+B^{\mathrm{fp}}_t(r,j).
\]
More generally, the Lean theorem is factor-parametric: any nonnegative exact
same-reflector growth factor \(c\) can be used in place of \(c_m\).
This is the floating-point handoff needed by the exact signed-pivot panel
sequence theorem; it still keeps the sorting-policy stage-budget field visible.
\end{corollary}


\begin{formalized}
\lean{Algorithms/QR/HouseholderApply.lean}
\leanfit{coxHigham_storedPanelStep_active_entry_bound_of_exact_growth_factor}
\leanfit{coxHigham_storedPanel_sequence_active_entry_bound_of_exact_growth_factor}
\leanfit{coxHigham_storedPanel_sequence_active_entry_bound_of_exact_active_growth}
\end{formalized}

\begin{corollary}[Stage-budget handoff for signed-pivot stored panels]
Let \(c_m=\max\{1+\sqrt2,\sqrt m\}\).  Suppose a stored rounded panel step uses
the signed-pivot Householder vector for the current pivot column, and suppose
the exact Cox--Higham stage fields hold at this step: positive active trailing
norm, column-pivot maximality, active-row budget
\[
  |A_t(r,\ell)|\le B_t,\qquad \ell\ge k,
\]
and active-column budget
\[
  |A_t(i,j)|\le B_t,\qquad i\ge p .
\]
Assume also the stored-panel convention for completed columns and the current
pivot column.  Then the rounded stored entry satisfies
\[
  |\widehat A_{t+1}(r,j)|
  \le
  c_m B_t + B^{\mathrm{fp}}_t(r,j),
\]
where \(B^{\mathrm{fp}}_t(r,j)\) is exactly the repository's compact
Householder component budget for the same vector, scalar, and column.
Consequently, if visible stage budgets satisfy
\[
  c_m B_t+B^{\mathrm{fp}}_t(r,j)\le B_{t+1},
  \qquad
  |\widehat A_0(r,j)|\le B_0,
\]
then the stored rounded sequence satisfies
\[
  |\widehat A_k(r,j)|\le B_k .
\]
\end{corollary}


\begin{formalized}
\lean{Algorithms/QR/HouseholderApply.lean}
\leanfit{coxHigham_storedPanelStep_active_entry_bound_of_exact_stage_budget_factor}
\leanfit{coxHigham_storedPanel_sequence_active_entry_bound_of_exact_stage_budgets_factor}
\leanfit{coxHigham_storedPanel_sequence_active_entry_bound_of_exact_active_stage_budgets}
\leanfit{coxHigham_storedPanelStep_active_entry_bound_of_signed_pivot_stage_bounds}
\end{formalized}

\begin{lemma}[Trailing-norm and pivot-budget scalar bridges]
Let \(x\) be the current pivot column and let \(x_{\mathrm{tail}}\) denote its
entries from row \(k\) onward.  If some active trailing entry is nonzero, then
\[
  \exists i\ge k,\ x_i\ne0
  \quad\Longrightarrow\quad
  \|x_{\mathrm{tail}}\|_2^2>0 .
\]
In particular, \(x_k\ne0\) is sufficient.  Moreover, if
\[
  \alpha_k^2=\|x_{\mathrm{tail}}\|_2^2,
\]
then
\[
  |\alpha_k|=\sqrt{\|x_{\mathrm{tail}}\|_2^2}.
\]
Consequently any bound
\[
  B_k<\sqrt{\|x_{\mathrm{tail}}\|_2^2}
\]
is exactly the stored-loop pivot-budget nonbreakdown condition
\[
  B_k<|\alpha_k|.
\]
Finally, for any active trailing index \(i\ge k\),
\[
  |x_i|\le \sqrt{\|x_{\mathrm{tail}}\|_2^2},
\]
so a concrete entrywise lower bound
\[
  B_k<|x_i|
\]
already implies the square-root trailing-norm budget.
\end{lemma}


\begin{formalized}
\lean{Algorithms/QR/HouseholderSpec.lean}
\leanfit{householderTrailingNorm2Sq_pos_of_exists_ne}
\leanfit{householderTrailingNorm2Sq_pos_of_pivot_ne_zero}
\leanfit{abs_alpha_eq_sqrt_trailingNorm2Sq}
\leanfit{budget_lt_abs_alpha_of_lt_sqrt_trailingNorm2Sq}
\leanfit{abs_entry_le_sqrt_householderTrailingNorm2Sq_of_pivot_le}
\leanfit{budget_lt_sqrt_householderTrailingNorm2Sq_of_lt_abs_active_entry}
\end{formalized}

\begin{lemma}[Positive trailing norm does not force a nonzero pivot]
The previous lemma has no converse in the unpivoted Householder setting.  There
is a two-entry active column \(x=(0,1)\) with pivot index \(0\) such that
\[
  \|x_{0:1}\|_2^2>0,
  \qquad x_0=0.
\]
Thus a proof route that tries to obtain the current leading pivot condition from
positive active trailing norm alone is invalid; it needs pivoting, a structured
leading-block invariant, or an explicit pivot/nonzero-domain assumption.
\end{lemma}


\begin{formalized}
\lean{Algorithms/QR/HouseholderSpec.lean}
\leanfit{householderTrailingPivotCounterexample2}
\leanfit{householderTrailingPivotCounterexample2_pivot_zero}
\leanfit{householderTrailingPivotCounterexample2_trailingNorm2Sq_pos}
\leanfit{not_forall_trailingNorm2Sq_pos_implies_pivot_ne_zero}
\end{formalized}

\begin{definition}[Prefix-span QR nonbreakdown invariant]
For the current stored QR panel \(A_k\), write \(a_j\) for its \(j\)-th
column.  The local prefix-span nonbreakdown invariant at pivot \(k\) has two
parts:
\[
  a_k\notin \operatorname{span}\{a_0,\ldots,a_{k-1}\},
\]
and every vector supported on the first \(k\) rows is in that same span:
\[
  \{y:\ y_i=0\ \text{for all } i\ge k\}
  \subseteq
  \operatorname{span}\{a_0,\ldots,a_{k-1}\}.
\]
\end{definition}

\begin{lemma}[Leading-block inverse gives prefix coefficients]
Let \(T_k\in\mathbb R^{k\times k}\) be the transpose of the leading block of
the first \(k\) rows and previous \(k\) columns,
\[
  T_k(j,s)=A_k(s,j),\qquad 0\le j,s<k.
\]
If \(C T_k=I\), then \(C\) is a valid prefix-basis coefficient matrix:
\[
  \sum_{j<k} C_{rj} A_k(s,j)=\delta_{sr},\qquad 0\le r,s<k.
\]
\end{lemma}


\begin{lemma}[Leading-block inverse gives ambient column duals]
Let \(S_k\in\mathbb R^{(k+1)\times(k+1)}\) be the leading block of the first
\(k+1\) rows and first \(k+1\) columns,
\[
  S_k(r,q)=A_k(r,q),\qquad 0\le r,q\le k.
\]
If \(D S_k=I\), define an ambient dual family by padding \(D\) with zeros:
\[
  L_{p i} =
  \begin{cases}
    D_{p i}, & i\le k,\\
    0, & i>k.
  \end{cases}
\]
Then \(L\) is a leading-column left inverse for the first \(k+1\) columns:
\[
  \sum_{i<m} L_{p i} A_k(i,q)=\delta_{q p},
  \qquad 0\le p,q\le k.
\]
\end{lemma}


\begin{lemma}[Leading-block coefficients imply prefix span]
Fix \(k\), and suppose the first \(k\) previous columns have already acquired
the QR lower-zero shape,
\[
  A_k(i,j)=0,\qquad i\ge k,\quad j<k.
\]
Assume there is a coefficient matrix \(C\in\mathbb R^{k\times k}\) such that
for every prefix coordinate basis vector \(e_r\),
\[
  \sum_{j<k} C_{rj}\, A_k(s,j)=\delta_{sr},
  \qquad 0\le r,s<k.
\]
Then the previous columns span every prefix-supported vector.  More precisely,
if \(y_i=0\) for all \(i\ge k\), then the coefficients
\[
  c_j=\sum_{r<k} y_r C_{rj}
\]
satisfy
\[
  y_i=\sum_{j<k} c_j A_k(i,j),\qquad 0\le i<m.
\]
\end{lemma}


\begin{lemma}[Leading-column left inverse gives column independence]
Suppose there is a dual coefficient family \(L\) for the first \(k+1\)
columns such that
\[
  \sum_i L_{p i} A_k(i,q)=\delta_{q p},
  \qquad 0\le p,q\le k.
\]
Then the current pivot column is not in the span of previous columns:
\[
  a_k\notin \operatorname{span}\{a_0,\ldots,a_{k-1}\}.
\]
\end{lemma}


\begin{lemma}[Leading-block witnesses give active trailing nonbreakdown]
Fix a pivot \(k\).  Suppose the previous columns have the QR lower-zero shape,
\[
  A_k(i,j)=0,\qquad i\ge k,\quad j<k,
\]
suppose a coefficient matrix \(C\) reproduces every prefix coordinate basis
vector using the previous columns,
\[
  \sum_{j<k} C_{rj}A_k(s,j)=\delta_{sr},\qquad 0\le r,s<k,
\]
and suppose a dual family \(L\) is a left inverse for the first \(k+1\)
columns,
\[
  \sum_i L_{p i}A_k(i,q)=\delta_{q p},\qquad 0\le p,q\le k.
\]
Then the active trailing part of the pivot column is nonzero:
\[
  \exists i\ge k,\quad A_k(i,k)\ne0.
\]
Consequently
\[
  \|A_k(k{:}m,k)\|_2^2>0.
\]
\end{lemma}


\begin{lemma}[Prefix-span nonbreakdown gives a nonzero active trailing entry]
If the prefix-span QR nonbreakdown invariant holds at pivot \(k\), then the
active trailing part of the pivot column is nonzero:
\[
  \exists i\ge k,\quad A_k(i,k)\ne0.
\]
Consequently
\[
  \|A_k(k{:}m,k)\|_2^2>0.
\]
\end{lemma}


\begin{corollary}[Stored QR diagonal nonbreakdown from prefix-span and square-root budget]
In the stored trailing Householder QR sequence, assume at each pivot \(k\)
the prefix-span QR nonbreakdown invariant, the sign-choice condition
\[
  \alpha_k A_k(k,k)\le0,
\]
the Householder normalization
\[
  \alpha_k^2=\|A_k(k{:}m,k)\|_2^2,
\]
and the explicit square-root pivot budget
\[
  B_k<\sqrt{\|A_k(k{:}m,k)\|_2^2}.
\]
Then the stored triangular diagonal is nonzero:
\[
  \widehat R_{ii}\ne0,\qquad 0\le i<n.
\]
\end{corollary}


\begin{corollary}[Stored QR diagonal nonbreakdown from prefix-span and active-entry budget]
In the stored trailing Householder QR sequence, assume at each pivot \(k\)
the prefix-span QR nonbreakdown invariant, the sign-choice condition, and the
Householder normalization
\[
  \alpha_k^2=\|A_k(k{:}m,k)\|_2^2 .
\]
Instead of assuming a square-root trailing-norm budget, suppose there is an
active trailing row \(i_k\ge k\) such that the compact diagonal update budget
is dominated by that entry,
\[
  B_k<|A_k(i_k,k)|.
\]
Then the stored triangular diagonal is nonzero:
\[
  \widehat R_{ii}\ne0,\qquad 0\le i<n.
\]
\end{corollary}


\begin{lemma}[Dimensioned norm margin gives an active entry]
Let \(x\in\mathbb R^m\), let \(p<m\), and let \(B\ge0\).  If
\[
  mB^2 < \|x(p{:}m)\|_2^2,
\]
then there exists an active index \(i\ge p\) such that
\[
  B<|x_i|.
\]
\end{lemma}


\begin{lemma}[Leading dual norm lower bound for the trailing pivot column]
Fix a pivot \(k\).  Let \(L_k\) be a leading-column dual for the first
\(k+1\) columns, and let \(L_{k,\ast}\) denote its row that selects the current
pivot column.  Assume:
\[
  L_{k,\ast}^T A_k(:,k)=1,\qquad
  L_{k,\ast}^T A_k(:,j)=0\quad(j<k),
\]
and assume the prefix-span invariant, so the prefix part of \(A_k(:,k)\) lies
in the span of the previous columns.  If
\[
  \|L_{k,\ast}\|_2^2\le K_k,\qquad K_k>0,
\]
then
\[
  \frac1{K_k}\le \|A_k(k{:}m,k)\|_2^2 .
\]
\end{lemma}


\begin{lemma}[Local inverse row norm lower bound]
Let \(C_k\) be a left inverse of the concrete leading block
\(S_k=A_k(0{:}k,0{:}k)\).  Pad the last row of \(C_k\) by zeros outside the
first \(k+1\) rows to obtain the ambient leading-column dual \(L_{k,\ast}\).
Then the padding preserves the squared norm:
\[
  \|L_{k,\ast}\|_2^2
    = \|C_k(k,:)\|_2^2 .
\]
Consequently, if
\[
  \|C_k(k,:)\|_2^2\le K_k,\qquad K_k>0,
\]
and the prefix-span invariant holds, then
\[
  \frac1{K_k}\le \|A_k(k{:}m,k)\|_2^2 .
\]
\end{lemma}


\begin{corollary}[Stored QR diagonal nonbreakdown from a leading-dual budget]
In the stored trailing Householder QR sequence, assume prefix-span, sign
choice, Householder normalization, and for every pivot \(k<n\) a leading
dual row satisfying
\[
  \|L_{k,\ast}\|_2^2\le K_k,\qquad K_k>0.
\]
If the compact diagonal update budget satisfies
\[
  mB_k^2 < \frac1{K_k},
\]
then the stored triangular diagonal is nonzero:
\[
  \widehat R_{ii}\ne0,\qquad 0\le i<n.
\]
\end{corollary}


\begin{corollary}[Stored QR diagonal nonbreakdown from a leading-block inverse row budget]
In the same stored trailing Householder QR sequence, replace the ambient dual
row by a local left inverse \(C_k\) of the leading block \(S_k\).  Suppose
\[
  \|C_k(k,:)\|_2^2\le K_k,\qquad K_k>0,
  \qquad
  mB_k^2 < \frac1{K_k}.
\]
Then the stored triangular diagonal is nonzero:
\[
  \widehat R_{ii}\ne0,\qquad 0\le i<n.
\]
\end{corollary}


\begin{corollary}[Stored QR diagonal nonbreakdown from a leading-block inverse Frobenius budget]
In the same stored trailing Householder QR sequence, assume a local left
inverse \(C_k\) of the leading block \(S_k\) at each pivot.  If
\[
  \|C_k\|_F^2\le K_k,\qquad K_k>0,
  \qquad
  mB_k^2 < \frac1{K_k},
\]
then the stored triangular diagonal is nonzero:
\[
  \widehat R_{ii}\ne0,\qquad 0\le i<n.
\]
\end{corollary}


\begin{corollary}[Stored QR diagonal nonbreakdown from a leading-block inverse infinity-norm budget]
In the same stored trailing Householder QR sequence, assume a local left
inverse \(C_k\) of the leading block \(S_k\) at each pivot.  If
\[
  (k+1)\|C_k\|_\infty^2\le K_k,\qquad K_k>0,
  \qquad
  mB_k^2 < \frac1{K_k},
\]
then the stored triangular diagonal is nonzero:
\[
  \widehat R_{ii}\ne0,\qquad 0\le i<n.
\]
\end{corollary}


\begin{corollary}[Stored QR diagonal nonbreakdown from prefix-span and norm-square budget]
In the stored trailing Householder QR sequence, assume the prefix-span QR
nonbreakdown invariant, sign choice, and Householder normalization at every
pivot \(k<n\).  If the compact diagonal update budget satisfies the dimensioned
margin
\[
  mB_k^2 < \|A_k(k{:}m,k)\|_2^2,
\]
then the stored triangular diagonal is nonzero:
\[
  \widehat R_{ii}\ne0,\qquad 0\le i<n.
\]
\end{corollary}


\begin{corollary}[Stored QR diagonal nonbreakdown from leading-block witnesses]
In the stored trailing Householder QR sequence, assume that at every pivot
\(k\) we have the leading-block coefficient witness \(C_k\), the leading-column
left-inverse witness \(L_k\), the QR lower-zero shape for the previous
columns, the sign-choice condition, the Householder normalization, and the
explicit square-root pivot budget
\[
  B_k<\sqrt{\|A_k(k{:}m,k)\|_2^2}.
\]
Then the final stored triangular diagonal is nonzero:
\[
  \widehat R_{ii}\ne0,\qquad 0\le i<n.
\]
\end{corollary}


\begin{corollary}[Stored QR diagonal nonbreakdown from local leading-block inverses]
In the stored trailing Householder QR sequence, assume that at every pivot
\(k\) the transposed previous leading block has a local left inverse,
\[
  C_k T_k=I,
\]
and the actual leading \((k+1)\times(k+1)\) block has a local left inverse,
\[
  D_k S_k=I.
\]
Assume also the QR lower-zero shape for the previous columns, the sign-choice
condition, the Householder normalization, and the explicit square-root pivot
budget
\[
  B_k<\sqrt{\|A_k(k{:}m,k)\|_2^2}.
\]
Then the final stored triangular diagonal is nonzero:
\[
  \widehat R_{ii}\ne0,\qquad 0\le i<n.
\]
\end{corollary}


\begin{lemma}[Nonzero determinant gives the local inverse predicate]
Let \(T\in\mathbb R^{d\times d}\).  If \(\det T\ne0\), then there exists a
matrix \(T^\dagger\) satisfying the repository predicate
\[
  T^\dagger T=I.
\]
The witness is Mathlib's nonsingular inverse, exported through the local
function-shaped matrix representation.
\end{lemma}

\begin{corollary}[Stored QR diagonal nonbreakdown from nonsingular leading blocks]
In the stored trailing Householder QR sequence, assume that for every pivot
\(k\) the transposed previous leading block \(T_k\) and the current leading
block \(S_k\in\mathbb R^{(k+1)\times(k+1)}\) are nonsingular:
\[
  \det T_k\ne0,\qquad \det S_k\ne0.
\]
Assume also the QR lower-zero shape for the previous columns, the sign-choice
condition, the Householder normalization, and the explicit square-root pivot
budget
\[
  B_k<\sqrt{\|A_k(k{:}m,k)\|_2^2}.
\]
Then the final stored triangular diagonal is nonzero:
\[
  \widehat R_{ii}\ne0,\qquad 0\le i<n.
\]
\end{corollary}


\begin{lemma}[Triangular local blocks have nonzero determinant]
Let \(U\in\mathbb R^{d\times d}\) be upper triangular in the convention
\[
  U_{ij}=0 \quad\text{whenever } j<i.
\]
If every diagonal entry is nonzero, then
\[
  \det U\ne0.
\]
The same statement holds for lower triangular matrices.  Consequently, in a
stored QR panel, the previous transposed leading block \(T_k\) and current
leading block \(S_k\) have nonzero determinant whenever their visible
triangular shape and local diagonal nonzero assumptions hold.
\end{lemma}


\begin{formalized}
\lean{Analysis/MatrixAlgebra.lean}
\leanfit{nonsingInv}
\leanfit{isLeftInverse_nonsingInv_of_det_isUnit}
\leanfit{exists_isLeftInverse_of_det_ne_zero}
\leanfit{det_ne_zero_of_upper_triangular_diag_ne_zero}
\leanfit{det_ne_zero_of_lower_triangular_diag_ne_zero}
\lean{Algorithms/QR/HouseholderQR.lean}
\leanfit{qrLeadingRow}
\leanfit{qrColumnNotInPreviousSpan}
\leanfit{qrPrefixSupportSpannedByPreviousColumns}
\leanfit{qrPrefixBasisCoefficientMatrix}
\leanfit{qrPreviousLeadingBlockTranspose}
\leanfit{qrLeadingBlock}
\leanfit{qrPrefixBasisCoefficientMatrix_of_leftInverse_previousLeadingBlockTranspose}
\leanfit{qrPrefixBasisCoefficientMatrix_of_det_ne_zero_previousLeadingBlockTranspose}
\leanfit{qrPrefixSupportSpannedByPreviousColumns_of_prefixBasisCoefficientMatrix}
\leanfit{qrPrefixSupportSpannedByPreviousColumns_of_leftInverse_previousLeadingBlockTranspose}
\leanfit{fl_householderStoredPanel_sequence_prefixSpan_of_leftInverse_previousLeadingBlockTranspose}
\leanfit{qrPrefixSupportSpannedByPreviousColumns_of_det_ne_zero_previousLeadingBlockTranspose}
\leanfit{qrLeadingColumnLeftInverse}
\leanfit{qrLeadingColumnLeftInverse_of_leftInverse_leadingBlock}
\leanfit{qrLeadingColumnLeftInverse_of_det_ne_zero_leadingBlock}
\leanfit{qrPreviousLeadingBlockTranspose_det_ne_zero_of_upper_triangular_diag_ne_zero}
\leanfit{qrLeadingBlock_det_ne_zero_of_upper_triangular_diag_ne_zero}
\leanfit{qrColumnNotInPreviousSpan_of_leadingColumnLeftInverse}
\leanfit{exists_active_trailing_entry_ne_of_column_notInPreviousSpan}
\leanfit{householderTrailingNorm2Sq_pos_of_column_notInPreviousSpan}
\leanfit{fl_householderStoredTrailingPanel_sequence_diag_nonzero_of_span_nonbreakdown_sqrt_budget}
\leanfit{fl_householderStoredTrailingPanel_sequence_diag_nonzero_of_span_nonbreakdown_active_entry_budget}
\leanfit{exists_active_entry_budget_lt_abs_of_dim_mul_budget_sq_lt_trailingNorm2Sq}
\leanfit{budget_lt_sqrt_householderTrailingNorm2Sq_of_dim_mul_budget_sq_lt_trailingNorm2Sq}
\leanfit{fl_householderStoredTrailingPanel_sequence_diag_nonzero_of_span_nonbreakdown_trailingNorm2Sq_budget}
\leanfit{householderTrailingNorm2Sq_ge_inv_leading_dual_norm_budget}
\leanfit{dim_mul_budget_sq_lt_trailingNorm2Sq_of_leading_dual_norm_budget}
\leanfit{fl_householderStoredTrailingPanel_sequence_diag_nonzero_of_span_nonbreakdown_leading_dual_norm_budget}
\leanfit{vecNorm2Sq_qrLeadingRow_padded_eq}
\leanfit{qrLeadingColumnLeftInverse_padded_row_norm_sq_eq}
\lean{Analysis/MatrixAlgebra.lean}
\leanfit{vecNorm2Sq_row_le_frobNormSq}
\leanfit{vecNorm2Sq_row_le_frobNorm_sq}
\leanfit{abs_coord_le_sum_abs}
\leanfit{vecNorm2Sq_le_sum_abs_sq}
\leanfit{frobNormSq_le_nat_mul_infNorm_sq}
\leanfit{frobNorm_sq_le_nat_mul_infNorm_sq}
\leanfit{infNorm_pos_of_det_ne_zero}
\leanfit{isRightInverse_nonsingInv_of_det_isUnit}
\leanfit{isInverse_nonsingInv_of_det_isUnit}
\leanfit{isInverse_nonsingInv_of_det_ne_zero}
\leanfit{exists_isInverse_of_det_ne_zero}
\lean{Algorithms/InverseBounds.lean}
\leanfit{triInv_infNorm_upperBound}
\leanfit{triInv_infNorm_sq_budget_of_diagDominantUpper}
\leanfit{triInv_infNorm_sq_budget_of_diagDominantUpper_det_ne_zero}
\lean{Algorithms/QR/HouseholderQR.lean}
\leanfit{householderTrailingNorm2Sq_ge_inv_leadingBlock_leftInverse_row_norm_budget}
\leanfit{dim_mul_budget_sq_lt_trailingNorm2Sq_of_leadingBlock_leftInverse_row_norm_budget}
\leanfit{householderTrailingNorm2Sq_ge_inv_leadingBlock_leftInverse_frobNorm_budget}
\leanfit{dim_mul_budget_sq_lt_trailingNorm2Sq_of_leadingBlock_leftInverse_frobNorm_budget}
\leanfit{householderTrailingNorm2Sq_ge_inv_leadingBlock_leftInverse_infNorm_budget}
\leanfit{dim_mul_budget_sq_lt_trailingNorm2Sq_of_leadingBlock_leftInverse_infNorm_budget}
\leanfit{fl_householderStoredTrailingPanel_sequence_diag_nonzero_of_span_nonbreakdown_leadingBlock_leftInverse_norm_budget}
\leanfit{fl_householderStoredTrailingPanel_sequence_diag_nonzero_of_span_nonbreakdown_leadingBlock_leftInverse_frobNorm_budget}
\leanfit{fl_householderStoredTrailingPanel_sequence_diag_nonzero_of_span_nonbreakdown_leadingBlock_leftInverse_infNorm_budget}
\leanfit{fl_householderStoredPanel_sequence_prefixSpan_of_leftInverse_previousLeadingBlockTranspose}
\leanfit{exists_active_trailing_entry_ne_of_leading_witnesses}
\leanfit{householderTrailingNorm2Sq_pos_of_leading_witnesses}
\leanfit{fl_householderStoredTrailingPanel_sequence_diag_nonzero_of_leading_witnesses_sqrt_budget}
\leanfit{exists_active_trailing_entry_ne_of_leading_block_leftInverses}
\leanfit{householderTrailingNorm2Sq_pos_of_leading_block_leftInverses}
\leanfit{fl_householderStoredTrailingPanel_sequence_diag_nonzero_of_leading_block_leftInverses_sqrt_budget}
\leanfit{exists_active_trailing_entry_ne_of_leading_block_det_ne_zero}
\leanfit{householderTrailingNorm2Sq_pos_of_leading_block_det_ne_zero}
\leanfit{fl_householderStoredTrailingPanel_sequence_diag_nonzero_of_leading_block_det_ne_zero_sqrt_budget}
\lean{Analysis/PerturbationTheory.lean}
\leanfit{infNorm_eq_sup_row_sum}
\leanfit{kappaInf_eq_infNorm_mul_infNorm}
\leanfit{infNorm_inv_le_of_kappaInf_le_and_norm_lower}
\leanfit{infNorm_sq_budget_of_kappaInf_le_and_norm_lower}
\leanfit{infNorm_inv_le_of_kappaInf_le_and_det_ne_zero}
\leanfit{infNorm_sq_budget_of_kappaInf_le_and_det_ne_zero}
\lean{Algorithms/LeastSquares/LSQRSolve.lean}
\leanfit{LSQRSolveBackwardError.of_stored_trailing_householder_sequence_topBlock_fl_backSub_gamma_bound_normBudget_of_span_nonbreakdown_active_entry_budget}
\leanfit{LSQRSolveBackwardError.of_stored_trailing_householder_sequence_topBlock_fl_backSub_gamma_bound_normBudget_of_span_nonbreakdown_trailingNorm2Sq_budget}
\leanfit{LSQRSolveBackwardError.of_stored_trailing_householder_sequence_topBlock_fl_backSub_gamma_bound_normBudget_of_span_nonbreakdown_leading_dual_norm_budget}
\leanfit{LSQRSolveBackwardError.of_stored_trailing_householder_sequence_topBlock_fl_backSub_gamma_bound_explicitNormBudget_of_span_nonbreakdown_leading_dual_norm_budget}
\leanfit{LSQRSolveBackwardError.of_stored_trailing_householder_sequence_topBlock_fl_backSub_gamma_bound_explicitCompactBudget_of_span_nonbreakdown_leading_dual_norm_budget}
\leanfit{LSQRSolveBackwardError.of_stored_trailing_householder_sequence_topBlock_fl_backSub_gamma_bound_normBudget_of_span_nonbreakdown_leadingBlock_leftInverse_norm_budget}
\leanfit{LSQRSolveBackwardError.of_stored_trailing_householder_sequence_topBlock_fl_backSub_gamma_bound_explicitNormBudget_of_span_nonbreakdown_leadingBlock_leftInverse_norm_budget}
\leanfit{LSQRSolveBackwardError.of_stored_trailing_householder_sequence_topBlock_fl_backSub_gamma_bound_explicitCompactBudget_of_span_nonbreakdown_leadingBlock_leftInverse_norm_budget}
\leanfit{LSQRSolveBackwardError.of_stored_trailing_householder_sequence_topBlock_fl_backSub_gamma_bound_explicitCompactBudget_of_previousLeadingBlock_leftInverse_leadingBlock_leftInverse_norm_budget}
\leanfit{LSQRSolveBackwardError.of_stored_trailing_householder_sequence_topBlock_fl_backSub_gamma_bound_normBudget_of_span_nonbreakdown_leadingBlock_leftInverse_frobNorm_budget}
\leanfit{LSQRSolveBackwardError.of_stored_trailing_householder_sequence_topBlock_fl_backSub_gamma_bound_explicitNormBudget_of_span_nonbreakdown_leadingBlock_leftInverse_frobNorm_budget}
\leanfit{LSQRSolveBackwardError.of_stored_trailing_householder_sequence_topBlock_fl_backSub_gamma_bound_explicitCompactBudget_of_span_nonbreakdown_leadingBlock_leftInverse_frobNorm_budget}
\leanfit{LSQRSolveBackwardError.of_stored_trailing_householder_sequence_topBlock_fl_backSub_gamma_bound_normBudget_of_span_nonbreakdown_leadingBlock_leftInverse_infNorm_budget}
\leanfit{LSQRSolveBackwardError.of_stored_trailing_householder_sequence_topBlock_fl_backSub_gamma_bound_explicitNormBudget_of_span_nonbreakdown_leadingBlock_leftInverse_infNorm_budget}
\leanfit{LSQRSolveBackwardError.of_stored_trailing_householder_sequence_topBlock_fl_backSub_gamma_bound_explicitCompactBudget_of_span_nonbreakdown_leadingBlock_leftInverse_infNorm_budget}
\leanfit{LSQRSolveBackwardError.of_stored_trailing_householder_sequence_topBlock_fl_backSub_gamma_bound_normBudget_of_span_nonbreakdown_leadingBlock_det_ne_zero_infNorm_budget}
\leanfit{LSQRSolveBackwardError.of_stored_trailing_householder_sequence_topBlock_fl_backSub_gamma_bound_explicitCompactBudget_of_previousLeadingBlock_leftInverse_leadingBlock_leftInverse_frobNorm_budget}
\leanfit{LSQRSolveBackwardError.of_stored_trailing_householder_sequence_topBlock_fl_backSub_gamma_bound_explicitCompactBudget_of_previousLeadingBlock_leftInverse_leadingBlock_leftInverse_infNorm_budget}
\leanfit{LSQRSolveBackwardError.of_stored_trailing_householder_sequence_topBlock_fl_backSub_gamma_bound_explicitCompactBudget_of_signed_alpha_previousLeadingBlock_leftInverse_leadingBlock_leftInverse_norm_budget}
\leanfit{LSQRSolveBackwardError.of_stored_trailing_householder_sequence_topBlock_fl_backSub_gamma_bound_explicitCompactBudget_of_signed_alpha_previousLeadingBlock_leftInverse_leadingBlock_leftInverse_frobNorm_budget}
\leanfit{LSQRSolveBackwardError.of_stored_trailing_householder_sequence_topBlock_fl_backSub_gamma_bound_explicitCompactBudget_of_signed_alpha_previousLeadingBlock_leftInverse_leadingBlock_leftInverse_infNorm_budget}
\leanfit{LSQRSolveBackwardError.of_stored_trailing_householder_sequence_topBlock_fl_backSub_gamma_bound_explicitCompactBudget_of_signed_alpha_previousLeadingBlock_det_ne_zero_leadingBlock_det_ne_zero_norm_budget}
\leanfit{LSQRSolveBackwardError.of_stored_trailing_householder_sequence_topBlock_fl_backSub_gamma_bound_explicitCompactBudget_of_signed_alpha_previousLeadingBlock_det_ne_zero_leadingBlock_det_ne_zero_frobNorm_budget}
\leanfit{LSQRSolveBackwardError.of_stored_trailing_householder_sequence_topBlock_fl_backSub_gamma_bound_explicitCompactBudget_of_signed_alpha_previousLeadingBlock_det_ne_zero_leadingBlock_det_ne_zero_infNorm_budget}
\leanfit{LSQRSolveBackwardError.of_stored_trailing_householder_sequence_topBlock_fl_backSub_gamma_bound_explicitCompactBudget_of_signed_alpha_previousLeadingBlock_det_ne_zero_leadingBlock_det_ne_zero_kappaInf_selfNorm_budget}
\leanfit{LSQRSolveBackwardError.of_stored_trailing_householder_sequence_topBlock_fl_backSub_gamma_bound_explicitCompactBudget_of_signed_alpha_upperTriangular_leadingDiag_ne_zero_kappaInf_selfNorm_budget}
\leanfit{LSQRSolveBackwardError.of_stored_trailing_householder_sequence_topBlock_fl_backSub_gamma_bound_normBudget_of_span_nonbreakdown_leadingBlock_det_ne_zero_kappaInf_budget}
\leanfit{LSQRSolveBackwardError.of_stored_trailing_householder_sequence_topBlock_fl_backSub_gamma_bound_normBudget_of_span_nonbreakdown_leadingBlock_det_ne_zero_kappaInf_selfNorm_budget}
\leanfit{LSQRSolveBackwardError.of_stored_trailing_householder_sequence_topBlock_fl_backSub_gamma_bound_normBudget_of_leading_blocks_det_ne_zero_kappaInf_selfNorm_budget}
\leanfit{LSQRSolveBackwardError.of_stored_trailing_householder_sequence_topBlock_fl_backSub_gamma_bound_normBudget_of_span_nonbreakdown_diagDominant_leadingBlock_invNorm_budget}
\leanfit{LSQRSolveBackwardError.of_stored_trailing_householder_sequence_topBlock_fl_backSub_gamma_bound_normBudget_of_span_nonbreakdown_diagDominant_leadingBlock_det_ne_zero_invNorm_budget}
\leanfit{LSQRSolveBackwardError.of_stored_trailing_householder_sequence_topBlock_fl_backSub_gamma_bound_normBudget_of_leading_block_det_ne_zero_sqrt_budget}
\leanfit{LSQRSolveBackwardError.of_stored_trailing_householder_sequence_topBlock_fl_backSub_gamma_bound_normBudget_of_triangular_leading_blocks_sqrt_budget}
\end{formalized}

\begin{lemma}[Full rank does not imply unpivoted leading-minor nonbreakdown]
There exists a real \(2\times2\) matrix \(A\) with \(\det A\ne0\) but
\(A_{00}=0\).  Therefore a no-pivot Householder QR route cannot replace the
current-pivot nonzero hypothesis by ordinary nonsingularity/full rank alone.
Nor can it replace the per-pivot leading-block determinant hypotheses by
whole-matrix nonsingularity: the first \(1\times1\) unpivoted leading block
of the same matrix has determinant zero.  A final no-pivot theorem must
therefore either assume no breakdown, use pivoting, or prove a stronger
structured invariant that implies the displayed pivot and leading-minor
conditions.
\end{lemma}


\begin{formalized}
\lean{Algorithms/QR/HouseholderQR.lean}
\leanfit{qrPivotCounterexample2_first_pivot_zero}
\leanfit{qrPivotCounterexample2_det_ne_zero}
\leanfit{not_forall_det_ne_zero_implies_first_pivot_ne_zero}
\leanfit{qrPivotCounterexample2_first_leadingBlock_det_zero}
\leanfit{not_forall_det_ne_zero_implies_all_leading_blocks_det_ne_zero}
\end{formalized}

\begin{corollary}[Structured current-pivot nonbreakdown from local leading determinants]
In the stored trailing Householder QR loop, let \(S_k\) denote the displayed
\((k+1)\times(k+1)\) leading block of the stored matrix before pivot \(k\).
The stored-loop lower-zero theorem makes \(S_k\) upper triangular.  If
\(\det S_k\ne 0\), then the current pivot \(A_k(k,k)\) is nonzero.
Consequently, the latest explicit compact QR least-squares certificate can
replace its bare current-pivot nonzero hypothesis by the structured local
condition
\[
  \det S_k\ne 0 \qquad (0\le k<n).
\]
Under the same signed-alpha rule, dimensioned norm-square compact pivot margin,
local \(\kappa_\infty\) hypotheses, dual/conditioning budget, and
\(\gamma\)-validity assumptions as before, the computed back-substitution
vector satisfies the same
\[
  \mathrm{LSQRSolveBackwardError}
    (A^TA,A^Tb,\widehat x;c_G^{\mathrm{QR}},c_g^{\mathrm{QR}}).
\]
\end{corollary}


\begin{formalized}
\lean{Analysis/MatrixAlgebra.lean}
\leanfit{diag_ne_zero_of_upper_triangular_det_ne_zero}
\lean{Algorithms/QR/HouseholderQR.lean}
\leanfit{qrLeadingBlock_current_pivot_ne_zero_of_local_upper_triangular_det_ne_zero}
\leanfit{fl_householderStoredTrailingPanel_sequence_current_pivot_ne_zero_of_leadingBlock_det_ne_zero}
\lean{Algorithms/LeastSquares/LSQRSolve.lean}
\leanfit{LSQRSolveBackwardError.of_stored_trailing_householder_sequence_topBlock_fl_backSub_gamma_bound_explicitCompactBudget_of_signed_alpha_prefix_lower_zero_leadingBlock_det_ne_zero_kappaInf_selfNorm_normSqBudget}
\end{formalized}

\begin{corollary}[Structured QR certificate from local \(\kappa_\infty\) and dual budgets]
In the setting of the preceding corollary, the separate dimensioned
norm-square compact pivot-margin hypothesis can be replaced by the following
local conditioning budget.  For every pivot \(0\le k<n\), let \(S_k\) be the
displayed local leading block.  Assume
\[
  \det S_k\ne 0,\qquad
  \kappa_\infty(S_k,\mathrm{nonsingInv}(S_k))\le \kappa_k,
\]
\[
  (k+1)\left(\frac{\kappa_k}{\|S_k\|_\infty}\right)^2\le K_k,
  \qquad
  m B_k^2 < \frac{1}{K_k},
\]
where \(B_k\) is the compact componentwise diagonal-update budget used at pivot
\(k\).  Under the same stored recurrences, signed-alpha rule, \(\gamma\)
validity, and final explicit compact/final-radius definitions as before, the
computed back-substitution vector satisfies
\[
  \mathrm{LSQRSolveBackwardError}
    (A^TA,A^Tb,\widehat x;c_G^{\mathrm{QR}},c_g^{\mathrm{QR}}).
\]
\end{corollary}


\begin{formalized}
\lean{Algorithms/QR/HouseholderQR.lean}
\leanfit{qrPrefixSupportSpannedByPreviousColumns_of_leadingBlock_upper_det_ne_zero}
\lean{Algorithms/LeastSquares/LSQRSolve.lean}
\leanfit{LSQRSolveBackwardError.of_stored_trailing_householder_sequence_topBlock_fl_backSub_gamma_bound_explicitCompactBudget_of_signed_alpha_prefix_lower_zero_leadingBlock_det_ne_zero_kappaInf_selfNorm_dualBudget}
\end{formalized}

\begin{corollary}[Structured QR certificate from a direct inverse-\(\infty\) budget]
In the same stored no-pivot QR setting, the local \(\kappa_\infty\) hypotheses
in the preceding corollary can be replaced by a direct inverse-budget
certificate.  For each pivot \(0\le k<n\), let \(S_k\) be the local leading
block and assume
\[
  \det S_k\ne 0,\qquad
  (k+1)\,\|\mathrm{nonsingInv}(S_k)\|_\infty^2\le K_k,
  \qquad
  mB_k^2<\frac{1}{K_k}.
\]
Under the same stored recurrences, signed-alpha rule, \(\gamma\)-validity
assumptions, and explicit compact/final-radius definitions, the computed
back-substitution vector satisfies
\[
  \mathrm{LSQRSolveBackwardError}
    (A^TA,A^Tb,\widehat x;c_G^{\mathrm{QR}},c_g^{\mathrm{QR}}).
\]
\end{corollary}


\begin{formalized}
\lean{Algorithms/LeastSquares/LSQRSolve.lean}
\leanfit{LSQRSolveBackwardError.of_stored_trailing_householder_sequence_topBlock_fl_backSub_gamma_bound_explicitCompactBudget_of_signed_alpha_prefix_lower_zero_leadingBlock_det_ne_zero_invNorm_dualBudget}
\end{formalized}

\begin{corollary}[Diagonal-dominant structured QR certificate]
In the setting of the preceding direct inverse-budget certificate, the displayed
inverse-\(\infty\) budget can itself be derived from Higham's triangular inverse
bound.  For each pivot \(0\le k<n\), let \(S_k\) be the local leading block.
Assume \(S_k\) is upper-triangular diagonally dominant, \(\det S_k\ne0\), and
\[
  (k+1)
  \left(
    2^k\,
    \frac{1}{\min_{0\le r\le k}|S_k(r,r)|}
  \right)^2
  \le K_k,
  \qquad
  mB_k^2<\frac{1}{K_k}.
\]
Then, under the same stored recurrences, signed-alpha rule,
\(\gamma\)-validity assumptions, and explicit compact/final-radius definitions,
the computed back-substitution vector satisfies
\[
  \mathrm{LSQRSolveBackwardError}
    (A^TA,A^Tb,\widehat x;c_G^{\mathrm{QR}},c_g^{\mathrm{QR}}).
\]
\end{corollary}


\begin{formalized}
\lean{Algorithms/LeastSquares/LSQRSolve.lean}
\leanfit{LSQRSolveBackwardError.of_stored_trailing_householder_sequence_topBlock_fl_backSub_gamma_bound_explicitCompactBudget_of_signed_alpha_prefix_lower_zero_diagDominant_leadingBlock_det_ne_zero_dualBudget}
\end{formalized}

\begin{corollary}[Concrete dual-budget diagonal-dominant QR certificate]
In the preceding diagonal-dominant certificate, define
\[
  D_k :=
  (k+1)
  \left(
    2^k\,
    \frac{1}{\min_{0\le r\le k}|S_k(r,r)|}
  \right)^2 .
\]
Assume the same stored recurrences, signed-alpha rule, \(\gamma\)-validity
conditions, local diagonal dominance, and \(\det S_k\ne0\).  If for every pivot
\[
  mB_k^2 < \frac{1}{2D_k},
\]
then the computed back-substitution vector satisfies the same
\[
  \mathrm{LSQRSolveBackwardError}
    (A^TA,A^Tb,\widehat x;c_G^{\mathrm{QR}},c_g^{\mathrm{QR}}).
\]
\end{corollary}


\begin{formalized}
\lean{Algorithms/InverseBounds.lean}
\leanfit{diagDominantUpperInvBudgetExpr}
\leanfit{diagDominantUpperInvBudgetExpr_pos}
\leanfit{triInv_infNorm_sq_budget_of_diagDominantUpper_det_ne_zero_twice_budget}
\lean{Algorithms/LeastSquares/LSQRSolve.lean}
\leanfit{LSQRSolveBackwardError.of_stored_trailing_householder_sequence_topBlock_fl_backSub_gamma_bound_explicitCompactBudget_of_signed_alpha_prefix_lower_zero_diagDominant_leadingBlock_det_ne_zero_concreteDualBudget}
\end{formalized}

\begin{corollary}[Product-form concrete dual-budget QR certificate]
Under the hypotheses of the preceding corollary, it is enough to assume the
dimensionless product inequality
\[
  2D_k\,mB_k^2 < 1
\]
for every pivot \(k\).  Then the same
\(\mathrm{LSQRSolveBackwardError}\) conclusion holds.
\end{corollary}


\begin{formalized}
\lean{Algorithms/InverseBounds.lean}
\leanfit{mul_sq_lt_inv_two_mul_of_two_mul_mul_sq_lt_one}
\lean{Algorithms/LeastSquares/LSQRSolve.lean}
\leanfit{LSQRSolveBackwardError.of_stored_trailing_householder_sequence_topBlock_fl_backSub_gamma_bound_explicitCompactBudget_of_signed_alpha_prefix_lower_zero_diagDominant_leadingBlock_det_ne_zero_concreteDualProductBudget}
\end{formalized}

\begin{corollary}[Stored-sequence product compact-smallness QR certificate]
Let \(c_{\mathrm{seq}}\) be the deterministic stored QR compact sequence
budget.  In the product-form concrete dual-budget certificate, the raw local
pivot compact budget \(B_k\) may be replaced by the computable upper bound
\[
  c_{\mathrm{seq}}\,
  \left\|(\widehat A_k)_{:,k}\right\|_2 .
\]
Thus it is enough to assume, for every pivot \(k\),
\[
  2D_k\,m
  \left(c_{\mathrm{seq}}
    \left\|(\widehat A_k)_{:,k}\right\|_2\right)^2 < 1.
\]
Then the same \(\mathrm{LSQRSolveBackwardError}\) conclusion holds.
\end{corollary}


\begin{formalized}
\lean{Algorithms/QR/HouseholderQR.lean}
\leanfit{storedQRCompactPivotBudget_le_sequence_column_norm}
\lean{Algorithms/InverseBounds.lean}
\leanfit{two_mul_mul_sq_lt_one_of_nonneg_le}
\lean{Algorithms/LeastSquares/LSQRSolve.lean}
\leanfit{LSQRSolveBackwardError.of_stored_trailing_householder_sequence_topBlock_fl_backSub_gamma_bound_explicitCompactBudget_of_signed_alpha_prefix_lower_zero_diagDominant_leadingBlock_det_ne_zero_concreteDualProductSequenceBudget}
\end{formalized}

\begin{lemma}[Budget positivity alone does not imply compact smallness]
The product-form certificate cannot be closed from \(D_k>0\) alone.  Even in
the scalar case \(D=1\), \(m=1\), and \(B=1\), the product condition would read
\[
  2\cdot1\cdot(1\cdot1^2)<1,
\]
which is false.  Thus a final QR theorem must either prove a real
compact-update budget bound from the computed loop or keep that product
smallness hypothesis visible.
\end{lemma}

\begin{formalized}
\lean{Algorithms/InverseBounds.lean}
\leanfit{not_forall_pos_implies_two_mul_mul_sq_lt_one}
\end{formalized}

\begin{lemma}[Diagonal dominance does not imply compact smallness]
Even local diagonal dominance and the displayed Higham inverse budget \(D_k\)
do not imply the product-form compact-smallness condition by themselves.  In
the scalar \(1\times1\) case, take the local block \(U=[1]\), so \(U\) is
diagonally dominant and \(D=1\).  With compact budget \(B=1\) and \(m=1\), the
same product condition again reads
\[
  2D\,mB^2=2<1,
\]
which is false.
\end{lemma}


\begin{formalized}
\lean{Algorithms/InverseBounds.lean}
\leanfit{not_forall_diagDominantUpper_implies_two_mul_budget_expr_mul_sq_lt_one}
\end{formalized}

\begin{lemma}[Product compact smallness does not imply diagonal dominance]
Even adding the product compact-smallness inequality to triangular
nonsingularity does not force diagonal dominance.  For the upper-triangular
matrix
\[
  U=\begin{pmatrix}1&2\\0&1\end{pmatrix},
\]
the displayed diagonal-dominant inverse budget \(D\) is finite.  Choosing, for
example, \(B=1/8\) and \(m=1\) makes
\[
  2D\,mB^2 < 1,
\]
but \(U\) is still not diagonally dominant, since \(2\nleq 1\).  Thus the
remaining QR assumptions cannot be collapsed into one another: diagonal
dominance and compact-product smallness each need a real computed-loop or
conditioning invariant, or else both must remain visible.
\end{lemma}

\begin{formalized}
\lean{Algorithms/InverseBounds.lean}
\leanfit{not_forall_upper_tri_det_ne_zero_product_budget_implies_diagDominant}
\end{formalized}

\begin{lemma}[Finite conditioning alone does not imply compact smallness]
Upper-triangular nonsingularity and a finite local
\(\kappa_\infty\) certificate also do not imply the product compact-smallness
condition.  For the same block
\[
  U=\begin{pmatrix}1&2\\0&1\end{pmatrix},
\]
choose the conditioning budget to be the exact formal value
\(\kappa_\infty(U)\).  The conditioning hypothesis is then satisfied.  But
with compact budget \(B=1\) and \(m=1\), the displayed product inequality
\[
  2D\,mB^2 < 1
\]
is false.  Thus local conditioning information does not control the
compact-update budget; that side condition still needs a computed-loop
compactness invariant or must remain visible as a domain hypothesis.
\end{lemma}

\begin{formalized}
\lean{Algorithms/InverseBounds.lean}
\leanfit{not_forall_upper_tri_det_ne_zero_kappaInf_bound_implies_two_mul_budget_expr_mul_sq_lt_one}
\end{formalized}

\begin{remark}[Current rectangular QR route choice]
After the preceding route eliminations, the remaining equation~(8)
QR/preconditioner gap is not another small adapter.  The formalized
counterexamples rule out deriving the visible assumptions from whole-matrix
full rank, positive trailing norm, diagonal dominance alone, product
compact-smallness alone, finite local \(\kappa_\infty\), or any implication
between diagonal dominance and product compact-smallness.  The next theorem
family must therefore be chosen explicitly: either prove a stronger
computed-loop/off-diagonal-control invariant, switch to a Cox--Higham style
pivoted or sorted row-wise weighted least-squares theorem, or keep the
remaining nonbreakdown, conditioning, and compact-product hypotheses visible
as domain assumptions.
\end{remark}

\begin{remark}[Cox--Higham is not a drop-in closure]
The Cox--Higham row-wise weighted least-squares route is deliberately not used
to close the current unpivoted stored-QR theorem.  Higham's QR notes state the
standard Householder result as a columnwise theorem, and then treat row-wise
weighted least-squares stability as a separate question: in general the
unpivoted answer is negative, while row-wise bounds require column pivoting
together with row pivoting or row sorting and the recommended sign convention.
Cox--Higham prove precisely this pivoted/sorted row-wise theorem family.
Consequently, switching to Cox--Higham would change the algorithmic hypotheses;
it is not a small adapter for the current equation~(8) implementation-backed
theorem.  For the present unpivoted route, the remaining honest alternatives
are to prove a stronger computed-loop/off-diagonal-control invariant or to
keep the local nonbreakdown, conditioning, diagonal-dominance, and
compact-product assumptions visible.
\end{remark}

\begin{lemma}[Triangular nonsingularity and finite conditioning alone do not imply diagonal dominance]
The diagonal-dominant branch above cannot be justified from upper-triangular
shape, nonzero diagonal entries, determinant nonzeroness, or the mere existence
of a finite \(\kappa_\infty\) certificate.  In Lean, the concrete matrix
\[
  \begin{pmatrix}1&2\\0&1\end{pmatrix}
\]
is upper triangular, has nonzero diagonal entries, and has determinant
nonzero, but it is not diagonally dominant in the repository's sense because
\(2\nleq 1\).  Taking the conditioning budget to be the matrix's own local
\(\kappa_\infty\) value gives a formal finite conditioning certificate for the
same counterexample.  Therefore, any final theorem using the diagonal-dominant
branch must either prove diagonal dominance from a stronger computed-loop or
conditioning invariant or keep it as a visible domain hypothesis.
\end{lemma}

\begin{formalized}
\lean{Algorithms/TriangularForwardBound.lean}
\leanfit{diagDominanceCounterexample2_det_ne_zero}
\leanfit{not_forall_upper_tri_diag_nonzero_implies_diagDominant}
\leanfit{not_forall_upper_tri_det_ne_zero_implies_diagDominant}
\leanfit{exists_upper_tri_det_ne_zero_kappaInf_bound_not_diagDominant}
\leanfit{not_forall_upper_tri_det_ne_zero_kappaInf_bound_implies_diagDominant}
\end{formalized}

\begin{lemma}[Exact QR shape alone does not imply diagonal dominance]
The obstruction above persists even if the matrix is already written in exact
QR form.  There exist matrices \(A,Q,R\in\mathbb R^{2\times2}\) such that
\[
  A = QR,\qquad Q \text{ is orthogonal},\qquad R
  \text{ is upper triangular with nonzero diagonal},
\]
but \(R\) is not diagonally dominant.  In Lean the witness is
\[
  Q=I,\qquad R=\begin{pmatrix}1&2\\0&1\end{pmatrix}.
\]
Thus a final theorem cannot obtain the diagonal-dominance hypothesis from the
final QR shape alone.  The positive unpivoted route must either prove a
stronger computed-loop/off-diagonal-control invariant or keep diagonal
dominance as a visible domain assumption.
\end{lemma}

\begin{formalized}
\lean{Algorithms/TriangularForwardBound.lean}
\leanfit{exists_orthogonal_upper_factorization_not_diagDominant}
\leanfit{not_forall_orthogonal_upper_factorization_implies_diagDominant}
\end{formalized}

\begin{lemma}[Exact no-pivot Householder QR need not give diagonal dominance]
The standard exact trailing Householder recurrence itself does not supply the
diagonal-dominance hypothesis.  There is a two-step unpivoted Householder QR
sequence \(A_0,A_1,A_2\in\mathbb R^{2\times2}\) with signed Householder
scalars \(\alpha_0,\alpha_1\) such that, for \(k=0,1\),
\[
  A_{k+1}
  =
  H_k A_k,\qquad
  \alpha_k^2=\|A_k(k{:}2,k)\|_2^2,\qquad
  v_k^Tv_k\ne0,
\]
where \(H_k\) is the exact trailing Householder reflector generated from the
\(k\)-th active column.  The concrete Lean witness starts from
\[
  A_0=\begin{pmatrix}1&2\\0&1\end{pmatrix}
\]
and reaches
\[
  A_2=\begin{pmatrix}-1&-2\\0&-1\end{pmatrix}.
\]
The final matrix is upper triangular and has nonzero diagonal, but it is not
diagonally dominant because \(|-2|>|-1|\).  Thus diagonal dominance is not a
generic consequence of the exact unpivoted Householder loop; a final theorem
must add a genuine off-diagonal-control/pivoting invariant or keep diagonal
dominance visible.  The strengthened universal theorem says the same
counterexample rules out deriving diagonal dominance together with any
additional final-block property, such as the scalar finite-max smallness side
condition, from the exact no-pivot recurrence alone.
\end{lemma}

\begin{formalized}
\lean{Algorithms/LeastSquares/LSQRSolve.lean}
\leanfit{exactHouseholderQRDiagDominanceCounterexample_step}
\leanfit{exactHouseholderQRDiagDominanceCounterexample_alpha_sq}
\leanfit{exactHouseholderQRDiagDominanceCounterexample_den_ne_zero}
\leanfit{exactHouseholderQRDiagDominanceCounterexampleFP}
\leanfit{exactHouseholderQRDiagDominanceCounterexample_stored_step}
\leanfit{exactHouseholderQRDiagDominanceCounterexample_signed_alpha_def}
\leanfit{exactHouseholderQRDiagDominanceCounterexample_not_diagDominant}
\leanfit{not_forall_exact_trailing_householder_sequence_implies_diagDominant}
\leanfit{not_forall_exact_trailing_householder_sequence_implies_diagDominant_and_property}
\leanfit{not_forall_signedAlphaDef_stored_trailing_sequence_actualU_sourceDen_implies_diagDominant}
\leanfit{exactHouseholderQRDiagDominanceCounterexample_not_activeMaxPivotChoice}
\leanfit{not_forall_signedAlphaDef_stored_trailing_sequence_actualU_sourceDen_implies_activeMaxPivotChoice}
\end{formalized}

\begin{lemma}[Active-block nonbreakdown does not imply off-diagonal control]
Even the stronger local data now available for the Cox--Higham raw-stage route
do not imply the source-shaped off-diagonal domination field.  There is a
two-by-two displayed leading-block sequence such that, at each stage, the
leading block is upper triangular, the leading-block determinant is nonzero,
and the active-block mass is positive:
\[
  \det S_k\ne0,
  \qquad
  \mathcal M_k(A_k)>0,
  \qquad k=0,1.
\]
Nevertheless the row-wise domination condition
\[
  |(S_k)_{ij}| \le |(S_k)_{ii}| \qquad (i<j)
\]
fails.  The Lean witness again uses
\[
  \begin{pmatrix}1&2\\0&1\end{pmatrix},
\]
whose first row has \(2>1\).  Therefore the rank/determinant
active-block-mass bridge above cannot be used as a hidden substitute for the
off-diagonal-control hypothesis required by
\(\lean{StoredQRSourceOffDiagonalControl}\).
\end{lemma}


\begin{formalized}
\lean{Algorithms/LeastSquares/LSQRSolve.lean}
\leanfit{not_forall_leadingBlock_upper_det_activeBlockPos_implies_offdiag_le_diag}
\end{formalized}

\begin{lemma}[Row-growth budgets do not imply diagonal lower bounds]
The diagonal lower-bound half of the Cox--Higham row-budget split is also
independent.  There is a two-by-two displayed leading-block sequence and a
row budget \(R_k=2\) such that, at each stage, the leading block is upper
triangular, the leading-block determinant is nonzero, the active-block mass is
positive, and every strict upper entry satisfies the row-growth budget
\[
  |(S_k)_{ij}| \le R_k \qquad (i<j).
\]
Nevertheless the matching diagonal lower-bound field
\[
  R_k \le |(S_k)_{ii}| \qquad (i<k)
\]
fails.  For the witness
\[
  S_1=\begin{pmatrix}1&2\\0&1\end{pmatrix},
  \qquad R_1=2,
\]
the strict upper entry is exactly bounded by \(R_1\), while the first displayed
diagonal has magnitude \(1\).  Hence \(2\le 1\) would be required.
The strengthened active-pivot variant shows that this obstruction persists
even if the sequence also satisfies the active max-pivot policy.  It uses
\[
  S_0=\begin{pmatrix}2&1\\0&1\end{pmatrix},
  \qquad
  S_1=\begin{pmatrix}1&2\\0&1\end{pmatrix},
\]
so the first stage has the displayed active column maximal, while the second
stage still violates the row-budget diagonal lower-bound field.
The active-block-budget strengthening adds that the same row budget bounds
every active trailing-block entry at each displayed stage:
\[
  |(A_k)_{r\ell}| \le R_k
  \qquad (k\le r,\; k\le \ell).
\]
The same two-stage witness satisfies this extra active-block magnitude
condition with \(R_k=2\), but still requires \(2\le 1\) in the displayed
diagonal lower-bound field.
Finally, the meta strengthening adds an arbitrary auxiliary side property
\(P(A_\bullet,R_\bullet)\).  Instantiating \(P\) with the always-true
property shows that no universal theorem of this shape can repair the
diagonal lower-bound failure by appending unrelated scalar or product
conditions.
\end{lemma}


\begin{formalized}
\lean{Algorithms/LeastSquares/LSQRSolve.lean}
\leanfit{not_forall_leadingBlock_upper_det_activeBlockPos_offdiagBudget_implies_rowBudget_diag}
\leanfit{activeMaxPivotRowBudgetDiagCounterexampleA0}
\leanfit{activeMaxPivotRowBudgetDiagCounterexampleSeq}
\leanfit{not_forall_leadingBlock_upper_det_activeBlockPos_activeMaxPivot_offdiagBudget_implies_rowBudget_diag}
\leanfit{not_forall_leadingBlock_upper_det_activeBlockPos_activeMaxPivot_activeBlockBudget_offdiagBudget_implies_rowBudget_diag}
\leanfit{not_forall_leadingBlock_upper_det_activeBlockPos_activeMaxPivot_activeBlockBudget_offdiagBudget_property_implies_rowBudget_diag}
\end{formalized}

\begin{remark}[Route choice after the exact-recurrence obstruction]
The exact-recurrence counterexample turns the remaining equation~(8)
QR/preconditioner gap into a theorem-scope choice rather than another local
adapter.  For the current unpivoted stored-QR theorem, the false generic routes
have been ruled out from full rank, positive trailing norm, finite local
conditioning, diagonal dominance/product smallness alone, final exact QR
shape, the standard exact no-pivot Householder recurrence, and positive
active-block mass plus local leading-block nonsingularity.  The row-budget
split obstruction above also rules out deriving the required diagonal
lower-bound field from a row-growth upper budget alone.  A final positive
theorem must therefore choose one of the following mathematical statements:
add and prove a stronger computed-loop/off-diagonal-control and diagonal
nonbreakdown invariant, change the algorithmic theorem family to a
pivoted/sorted row-wise route, or keep the remaining nonbreakdown,
conditioning, diagonal-dominance, off-diagonal-control, diagonal lower-bound,
and compact-product assumptions visible as domain hypotheses.
\end{remark}

\begin{corollary}[QR solve certificate from an off-diagonal-control invariant]
Fix \(n\le m\), a floating-point model, and the stored trailing Householder
recurrences for \(\widehat A_k,\widehat b_k\) with the signed Householder
choice \(\alpha_k\).  Suppose that the loop satisfies the invariant
\(\mathcal I_{\mathrm{off}}\) consisting, for every \(k<n\), of:
\[
  \det S_k\ne0,\qquad
  S_k\ \hbox{is upper-triangular diagonally dominant},
\]
where \(S_k\) is the \((k+1)\times(k+1)\) displayed leading block of
\(\widehat A_k\), and
\[
  2D_k\, m\,
  \bigl(B\,\|\widehat A_k(:,k)\|_2\bigr)^2 < 1 .
\]
Here \(D_k\) is the formal Higham diagonal-dominant triangular inverse budget
for \(S_k\), and \(B\) is the repository's deterministic stored QR compact
sequence budget for the whole loop.  Then the computed triangular solve
\[
  \widehat x
  =
  \mathrm{fl\_backSub}\bigl(\widehat R,\widehat c\bigr)
\]
satisfies
\[
  \mathrm{LSQRSolveBackwardError}
  \bigl(n,A^TA,A^Tb,\widehat x,
        c_G(A,\widehat R,B),c_g(A,b,\widehat R,B)\bigr),
\]
with \(c_G=\lean{qrSolveFinalGramBudget}\) and
\(c_g=\lean{qrSolveFinalRhsBudget}\).
\end{corollary}


\begin{formalized}
\lean{Algorithms/LeastSquares/LSQRSolve.lean}
\leanfit{StoredQROffDiagonalControlInvariant}
\leanfit{LSQRSolveBackwardError.of_stored_trailing_householder_sequence_topBlock_fl_backSub_gamma_bound_explicitCompactBudget_of_signed_alpha_offDiagonalControl}
\end{formalized}

\begin{corollary}[Finite-max diagonal-dominant route to the packaged invariant]
In the setting above, suppose every displayed leading block \(S_k\) is
upper-triangular diagonally dominant.  Let \(D_*\) be the canonical maximum of
the diagonal-dominant inverse factors, let \(N_*\) be the canonical maximum of
the pivot-column norms, and let \(B\) be the stored QR compact sequence
relative budget.  These three quantities are the repository definitions listed
in the formalization block below.
If the scalar finite-max smallness condition
\[
  2D_*\,m\,(B N_*)^2 < 1
\]
holds, then the whole stored QR loop satisfies
\(\lean{StoredQROffDiagonalControlInvariant}\).  Consequently, the preceding
QR solve certificate applies with the repository perturbation radii
\(\lean{qrSolveFinalGramBudget}\) and \(\lean{qrSolveFinalRhsBudget}\).
\end{corollary}


\begin{formalized}
\lean{Algorithms/LeastSquares/LSQRSolve.lean}
\leanfit{StoredQROffDiagonalControlInvariant.of_diagDominant_finiteMaxSmallness}
\leanfit{storedQRCompactSequenceProductBudget_lt_one_of_diagDominant_finite_max_smallness}
\leanfit{storedQRDiagDominantInvFactorBudget}
\leanfit{storedQRPivotColumnNormBudget}
\leanfit{storedQRCompactSequenceRelativeBudget}
\leanfit{det_ne_zero_of_diagDominantUpper}
\end{formalized}

\begin{corollary}[Source-shaped off-diagonal control implies the QR certificate]
In the setting of the preceding corollary, it is enough to verify the following
more primitive local control data.  For every \(k<n\), the displayed leading
block \(S_k\) is upper triangular, its displayed diagonal is nonzero, and its
off-diagonal entries in each row are dominated by the row diagonal:
\[
  (S_k)_{ij}=0\quad(j<i),\qquad
  (S_k)_{ii}\ne0,\qquad
  |(S_k)_{ij}|\le |(S_k)_{ii}|\quad(i<j).
\]
Assume also the same stored-sequence compact-product condition
\[
  2D_k\,m\,
  \bigl(B\,\|\widehat A_k(:,k)\|_2\bigr)^2 < 1 .
\]
Then the stored QR loop satisfies \(\mathcal I_{\mathrm{off}}\), and therefore
the same computed triangular solve satisfies the
\(\mathrm{LSQRSolveBackwardError}\) certificate with
\(\lean{qrSolveFinalGramBudget}\) and \(\lean{qrSolveFinalRhsBudget}\).
\end{corollary}


\begin{formalized}
\lean{Algorithms/LeastSquares/LSQRSolve.lean}
\leanfit{StoredQRSourceOffDiagonalControl}
\leanfit{StoredQROffDiagonalControlInvariant.of_sourceOffDiagonalControl}
\leanfit{LSQRSolveBackwardError.of_stored_trailing_householder_sequence_topBlock_fl_backSub_gamma_bound_explicitCompactBudget_of_signed_alpha_sourceOffDiagonalControl}
\end{formalized}

\begin{corollary}[Stored recurrence supplies the triangular source field]
In the same stored trailing Householder recurrence, the lower-zero shape theorem
for completed columns implies the upper-triangular leading-block equalities
\[
  (S_k)_{ij}=0\qquad (j<i)
\]
at every displayed leading block \(S_k\).  Therefore, to instantiate the
source-shaped off-diagonal-control certificate above, it is enough to supply
only the remaining three source-specific obligations:
\[
  (S_k)_{ii}\ne0,\qquad
  |(S_k)_{ij}|\le |(S_k)_{ii}|\quad(i<j),\qquad
  2D_k\,m\,
  \bigl(B\,\|\widehat A_k(:,k)\|_2\bigr)^2 < 1 .
\]
Under these three obligations, the same computed triangular solve satisfies
the \(\mathrm{LSQRSolveBackwardError}\) certificate with
\(\lean{qrSolveFinalGramBudget}\) and
\(\lean{qrSolveFinalRhsBudget}\).
\end{corollary}


\begin{formalized}
\lean{Algorithms/LeastSquares/LSQRSolve.lean}
\leanfit{StoredQRSourceOffDiagonalControl.of_stored_trailing_sequence_diag_offdiag_product}
\leanfit{LSQRSolveBackwardError.of_stored_trailing_householder_sequence_topBlock_fl_backSub_gamma_bound_explicitCompactBudget_of_signed_alpha_diag_offdiag_product}
\end{formalized}

\begin{corollary}[Stored recurrence supplies previous diagonal nonzeros]
Assume the same stored trailing Householder recurrence and signed-alpha rule.
Assume also the square-root nonbreakdown budget
\[
  B^{\rm diag}_k
  < \sqrt{\|A_k(k:m,k)\|_2^2},
  \qquad k<n,
\]
where \(B^{\rm diag}_k\) is the local compact component budget for the
displayed pivot entry.  If the current pivot entries \(A_k(k,k)\) are nonzero,
and if the row-wise off-diagonal domination and compact-product obligations of
the source-shaped route hold, then
\[
  \lean{StoredQRSourceOffDiagonalControl}
\]
holds.  Consequently the corresponding computed triangular solve satisfies the
same \(\lean{LSQRSolveBackwardError}\) certificate with
\(\lean{qrSolveFinalGramBudget}\) and
\(\lean{qrSolveFinalRhsBudget}\).
\end{corollary}


\begin{formalized}
\lean{Algorithms/LeastSquares/LSQRSolve.lean}
\leanfit{StoredQRSourceOffDiagonalControl.of_stored_trailing_sequence_pivot_sqrtBudget_offdiag_product}
\leanfit{LSQRSolveBackwardError.of_stored_trailing_householder_sequence_topBlock_fl_backSub_gamma_bound_explicitCompactBudget_of_signed_alpha_pivot_sqrtBudget_offdiag_product}
\end{formalized}

\begin{corollary}[Stored recurrence supplies current pivots from leading determinants]
Assume the same stored trailing Householder recurrence, signed-alpha rule, and
square-root nonbreakdown budget as in the preceding corollary.  Replace the
raw current-pivot condition \(A_k(k,k)\ne0\) by the structured local
leading-block condition
\[
  \det S_k \ne 0,\qquad k<n,
\]
where \(S_k\) is the displayed \((k+1)\times(k+1)\) leading block of
\(A_k\).  If the row-wise off-diagonal domination and compact-product
obligations of the source-shaped route also hold, then
\[
  \lean{StoredQRSourceOffDiagonalControl}
\]
holds, and the corresponding computed triangular solve satisfies the same
\(\lean{LSQRSolveBackwardError}\) certificate with
\(\lean{qrSolveFinalGramBudget}\) and
\(\lean{qrSolveFinalRhsBudget}\).
\end{corollary}


\begin{formalized}
\lean{Algorithms/QR/HouseholderQR.lean}
\leanfit{fl_householderStoredTrailingPanel_sequence_current_pivot_ne_zero_of_leadingBlock_det_ne_zero}
\lean{Algorithms/LeastSquares/LSQRSolve.lean}
\leanfit{StoredQRSourceOffDiagonalControl.of_stored_trailing_sequence_leadingBlock_det_ne_zero_sqrtBudget_offdiag_product}
\leanfit{LSQRSolveBackwardError.of_stored_trailing_householder_sequence_topBlock_fl_backSub_gamma_bound_explicitCompactBudget_of_signed_alpha_leadingBlock_det_ne_zero_sqrtBudget_offdiag_product}
\end{formalized}

\begin{corollary}[Norm-square nonbreakdown form of the source-shaped QR route]
In the preceding corollary, the square-root nonbreakdown budget may be replaced
by the dimensioned trailing-norm-square margin
\[
  m\,(B^{\rm diag}_k)^2
  < \|A_k(k:m,k)\|_2^2,\qquad k<n,
\]
where \(B^{\rm diag}_k\) is the same local compact component budget for the
displayed pivot entry.  Under the same local leading-block nonsingularity,
row-wise off-diagonal domination, and compact-product hypotheses, the stored
QR loop again satisfies
\[
  \lean{StoredQRSourceOffDiagonalControl},
\]
and the corresponding computed triangular solve satisfies the same
\(\lean{LSQRSolveBackwardError}\) certificate with
\(\lean{qrSolveFinalGramBudget}\) and
\(\lean{qrSolveFinalRhsBudget}\).
\end{corollary}


\begin{formalized}
\lean{Algorithms/QR/HouseholderSpec.lean}
\leanfit{budget_lt_sqrt_householderTrailingNorm2Sq_of_dim_mul_budget_sq_lt_trailingNorm2Sq}
\lean{Algorithms/LeastSquares/LSQRSolve.lean}
\leanfit{StoredQRSourceOffDiagonalControl.of_stored_trailing_sequence_leadingBlock_det_ne_zero_normSqBudget_offdiag_product}
\leanfit{LSQRSolveBackwardError.of_stored_trailing_householder_sequence_topBlock_fl_backSub_gamma_bound_explicitCompactBudget_of_signed_alpha_leadingBlock_det_ne_zero_normSqBudget_offdiag_product}
\end{formalized}

\begin{corollary}[Row-budget form of the source-shaped QR route]
Keep the determinant, norm-square nonbreakdown, and compact-product hypotheses
of the preceding corollary.  Instead of assuming row-wise off-diagonal
domination directly, suppose that for every displayed leading block \(S_k\)
there are row budgets \(R_{k,i}\) such that
\[
  |(S_k)_{ij}| \le R_{k,i} \le |(S_k)_{ii}|,
  \qquad 0\le i<j\le k<n .
\]
Then \(S_k\) satisfies the row-wise off-diagonal domination field required by
\(\lean{StoredQRSourceOffDiagonalControl}\), and the same computed triangular
solve satisfies the local \(\lean{LSQRSolveBackwardError}\) certificate.
\end{corollary}


\begin{formalized}
\lean{Algorithms/LeastSquares/LSQRSolve.lean}
\leanfit{StoredQRSourceOffDiagonalControl.of_stored_trailing_sequence_leadingBlock_det_ne_zero_normSqBudget_rowBudget_product}
\leanfit{LSQRSolveBackwardError.of_stored_trailing_householder_sequence_topBlock_fl_backSub_gamma_bound_explicitCompactBudget_of_signed_alpha_leadingBlock_det_ne_zero_normSqBudget_rowBudget_product}
\leanfit{StoredQRSourceOffDiagonalControl.of_stored_trailing_sequence_leadingBlock_det_ne_zero_normSqBudget_rowBudget_product_of_offdiag_rows}
\leanfit{LSQRSolveBackwardError.of_stored_trailing_householder_sequence_topBlock_fl_backSub_gamma_bound_explicitCompactBudget_of_signed_alpha_leadingBlock_det_ne_zero_normSqBudget_rowBudget_product_of_offdiag_rows}
\end{formalized}

\begin{corollary}[Packaged displayed row-budget control]
In the preceding row-budget setting, suppose the two row-budget obligations are
provided as a single certificate:
\[
  \mathcal R(A,R) =
  \left[
  |(S_k)_{ij}| \le R_{k,i}\ \text{for all } i<j,\qquad
  R_{k,i}\le |(S_k)_{ii}|\ \text{for all } i<k
  \right].
\]
Then the stored QR source-control theorem and the solver-level backward-error
certificate follow from this packaged assumption, together with the same
leading-block nonsingularity, norm-square nonbreakdown, and compact-product
smallness hypotheses as above.
\end{corollary}


\begin{formalized}
\lean{Algorithms/LeastSquares/LSQRSolve.lean}
\lean{StoredQRDisplayedRowBudgetControl}
\leanfit{StoredQRDisplayedRowBudgetControl.of_signed_stage_budgets_factor_of_normSqBudget}
\leanfit{StoredQRDisplayedRowBudgetControl.of_signed_stage_uniformBudget_factor_of_normSqBudget}
\leanfit{StoredQRDisplayedRowBudgetControl.of_signed_stage_uniformBudget_globalCompactBudget_of_normSqBudget}
\leanfit{StoredQRDisplayedRowBudgetControl.of_signed_stage_uniformBudget_globalCompactBudget_kappaInf_dualBudget}
\leanfit{StoredQRDisplayedRowBudgetControl.of_signed_stage_uniformBudget_globalCompactBudget_activeMaxPivot_of_normSqBudget}
\leanfit{StoredQRDisplayedRowBudgetControl.of_signed_stage_uniformBudget_globalCompactBudget_activeMaxPivot_kappaInf_dualBudget}
\leanfit{StoredQRSourceOffDiagonalControl.of_stored_trailing_sequence_leadingBlock_det_ne_zero_normSqBudget_rowBudgetControl_product}
\leanfit{LSQRSolveBackwardError.of_stored_trailing_householder_sequence_topBlock_fl_backSub_gamma_bound_explicitCompactBudget_of_signed_alpha_leadingBlock_det_ne_zero_normSqBudget_rowBudgetControl_product}
\end{formalized}

\begin{corollary}[Canonical row-max row-budget certificate]
Assume the stronger source-shaped off-diagonal-control field is available:
for every displayed leading block \(S_k\),
\[
  |(S_k)_{ij}| \le |(S_k)_{ii}|,\qquad i<j\le k .
\]
For each row \(i<k\), define the canonical row budget
\[
  R_{k,i}^{\max}=\max_{i<j\le k}|(S_k)_{ij}|.
\]
Then the packaged displayed row-budget certificate holds with
\(R_{k,i}=R_{k,i}^{\max}\).
\end{corollary}


\begin{formalized}
\lean{Algorithms/LeastSquares/LSQRSolve.lean}
\lean{qrLeadingStrictUpperRowMaxBudget}
\lean{qrLeadingStrictUpperRowMaxBudget_entry_le}
\lean{qrLeadingStrictUpperRowMaxBudget_le_diag_of_offdiag}
\leanfit{StoredQRDisplayedRowBudgetControl.of_sourceOffDiagonalControl_rowMaxBudget}
\end{formalized}

\begin{corollary}[Route-1 row-max contraction handoff]
Let \(S_k\) be the displayed \((k+1)\times(k+1)\) leading block of the
stored QR data at stage \(k\), and let
\[
  R^{\max}_{k,i}=\max_{i<j\le k}|(S_k)_{ij}|
  \qquad (i<k)
\]
be the canonical strict-upper row maximum.  Suppose a computed-loop invariant
supplies one scalar \(\rho\le 1\) such that
\[
  R^{\max}_{k,i}\le \rho\, |(S_k)_{ii}|
  \qquad\text{for all } i<k<n .
\]
Then the packaged displayed row-budget certificate holds with
\(R_{k,i}=R^{\max}_{k,i}\).  Consequently, under the same stored recurrence,
leading-block nonsingularity, norm-square nonbreakdown, and finite global
compact-product hypotheses used in the row-budget route, the
\(\lean{StoredQRSourceOffDiagonalControl}\) certificate and the local
\(\lean{LSQRSolveBackwardError}\) theorem follow.
\end{corollary}


\begin{formalized}
\lean{Algorithms/LeastSquares/LSQRSolve.lean}
\leanfit{StoredQRDisplayedRowBudgetControl.of_rowMaxBudget_le_diag_factor}
\leanfit{StoredQRSourceOffDiagonalControl.of_stored_trailing_sequence_leadingBlock_det_ne_zero_normSqBudget_rowMaxBudgetFactor_globalProduct}
\leanfit{LSQRSolveBackwardError.of_stored_trailing_householder_sequence_topBlock_fl_backSub_gamma_bound_explicitCompactBudget_of_signed_alpha_leadingBlock_det_ne_zero_normSqBudget_rowMaxBudgetFactor_globalProduct}
\end{formalized}

\begin{corollary}[Scalar row-max defect handoff]
With the notation of the previous corollary, define the finite scalar defect
\[
  \Delta_{\max}
    =\max_{0\le k<n}\max_{0\le i<k}
      \bigl(R^{\max}_{k,i}- |(S_k)_{ii}|\bigr),
\]
with value \(0\) if the index set is empty.  If
\[
  \Delta_{\max}\le 0,
\]
then the packaged displayed row-budget certificate holds with
\(R_{k,i}=R^{\max}_{k,i}\).  Under the same stored recurrence,
leading-block nonsingularity, norm-square nonbreakdown, and finite global
compact-product hypotheses, the source-control certificate and the local
QR solve backward-error theorem follow.
\end{corollary}


\begin{formalized}
\lean{Algorithms/LeastSquares/LSQRSolve.lean}
\lean{storedQRRowMaxDiagDefectBudget}
\lean{storedQRRowMaxDiagDefect_le_budget}
\leanfit{StoredQRDisplayedRowBudgetControl.of_rowMaxDiagDefectBudget_nonpos}
\leanfit{StoredQRSourceOffDiagonalControl.of_stored_trailing_sequence_leadingBlock_det_ne_zero_normSqBudget_rowMaxDiagDefect_globalProduct}
\leanfit{LSQRSolveBackwardError.of_stored_trailing_householder_sequence_topBlock_fl_backSub_gamma_bound_explicitCompactBudget_of_signed_alpha_leadingBlock_det_ne_zero_normSqBudget_rowMaxDiagDefect_globalProduct}
\end{formalized}

\begin{corollary}[Actual-unit-roundoff validity for scalar row-max-defect route]
The scalar row-max-defect/global-product QR certificate can be stated without
an abstract \(\lean{gammaValid}\) hypothesis.  The local solver theorem derives
both \(\lean{gammaValid fp m}\) and \(\lean{gammaValid fp n}\) from the
displayed actual-unit-roundoff guard
\[
  m\,\lean{fp.u} < 1,
\]
and the probability-level equation (8) theorem derives
\(\lean{gammaValid fp s}\) and \(\lean{gammaValid fp n}\) from
\[
  s\,\lean{fp.u} < 1.
\]
The determinant, norm-square nonbreakdown, scalar row-defect, and global
compact-product hypotheses remain visible.
\end{corollary}

\begin{remark}
This is a validity-surface reduction only.  It does not prove the scalar
row-defect condition, local determinant or nonbreakdown facts, the global
compact-product smallness condition, or the generic implementation-backed
rectangular QR/preconditioner theorem.
\end{remark}

\begin{formalized}
\lean{Algorithms/LeastSquares/LSQRSolve.lean}
\leanfit{LSQRSolveBackwardError.of_stored_trailing_householder_sequence_topBlock_fl_backSub_gamma_bound_explicitCompactBudget_of_signed_alpha_leadingBlock_det_ne_zero_normSqBudget_rowMaxDiagDefect_globalProduct_of_actualUnitRoundoff_no_gammaValid}
\lean{Algorithms/RandNLA/LeastSquaresSketch.lean}
\leanfit{leverageTraceProbability_eventProb_fl_rowSampleLSObjective_le_one_add_eta_of_augmentedSpan_sample_budget_of_objective_error_and_stored_qr_rowMaxDiagDefect_globalProduct_solver_of_actualUnitRoundoff_no_gammaValid}
\end{formalized}

\begin{corollary}[Finite stage-diagonal lower-bound defect]
For the active/prefix Cox--Higham QR route, let \(B_k\) be the monotone
stage budget and let \(S_k\) be the displayed \((k+1)\times(k+1)\) leading
block.  Define the finite scalar defect
\[
  \Delta_{\mathrm{stage}}
    =\max_{0\le k<n}\max_{0\le i<k}
      \bigl(B_k-| (S_k)_{ii} |\bigr),
\]
with value \(0\) if the index set is empty.  If
\[
  \Delta_{\mathrm{stage}}\le 0,
\]
then every displayed off-diagonal row satisfies the diagonal lower-bound
condition
\[
  B_k\le |(S_k)_{ii}| \qquad (i<k).
\]
Conversely, if these displayed pointwise lower bounds hold for every
off-diagonal row \(i<k\), then the same finite maximum satisfies
\(\Delta_{\mathrm{stage}}\le 0\).  Thus the scalar defect is now formalized as
an exact finite package for the displayed stage-diagonal lower-bound family.
Consequently the active/prefix global compact-step and finite global-product
QR source-control theorem, and its local QR solve wrapper, can consume the
single scalar condition \(\Delta_{\mathrm{stage}}\le 0\) instead of the full
pointwise diagonal-lower-bound family.
\end{corollary}

\begin{remark}
This is a bottleneck dependency reduction, not a proof of diagonal growth.
The scalar condition \(\Delta_{\mathrm{stage}}\le 0\) still has to be proved
from a concrete pivoted/sorted/off-diagonal-controlled loop, or retained as a
visible domain assumption.
\end{remark}

\begin{formalized}
\lean{Algorithms/LeastSquares/LSQRSolve.lean}
\lean{storedQRStageDiagLowerDefectBudget}
\lean{storedQRStageDiagLowerDefect_le_budget}
\leanfit{storedQRStageBudget_le_diag_of_stageDiagLowerDefectBudget_nonpos}
\leanfit{storedQRStageDiagLowerDefectBudget_nonpos_of_stageBudget_le_diag}
\leanfit{StoredQRSourceOffDiagonalControl.of_stored_trailing_sequence_leadingBlock_det_ne_zero_normSqBudget_signedStageUniformBudget_product_of_normSqBudget_activePrefix_activeBlockRecurrence_globalCompactBudget_completedColumns_globalProduct_stageDiagDefect_offdiag_rows}
\leanfit{LSQRSolveBackwardError.of_stored_trailing_householder_sequence_topBlock_fl_backSub_gamma_bound_explicitCompactBudget_of_signed_alpha_leadingBlock_det_ne_zero_normSqBudget_signedStageUniformBudget_product_of_normSqBudget_activePrefix_activeBlockRecurrence_globalCompactBudget_completedColumns_globalProduct_stageDiagDefect_offdiag_rows}
\leanfit{StoredQRSourceOffDiagonalControl.of_stored_trailing_sequence_leadingBlock_det_ne_zero_kappaInf_selfNorm_dualBudget_signedStageUniformBudget_globalCompactBudget_completedColumns_globalProduct_stageDiagDefect_offdiag_rows}
\leanfit{LSQRSolveBackwardError.of_stored_trailing_householder_sequence_topBlock_fl_backSub_gamma_bound_explicitCompactBudget_of_signed_alpha_leadingBlock_det_ne_zero_kappaInf_selfNorm_dualBudget_signedStageUniformBudget_globalCompactBudget_completedColumns_globalProduct_stageDiagDefect_offdiag_rows}
\end{formalized}

\begin{corollary}[Row-max bridge to stage-diagonal scalar nonpositivity]
Suppose the row-max scalar defect satisfies \(\Delta_{\max}\le 0\), and suppose
the uniform stage budget is bounded by the displayed strict-upper row maximum:
\[
  B_k \le R^{\max}_{k,i}\qquad (i<k).
\]
Then the stage-diagonal scalar defect also satisfies
\[
  \Delta_{\mathrm{stage}}\le 0.
\]
In particular, local displayed diagonal dominance implies the same conclusion
under the same \(B_k\le R^{\max}_{k,i}\) comparison.
\end{corollary}


\begin{formalized}
\lean{Algorithms/LeastSquares/LSQRSolve.lean}
\leanfit{storedQRStageDiagLowerDefectBudget_nonpos_of_rowMaxDiagDefectBudget_nonpos_stageBudget_le_rowMax}
\leanfit{storedQRStageDiagLowerDefectBudget_nonpos_of_diagDominant_stageBudget_le_rowMax}
\end{formalized}

\begin{corollary}[Finite stage-budget/row-max comparison defect]
For the same displayed off-diagonal rows \(i<k\), define
\[
  \Delta_{\mathrm{cmp}}
    =\max_{0\le k<n}\max_{0\le i<k}
      \bigl(B_k-R^{\max}_{k,i}\bigr),
\]
again with value \(0\) if the index set is empty.  Nonpositivity of this scalar
defect is equivalent, as a finite package, to the displayed comparison
\[
  B_k\le R^{\max}_{k,i}\qquad (i<k).
\]
Consequently the row-max bridge can be stated using two scalar conditions:
\[
  \Delta_{\max}\le 0,\qquad \Delta_{\mathrm{cmp}}\le 0.
\]
These imply the scalar stage-diagonal condition
\[
  \Delta_{\mathrm{stage}}\le 0.
\]
\end{corollary}

\begin{remark}
This packages a listed dependency; it does not prove the comparison defect
from the concrete pivoted/sorted QR loop.  That scalar condition remains a
visible invariant unless supplied by a stronger row-growth or ordering proof.
\end{remark}

\begin{formalized}
\lean{Algorithms/LeastSquares/LSQRSolve.lean}
\lean{storedQRStageRowMaxComparisonDefectBudget}
\lean{storedQRStageRowMaxComparisonDefect_le_budget}
\leanfit{storedQRStageBudget_le_rowMax_of_stageRowMaxComparisonDefectBudget_nonpos}
\leanfit{storedQRStageRowMaxComparisonDefectBudget_nonpos_of_stageBudget_le_rowMax}
\leanfit{storedQRStageDiagLowerDefectBudget_nonpos_of_rowMaxDiagDefectBudget_nonpos_stageRowMaxComparisonDefectBudget_nonpos}
\leanfit{storedQRStageDiagLowerDefectBudget_nonpos_of_diagDominant_stageRowMaxComparisonDefectBudget_nonpos}
\end{formalized}

\begin{corollary}[Scalar-comparison active-pivot row-max QR surface]
The active-pivot row-max source-control, solver, and sampled equation~(8)
wrappers can now be stated using the two scalar row-max conditions
\[
  \Delta_{\max}\le 0,\qquad \Delta_{\mathrm{cmp}}\le 0,
\]
instead of the displayed comparison family
\(B_k\le R^{\max}_{k,i}\).  Internally, the proofs extract that family from
\(\Delta_{\mathrm{cmp}}\le 0\) and then reuse the already-formalized visible
row-max active-pivot theorems.  The explicit-validity horizon sampled sibling
keeps the sampled \(\gamma\)-validity fields visible, but it calls the horizon
row-max probability wrapper after the same finite comparison extraction, so
the two-scalar surface no longer exposes samplewise global budget monotonicity.
The actual-unit horizon sampled sibling derives those validity fields from
\((s : \mathbb{R})u < 1\) and then calls the explicit-validity horizon
wrapper, so it also removes sampled validity fields.
\end{corollary}

\begin{remark}
This is a theorem-interface reduction, not a proof of either scalar defect
from a concrete pivoted or sorted loop.  The remaining route still has to
establish \(\Delta_{\max}\le 0\), \(\Delta_{\mathrm{cmp}}\le 0\),
determinant/conditioning fields, and compact-product smallness, or retain
them as explicit assumptions.
\end{remark}

\begin{formalized}
\lean{Algorithms/LeastSquares/LSQRSolve.lean}
\lean{Algorithms/RandNLA/LeastSquaresSketch.lean}
\leanfit{StoredQRSourceOffDiagonalControl.of_stored_trailing_sequence_leadingBlock_det_ne_zero_kappaInf_selfNorm_dualBudget_signedStageUniformBudget_globalCompactBudget_activeMaxPivot_completedColumns_globalProduct_rowMaxDiagDefect_stageRowMaxComparisonDefect_offdiag_rows}
\leanfit{StoredQRSourceOffDiagonalControl.of_stored_trailing_sequence_leadingBlock_det_ne_zero_kappaInf_selfNorm_dualBudget_signedStageUniformBudget_globalCompactBudget_activeMaxPivot_completedColumns_globalProduct_rowMaxDiagDefect_stageRowMaxComparisonDefect_offdiag_rows_of_actualUnitRoundoff_no_gammaValid}
\leanfit{LSQRSolveBackwardError.of_stored_trailing_householder_sequence_topBlock_fl_backSub_gamma_bound_explicitCompactBudget_of_signed_alpha_leadingBlock_det_ne_zero_kappaInf_selfNorm_dualBudget_signedStageUniformBudget_globalCompactBudget_activeMaxPivot_completedColumns_globalProduct_rowMaxDiagDefect_stageRowMaxComparisonDefect_offdiag_rows}
\leanfit{LSQRSolveBackwardError.of_stored_trailing_householder_sequence_topBlock_fl_backSub_gamma_bound_explicitCompactBudget_of_signed_alpha_leadingBlock_det_ne_zero_kappaInf_selfNorm_dualBudget_signedStageUniformBudget_globalCompactBudget_activeMaxPivot_completedColumns_globalProduct_rowMaxDiagDefect_stageRowMaxComparisonDefect_offdiag_rows_of_actualUnitRoundoff_no_gammaValid}
\leanfit{leverageTraceProbability_eventProb_fl_rowSampleLSObjective_le_one_add_eta_of_augmentedSpan_sample_budget_of_objective_error_and_stored_qr_activePrefix_globalProduct_activeMaxPivot_rowMaxDiagDefect_stageRowMaxComparisonDefect_kappaInf_dualBudget_solver}
\leanfit{leverageTraceProbability_eventProb_fl_rowSampleLSObjective_le_one_add_eta_of_augmentedSpan_sample_budget_of_objective_error_and_stored_qr_activePrefix_globalProduct_activeMaxPivot_rowMaxDiagDefect_stageRowMaxComparisonDefect_kappaInf_dualBudget_solver_of_horizonBudget}
\leanfit{leverageTraceProbability_eventProb_fl_rowSampleLSObjective_le_one_add_eta_of_augmentedSpan_sample_budget_of_objective_error_and_stored_qr_activePrefix_globalProduct_activeMaxPivot_rowMaxDiagDefect_stageRowMaxComparisonDefect_kappaInf_dualBudget_solver_of_actualUnitRoundoff_no_gammaValid}
\leanfit{leverageTraceProbability_eventProb_fl_rowSampleLSObjective_le_one_add_eta_of_augmentedSpan_sample_budget_of_objective_error_and_stored_qr_activePrefix_globalProduct_activeMaxPivot_rowMaxDiagDefect_stageRowMaxComparisonDefect_kappaInf_dualBudget_solver_of_actualUnitRoundoff_no_gammaValid_of_horizonBudget}
\end{formalized}

\begin{corollary}[Diagonal-dominant scalar-comparison active-pivot surface]
For the local active-pivot source-control and solver routes, local displayed
diagonal dominance can replace the separate leading-block determinant and
row-max-defect fields.  The proof first obtains leading-block nonsingularity
from diagonal dominance, and then combines diagonal dominance with
\(\Delta_{\mathrm{cmp}}\le0\) to obtain the scalar stage-diagonal condition.
\end{corollary}

\begin{remark}
This is a local theorem-surface reduction.  It does not prove local diagonal
dominance, the scalar comparison defect, conditioning or dual compact-budget
fields, active-pivot policy, compact-product smallness, or the final generic
QR/preconditioner theorem from a concrete computed loop.
\end{remark}

\begin{formalized}
\lean{Algorithms/LeastSquares/LSQRSolve.lean}
\leanfit{StoredQRSourceOffDiagonalControl.of_stored_trailing_sequence_diagDominant_kappaInf_selfNorm_dualBudget_signedStageUniformBudget_globalCompactBudget_activeMaxPivot_completedColumns_globalProduct_stageRowMaxComparisonDefect_offdiag_rows}
\leanfit{LSQRSolveBackwardError.of_stored_trailing_householder_sequence_topBlock_fl_backSub_gamma_bound_explicitCompactBudget_of_signed_alpha_diagDominant_kappaInf_selfNorm_dualBudget_signedStageUniformBudget_globalCompactBudget_activeMaxPivot_completedColumns_globalProduct_stageRowMaxComparisonDefect_offdiag_rows}
\end{formalized}

\begin{corollary}[Probability-level diagonal-dominant scalar-comparison surface]
The sampled equation~(8) active-pivot theorem has the same diagonal-dominant
scalar-comparison interface.  For each sampled trace, local displayed diagonal
dominance supplies the leading-block determinant and row-max scalar-defect
fields, while the scalar comparison condition
\(\Delta_{\mathrm{cmp}}\le0\) supplies the displayed
stage-budget/row-max comparison through the finite-max extractor.  The theorem
then reuses the scalar-comparison active-pivot probability wrapper.
\end{corollary}

\begin{remark}
This lifts the local dependency reduction to the high-probability rounded
sampled-row objective theorem.  The remaining visible obligations are unchanged:
local diagonal dominance, the scalar comparison defect, conditioning and dual
compact-budget data, active-pivot policy, compact-product smallness, and the
final concrete QR/preconditioner theorem.
\end{remark}

\begin{formalized}
\lean{Algorithms/RandNLA/LeastSquaresSketch.lean}
\leanfit{leverageTraceProbability_eventProb_fl_rowSampleLSObjective_le_one_add_eta_of_augmentedSpan_sample_budget_of_objective_error_and_stored_qr_activePrefix_globalProduct_activeMaxPivot_diagDominant_stageRowMaxComparisonDefect_kappaInf_dualBudget_solver}
\end{formalized}

\begin{lemma}[Diagonal dominance does not force the scalar comparison defect]
Local displayed diagonal dominance alone does not imply
\[
  \Delta_{\mathrm{cmp}}\le0.
\]
There is a checked two-by-two witness whose displayed leading blocks are
diagonally dominant, but whose uniform stage budget still exceeds the
displayed strict-upper row maximum.
\end{lemma}


\begin{formalized}
\lean{Algorithms/LeastSquares/LSQRSolve.lean}
\leanfit{activeMaxPivotRowMaxComparisonCounterexample_diagDominant}
\leanfit{not_forall_diagDominant_implies_stageRowMaxComparisonDefectBudget_nonpos}
\end{formalized}

\begin{lemma}[Product smallness does not rescue diagonal dominance]
Adding the compact-product smallness inequality to local displayed diagonal
dominance still does not imply
\[
  \Delta_{\mathrm{cmp}}\le0.
\]
\end{lemma}


\begin{formalized}
\lean{Algorithms/LeastSquares/LSQRSolve.lean}
\leanfit{activeMaxPivotRowMaxComparisonCounterexample_productBudget}
\leanfit{not_forall_diagDominant_product_budget_implies_stageRowMaxComparisonDefectBudget_nonpos}
\end{formalized}

\begin{lemma}[Active-pivot budget data does not rescue the comparison scalar]
Even adding the active-pivot budget surface to local displayed diagonal
dominance and compact-product smallness still does not imply
\[
  \Delta_{\mathrm{cmp}}\le0.
\]
\end{lemma}


\begin{formalized}
\lean{Algorithms/LeastSquares/LSQRSolve.lean}
\leanfit{not_forall_diagDominant_product_activeBlockPos_activeMaxPivot_activeBlockBudget_offdiagBudget_implies_stageRowMaxComparisonDefectBudget_nonpos}
\end{formalized}

\begin{corollary}[Finite-max product handoff for the active diagonal-dominant branch]
On the active-pivot diagonal-dominant scalar-comparison route, the raw product
hypothesis
\[
  \Pi_{\max}<1
\]
may be replaced by the canonical finite-max smallness scalar involving the
diagonal-dominant inverse-factor budget, the compact-sequence relative budget,
and the pivot-column norm budget.
\end{corollary}


\begin{formalized}
\lean{Algorithms/LeastSquares/LSQRSolve.lean}
\leanfit{StoredQRSourceOffDiagonalControl.of_stored_trailing_sequence_diagDominant_kappaInf_selfNorm_dualBudget_signedStageUniformBudget_globalCompactBudget_activeMaxPivot_completedColumns_finiteMaxSmallness_stageRowMaxComparisonDefect_offdiag_rows}
\leanfit{LSQRSolveBackwardError.of_stored_trailing_householder_sequence_topBlock_fl_backSub_gamma_bound_explicitCompactBudget_of_signed_alpha_diagDominant_kappaInf_selfNorm_dualBudget_signedStageUniformBudget_globalCompactBudget_activeMaxPivot_completedColumns_finiteMaxSmallness_stageRowMaxComparisonDefect_offdiag_rows}
\lean{Algorithms/RandNLA/LeastSquaresSketch.lean}
\leanfit{leverageTraceProbability_eventProb_fl_rowSampleLSObjective_le_one_add_eta_of_augmentedSpan_sample_budget_of_objective_error_and_stored_qr_activePrefix_finiteMaxSmallness_activeMaxPivot_diagDominant_stageRowMaxComparisonDefect_kappaInf_dualBudget_solver}
\end{formalized}

\begin{corollary}[Concrete-dual finite-max handoff for the active diagonal-dominant branch]
On the active-pivot diagonal-dominant scalar-comparison route, the auxiliary
\(\kappa_k\), \(K_k\), and dual compact-budget hypotheses can be removed from
the finite-max theorem surface.  It is enough to supply local diagonal
dominance, the signed-stage recurrence budget, active-pivot policy, the scalar
comparison defect, and the canonical finite-max product-smallness scalar.
\end{corollary}


\begin{formalized}
\lean{Algorithms/LeastSquares/LSQRSolve.lean}
\leanfit{storedQRSignedStage_normSqBudget_of_leadingBlock_det_ne_zero_diagDominant_concreteDualProductSequenceBudget}
\leanfit{StoredQRSourceOffDiagonalControl.of_stored_trailing_sequence_diagDominant_signedStageUniformBudget_globalCompactBudget_activeMaxPivot_completedColumns_finiteMaxSmallness_stageRowMaxComparisonDefect_concreteDual_offdiag_rows}
\leanfit{LSQRSolveBackwardError.of_stored_trailing_householder_sequence_topBlock_fl_backSub_gamma_bound_explicitCompactBudget_of_signed_alpha_diagDominant_signedStageUniformBudget_globalCompactBudget_activeMaxPivot_completedColumns_finiteMaxSmallness_stageRowMaxComparisonDefect_concreteDual_offdiag_rows}
\lean{Algorithms/RandNLA/LeastSquaresSketch.lean}
\leanfit{leverageTraceProbability_eventProb_fl_rowSampleLSObjective_le_one_add_eta_of_augmentedSpan_sample_budget_of_objective_error_and_stored_qr_activePrefix_finiteMaxSmallness_activeMaxPivot_diagDominant_stageRowMaxComparisonDefect_concreteDual_solver}
\end{formalized}

\begin{corollary}[Actual-unit validity for the active concrete-dual finite-max handoff]
Let \(n\le m\).  On the local active concrete-dual finite-max QR surface,
assume the computed stored-QR data
\[
  \widehat A_{k+1}=\fl_{\rm panel}(\widehat A_k,\alpha_k),
  \qquad
  \alpha_k=\alpha_{\rm sign}
    \Bigl(\sqrt{\|(\widehat A_k)_{k:m,k}\|_2^2},
           (\widehat A_k)_{kk}\Bigr)
\]
for \(k<n\), local diagonal dominance of every leading block
\[
  {\rm IsDiagDominantUpper}\bigl((\widehat A_k)_{0:k,0:k}\bigr),
\]
the initial displayed-entry and full-entry bounds by \(B_0\), the signed
global recurrence
\[
  \rho_m B_t+
  {\rm storedQRSignedStageGlobalCompactBudget}
      (fp,\widehat A,\alpha,t)
  \le B_{t+1},
\]
nonnegativity and monotonicity of \(B_t\), the active-pivot policy
\[
  k={\rm householderActiveMaxPivotColumn}
       (\widehat A_k;k),
\]
the scalar comparison defect
\[
  {\rm storedQRStageRowMaxComparisonDefectBudget}
       (\widehat A,B)\le 0,
\]
and the finite-max concrete-dual smallness condition
\[
  2\,D_{\rm dd}(\widehat A)\,
  \Bigl(m\bigl(C_{\rm seq}(fp,\widehat A,\widehat b,\alpha)
              N_{\rm piv}(\widehat A)\bigr)^2\Bigr)<1.
\]
If the actual unit roundoff satisfies
\[
  m\,fp.u<1,
\]
then the source-control conclusion
\[
  {\rm StoredQRSourceOffDiagonalControl}
       (n\le m,fp,\widehat A,\widehat b,\alpha)
\]
holds without an explicit \({\rm gammaValid}(fp,m)\) hypothesis.

If, in addition, the right-hand side recurrence
\[
  \widehat b_{k+1}=\fl_{\rm rhs}(\widehat b_k,\widehat A_k,\alpha_k),
  \qquad
  \widehat A_0=A,\qquad \widehat b_0=b,
\]
is supplied, then the rounded QR/back-substitution solution satisfies the
same local backward-error theorem as the explicit-validity concrete-dual
finite-max handoff:
\[
  {\rm LSQRSolveBackwardError}\!\left(
    n,\ A^{\mathsf T}A,\ A^{\mathsf T}b,\
    \fl_{\rm backSub}\bigl(fp,\widehat R,\widehat c\bigr),\
    c_G,\ c_g
  \right).
\]
The triangular validity guard \({\rm gammaValid}(fp,n)\) is derived from
\(n\le m\) by monotonicity.

For the sampled equation~(8) theorem, replace \(m\) by the sample size \(s\)
in the actual-unit guard and require the same displayed QR hypotheses
samplewise for the rounded sampled matrix and right-hand side.  If
\[
  s\,fp.u<1,
\]
then the high-probability rounded objective conclusion remains
\[
  1-\delta \le
  {\mathbb P}\!\left\{
    \|Ax(\omega)-b\|_2^2
    \le (1+\eta)\|Ax_\star-b\|_2^2
  \right\},
\]
with the same objective-budget, exact sampling law, and exact probability
measure as in the explicit-validity theorem.  No probability-construction,
normalization, alias-table, repair, or law-transfer error is introduced.
\end{corollary}


\begin{formalized}
\lean{Algorithms/LeastSquares/LSQRSolve.lean}
\leanfit{StoredQRSourceOffDiagonalControl.of_stored_trailing_sequence_diagDominant_signedStageUniformBudget_globalCompactBudget_activeMaxPivot_completedColumns_finiteMaxSmallness_stageRowMaxComparisonDefect_concreteDual_offdiag_rows_of_actualUnitRoundoff_no_gammaValid}
\leanfit{LSQRSolveBackwardError.of_stored_trailing_householder_sequence_topBlock_fl_backSub_gamma_bound_explicitCompactBudget_of_signed_alpha_diagDominant_signedStageUniformBudget_globalCompactBudget_activeMaxPivot_completedColumns_finiteMaxSmallness_stageRowMaxComparisonDefect_concreteDual_offdiag_rows_of_actualUnitRoundoff_no_gammaValid}
\lean{Algorithms/RandNLA/LeastSquaresSketch.lean}
\leanfit{leverageTraceProbability_eventProb_fl_rowSampleLSObjective_le_one_add_eta_of_augmentedSpan_sample_budget_of_objective_error_and_stored_qr_activePrefix_finiteMaxSmallness_activeMaxPivot_diagDominant_stageRowMaxComparisonDefect_concreteDual_solver_of_actualUnitRoundoff_no_gammaValid}
\end{formalized}

\begin{corollary}[Active scalar-comparison source-denominator cap handoff]
On the active-pivot diagonal-dominant scalar-comparison route, the local
source-control and solver surfaces need not expose the assembled finite-max
compact-product scalar involving the compact-sequence relative budget.  It is
enough to supply local diagonal dominance, the signed-stage recurrence budget,
active-pivot policy, the scalar comparison defect, source-denominator
nonbreakdown, \(fp.u\le U_{\mathrm{cap}}\),
\((m:\mathbb R)U_{\mathrm{cap}}<1\), and the canonical rational-gamma
cap-smallness inequality.
\end{corollary}


\begin{formalized}
\lean{Algorithms/LeastSquares/LSQRSolve.lean}
\leanfit{StoredQRSourceOffDiagonalControl.of_stored_trailing_sequence_diagDominant_signedStageUniformBudget_globalCompactBudget_activeMaxPivot_completedColumns_finiteMaxSourceDenURationalGammaCanonicalBounds_stageRowMaxComparisonDefect_offdiag_rows}
\leanfit{LSQRSolveBackwardError.of_stored_trailing_householder_sequence_topBlock_fl_backSub_gamma_bound_explicitCompactBudget_of_signed_alpha_diagDominant_signedStageUniformBudget_globalCompactBudget_activeMaxPivot_completedColumns_finiteMaxSourceDenURationalGammaCanonicalBounds_stageRowMaxComparisonDefect_offdiag_rows}
\end{formalized}

\begin{corollary}[Sampled active scalar-comparison source-denominator cap handoff]
The same source-denominator rational-gamma compact-product handoff lifts to the
sampled equation~(8) theorem.  On each sampled stored-QR trace, it is enough to
provide the active-pivot policy, local diagonal dominance, the scalar
comparison defect, signed-stage recurrence budgets, source-denominator
nonbreakdown, \(fp.u\le U_{\mathrm{cap}}\),
\((s:\mathbb R)U_{\mathrm{cap}}<1\), and the canonical rational-gamma
cap-smallness inequality.
\end{corollary}


\begin{formalized}
\lean{Algorithms/RandNLA/LeastSquaresSketch.lean}
\leanfit{leverageTraceProbability_eventProb_fl_rowSampleLSObjective_le_one_add_eta_of_augmentedSpan_sample_budget_of_objective_error_and_stored_qr_activePrefix_finiteMaxSourceDenURationalGammaCanonicalBounds_activeMaxPivot_diagDominant_stageRowMaxComparisonDefect_solver}
\end{formalized}

\begin{corollary}[Actual-unit source-denominator validity surface]
The active source-denominator handoff also has an actual-unit-roundoff form.
On the local source-control and solver surfaces, the auxiliary cap may be
specialized to \(U_{\mathrm{cap}}=fp.u\), and the validity guard is derived
from \((m:\mathbb R)fp.u<1\).  On the sampled equation~(8) surface, the same
validity guard is derived from \((s:\mathbb R)fp.u<1\).
\end{corollary}


\begin{formalized}
\lean{Algorithms/LeastSquares/LSQRSolve.lean}
\leanfit{StoredQRSourceOffDiagonalControl.of_stored_trailing_sequence_diagDominant_signedStageUniformBudget_globalCompactBudget_activeMaxPivot_completedColumns_finiteMaxSourceDenURationalGammaCanonicalBounds_stageRowMaxComparisonDefect_offdiag_rows_of_actualUnitRoundoff_no_gammaValid}
\leanfit{LSQRSolveBackwardError.of_stored_trailing_householder_sequence_topBlock_fl_backSub_gamma_bound_explicitCompactBudget_of_signed_alpha_diagDominant_signedStageUniformBudget_globalCompactBudget_activeMaxPivot_completedColumns_finiteMaxSourceDenURationalGammaCanonicalBounds_stageRowMaxComparisonDefect_offdiag_rows_of_actualUnitRoundoff_no_gammaValid}
\lean{Algorithms/RandNLA/LeastSquaresSketch.lean}
\leanfit{leverageTraceProbability_eventProb_fl_rowSampleLSObjective_le_one_add_eta_of_augmentedSpan_sample_budget_of_objective_error_and_stored_qr_activePrefix_finiteMaxSourceDenURationalGammaCanonicalBounds_activeMaxPivot_diagDominant_stageRowMaxComparisonDefect_solver_of_actualUnitRoundoff_no_gammaValid}
\end{formalized}

\begin{corollary}[Stored-lower active source-denominator nonbreakdown]
On the actual-unit active source-denominator branch, the raw Householder
source-denominator nonbreakdown hypothesis can be derived from the stored QR
trace.  It is enough to provide the stored panel recurrence, the signed-alpha
definition, and local diagonal dominance.
\end{corollary}


\begin{formalized}
\lean{Algorithms/LeastSquares/LSQRSolve.lean}
\leanfit{storedQRSourceDenominator_ne_zero_of_diagDominant_signedAlphaDef_stored_trailing_sequence}
\leanfit{StoredQRSourceOffDiagonalControl.of_stored_trailing_sequence_diagDominant_storedLower_signedStageUniformBudget_globalCompactBudget_activeMaxPivot_completedColumns_finiteMaxSourceDenURationalGammaCanonicalBounds_stageRowMaxComparisonDefect_offdiag_rows_of_actualUnitRoundoff_no_gammaValid}
\leanfit{LSQRSolveBackwardError.of_stored_trailing_householder_sequence_topBlock_fl_backSub_gamma_bound_explicitCompactBudget_of_signed_alpha_diagDominant_storedLower_signedStageUniformBudget_globalCompactBudget_activeMaxPivot_completedColumns_finiteMaxSourceDenURationalGammaCanonicalBounds_stageRowMaxComparisonDefect_offdiag_rows_of_actualUnitRoundoff_no_gammaValid}
\lean{Algorithms/RandNLA/LeastSquaresSketch.lean}
\leanfit{leverageTraceProbability_eventProb_fl_rowSampleLSObjective_le_one_add_eta_of_augmentedSpan_sample_budget_of_objective_error_and_stored_qr_activePrefix_storedLower_finiteMaxSourceDenURationalGammaCanonicalBounds_activeMaxPivot_diagDominant_stageRowMaxComparisonDefect_solver_of_actualUnitRoundoff_no_gammaValid}
\end{formalized}

\begin{lemma}[Active-pivot budgets do not imply nonpositive scalar comparison defect]
The scalar comparison condition
\[
  \Delta_{\mathrm{cmp}}\le 0
\]
is not a consequence of the current active-pivot active/off-diagonal budget
surface.  Adding the product compact-smallness hypotheses used by the
rectangular QR route still does not imply it.
\end{lemma}


\begin{formalized}
\lean{Algorithms/LeastSquares/LSQRSolve.lean}
\leanfit{not_forall_leadingBlock_upper_det_activeBlockPos_activeMaxPivot_activeBlockBudget_offdiagBudget_implies_stageRowMaxComparisonDefectBudget_nonpos}
\leanfit{not_forall_leadingBlock_upper_det_product_activeBlockPos_activeMaxPivot_activeBlockBudget_offdiagBudget_implies_stageRowMaxComparisonDefectBudget_nonpos}
\end{formalized}

\begin{lemma}[Row-max defect does not force the scalar comparison defect]
Even after the scalar row-max condition
\[
  \Delta_{\max}\le 0
\]
is supplied, the active-pivot active/off-diagonal budget surface does not imply
\(\Delta_{\mathrm{cmp}}\le 0\).  The same conclusion holds after adding the
product compact-smallness hypotheses.
\end{lemma}


\begin{formalized}
\lean{Algorithms/LeastSquares/LSQRSolve.lean}
\leanfit{activeMaxPivotRowMaxComparisonCounterexample_rowMaxDiagDefectBudget_nonpos}
\leanfit{activeMaxPivotRowMaxComparisonCounterexample_stageRowMaxComparisonDefectBudget_pos}
\leanfit{not_forall_leadingBlock_upper_det_activeBlockPos_activeMaxPivot_activeBlockBudget_offdiagBudget_rowMaxDiagDefect_implies_stageRowMaxComparisonDefectBudget_nonpos}
\leanfit{not_forall_leadingBlock_upper_det_product_activeBlockPos_activeMaxPivot_activeBlockBudget_offdiagBudget_rowMaxDiagDefect_implies_stageRowMaxComparisonDefectBudget_nonpos}
\end{formalized}

\begin{lemma}[Row-max defect alone does not force the stage-diagonal defect]
The scalar row-max condition \(\Delta_{\max}\le 0\) does not by itself imply
the downstream scalar stage-diagonal condition \(\Delta_{\mathrm{diag}}\le 0\).
\end{lemma}


\begin{formalized}
\lean{Algorithms/LeastSquares/LSQRSolve.lean}
\leanfit{activeMaxPivotRowMaxComparisonCounterexample_stageDiagLowerDefectBudget_pos}
\leanfit{not_forall_rowMaxDiagDefectBudget_implies_stageDiagLowerDefectBudget_nonpos}
\end{formalized}

\begin{lemma}[Product-active row-max data still do not force the stage-diagonal defect]
The active-pivot budget surface, even strengthened by compact-product
smallness and a granted scalar row-max defect \(\Delta_{\max}\le 0\), does not
imply the downstream scalar stage-diagonal condition
\(\Delta_{\mathrm{diag}}\le 0\).
\end{lemma}


\begin{formalized}
\lean{Algorithms/LeastSquares/LSQRSolve.lean}
\leanfit{not_forall_leadingBlock_upper_det_activeBlockPos_activeMaxPivot_activeBlockBudget_offdiagBudget_rowMaxDiagDefect_implies_stageDiagLowerDefectBudget_nonpos}
\leanfit{not_forall_leadingBlock_upper_det_product_activeBlockPos_activeMaxPivot_activeBlockBudget_offdiagBudget_rowMaxDiagDefect_implies_stageDiagLowerDefectBudget_nonpos}
\end{formalized}

\begin{lemma}[Active-pivot budgets do not imply nonpositive row-max defect]
The scalar row-max defect condition
\[
  \Delta_{\max}\le 0
\]
is not a consequence of the current active-pivot budget surface.  Even if
displayed leading blocks are upper triangular and nonsingular, the active block
has positive mass, the displayed pivot satisfies active max-pivoting, and one
nonnegative budget controls both active-block entries and displayed
strict-upper entries, the finite row-max defect can be positive.
\end{lemma}


\begin{formalized}
\lean{Algorithms/LeastSquares/LSQRSolve.lean}
\leanfit{activeMaxPivotRowBudgetDiagCounterexample_rowMaxDiagDefectBudget_pos}
\leanfit{not_forall_leadingBlock_upper_det_activeBlockPos_activeMaxPivot_activeBlockBudget_offdiagBudget_implies_rowMaxDiagDefectBudget_nonpos}
\end{formalized}

\begin{lemma}[Product smallness does not imply nonpositive row-max defect]
The product compact-smallness condition used by the diagonal-dominant
rectangular QR route cannot replace the row-max scalar defect condition.  Even
with upper-triangular nonsingular displayed leading blocks and a satisfiable
product compact-smallness inequality, it is possible that
\[
  \Delta_{\max}>0.
\]
\end{lemma}


\begin{formalized}
\lean{Algorithms/LeastSquares/LSQRSolve.lean}
\leanfit{not_forall_upper_tri_det_ne_zero_product_budget_implies_rowMaxDiagDefectBudget_nonpos}
\end{formalized}

\begin{lemma}[Product smallness does not imply active row-max defect control]
Even after adding product compact-smallness to the active-pivot budget surface,
the nonpositive row-max scalar defect condition can fail.  Thus the product
condition and the current active-pivot budget fields cannot be combined as a
hidden proof of
\[
  \Delta_{\max}\le 0.
\]
\end{lemma}


\begin{formalized}
\lean{Algorithms/LeastSquares/LSQRSolve.lean}
\leanfit{not_forall_leadingBlock_upper_det_product_activeBlockPos_activeMaxPivot_activeBlockBudget_offdiagBudget_implies_rowMaxDiagDefectBudget_nonpos}
\end{formalized}

\begin{lemma}[Product smallness does not imply nonpositive stage-diagonal defect]
The same product compact-smallness condition also cannot replace the scalar
stage-diagonal lower-bound condition.  With the constant stage budget \(B_k=2\),
there are upper-triangular nonsingular displayed leading blocks satisfying the
compact-product inequality while
\[
  \Delta_{\mathrm{stage}}>0.
\]
\end{lemma}


\begin{formalized}
\lean{Algorithms/LeastSquares/LSQRSolve.lean}
\leanfit{activeMaxPivotRowBudgetDiagCounterexample_stageDiagLowerDefectBudget_pos}
\leanfit{not_forall_upper_tri_det_ne_zero_product_budget_implies_stageDiagLowerDefectBudget_nonpos}
\end{formalized}

\begin{lemma}[Product smallness does not imply active stage-diagonal defect control]
Even after adding product compact-smallness to the active-pivot budget surface,
the scalar stage-diagonal lower-bound condition can fail.  Thus the product
condition and the current active-pivot budget fields cannot be combined as a
hidden proof of
\[
  \Delta_{\mathrm{stage}}\le 0.
\]
\end{lemma}


\begin{formalized}
\lean{Algorithms/LeastSquares/LSQRSolve.lean}
\leanfit{not_forall_leadingBlock_upper_det_product_activeBlockPos_activeMaxPivot_activeBlockBudget_offdiagBudget_implies_stageDiagLowerDefectBudget_nonpos}
\end{formalized}

\begin{lemma}[Active-pivot budgets do not imply nonpositive stage-diagonal defect]
The active-pivot budget surface alone also cannot imply the scalar
stage-diagonal lower-bound condition.  Active max-pivoting and
active/off-diagonal budget control therefore cannot be used as a hidden proof
of
\[
  \Delta_{\mathrm{stage}}\le 0.
\]
\end{lemma}


\begin{formalized}
\lean{Algorithms/LeastSquares/LSQRSolve.lean}
\leanfit{not_forall_leadingBlock_upper_det_activeBlockPos_activeMaxPivot_activeBlockBudget_offdiagBudget_implies_stageDiagLowerDefectBudget_nonpos}
\end{formalized}

\begin{lemma}[Product smallness does not imply the active stage-budget row-max comparison]
Even after adding product compact-smallness to the active-pivot budget surface,
the comparison \(B_k\le R^{\max}_{k,i}\) required by the row-max bridge can
fail.  Thus the product condition cannot be used as a hidden proof of this
comparison.
\end{lemma}


\begin{formalized}
\lean{Algorithms/LeastSquares/LSQRSolve.lean}
\leanfit{not_forall_leadingBlock_upper_det_product_activeBlockPos_activeMaxPivot_activeBlockBudget_offdiagBudget_implies_stageBudget_le_rowMax}
\end{formalized}

\begin{lemma}[Active-pivot budgets do not imply the stage-budget row-max comparison]
The comparison \(B_k\le R^{\max}_{k,i}\) required by the preceding row-max
bridge is not a consequence of the current active-pivot budget surface.  Even
if displayed leading blocks are upper triangular and nonsingular, the active
block has positive mass, the displayed pivot satisfies active max-pivoting, and
one uniform stage budget controls both active-block entries and displayed
strict-upper entries, the conclusion
\[
  B_k\le R^{\max}_{k,i}
\]
can fail.
\end{lemma}


\begin{formalized}
\lean{Algorithms/LeastSquares/LSQRSolve.lean}
\leanfit{not_forall_leadingBlock_upper_det_activeBlockPos_activeMaxPivot_activeBlockBudget_offdiagBudget_implies_stageBudget_le_rowMax}
\end{formalized}

\begin{corollary}[Visible row-max assumptions for the active-pivot stage-diagonal QR route]
In the active max-pivot setting, the local source-control and solver theorems
can expose row-max data instead of exposing the scalar stage-diagonal condition
directly.  Assume the finite row-max defect is nonpositive and that the uniform
stage budget is bounded by the displayed strict-upper row maximum on every
displayed off-diagonal row:
\[
  B_k\le R^{\max}_{k,i}\qquad (i<k).
\]
Then the scalar stage-diagonal lower-bound condition follows internally, and
the active-pivot stage-diagonal QR route applies with the same stored
recurrence, completed-column, global compact-step, determinant,
\(\kappa_\infty\)/dual-budget, and global-product hypotheses.
\end{corollary}

\begin{remark}
This corollary does not prove the row-max scalar defect or the comparison
\(B_k\le R^{\max}_{k,i}\) for a concrete pivoted loop.  It makes those two
fields visible on the theorem surface and prevents the downstream
source-control/solver proof from re-assuming the stronger scalar
\(\Delta_{\mathrm{stage}}\le 0\) condition.  Its actual-unit-roundoff
source-control and solver siblings also remove the local abstract validity
fields, deriving them from \((m:\mathbb R)u<1\), while leaving the same
row-max, determinant, conditioning, dual-budget, and product-smallness fields
visible.
\end{remark}

\begin{formalized}
\lean{Algorithms/LeastSquares/LSQRSolve.lean}
\leanfit{StoredQRSourceOffDiagonalControl.of_stored_trailing_sequence_leadingBlock_det_ne_zero_kappaInf_selfNorm_dualBudget_signedStageUniformBudget_globalCompactBudget_activeMaxPivot_completedColumns_globalProduct_rowMaxDiagDefect_stageBudgetLeRowMax_offdiag_rows}
\leanfit{StoredQRSourceOffDiagonalControl.of_stored_trailing_sequence_leadingBlock_det_ne_zero_kappaInf_selfNorm_dualBudget_signedStageUniformBudget_globalCompactBudget_activeMaxPivot_completedColumns_globalProduct_rowMaxDiagDefect_stageBudgetLeRowMax_offdiag_rows_of_actualUnitRoundoff_no_gammaValid}
\leanfit{LSQRSolveBackwardError.of_stored_trailing_householder_sequence_topBlock_fl_backSub_gamma_bound_explicitCompactBudget_of_signed_alpha_leadingBlock_det_ne_zero_kappaInf_selfNorm_dualBudget_signedStageUniformBudget_globalCompactBudget_activeMaxPivot_completedColumns_globalProduct_rowMaxDiagDefect_stageBudgetLeRowMax_offdiag_rows}
\leanfit{LSQRSolveBackwardError.of_stored_trailing_householder_sequence_topBlock_fl_backSub_gamma_bound_explicitCompactBudget_of_signed_alpha_leadingBlock_det_ne_zero_kappaInf_selfNorm_dualBudget_signedStageUniformBudget_globalCompactBudget_activeMaxPivot_completedColumns_globalProduct_rowMaxDiagDefect_stageBudgetLeRowMax_offdiag_rows_of_actualUnitRoundoff_no_gammaValid}
\end{formalized}

\begin{corollary}[Active max-pivot stage-diagonal QR solve]
In the setting of the preceding corollary, suppose additionally that every
displayed pivot column is chosen by the finite active max-pivot selector:
\[
  k=\operatorname{householderActiveMaxPivotColumn}(k,A_k).
\]
Then the raw pivot-maximality inequality required by the active/prefix QR
source-control theorem follows internally from that policy equation.  Hence the
local solver can consume the scalar stage-diagonal condition
\(\Delta_{\mathrm{stage}}\le 0\), the finite active-pivot policy, the stored
recurrence, the global compact-step recurrence, the local determinant and
\(\kappa_\infty\)/dual-budget nonbreakdown fields, and the scalar global-product
condition, without a separate raw pivot-maximality hypothesis.
\end{corollary}

\begin{remark}
This is a pivot-policy reduction only.  It does not prove
\(\Delta_{\mathrm{stage}}\le 0\), determinant or conditioning facts,
nonbreakdown, compact-product smallness, or the final QR/preconditioner theorem
from a concrete implementation.
\end{remark}

\begin{formalized}
\lean{Algorithms/LeastSquares/LSQRSolve.lean}
\lean{storedQRActiveMaxPivotColumn_pivotMax}
\leanfit{StoredQRSourceOffDiagonalControl.of_stored_trailing_sequence_leadingBlock_det_ne_zero_normSqBudget_signedStageUniformBudget_product_of_normSqBudget_activePrefix_activeBlockRecurrence_globalCompactBudget_activeMaxPivot_completedColumns_globalProduct_stageDiagDefect_offdiag_rows}
\leanfit{StoredQRSourceOffDiagonalControl.of_stored_trailing_sequence_leadingBlock_det_ne_zero_kappaInf_selfNorm_dualBudget_signedStageUniformBudget_globalCompactBudget_activeMaxPivot_completedColumns_globalProduct_stageDiagDefect_offdiag_rows}
\leanfit{LSQRSolveBackwardError.of_stored_trailing_householder_sequence_topBlock_fl_backSub_gamma_bound_explicitCompactBudget_of_signed_alpha_leadingBlock_det_ne_zero_kappaInf_selfNorm_dualBudget_signedStageUniformBudget_globalCompactBudget_activeMaxPivot_completedColumns_globalProduct_stageDiagDefect_offdiag_rows}
\end{formalized}

\begin{lemma}[Exact no-pivot QR does not imply nonpositive stage-diagonal defect]
There is a two-by-two exact unpivoted trailing Householder QR sequence
\(\widehat A_k\), with valid signed Householder squared-norm identities and
nonzero Householder denominators at each of its two stages, for which the
constant stage budget \(B_k=2\) has
\[
  \Delta_{\mathrm{stage}}(\widehat A,B)>0 .
\]
Consequently, there is no universal theorem of the form
\[
\begin{aligned}
&\text{exact no-pivot trailing Householder recurrence} \\
&\quad+\ \text{valid Householder squared norms} \\
&\quad+\ \text{nonzero denominators} \\
&\Longrightarrow\ \Delta_{\mathrm{stage}}(\widehat A,B)\le 0
   \qquad(B_k\equiv 2).
\end{aligned}
\]
\end{lemma}


\begin{formalized}
\lean{Algorithms/LeastSquares/LSQRSolve.lean}
\leanfit{exactHouseholderQRDiagDominanceCounterexample_stageDiagLowerDefectBudget_pos}
\leanfit{not_forall_exact_trailing_householder_sequence_implies_stageDiagLowerDefectBudget_nonpos}
\end{formalized}

\begin{lemma}[Exact no-pivot recurrence does not force the comparison scalar]
For the same exact no-pivot two-stage Householder witness, the constant stage
budget \(B_k\equiv 3\) gives
\[
  \Delta_{\mathrm{cmp}}(\widehat A,B)>0 .
\]
Consequently, there is no universal theorem of the form
\[
\begin{aligned}
&\text{exact no-pivot trailing Householder recurrence} \\
&\quad+\ \text{valid Householder squared norms} \\
&\quad+\ \text{nonzero denominators} \\
&\quad+\ \text{nonnegative stage budgets} \\
&\Longrightarrow\ \Delta_{\mathrm{cmp}}(\widehat A,B)\le 0.
\end{aligned}
\]
\end{lemma}


\begin{formalized}
\lean{Algorithms/LeastSquares/LSQRSolve.lean}
\leanfit{exactHouseholderQRDiagDominanceCounterexample_stageRowMaxComparisonDefectBudget_pos}
\leanfit{not_forall_exact_trailing_householder_sequence_implies_stageRowMaxComparisonDefectBudget_nonpos}
\end{formalized}

\begin{corollary}[Diagonal dominance implies nonpositive scalar defect]
If every displayed leading block \(S_k\) is locally diagonally dominant in the
repository sense \(\lean{IsDiagDominantUpper}\), then
\[
  \Delta_{\max}\le 0.
\]
Consequently the canonical row-max budget gives the packaged
\(\lean{StoredQRDisplayedRowBudgetControl}\) certificate.
\end{corollary}


\begin{formalized}
\lean{Algorithms/LeastSquares/LSQRSolve.lean}
\leanfit{storedQRRowMaxDiagDefectBudget_nonpos_of_diagDominant}
\leanfit{StoredQRDisplayedRowBudgetControl.of_diagDominant}
\end{formalized}

\begin{lemma}[Exact no-pivot QR does not imply nonpositive scalar defect]
There is a two-by-two exact unpivoted trailing Householder QR sequence
\(\widehat A_k\), with valid signed Householder squared-norm identities and
nonzero Householder denominators at each of its two stages, for which
\[
  \Delta_{\max}(\widehat A)>0 .
\]
Consequently, there is no universal theorem of the form
\[
\begin{aligned}
&\text{exact no-pivot trailing Householder recurrence} \\
&\quad+\ \text{valid Householder squared norms} \\
&\quad+\ \text{nonzero denominators} \\
&\Longrightarrow\ \Delta_{\max}(\widehat A)\le 0 .
\end{aligned}
\]
\end{lemma}


\begin{formalized}
\lean{Algorithms/LeastSquares/LSQRSolve.lean}
\leanfit{exactHouseholderQRDiagDominanceCounterexample_rowMaxDiagDefectBudget_pos}
\leanfit{not_forall_exact_trailing_householder_sequence_implies_rowMaxDiagDefectBudget_nonpos}
\end{formalized}

\begin{corollary}[Global compact-budget row-budget certificate]
Let \(B_t\) be a monotone signed-stage budget sequence for the stored QR
panel.  Suppose the initial sampled matrix entries are bounded by \(B_0\), and
suppose the single finite compact-step recurrence
\[
  c_m B_t+\Gamma_t \le B_{t+1}
\]
holds at every stage \(t<n\), where \(c_m\) is the Cox--Higham active-row
growth factor and \(\Gamma_t\) is the finite maximum of the rounded
Householder compact-update component budget at stage \(t\).  If the signed
Householder stages satisfy pivot maximality and the norm-square nonbreakdown
margin, and if the same diagonal lower-bound field
\[
  B_k \le |(S_k)_{ii}|,\qquad i<k,
\]
holds, then the packaged displayed row-budget certificate
\(\mathcal R(A,B)\) follows.

Moreover, the norm-square nonbreakdown margin can be replaced by the local
\(\kappa_\infty\)/dual-budget hypotheses for the leading block \(S_k\): local
leading-block nonsingularity, a bound on \(\kappa_\infty(S_k)\), a self-norm
budget \(K_k\), and the displayed dual compact-budget inequality.

There is also an active-pivot-policy version.  Instead of assuming the raw
pivot-maximality inequality at each stage, assume the displayed pivot column is
the finite active max-pivot selector for the current active block.  The same
row-budget certificates follow, with or without the \(\kappa_\infty\)/dual
budget replacement of the norm-square nonbreakdown margin.
\end{corollary}


\begin{formalized}
\lean{Algorithms/LeastSquares/LSQRSolve.lean}
\leanfit{StoredQRDisplayedRowBudgetControl.of_signed_stage_uniformBudget_globalCompactBudget_of_normSqBudget}
\leanfit{StoredQRDisplayedRowBudgetControl.of_signed_stage_uniformBudget_globalCompactBudget_kappaInf_dualBudget}
\leanfit{StoredQRDisplayedRowBudgetControl.of_signed_stage_uniformBudget_globalCompactBudget_activeMaxPivot_of_normSqBudget}
\leanfit{StoredQRDisplayedRowBudgetControl.of_signed_stage_uniformBudget_globalCompactBudget_activeMaxPivot_kappaInf_dualBudget}
\end{formalized}

\begin{lemma}[Cox--Higham stage budgets give displayed off-diagonal row budgets]
Fix a stored Householder panel sequence \(A_t\), a nonnegative growth factor
\(c\), and displayed row budgets \(R_{k,i}\).  For each upper displayed entry
\((i,j)\) of the \(k\)th leading block, let
\[
  \tau(k,j)=
  \begin{cases}
    j+1, & j<k,\\
    k, & j=k.
  \end{cases}
\]
Assume a scalar Cox--Higham stage-budget recurrence for the same entry through
stage \(\tau(k,j)\): the initial entry is below \(B_0\), each exact
same-reflector step is below \(cB_t\), the concrete floating-point compact
Householder component budget is absorbed by \(B_{t+1}\), and the pivot-column
zeroing equation holds when the entry's column is the active pivot.  If
\(B_{\tau(k,j)}\le R_{k,i}\), then
\[
  |(S_k)_{ij}| \le R_{k,i}, \qquad 0\le i<j\le k<n .
\]
\end{lemma}


\begin{formalized}
\lean{Algorithms/QR/HouseholderQR.lean}
\leanfit{qrLeadingOffdiagStop}
\leanfit{storedQRSignedStageVector}
\leanfit{storedQRSignedStageBeta}
\leanfit{one_le_coxHighamActiveRowGrowthFactor}
\leanfit{storedQRSignedStage_exact_same_reflector_bound_of_prefix_or_active_stage_bounds}
\leanfit{fl_householderStoredPanel_sequence_completed_column_eq_pivot_succ}
\leanfit{fl_householderStoredPanel_sequence_leadingBlock_offdiag_budget_of_exact_stage_budgets_factor}
\leanfit{fl_householderStoredPanel_sequence_leadingBlock_offdiag_budget_of_signed_stage_budgets_factor}
\leanfit{storedQRSignedStage_pivot_column_zero_below_of_trailingNorm_pos}
\leanfit{storedQRSignedStage_pivot_zeroing_field_of_trailingNorm_pos}
\leanfit{storedQRSignedStage_pivot_zeroing_field_of_normSqBudget}
\leanfit{StoredQRSourceOffDiagonalControl.of_stored_trailing_sequence_leadingBlock_det_ne_zero_normSqBudget_signedStageUniformBudget_product_of_normSqBudget_stage_entry_bounds}
\leanfit{LSQRSolveBackwardError.of_stored_trailing_householder_sequence_topBlock_fl_backSub_gamma_bound_explicitCompactBudget_of_signed_alpha_leadingBlock_det_ne_zero_normSqBudget_signedStageUniformBudget_product_of_normSqBudget_stage_entry_bounds}
\end{formalized}

\begin{corollary}[Source-control certificate from stage budgets and diagonal lower bounds]
Keep the determinant, norm-square nonbreakdown, and compact-product hypotheses
of the norm-square source-shaped QR route.  Suppose also that the displayed
upper entries of each leading block are equipped with Cox--Higham stage
budgets \(B^{(k,i,j)}_t\) satisfying the recurrence in the previous lemma, and
that their terminal values are dominated by row budgets \(R_{k,i}\):
\[
  B^{(k,i,j)}_{\tau(k,j)}\le R_{k,i},\qquad 0\le i<j\le k<n .
\]
If the row budgets satisfy the diagonal lower-bound condition
\[
  R_{k,i}\le |(S_k)_{ii}|,\qquad 0\le i<k<n ,
\]
then the stored QR sequence satisfies
\[
  \lean{StoredQRSourceOffDiagonalControl}.
\]
Consequently, the corresponding computed triangular solve satisfies the local
\(\lean{LSQRSolveBackwardError}\) certificate with
\(\lean{qrSolveFinalGramBudget}\) and
\(\lean{qrSolveFinalRhsBudget}\).
\end{corollary}


\begin{formalized}
\lean{Algorithms/LeastSquares/LSQRSolve.lean}
\leanfit{StoredQRSourceOffDiagonalControl.of_stored_trailing_sequence_leadingBlock_det_ne_zero_normSqBudget_stageBudget_rowBudget_product}
\leanfit{LSQRSolveBackwardError.of_stored_trailing_householder_sequence_topBlock_fl_backSub_gamma_bound_explicitCompactBudget_of_signed_alpha_leadingBlock_det_ne_zero_normSqBudget_stageBudget_rowBudget_product}
\leanfit{StoredQRSourceOffDiagonalControl.of_stored_trailing_sequence_leadingBlock_det_ne_zero_normSqBudget_stageBudget_rowBudget_product_of_offdiag_rows}
\leanfit{LSQRSolveBackwardError.of_stored_trailing_householder_sequence_topBlock_fl_backSub_gamma_bound_explicitCompactBudget_of_signed_alpha_leadingBlock_det_ne_zero_normSqBudget_stageBudget_rowBudget_product_of_offdiag_rows}
\leanfit{StoredQRSourceOffDiagonalControl.of_stored_trailing_sequence_leadingBlock_det_ne_zero_normSqBudget_signedStageBudget_rowBudget_product}
\leanfit{LSQRSolveBackwardError.of_stored_trailing_householder_sequence_topBlock_fl_backSub_gamma_bound_explicitCompactBudget_of_signed_alpha_leadingBlock_det_ne_zero_normSqBudget_signedStageBudget_rowBudget_product}
\leanfit{StoredQRSourceOffDiagonalControl.of_stored_trailing_sequence_leadingBlock_det_ne_zero_normSqBudget_signedStageBudget_rowBudget_product_of_offdiag_rows}
\leanfit{LSQRSolveBackwardError.of_stored_trailing_householder_sequence_topBlock_fl_backSub_gamma_bound_explicitCompactBudget_of_signed_alpha_leadingBlock_det_ne_zero_normSqBudget_signedStageBudget_rowBudget_product_of_offdiag_rows}
\leanfit{StoredQRSourceOffDiagonalControl.of_stored_trailing_sequence_leadingBlock_det_ne_zero_normSqBudget_signedStageBudget_rowBudget_product_of_normSqBudget}
\leanfit{LSQRSolveBackwardError.of_stored_trailing_householder_sequence_topBlock_fl_backSub_gamma_bound_explicitCompactBudget_of_signed_alpha_leadingBlock_det_ne_zero_normSqBudget_signedStageBudget_rowBudget_product_of_normSqBudget}
\leanfit{StoredQRSourceOffDiagonalControl.of_stored_trailing_sequence_leadingBlock_det_ne_zero_normSqBudget_signedStageBudget_rowBudget_product_of_normSqBudget_offdiag_rows}
\leanfit{LSQRSolveBackwardError.of_stored_trailing_householder_sequence_topBlock_fl_backSub_gamma_bound_explicitCompactBudget_of_signed_alpha_leadingBlock_det_ne_zero_normSqBudget_signedStageBudget_rowBudget_product_of_normSqBudget_offdiag_rows}
\leanfit{StoredQRSourceOffDiagonalControl.of_stored_trailing_sequence_leadingBlock_det_ne_zero_normSqBudget_rowBudgetControl_globalProduct}
\leanfit{qrLeadingOffdiagStop_le}
\leanfit{StoredQRSourceOffDiagonalControl.of_stored_trailing_sequence_leadingBlock_det_ne_zero_normSqBudget_signedStageUniformBudget_product_of_normSqBudget}
\leanfit{LSQRSolveBackwardError.of_stored_trailing_householder_sequence_topBlock_fl_backSub_gamma_bound_explicitCompactBudget_of_signed_alpha_leadingBlock_det_ne_zero_normSqBudget_signedStageUniformBudget_product_of_normSqBudget}
\leanfit{StoredQRSourceOffDiagonalControl.of_stored_trailing_sequence_leadingBlock_det_ne_zero_normSqBudget_signedStageUniformBudget_product_of_normSqBudget_offdiag_rows}
\leanfit{LSQRSolveBackwardError.of_stored_trailing_householder_sequence_topBlock_fl_backSub_gamma_bound_explicitCompactBudget_of_signed_alpha_leadingBlock_det_ne_zero_normSqBudget_signedStageUniformBudget_product_of_normSqBudget_offdiag_rows}
\leanfit{StoredQRSourceOffDiagonalControl.of_stored_trailing_sequence_leadingBlock_det_ne_zero_normSqBudget_signedStageUniformBudget_product_of_normSqBudget_stage_entry_bounds_offdiag_rows}
\leanfit{LSQRSolveBackwardError.of_stored_trailing_householder_sequence_topBlock_fl_backSub_gamma_bound_explicitCompactBudget_of_signed_alpha_leadingBlock_det_ne_zero_normSqBudget_signedStageUniformBudget_product_of_normSqBudget_stage_entry_bounds_offdiag_rows}
\end{formalized}

\begin{corollary}[Active/prefix split for the stage-entry QR certificate]
In the setting of the preceding stage-entry certificate, suppose the stage
entry control is supplied in the following source-shaped form.  For every
stage \(t<n\), all entries in the active suffix obey the single block budget
\[
  |A_t(r,\ell)|\le B_t
  \qquad (t\le r,\; t\le \ell).
\]
For a displayed strict upper entry \(i<j\) in the \(k\)th leading block, the
only separate row-growth hypothesis is the prefix case
\[
  |A_t(\rho_{k,i},\ell)|\le B_t
  \qquad
  (t<\sigma_{k,j},\; t\le \ell,\; \rho_{k,i}<t),
\]
where \(\rho_{k,i}\) is the stored row index of displayed row \(i\), and
\(\sigma_{k,j}=\lean{qrLeadingOffdiagStop}(j)\).  Together with the same
signed-alpha recurrence, norm-square nonbreakdown, pivot-maximality, monotone
stage budget, offdiag-row-only diagonal lower bounds \(B_k\le |(S_k)_{ii}|\)
for \(i<k\), and compact-product smallness as above, these hypotheses imply
the same stored-QR source-control certificate and hence the same
\(\lean{LSQRSolveBackwardError}\) conclusion.
\end{corollary}


\begin{formalized}
\lean{Algorithms/LeastSquares/LSQRSolve.lean}
\leanfit{StoredQRSourceOffDiagonalControl.of_stored_trailing_sequence_leadingBlock_det_ne_zero_normSqBudget_signedStageUniformBudget_product_of_normSqBudget_activePrefix_stage_entry_bounds_offdiag_rows}
\leanfit{LSQRSolveBackwardError.of_stored_trailing_householder_sequence_topBlock_fl_backSub_gamma_bound_explicitCompactBudget_of_signed_alpha_leadingBlock_det_ne_zero_normSqBudget_signedStageUniformBudget_product_of_normSqBudget_activePrefix_stage_entry_bounds_offdiag_rows}
\end{formalized}

\begin{corollary}[Active-block recurrence for the active/prefix QR certificate]
In the setting of the active/prefix stage-entry certificate, the active-suffix
block budget
\[
  |A_t(r,\ell)|\le B_t
  \qquad (t\le r,\;t\le \ell)
\]
need not be supplied as an independent hypothesis.  It follows from the stored
signed-stage recurrence, norm-square nonbreakdown, pivot-column maximality,
completed old-column preservation, the initial full-block bound
\[
  |A_0(r,\ell)|\le B_0,
\]
and the Cox--Higham active-block recurrence inequality
\[
  \Gamma_m B_t+
  \Delta_t(r,\ell)\le B_{t+1}
  \qquad (t+1<n,\;t+1\le r,\;t+1\le \ell),
\]
where \(\Gamma_m=\lean{coxHighamActiveRowGrowthFactor}(m)\) and
\(\Delta_t(r,\ell)\) is the corresponding compact Householder component budget.
Consequently the same stored-QR source-control certificate and
\(\lean{LSQRSolveBackwardError}\) conclusion follow with the active block
propagated by recurrence.  The separate prefix displayed-row budget, diagonal
lower bounds, determinant/norm-square nonbreakdown, and compact-product
smallness remain explicit.
\end{corollary}


\begin{formalized}
\lean{Algorithms/LeastSquares/LSQRSolve.lean}
\leanfit{StoredQRSourceOffDiagonalControl.of_stored_trailing_sequence_leadingBlock_det_ne_zero_normSqBudget_signedStageUniformBudget_product_of_normSqBudget_activePrefix_activeBlockRecurrence_offdiag_rows}
\leanfit{LSQRSolveBackwardError.of_stored_trailing_householder_sequence_topBlock_fl_backSub_gamma_bound_explicitCompactBudget_of_signed_alpha_leadingBlock_det_ne_zero_normSqBudget_signedStageUniformBudget_product_of_normSqBudget_activePrefix_activeBlockRecurrence_offdiag_rows}
\end{formalized}

\begin{corollary}[Prefix-row recurrence for the active/prefix QR certificate]
In the setting of the active-block recurrence certificate, the prefix displayed
row budget
\[
  |A_t(\rho,\ell)|\le B_t
  \qquad (\rho<t,\;t\le \ell)
\]
need not be supplied as a raw stage hypothesis.  It follows by induction from
the active-block bound, the signed-stage exact same-reflector prefix/active
split, and the one-step prefix recurrence
\[
  \Gamma_m B_t+\Delta_t(\rho,\ell)\le B_{t+1}
  \qquad (\rho<t+1,\;t\le \ell),
\]
where \(\Gamma_m=\lean{coxHighamActiveRowGrowthFactor}(m)\).  Consequently the
stored-QR source-control certificate and the corresponding
\(\lean{LSQRSolveBackwardError}\) conclusion follow with both active and prefix
stage-entry fields generated from one-step recurrences.
\end{corollary}


\begin{formalized}
\lean{Algorithms/LeastSquares/LSQRSolve.lean}
\leanfit{storedQRSignedStage_active_block_bound_of_signed_stage_budget}
\leanfit{storedQRSignedStage_prefix_row_bound_of_active_block_and_prefix_budget}
\leanfit{StoredQRSourceOffDiagonalControl.of_stored_trailing_sequence_leadingBlock_det_ne_zero_normSqBudget_signedStageUniformBudget_product_of_normSqBudget_activePrefix_activeBlockRecurrence_prefixRowRecurrence_offdiag_rows}
\leanfit{LSQRSolveBackwardError.of_stored_trailing_householder_sequence_topBlock_fl_backSub_gamma_bound_explicitCompactBudget_of_signed_alpha_leadingBlock_det_ne_zero_normSqBudget_signedStageUniformBudget_product_of_normSqBudget_activePrefix_activeBlockRecurrence_prefixRowRecurrence_offdiag_rows}
\end{formalized}

\begin{corollary}[Finite global compact-step budget for active/prefix QR]
For a signed stored-QR stage \(t<n\), define
\[
  \Delta_t^{\max}=\max_{r,\ell}
    \Delta_t(r,\ell),
\]
where \(\Delta_t(r,\ell)\) is the compact Householder floating-point component
budget for entry \((r,\ell)\) at stage \(t\).  If the scalar recurrence
\[
  \Gamma_m B_t+\Delta_t^{\max}\le B_{t+1}
\]
holds at every stage, then all three one-step budget fields used by the
active/prefix QR certificate follow: the displayed off-diagonal row recurrence,
the active-suffix block recurrence, and the prefix displayed-row recurrence.
Consequently the stored-QR source-control certificate and the corresponding
\(\lean{LSQRSolveBackwardError}\) conclusion follow without separately
assuming these three one-step budget families.
\end{corollary}


\begin{formalized}
\lean{Algorithms/LeastSquares/LSQRSolve.lean}
\leanfit{storedQRSignedStageGlobalCompactBudget}
\leanfit{storedQRSignedStage_compact_component_le_globalBudget}
\leanfit{storedQRSignedStage_active_prefix_budgets_of_globalCompactBudget}
\leanfit{StoredQRSourceOffDiagonalControl.of_stored_trailing_sequence_leadingBlock_det_ne_zero_normSqBudget_signedStageUniformBudget_product_of_normSqBudget_activePrefix_activeBlockRecurrence_globalCompactBudget_offdiag_rows}
\leanfit{LSQRSolveBackwardError.of_stored_trailing_householder_sequence_topBlock_fl_backSub_gamma_bound_explicitCompactBudget_of_signed_alpha_leadingBlock_det_ne_zero_normSqBudget_signedStageUniformBudget_product_of_normSqBudget_activePrefix_activeBlockRecurrence_globalCompactBudget_offdiag_rows}
\end{formalized}

\begin{corollary}[Monotone signed-stage budgets from the global recurrence]
Assume the signed-stage finite global compact-step recurrence
\[
  \Gamma_m B_t+\Delta_t^{\max}\le B_{t+1},\qquad t<n,
\]
and \(B_t\ge 0\) for all \(t\).  Then the stage budget is monotone on the
algorithmic QR horizon:
\[
  a\le b\le n \quad\Longrightarrow\quad B_a\le B_b .
\]
This is only a stage-range consequence of the recurrence; it does not assert
any behavior for budget indices beyond \(n\).
\end{corollary}


\begin{formalized}
\lean{Algorithms/LeastSquares/LSQRSolve.lean}
\leanfit{storedQRSignedStageGlobalCompactBudget_nonneg}
\leanfit{storedQRSignedStageBudget_mono_on_stages_of_globalCompactBudget}
\end{formalized}

\begin{corollary}[Horizon-clamped global-product QR budget surface]
In the completed-column, global-product active/prefix QR route, the separate
hypothesis that the stage budget is monotone for all natural indices may be
removed.  Define the horizon-clamped budget
\[
  \widetilde B_t =
  \begin{cases}
    B_t, & t\le n,\\
    B_n, & t>n.
  \end{cases}
\]
The recurrence above proves that \(\widetilde B_t\) is globally monotone, while
\(\widetilde B_t=B_t\) at every algorithmic QR stage.  Therefore the existing
completed-column global-product source-control and solver certificates apply
with no assumption about budget behavior beyond the QR horizon.
\end{corollary}


\begin{formalized}
\lean{Algorithms/LeastSquares/LSQRSolve.lean}
\leanfit{qrStageHorizonBudget}
\leanfit{qrStageHorizonBudget_eq_of_le}
\leanfit{qrStageHorizonBudget_nonneg}
\leanfit{qrStageHorizonBudget_mono_of_mono_on_stages}
\leanfit{StoredQRSourceOffDiagonalControl.of_stored_trailing_sequence_leadingBlock_det_ne_zero_normSqBudget_signedStageUniformBudget_product_of_normSqBudget_activePrefix_activeBlockRecurrence_globalCompactBudget_completedColumns_globalProduct_offdiag_rows_of_horizonBudget}
\leanfit{LSQRSolveBackwardError.of_stored_trailing_householder_sequence_topBlock_fl_backSub_gamma_bound_explicitCompactBudget_of_signed_alpha_leadingBlock_det_ne_zero_normSqBudget_signedStageUniformBudget_product_of_normSqBudget_activePrefix_activeBlockRecurrence_globalCompactBudget_completedColumns_globalProduct_offdiag_rows_of_horizonBudget}
\end{formalized}

\begin{corollary}[Completed-column preservation for the active/prefix QR route]
In the same signed stored-QR loop, assume the rounded stored-panel recurrence
for \(A_t\).  If \(j<t\), then the exact signed Householder reflector used at
stage \(t\) leaves the already completed column \(j\) unchanged:
\[
  H_t\,A_t(:,j)=A_t(:,j).
\]
Consequently the global compact-step source-control and solver certificates
above no longer need a separate completed-column hypothesis; it is derived
from the stored-loop recurrence itself.
\end{corollary}


\begin{formalized}
\lean{Algorithms/QR/HouseholderQR.lean}
\leanfit{storedQRSignedStage_completed_column_preservation}
\lean{Algorithms/LeastSquares/LSQRSolve.lean}
\leanfit{StoredQRSourceOffDiagonalControl.of_stored_trailing_sequence_leadingBlock_det_ne_zero_normSqBudget_signedStageUniformBudget_product_of_normSqBudget_activePrefix_activeBlockRecurrence_globalCompactBudget_completedColumns_offdiag_rows}
\leanfit{LSQRSolveBackwardError.of_stored_trailing_householder_sequence_topBlock_fl_backSub_gamma_bound_explicitCompactBudget_of_signed_alpha_leadingBlock_det_ne_zero_normSqBudget_signedStageUniformBudget_product_of_normSqBudget_activePrefix_activeBlockRecurrence_globalCompactBudget_completedColumns_offdiag_rows}
\end{formalized}

\begin{corollary}[Finite global compact-product budget for active/prefix QR]
In the same signed stored-QR loop, define for each pivot \(k<n\)
\[
  \Pi_k
  =
  2\,D_k
  \left(m\,\bigl(c_{\rm seq}\,\|A_k(:,k)\|_2\bigr)^2\right),
\]
where
\[
  D_k =
  \operatorname{diagDominantUpperInvBudgetExpr}(S_k,k),
  \qquad
  c_{\rm seq} =
  \operatorname{storedQRCompactSequenceRelativeBudget}.
\]
Let
\[
  \Pi_{\max}=\max_{0\le k<n}\Pi_k
\]
with the vacuous value \(0\) when \(n=0\).  If
\[
  \Pi_{\max}<1,
\]
then the per-pivot compact-product condition
\[
  \Pi_k<1,\qquad 0\le k<n,
\]
used by the active/prefix source-control and solver certificates follows.
Consequently the global compact-step and completed-column QR certificates no
longer need a separate compact-product hypothesis at every pivot; they consume
the single scalar smallness condition \(\Pi_{\max}<1\).

Conversely, the finite-maximum bookkeeping is now formalized in the other
direction as well: if every \(\Pi_k<1\), then \(\Pi_{\max}<1\).  Thus the
global scalar assumption is exactly a finite packaging of the per-pivot
compact-product smallness field; the remaining mathematical task is to prove
those smallness inequalities from a concrete pivoted/sorted QR invariant or to
keep them visible as domain assumptions.

There is one exact-arithmetic endpoint where the product-smallness side
condition is automatic.  If the stored compact relative budget vanishes,
\[
  c_{\rm seq}=0,
\]
then every \(\Pi_k=0\), and hence \(\Pi_{\max}<1\).  This closes the
compact-product dependency for the exact/zero-compact regime.  It is not a
floating-point smallness theorem: when \(c_{\rm seq}>0\), the positive-budget
smallness inequality must still be proved from conditioning, pivoting, or a
source/domain assumption.

A further scalar sufficient condition is also formalized.  If
\[
  D_k\le D_{\max},
  \qquad
  \|A_k(:,k)\|_2\le N_{\max}
\]
for every pivot \(k\), each local leading block is diagonally dominant, and
\[
  2D_{\max}\left(m\,(c_{\rm seq}N_{\max})^2\right)<1,
\]
then \(\Pi_{\max}<1\).  This records a single source-facing product budget;
it does not prove the numerical smallness of \(D_{\max}\), \(N_{\max}\), or
\(c_{\rm seq}\) from the computed QR loop.

The global bounds themselves are now packaged canonically.  Let
\[
  D_*=\max_{0\le k<n}D_k,
  \qquad
  N_*=\max_{0\le k<n}\|A_k(:,k)\|_2,
\]
again with the vacuous value \(0\) for an empty pivot set.  Under local
diagonal dominance, \(D_*\ge 0\), \(N_*\ge 0\), and every \(D_k\le D_*\) and
\(\|A_k(:,k)\|_2\le N_*\).  Therefore the single canonical condition
\[
  2D_*\left(m\,(c_{\rm seq}N_*)^2\right)<1
\]
implies \(\Pi_{\max}<1\).  The remaining product-smallness task is no longer a
family of global-bound hypotheses; it is the scalar inequality for these
canonical finite maxima.  This scalar condition is also composed into the
main diagonal-dominant QR source-control and solver surfaces: those theorems
now consume
\[
  2D_*\left(m\,(c_{\rm seq}N_*)^2\right)<1
\]
directly rather than asking the caller to separately provide
\(\Pi_{\max}<1\).

There is also a positive-budget threshold form.  Let \(c_{\max}\ge 0\) be an
explicit cap such that
\[
  c_{\rm seq}\le c_{\max}.
\]
If
\[
  2D_*\left(m\,(c_{\max}N_*)^2\right)<1,
\]
then again \(\Pi_{\max}<1\).  This is the finite-max monotonicity step needed
to replace the exact compact relative budget by a source-facing upper bound.
It does not prove \(c_{\rm seq}\le c_{\max}\), nor does it prove the scalar
smallness inequality from machine precision or conditioning data.

One useful specialization is now also formalized at the stored-QR layer.  Let
\(c_{\rm step}\) be a uniform one-step compact-panel relative budget:
\[
  \operatorname{storedQRCompactStepRelativeBudget}_k\le c_{\rm step},
  \qquad 0\le k<n.
\]
Since the sequence budget is the finite sum of these one-step budgets,
\[
  c_{\rm seq}\le n\,c_{\rm step}.
\]
Thus, if \(c_{\rm step}\ge 0\) and
\[
  2D_*\left(m\,((n\,c_{\rm step})N_*)^2\right)<1,
\]
then \(\Pi_{\max}<1\).  This closes the finite-sum reduction from a uniform
per-step cap to the positive-budget compact-product threshold.  It still does
not prove the one-step cap \(c_{\rm step}\), local diagonal dominance, or the
final scalar inequality from a concrete computed QR recurrence.

The one-step cap itself has also been reduced to vector-level compact
Householder caps.  For a stored QR stage \(k\), let
\[
  \begin{aligned}
  \rho_{k,j}^{\rm col}
  &=
  \operatorname{householderCompactRelativeBudget}
    (v_k,\beta_k,\widehat A_k(:,j)),\\
  \rho_k^{\rm rhs}
  &=
  \operatorname{householderCompactRelativeBudget}
    (v_k,\beta_k,\widehat b_k).
  \end{aligned}
\]
If
\[
  \rho_{k,j}^{\rm col}\le c_{\rm col}\quad\text{for all }j,
  \qquad
  \rho_k^{\rm rhs}\le c_{\rm rhs},
\]
then the stored-step panel budget satisfies
\[
  \operatorname{storedQRCompactStepRelativeBudget}_k
  \le n\,c_{\rm col}+c_{\rm rhs}.
\]
Consequently
\[
  c_{\rm seq}\le n\,(n\,c_{\rm col}+c_{\rm rhs}).
\]
Thus, if \(c_{\rm col},c_{\rm rhs}\ge0\) and
\[
  2D_*\left(m\,\bigl((n\,(n\,c_{\rm col}+c_{\rm rhs}))N_*\bigr)^2\right)<1,
\]
then \(\Pi_{\max}<1\).  This closes the finite-sum bridge from vector-level
compact caps to the product threshold.  It still does not prove the
vector-level caps, local diagonal dominance, or the scalar product-smallness
inequality from a concrete pivoted/sorted QR loop.

There is now a primitive norm-budget version of the same bridge.  Define
\[
  \begin{aligned}
  B_{k,j}^{\rm col}
  &=
  \operatorname{householderCompactNormBudget}
    (v_k,\beta_k,\widehat A_k(:,j)),\\
  B_k^{\rm rhs}
  &=
  \operatorname{householderCompactNormBudget}
    (v_k,\beta_k,\widehat b_k).
  \end{aligned}
\]
If \(c_{\rm col},c_{\rm rhs}\ge0\) and
\[
  B_{k,j}^{\rm col}
  \le c_{\rm col}\,\|\widehat A_k(:,j)\|_2
  \quad\text{for all }j,
  \qquad
  B_k^{\rm rhs}\le c_{\rm rhs}\,\|\widehat b_k\|_2,
\]
then the definition of relative compact budget gives
\[
  \rho_{k,j}^{\rm col}\le c_{\rm col},
  \qquad
  \rho_k^{\rm rhs}\le c_{\rm rhs}.
\]
Consequently the same sequence cap
\[
  c_{\rm seq}\le n\,(n\,c_{\rm col}+c_{\rm rhs})
\]
and the same finite-max product-smallness conclusion hold.  This is the
primitive-budget bridge: it removes the need to assume already-normalized
relative compact caps, but it still leaves the norm-budget inequalities and
the scalar product-smallness inequality visible.

There is also a componentwise primitive version.  Define the entrywise
compact-budget quantities
\[
  \begin{aligned}
  E_{k,j,i}^{\rm col}
  &=
  \operatorname{householderCompactComponentBudget}
    (v_k,\beta_k,\widehat A_k(:,j))_i,\\
  E_{k,i}^{\rm rhs}
  &=
  \operatorname{householderCompactComponentBudget}
    (v_k,\beta_k,\widehat b_k)_i.
  \end{aligned}
\]
If \(c_{\rm col},c_{\rm rhs}\ge0\) and
\[
  E_{k,j,i}^{\rm col}\le c_{\rm col}\,|\widehat A_k(i,j)|
  \quad\text{for all }i,j,
  \qquad
  E_{k,i}^{\rm rhs}\le c_{\rm rhs}\,|\widehat b_k(i)|
  \quad\text{for all }i,
\]
then local norm monotonicity and homogeneity give
\[
  B_{k,j}^{\rm col}
  \le c_{\rm col}\,\|\widehat A_k(:,j)\|_2,
  \qquad
  B_k^{\rm rhs}\le c_{\rm rhs}\,\|\widehat b_k\|_2.
\]
Thus the same sequence cap
\[
  c_{\rm seq}\le n\,(n\,c_{\rm col}+c_{\rm rhs})
\]
and the same finite-max product-smallness conclusion hold.  This closes the
bookkeeping step from entrywise compact budgets to the product threshold, but
it still does not prove the entrywise inequalities or the scalar smallness
inequality from a concrete floating-point QR loop.

The valid primitive floating-point route is normwise rather than entrywise
relative.  For a reflector \((v,\beta)\), let
\[
  V=\|v\|_2,\qquad
  \sigma_*=(1+\gamma_m)V,\qquad
  \tau_*=|\beta|\,\sigma_*(1+u),
\]
and let
\[
  e_* = u\,\tau_* + u\,|\beta|\,\sigma_* + |\beta|\,\gamma_m V,
  \qquad
  z_*=\tau_*(1+u).
\]
The formal coefficient is
\[
  c_{\rm HH}(v,\beta)
  =
  u + (e_*+u z_*)\,V .
\]
For every vector \(x\),
\[
  \operatorname{householderCompactNormBudget}(v,\beta,x)
  \le c_{\rm HH}(v,\beta)\,\|x\|_2.
\]
The key component estimate behind this norm statement is
\[
  E_i(v,\beta,x)
  \le (e_*+u z_*)\,\|x\|_2\,|v_i|+u\,|x_i|.
\]
Thus the nonlocal dot/scale/update error is charged to
\(\|x\|_2 |v_i|\), while only the final subtraction is charged to
\(|x_i|\).  This is why the theorem does not require the generally false
universal cap \(E_i\le c|x_i|\).

For a stored QR sequence define
\[
  c_k^{\rm HH}=c_{\rm HH}(v_k,\beta_k).
\]
The Lean library first proves the conditional version: if
\(c_k^{\rm HH}\le c_{\rm HH}^{\max}\) for all \(k\), then each compact
stored step has relative budget at most
\[
  n\,c_{\rm HH}^{\max}+c_{\rm HH}^{\max},
\]
and hence the global sequence budget is bounded by
\[
  c_{\rm seq}\le
  n\,(n\,c_{\rm HH}^{\max}+c_{\rm HH}^{\max}).
\]
Therefore, under local diagonal dominance, if
\[
  2D_*\left(m\,\bigl((n\,(n\,c_{\rm HH}^{\max}
      +c_{\rm HH}^{\max}))N_*\bigr)^2\right)<1,
\]
then \(\Pi_{\max}<1\).  This closes the concrete floating-point
norm-coefficient reduction for the compact-product route.

The finite-max version removes the arbitrary uniform-coefficient field.
Define
\[
  C_{\rm HH}^*
  =
  \begin{cases}
  \max_{k\in\operatorname{Fin}(n)} c_k^{\rm HH}, & n>0,\\
  0, & n=0.
  \end{cases}
\]
The formalization proves \(0\le C_{\rm HH}^*\) and
\(c_k^{\rm HH}\le C_{\rm HH}^*\) for every stored stage.  Hence, under local
diagonal dominance, the scalar condition
\[
  2D_*\left(m\,\bigl((n\,(n\,C_{\rm HH}^*+C_{\rm HH}^*))N_*\bigr)^2\right)<1
\]
implies \(\Pi_{\max}<1\).  This is still not a machine-precision-only
smallness theorem: diagonal dominance/off-diagonal control and the displayed
scalar inequality for the canonical coefficient maximum remain visible
numerical assumptions.

The coefficient maximum has a new per-pivot specialization.  Define
\[
  F_m(u,\gamma_m)
  =
  u(1+\gamma_m)(1+u)
  + u(1+\gamma_m)
  + \gamma_m
  + u(1+\gamma_m)(1+u)^2 .
\]
The compact coefficient has the exact form
\[
  c_{\rm HH}(v,\beta)
  =
  u + \bigl(|\beta|\,\|v\|_2^2\bigr)F_m(u,\gamma_m).
\]
Moreover \(F_m(u,\gamma_m)\ge 0\) under the repository's
\(\gamma\)-validity hypothesis.  Hence a uniform signed-stage estimate
\[
  |\beta_k|\,\|v_k\|_2^2\le B_{\max}
\]
implies
\[
  C_{\rm HH}^*
  \le
  u+B_{\max}F_m(u,\gamma_m).
\]
Consequently, the finite-max product theorem may be checked with
\[
  C_{\max}=u+B_{\max}F_m(u,\gamma_m)
\]
instead of an abstract coefficient maximum.  The resulting scalar condition is
\[
  2D_{\max}
  \left(
    m\left(
      n\bigl(n(u+B_{\max}F_m)+(u+B_{\max}F_m)\bigr)N_{\max}
    \right)^2
  \right)<1.
\]
This reduces the open coefficient-bound field to the familiar Householder
normalization quantity \(|\beta_k|\|v_k\|_2^2\).  It still does not prove the
\(B_{\max}\) estimate, local diagonal dominance, or the displayed scalar
smallness inequality from the concrete stored QR loop.

That \(B_{\max}\) field is now closed for the exact Householder normalization
used by the stored signed stages.  If
\[
  \beta=\frac{2}{v^Tv}
  \qquad\text{and}\qquad
  v^Tv\ne 0,
\]
then
\[
  |\beta|\,\|v\|_2^2 = 2.
\]
Consequently, for every stored signed stage whose source-shaped Householder
denominator is nonzero,
\[
  |\beta_k|\,\|v_k\|_2^2 = 2,
  \qquad
  c_k^{\rm HH}
  \le
  u+2F_m(u,\gamma_m),
\]
and therefore
\[
  C_{\rm HH}^*
  \le
  u+2F_m(u,\gamma_m).
\]
The finite-max product and off-diagonal-control invariant may thus be checked
with the concrete coefficient maximum
\[
  C_{\max}=u+2F_m(u,\gamma_m),
\]
provided the signed-stage denominators are nonzero.  The remaining fields are
the scalar smallness inequality with this concrete \(C_{\max}\), local diagonal
dominance or another off-diagonal-control invariant, inverse-factor and
pivot-column estimates, and determinant/conditioning data needed by the chosen
QR/preconditioner route.

The same active/prefix route now also has a structured nonbreakdown variant:
instead of assuming the norm-square pivot margin directly, one may assume
local leading-block nonsingularity together with a displayed \(\kappa_\infty\)
self-norm budget and a dual compact-update budget.  The previously formalized
leading-block inverse-budget theorem then gives the dimensioned norm-square
margin used by the signed-stage QR source-control proof.

The row-budget-control solver surface has the same finite-global-product
packaging.  A local \(\lean{LSQRSolveBackwardError}\) theorem consumes the
packaged displayed row-budget certificate and the scalar condition
\(\Pi_{\max}<1\); its \(\kappa_\infty\)/dual-budget sibling derives the
norm-square nonbreakdown margin from the local leading-block inverse-budget
route before applying the finite-global-product solver theorem.

There is now also a solver-facing active-pivot-policy version of this
\(\kappa_\infty\) route.  Instead of assuming the raw column-pivot maximality
inequality, it assumes the displayed pivot column is the finite active
max-pivot selector \(\lean{householderActiveMaxPivotColumn}\).  The raw
inequality is then recovered locally from
\(\lean{householderActiveMaxPivotColumn_pivot_max}\) before applying the
existing global-product solver certificate.  This closes the pivot-policy
handoff at the local least-squares layer only; the diagonal lower-bound,
local determinant/conditioning, dual compact-budget, and compact-product
smallness assumptions remain explicit.
\end{corollary}


\begin{corollary}[Source-facing scalar certificate]
With the hypotheses of the exact-normalization branch above, set
\[
  C_{\max}=u+2F_m(u,\gamma_m).
\]
If the displayed route constants satisfy
\[
  2D_{\max}\,m\,
    \bigl(n(n+1)C_{\max}N_{\max}\bigr)^2 < 1,
\]
then the canonical compact-product smallness condition holds, and hence the
same packaged off-diagonal-control invariant follows under the same local
diagonal-dominance, source-nonbreakdown, \(D_{\max}\), and \(N_{\max}\)
hypotheses.

This corollary is an algebraic normalization of the scalar side condition. It
does not prove the numerical inequality, diagonal dominance, inverse-factor or
pivot-column bounds, or determinant/conditioning assumptions.
\end{corollary}

\begin{corollary}[Capped source-facing scalar certificate]
Keep the hypotheses of the preceding corollary, but suppose that the
unit-roundoff and coefficient factor are bounded by explicit caps
\[
  u\le U_{\max},\qquad F_m(u,\gamma_m)\le F_{\max},
\]
with \(U_{\max},F_{\max}\ge 0\).  If
\[
  2D_{\max}\,m\,
    \bigl(n(n+1)(U_{\max}+2F_{\max})N_{\max}\bigr)^2 < 1,
\]
then the same canonical compact-product smallness condition holds, and hence
the same packaged off-diagonal-control invariant follows under the same local
diagonal-dominance, source-nonbreakdown, \(D_{\max}\), and \(N_{\max}\)
hypotheses.

This is the next scalar-smallness reduction in the LS.2g red-bottleneck route:
it replaces the formula involving the concrete floating-point constants by a
monotone cap certificate.  It does not prove the caps, the \(D_{\max}\) or
\(N_{\max}\) pointwise estimates, diagonal dominance, source nonbreakdown, or
the final QR/preconditioner theorem.
\end{corollary}

\begin{lemma}[Householder coefficient-factor cap]
Let
\[
  F_m(u,\gamma_m)
  =
  u(1+\gamma_m)(1+u)
  +u(1+\gamma_m)
  +\gamma_m
  +u(1+\gamma_m)(1+u)^2
\]
be the compact Householder coefficient factor used above.  Suppose
\(\gamma_m\) is valid in the local floating-point model, and suppose
\[
  0\le U_{\max},\qquad 0\le G_{\max},\qquad
  u\le U_{\max},\qquad \gamma_m\le G_{\max}.
\]
Then
\[
  F_m(u,\gamma_m)\le
  U_{\max}(1+G_{\max})(1+U_{\max})
  +U_{\max}(1+G_{\max})
  +G_{\max}
  +U_{\max}(1+G_{\max})(1+U_{\max})^2.
\]

This lemma supplies one concrete source for the \(F_{\max}\) cap in the
preceding corollary.  It does not prove the primitive caps on \(u\) or
\(\gamma_m\), nor the final scalar cap inequality.
\end{lemma}

\begin{lemma}[Gamma cap from a unit-roundoff cap]
Suppose the unit roundoff satisfies \(u\le U_{\max}\), and suppose
\[
  mU_{\max}<1.
\]
Then
\[
  mu<1,
  \qquad
  \gamma_m
  =
  \frac{mu}{1-mu}
  \le
  \frac{mU_{\max}}{1-mU_{\max}}.
\]
Consequently, if a displayed constant \(G_{\max}\) satisfies
\[
  \frac{mU_{\max}}{1-mU_{\max}}\le G_{\max},
\]
then \(\gamma_m\le G_{\max}\).

The first inequality is the Lean validity guard
\(\lean{gammaValid fp m}\).  Together with the preceding lemma, this supplies
the gamma side of the coefficient-factor cap
\(F_m(u,\gamma_m)\le F_{\max}\).  It still leaves the actual unit-roundoff cap
and the final scalar smallness inequality as explicit route obligations.
\end{lemma}

\begin{corollary}[Coefficient-factor cap from a unit-roundoff cap]
Assume \(0\le U_{\max}\), \(u\le U_{\max}\), and
\[
  mU_{\max}<1,\qquad
  \frac{mU_{\max}}{1-mU_{\max}}\le G_{\max}.
\]
Then the same displayed polynomial bound for \(F_m(u,\gamma_m)\) in the
Householder coefficient-factor cap holds with \(U_{\max}\) and \(G_{\max}\).

This corollary is the direct composition of the two preceding lemmas.  It
removes the independent hypothesis \(\gamma_m\le G_{\max}\) from the
coefficient-factor cap route, but it still leaves the primitive unit-roundoff
cap and the downstream scalar product-smallness inequality as explicit
obligations.
\end{corollary}

\begin{corollary}[Source-denominator cap route from unit-roundoff and gamma caps]
In the source-denominator scalar-smallness route above, replace the abstract
factor cap \(F_{\max}\) by the displayed polynomial bound from Corollary~322.
Then the same scalar-smallness, compact-product, and packaged
off-diagonal-control conclusions hold from
\[
  u\le U_{\max},\qquad
  mU_{\max}<1,\qquad
  \frac{mU_{\max}}{1-mU_{\max}}\le G_{\max},
\]
and the displayed scalar inequality using
\[
  F_{\max}
  =
  U_{\max}(1+G_{\max})(1+U_{\max})
  +U_{\max}(1+G_{\max})
  +G_{\max}
  +U_{\max}(1+G_{\max})(1+U_{\max})^2.
\]

This is a theorem-surface reduction.  It removes the separate
\(F_m(u,\gamma_m)\le F_{\max}\) proof field from the least-squares
source-denominator route; it does not prove the primitive unit-roundoff cap,
the rational \(G_{\max}\) domination, local diagonal dominance, pointwise
inverse/pivot-column bounds, source nonbreakdown, or the final equation~(8)
QR/preconditioner theorem.
\end{corollary}

\begin{corollary}[Rational gamma cap specialization]
In Corollary~323, choose
\[
  G_{\max}=\frac{mU_{\max}}{1-mU_{\max}}.
\]
Then the same source-denominator scalar-smallness, compact-product, and
packaged off-diagonal-control conclusions hold under the unit-roundoff cap
\(u\le U_{\max}\), the validity condition \(mU_{\max}<1\), the pointwise
\(D_{\max}\) and \(N_{\max}\) route bounds, source nonbreakdown, and the
displayed scalar-smallness inequality with this value of \(G_{\max}\)
substituted into \(F_{\max}\).  Thus the rational \(G_{\max}\)-domination
hypothesis is not a separate route obligation when the rational cap itself is
used.

This specialization does not prove the primitive unit-roundoff cap, the
scalar-smallness inequality, local diagonal dominance, pointwise inverse or
pivot-column estimates, source nonbreakdown, or the final equation~(8)
QR/preconditioner theorem.
\end{corollary}

\begin{corollary}[Canonical finite-max rational gamma cap route]
In Corollary~324, choose
\[
  D_{\max}=D_\ast(A_{\mathrm{fl}}),
  \qquad
  N_{\max}=N_\ast(A_{\mathrm{fl}}),
\]
where \(D_\ast\) is the repository finite maximum of the diagonal-dominant
inverse-factor budgets and \(N_\ast\) is the repository finite maximum of the
displayed pivot-column norms.  Then the same source-denominator
scalar-smallness, compact-product, and packaged off-diagonal-control
conclusions hold with the scalar-smallness inequality written directly in
terms of \(D_\ast\), \(N_\ast\), and the rational
\(G_{\max}=mU_{\max}/(1-mU_{\max})\).

Thus the rational-gamma route no longer needs separate pointwise certificates
showing every local inverse factor is bounded by \(D_{\max}\) and every pivot
column norm is bounded by \(N_{\max}\); those proofs are discharged by the
finite-maximum lemmas.  The result still assumes local diagonal dominance,
source nonbreakdown, \(u\le U_{\max}\), \(mU_{\max}<1\), and the displayed
scalar-smallness inequality, and it does not prove the final equation~(8)
QR/preconditioner theorem.
\end{corollary}

\begin{corollary}[Solver and probability handoff for the canonical rational-gamma route]
Assume the hypotheses of Corollary~325 samplewise for the stored QR sequence:
local diagonal dominance, source-shaped Householder denominator
nonbreakdown, \(u\le U_{\max}\), \(mU_{\max}<1\), and the canonical
finite-max scalar inequality with
\[
  G_{\max}=\frac{mU_{\max}}{1-mU_{\max}},
  \qquad
  D_{\max}=D_\ast(A_{\mathrm{fl}}),
  \qquad
  N_{\max}=N_\ast(A_{\mathrm{fl}}).
\]
Then the stored-QR solve certificate used by the equation~(8) least-squares
route holds.  Composed with the already formalized Bennett sample-budget
least-squares theorem, this gives the corresponding high-probability rounded
sampled-row objective bound for Algorithm~2.

This corollary is a solver/probability assembly theorem.  It does not prove
local diagonal dominance, source denominator nonbreakdown, the primitive
unit-roundoff cap, or the canonical scalar finite-max smallness inequality.
Those remain the visible QR-domain obligations for this branch.
\end{corollary}

\begin{lemma}[No universal unit-roundoff cap for the abstract model]
For every real number \(U_{\max}\),
\[
  \neg\bigl(\forall \mathsf{fp}:\mathrm{FPModel},\;
    u(\mathsf{fp})\le U_{\max}\bigr).
\]
Indeed, exact arithmetic can be packaged as an \(\mathrm{FPModel}\) with any
chosen nonnegative declared unit roundoff.  Therefore the hypothesis
\(u\le U_{\max}\) in the preceding cap route is not a hidden consequence of
the abstract floating-point model.  It must be supplied as a domain assumption
or derived from a more concrete machine model.
\end{lemma}

\begin{formalized}
\lean{Model.lean}
\leanfit{FPModel.exactWithUnitRoundoff}
\leanfit{FPModel.not_forall_u_le_cap}
\lean{Analysis/Rounding.lean}
\leanfit{gamma_le_of_u_le_cap}
\leanfit{gamma_le_Gcap_of_u_le_cap}
\lean{Algorithms/QR/HouseholderSpec.lean}
\leanfit{abs_householderBeta_mul_vecNorm2_sq_eq_two}
\leanfit{abs_householderBeta_mul_vecNorm2_sq_le_two}
\lean{Algorithms/QR/HouseholderApply.lean}
\leanfit{householderAbsDotBudget_le_vecNorm2_mul}
\leanfit{householderCompactUpdateCoeff}
\leanfit{householderCompactNormBudgetCoeff}
\leanfit{householderCompactNormBudgetCoeffFactor}
\leanfit{householderCompactNormBudgetCoeff_eq_u_add_abs_beta_norm_sq_mul_factor}
\leanfit{householderCompactNormBudgetCoeffFactor_nonneg}
\leanfit{householderCompactNormBudgetCoeffFactor_le_of_u_gamma_le}
\leanfit{householderCompactNormBudgetCoeffFactor_le_of_u_cap_gamma_cap}
\leanfit{householderCompactNormBudgetCoeff_le_of_abs_beta_norm_sq_le}
\leanfit{householderCompactUpdateCoeff_nonneg}
\leanfit{householderCompactNormBudgetCoeff_nonneg}
\leanfit{householderCompactComponentBudget_le_updateCoeff_mul_norm}
\leanfit{householderCompactNormBudget_le_normBudgetCoeff_mul}
\leanfit{householderCompactRelativeBudget_le_normBudgetCoeff}
\leanfit{householderCompactPanelRelativeBudget_le_mul_add_normBudgetCoeff}
\leanfit{householderCompactNormBudget_le_mul_of_componentBudget_le_mul_abs}
\leanfit{householderCompactRelativeBudget_le_of_normBudget_le_mul}
\leanfit{householderCompactRelativeBudget_le_of_componentBudget_le_mul_abs}
\leanfit{householderCompactPanelRelativeBudget_le_mul_add_of_column_rhs_le}
\leanfit{householderCompactPanelRelativeBudget_le_mul_add_of_normBudget_le_mul}
\leanfit{householderCompactPanelRelativeBudget_le_mul_add_of_componentBudget_le_mul_abs}
\lean{Algorithms/QR/HouseholderQR.lean}
\leanfit{storedQRCompactStepNormBudgetCoeff}
\leanfit{storedQRCompactStepNormBudgetCoeff_nonneg}
\leanfit{storedQRCompactStepNormBudgetCoeff_le_of_abs_beta_norm_sq_le}
\leanfit{storedQRCompactStepNormBudgetCoeff_le_of_forall_abs_beta_norm_sq_le}
\leanfit{storedQRSignedStage_abs_beta_norm_sq_eq_two_of_den_ne_zero}
\leanfit{storedQRSignedStage_abs_beta_norm_sq_le_two_of_den_ne_zero}
\leanfit{storedQRCompactStepNormBudgetCoeff_le_of_den_ne_zero}
\leanfit{storedQRCompactStepNormBudgetCoeff_le_of_forall_den_ne_zero}
\leanfit{storedQRCompactStepNormBudgetCoeff_le_of_forall_source_den_ne_zero}
\leanfit{storedQRCompactStepRelativeBudget_le_mul_add_of_column_rhs_le}
\leanfit{storedQRCompactStepRelativeBudget_le_mul_add_of_normBudget_le_mul}
\leanfit{storedQRCompactStepRelativeBudget_le_mul_add_of_componentBudget_le_mul_abs}
\leanfit{storedQRCompactStepRelativeBudget_le_mul_add_normBudgetCoeff}
\leanfit{storedQRCompactSequenceRelativeBudget_le_mul_of_step_le}
\lean{Algorithms/LeastSquares/LSQRSolve.lean}
\leanfit{storedQRCompactSequenceProductExpr}
\leanfit{storedQRCompactSequenceProductBudget}
\leanfit{storedQRCompactSequenceProductExpr_le_budget}
\leanfit{storedQRCompactSequenceProductBudget_lt_one_of_forall_expr_lt}
\leanfit{storedQRCompactSequenceProductBudget_lt_one_of_forall_pivot_product}
\leanfit{storedQRCompactSequenceProductBudget_lt_one_of_relativeBudget_eq_zero}
\leanfit{StoredQRSourceOffDiagonalControl.of_stored_trailing_sequence_leadingBlock_det_ne_zero_normSqBudget_rowBudgetControl_globalProduct}
\leanfit{StoredQRSourceOffDiagonalControl.of_stored_trailing_sequence_leadingBlock_det_ne_zero_normSqBudget_rowMaxBudgetFactor_globalProduct}
\leanfit{LSQRSolveBackwardError.of_stored_trailing_householder_sequence_topBlock_fl_backSub_gamma_bound_explicitCompactBudget_of_signed_alpha_leadingBlock_det_ne_zero_normSqBudget_rowBudgetControl_globalProduct}
\leanfit{LSQRSolveBackwardError.of_stored_trailing_householder_sequence_topBlock_fl_backSub_gamma_bound_explicitCompactBudget_of_signed_alpha_leadingBlock_det_ne_zero_normSqBudget_rowMaxBudgetFactor_globalProduct}
\leanfit{LSQRSolveBackwardError.of_stored_trailing_householder_sequence_topBlock_fl_backSub_gamma_bound_explicitCompactBudget_of_signed_alpha_leadingBlock_det_ne_zero_kappaInf_selfNorm_dualBudget_rowBudgetControl_globalProduct}
\leanfit{storedQRCompactSequenceProductBudget_lt_one_of_diagDominant_factor_norm_bounds}
\leanfit{storedQRDiagDominantInvFactorBudget}
\leanfit{storedQRDiagDominantInvFactor_le_budget}
\leanfit{storedQRDiagDominantInvFactorBudget_nonneg}
\leanfit{storedQRPivotColumnNormBudget}
\leanfit{storedQRPivotColumnNorm_le_budget}
\leanfit{storedQRPivotColumnNormBudget_nonneg}
\leanfit{storedQRCompactSequenceProductBudget_lt_one_of_diagDominant_finite_max_smallness}
\leanfit{storedQRCompactSequenceProductBudget_lt_one_of_diagDominant_finite_max_relativeBudget_le}
\leanfit{storedQRCompactSequenceProductBudget_lt_one_of_diagDominant_finite_max_stepBudget_le}
\leanfit{storedQRCompactSequenceProductBudget_lt_one_of_diagDominant_finite_max_columnRhsBudget_le}
\leanfit{storedQRCompactSequenceProductBudget_lt_one_of_diagDominant_finite_max_columnRhsNormBudget_le}
\leanfit{storedQRCompactSequenceProductBudget_lt_one_of_diagDominant_finite_max_columnRhsComponentBudget_le}
\leanfit{storedQRCompactSequenceProductBudget_lt_one_of_diagDominant_finite_max_normBudgetCoeff_le}
\leanfit{storedQRCompactStepNormBudgetCoeffBudget}
\leanfit{storedQRCompactStepNormBudgetCoeff_le_budget}
\leanfit{storedQRCompactStepNormBudgetCoeffBudget_nonneg}
\leanfit{storedQRDiagDominantInvFactorBudget_le_of_forall_le}
\leanfit{storedQRPivotColumnNormBudget_le_of_forall_le}
\leanfit{storedQRCompactStepNormBudgetCoeffBudget_le_of_forall_le}
\leanfit{storedQRCompactStepNormBudgetCoeffBudget_le_of_forall_abs_beta_norm_sq_le}
\leanfit{storedQRCompactStepNormBudgetCoeffBudget_le_of_forall_source_den_ne_zero}
\leanfit{storedQRCompactSequenceProductBudget_lt_one_of_diagDominant_finite_max_normBudgetCoeffBudget}
\leanfit{storedQRCompactNormBudgetCoeffSmallness_of_bounds}
\leanfit{storedQRCompactSequenceProductBudget_lt_one_of_diagDominant_finite_max_normBudgetCoeffBudget_bounds}
\leanfit{storedQRCompactNormBudgetCoeffSmallness_of_pointwise_bounds}
\leanfit{storedQRCompactSequenceProductBudget_lt_one_of_diagDominant_finite_max_normBudgetCoeffBudget_pointwise_bounds}
\leanfit{storedQRCompactNormBudgetCoeffSmallness_of_pointwise_absBetaNormSq_bounds}
\leanfit{storedQRCompactNormBudgetCoeffSmallness_of_pointwise_source_den_ne_zero_bounds}
\leanfit{storedQRCompactNormBudgetCoeffSmallness_of_pointwise_source_den_ne_zero_simple_bounds}
\leanfit{storedQRCompactNormBudgetCoeffSmallness_of_pointwise_source_den_ne_zero_cap_bounds}
\leanfit{storedQRCompactNormBudgetCoeffSmallness_of_pointwise_source_den_ne_zero_uGammaCapBounds}
\leanfit{storedQRCompactNormBudgetCoeffSmallness_of_pointwise_source_den_ne_zero_uRationalGammaCapBounds}
\leanfit{storedQRCompactNormBudgetCoeffSmallness_of_source_den_ne_zero_uRationalGammaCanonicalBounds}
\leanfit{storedQRCompactSequenceProductBudget_lt_one_of_diagDominant_finite_max_absBetaNormSqPointwiseBounds}
\leanfit{storedQRCompactSequenceProductBudget_lt_one_of_diagDominant_finite_max_sourceDenPointwiseBounds}
\leanfit{storedQRCompactSequenceProductBudget_lt_one_of_diagDominant_finite_max_sourceDenSimpleBounds}
\leanfit{storedQRCompactSequenceProductBudget_lt_one_of_diagDominant_finite_max_sourceDenCapBounds}
\leanfit{storedQRCompactSequenceProductBudget_lt_one_of_diagDominant_finite_max_sourceDenUGammaCapBounds}
\leanfit{storedQRCompactSequenceProductBudget_lt_one_of_diagDominant_finite_max_sourceDenURationalGammaCapBounds}
\leanfit{storedQRCompactSequenceProductBudget_lt_one_of_diagDominant_finite_max_sourceDenURationalGammaCanonicalBounds}
\leanfit{StoredQROffDiagonalControlInvariant.of_diagDominant_finiteMaxSmallness}
\leanfit{StoredQROffDiagonalControlInvariant.of_diagDominant_finiteMaxNormBudgetCoeffSmallness}
\leanfit{StoredQROffDiagonalControlInvariant.of_diagDominant_finiteMaxNormBudgetCoeffBoundedSmallness}
\leanfit{StoredQROffDiagonalControlInvariant.of_diagDominant_finiteMaxNormBudgetCoeffPointwiseBounds}
\leanfit{StoredQROffDiagonalControlInvariant.of_diagDominant_finiteMaxAbsBetaNormSqPointwiseBounds}
\leanfit{StoredQROffDiagonalControlInvariant.of_diagDominant_finiteMaxSourceDenPointwiseBounds}
\leanfit{StoredQROffDiagonalControlInvariant.of_diagDominant_finiteMaxSourceDenSimpleBounds}
\leanfit{StoredQROffDiagonalControlInvariant.of_diagDominant_finiteMaxSourceDenCapBounds}
\leanfit{StoredQROffDiagonalControlInvariant.of_diagDominant_finiteMaxSourceDenUGammaCapBounds}
\leanfit{StoredQROffDiagonalControlInvariant.of_diagDominant_finiteMaxSourceDenURationalGammaCapBounds}
\leanfit{StoredQROffDiagonalControlInvariant.of_diagDominant_finiteMaxSourceDenURationalGammaCanonicalBounds}
\leanfit{LSQRSolveBackwardError.of_stored_trailing_householder_sequence_topBlock_fl_backSub_gamma_bound_explicitCompactBudget_of_signed_alpha_diagDominant_finiteMaxSourceDenURationalGammaCanonicalBounds}
\lean{Algorithms/RandNLA/LeastSquaresSketch.lean}
\leanfit{leverageTraceProbability_eventProb_fl_rowSampleLSObjective_le_one_add_eta_of_augmentedSpan_sample_budget_of_objective_error_and_stored_qr_diagDominant_finiteMaxSourceDenURationalGammaCanonicalBounds_solver}
\lean{Algorithms/LeastSquares/LSQRSolve.lean}
\leanfit{LSQRSolveBackwardError.of_stored_trailing_householder_sequence_topBlock_fl_backSub_gamma_bound_explicitCompactBudget_of_signed_alpha_leadingBlock_det_ne_zero_diagDominant_finiteMaxNormBudgetCoeffPointwiseBounds_concreteDual}
\leanfit{LSQRSolveBackwardError.of_stored_trailing_householder_sequence_topBlock_fl_backSub_gamma_bound_explicitCompactBudget_of_signed_alpha_diagDominant_finiteMaxNormBudgetCoeffPointwiseBounds_concreteDual}
\leanfit{StoredQRSourceOffDiagonalControl.of_stored_trailing_sequence_leadingBlock_det_ne_zero_normSqBudget_diagDominant_finiteMaxSmallness}
\leanfit{StoredQRSourceOffDiagonalControl.of_stored_trailing_sequence_leadingBlock_det_ne_zero_kappaInf_selfNorm_dualBudget_diagDominant_finiteMaxSmallness}
\leanfit{LSQRSolveBackwardError.of_stored_trailing_householder_sequence_topBlock_fl_backSub_gamma_bound_explicitCompactBudget_of_signed_alpha_leadingBlock_det_ne_zero_kappaInf_selfNorm_dualBudget_diagDominant_finiteMaxSmallness}
\leanfit{storedQRSignedStage_normSqBudget_of_leadingBlock_det_ne_zero_kappaInf_selfNorm_dualBudget}
\leanfit{StoredQRSourceOffDiagonalControl.of_stored_trailing_sequence_leadingBlock_det_ne_zero_normSqBudget_signedStageUniformBudget_product_of_normSqBudget_activePrefix_activeBlockRecurrence_globalCompactBudget_completedColumns_globalProduct_offdiag_rows}
\leanfit{LSQRSolveBackwardError.of_stored_trailing_householder_sequence_topBlock_fl_backSub_gamma_bound_explicitCompactBudget_of_signed_alpha_leadingBlock_det_ne_zero_normSqBudget_signedStageUniformBudget_product_of_normSqBudget_activePrefix_activeBlockRecurrence_globalCompactBudget_completedColumns_globalProduct_offdiag_rows}
\leanfit{StoredQRSourceOffDiagonalControl.of_stored_trailing_sequence_leadingBlock_det_ne_zero_kappaInf_selfNorm_dualBudget_signedStageUniformBudget_globalCompactBudget_completedColumns_globalProduct_offdiag_rows}
\leanfit{LSQRSolveBackwardError.of_stored_trailing_householder_sequence_topBlock_fl_backSub_gamma_bound_explicitCompactBudget_of_signed_alpha_leadingBlock_det_ne_zero_kappaInf_selfNorm_dualBudget_signedStageUniformBudget_globalCompactBudget_completedColumns_globalProduct_offdiag_rows}
\leanfit{LSQRSolveBackwardError.of_stored_trailing_householder_sequence_topBlock_fl_backSub_gamma_bound_explicitCompactBudget_of_signed_alpha_leadingBlock_det_ne_zero_kappaInf_selfNorm_dualBudget_signedStageUniformBudget_globalCompactBudget_activeMaxPivot_completedColumns_globalProduct_offdiag_rows}
\lean{Algorithms/RandNLA/LeastSquaresSketch.lean}
\leanfit{leverageTraceProbability_eventProb_fl_rowSampleLSObjective_le_one_add_eta_of_augmentedSpan_sample_budget_of_objective_error_and_stored_qr_offDiagonalControl_solver_of_diagDominant_finiteMaxNormBudgetCoeffSmallness}
\lean{Algorithms/LeastSquares/LSQRSolve.lean}
\leanfit{not_forall_diagDominant_sourceDen_uCap_implies_uRationalGammaCanonicalSmallness}
\leanfit{not_forall_diagDominant_sourceDen_actualU_implies_uRationalGammaCanonicalSmallness}
\end{formalized}

\begin{corollary}[Canonical scalar smallness is not automatic]
The canonical rational-gamma route still has a genuine numerical smallness
assumption.  There is a one-dimensional witness with exact arithmetic packaged
as \(\lean{FPModel.exactWithUnitRoundoff}(1/2)\), \(A=[1]\),
\(\alpha=0\), and \(U_{\max}=1/2\), for which the leading block is
diagonally dominant, source nonbreakdown holds, \(u\le U_{\max}\), and
\(mU_{\max}<1\), but the displayed canonical scalar condition is false.
Consequently, diagonal dominance, source nonbreakdown, the unit-roundoff cap,
and \(mU_{\max}<1\) do not by themselves discharge the scalar smallness
obligation.

The same obstruction remains if the displayed cap \(U_{\max}\) is removed and
the rational-gamma expression is written directly with the actual model
roundoff \(u\).  Thus this is not just an artifact of introducing an explicit
cap parameter: the canonical finite-max scalar inequality needs an independent
scale, conditioning, or machine-model smallness argument.
\end{corollary}


\begin{formalized}
\lean{Algorithms/LeastSquares/LSQRSolve.lean}
\leanfit{not_forall_diagDominant_sourceDen_uCap_implies_uRationalGammaCanonicalSmallness}
\leanfit{not_forall_diagDominant_sourceDen_actualU_implies_uRationalGammaCanonicalSmallness}
\end{formalized}

\begin{corollary}[Source nonbreakdown from signed trailing norms]
In the canonical rational-gamma QR route, the raw Householder denominator
condition \(v_k^\top v_k\ne0\) may be replaced by the signed-alpha source
facts
\[
  \alpha_k^2=\|A_k(k:m,k)\|_2^2,\qquad
  \|A_k(k:m,k)\|_2^2>0,\qquad
  \alpha_k A_k(k,k)\le0 .
\]
Under these hypotheses, the canonical off-diagonal-control invariant, the
stored-QR backward-error certificate, and the high-probability rounded
sampled-row objective theorem are the same as in the previous canonical route,
but no longer expose raw source denominator nonbreakdown.
\end{corollary}


\begin{formalized}
\lean{Algorithms/LeastSquares/LSQRSolve.lean}
\leanfit{storedQRSourceDenominator_ne_zero_of_trailingNorm2Sq_pos_mul_nonpos}
\leanfit{StoredQROffDiagonalControlInvariant.of_diagDominant_trailingNormPos_finiteMaxSourceDenURationalGammaCanonicalBounds}
\leanfit{LSQRSolveBackwardError.of_stored_trailing_householder_sequence_topBlock_fl_backSub_gamma_bound_explicitCompactBudget_of_signed_alpha_diagDominant_trailingNormPos_finiteMaxSourceDenURationalGammaCanonicalBounds}
\lean{Algorithms/RandNLA/LeastSquaresSketch.lean}
\leanfit{leverageTraceProbability_eventProb_fl_rowSampleLSObjective_le_one_add_eta_of_augmentedSpan_sample_budget_of_objective_error_and_stored_qr_diagDominant_trailingNormPos_finiteMaxSourceDenURationalGammaCanonicalBounds_solver}
\end{formalized}

\begin{corollary}[Source nonbreakdown from leading-block determinants]
In the previous canonical rational-gamma route, the positive trailing-norm
hypothesis
\[
  \|A_k(k:m,k)\|_2^2>0
\]
may itself be derived from local determinant data.  It is enough that, at each
stage \(k\), the previous transposed leading block and the current leading
block have nonzero determinant, and the already-completed previous columns are
zero below row \(k\).  Under these determinant and lower-zero hypotheses, the
same canonical off-diagonal-control invariant, stored-QR backward-error
certificate, and high-probability rounded sampled-row objective theorem hold.
\end{corollary}


\begin{formalized}
\lean{Algorithms/LeastSquares/LSQRSolve.lean}
\leanfit{StoredQROffDiagonalControlInvariant.of_diagDominant_leadingBlock_det_ne_zero_finiteMaxSourceDenURationalGammaCanonicalBounds}
\leanfit{StoredQROffDiagonalControlInvariant.of_diagDominant_signedAlphaDef_leadingBlock_det_ne_zero_finiteMaxSourceDenURationalGammaCanonicalBounds}
\leanfit{LSQRSolveBackwardError.of_stored_trailing_householder_sequence_topBlock_fl_backSub_gamma_bound_explicitCompactBudget_of_signed_alpha_diagDominant_leadingBlock_det_ne_zero_finiteMaxSourceDenURationalGammaCanonicalBounds}
\lean{Algorithms/RandNLA/LeastSquaresSketch.lean}
\leanfit{leverageTraceProbability_eventProb_fl_rowSampleLSObjective_le_one_add_eta_of_augmentedSpan_sample_budget_of_objective_error_and_stored_qr_diagDominant_leadingBlock_det_ne_zero_finiteMaxSourceDenURationalGammaCanonicalBounds_solver}
\end{formalized}

\begin{corollary}[Current determinant from diagonal dominance]
In the determinant-facing canonical rational-gamma route above, the current
leading-block determinant hypothesis may be omitted when the same current
leading block is already assumed to satisfy \(\operatorname{IsDiagDominantUpper}\).
The remaining determinant condition is only the previous transposed
leading-block determinant; the current determinant follows from the local
upper-triangular/nonzero-diagonal content of diagonal dominance.
\end{corollary}


\begin{formalized}
\lean{Algorithms/LeastSquares/LSQRSolve.lean}
\leanfit{StoredQROffDiagonalControlInvariant.of_diagDominant_signedAlphaDef_previousLeadingBlock_det_ne_zero_finiteMaxSourceDenURationalGammaCanonicalBounds}
\leanfit{LSQRSolveBackwardError.of_stored_trailing_householder_sequence_topBlock_fl_backSub_gamma_bound_explicitCompactBudget_of_signed_alpha_diagDominant_previousLeadingBlock_det_ne_zero_finiteMaxSourceDenURationalGammaCanonicalBounds}
\lean{Algorithms/RandNLA/LeastSquaresSketch.lean}
\leanfit{leverageTraceProbability_eventProb_fl_rowSampleLSObjective_le_one_add_eta_of_augmentedSpan_sample_budget_of_objective_error_and_stored_qr_diagDominant_previousLeadingBlock_det_ne_zero_finiteMaxSourceDenURationalGammaCanonicalBounds_solver}
\end{formalized}

\begin{corollary}[Previous determinant from diagonal dominance]
In the same determinant-facing canonical rational-gamma route, the previous
transposed leading-block determinant hypothesis may also be omitted.  The
previous \(k\times k\) block is the transpose orientation of the top-left part
of the current \((k+1)\times(k+1)\) leading block, so the upper-triangular and
nonzero-diagonal components of
\(\operatorname{IsDiagDominantUpper}\) imply its nonsingularity.  The remaining
local shape assumption is the previous-column lower-zero field used by the
trailing-norm nonbreakdown bridge.
\end{corollary}


\begin{formalized}
\lean{Algorithms/LeastSquares/LSQRSolve.lean}
\leanfit{qrPreviousLeadingBlockTranspose_det_ne_zero_of_diagDominant_leadingBlock}
\leanfit{StoredQROffDiagonalControlInvariant.of_diagDominant_signedAlphaDef_lowerPrev_finiteMaxSourceDenURationalGammaCanonicalBounds}
\leanfit{LSQRSolveBackwardError.of_stored_trailing_householder_sequence_topBlock_fl_backSub_gamma_bound_explicitCompactBudget_of_signed_alpha_diagDominant_lowerPrev_finiteMaxSourceDenURationalGammaCanonicalBounds}
\lean{Algorithms/RandNLA/LeastSquaresSketch.lean}
\leanfit{leverageTraceProbability_eventProb_fl_rowSampleLSObjective_le_one_add_eta_of_augmentedSpan_sample_budget_of_objective_error_and_stored_qr_diagDominant_lowerPrev_finiteMaxSourceDenURationalGammaCanonicalBounds_solver}
\end{formalized}

\begin{corollary}[Previous lower-zero shape from the stored QR recurrence]
In the same canonical signed-alpha rational-gamma route, assume the stored
Householder panel recurrence
\[
  A_{k+1}=\operatorname{fl\_householderStoredPanelStep}(A_k,v_k,\beta_k),
  \qquad k=0,\ldots,n-1 .
\]
Then the previous-column lower-zero field used above is not a separate
mathematical hypothesis:
\[
  i\ge k,\ j<k
  \quad\Longrightarrow\quad
  (A_k)_{i,j}=0 .
\]
Consequently the determinant-facing solver and probability wrappers may be
stated without the explicit lower-zero assumption.  The remaining visible
QR-domain assumptions on this route are local diagonal dominance/off-diagonal
control, the displayed scalar smallness inequality, primitive unit-roundoff
caps, conditioning fields, and the final concrete rectangular QR/preconditioner
theorem.
\end{corollary}


\begin{formalized}
\lean{Algorithms/LeastSquares/LSQRSolve.lean}
\leanfit{storedQRPreviousColumnLowerZero_of_stored_trailing_householder_sequence}
\leanfit{LSQRSolveBackwardError.of_stored_trailing_householder_sequence_topBlock_fl_backSub_gamma_bound_explicitCompactBudget_of_signed_alpha_diagDominant_storedLower_finiteMaxSourceDenURationalGammaCanonicalBounds}
\lean{Algorithms/RandNLA/LeastSquaresSketch.lean}
\leanfit{leverageTraceProbability_eventProb_fl_rowSampleLSObjective_le_one_add_eta_of_augmentedSpan_sample_budget_of_objective_error_and_stored_qr_diagDominant_storedLower_finiteMaxSourceDenURationalGammaCanonicalBounds_solver}
\end{formalized}

\begin{corollary}[Unit-roundoff cap nonnegativity is automatic]
In the stored-lower canonical rational-gamma route, the separate assumption
\[
  0\le U_{\rm cap}
\]
is not an independent domain hypothesis.  The FP model already contains
\(0\le u\), and therefore any primitive cap
\[
  u\le U_{\rm cap}
\]
implies \(0\le U_{\rm cap}\).  Consequently the solver and probability wrappers
for this stored-lower route may be stated without an explicit nonnegativity
field for \(U_{\rm cap}\).
\end{corollary}

\begin{remark}
This corollary removes only a proof-artifact hypothesis.  It does not prove the
primitive unit-roundoff cap \(u\le U_{\rm cap}\), the scalar smallness
inequality, local diagonal dominance/off-diagonal control, conditioning fields,
or the final concrete rectangular QR/preconditioner theorem.
\end{remark}

\begin{formalized}
\lean{Algorithms/LeastSquares/LSQRSolve.lean}
\leanfit{LSQRSolveBackwardError.of_stored_trailing_householder_sequence_topBlock_fl_backSub_gamma_bound_explicitCompactBudget_of_signed_alpha_diagDominant_storedLower_finiteMaxSourceDenURationalGammaCanonicalBounds_of_uCap}
\lean{Algorithms/RandNLA/LeastSquaresSketch.lean}
\leanfit{leverageTraceProbability_eventProb_fl_rowSampleLSObjective_le_one_add_eta_of_augmentedSpan_sample_budget_of_objective_error_and_stored_qr_diagDominant_storedLower_finiteMaxSourceDenURationalGammaCanonicalBounds_solver_of_uCap}
\end{formalized}

\begin{corollary}[Gamma-validity guards are automatic from the displayed cap]
In the same stored-lower canonical rational-gamma route, the cap hypotheses
\[
  u\le U_{\rm cap},
  \qquad
  mU_{\rm cap}<1
\]
imply \(\lean{gammaValid fp m}\).  Since the rectangular solve also assumes
\(n\le m\), monotonicity gives \(\lean{gammaValid fp n}\).  At the sampled
row-problem level the identical argument uses
\[
  u\le U_{\rm cap},
  \qquad
  sU_{\rm cap}<1,
  \qquad
  n\le s .
\]
Thus the cap-based solver and probability statements need not list
\(\lean{gammaValid fp m}\), \(\lean{gammaValid fp s}\), or
\(\lean{gammaValid fp n}\) as independent hypotheses.
\end{corollary}

\begin{remark}
This is still a hypothesis-surface cleanup.  The primitive machine cap
\(u\le U_{\rm cap}\), the displayed scalar smallness inequality, local diagonal
dominance/off-diagonal control, conditioning fields, and the final generic
QR/preconditioner theorem remain visible or open.
\end{remark}

\begin{formalized}
\lean{Analysis/Rounding.lean}
\leanfit{gammaValid_of_u_le_cap}
\lean{Algorithms/LeastSquares/LSQRSolve.lean}
\leanfit{LSQRSolveBackwardError.of_stored_trailing_householder_sequence_topBlock_fl_backSub_gamma_bound_explicitCompactBudget_of_signed_alpha_diagDominant_storedLower_finiteMaxSourceDenURationalGammaCanonicalBounds_of_uCap_no_gammaValid}
\lean{Algorithms/RandNLA/LeastSquaresSketch.lean}
\leanfit{leverageTraceProbability_eventProb_fl_rowSampleLSObjective_le_one_add_eta_of_augmentedSpan_sample_budget_of_objective_error_and_stored_qr_diagDominant_storedLower_finiteMaxSourceDenURationalGammaCanonicalBounds_solver_of_uCap_no_gammaValid}
\end{formalized}

\begin{corollary}[Stored-lower route at the actual unit roundoff]
In the stored-lower canonical rational-gamma route one may choose the displayed
cap to be the model's actual unit roundoff:
\[
  U_{\rm cap}=u .
\]
Then the primitive cap hypothesis \(u\le U_{\rm cap}\) is discharged by
reflexivity, and all \(U_{\rm cap}\)-notation disappears from the theorem
surface.  The solver-level statement assumes \(\lean{gammaValid fp m}\); the
sampled least-squares statement assumes \(\lean{gammaValid fp s}\).  The
remaining scalar smallness condition is the same canonical finite-max
inequality, but with
\[
  G_m(u)=\frac{m u}{1-mu},
  \qquad
  F_m(u)=u(1+G_m(u))(1+u)+u(1+G_m(u))+G_m(u)
        +u(1+G_m(u))(1+u)^2
\]
in the solver theorem, and with \(m\) replaced by \(s\) in the sampled
probability theorem.
\end{corollary}

\begin{remark}
This corollary is a theorem-surface reduction, not a proof of scalar smallness.
The repository also contains a checked counterexample showing that the
actual-\(u\) scalar smallness condition is not automatic from diagonal
dominance, source-denominator nonbreakdown, and \(m u<1\).  Local
diagonal/off-diagonal control, conditioning fields, scalar smallness, and the
final generic QR/preconditioner theorem therefore remain visible or open.
\end{remark}

\begin{formalized}
\lean{Algorithms/LeastSquares/LSQRSolve.lean}
\leanfit{LSQRSolveBackwardError.of_stored_trailing_householder_sequence_topBlock_fl_backSub_gamma_bound_explicitCompactBudget_of_signed_alpha_diagDominant_storedLower_finiteMaxSourceDenURationalGammaCanonicalBounds_of_actualUnitRoundoff}
\lean{Algorithms/RandNLA/LeastSquaresSketch.lean}
\leanfit{leverageTraceProbability_eventProb_fl_rowSampleLSObjective_le_one_add_eta_of_augmentedSpan_sample_budget_of_objective_error_and_stored_qr_diagDominant_storedLower_finiteMaxSourceDenURationalGammaCanonicalBounds_solver_of_actualUnitRoundoff}
\end{formalized}

\begin{corollary}[Actual-unit-roundoff route with scalar validity guard]
In the stored-lower canonical rational-gamma QR route above, the ordinary
\(\lean{gammaValid}\) hypothesis can be replaced by the displayed scalar guard
\[
  m u < 1
\]
for the local solver theorem, and by \(s u < 1\) for the sampled-row
probability theorem.  The proof chooses the actual unit roundoff
\(u=\lean{fp.u}\)
and reuses the cap-derived gamma-validity theorem with \(U_{\max}=u\).
The canonical finite-max scalar smallness inequality remains a separate
visible hypothesis.
\end{corollary}

\begin{remark}
This is only a hypothesis-surface reduction: it does not prove diagonal
dominance/off-diagonal control, scalar smallness, conditioning fields, or the
generic implementation-backed rectangular QR/preconditioner theorem.
\end{remark}

\begin{formalized}
\lean{Algorithms/LeastSquares/LSQRSolve.lean}
\leanfit{LSQRSolveBackwardError.of_stored_trailing_householder_sequence_topBlock_fl_backSub_gamma_bound_explicitCompactBudget_of_signed_alpha_diagDominant_storedLower_finiteMaxSourceDenURationalGammaCanonicalBounds_of_actualUnitRoundoff_no_gammaValid}
\lean{Algorithms/RandNLA/LeastSquaresSketch.lean}
\leanfit{leverageTraceProbability_eventProb_fl_rowSampleLSObjective_le_one_add_eta_of_augmentedSpan_sample_budget_of_objective_error_and_stored_qr_diagDominant_storedLower_finiteMaxSourceDenURationalGammaCanonicalBounds_solver_of_actualUnitRoundoff_no_gammaValid}
\end{formalized}

\begin{corollary}[Signed-alpha scalar-smallness route elimination after validity reduction]
The scalar smallness hypothesis in the preceding actual-unit-roundoff route is
not a consequence of the remaining signed Householder source fields.  Even in a
\(1\times 1\) exact-with-unit-roundoff model, local diagonal dominance, the
concrete signed-alpha equation, source denominator nonbreakdown, and
\[
  u < 1
\]
do not force the canonical finite-max rational-gamma inequality.  The
no-\(\lean{gammaValid}\) theorem states this after the validity guard has been
replaced by the scalar condition \(u=\lean{fp.u}\), matching the post-LS.2g-gc
surface.
\end{corollary}

\begin{remark}
This is a route-elimination theorem, not a positive scalar-smallness theorem:
future route-1 progress must add stronger scale/conditioning hypotheses, prove a
concrete machine-model smallness fact, or keep the finite-max scalar inequality
visible.
\end{remark}

\begin{formalized}
\lean{Algorithms/LeastSquares/LSQRSolve.lean}
\leanfit{not_forall_diagDominant_signedAlphaDef_sourceDen_actualU_implies_uRationalGammaCanonicalSmallness}
\leanfit{not_forall_diagDominant_signedAlphaDef_sourceDen_actualU_no_gammaValid_implies_uRationalGammaCanonicalSmallness}
\end{formalized}

\begin{corollary}[Stored-recurrence scalar-smallness route elimination]
The scalar smallness hypothesis is still not forced after replacing the raw
source-denominator field by the actual stored panel recurrence.  There is a
\(1\times 1\) exact-with-unit-roundoff witness satisfying the stored
Householder panel recurrence, local diagonal dominance, the signed-alpha
equation, and \(u<1\), while the canonical finite-max rational-gamma scalar is
larger than \(1\).  Thus the post-stored-lower route still needs an independent
scale, conditioning, or concrete machine-model smallness theorem, or else the
smallness condition must stay visible.
\end{corollary}

\begin{formalized}
\lean{Algorithms/LeastSquares/LSQRSolve.lean}
\leanfit{not_forall_diagDominant_signedAlphaDef_stored_trailing_sequence_actualU_no_gammaValid_implies_uRationalGammaCanonicalSmallness}
\end{formalized}

\begin{corollary}[Stored-recurrence diagonal-dominance route elimination]
The local diagonal-dominance hypothesis is also not forced by the stored-panel
actual-unit data.  The existing \(2\times2\) Householder diagonal-dominance
counterexample follows the stored recurrence under exact floating-point
arithmetic with \(u=0\), satisfies the signed-alpha equation and nonzero source
denominators, and has \((2:\mathbb R)u<1\).  Its final displayed triangular
block is still not diagonally dominant.  Thus local diagonal dominance, or an
equivalent off-diagonal-control invariant, must be proved from a stronger
pivoted/sorted loop invariant or kept visible as a domain assumption.
\end{corollary}

\begin{formalized}
\lean{Algorithms/LeastSquares/LSQRSolve.lean}
\leanfit{exactHouseholderQRDiagDominanceCounterexampleFP}
\leanfit{exactHouseholderQRDiagDominanceCounterexample_stored_step}
\leanfit{exactHouseholderQRDiagDominanceCounterexample_signed_alpha_def}
\leanfit{not_forall_signedAlphaDef_stored_trailing_sequence_actualU_sourceDen_implies_diagDominant}
\end{formalized}

\begin{corollary}[Stored-recurrence active-pivot route elimination]
The displayed active-pivot policy is likewise not forced by the stored-panel
actual-unit data.  At the initial state of the same \(2\times2\) witness,
namely \(\begin{bmatrix}1&2\\0&1\end{bmatrix}\), the second column has active
trailing norm square \(5\), while the displayed first pivot has active trailing
norm square \(1\).  Hence the finite selector theorem
\(\lean{householderActiveMaxPivotColumn\_pivot\_max}\) contradicts the claim
that the displayed first pivot is
\(\lean{householderActiveMaxPivotColumn}\).  The universal stored-surface
theorem therefore rules out deriving the active-pivot-policy equation from the
stored recurrence, signed-alpha equation, nonzero source denominators, and
\((2:\mathbb R)u<1\).  A positive proof must use a concrete pivoted/sorted loop
invariant, or keep the active-pivot equation visible.
\end{corollary}

\begin{formalized}
\lean{Algorithms/LeastSquares/LSQRSolve.lean}
\leanfit{exactHouseholderQRDiagDominanceCounterexample_not_activeMaxPivotChoice}
\leanfit{not_forall_signedAlphaDef_stored_trailing_sequence_actualU_sourceDen_implies_activeMaxPivotChoice}
\end{formalized}

\begin{corollary}[Stored-recurrence scalar row-budget route elimination]
The scalar row-budget fields in the active/prefix QR route are not forced by
the stored-panel actual-unit data either.  The same exact floating-point
witness with \(u=0\) follows the stored recurrence, satisfies the signed-alpha
equation and nonzero source denominators, and has \((2:\mathbb R)u<1\).  With
the constant stage budget \(B_k=2\), the scalar stage-diagonal lower-bound
defect is positive.  With the nonnegative constant stage budget \(B_k=3\), the
scalar stage-budget/row-max comparison defect is positive.  Thus these scalar
fields must be proved from a stronger pivoted/sorted/off-diagonal-control loop
invariant or kept visible as source assumptions.
\end{corollary}

\begin{formalized}
\lean{Algorithms/LeastSquares/LSQRSolve.lean}
\leanfit{not_forall_signedAlphaDef_stored_trailing_sequence_actualU_sourceDen_implies_stageDiagLowerDefectBudget_nonpos}
\leanfit{not_forall_signedAlphaDef_stored_trailing_sequence_actualU_sourceDen_nonnegStageBudget_implies_stageRowMaxComparisonDefectBudget_nonpos}
\end{formalized}

\begin{corollary}[Diagonal-dominant stored-recurrence scalar row-budget route elimination]
Even adding local diagonal dominance to the stored-panel actual-unit surface
does not force the scalar row-budget fields.  The exact stored sequence
starting from \(\begin{bmatrix}3&2\\0&1\end{bmatrix}\) follows the stored
recurrence under exact floating-point arithmetic with \(u=0\), satisfies the
signed-alpha equation and nonzero source denominators, has \((2:\mathbb R)u<1\),
and is locally diagonally dominant at every displayed leading block.  Still,
the constant stage budget \(B_k=4\) makes the scalar stage-diagonal lower-bound
defect positive, and the nonnegative constant stage budget \(B_k=3\) makes the
scalar stage-budget/row-max comparison defect positive.  Thus diagonal
dominance does not remove the need to prove these scalar fields from a
concrete loop invariant or keep them visible.
\end{corollary}

\begin{formalized}
\lean{Algorithms/LeastSquares/LSQRSolve.lean}
\leanfit{storedDiagDominantComparisonCounterexample_stored_step}
\leanfit{storedDiagDominantComparisonCounterexample_diagDominant}
\leanfit{not_forall_diagDominant_signedAlphaDef_stored_trailing_sequence_actualU_sourceDen_nonnegStageBudget_implies_stageDiagLowerDefectBudget_nonpos}
\leanfit{not_forall_diagDominant_signedAlphaDef_stored_trailing_sequence_actualU_sourceDen_nonnegStageBudget_implies_stageRowMaxComparisonDefectBudget_nonpos}
\end{formalized}

\begin{corollary}[Active-pivot diagonal-dominant stored-recurrence comparison route elimination]
The same diagonal-dominant stored witness also satisfies the finite active
max-pivot selector policy.  Starting from
\(\begin{bmatrix}3&2\\0&1\end{bmatrix}\), the exact stored sequence follows
the stored recurrence with \(u=0\), satisfies the signed-alpha equation,
source-denominator nonzero facts, actual-unit validity, local diagonal
dominance, the active-pivot policy, and nonnegative stage budget \(B_k=3\).
Nevertheless the scalar stage-budget/row-max comparison defect is positive.
Thus active max-pivoting does not turn diagonal dominance into the scalar
comparison invariant; that invariant must be proved from a stronger concrete
loop argument or kept visible.
\end{corollary}

\begin{formalized}
\lean{Algorithms/LeastSquares/LSQRSolve.lean}
\leanfit{storedDiagDominantComparisonCounterexample_activeMaxPivotChoice}
\leanfit{not_forall_diagDominant_activeMaxPivot_signedAlphaDef_stored_trailing_sequence_actualU_sourceDen_nonnegStageBudget_implies_stageRowMaxComparisonDefectBudget_nonpos}
\end{formalized}

\begin{corollary}[Compact-product active-pivot diagonal-dominant stored-recurrence comparison route elimination]
Adding the stored right-hand-side recurrence and global compact-product
smallness still does not force the scalar comparison invariant.  Pair the same
stored sequence from \(\begin{bmatrix}3&2\\0&1\end{bmatrix}\) with the zero
right-hand-side trace.  In exact arithmetic the RHS trace follows the stored
RHS recurrence, the compact relative budget is zero, and hence the finite
compact-product budget is strictly below one.  The local diagonal dominance,
active-pivot policy, signed-alpha equation, source-denominator facts, and
actual-unit validity all remain available.  Nevertheless \(B_k=3\) still makes
the scalar stage-budget/row-max comparison defect positive.
\end{corollary}

\begin{formalized}
\lean{Algorithms/LeastSquares/LSQRSolve.lean}
\leanfit{storedDiagDominantComparisonCounterexample_rhs_step}
\leanfit{storedDiagDominantComparisonCounterexample_compactSequenceRelativeBudget_eq_zero}
\leanfit{storedDiagDominantComparisonCounterexample_compactSequenceProductBudget_lt_one}
\leanfit{not_forall_diagDominant_activeMaxPivot_product_signedAlphaDef_stored_trailing_sequence_actualU_sourceDen_nonnegStageBudget_implies_stageRowMaxComparisonDefectBudget_nonpos}
\end{formalized}

\begin{corollary}[Finite-max active-pivot diagonal-dominant stored-recurrence comparison route elimination]
The canonical finite-max compact-product scalar also does not force the scalar
comparison invariant.  For the same exact stored sequence and zero RHS trace,
the compact relative budget is zero, so the displayed finite-max scalar used by
the active branch reduces to \(0<1\).  The stored panel recurrence, stored RHS
recurrence, signed-alpha equation, source-denominator facts, actual-unit
validity, local diagonal dominance, active-pivot policy, and nonnegative stage
budgets still hold.  Nevertheless \(B_k=3\) leaves the scalar
stage-budget/row-max comparison defect positive.
\end{corollary}

\begin{formalized}
\lean{Algorithms/LeastSquares/LSQRSolve.lean}
\leanfit{storedDiagDominantComparisonCounterexample_finiteMaxSmallness}
\leanfit{not_forall_diagDominant_activeMaxPivot_finiteMaxSmallness_signedAlphaDef_stored_trailing_sequence_actualU_sourceDen_nonnegStageBudget_implies_stageRowMaxComparisonDefectBudget_nonpos}
\end{formalized}

\begin{corollary}[Signed-stage global-budget finite-max active-pivot comparison route elimination]
Even adding the signed-stage finite global compact-budget recurrence does not
force the scalar comparison invariant.  Use the same exact stored sequence and
zero RHS trace, but take the nonconstant nonnegative stage budget
\[
  B_0=0,\qquad B_1=3,\qquad B_k=10\quad(k\ge 2).
\]
The exact global compact-update budget is zero at both displayed stages, and
the two-dimensional Cox--Higham factor is at most \(3\), so the finite global
compact-budget recurrence holds.  The stored panel recurrence, stored RHS
recurrence, signed-alpha equation, source-denominator facts, actual-unit
validity, local diagonal dominance, active-pivot policy, canonical finite-max
compact-product scalar, and nonnegative stage budgets are all present.
Nevertheless the stage-one scalar stage-budget/row-max comparison defect is
positive.
\end{corollary}

\begin{formalized}
\lean{Algorithms/LeastSquares/LSQRSolve.lean}
\leanfit{storedDiagDominantComparisonCounterexample_globalBudgetStageBudget_nonneg}
\leanfit{storedDiagDominantComparisonCounterexample_globalCompactBudget_eq_zero}
\leanfit{storedDiagDominantComparisonCounterexample_globalCompactBudget_recurrence}
\leanfit{storedDiagDominantComparisonCounterexample_stageRowMaxComparisonDefectBudget_pos_globalBudgetStageBudget}
\leanfit{not_forall_diagDominant_activeMaxPivot_globalCompactBudget_finiteMaxSmallness_signedAlphaDef_stored_trailing_sequence_actualU_sourceDen_nonnegStageBudget_implies_stageRowMaxComparisonDefectBudget_nonpos}
\end{formalized}

\begin{corollary}[Row-max-granted signed-stage global-budget comparison route elimination]
The finite row-max/diagonal defect scalar is not enough to recover the missing
stage-budget/row-max comparison scalar.  For the same exact stored sequence,
zero RHS trace, and nonconstant stage budget \(B_0=0\), \(B_1=3\),
\(B_k=10\) for \(k\ge 2\), local diagonal dominance packages
\(\mathsf{storedQRRowMaxDiagDefectBudget}\le 0\).  The stored panel and RHS
recurrences, signed-alpha equation, source-denominator facts, actual-unit
validity, active-pivot policy, signed-stage global compact-budget recurrence,
canonical finite-max compact-product scalar, row-max defect scalar, and
nonnegative stage budgets all hold together.  Nevertheless the stage-one
scalar stage-budget/row-max comparison defect is positive.
\end{corollary}

\begin{formalized}
\lean{Algorithms/LeastSquares/LSQRSolve.lean}
\leanfit{storedDiagDominantComparisonCounterexample_rowMaxDiagDefectBudget_nonpos}
\leanfit{not_forall_diagDominant_activeMaxPivot_globalCompactBudget_finiteMaxSmallness_rowMaxDiagDefect_signedAlphaDef_stored_trailing_sequence_actualU_sourceDen_nonnegStageBudget_implies_stageRowMaxComparisonDefectBudget_nonpos}
\end{formalized}

\begin{corollary}[Determinant-granted signed-stage global-budget comparison route elimination]
Nonsingular displayed leading blocks are not enough to recover the missing
stage-budget/row-max comparison scalar either.  For the same exact stored
sequence, zero RHS trace, and nonconstant stage budget \(B_0=0\), \(B_1=3\),
\(B_k=10\) for \(k\ge 2\), the displayed leading blocks have determinants
\(3\) and \(-3\).  The stored panel and RHS recurrences, signed-alpha equation,
source-denominator facts, actual-unit validity, local diagonal dominance,
active-pivot policy, row-max defect scalar, signed-stage global compact-budget
recurrence, canonical finite-max compact-product scalar, determinant
nonbreakdown, and nonnegative stage budgets all hold together.  Nevertheless
the stage-one scalar stage-budget/row-max comparison defect is positive.
\end{corollary}

\begin{formalized}
\lean{Algorithms/LeastSquares/LSQRSolve.lean}
\leanfit{storedDiagDominantComparisonCounterexample_leadingBlock_det_ne_zero}
\leanfit{not_forall_diagDominant_activeMaxPivot_leadingBlock_det_ne_zero_globalCompactBudget_finiteMaxSmallness_rowMaxDiagDefect_signedAlphaDef_stored_trailing_sequence_actualU_sourceDen_nonnegStageBudget_implies_stageRowMaxComparisonDefectBudget_nonpos}
\end{formalized}

\begin{corollary}[Conditioning-granted signed-stage global-budget comparison route elimination]
Local conditioning and dual compact-budget data are not enough to recover the
missing stage-budget/row-max comparison scalar either.  For the same exact
stored sequence, zero RHS trace, and nonconstant stage budget \(B_0=0\),
\(B_1=3\), \(B_k=10\) for \(k\ge 2\), the compact component budgets vanish in
the exact arithmetic model.  Taking the local \(\kappa_\infty\) budget to be
the exact displayed condition expression and taking \(K\) to be the exact
self-norm squared expression plus one supplies positive local conditioning and
dual compact-budget fields.  The stored panel and RHS recurrences,
signed-alpha equation, source-denominator facts, actual-unit validity, local
diagonal dominance, active-pivot policy, determinant nonbreakdown,
row-max defect scalar, signed-stage global compact-budget recurrence,
canonical finite-max compact-product scalar, and nonnegative stage budgets all
hold together.  Nevertheless the stage-one scalar stage-budget/row-max
comparison defect is positive.
\end{corollary}

\begin{formalized}
\lean{Algorithms/LeastSquares/LSQRSolve.lean}
\leanfit{storedDiagDominantComparisonCounterexample_compactComponentBudget_eq_zero}
\leanfit{storedDiagDominantComparisonCounterexampleKappaBudget}
\leanfit{storedDiagDominantComparisonCounterexampleKappaNormSqBudget}
\leanfit{storedDiagDominantComparisonCounterexample_kappaBudget_le}
\leanfit{storedDiagDominantComparisonCounterexample_kappaNormSqBudget_pos}
\leanfit{storedDiagDominantComparisonCounterexample_kappaNormSqBudget}
\leanfit{storedDiagDominantComparisonCounterexample_dualBudget}
\leanfit{not_forall_diagDominant_activeMaxPivot_leadingBlock_det_ne_zero_kappaInf_selfNorm_dualBudget_globalCompactBudget_finiteMaxSmallness_rowMaxDiagDefect_signedAlphaDef_stored_trailing_sequence_actualU_sourceDen_nonnegStageBudget_implies_stageRowMaxComparisonDefectBudget_nonpos}
\end{formalized}

\begin{corollary}[Final-surface signed-stage global-budget comparison route elimination]
Even the final active/prefix global-product surface does not recover the
missing stage-budget/row-max comparison scalar automatically.  Reuse the same
exact stored sequence and zero RHS trace, but take
\[
  B_0=3,\qquad B_1=10,\qquad B_k=30\quad(k\ge 2).
\]
This final-surface budget satisfies the initial full-block and displayed
row-budget fields, is nonnegative and monotone, and still satisfies the exact
signed-stage global compact-budget recurrence and global compact-product
smallness.  The determinant/nonbreakdown facts, local \(\kappa_\infty\) and
self-norm budgets, dual compact-budget inequality, row-max defect scalar,
active-pivot policy, and local diagonal dominance all remain available.
Nevertheless the stage-one comparison scalar contains the positive displayed
gap \(10-2\).  Thus the final surface must still prove
\(\mathsf{storedQRStageRowMaxComparisonDefectBudget}\le 0\) from a genuine
loop invariant, or expose it as a source/domain assumption.
\end{corollary}

\begin{formalized}
\lean{Algorithms/LeastSquares/LSQRSolve.lean}
\leanfit{storedDiagDominantComparisonCounterexampleFinalSurfaceStageBudget}
\leanfit{storedDiagDominantComparisonCounterexample_finalSurfaceStageBudget_nonneg}
\leanfit{storedDiagDominantComparisonCounterexample_finalSurfaceStageBudget_mono}
\leanfit{storedDiagDominantComparisonCounterexample_finalSurface_init}
\leanfit{storedDiagDominantComparisonCounterexample_finalSurface_initBlock}
\leanfit{storedDiagDominantComparisonCounterexample_finalSurface_globalCompactBudget_recurrence}
\leanfit{storedDiagDominantComparisonCounterexample_stageRowMaxComparisonDefectBudget_pos_finalSurfaceStageBudget}
\leanfit{not_forall_diagDominant_activeMaxPivot_leadingBlock_det_ne_zero_kappaInf_selfNorm_dualBudget_init_globalProduct_globalCompactBudget_rowMaxDiagDefect_signedAlphaDef_stored_trailing_sequence_actualU_sourceDen_nonneg_monoStageBudget_implies_stageRowMaxComparisonDefectBudget_nonpos}
\end{formalized}

\begin{corollary}[Final-surface source-denominator comparison route elimination]
The source-denominator rational-gamma finite-max route also does not recover
the missing stage-budget/row-max comparison scalar automatically.  On the same
exact stored sequence and zero RHS trace, choose the source cap \(U_{\max}=0\).
Then the rational-gamma term and the Householder-factor cap term vanish, so
the canonical source-denominator scalar smallness inequality is true
independently of the finite maximum budgets.  All final-surface fields from
the preceding corollary remain available, with the raw global-product
smallness replaced by the displayed source-denominator cap-smallness fields.
Nevertheless the stage-one comparison scalar still contains the positive gap
\(10-2\).  Thus the canonical source-denominator scalar can be a valid product
route assumption, but it cannot replace the separate
\(\mathsf{storedQRStageRowMaxComparisonDefectBudget}\le 0\) invariant.
\end{corollary}

\begin{formalized}
\lean{Algorithms/LeastSquares/LSQRSolve.lean}
\leanfit{storedDiagDominantComparisonCounterexample_finalSurface_sourceDenURationalGammaCanonicalSmallness}
\leanfit{not_forall_diagDominant_activeMaxPivot_leadingBlock_det_ne_zero_kappaInf_selfNorm_dualBudget_init_sourceDenURationalGammaCanonical_globalCompactBudget_rowMaxDiagDefect_signedAlphaDef_stored_trailing_sequence_actualU_sourceDen_nonneg_monoStageBudget_implies_stageRowMaxComparisonDefectBudget_nonpos}
\end{formalized}

\begin{corollary}[Horizon-clamped source-denominator QR handoff]
For the active/prefix diagonal-dominant source-denominator route, the
stage-budget sequence need not be globally monotone on all natural indices.
If the signed-stage compact recurrence and nonnegativity hold through the QR
horizon, and if the scalar comparison defect
\[
  \mathsf{storedQRStageRowMaxComparisonDefectBudget}\le 0
\]
holds for the original stage budget, then the same comparison certificate
holds after replacing the budget by
\(\mathsf{qrStageHorizonBudget}\ n\).  The clamped budget is globally
monotone by construction and agrees with the original budget on every stage
used by the QR loop.  Consequently the source-denominator rational-gamma
source-control theorem, and its sampled equation-(8), actual-unit, and
stored-lower probability siblings, can be stated without a samplewise global
stage-budget monotonicity hypothesis.  The scalar comparison defect, local
diagonal dominance, active-pivot policy, signed-stage recurrence, and
source-denominator cap-smallness assumptions remain visible.
\end{corollary}

\begin{formalized}
\lean{Algorithms/LeastSquares/LSQRSolve.lean}
\leanfit{storedQRStageRowMaxComparisonDefectBudget_nonpos_of_horizonBudget}
\leanfit{StoredQRSourceOffDiagonalControl.of_stored_trailing_sequence_diagDominant_signedStageUniformBudget_globalCompactBudget_activeMaxPivot_completedColumns_finiteMaxSourceDenURationalGammaCanonicalBounds_stageRowMaxComparisonDefect_offdiag_rows_of_horizonBudget}
\lean{Algorithms/RandNLA/LeastSquaresSketch.lean}
\leanfit{leverageTraceProbability_eventProb_fl_rowSampleLSObjective_le_one_add_eta_of_augmentedSpan_sample_budget_of_objective_error_and_stored_qr_activePrefix_finiteMaxSourceDenURationalGammaCanonicalBounds_activeMaxPivot_diagDominant_stageRowMaxComparisonDefect_solver_of_horizonBudget}
\leanfit{leverageTraceProbability_eventProb_fl_rowSampleLSObjective_le_one_add_eta_of_augmentedSpan_sample_budget_of_objective_error_and_stored_qr_activePrefix_finiteMaxSourceDenURationalGammaCanonicalBounds_activeMaxPivot_diagDominant_stageRowMaxComparisonDefect_solver_of_actualUnitRoundoff_no_gammaValid_of_horizonBudget}
\leanfit{leverageTraceProbability_eventProb_fl_rowSampleLSObjective_le_one_add_eta_of_augmentedSpan_sample_budget_of_objective_error_and_stored_qr_activePrefix_storedLower_finiteMaxSourceDenURationalGammaCanonicalBounds_activeMaxPivot_diagDominant_stageRowMaxComparisonDefect_solver_of_actualUnitRoundoff_no_gammaValid_of_horizonBudget}
\end{formalized}

\begin{corollary}[Equation (8) assembly from active/prefix global-product QR]
Let \(A,b\) be the original least-squares data and let the row sampler use
leverage probabilities from the augmented span basis of \([A\ b]\).  Assume
the Bennett sample-budget hypotheses for the two-sided finite-Loewner event,
the positive-probability support event, and the rounded sampled-row objective
slack condition.  For each sample trace, suppose the rounded stored QR loop
satisfies the signed-stage recurrences, either the norm-square nonbreakdown
margin or the local \(\kappa_\infty\)/dual-budget hypotheses that imply it,
local leading-block nonsingularity, the scalar global compact-step recurrence,
the off-diagonal-row diagonal lower bounds, pivot maximality, and the global
compact-product condition
\[
  \Pi_{\max}<1 .
\]
If the reported solution is the back-substitution solution from the final
stored QR block, and the reported perturbation radii are the repository
\(\lean{qrSolveFinalGramBudget}\) and \(\lean{qrSolveFinalRhsBudget}\), then
with probability at least \(1-\delta\),
\[
  \|Ax_{\mathrm{comp}}-b\|_2^2
  \le
  (1+\eta)\,\|Ax_\star-b\|_2^2 .
\]
\end{corollary}

\begin{corollary}[Equation (8) horizon-clamped active/prefix QR assembly]
In the active/prefix global-product equation~(8) assembly above, the samplewise
global stage-budget monotonicity hypothesis may be removed from the theorem
surface.  It is enough to assume the finite compact-step recurrence and
nonnegative stage budgets for each sample trace, together with the same
diagonal lower-bound, pivot-maximality, nonbreakdown, and global-product
fields.  The proof applies the horizon-clamped stored-QR source-control wrapper
to each sample, so no budget behavior beyond the QR horizon is asserted or
needed by the probability theorem.  The same reduction also composes with the
local \(\kappa_\infty\)/dual-budget adapter that derives the norm-square
nonbreakdown margin, and with the active-max-pivot wrapper that recovers raw
pivot maximality from the finite selector.
\end{corollary}

\begin{formalized}
\lean{Algorithms/RandNLA/LeastSquaresSketch.lean}
\leanfit{leverageTraceProbability_eventProb_fl_rowSampleLSObjective_le_one_add_eta_of_augmentedSpan_sample_budget_of_objective_error_and_stored_qr_activePrefix_globalProduct_solver_of_horizonBudget}
\leanfit{leverageTraceProbability_eventProb_fl_rowSampleLSObjective_le_one_add_eta_of_augmentedSpan_sample_budget_of_objective_error_and_stored_qr_activePrefix_globalProduct_kappaInf_dualBudget_solver_of_horizonBudget}
\leanfit{leverageTraceProbability_eventProb_fl_rowSampleLSObjective_le_one_add_eta_of_augmentedSpan_sample_budget_of_objective_error_and_stored_qr_activePrefix_globalProduct_activeMaxPivot_kappaInf_dualBudget_solver_of_horizonBudget}
\leanfit{leverageTraceProbability_eventProb_fl_rowSampleLSObjective_le_one_add_eta_of_augmentedSpan_sample_budget_of_objective_error_and_stored_qr_activePrefix_globalProduct_activeMaxPivot_rowMaxDiagDefect_stageBudgetLeRowMax_kappaInf_dualBudget_solver_of_horizonBudget}
\leanfit{leverageTraceProbability_eventProb_fl_rowSampleLSObjective_le_one_add_eta_of_augmentedSpan_sample_budget_of_objective_error_and_stored_qr_activePrefix_globalProduct_activeMaxPivot_rowMaxDiagDefect_stageBudgetLeRowMax_kappaInf_dualBudget_solver_of_actualUnitRoundoff_no_gammaValid_of_horizonBudget}
\end{formalized}


\begin{formalized}
\lean{Algorithms/RandNLA/LeastSquaresSketch.lean}
\leanfit{leverageTraceProbability_eventProb_fl_rowSampleLSObjective_le_one_add_eta_of_augmentedSpan_sample_budget_of_objective_error_and_stored_qr_activePrefix_globalProduct_solver}
\leanfit{leverageTraceProbability_eventProb_fl_rowSampleLSObjective_le_one_add_eta_of_augmentedSpan_sample_budget_of_objective_error_and_stored_qr_activePrefix_globalProduct_solver_of_horizonBudget}
\leanfit{leverageTraceProbability_eventProb_fl_rowSampleLSObjective_le_one_add_eta_of_augmentedSpan_sample_budget_of_objective_error_and_stored_qr_activePrefix_globalProduct_kappaInf_dualBudget_solver}
\leanfit{leverageTraceProbability_eventProb_fl_rowSampleLSObjective_le_one_add_eta_of_augmentedSpan_sample_budget_of_objective_error_and_stored_qr_activePrefix_globalProduct_kappaInf_dualBudget_solver_of_horizonBudget}
\leanfit{leverageTraceProbability_eventProb_fl_rowSampleLSObjective_le_one_add_eta_of_augmentedSpan_sample_budget_of_objective_error_and_stored_qr_activePrefix_globalProduct_activeMaxPivot_kappaInf_dualBudget_solver}
\leanfit{leverageTraceProbability_eventProb_fl_rowSampleLSObjective_le_one_add_eta_of_augmentedSpan_sample_budget_of_objective_error_and_stored_qr_activePrefix_globalProduct_activeMaxPivot_kappaInf_dualBudget_solver_of_horizonBudget}
\leanfit{leverageTraceProbability_eventProb_fl_rowSampleLSObjective_le_one_add_eta_of_augmentedSpan_sample_budget_of_objective_error_and_stored_qr_activePrefix_globalProduct_activeMaxPivot_rowMaxDiagDefect_stageBudgetLeRowMax_kappaInf_dualBudget_solver}
\leanfit{leverageTraceProbability_eventProb_fl_rowSampleLSObjective_le_one_add_eta_of_augmentedSpan_sample_budget_of_objective_error_and_stored_qr_activePrefix_globalProduct_activeMaxPivot_rowMaxDiagDefect_stageBudgetLeRowMax_kappaInf_dualBudget_solver_of_horizonBudget}
\leanfit{leverageTraceProbability_eventProb_fl_rowSampleLSObjective_le_one_add_eta_of_augmentedSpan_sample_budget_of_objective_error_and_stored_qr_activePrefix_globalProduct_activeMaxPivot_rowMaxDiagDefect_stageBudgetLeRowMax_kappaInf_dualBudget_solver_of_actualUnitRoundoff_no_gammaValid}
\leanfit{leverageTraceProbability_eventProb_fl_rowSampleLSObjective_le_one_add_eta_of_augmentedSpan_sample_budget_of_objective_error_and_stored_qr_activePrefix_globalProduct_activeMaxPivot_rowMaxDiagDefect_stageBudgetLeRowMax_kappaInf_dualBudget_solver_of_actualUnitRoundoff_no_gammaValid_of_horizonBudget}
\leanfit{leverageTraceProbability_eventProb_fl_rowSampleLSObjective_le_one_add_eta_of_augmentedSpan_sample_budget_of_objective_error_and_stored_qr_rowBudgetControl_globalProduct_kappaInf_dualBudget_solver}
\leanfit{leverageTraceProbability_eventProb_fl_rowSampleLSObjective_le_one_add_eta_of_augmentedSpan_sample_budget_of_objective_error_and_stored_qr_rowBudgetControl_globalProduct_activeMaxPivot_kappaInf_dualBudget_solver}
\end{formalized}

\begin{corollary}[Equation (8) assembly from diagonal-dominant QR]
Let \(A,b\) be the original least-squares data and use the same leverage
sampling law on the augmented span basis of \([A\ b]\).  Assume the Bennett
sample-budget hypotheses, the positive-probability support event, and the
rounded sampled-row objective slack condition from the preceding equation~(8)
theorem.  For each sample trace, suppose the rounded stored QR loop satisfies
the signed-stage recurrences.  The compact-product assumption may be supplied
either as the global condition \(\Pi_{\max}<1\), or in the canonical finite-max
form
\[
  2D_*\left(s\,(c_{\rm seq}N_*)^2\right)<1,
\]
where \(s\) is the sampled row count, \(D_*\) is the maximum local
diagonal-dominant inverse factor, \(N_*\) is the maximum displayed pivot-column
norm, and \(c_{\rm seq}\) is the stored QR sequence budget for that sample
trace.
Suppose further that every displayed local leading block \(S_k\) is
upper-triangular and diagonally dominant in the repository sense
\[
  \lean{IsDiagDominantUpper}\;(k+1)\;S_k,
\]
and that the local determinant and compact-smallness hypotheses are available.
There are two formalized theorem surfaces.  The first keeps the local
\(\kappa_\infty\) self-norm and dual compact-budget hypotheses explicit, so
that the dimensioned norm-square nonbreakdown margin follows from the
condition-number route.  The second is the concrete-dual finite-max variant:
it uses the repository's diagonal-dominant inverse-budget theorem to choose
the auxiliary inverse budget internally, and therefore does not expose
separate \(\kappa_\infty\), \(K_k\), or dual compact-budget fields.  In that
concrete-dual form, the remaining QR-domain assumptions are local leading
block nonsingularity, local diagonal dominance, and the displayed scalar
finite-max smallness inequality.  The final determinant-free concrete-dual
form observes that \(\lean{IsDiagDominantUpper}\) already contains upper
triangular shape and nonzero diagonal entries, so local leading-block
nonsingularity follows from diagonal dominance itself.  In that form, the
visible QR-domain assumptions are local diagonal dominance and the displayed
scalar finite-max smallness inequality.  Under any of these theorem surfaces,
the stored-QR back-substitution solution \(\widehat x\) satisfies, with
probability at least \(1-\delta\),
\[
  \|A\widehat x-b\|_2^2
  \le
  (1+\eta)\,\|Ax_\star-b\|_2^2 .
\]
\end{corollary}


\begin{formalized}
\lean{Algorithms/LeastSquares/LSQRSolve.lean}
\leanfit{StoredQRSourceOffDiagonalControl.of_stored_trailing_sequence_leadingBlock_det_ne_zero_normSqBudget_diagDominant_globalProduct}
\leanfit{StoredQRSourceOffDiagonalControl.of_stored_trailing_sequence_leadingBlock_det_ne_zero_kappaInf_selfNorm_dualBudget_diagDominant_globalProduct}
\leanfit{LSQRSolveBackwardError.of_stored_trailing_householder_sequence_topBlock_fl_backSub_gamma_bound_explicitCompactBudget_of_signed_alpha_leadingBlock_det_ne_zero_kappaInf_selfNorm_dualBudget_diagDominant_globalProduct}
\leanfit{StoredQRSourceOffDiagonalControl.of_stored_trailing_sequence_leadingBlock_det_ne_zero_normSqBudget_diagDominant_finiteMaxSmallness}
\leanfit{StoredQRSourceOffDiagonalControl.of_stored_trailing_sequence_leadingBlock_det_ne_zero_kappaInf_selfNorm_dualBudget_diagDominant_finiteMaxSmallness}
\leanfit{LSQRSolveBackwardError.of_stored_trailing_householder_sequence_topBlock_fl_backSub_gamma_bound_explicitCompactBudget_of_signed_alpha_leadingBlock_det_ne_zero_kappaInf_selfNorm_dualBudget_diagDominant_finiteMaxSmallness}
\leanfit{LSQRSolveBackwardError.of_stored_trailing_householder_sequence_topBlock_fl_backSub_gamma_bound_explicitCompactBudget_of_signed_alpha_leadingBlock_det_ne_zero_diagDominant_finiteMaxSmallness_concreteDual}
\leanfit{det_ne_zero_of_diagDominantUpper}
\leanfit{LSQRSolveBackwardError.of_stored_trailing_householder_sequence_topBlock_fl_backSub_gamma_bound_explicitCompactBudget_of_signed_alpha_diagDominant_finiteMaxSmallness_concreteDual}
\lean{Algorithms/RandNLA/LeastSquaresSketch.lean}
\leanfit{leverageTraceProbability_eventProb_fl_rowSampleLSObjective_le_one_add_eta_of_augmentedSpan_sample_budget_of_objective_error_and_stored_qr_diagDominant_globalProduct_kappaInf_dualBudget_solver}
\leanfit{leverageTraceProbability_eventProb_fl_rowSampleLSObjective_le_one_add_eta_of_augmentedSpan_sample_budget_of_objective_error_and_stored_qr_offDiagonalControl_solver_of_diagDominant_finiteMaxSmallness}
\leanfit{leverageTraceProbability_eventProb_fl_rowSampleLSObjective_le_one_add_eta_of_augmentedSpan_sample_budget_of_objective_error_and_stored_qr_diagDominant_finiteMaxSmallness_kappaInf_dualBudget_solver}
\leanfit{leverageTraceProbability_eventProb_fl_rowSampleLSObjective_le_one_add_eta_of_augmentedSpan_sample_budget_of_objective_error_and_stored_qr_diagDominant_finiteMaxSmallness_concreteDual_solver}
\leanfit{leverageTraceProbability_eventProb_fl_rowSampleLSObjective_le_one_add_eta_of_augmentedSpan_sample_budget_of_objective_error_and_stored_qr_diagDominant_finiteMaxSmallness_concreteDual_solver_of_diagDominant}
\end{formalized}

\begin{corollary}[Least-squares QR certificate from prefix-span and active-entry budget]
In the stored trailing Householder QR loop, assume the same rounded matrix and
right-hand-side recurrences used in the stored-loop solver certificate.  At
every pivot \(k<n\), suppose the prefix-span QR nonbreakdown invariant holds,
the Householder normalization and sign choice hold, and there is an active row
\(i_k\ge k\) whose pivot-column magnitude dominates the compact diagonal update
budget,
\[
  B_k<|A_k(i_k,k)|.
\]
Assume also the columnwise compact-update budget domination for the matrix
panel and right-hand side, the \(\gamma\)-validity hypotheses, and the final
Gram/RHS norm-budget bounds.  Then the computed back-substitution vector
\[
  \widehat x=\mathrm{fl\_backSub}(\widehat R,\widehat c)
\]
satisfies the repository's rectangular QR least-squares backward-error
certificate,
\[
  \mathrm{LSQRSolveBackwardError}(A^TA,A^Tb,\widehat x;c_G,c_g).
\]
\end{corollary}


\begin{corollary}[Least-squares QR certificate from prefix-span and norm-square budget]
In the same stored trailing Householder QR loop, assume prefix-span
nonbreakdown, sign choice, Householder normalization, columnwise compact-update
budget domination, \(\gamma\)-validity, and final Gram/RHS norm-budget bounds.
Replace the active-entry witness by the dimensioned per-pivot margin
\[
  mB_k^2 < \|A_k(k{:}m,k)\|_2^2 .
\]
Then the computed back-substitution vector
\[
  \widehat x=\mathrm{fl\_backSub}(\widehat R,\widehat c)
\]
satisfies the same rectangular QR least-squares backward-error certificate,
\[
  \mathrm{LSQRSolveBackwardError}(A^TA,A^Tb,\widehat x;c_G,c_g).
\]
\end{corollary}


\begin{corollary}[Least-squares QR certificate from prefix-span and leading-dual budget]
In the same stored trailing Householder QR loop, assume prefix-span,
Householder normalization, sign choice, columnwise compact-update budget
domination, \(\gamma\)-validity, and final Gram/RHS norm-budget bounds.  For
each pivot \(k<n\), suppose there is a leading dual row satisfying
\[
  \|L_{k,\ast}\|_2^2\le K_k,\qquad K_k>0,
\]
and the compact diagonal update budget obeys
\[
  mB_k^2 < \frac1{K_k}.
\]
Then the computed back-substitution vector
\[
  \widehat x=\mathrm{fl\_backSub}(\widehat R,\widehat c)
\]
satisfies the same rectangular QR least-squares backward-error certificate,
\[
  \mathrm{LSQRSolveBackwardError}(A^TA,A^Tb,\widehat x;c_G,c_g).
\]
\end{corollary}


\begin{corollary}[Prefix-span leading-dual QR certificate with repository budgets]
In the preceding leading-dual setting, the final perturbation radii may be
chosen to be the repository radii
\[
  c_G=\mathrm{qrSolveFinalGramBudget}(A,\widehat R,c_{\mathrm{step}}),
  \qquad
  c_g=\mathrm{qrSolveFinalRhsBudget}(A,b,\widehat R,c_{\mathrm{step}}).
\]
Moreover, if
\[
  c_{\mathrm{step}}
    = \mathrm{storedQRCompactSequenceRelativeBudget}
        (A_0,A_1,\ldots,A_n;b_0,b_1,\ldots,b_n;\alpha),
\]
then the columnwise panel and RHS compact-update domination hypotheses are
discharged by the local stored-loop budget lemmas.  Thus the computed vector
\(\widehat x=\mathrm{fl\_backSub}(\widehat R,\widehat c)\) satisfies
\[
  \mathrm{LSQRSolveBackwardError}
    (A^TA,A^Tb,\widehat x;
     \mathrm{qrSolveFinalGramBudget},
     \mathrm{qrSolveFinalRhsBudget})
\]
under the source-faithful prefix-span, leading-dual, sign-choice, and
dual compact-smallness hypotheses.
\end{corollary}


\begin{corollary}[Least-squares QR certificate from prefix-span and local inverse row budget]
In the same stored trailing Householder QR loop, replace the supplied ambient
dual row by a local left inverse \(C_k\) of the leading block \(S_k\).  Assume
prefix-span, Householder normalization, sign choice, columnwise compact-update
budget domination, \(\gamma\)-validity, and final Gram/RHS norm-budget bounds.
For every pivot \(k<n\), suppose
\[
  \|C_k(k,:)\|_2^2\le K_k,\qquad K_k>0,
  \qquad
  mB_k^2 < \frac1{K_k}.
\]
Then the computed back-substitution vector
\[
  \widehat x=\mathrm{fl\_backSub}(\widehat R,\widehat c)
\]
satisfies the same rectangular QR least-squares backward-error certificate,
\[
  \mathrm{LSQRSolveBackwardError}(A^TA,A^Tb,\widehat x;c_G,c_g).
\]
\end{corollary}


\begin{corollary}[Local inverse row-budget QR certificate with repository budgets]
In the preceding local-leading-block inverse row setting, the final radii may
again be chosen as
\(\mathrm{qrSolveFinalGramBudget}\) and
\(\mathrm{qrSolveFinalRhsBudget}\).  If the step constant is chosen as
\[
  c_{\mathrm{step}}
    = \mathrm{storedQRCompactSequenceRelativeBudget}
        (A_0,A_1,\ldots,A_n;b_0,b_1,\ldots,b_n;\alpha),
\]
then the columnwise panel and RHS compact-update domination hypotheses are
also supplied by the stored compact-budget lemmas.  Hence the local left
inverse row-norm budget
\[
  \|C_k(k,:)\|_2^2\le K_k,\qquad K_k>0,
  \qquad
  mB_k^2<\frac1{K_k}
\]
implies the same repository-budgeted
\(\mathrm{LSQRSolveBackwardError}\) certificate.
\end{corollary}


\begin{corollary}[Stored-prefix-span local inverse row-budget QR certificate]
In the preceding row-budget setting, the separate prefix-span hypothesis can
be replaced by a previous-block local inverse hypothesis.  For every pivot
\(k<n\), assume that the transposed previous leading block \(T_k\) has a local
left inverse \(C^{\mathrm{prev}}_k\), and assume the current leading block
\(S_k\) has the local left inverse \(C_k\) used in the row-budget condition.
Then the stored panel recurrence supplies the lower-zero shape of the previous
columns, and this shape plus \(C^{\mathrm{prev}}_kT_k=I\) implies the
prefix-span invariant.  Hence the repository-budgeted
\[
  \mathrm{LSQRSolveBackwardError}
\]
certificate follows from the row-norm budget
\[
  \|C_k(k,:)\|_2^2\le K_k,\qquad K_k>0,
  \qquad mB_k^2<\frac1{K_k},
\]
sign choice, \(\gamma\)-validity, and the visible compact-smallness
hypothesis, without assuming prefix-span separately.
\end{corollary}


\begin{corollary}[Least-squares QR certificate from prefix-span and local inverse Frobenius budget]
In the same stored trailing Householder QR loop, suppose each leading block has
a local left inverse \(C_k\), and replace the row-norm condition above by
\[
  \|C_k\|_F^2\le K_k,\qquad K_k>0,
  \qquad
  mB_k^2 < \frac1{K_k}.
\]
Under the same prefix-span, Householder, compact-update, \(\gamma\)-validity,
and final Gram/RHS norm-budget hypotheses, the computed vector
\[
  \widehat x=\mathrm{fl\_backSub}(\widehat R,\widehat c)
\]
satisfies
\[
  \mathrm{LSQRSolveBackwardError}(A^TA,A^Tb,\widehat x;c_G,c_g).
\]
\end{corollary}


\begin{corollary}[Least-squares QR certificate from prefix-span and local inverse infinity-norm budget]
In the same stored trailing Householder QR loop, suppose each leading block has
a local left inverse \(C_k\), and replace the Frobenius condition above by
\[
  (k+1)\|C_k\|_\infty^2\le K_k,\qquad K_k>0,
  \qquad
  mB_k^2 < \frac1{K_k}.
\]
Under the same prefix-span, Householder, compact-update, \(\gamma\)-validity,
and final Gram/RHS norm-budget hypotheses, the computed vector
\[
  \widehat x=\mathrm{fl\_backSub}(\widehat R,\widehat c)
\]
satisfies
\[
  \mathrm{LSQRSolveBackwardError}(A^TA,A^Tb,\widehat x;c_G,c_g).
\]
\end{corollary}


\begin{corollary}[Repository-budget inverse-norm QR certificates]
In the Frobenius and infinity inverse-budget corollaries above, the final
Gram/RHS radii may be chosen as
\(\mathrm{qrSolveFinalGramBudget}\) and
\(\mathrm{qrSolveFinalRhsBudget}\).  If the step constant is chosen as
\[
  c_{\mathrm{step}}
    = \mathrm{storedQRCompactSequenceRelativeBudget}
        (A_0,A_1,\ldots,A_n;b_0,b_1,\ldots,b_n;\alpha),
\]
then the stored compact-budget lemmas also supply the columnwise panel and RHS
compact-update domination hypotheses.  Thus either
\[
  \|C_k\|_F^2\le K_k
  \quad\text{or}\quad
  (k+1)\|C_k\|_\infty^2\le K_k
\]
together with \(K_k>0\), prefix-span, sign choice, and
\(mB_k^2<1/K_k\) gives the same repository-budgeted
\(\mathrm{LSQRSolveBackwardError}\) certificate.
\end{corollary}


\begin{corollary}[Repository-budget inverse-norm certificates without separate prefix-span]
For the Frobenius and infinity branches in the previous corollary, the
separate prefix-span hypothesis can also be removed when the stored panel
recurrence is available and each previous transposed leading block has a local
left inverse.  Under those previous-block local inverse hypotheses, the same
Frobenius or infinity inverse-budget assumptions imply the repository-budgeted
\(\mathrm{LSQRSolveBackwardError}\) certificate.
\end{corollary}


\begin{corollary}[Signed-alpha local-inverse row certificate]
In the stored-prefix-span local-inverse row certificate, the separate
assumptions
\[
  \alpha_k^2=\|(\widehat A_k)_{k:m,k}\|_2^2,
  \qquad
  \alpha_k(\widehat A_k)_{kk}\le0
\]
may be replaced by the source Householder choice
\[
  \alpha_k =
  \mathrm{signedHouseholderAlpha}
    \bigl(\|(\widehat A_k)_{k:m,k}\|_2,(\widehat A_k)_{kk}\bigr).
\]
Under the same stored recurrence, previous/current local left-inverse,
row-norm budget, and compact-smallness hypotheses, the repository-budgeted
\(\mathrm{LSQRSolveBackwardError}\) certificate follows.
\end{corollary}


\begin{corollary}[Signed-alpha local-inverse inverse-norm certificates]
In the stored-prefix-span Frobenius and infinity local-inverse certificates,
the same source Householder choice
\[
  \alpha_k =
  \mathrm{signedHouseholderAlpha}
    \bigl(\|(\widehat A_k)_{k:m,k}\|_2,(\widehat A_k)_{kk}\bigr)
\]
may replace the separate squared-alpha and sign-choice assumptions.  Thus
either inverse-norm budget
\[
  \|C_k\|_F^2\le K_k
  \quad\text{or}\quad
  (k+1)\|C_k\|_\infty^2\le K_k
\]
together with the previous/current local left-inverse hypotheses,
\(K_k>0\), and \(mB_k^2<1/K_k\), gives the repository-budgeted
\(\mathrm{LSQRSolveBackwardError}\) certificate without assuming prefix-span,
the squared-alpha identity, or the sign inequality separately.
\end{corollary}


\begin{corollary}[Signed-alpha determinant local-inverse certificates]
Let \(P_k\) be the previous \(k\times k\) transposed leading block and let
\(S_k\) be the current \((k+1)\times(k+1)\) leading block of the stored
Householder panel at step \(k\).  In the signed-alpha stored-prefix-span
certificates above, the previous/current local left-inverse witnesses may be
replaced by
\[
  \det P_k\ne0,\qquad \det S_k\ne0,
\]
using the repository nonsingular inverses
\[
  C^{\mathrm{prev}}_k=\mathrm{nonsingInv}(P_k),
  \qquad
  C_k=\mathrm{nonsingInv}(S_k).
\]
Thus any one of the visible inverse-budget assumptions
\[
  \|C_k(k,:)\|_2^2\le K_k,\qquad
  \|C_k\|_F^2\le K_k,\qquad
  (k+1)\|C_k\|_\infty^2\le K_k
\]
together with \(K_k>0\), the signed Householder alpha rule, the stored panel
recurrence, and \(mB_k^2<1/K_k\), gives the repository-budgeted
\(\mathrm{LSQRSolveBackwardError}\) certificate.  The theorem does not derive
the determinant facts, inverse-budget inequality, or compact-smallness
inequality from the computed loop.
\end{corollary}


\begin{corollary}[Signed-alpha determinant certificate from \(\kappa_\infty\) budgets]
In the infinity-norm branch of the preceding corollary, the direct inverse
budget may be replaced by a visible local condition-number budget.  Suppose
\[
  \kappa_\infty(S_k,\mathrm{nonsingInv}(S_k))\le \kappa_k,
  \qquad
  (k+1)
  \left(
    \frac{\kappa_k}{\|S_k\|_\infty}
  \right)^2
  \le K_k .
\]
Together with \(\det P_k\ne0\), \(\det S_k\ne0\), \(K_k>0\), the signed
Householder alpha rule, the stored panel recurrence, and \(mB_k^2<1/K_k\),
these hypotheses imply the same repository-budgeted
\(\mathrm{LSQRSolveBackwardError}\) certificate.
\end{corollary}


\begin{corollary}[Signed-alpha triangular leading-block certificate]
In the previous corollary, the determinant hypotheses may be replaced by a
visible triangular leading-block invariant.  Suppose the stored panel at each
step is upper triangular on its leading part and all displayed leading diagonal
entries are nonzero:
\[
  (\widehat A_k)_{ij}=0\quad(j<i),
  \qquad
  (\widehat A_k)_{rr}\ne0
  \quad(0\le r\le k).
\]
Then the previous transposed leading block \(P_k\) and current leading block
\(S_k\) are nonsingular, and the signed-alpha \(\kappa_\infty\) certificate
above applies under the same visible condition-number, self-norm, and
compact-smallness hypotheses.
\end{corollary}


\begin{corollary}[Least-squares QR certificate from nonsingular local inverse-norm budgets]
In the previous corollary, the explicit local left-inverse witness \(C_k\) can
be replaced by a determinant hypothesis and the repository nonsingular inverse.
Let \(S_k\) be the local \((k+1)\times(k+1)\) leading block.  Suppose
\[
  \det S_k\ne0,\qquad
  (k+1)\|\mathrm{nonsingInv}(S_k)\|_\infty^2\le K_k,\qquad
  K_k>0,\qquad mB_k^2<\frac1{K_k}.
\]
Under the same prefix-span, Householder sign, compact-update,
\(\gamma\)-validity, and final Gram/RHS norm-budget hypotheses as above, the
computed vector
\[
  \widehat x=\mathrm{fl\_backSub}(\widehat R,\widehat c)
\]
satisfies
\[
  \mathrm{LSQRSolveBackwardError}(A^TA,A^Tb,\widehat x;c_G,c_g).
\]
\end{corollary}


\begin{corollary}[Least-squares QR certificate from local \(\kappa_\infty\) budgets]
In the previous corollary, the inverse-\(\infty\) budget may be derived from a
condition-number route.  Let \(S_k\) be the local \((k+1)\times(k+1)\) leading
block.  Suppose
\[
  \det S_k\ne0,\qquad 0<\rho_k\le \|S_k\|_\infty,\qquad
  \kappa_\infty(S_k,\mathrm{nonsingInv}(S_k))\le \kappa_k,
\]
and suppose the visible squared budget satisfies
\[
  (k+1)\left(\frac{\kappa_k}{\rho_k}\right)^2\le K_k,\qquad
  K_k>0,\qquad mB_k^2<\frac1{K_k}.
\]
Under the same prefix-span, Householder sign, compact-update,
\(\gamma\)-validity, and final Gram/RHS norm-budget hypotheses as above, the
computed vector
\[
  \widehat x=\mathrm{fl\_backSub}(\widehat R,\widehat c)
\]
satisfies
\[
  \mathrm{LSQRSolveBackwardError}(A^TA,A^Tb,\widehat x;c_G,c_g).
\]
\end{corollary}


\begin{corollary}[Least-squares QR certificate from local \(\kappa_\infty\) budgets without an auxiliary \(\rho_k\)]
In the preceding condition-number corollary, the auxiliary lower-bound
parameter can be removed.  Let \(S_k\) be the local
\((k+1)\times(k+1)\) leading block.  Suppose
\[
  \det S_k\ne0,\qquad
  \kappa_\infty(S_k,\mathrm{nonsingInv}(S_k))\le \kappa_k,
\]
and suppose the visible squared budget satisfies
\[
  (k+1)\left(\frac{\kappa_k}{\|S_k\|_\infty}\right)^2\le K_k,\qquad
  K_k>0,\qquad mB_k^2<\frac1{K_k}.
\]
Under the same prefix-span, Householder sign, compact-update,
\(\gamma\)-validity, and final Gram/RHS norm-budget hypotheses as above, the
computed vector
\[
  \widehat x=\mathrm{fl\_backSub}(\widehat R,\widehat c)
\]
satisfies
\[
  \mathrm{LSQRSolveBackwardError}(A^TA,A^Tb,\widehat x;c_G,c_g).
\]
\end{corollary}


\begin{corollary}[Self-norm \(\kappa_\infty\) QR certificate without abstract prefix span]
In the preceding self-norm condition-number corollary, the abstract
prefix-span invariant may be replaced by determinant and shape hypotheses on
the completed leading block.  Let \(T_k\) be the transposed leading block of
the first \(k\) rows and previous \(k\) columns, and let \(S_k\) be the current
\((k+1)\times(k+1)\) leading block.  Suppose
\[
  \det T_k\ne0,\qquad
  A_k(i,j)=0\quad (i\ge k,\ j<k),
\]
and also
\[
  \det S_k\ne0,\qquad
  \kappa_\infty(S_k,\mathrm{nonsingInv}(S_k))\le \kappa_k,\qquad
  (k+1)\left(\frac{\kappa_k}{\|S_k\|_\infty}\right)^2\le K_k .
\]
Under the same Householder sign, compact-update, \(\gamma\)-validity, and
final Gram/RHS norm-budget hypotheses as above, with \(K_k>0\) and
\[
  mB_k^2<\frac1{K_k},
\]
the computed vector
\[
  \widehat x=\mathrm{fl\_backSub}(\widehat R,\widehat c)
\]
satisfies
\[
  \mathrm{LSQRSolveBackwardError}(A^TA,A^Tb,\widehat x;c_G,c_g).
\]
\end{corollary}


\begin{formalized}
\lean{Algorithms/QR/HouseholderQR.lean}
\leanfit{qrPrefixSupportSpannedByPreviousColumns_of_det_ne_zero_previousLeadingBlockTranspose}
\lean{Algorithms/LeastSquares/LSQRSolve.lean}
\leanfit{LSQRSolveBackwardError.of_stored_trailing_householder_sequence_topBlock_fl_backSub_gamma_bound_normBudget_of_leading_blocks_det_ne_zero_kappaInf_selfNorm_budget}
\end{formalized}

\begin{corollary}[Triangular self-norm \(\kappa_\infty\) QR certificate]
In the preceding corollary, the hypotheses
\[
  \det T_k\ne0,\qquad \det S_k\ne0,\qquad
  A_k(i,j)=0\quad(i\ge k,\ j<k)
\]
may be replaced by the following visible triangular-domain assumptions.  For
each pivot \(k\), suppose the current stored panel satisfies
\[
  A_k(i,j)=0\quad(j<i),
\]
and suppose the diagonal entries of the previous and current leading blocks are
nonzero.  Then the same self-norm condition-number hypotheses
\[
  \kappa_\infty(S_k,\mathrm{nonsingInv}(S_k))\le\kappa_k,\qquad
  (k+1)\left(\frac{\kappa_k}{\|S_k\|_\infty}\right)^2\le K_k,
  \qquad mB_k^2<\frac1{K_k}
\]
imply the same
\[
  \mathrm{LSQRSolveBackwardError}(A^TA,A^Tb,\widehat x;c_G,c_g)
\]
conclusion under the same Householder sign, compact-update, \(\gamma\)-validity,
and final Gram/RHS norm-budget hypotheses.
\end{corollary}


\begin{formalized}
\lean{Algorithms/LeastSquares/LSQRSolve.lean}
\leanfit{LSQRSolveBackwardError.of_stored_trailing_householder_sequence_topBlock_fl_backSub_gamma_bound_normBudget_of_triangular_leading_blocks_kappaInf_selfNorm_budget}
\end{formalized}

\begin{corollary}[Computed prefix-zero triangular self-norm QR certificate]
In the preceding triangular self-norm certificate, the broad intermediate-panel
triangular-shape hypothesis can be replaced by the stored-loop prefix
lower-zero invariant.  For the stored Householder panel recurrence,
\[
  A_{k+1}=\mathrm{fl\_householderStoredPanelStep}(A_k),
\]
the repository proves that after \(k\) steps all entries below the diagonal in
the first \(k\) completed columns are exactly zero:
\[
  j<k,\ j<i \quad\Longrightarrow\quad A_k(i,j)=0.
\]
Consequently, if the previous and current leading diagonal entries are
nonzero, and if the same local \(\kappa_\infty\), sign-choice,
compact-update, and final Gram/RHS budget hypotheses hold, then
\[
  \mathrm{LSQRSolveBackwardError}(A^TA,A^Tb,\widehat x;c_G,c_g)
\]
again follows.
\end{corollary}


\begin{formalized}
\lean{Algorithms/QR/HouseholderQR.lean}
\leanfit{fl_householderStoredPanel_sequence_prefix_lower_zero}
\leanfit{qrPreviousLeadingBlockTranspose_det_ne_zero_of_local_lower_triangular_diag_ne_zero}
\leanfit{qrLeadingBlock_det_ne_zero_of_local_upper_triangular_diag_ne_zero}
\lean{Algorithms/LeastSquares/LSQRSolve.lean}
\leanfit{LSQRSolveBackwardError.of_stored_trailing_householder_sequence_topBlock_fl_backSub_gamma_bound_normBudget_of_prefix_lower_zero_triangular_leading_blocks_kappaInf_selfNorm_budget}
\end{formalized}

\begin{corollary}[Signed-alpha computed prefix-zero QR certificate]
Define the signed Householder scalar by
\[
  \alpha(x_p,\rho)=
  \begin{cases}
    -\rho, & 0\le x_p,\\
     \rho, & x_p<0,
  \end{cases}
  \qquad
  \rho=\sqrt{\sum_{i=k}^m A_k(i,k)^2}.
\]
Then \(\alpha^2=\rho^2\) and
\(\alpha A_k(k,k)\le0\).  Hence, in the computed prefix-zero triangular
self-norm QR certificate above, the independent Householder squared-norm and
sign-choice hypotheses may be replaced by the explicit source convention
\[
  \alpha_k =
  \alpha\!\left(A_k(k,k),
    \sqrt{\sum_{i=k}^m A_k(i,k)^2}\right).
\]
The resulting least-squares QR backward-error certificate still keeps the
remaining assumptions visible: nonzero local leading diagonals, local
\(\kappa_\infty\) budgets, compact-update budgets, and final Gram/RHS solver
budgets.
\end{corollary}


\begin{formalized}
\lean{Algorithms/QR/HouseholderSpec.lean}
\leanfit{signedHouseholderAlpha}
\leanfit{signedHouseholderAlpha_sqrt_trailingNorm2Sq_sq}
\leanfit{signedHouseholderAlpha_sqrt_trailingNorm2Sq_mul_pivot_nonpos}
\lean{Algorithms/QR/HouseholderQR.lean}
\leanfit{fl_householderStoredTrailingPanel_sequence_diag_nonzero_of_signed_alpha_trailingNorm_pos_sqrt_budget}
\lean{Algorithms/LeastSquares/LSQRSolve.lean}
\leanfit{LSQRSolveBackwardError.of_stored_trailing_householder_sequence_topBlock_fl_backSub_gamma_bound_normBudget_of_signed_alpha_prefix_lower_zero_triangular_leading_blocks_kappaInf_selfNorm_budget}
\end{formalized}

\begin{corollary}[Prefix-local diagonal nonbreakdown from the stored loop]
In the signed-alpha stored trailing Householder QR loop above, suppose each
completed pivot step has positive trailing norm and a square-root
compact-update budget small enough to keep the written pivot nonzero.  Then,
after \(k\) stored steps, all previously written diagonal entries remain
nonzero:
\[
  A_k(r,r)\ne0,\qquad 0\le r<k .
\]
Consequently, in the computed prefix-zero triangular self-norm QR certificate,
the previous local diagonal hypotheses may be generated by the stored loop.
Only the current leading pivot \(A_k(k,k)\ne0\) remains as the visible
diagonal-domain condition for the current \((k+1)\times(k+1)\) principal
block.  The local \(\kappa_\infty\), compact-update, and final Gram/RHS
budgets remain explicit.
\end{corollary}


\begin{formalized}
\lean{Algorithms/QR/HouseholderQR.lean}
\leanfit{fl_householderStoredPanel_sequence_prefix_diag_nonzero_of_step_diag_nonzero}
\leanfit{fl_householderStoredTrailingPanel_sequence_prefix_diag_nonzero_of_signed_alpha_trailingNorm_pos_sqrt_budget}
\lean{Algorithms/LeastSquares/LSQRSolve.lean}
\leanfit{LSQRSolveBackwardError.of_stored_trailing_householder_sequence_topBlock_fl_backSub_gamma_bound_normBudget_of_signed_alpha_prefix_lower_zero_pivot_nonzero_kappaInf_selfNorm_budget}
\end{formalized}

\begin{corollary}[Explicit final Gram/RHS radii for the prefix-local QR certificate]
Let \(\widehat R=A_n(1{:}n,1{:}n)\), and define the final rectangular data and
right-hand-side perturbation radii
\[
  c_A^{\mathrm{QR}}
  =
  \big((1+c_{\mathrm{step}})^n-1\big)\|A\|_F
  +\gamma_n\|[\widehat R;0]\|_F,
  \qquad
  c_b^{\mathrm{QR}}
  =
  \big((1+c_{\mathrm{step}})^n-1\big)\|b\|_2 .
\]
The induced final Gram and right-hand-side radii are
\[
  c_G^{\mathrm{QR}}
  =
  \left\|
    \mathrm{rectLSGramPerturbationNormBudget}(A,c_A^{\mathrm{QR}})
  \right\|_F .
\]
\[
  c_g^{\mathrm{QR}}
  =
  \sum_{j=1}^n
    \mathrm{rectLSRhsPerturbationNormBudget}
      (A,b,c_A^{\mathrm{QR}},c_b^{\mathrm{QR}})_j .
\]
Under the hypotheses of the preceding prefix-local signed-alpha certificate,
the computed vector satisfies
\[
  \mathrm{LSQRSolveBackwardError}
    (A^TA,A^Tb,\widehat x;c_G^{\mathrm{QR}},c_g^{\mathrm{QR}}).
\]
Thus the final Gram/RHS budget-domination assumptions are no longer separate
side hypotheses in this route; they are instantiated by the displayed
repository radii.  The current-pivot nonzero condition, local
\(\kappa_\infty\) budgets, compact-update budgets, square-root pivot budgets,
and \(\gamma\)-validity remain explicit.
\end{corollary}


\begin{formalized}
\lean{Algorithms/LeastSquares/LSQRSolve.lean}
\leanfit{qrSolveFinalDataPerturbationBudget}
\leanfit{qrSolveFinalRhsPerturbationBudget}
\leanfit{qrSolveFinalGramBudget}
\leanfit{qrSolveFinalRhsBudget}
\leanfit{rectLSRhsPerturbationNormBudget_le_sum}
\leanfit{LSQRSolveBackwardError.of_stored_trailing_householder_sequence_topBlock_fl_backSub_gamma_bound_explicitNormBudget_of_signed_alpha_prefix_lower_zero_pivot_nonzero_kappaInf_selfNorm_budget}
\end{formalized}

\begin{corollary}[Explicit compact-update and final radii for the prefix-local QR certificate]
For each stored QR step \(k\), let \(d_k^{\mathrm{cmp}}\) be the
repository-defined compact-panel relative budget
\[
  d_k^{\mathrm{cmp}}
  =
  \mathrm{householderCompactPanelRelativeBudget}
  (\widehat A_k,\widehat b_k,v_k,\beta_k).
\]
Set
\[
  c_{\mathrm{cmp}}^{\mathrm{QR}}=\sum_{k=0}^{n-1} d_k^{\mathrm{cmp}} .
\]
Then every stored panel-column and right-hand-side compact-update budget is
dominated by \(c_{\mathrm{cmp}}^{\mathrm{QR}}\) times the corresponding input
vector norm.  Consequently, in the preceding prefix-local signed-alpha
certificate one may take
\[
  c_{\mathrm{step}}=c_{\mathrm{cmp}}^{\mathrm{QR}}
\]
and the final radii \(c_G^{\mathrm{QR}},c_g^{\mathrm{QR}}\) from
Corollary~226.  Under the same current-pivot nonzero, local
\(\kappa_\infty\), square-root/diagonal budget, and \(\gamma\)-validity
hypotheses, the computed vector satisfies
\[
  \mathrm{LSQRSolveBackwardError}
    (A^TA,A^Tb,\widehat x;c_G^{\mathrm{QR}},c_g^{\mathrm{QR}}).
\]
Thus neither the compact-update domination constant nor the final Gram/RHS
radii are separate user-supplied side hypotheses in this route.
\end{corollary}


\begin{formalized}
\lean{Algorithms/QR/HouseholderSpec.lean}
\leanfit{householderTrailingNorm2Sq_pos_of_nonneg_budget_lt_sqrt}
\lean{Algorithms/QR/HouseholderApply.lean}
\leanfit{householderCompactRelativeBudget}
\leanfit{householderCompactPanelRelativeBudget}
\lean{Algorithms/QR/HouseholderQR.lean}
\leanfit{storedQRCompactStepRelativeBudget}
\leanfit{storedQRCompactSequenceRelativeBudget}
\leanfit{storedQRCompactSequenceRelativeBudget_column_bound}
\leanfit{storedQRCompactSequenceRelativeBudget_rhs_bound}
\lean{Algorithms/LeastSquares/LSQRSolve.lean}
\leanfit{LSQRSolveBackwardError.of_stored_trailing_householder_sequence_topBlock_fl_backSub_gamma_bound_explicitCompactBudget_of_signed_alpha_prefix_lower_zero_pivot_nonzero_kappaInf_selfNorm_budget}
\end{formalized}

\begin{corollary}[Explicit compact QR certificate from norm-square pivot margins]
In the setting of Corollary~227, replace the visible square-root pivot-budget
condition by the dimensioned margin
\[
  mB_k^2<\|A_k(k{:}m,k)\|_2^2,\qquad 0\le B_k ,
\]
where \(B_k\) is the repository compact diagonal-update budget for pivot \(k\).
Under the same current-pivot nonzero, local \(\kappa_\infty\), dual/conditioning
budget, and \(\gamma\)-validity hypotheses, the computed vector satisfies
\[
  \mathrm{LSQRSolveBackwardError}
    (A^TA,A^Tb,\widehat x;c_G^{\mathrm{QR}},c_g^{\mathrm{QR}}).
\]
\end{corollary}


\begin{formalized}
\lean{Algorithms/QR/HouseholderSpec.lean}
\leanfit{budget_lt_sqrt_householderTrailingNorm2Sq_of_dim_mul_budget_sq_lt_trailingNorm2Sq}
\lean{Algorithms/LeastSquares/LSQRSolve.lean}
\leanfit{LSQRSolveBackwardError.of_stored_trailing_householder_sequence_topBlock_fl_backSub_gamma_bound_explicitCompactBudget_of_signed_alpha_prefix_lower_zero_pivot_nonzero_kappaInf_selfNorm_normSqBudget}
\end{formalized}

\begin{corollary}[Least-squares QR certificate from diagonal-dominant local inverse budgets]
In the same stored trailing Householder QR loop, let \(S_k\) be the
\((k+1)\times(k+1)\) local leading block at pivot \(k\), and let \(C_k\) be a
full inverse of \(S_k\).  Suppose \(S_k\) is upper triangular and diagonally
dominant in the Higham sense,
\[
  S_k(i,j)=0\quad(j<i),\qquad
  |S_k(i,j)|\le |S_k(i,i)|\quad(i<j),
\]
with nonzero diagonal.  If the visible budget satisfies
\[
  (k+1)\left(2^k\,
  \frac{1}{\min_{0\le r\le k}|S_k(r,r)|}\right)^2 \le K_k,
  \qquad K_k>0,\qquad
  mB_k^2<\frac1{K_k},
\]
then, under the same prefix-span, Householder, compact-update,
\(\gamma\)-validity, and final Gram/RHS norm-budget hypotheses as above, the
computed vector
\[
  \widehat x=\mathrm{fl\_backSub}(\widehat R,\widehat c)
\]
satisfies
\[
  \mathrm{LSQRSolveBackwardError}(A^TA,A^Tb,\widehat x;c_G,c_g).
\]
\end{corollary}


\begin{corollary}[Diagonal-dominant QR certificate from determinant]
In the previous corollary, the explicit full inverse witness \(C_k\) can be
removed.  It is enough to assume
\[
  \det S_k\ne 0
\]
for every local leading block, while keeping the same upper-triangular
diagonal-dominance and diagonal-minimum budget.  Lean instantiates
\(C_k=S_k^{-1}\) using the repository nonsingular inverse
\(\mathrm{nonsingInv}\), proves it is a two-sided inverse from
\(\det S_k\ne0\), and obtains the same
\[
  \mathrm{LSQRSolveBackwardError}(A^TA,A^Tb,\widehat x;c_G,c_g)
\]
conclusion.
\end{corollary}


\begin{corollary}[Least-squares QR certificate from nonsingular local leading blocks]
In the stored trailing Householder QR loop, assume the same rounded matrix and
right-hand-side recurrences used in the stored-loop solver certificate.  For
each pivot \(k<n\), suppose that the previous transposed leading block
\(T_k\) and the current leading block \(S_k\) are nonsingular,
\[
  \det T_k\ne0,\qquad \det S_k\ne0,
\]
that the previous columns have the QR lower-zero shape, that the Householder
sign choice satisfies
\[
  \alpha_k A_k(k,k)\le 0,
\]
and that the compact rounded diagonal update budget satisfies
\[
  B_k<\sqrt{\|A_k(k{:}m,k)\|_2^2}.
\]
Assume also the columnwise compact-update budget domination for the matrix
panel and right-hand side, the \(\gamma\)-validity hypotheses, and the final
Gram/RHS norm-budget bounds.  Then the computed back-substitution vector
\[
  \widehat x=\mathrm{fl\_backSub}(\widehat R,\widehat c)
\]
satisfies the repository's rectangular QR least-squares backward-error
certificate,
\[
  \mathrm{LSQRSolveBackwardError}(A^TA,A^Tb,\widehat x;c_G,c_g).
\]
\end{corollary}


\begin{corollary}[Least-squares QR certificate from triangular leading blocks]
In the stored trailing Householder QR loop, assume the same rounded matrix and
right-hand-side recurrences, Householder normalizations, sign choices, compact
panel/RHS budget domination, square-root pivot budgets, and final Gram/RHS
norm-budget bounds as in the previous corollary.  Instead of assuming
\(\det T_k\ne0\) and \(\det S_k\ne0\) directly, assume that at every pivot
\(k<n\) the local panel has the visible upper-triangular shape used by the
leading blocks and that the previous and current leading principal diagonals
are nonzero:
\[
  A_k(i,j)=0 \quad (j<i),\qquad
  A_k(r,r)\ne0
\]
for the relevant previous block \(T_k\) and current block \(S_k\).
Then the computed back-substitution vector again satisfies
\[
  \mathrm{LSQRSolveBackwardError}(A^TA,A^Tb,\widehat x;c_G,c_g).
\]
\end{corollary}


\paragraph{Role in the bottleneck.}
This closes both the exact trailing-reflector shape theorem and the stored
rounded loop shape theorem, and it closes the stored-step perturbation contract
under explicit preservation, pivot-zeroing, RHS-prefix, and budget-domination
hypotheses.  It also closes the one-step trailing-reflector discharge of the
exact preservation and zeroing hypotheses.  The multi-step stored trailing
sequence theorem and the stored-loop solver-spec handoff into
\(\mathrm{LSQRSolveBackwardError}\) are also closed.  The diagonal condition is
narrowed to the concrete per-pivot budget inequalities \(B_k<|\alpha_k|\), and
the Householder denominator condition is narrowed to the scalar pivot
condition \(A_k(k,k)\ne\alpha_k\), and then to the standard sign-choice
conditions \(\|A_k(k{:}m,k)\|_2^2>0\) and
\(\alpha_kA_k(k,k)\le0\).  The positive-norm obligation is now reduced further:
a left inverse for the transposed leading block gives the leading-block
coefficient witness, while a left inverse for the actual leading
\((k+1)\times(k+1)\) block pads to the ambient leading-column dual witness.
The coefficient witness plus the QR lower-zero shape gives the prefix-span
invariant; the leading-column dual witness gives column independence;
prefix-span plus column independence gives a nonzero active trailing-column
entry; and that entry gives the positive norm.  One stored-loop wrapper
consumes the concrete coefficient/left-inverse witnesses directly, and a
second wrapper consumes the two local \(\mathrm{IsLeftInverse}\) block
witnesses directly.  A third wrapper now consumes nonzero determinants for
those two leading blocks and generates the local inverse witnesses internally.
The new solver-facing wrapper feeds those determinant facts and the
square-root pivot budget all the way into
\(\mathrm{LSQRSolveBackwardError}\).  The triangular-leading-block wrapper now
pushes the visible triangular-principal-minor route through that same solver
certificate.  The prefix-span active-entry solver wrapper also pushes a
different visible domain route through \(\mathrm{LSQRSolveBackwardError}\):
it replaces the square-root budget by \(B_k<|A_k(i_k,k)|\) for an active
trailing entry.  The diagonal-dominant triangular inverse wrapper now replaces
the abstract inverse-\(\infty\)-norm budget by Higham's explicit
\(2^k/\min_r |S_k(r,r)|\) estimate when each local leading block is
diagonally dominant and a full inverse witness is supplied.  The remaining
source-faithful
implementation work is either to keep the triangular/nonzero-diagonal
and diagonal-dominance conditions visible as domain assumptions, or to prove
the needed determinant/rank facts for the computed leading blocks and derive
the usable budget lower bound from a formal full-rank, nonbreakdown, or
conditioning assumption for the computed loop.

\begin{definition}[Common rounded Householder panel/update contract]
For a rectangular matrix panel \(A\in\mathbb R^{m\times n}\), a right-hand side
\(b\in\mathbb R^m\), an exact reflector \(P\in\mathbb R^{m\times m}\), and
rounded outputs \(\widehat A,\widehat b\), the contract
\(\mathrm{HouseholderPanelAppError}(P,A,\widehat A,b,\widehat b;c)\) means:
\[
  P^TP=I
\]
and there is one matrix \(\Delta P\) such that
\[
  \|\Delta P\|_F\le c,\qquad
  \widehat A=(P+\Delta P)A,\qquad
  \widehat b=(P+\Delta P)b.
\]
\end{definition}

\begin{theorem}[Common panel contracts accumulate with one \(Q\)]
Let \(A_0=A\), \(b_0=b\).  Suppose that for every step \(k<r\), the rounded
Householder panel update from \((A_k,b_k)\) to \((A_{k+1},b_{k+1})\) satisfies
\[
  \mathrm{HouseholderPanelAppError}
  (P_k,A_k,A_{k+1},b_k,b_{k+1};c),
  \qquad c\ge0.
\]
Then there are a common orthogonal matrix \(Q\), a matrix perturbation
\(\Delta A\), and a vector perturbation \(\Delta b\) such that
\[
  A_r = Q^T(A+\Delta A),\qquad
  b_r = Q^T(b+\Delta b),
\]
and
\[
  \|\Delta A\|_F
  \le ((1+c)^r-1)\|A\|_F,\qquad
  \|\Delta b\|_2
  \le ((1+c)^r-1)\|b\|_2.
\]
\end{theorem}


\begin{formalized}
\lean{Algorithms/QR/HouseholderSpec.lean}
\leanfit{HouseholderPanelAppError}
\lean{Algorithms/QR/HouseholderQR.lean}
\leanfit{householderPanelAppError_rect_orthogonal_matrix_vector_sequence_geometric}
\end{formalized}

\paragraph{Role in the bottleneck.}
This is a strong route-A adapter, but it is not the source-faithful standard
Householder QR route.  Higham's columnwise Householder QR analysis allows the
per-step perturbation matrix to depend on the panel column.  Therefore the
active route has been corrected to the following columnwise theorem.

\begin{definition}[Columnwise rounded Householder panel/update contract]
For \(A\in\mathbb R^{m\times n}\), \(b\in\mathbb R^m\), exact reflector
\(P\), and rounded outputs \(\widehat A,\widehat b\),
\(\mathrm{HouseholderColumnwisePanelAppError}
(P,A,\widehat A,b,\widehat b;c)\) means:
\[
  P^TP=I,
\]
and for every panel column \(j\) there is a perturbation matrix
\(\Delta P_j\) such that
\[
  \|\Delta P_j\|_F\le c,\qquad
  \widehat A(:,j)=(P+\Delta P_j)A(:,j),
\]
while the right-hand side has its own perturbation matrix \(\Delta P_b\)
satisfying
\[
  \|\Delta P_b\|_F\le c,\qquad
  \widehat b=(P+\Delta P_b)b.
\]
\end{definition}

\begin{theorem}[Columnwise panel contracts accumulate with one \(Q\)]
Let \(A_0=A\), \(b_0=b\).  Suppose that for every step \(k<r\), the rounded
Householder panel update from \((A_k,b_k)\) to \((A_{k+1},b_{k+1})\) satisfies
\[
  \mathrm{HouseholderColumnwisePanelAppError}
  (P_k,A_k,A_{k+1},b_k,b_{k+1};c),
  \qquad c\ge0.
\]
Then there are an orthogonal matrix \(Q\), a matrix perturbation \(\Delta A\),
and a vector perturbation \(\Delta b\) such that
\[
  A_r = Q^T(A+\Delta A),\qquad
  b_r = Q^T(b+\Delta b),
\]
with columnwise and right-hand-side radii
\[
  \|\Delta A(:,j)\|_2
  \le ((1+c)^r-1)\|A(:,j)\|_2
  \quad\text{for all }j,
\]
\[
  \|\Delta b\|_2
  \le ((1+c)^r-1)\|b\|_2.
\]
\end{theorem}


\begin{formalized}
\lean{Algorithms/QR/HouseholderSpec.lean}
\leanfit{HouseholderColumnwisePanelAppError}
\leanfit{HouseholderColumnwisePanelAppError.of_vector_applications}
\lean{Algorithms/QR/HouseholderQR.lean}
\leanfit{orthogonal_vector_sequence_one_step_fixedQ}
\leanfit{rect_orthogonal_columnwise_vector_sequence_geometric}
\leanfit{householderColumnwisePanelAppError_rect_orthogonal_columnwise_vector_sequence_geometric}
\end{formalized}

\begin{lemma}[Rank-one norm bridge for vector residuals]
Let \(e,b\in\mathbb R^m\) and suppose \(\|b\|_2\ne0\).  Then
\[
  \left\|\frac{e b^T}{\|b\|_2^2}\right\|_F
  =
  \frac{\|e\|_2}{\|b\|_2}.
\]
The proof first factors \(\|xy^T\|_F=\|x\|_2\|y\|_2\), then applies the scalar
homogeneity of the Frobenius norm.
\end{lemma}

\begin{theorem}[Forward error gives a Householder backward-error contract]
Let \(P\in\mathbb R^{m\times m}\) be orthogonal, let \(b,\widehat y\in\mathbb R^m\),
and let \(c\ge0\).  If
\[
  \|\widehat y-Pb\|_2\le c\|b\|_2,
\]
then \(\mathrm{HouseholderAppError}(P,b,\widehat y;c)\) holds.  When
\(b\ne0\), the perturbation witness is
\[
  \Delta P = \frac{(\widehat y-Pb)b^T}{\|b\|_2^2}.
\]
When \(b=0\), the forward-error hypothesis forces \(\widehat y-Pb=0\), and
the witness is \(\Delta P=0\).  Thus the Lean theorem has no hidden nonzero
input hypothesis.
\end{theorem}

\begin{corollary}[Columnwise panel contract from forward-error bounds]
Let \(P\) be orthogonal, \(c\ge0\), and suppose each rounded panel column and
the rounded right-hand side satisfy
\[
  \|\widehat A(:,j)-P A(:,j)\|_2\le c\|A(:,j)\|_2,
  \qquad
  \|\widehat b-Pb\|_2\le c\|b\|_2.
\]
Then
\[
  \mathrm{HouseholderColumnwisePanelAppError}
  (P,A,\widehat A,b,\widehat b;c)
\]
holds.  This is the adapter that turns future low-level rounded
dot/scale/subtract forward-error estimates into the source-faithful columnwise
QR accumulation theorem above.
\end{corollary}

\begin{theorem}[Explicit-matrix rounded Householder application]
Let \(P\in\mathbb R^{m\times m}\), \(b\in\mathbb R^m\), and let
\(\widehat y=\operatorname{fl\_matVec}(P,b)\).  Under the usual
\(\gamma_m\)-validity condition for the floating-point model,
\[
  \|\widehat y-Pb\|_2
  \le
  \gamma_m\,\Frob{\operatorname{abs}(P)}\,\|b\|_2 .
\]
Equivalently, this is the normwise consequence of the repository's existing
componentwise matrix-vector theorem
\[
  |\widehat y_i-(Pb)_i|
  \le
  \gamma_m \sum_j |P_{ij}|\,|b_j|.
\]
The proof then uses
\[
  \Norm{\operatorname{abs}(P)\operatorname{abs}(b)}
  \le
  \Frob{\operatorname{abs}(P)}\Norm{\operatorname{abs}(b)}
  =
  \Frob{\operatorname{abs}(P)}\Norm{b}.
\]
\end{theorem}

\begin{corollary}[Explicit-matrix route to the columnwise Householder contract]
If \(P\) is orthogonal and every vector application in a panel is computed by
the explicit-matrix routine above, then
\[
  \mathrm{HouseholderAppError}
  \bigl(P,b,\operatorname{fl\_matVec}(P,b);
        \gamma_m\Frob{\operatorname{abs}(P)}\bigr)
\]
holds, and the corresponding columnwise panel computation satisfies
\[
  \mathrm{HouseholderColumnwisePanelAppError}
  \bigl(P,A,\widehat A,b,\widehat b;
        \gamma_m\Frob{\operatorname{abs}(P)}\bigr).
\]
This closes a concrete formed-reflector implementation route.  It does not
claim to be the compact Householder update that computes
\(b-\beta v(v^Tb)\) by rounded dot, scale, and subtract operations.
\end{corollary}

\begin{definition}[Compact Householder arithmetic budget]
For the compact update, define
\[
  \sigma=v^Tb,\qquad
  S=\sum_j |v_j|\,|b_j|,\qquad
  D=\gamma_m S.
\]
Let
\[
  B_\sigma=S+D,\qquad
  B_\tau=|\beta|\,B_\sigma(1+u),\qquad
  B_{z,i}=B_\tau |v_i|(1+u).
\]
The deterministic component budget is
\[
  E_i =
  u B_\tau |v_i|
  +u|\beta|B_\sigma |v_i|
  +|\beta|\,|v_i|D
  +u\bigl(|b_i|+B_{z,i}\bigr).
\]
\end{definition}

\begin{theorem}[Compact rounded Householder dot--scale--subtract update]
Let
\[
  \widehat\sigma=\operatorname{fl\_dotProduct}(v,b),\qquad
  \widehat\tau=\operatorname{fl\_mul}(\beta,\widehat\sigma),
\]
and let the computed compact Householder application be
\[
  \widehat y_i
  =
  \operatorname{fl\_sub}
  \bigl(b_i,\operatorname{fl\_mul}(\widehat\tau,v_i)\bigr).
\]
If \(P=I-\beta vv^T\) and the \(\gamma_m\)-validity condition holds, then
\[
  Pb=b-\beta v(v^Tb)
\]
and, componentwise,
\[
  |\widehat y_i-(Pb)_i| \le E_i .
\]
Consequently,
\[
  \|\widehat y-Pb\|_2 \le \|(E_i)_{i=1}^m\|_2 .
\]
The proof reuses the local dot-product theorem
\[
  |\widehat\sigma-\sigma|\le \gamma_m\sum_j |v_j||b_j|
\]
and the primitive floating-point model for the two multiplications and final
subtraction.
\end{theorem}

\begin{corollary}[Compact route to the columnwise Householder contract]
Suppose \(P=I-\beta vv^T\) is orthogonal.  If, for a vector \(b\), the visible
deterministic budget satisfies
\[
  \|(E_i(b))_{i=1}^m\|_2\le c\|b\|_2,
\]
then the compact rounded update satisfies
\[
  \mathrm{HouseholderAppError}(P,b,\widehat y;c).
\]
If the same budget-domination condition holds for every panel column and for
the right-hand side, the compact columnwise panel routine satisfies
\[
  \mathrm{HouseholderColumnwisePanelAppError}
  (P,A,\widehat A,b,\widehat b;c).
\]
This closes the low-level compact vector-application primitive.  The
budget-domination condition is an explicit deterministic side condition, not a
probabilistic event or an assumed stability theorem.
\end{corollary}

\begin{corollary}[Compact Householder sequence accumulation]
Let \(P_k=I-\beta_k v_kv_k^T\) be a supplied sequence of orthogonal reflectors.
For readability, write
\[
  \begin{aligned}
  C_A(v,\beta,A)
  &:=
  \operatorname{fl\_householderApplyCompactPanel}(v,\beta,A),\\
  C_b(v,\beta,b)
  &:=
  \operatorname{fl\_householderApplyCompact}(v,\beta,b).
  \end{aligned}
\]
Assume each rounded panel update and right-hand-side update is computed by these
compact routines:
\[
  \widehat A_{k+1}=C_A(v_k,\beta_k,\widehat A_k),
  \qquad
  \widehat b_{k+1}=C_b(v_k,\beta_k,\widehat b_k),
\]
and assume the explicit compact budget is dominated by the same \(c\) for every
panel column and for the right-hand side.  Then the existing columnwise
geometric theorem gives \(Q,\Delta A,\Delta b\) such that
\[
  \widehat A_r
  =
  Q^T(A+\Delta A),
  \qquad
  \widehat b_r
  =
  Q^T(b+\Delta b),
\]
with \(Q\) orthogonal and
\[
  \|\Delta A_{:j}\|_2
  \le ((1+c)^r-1)\|A_{:j}\|_2,
  \qquad
  \|\Delta b\|_2
  \le ((1+c)^r-1)\|b\|_2 .
\]
This is still a supplied-sequence theorem: the concrete QR loop must prove the
reflector construction, transformed \([R;0]\) shape, and top-RHS linkage.
\end{corollary}

\begin{corollary}[Compact Householder sequence to the LS QR solver specification]
Let \(\eta_r=(1+c)^r-1\).  Under the hypotheses of Corollary~184, assume in
addition that the final transformed data have the QR least-squares shape
\[
  \widehat A_r=
  \begin{bmatrix}R\\0\end{bmatrix},
  \qquad
  (\widehat b_r)_{1:n}=c_{\rm top},
\]
that \(R\) is upper triangular with nonzero diagonal, and that the induced
Gram/RHS perturbation norm budgets are bounded by \(c_G,c_g\) after adding the
rounded triangular-solve top-block radius
\[
  \eta_r\|A\|_F+\gamma_n\left\|\begin{bmatrix}R\\0\end{bmatrix}\right\|_F .
\]
Then the computed vector
\[
  \widehat x=\operatorname{fl\_backSub}(R,c_{\rm top})
\]
satisfies the repository's QR least-squares backward-error specification
\[
  \mathrm{LSQRSolveBackwardError}
  (A^TA,A^Tb,\widehat x;c_G,c_g).
\]
The proof first converts the columnwise bounds
\(\|\Delta A_{:j}\|_2\le\eta_r\|A_{:j}\|_2\) into the Frobenius bound
\(\|\Delta A\|_F\le\eta_r\|A\|_F\), then invokes the existing top-block
triangular-solve pullback theorem.
\end{corollary}

\begin{formalized}
\lean{Analysis/MatrixAlgebra.lean}
\leanfit{vecNorm2Sq_abs}
\leanfit{vecNorm2_abs}
\leanfit{frobNormSqRect_eq_sum_vecNorm2Sq_cols}
\leanfit{frobNormRect_le_of_col_vecNorm2_le}
\leanfit{frobNormSq_rankOne}
\leanfit{frobNorm_rankOne}
\leanfit{frobNorm_rankOne_smul}
\leanfit{frobNorm_rankOne_div_vecNorm2Sq}
\lean{Algorithms/QR/HouseholderSpec.lean}
\leanfit{HouseholderAppError.of_forward_error}
\leanfit{HouseholderColumnwisePanelAppError.of_forward_errors}
\lean{Algorithms/QR/HouseholderApply.lean}
\leanfit{fl_householderApplyExplicit_forward_error_bound}
\leanfit{fl_householderApplyExplicit_HouseholderAppError}
\leanfit{fl_householderApplyExplicitPanel_HouseholderColumnwisePanelAppError}
\leanfit{matMulVec_householder_eq_compact}
\leanfit{fl_householderApplyCompact_componentwise_error_bound}
\leanfit{fl_householderApplyCompact_forward_error_bound}
\leanfit{fl_householderApplyCompact_HouseholderAppError_of_budget}
\leanfit{fl_householderApplyCompactPanel_HouseholderColumnwisePanelAppError_of_budget}
\lean{Algorithms/QR/HouseholderQR.lean}
\leanfit{fl_householderApplyCompactPanel_rect_orthogonal_columnwise_vector_sequence_geometric}
\lean{Algorithms/LeastSquares/LSQRSolve.lean}
\leanfit{LSQRSolveBackwardError.of_compact_householder_sequence_topBlock_fl_backSub_gamma_bound_normBudget}
\end{formalized}

\paragraph{Source and remaining obligation.}
This theorem matches the columnwise structure in Higham's QR notes,
Section~4.1, Lemmas~4.3--4.4 and Theorem~4.5.  The source is used only to
choose the route; the local Lean theorems above prove the accumulation step.
The front-loaded proof-source ledger also records the standard least-squares
QR route in Higham's \emph{Accuracy and Stability of Numerical Algorithms},
second edition, Chapter~19, especially Lemma~19.3 and Theorems~19.4--19.6.
That source is advisory here: it guides the remaining source-faithful route,
but no paper-level theorem is closed by citation alone.
The exact forward-error adapter, compact dot--scale--subtract primitive,
compact supplied-sequence adapter, Frobenius/top-block bridge, and exact
trailing-shape recurrence theorem are all closed.  The remaining obligation is
now narrower: assemble the rounded stored-\(R\) and transformed right-hand-side
loop around the compact update, and prove or expose the nonzero-pivot/rank
condition required by the triangular solve.

\begin{corollary}[High-probability literal rounded sampled-row solver from normal equations and Cholesky]
Use the notation and Bennett sample-budget assumptions of Corollary~157, but
compute the rounded sampled least-squares solution by the repository's
normal-equations/Cholesky pipeline.  For each sample trace \(\omega\), let
\[
  G_\omega=\operatorname{fl}(S_\omega A)^T\operatorname{fl}(S_\omega A),
  \qquad
  g_\omega=\operatorname{fl}(S_\omega A)^T\operatorname{fl}(S_\omega b),
\]
and let \(\widehat C_\omega,\widehat c_\omega,\widehat R_\omega\) be the
computed Gram matrix, computed right-hand side, and Cholesky factor.  Suppose
the local theorems
\[
  \lean{GramProductError},\qquad
  \lean{GramVecError},\qquad
  \lean{CholeskyBackwardError}
\]
hold for these objects, with the \(s\)-term Gram/vector product radius
\(\gamma_s\) and the Cholesky radius \(\gamma_{n+1}\).  Define the computed
solution by
\[
  \widehat y_\omega
    = \operatorname{flSolve}(\widehat R_\omega^T,\widehat c_\omega),
  \qquad
  \widehat x_\omega
    = \operatorname{flSolve}(\widehat R_\omega,\widehat y_\omega).
\]
The local normal-equations theorem supplies perturbations satisfying
\[
  (G_\omega+\Delta G_\omega)\widehat x_\omega
  = g_\omega+\Delta g_\omega,
\]
with the componentwise bounds
\[
  |(\Delta G_\omega)_{jk}|
  \le
  \gamma_s\,B^\mathrm{ATA}_{\omega,jk}
  +
  \bigl(\gamma_{n+1}+2\gamma_n+\gamma_n^2\bigr)
  \sum_{\ell=1}^n
    |(\widehat R_\omega)_{\ell j}|
    |(\widehat R_\omega)_{\ell k}|,
\]
and
\[
  |(\Delta g_\omega)_j|\le \gamma_s B^\mathrm{ATb}_{\omega,j}.
\]
Consequently the componentwise certificate
\[
  d_{\omega,i}
  =
  \sum_{j=1}^n |(H_\omega)_{ij}|
  \left(
    \sum_{k=1}^n B^G_{\omega,jk}|\widehat x_{\omega,k}|
    +
    B^g_{\omega,j}
  \right)
\]
is valid, where \(H_\omega=G_\omega^{-1}\) and \(B^G,B^g\) denote the two
displayed perturbation radii.  Hence, if the same support-event budget
condition holds with this \(d_\omega\), then
\[
  \mathbb P\!\left[
    \|A\widehat x_\omega-b\|_2^2
    \le (1+\eta)\|Ax_{\mathrm{opt}}-b\|_2^2
  \right]\ge 1-\delta .
\]
\end{corollary}


\begin{formalized}
\leanfit{normalEqCholeskyXHat}
\leanfit{normalEqCholeskyGramBound}
\leanfit{normalEqCholeskyRhsBound}
\leanfit{normalEqCholeskySolverDx}
\leanfit{normalEqCholeskySolverDx_nonneg}
\leanfit{normal_equations_cholesky_forward_error_certificate}
\leanfit{leverageTraceProbability_eventProb_fl_rowSampleLSObjective_le_one_add_eta_of_augmentedSpan_sample_budget_of_objective_error_and_normal_eq_cholesky_solver}
\end{formalized}

\begin{lemma}[Column representation implies residual representation]
Suppose the columns of \(A\) and the right-hand side \(b\) have coordinates in
the columns of \(U\):
\[
  A_{ij}=\sum_{\ell=1}^d U_{i\ell}c_{j\ell},
  \qquad
  b_i=\sum_{\ell=1}^d U_{i\ell}d_\ell .
\]
Then every residual has the induced coordinates
\[
  y_\ell(x)=\sum_{j=1}^n c_{j\ell}x_j-d_\ell,
  \qquad
  (Ax-b)_i=\sum_{\ell=1}^d U_{i\ell}y_\ell(x).
\]
If \(U^TU=I\), the canonical coordinates \(U^T(Ax-b)\) reconstruct the same
residual.
\end{lemma}

\begin{corollary}[Augmented-span leverage basis]
Let
\[
  \mathcal S=\operatorname{span}\{A_{*,1},\ldots,A_{*,n},b\}\subseteq \R^m,
  \qquad d=\dim \mathcal S .
\]
The repository constructs an orthonormal basis \(u_1,\ldots,u_d\) of
\(\mathcal S\) using Mathlib's finite-dimensional orthonormal-basis theorem,
and forms \(U=[u_1\ \cdots\ u_d]\).  It also constructs coordinates
\[
  A_{*,j}=\sum_{\ell=1}^d u_\ell c_{j\ell},
  \qquad
  b=\sum_{\ell=1}^d u_\ell d_\ell .
\]
Consequently every residual \(Ax-b\) lies in the column span of \(U\).  If
\(d>0\), equation~(6) is well-defined for this basis and the concrete
row-sampled least-squares theorem becomes
\[
  \mathbb P\!\left[
    \Norm{A\widehat x(\omega)-b}^2
    \le
    (1+\eta)\Norm{Ax_{\rm opt}-b}^2
  \right]
  \ge
  1-\frac{1}{s(\varepsilon/d)^2}.
\]
The positive-dimension hypothesis is the expected nondegeneracy condition for
the leverage denominator.  If \(\mathcal S=\{0\}\), the least-squares problem
is degenerate and the equation~(6) probabilities are not defined by this basis.
The source-aligned Bennett sample-budget theorem uses the sharper condition
\(d>1\), matching the explicit \((d-1)\) variance denominator in the
formalized equation~(7) concentration theorem.

The fully floating-point rounded-Gram transfer also specializes to this same
augmented-span basis.  In that theorem the preservation radius is
\[
  \varepsilon+
  \operatorname{rowSampleGramFullFpPerturbBudget}(\mathrm{fp},s,U),
\]
and the coordinate/original-objective hypotheses are discharged by the
augmented-span construction.  The theorem still explicitly requires the
rounded sketched objective to have the rounded-Gram quadratic representation;
it is an abstract objective transfer, not yet the theorem for the literal
rounded sampled \(A,b\) implementation above.
\end{corollary}

\begin{formalized}
\leanfit{augmentedDataVector}
\leanfit{augmentedDataSpan}
\leanfit{augmentedDataVector_mem_span}
\leanfit{augmentedSpanBasisMatrix}
\leanfit{augmentedSpanColumnCoords}
\leanfit{augmentedSpanRhsCoords}
\leanfit{hasOrthonormalColumns_augmentedSpanBasisMatrix}
\leanfit{columnsAndRhsInColumnSpace_augmentedSpanBasisMatrix}
\leanfit{residualsInColumnSpace_augmentedSpanBasisMatrix}
\leanfit{leverageTraceProbability_eventProb_rowSampleLSObjective_le_one_add_eta_of_augmentedSpan}
\leanfit{leverageTraceProbability_eventProb_rowSampleLSObjective_le_one_add_eta_of_augmentedSpan_sample_budget}
\leanfit{leverageTraceProbability_eventProb_fl_lsObjective_le_one_add_eta_of_augmentedSpan}
\leanfit{leverageTraceProbability_eventProb_fl_lsObjective_le_one_add_eta_of_augmentedSpan_sample_budget}
\end{formalized}

\begin{corollary}[Identity-basis fallback]
Taking \(U=I_m\) gives a completely concrete equation~(6) instance:
\[
  p_i=\frac{\|e_i^T\|_2^2}{m}=\frac1m .
\]
Every residual lies in the span of \(I_m\), and the canonical residual
coordinates are just the residual itself.  Therefore the previous corollary
specializes to uniform row sampling with \(d=m\):
\[
  \mathbb P\!\left[
    \Norm{A\widehat x(\omega)-b}^2
    \le
    (1+\eta)\Norm{Ax_{\rm opt}-b}^2
  \right]
  \ge
  1-\frac{1}{s(\varepsilon/m)^2}.
\]
This theorem is useful as a no-hidden-coordinate-assumption fallback.  It is
not the sharp survey theorem, because the paper's leverage basis should have
dimension equal to the residual subspace rank, not necessarily \(m\).
The analogous fully floating-point rounded-Gram objective transfer is also
formalized for \(U=I_m\), with the same \(d=m\) denominator and the same
explicit rounded-Gram representation hypothesis.  The implementation-backed
literal rounded \(A,b\) transfer is now proved for the augmented-span leverage
law above under its explicit objective-budget slack condition; an \(I_m\)
specialization can be derived separately if it becomes useful.
\end{corollary}

\begin{formalized}
\leanfit{hasOrthonormalColumns_idMatrix}
\leanfit{residualCoordinates_idMatrix}
\leanfit{residualsInColumnSpace_idMatrix}
\leanfit{leverageTraceProbability_eventProb_rowSampleLSObjective_le_one_add_eta_of_idBasis}
\leanfit{leverageTraceProbability_eventProb_fl_lsObjective_le_one_add_eta_of_idBasis}
\end{formalized}

\begin{remark}
This section now proves the deterministic least-squares implication, the
coordinate quadratic-error and finite-Loewner bridges, the probability transfer
to a random sketched minimizer, exact/FP leverage-score adapters, and the
concrete sampled-row least-squares objective algebra using canonical residual
coordinates.  The augmented-span theorem constructs a low-dimensional
orthonormal basis of \(\operatorname{span}\{A_{*,j},b\}\) and closes the
residual-coordinate hypothesis for that basis in the exact theorem, the
Bennett sample-budget theorem, and the rounded-Gram FP transfer, while the
identity-basis specialization remains a valid full-row-space exact/FP fallback.
The note also proves the first implementation-backed rounded sampled-row
foundation: rowwise residual perturbation and deterministic objective
perturbation bounds for the literal rounded sampled/scaled \(A,b\)
construction, and the corresponding support-aware high-probability rounded
minimizer transfer under an explicit objective-budget slack condition.  The
same support-aware theorem is also proved for additive-gap approximate rounded
minimizers, with the solver objective gap included explicitly in the slack
budget.  It does not yet prove a random-projection version or derive that gap
from a concrete floating-point solver/preconditioner pipeline.
\end{remark}

\section{Hypotheses and Scope}

\begin{itemize}
\item Algorithm 1 exact unbiasedness is proved for all entries under the
  canonical independent squared-magnitude trace law with \(steps=s\) and
  \(s\ne0\).
\item Algorithm 1 stability is entrywise and hit-counter based.  The Chernoff
  final API proves the hit-count moment-generating-function bound from the
  canonical product trace law; it is not assumed.
\item Algorithm 1 exact residuals are now decomposed as sums of one-sample
  mean-zero residual increments.  Entrywise, Frobenius, and fixed-vector
  second-moment bounds are proved for these increments, and exact and
  floating-point fixed-vector, finite-test-set, and supplied finite-cover
  Markov tails are proved for the full residual.  This is a
  matrix-concentration foundation, not by itself a uniform Bernstein or
  Khintchine spectral tail bound.
\item Algorithm 1 also has a self-adjoint dilation layer: the rectangular
  residual is embedded as \(\mathcal D(R)=\begin{bmatrix}0&R\\R^T&0\end{bmatrix}\),
  dilation events transfer to rectangular exact and floating-point residual
  events, and the dilation residual has a proved zero-mean increment
  decomposition with squared-Frobenius, quadratic-form, Loewner-order, and
  PSD variance proxies, plus trace variance proxies with the explicit
  dilation dimension \(m+n\).
  The finite square matrix layer also has product/trace vocabulary, PSD and
  Loewner-order vocabulary, basic trace algebra, trace monotonicity, Loewner
  sum/scaling algebra, symmetric-square PSD facts, the deterministic
  \(M^2\preceq L^2I\) consequence of an operator bound, and the converse
  adapter from a squared Loewner event to an operator event.
  It also proves \(\operatorname{tr}(\mathcal D(M)^2)=2\Frob{M}^2\).  These
  facts are now used by the source-aligned matrix-Bernstein proof route, while
  the generic rectangular Bernstein/Khintchine library remains a separate
  foundation target.
\item Algorithm 1 has a deterministic floating-point transfer theorem in
  rectangular vector-action form.  The repository also has a closed
  source-aligned square high-probability theorem for the hard-thresholded
  squared-magnitude distribution \(\hat A_{ij}^2/\|\hat A\|_F^2\), which is
  not the literal CACM Algorithm~1 distribution.  The theorem derives the
  entrywise
  floating-point perturbation budget from the local gamma/hit-count stability
  library and the sampler's probability-one support event; it does not assume
  the concentration or perturbation event as a hypothesis.  Its corollary
  expands that internal budget matrix into the explicit radius
  \(\varepsilon+(2n^2\Frob{A}^2/\varepsilon)\gamma_{s+1}\), where
  \(\gamma_{s+1}=((s+1)u)/(1-(s+1)u)\).  The repository also has a faithful
  nontruncated Frobenius/Markov FP corollary under the literal
  \(A_{ij}^2/\|A\|_F^2\) law, with explicit radius
  \(\varepsilon+n(\Frob{A}^2/\alpha)\gamma_{s+1}\) under a visible
  nonzero-entry floor \(\alpha\).  In addition, the literal support-radius
  theorem is closed with no nonzero-entry floor; it pays the reciprocal-entry
  term \(G(A)=\sum_{A_{ij}\ne0}\|A\|_F^2/|A_{ij}|\) explicitly in the exact
  threshold and in the FP radius.  The literal source-rate
  specialization is closed when every nonzero entry satisfies
  \(\varepsilon/(2n)\le |A_{ij}|\), so the source truncation is the identity.
  The source-uniform sharp literal equation~(2) matrix-Bernstein law and a
  fully rectangular CACM-prose variant remain separate theorem targets.
\item The Drineas--Zouzias truncation step for elementwise sparsification is
  now explicit: thresholding a square \(n\times n\) matrix at
  \(\varepsilon/(2n)\) costs at most \(\varepsilon/2\) in Frobenius/operator
  norm, and exact or floating-point half-budget theorems for sampling
  \(\hat A\) transfer to original-matrix residual events.  This truncation
  layer is composed below with the source-aligned matrix-Bernstein theorem and
  the derived floating-point gamma budget.
\item Algorithm 1 also has two older nonconditional weaker probability routes:
  Frobenius/Markov and scalar-entry/union-bound.  They remain useful sanity
  checks, but the advertised source-aligned equation~(2) theorem now uses the
  sharper trace-MGF/Bernstein route.
\item The product-law scalar MGF and finite-family self-adjoint-dilation
  quadratic-form theorems are support results for supplied scalar statistics or
  supplied test vectors.  The source-aligned equation~(2) theorem uses the
  noncommutative trace-MGF/largest-eigenvalue route formalized later in the
  note; a reusable generic Bernstein/Khintchine API beyond that specialized
  route remains open.
\item Algorithm 2 equation (5) is a Frobenius-norm high-probability theorem
  derived from a proved finite-probability second-moment bound and Markov's
  inequality.
\item Algorithm 2 equation (7) is stated as a vector-action operator-2 bound:
  for all \(x\), \(\Norm{Mx}\le c\Norm{x}\).  The repository intentionally does
  not introduce a separate supremum-valued spectral norm object.  The
  source-sharp Oliveira/Tropp route now has the one-step leverage covariance
  facts, row-trace MGF iteration, positive/negative centered log-CGF bounds,
  a two-sided finite-Loewner high-probability theorem with explicit exponential
  tails, Bennett sample-budget corollaries, and a matching floating-point
  finite-Loewner transfer.  The theorem keeps the upper and lower scalar
  Bennett denominators explicit.
\item Algorithm 3 random-projection concentration is distribution-specific.
  The formalization proves exact orthogonal Frobenius preservation,
  deterministic preservation of the orthonormal-column leverage basis and
  equation~(6) denominator, the signed-diagonal orthogonality prerequisite for
  the SRHT route, the finite Rademacher sign-vector law with a probability-one
  signed-preprocessing event, finite Rademacher first/second moments, and the
  exact scalar signed-linear-form MGF factorization, plus the one-sided
  exponential-Markov tail skeleton, scalar Hoeffding MGF/tail bounds, and a
  weaker coordinate-Hoeffding row-norm theorem.  The coordinate-Hoeffding
  theorem is also lifted to an equation~(6) leverage-probability bound
  \(p_i(HD_\omega U)\le B^2\) under the same finite-entry failure budget.
  The deterministic-after-preconditioning iid uniform sample-average
  concentration theorem is also formalized in finite-Loewner tail-budget form.
  The signed-Hadamard preprocessing event and the uniform row-sampling event
  are then composed on the product probability law, giving the exact
  two-sided sample-Gram event with probability
  \(1-\delta_{\mathrm{pre}}-\delta_{\mathrm{sample}}\) for the scoped
  coordinate-Hoeffding route.  It also proves the
  flat-Hadamard expected row-norm identity
  \(\mathbb E\Norm{(HD_\omega U)_{i,*}}^2=n/m\), together with the weak
  Markov/union consequence
  \(\mathbb P\{\forall i,\Norm{(HD_\omega U)_{i,*}}^2\le T\}
  \ge 1-m((n/m)/T)\).  It also proves deterministic floating-point stability
  for the preprocessing products.  The source-sharp SRHT row-norm branch is now
  closed by a positive-drop finite-cube route: one obtains the all-row cap
  \(\Norm{(HD_\omega U)_{i,*}}\le \sqrt{n/m}+t\), the induced leverage cap
  \((\sqrt{n/m}+t)^2/n\), logarithmic preprocessing-budget wrappers, exact
  SRHT-plus-uniform-row composition, and a source-sharp FP constant-budget
  transfer for rounded uniform sketches.  It does not prove a universal
  Gaussian/FJLT/input-sparsity random-projection uniformization theorem, nor
  concrete implementation certificates for generated/stored transforms, FHT or
  sign storage, computed bases/singular vectors/projectors, or scale
  denominators.
\item Equation~(8) now has a deterministic sketched least-squares objective
  theorem and a coordinate quadratic-error bridge from an operator-2
  Gram-error event to residual-objective preservation.  It also has the
  concrete sampled-row objective identity for rows of \(A\) and \(b\) scaled by
  leverage probabilities from \(U\), using the canonical coordinates
  \(U^T(Ax-b)\).  The augmented-span theorem constructs an orthonormal basis
  of \(\operatorname{span}\{A_{*,j},b\}\), proves coordinates for \([A\ b]\),
  and discharges the residual-column-space condition with dimension
  \(d=\dim\operatorname{span}\{A_{*,j},b\}\) when \(d>0\).  The same
  augmented-span and identity-basis coordinate discharges are now packaged for
  the fully floating-point rounded-Gram objective transfer.  The identity-basis
  fallback also discharges the condition for uniform row sampling with
  \(d=m\).  The leverage-score Bennett sample-budget theorem is now closed in
  exact arithmetic and as a rounded-Gram FP transfer.  The literal rounded
  sampled/scaled \(A,b\) construction now has residual and objective
  perturbation bounds and a support-aware high-probability rounded-minimizer
  transfer under an explicit objective-budget slack condition, including an
  additive solver-objective-gap variant.  A concrete normal-equations/Cholesky
  solver transfer is formalized, and the rectangular stored-QR route now reaches
  equation~(8) through visible-domain certificates: local diagonal dominance or
  row-budget control, source-denominator/nonbreakdown, compact-product
  smallness, active-pivot/scalar-comparison fields, and rounded triangular
  solve budgets are explicit hypotheses rather than hidden conclusions.
  Random-projection equation~(8) variants and the generic claim that an
  arbitrary computed QR/preconditioner loop supplies those source/domain fields
  remain open.
\item The full CACM article also discusses matrix completion, Laplacian-based
  solvers, kernel methods, and other applications.  Those are not explicit
  meta-algorithms like Algorithms 1--3, and their application-specific
  guarantees require additional foundations before they can be stated as Lean
  theorems.
\end{itemize}

\section{Proof-Source Ledger}

The full-paper automation now maintains a separate proof-source ledger:
\[
  \texttt{docs/RANDNLA\_CACM\_PROOF\_SOURCE\_LEDGER.md}.
\]
Its purpose is to prevent repeated proof work around the same missing
mathematical engine.  For every source claim whose proof is incomplete in the
CACM article, the ledger records the source location, the missing proof step,
the external primary source chain, the intended Lean target, and whether the
source has been formalized locally, rejected, kept advisory, or left open.

\begin{definition}[Proof-source status]
A source theorem is \emph{formalized} only after the corresponding Lean theorem
is proved locally or reused from an already formalized local dependency.  A
source theorem is \emph{advisory} when it guides the statement, proof order, or
constants but is not available as a Lean theorem.  Advisory sources never close
a paper-level theorem.
\end{definition}

\begin{theorem}[Current source chain for Algorithm 1 equation (2)]
The active proof-source chain for the Algorithm~1 spectral concentration
claim is:
\[
\begin{array}{c}
\text{CACM Algorithm 1 and equation (2)}\\
\Downarrow\\
\text{Drineas--Zouzias Algorithm 1, Theorem 1, and Lemmas 1--4}\\
\Downarrow\\
\text{matrix-valued Bernstein for the sampled truncated residual}\\
\Downarrow\\
\text{Tropp Theorem 3.2, Corollary 3.3, Theorem 3.6, Corollary 3.7,}\\
\text{Theorem 1.4, Theorem 1.6, and the relative-entropy/Lieb route in}\\
\text{Tropp's 2012 note and monograph, with Hansen--Pedersen Jensen,}\\
\text{Effros's matrix perspective,}\\
\text{and Lindblad's relative-entropy convexity theorem as source-chain}\\
\text{references for the joint-convexity step.}
\end{array}
\]
The local formalization has already proved the truncation transfer,
self-adjoint dilation, variance/boundedness prerequisites, trace/eigenvalue
Markov interfaces, C\(^*\)-matrix expectation/order bridges, and finite Jensen
adapter.  It has also named the strictly positive cone, proved that this cone
is convex, named the local Lieb trace functional, and proved the local
\(\log(\exp X)=X\) and \(\exp(\log A)=A\) functional-calculus inverse
bridges needed by that route.  The chosen relative-entropy route also now has
scalar and finite-vector nonnegativity, the finite log-sum inequality,
scalar and finite-vector joint convexity, local vocabulary \(D(A;B)\), the
diagonal normalization \(D(A;A)=0\), and the real scalar-identity case
\(D(aI;bI)\ge 0\) for \(a,b>0\), together with the real diagonal case
\(D(\Delta(a);\Delta(b))=D_{\mathrm{vec}}(a;b)\) and its positivity
corollary for positive diagonals, plus the real-diagonal joint-convexity
subcase, and the left/right multiplication endomorphism layer \(L_A,R_A\)
with product/power algebra, and the explicit ratio endomorphism \(L_XR_A^{-1}\)
needed for the Effros
perspective route, as well as the finite Kronecker lifts \(A\otimes I\) and
\(I\otimes H\), their affine laws, commutation/product identities, and
  positive-definiteness, trace-normalization, product-index vectorization
  action, vectorized-identity trace-pairing facts, and the polynomial
  powered trace-pairing identity
  \(v_I^*(A\otimes B^{\mathsf T})^k v_I=\operatorname{tr}(A^kB^k)\), and
  the finite-polynomial version of this trace identity, plus continuity of the
  finite quadratic form and finite polynomial evaluation in the matrix argument.
  It also now splits the
  Hansen--Pedersen/Effros source route into ordinary positive-cone operator
  convexity, Hansen--Pedersen transfer, and assembled two-point Jensen targets;
  it proves the identity-function sanity cases, the assembly adapters, the
  direct \(x\log x\) operator-convexity input, and the concrete two-point
  \(x\log x\) Jensen theorem by a reflection-average corner argument.  It also
  proves the conditional
relative-entropy route reduction: joint convexity of \(D\) plus the normalized
entropy variational formula imply the finite-dimensional Lieb trace concavity
theorem for
\[
  A \mapsto
  \operatorname{Re}\operatorname{tr}
  \exp\!\bigl(H+\log A\bigr)
\]
on the strictly positive cone.  For that normalized variational formula, the
optimizer-candidate equality
\(\Psi_H(\exp(H+\log A),A)=\Phi_H(A)\) is also formalized, and the maximality
part is reduced to matrix relative-entropy nonnegativity.  On the
nonnegativity side, the generalized Klein/nonnegativity equivalence, the
Hermitian spectral-overlap expansion, and the positive-spectrum scalar entropy
kernel are now formalized.  The bridge from the separated Hermitian CFC
theorem to the compact complex C\(^*\)-matrix logarithm form of generalized
Klein is also formalized, so local matrix relative-entropy nonnegativity and
the normalized variational formula are closed.  Matrix relative-entropy joint
convexity is also now formalized through the finite-dimensional Effros
perspective route, and this gives finite-dimensional Lieb trace concavity and
the one-step Tropp trace-MGF inequality without an explicit Lieb hypothesis.
The Algorithm~1 product-law trace-MGF iteration, finite-real trace-exponential
adapter, residual-dilation specialization, and trace-exponential
Markov-to-eigenvalue step are also formalized.  The remaining concentration
foundation is the scalar matrix-CGF/log-MGF bound, with explicit
Bernstein/Khintchine constants, needed for the Algorithm~1 equation~(2)
high-probability spectral-norm theorem.
The proof-source ledger records
Effros's matrix-perspective proof route, the Hansen--Pedersen operator Jensen
source behind that route, Lindblad's original relative-entropy convexity
lineage, and Bendat--Sherman's operator monotonicity/convexity route as
advisory primary sources.  The concrete Hansen--Pedersen `x log x` theorem is
now represented by local Lean proofs, the CFC square-root algebra needed for
perspective statements is local, and the ordinary finite perspective theorem
for \(f(x)=x\log x-(x-1)\) is formalized.  The finite-polynomial Kronecker
trace-representation layer is also local, and an explicit uniform polynomial
approximation hypothesis is now enough to transfer that layer to the normalized
entropy-kernel CFC trace term.  The positive-definite spectrum case now also
constructs the required uniform polynomial approximation and specializes the
trace convergence theorem to \(X\otimes(A^{-1})^{\mathsf T}\).  The common
polynomial sequence also converges on every eigenvalue ratio
\(\lambda_i(X)\lambda_j(A)^{-1}\), so the limiting overlap theorem and the
source-faithful Umegaki superoperator trace representation are now formalized.
The earlier ordinary source-matrix perspective route is not used to close this
step, because it is not the \(L_XR_A^{-1}\) perspective required by
Tropp--Effros.  No theorem is closed by citation.
\end{theorem}


\begin{corollary}[Current full-paper gate implication]
Matrix relative-entropy joint convexity, finite-dimensional Lieb trace
concavity, and the one-step Tropp trace-MGF inequality are now formalized
locally.  The iid independent-sum trace-MGF domination theorem, its finite-real
trace-exponential adapter, the actual Algorithm~1 self-adjoint dilation
residual specialization, and the trace-exponential Markov-to-eigenvalue step
are also formalized locally.  The scalar Bernstein parabola with constants and
the CFC step that lifts it under a real-spectrum upper bound to the operator
quadratic bound
\(\exp(\theta X)\preceq I+\theta X+g(\theta,R)X^2\) are formalized as well.
The generic centered one-sample log-CGF variance proxy
\[
  \log\mathbb E\exp(\theta X)\preceq g(\theta,R)\mathbb E[X^2]
\]
is now also formalized locally, including a support-aware version.  The
truncated Algorithm~1 dilation-increment instantiation of that theorem is now
formalized locally as well.  The one-sample dilation variance proxy, the
negative-support lower-tail companion, and the positive/negative scalar
trace-MGF bounds now give parameterized upper-tail and two-sided
trace-exponential eigenvalue bounds for the truncated dilation residual, plus
an explicit \(1-\delta\) scaled-eigenvalue corollary obtained by choosing
\(T=\log(2B/\delta)\).  The deterministic bridge from that scaled eigenvalue
event to a rectangular spectral event at radius \(\log(2B/\delta)/\theta\) is
also formalized, and the scalar Bennett choice of \(\theta\) now gives a
radius-\(r\) high-probability theorem under the explicit Bennett budget.
The proof-source ledger now marks the source denominator, source sample budget,
deterministic truncation transfer, and support-aware floating-point transfer as
closed for the hard-thresholded element-wise variant.  The explicit entrywise
perturbation budget used by the FP transfer theorem is now derived locally from
the \(\gamma\)- and hit-count stability infrastructure on the sampler's
probability-one support.  The same support-aware mechanism is now composed with
the literal, nontruncated Frobenius/Markov route and with the source-rate
no-small-entry specialization recorded at the end of the Algorithm~1 section.
\end{corollary}

\section{Low-Rank Approximation Foundations}

\begin{corollary}[Exact rank certificates for basis-product approximants]
Let \(A\in\R^{m\times n}\), \(V\in\R^{n\times q}\), and
\(U\in\R^{m\times q}\).  Define the exact basis-product matrices
\[
  P_V = V V^T,\qquad P_U = U U^T,
\]
and the projected approximants
\[
  B_R = A P_V,\qquad B_L = P_U A .
\]
Then \(B_R\) has the explicit rank factorization
\[
  B_R(i,j)
    =
    \sum_{a=1}^q
      \left(\sum_{k=1}^n A(i,k)V(k,a)\right)V(j,a),
\]
so \(B_R\) has repository rank at most \(q\).  Similarly \(B_L\) has the
explicit rank factorization
\[
  B_L(i,j)
    =
    \sum_{a=1}^q
      U(i,a)\left(\sum_{k=1}^m U(k,a)A(k,j)\right),
\]
so \(B_L\) has repository rank at most \(q\).

This is an exact-object foundation for the equation~(9) low-rank route.  It
does not assume that \(U\) or \(V\) was computed, and it introduces no
probability-construction or sampling-law error.  If an implementation computes
the bases, singular vectors, projectors, pseudoinverses, or downstream products,
those non-probability computations remain subject to the existing computed-basis
and computed-projector certificate rules.
\end{corollary}

\begin{formalized}
\leanfit{RectRankFactorization}
\leanfit{RectRankAtMost}
\leanfit{lowRankResidualFrob}
\leanfit{IsBestRankApproxFrob}
\leanfit{rightBasisProjectorApprox}
\leanfit{leftBasisProjectorApprox}
\leanfit{rightBasisProjectorApproxFactorization}
\leanfit{leftBasisProjectorApproxFactorization}
\leanfit{rightBasisProjectorApprox_rankAtMost}
\leanfit{leftBasisProjectorApprox_rankAtMost}
\end{formalized}

\begin{corollary}[Best-rank comparison for exact basis-product approximants]
Let \(A,A_k\in\R^{m\times n}\), and suppose \(A_k\) carries a repository
certificate \(\mathrm{BestRankFrob}(A,A_k,k)\), meaning that
\[
  \operatorname{rank}_{\mathrm{cert}}(A_k)\le k
  \quad\text{and}\quad
  \lVert A-A_k\rVert_F \le \lVert A-B\rVert_F
  \quad
  \text{for every }B\text{ with }\operatorname{rank}_{\mathrm{cert}}(B)\le k .
\]
Then every exact repository-rank-\(k\) comparison matrix satisfies
\[
  \lVert A-A_k\rVert_F \le \lVert A-B\rVert_F .
\]
In particular, for \(V\in\R^{n\times k}\) and \(U\in\R^{m\times k}\), the
rank certificates above give
\[
  \lVert A-A_k\rVert_F
    \le
  \lVert A-A(VV^T)\rVert_F,
  \qquad
  \lVert A-A_k\rVert_F
    \le
  \lVert A-(UU^T)A\rVert_F .
\]

This is still an exact-object theorem.  It compares against exact analysis
matrices \(A(VV^T)\) and \((UU^T)A\); if an algorithm computes \(U\), \(V\),
projectors, singular vectors, pseudoinverses, products, or downstream solver
inputs, those non-probability computations require the explicit
floating-point/inexact-arithmetic certificates recorded elsewhere.  Sampling
probabilities and laws remain exact mathematical inputs by the current
convention.
\end{corollary}

\begin{formalized}
\leanfit{IsBestRankApproxFrob.residual_le_of_rankAtMost}
\leanfit{IsBestRankApproxFrob.residual_le_rightBasisProjectorApprox}
\leanfit{IsBestRankApproxFrob.residual_le_leftBasisProjectorApprox}
\end{formalized}

\begin{corollary}[Column-sketch projector rank certificate for equation (9)]
Let \(A\in\R^{m\times n}\), \(Z\in\R^{n\times r}\), and define the exact column
sketch
\[
  B = A Z \in \R^{m\times r}.
\]
Let \(P_{AZ}\in\R^{m\times m}\) be an exact left multiplier that factors through
the displayed columns of \(B\): assume there is a coefficient matrix
\(C\in\R^{r\times m}\) such that
\[
  P_{AZ}(i,\ell)=\sum_{a=1}^r B(i,a)C(a,\ell).
\]
Then the exact projected approximation \(P_{AZ}A\) has the rank factorization
\[
  (P_{AZ}A)(i,j)
    =
    \sum_{a=1}^r
      B(i,a)\left(\sum_{\ell=1}^m C(a,\ell)A(\ell,j)\right),
\]
and hence has repository rank at most \(r\).

This is the exact rank surface needed for the source equation~(9) projector
\(P_{AZ}\).  It does not prove that \(P_{AZ}\) is the orthogonal projector onto
\(\operatorname{range}(AZ)\), does not construct \((V_k^T Z)^+\), does not prove
the full-rank condition on \(V_k^T Z\), and does not prove the source residual
bound involving \(\Sigma_{k,\perp}(V_{k,\perp}^T Z)(V_k^T Z)^+\).  If an
algorithm computes \(Z\), \(A Z\), \(P_{AZ}\), a pseudoinverse, or the product
\(P_{AZ}A\), those non-probability computations require explicit FP or
inexact-arithmetic certificates.  Sampling probabilities and laws remain exact
mathematical inputs.
\end{corollary}

\begin{formalized}
\leanfit{columnSketch}
\leanfit{LeftFactorThrough}
\leanfit{leftProductFactorizationOfLeftFactorThrough}
\leanfit{preconditionRows_rankAtMost_of_leftFactorThrough}
\leanfit{sketchColumnProjectorApprox_rankAtMost}
\end{formalized}

\begin{corollary}[Equation (9) residual-certificate surface]
Let \(A\in\R^{m\times n}\), \(Z\in\R^{n\times r}\), and
\(P_{AZ}\in\R^{m\times m}\).  Assume first that \(P_{AZ}\) factors through the
exact column sketch \(AZ\), in the sense of the preceding corollary.  Assume
also that there are exact analysis scalars \(T,C\ge 0\) such that the displayed
equation~(9) residual certificate holds:
\[
  \lVert A-P_{AZ}A\rVert_F \le T+C .
\]
Then the exact candidate \(P_{AZ}A\) has repository rank at most \(r\), and
\[
  \lVert A-P_{AZ}A\rVert_F \le T+C .
\]
Moreover, if \(A_k\) carries a repository certificate
\(\operatorname{BestRankFrob}(A,A_k,k)\), and if the source-side scalar
comparison
\[
  T+C \le \rho\,\lVert A-A_k\rVert_F
\]
is supplied, then
\[
  \operatorname{rankcert}(A_k)\le k,\qquad
  \operatorname{rankcert}(P_{AZ}A)\le r,\qquad
  \lVert A-P_{AZ}A\rVert_F
    \le \rho\,\lVert A-A_k\rVert_F .
\]

This is not a proof of the CACM equation~(9) SVD/pseudoinverse inequality.
Rather, it makes the residual inequality an explicit exact-object certificate:
the missing source instantiation must still prove that the tail and coupling
terms \(T\) and \(C\) can be chosen from the singular-subspace expression
\(\lVert\Sigma_{k,\perp}\rVert\) and
\(\lVert\Sigma_{k,\perp}(V_{k,\perp}^T Z)(V_k^T Z)^+\rVert\).  It therefore
keeps rectangular SVD, the full-rank condition on \(V_k^T Z\), pseudoinverse
algebra, unitarily invariant norm infrastructure, and every computed
non-probability object as separate obligations.  Sampling probabilities and
laws remain exact mathematical inputs.
\end{corollary}

\begin{formalized}
\leanfit{Equation9ResidualCertificate}
\leanfit{Equation9ResidualCertificate.tail_add_coupling_nonneg}
\leanfit{equation9RankResidualSurface}
\leanfit{equation9RelativeResidualSurface}
\end{formalized}

\begin{corollary}[Explicit coefficient handoff for the equation (9) projector]
Let \(A\in\R^{m\times n}\), \(Z\in\R^{n\times r}\), and let
\(C\in\R^{r\times m}\) be any exact coefficient table.  Define
\[
  P_C := (AZ)C,\qquad
  P_C(i,j)=\sum_{a=1}^r (AZ)(i,a)C(a,j).
\]
Then \(P_C\) factors through the exact column sketch \(AZ\), and therefore
\[
  \operatorname{rankcert}(P_C A)\le r .
\]
If, in addition, \(P_C\) carries an equation~(9) residual certificate with
exact terms \(T,C_0\ge 0\), then
\[
  \lVert A-P_CA\rVert_F \le T+C_0 .
\]
Finally, if \(A_k\) is certified best rank \(k\) and \(T+C_0\le
\rho\lVert A-A_k\rVert_F\), then the relative residual conclusion
\[
  \lVert A-P_CA\rVert_F \le \rho\lVert A-A_k\rVert_F
\]
holds together with the displayed rank certificates.

This corollary is the algebraic handoff for the future choice
\(C=(AZ)^+\).  It does not prove that the supplied \(C\) is a pseudoinverse,
that \(P_C\) is the orthogonal projector onto \(\operatorname{range}(AZ)\), or
that the equation~(9) residual certificate follows from the source SVD
hypotheses.  Computed construction of \(C\), \(P_C\), a pseudoinverse, or
products remains a non-probability FP/certificate obligation.
\end{corollary}

\begin{formalized}
\leanfit{columnSketchLeftMultiplier}
\leanfit{columnSketchLeftMultiplier_leftFactorThrough}
\leanfit{columnSketchLeftMultiplier_rankAtMost}
\leanfit{columnSketchLeftMultiplier_equation9RankResidualSurface}
\leanfit{columnSketchLeftMultiplier_equation9RelativeResidualSurface}
\end{formalized}

\begin{corollary}[Generalized-inverse projector surface for the equation (9) sketch]
Let \(A\in\R^{m\times n}\), \(Z\in\R^{n\times r}\), and write
\[
  B=AZ,\qquad P_C=BC
\]
for an exact coefficient table \(C\in\R^{r\times m}\).  Assume the displayed
generalized-inverse condition
\[
  P_C B = B,
  \qquad\text{equivalently}\qquad
  \sum_{\ell=1}^m
    \left(\sum_{a=1}^r B(i,a)C(a,\ell)\right)B(\ell,b)
  =
  B(i,b)
\]
for every row \(i\) and sketch column \(b\).  Then \(P_C\) reproduces the exact
column sketch and is idempotent:
\[
  P_C B = B,\qquad P_C^2=P_C .
\]
Moreover \(P_C\) factors through the exact sketch columns, and hence the
projected approximation satisfies
\[
  \operatorname{rankcert}(P_C A)\le r .
\]

This is the exact algebraic projector surface needed before a future
pseudoinverse instantiation \(C=(AZ)^+\).  It does not prove the Moore--Penrose
conditions, orthogonality or symmetry of \(P_C\), existence of \(C=(AZ)^+\),
the full-rank condition on \(V_k^T Z\), or the equation~(9) SVD residual
certificate.  Computed construction of \(C\), \(P_C\), pseudoinverses, or
products remains a non-probability FP/certificate obligation.
\end{corollary}

\begin{formalized}
\leanfit{ColumnSketchGeneralizedInverse}
\leanfit{columnSketchLeftMultiplier_reproducesSketch_of_generalizedInverse}
\leanfit{columnSketchLeftMultiplier_idempotent_of_generalizedInverse}
\leanfit{columnSketchLeftMultiplier_projectorSurface_of_generalizedInverse}
\end{formalized}

\begin{corollary}[Symmetric-idempotent projector surface for the equation (9) sketch]
In the setting of the preceding corollary, assume in addition that the exact
multiplier \(P_C=BC\) is symmetric:
\[
  P_C(i,j)=P_C(j,i)\qquad\text{for all }i,j.
\]
Then \(P_C\) is a symmetric idempotent multiplier that reproduces the sketch:
\[
  P_C^T=P_C,\qquad P_C^2=P_C,\qquad P_C B=B .
\]
It also factors through the sketch columns, and so
\[
  \operatorname{rankcert}(P_C A)\le r .
\]

This is the exact symmetric-idempotent surface expected of the source
orthogonal projector onto \(\operatorname{range}(AZ)\).  It does not construct
the Moore--Penrose pseudoinverse, prove that \(C=(AZ)^+\) satisfies the
certificate, prove full rank of \(V_k^T Z\), or supply the equation~(9)
residual certificate.  Computed construction of \(C\), \(P_C\), pseudoinverses,
or products remains a non-probability FP/certificate obligation.
\end{corollary}

\begin{formalized}
\leanfit{ColumnSketchOrthogonalProjectorCertificate}
\leanfit{columnSketchLeftMultiplier_symmetric_of_orthogonalProjectorCertificate}
\leanfit{columnSketchLeftMultiplier_idempotent_of_orthogonalProjectorCertificate}
\leanfit{columnSketchLeftMultiplier_orthogonalProjectorSurface}
\end{formalized}

\begin{corollary}[Moore--Penrose certificate handoff for the equation (9) sketch]
Let \(A\in\mathbb R^{m\times n}\), \(Z\in\mathbb R^{n\times r}\), and
\[
  B=AZ,\qquad P_C=BC,\qquad Q_C=CB
\]
for an exact coefficient table \(C\in\mathbb R^{r\times m}\).  Assume the four
exact Moore--Penrose certificate equations
\[
  P_C B=B,\qquad Q_C C=C,\qquad P_C^T=P_C,\qquad Q_C^T=Q_C .
\]
Equivalently, the first two equations are \(BCB=B\) and \(CBC=C\).  Then the
first and third fields supply the generalized-inverse and symmetric-projector
certificates for \(P_C=(AZ)C\).  Consequently
\[
  P_C B=B,\qquad P_C^2=P_C,\qquad P_C^T=P_C,\qquad
  \operatorname{rankcert}(P_C A)\le r .
\]
The second and fourth fields are retained as the coefficient-side
Moore--Penrose obligations \(CBC=C\) and \((CB)^T=CB\), so later
pseudoinverse algebra can use the full four-equation surface.

This corollary is a certificate handoff only.  It does not construct
\((AZ)^+\), prove that such a \(C\) exists, prove full rank of \(V_k^T Z\), or
supply the equation~(9) residual certificate.  Computed construction of \(C\),
\(P_C\), \(Q_C\), pseudoinverses, or products remains a non-probability
FP/certificate obligation.
\end{corollary}

\begin{formalized}
\leanfit{columnSketchRightMultiplier}
\leanfit{ColumnSketchMoorePenroseCertificate}
\leanfit{ColumnSketchMoorePenroseCertificate.to_generalizedInverse}
\leanfit{ColumnSketchMoorePenroseCertificate.to_orthogonalProjectorCertificate}
\leanfit{columnSketchRightMultiplier_reproducesCoeff_of_moorePenroseCertificate}
\leanfit{columnSketchRightMultiplier_symmetric_of_moorePenroseCertificate}
\leanfit{columnSketchLeftMultiplier_orthogonalProjectorSurface_of_moorePenroseCertificate}
\end{formalized}

\begin{corollary}[Gram-inverse Moore--Penrose instantiation for the equation (9) sketch]
Let \(B=AZ\), let
\[
  G=B^T B,\qquad C=G_{\mathrm{inv}}B^T,
  \qquad C(a,i)=\sum_b G_{\mathrm{inv}}(a,b)B(i,b).
\]
Assume that \(G_{\mathrm{inv}}\) is an exact two-sided inverse certificate for
\(G\), and assume that \(G_{\mathrm{inv}}\) is symmetric.  Then
\[
  C B=I,\qquad C B C=C.
\]
Moreover \(B C\) and \(C B\) are symmetric, and the generalized-inverse
equation also holds:
\[
  B C B=B,\qquad (BC)^T=BC,\qquad (CB)^T=CB.
\]
Therefore the coefficient table \(C=G_{\mathrm{inv}}B^T\) satisfies the
four-equation Moore--Penrose certificate of the preceding corollary.  In
particular, for \(P_C=BC\),
\[
  P_C B=B,\qquad P_C^2=P_C,\qquad P_C^T=P_C,\qquad
  \operatorname{rankcert}(P_C A)\le r .
\]

This is still exact-object certificate algebra.  It does not prove that
\(G=B^TB\) is invertible, construct \(G_{\mathrm{inv}}\), prove the full-rank
condition for \(AZ\) or \(V_k^T Z\), or supply the equation~(9) residual
certificate.  If an implementation computes \(G\), \(G_{\mathrm{inv}}\), \(C\),
\(P_C\), or downstream products, those non-probability computations require
explicit FP or inexact-arithmetic certificates.  Sampling probabilities and
laws remain exact mathematical inputs.
\end{corollary}

\begin{formalized}
\leanfit{columnSketchGram}
\leanfit{columnSketchGramInverseCoefficient}
\leanfit{ColumnSketchGramInverseCertificate}
\leanfit{columnSketchRightMultiplier_eq_id_of_gramInverseCertificate}
\leanfit{columnSketchGramInverseCoefficient_reproducesCoeff}
\leanfit{columnSketchLeftMultiplier_symmetric_of_gramInverseCertificate}
\leanfit{columnSketchRightMultiplier_symmetric_of_gramInverseCertificate}
\leanfit{columnSketchGramInverseCoefficient_generalizedInverse}
\leanfit{columnSketchGramInverseCoefficient_moorePenroseCertificate}
\leanfit{columnSketchLeftMultiplier_orthogonalProjectorSurface_of_gramInverseCertificate}
\end{formalized}

\begin{corollary}[Determinant-facing Gram-inverse route for the equation (9) sketch]
Let \(B=AZ\) and \(G=B^TB\).  The exact Gram matrix is symmetric:
\[
  G^T=G .
\]
Assume the exact determinant condition
\[
  \det(G)\ne 0,
\]
and define
\[
  G_{\mathrm{ns}}^{-1}:=\operatorname{nonsingInv}(G),\qquad
  C=G_{\mathrm{ns}}^{-1}B^T .
\]
Then \(G_{\mathrm{ns}}^{-1}\) is a two-sided inverse certificate for \(G\), and
\[
  (G_{\mathrm{ns}}^{-1})^T=G_{\mathrm{ns}}^{-1}.
\]
Consequently the coefficient table \(C\) satisfies the same four
Moore--Penrose equations as in the preceding corollary:
\[
  B C B=B,\qquad C B C=C,\qquad (BC)^T=BC,\qquad (CB)^T=CB.
\]
For \(P_C=BC\), it follows that
\[
  P_C B=B,\qquad P_C^2=P_C,\qquad P_C^T=P_C,\qquad
  \operatorname{rankcert}(P_C A)\le r .
\]

This is an exact-object determinant theorem.  It does not prove the determinant
condition from full column rank of \(AZ\), does not connect that condition to
the source hypothesis on \(V_k^T Z\), and does not supply the equation~(9)
residual certificate.  If an implementation computes \(G\), \(G_{\mathrm{ns}}^{-1}\),
\(C\), \(P_C\), or downstream products, those non-probability computations
require explicit FP or inexact-arithmetic certificates.  Sampling probabilities
and laws remain exact mathematical inputs.
\end{corollary}

\begin{formalized}
\leanfit{nonsingInv_symmetric_of_symmetric}
\leanfit{columnSketchGram_symmetric}
\leanfit{columnSketchGramInverseCertificate_of_det_ne_zero}
\leanfit{columnSketchGramInverseCoefficient_moorePenroseCertificate_of_det_ne_zero}
\leanfit{columnSketchLeftMultiplier_orthogonalProjectorSurface_of_det_ne_zero}
\end{formalized}

\begin{corollary}[Thin-factor Gram determinant bridge for the equation (9) sketch]
Let \(B=AZ\).  Suppose there are exact matrices
\[
  U\in\mathbb R^{m\times r},\qquad R\in\mathbb R^{r\times r}
\]
such that
\[
  B=UR,\qquad U^T U=I,\qquad \det(R)\ne0 .
\]
Then the exact sketch Gram matrix satisfies
\[
  B^T B = R^T R,
\]
and therefore
\[
  \det(B^T B)=\det(R)^2\ne0 .
\]
Consequently the determinant-facing corollary applies with
\[
  G=B^TB,\qquad G_{\mathrm{ns}}^{-1}:=\operatorname{nonsingInv}(G),
  \qquad C=G_{\mathrm{ns}}^{-1}B^T .
\]
Thus \(C\) satisfies
\[
  B C B=B,\qquad C B C=C,\qquad (BC)^T=BC,\qquad (CB)^T=CB,
\]
and for \(P_C=BC\),
\[
  P_C B=B,\qquad P_C^2=P_C,\qquad P_C^T=P_C,\qquad
  \operatorname{rankcert}(P_C A)\le r .
\]

This is an exact-object thin-factor theorem.  It does not construct such a
thin factorization from rectangular QR or SVD, does not connect it to the
source hypothesis on \(V_k^T Z\), and does not supply the equation~(9) residual
certificate.  If an implementation computes \(U\), \(R\), \(G\),
\(G_{\mathrm{ns}}^{-1}\), \(C\), \(P_C\), or downstream products, those
non-probability computations require explicit FP or inexact-arithmetic
certificates.  Sampling probabilities and laws remain exact mathematical inputs.
\end{corollary}

\begin{formalized}
\leanfit{ColumnSketchThinFactorCertificate}
\leanfit{columnSketchGram_eq_factorGram_of_thinFactorCertificate}
\leanfit{columnSketchGram_det_ne_zero_of_thinFactorCertificate}
\leanfit{columnSketchGramInverseCertificate_of_thinFactorCertificate}
\leanfit{columnSketchGramInverseCoefficient_moorePenroseCertificate_of_thinFactorCertificate}
\leanfit{columnSketchLeftMultiplier_orthogonalProjectorSurface_of_thinFactorCertificate}
\end{formalized}

\begin{corollary}[Source-SVD head rank certificate for equation~(9)]
Let
\[
  U\in\mathbb R^{m\times r},\qquad
  \Sigma\in\mathbb R^{r\times r},\qquad
  V\in\mathbb R^{n\times r},
\]
and define the exact source head
\[
  H = U\Sigma V^T,\qquad
  H_{ij}=\sum_{a=1}^r U_{ia}\sum_{b=1}^r \Sigma_{ab}V_{jb}.
\]
Then \(H\) factors through the displayed source dimension \(r\):
\[
  H_{ij}=\sum_{a=1}^r U_{ia}R_{aj},
  \qquad
  R_{aj}=\sum_{b=1}^r \Sigma_{ab}V_{jb}.
\]
Consequently
\[
  \operatorname{rankcert}(H)\le r .
\]

This is an exact-object D3 source-SVD theorem.  It supplies the rank
certificate for the displayed source head only.  It does not construct the
rectangular SVD, prove source-tail orthogonality, prove singular-value
comparison or Eckart--Young optimality, or certify computed
SVD/singular-vector/projector/Gram data.  Sampling probabilities and laws
remain exact mathematical inputs.
\end{corollary}

\begin{formalized}
\leanfit{sourceSVDFactorMatrixRankFactorization}
\leanfit{sourceSVDFactorMatrix_rankAtMost}
\end{formalized}

\begin{corollary}[Source-split residual and best-rank handoff for equation~(9)]
Let \(A,H,T\in\mathbb R^{m\times n}\), let
\[
  H=U\Sigma V^T,
\]
and suppose the exact source split
\[
  A = H+T
\]
holds entrywise.  Then the source-head residual is exactly the tail norm:
\[
  \rho_F(A,H)
  :=\|A-H\|_F
  =\|T\|_F .
\]
More generally, for any supplied rectangular norm certificate \(\xi\),
\[
  \rho_\xi(A,H)
  :=\xi(A-H)
  =\xi(T).
\]
If, in addition, the Frobenius tail-optimality inequality
\[
  \|T\|_F\le \rho_F(A,B)
  \qquad
  \text{for every }B\text{ with }\operatorname{rankcert}(B)\le r
\]
is supplied, then \(H\) is a Frobenius best rank-\(r\) approximation to \(A\):
\[
  \operatorname{BestRankFrob}_r(A,H).
\]
Likewise, if
\[
  \xi(T)\le \rho_\xi(A,B)
  \qquad
  \text{for every }B\text{ with }\operatorname{rankcert}(B)\le r,
\]
then \(H\) is a best rank-\(r\) approximation for the norm \(\xi\):
\[
  \operatorname{BestRank}_{\xi,r}(A,H).
\]

This is an exact-object D3/D4 handoff.  The residual equalities are proved
from the displayed source split, and the rank side uses the preceding source
head rank certificate.  The Eckart--Young/tail-optimality inequalities remain
explicit hypotheses until the rectangular SVD and singular-value comparison
foundations are proved.  Any implementation that computes \(U,\Sigma,V,T\), or
the norm data must supply non-probability FP or inexact-arithmetic
certificates.  Sampling probabilities and laws remain exact mathematical
inputs.
\end{corollary}

\begin{formalized}
\leanfit{lowRankResidualFrob_sourceSVDFactorMatrix_eq_tail}
\leanfit{lowRankResidualNorm_sourceSVDFactorMatrix_eq_tail}
\leanfit{sourceSVDFactorMatrix_isBestRankApproxFrob_of_tail_optimal}
\leanfit{sourceSVDFactorMatrix_isBestRankApproxNorm_of_tail_optimal}
\end{formalized}

\begin{corollary}[Source-SVD Gram determinant bridge for the equation (9) sketch]
Let \(A\in\mathbb R^{m\times n}\), \(Z\in\mathbb R^{n\times r}\), and suppose
there are exact matrices
\[
  U\in\mathbb R^{m\times r},\qquad
  \Sigma\in\mathbb R^{r\times r},\qquad
  V\in\mathbb R^{n\times r}
\]
such that
\[
  A = U\Sigma V^T,\qquad U^T U=I.
\]
Define the exact source cross factor and sketch right factor by
\[
  M:=V^T Z,\qquad R:=\Sigma M=\Sigma(V^T Z).
\]
If
\[
  \det(\Sigma)\ne0,\qquad \det(M)\ne0,
\]
then
\[
  AZ=U R,\qquad \det(R)=\det(\Sigma)\det(M)\ne0 .
\]
Consequently the thin-factor corollary applies, and hence
\[
  \det\!\bigl((AZ)^T(AZ)\bigr)\ne0 .
\]
With
\[
  G=(AZ)^T(AZ),\qquad G_{\mathrm{ns}}^{-1}:=\operatorname{nonsingInv}(G),
  \qquad C=G_{\mathrm{ns}}^{-1}(AZ)^T ,
\]
the coefficient table \(C\) satisfies
\[
  (AZ)C(AZ)=AZ,\qquad C(AZ)C=C,\qquad
  ((AZ)C)^T=(AZ)C,\qquad (C(AZ))^T=C(AZ).
\]
For \(P_C=(AZ)C\),
\[
  P_C(AZ)=AZ,\qquad P_C^2=P_C,\qquad P_C^T=P_C,\qquad
  \operatorname{rankcert}(P_CA)\le r .
\]

This is an exact-object source-factor theorem.  It assumes the exact
factorization \(A=U\Sigma V^T\), exact orthonormality of \(U\), and exact
determinant hypotheses on \(\Sigma\) and \(V^T Z\).  It does not prove
rectangular SVD existence, Eckart--Young best-rank optimality, the source
equation~(9) residual inequality, or any computed SVD/singular-vector/
projector/Gram/inverse/product routine.  If an implementation computes
\(U\), \(\Sigma\), \(V\), \(M\), \(R\), \(G\), \(G_{\mathrm{ns}}^{-1}\),
\(C\), \(P_C\), or downstream products, those non-probability computations
require explicit FP or inexact-arithmetic certificates.  Sampling
probabilities and laws remain exact mathematical inputs.
\end{corollary}

\begin{formalized}
\leanfit{sourceSVDFactorMatrix}
\leanfit{rightSketchCrossGram}
\leanfit{sourceSVDSketchRightFactor}
\leanfit{columnSketch_eq_sourceSVDFactorization}
\leanfit{sourceSVDSketchRightFactor_det_ne_zero_of_det_ne_zero}
\leanfit{columnSketchThinFactorCertificate_of_sourceSVD}
\leanfit{columnSketchThinFactorCertificate_of_sourceSVD_det_factors}
\leanfit{columnSketchGram_det_ne_zero_of_sourceSVD_det_factors}
\leanfit{columnSketchGramInverseCertificate_of_sourceSVD_det_factors}
\leanfit{columnSketchGramInverseCoefficient_moorePenroseCertificate_of_sourceSVD_det_factors}
\leanfit{columnSketchLeftMultiplier_orthogonalProjectorSurface_of_sourceSVD_det_factors}
\end{formalized}

\begin{corollary}[Source-SVD head-tail residual bridge for equation (9)]
Let \(A\in\mathbb R^{m\times n}\), \(Z\in\mathbb R^{n\times r}\), and define
the exact Gram-inverse projector
\[
  B:=AZ,\qquad
  G:=B^TB,\qquad
  P_G:=B\,\operatorname{nonsingInv}(G)\,B^T .
\]
Assume the source-SVD determinant hypotheses of the previous corollary:
\[
  A=U\Sigma V^T,\qquad U^TU=I,\qquad
  \det(\Sigma)\ne0,\qquad \det(V^TZ)\ne0 .
\]
Then \(P_G\) is symmetric, reproduces the sketch, is idempotent on itself, and
\(\operatorname{rankcert}(P_GA)\le r\).

Now suppose there are exact matrices
\[
  H,T\in\mathbb R^{m\times n},\qquad W\in\mathbb R^{r\times n},
\]
and exact nonnegative radii \(\tau,\kappa\ge0\), such that
\[
  A=H+T,\qquad H=BW,\qquad
  \|T\|_F\le \tau,\qquad \|P_GT\|_F\le \kappa .
\]
Since \(P_GB=B\), the head is reproduced:
\[
  P_GH=P_GBW=BW=H .
\]
Therefore
\[
  A-P_GA = T-P_GT,
\]
and the rectangular Frobenius subtraction inequality gives
\[
  \|A-P_GA\|_F
    =\|T-P_GT\|_F
    \le \|T\|_F+\|P_GT\|_F
    \le \tau+\kappa .
\]
Equivalently, the data above instantiate the equation~(9) residual certificate
with tail radius \(\tau\) and coupling radius \(\kappa\).  In particular,
\[
  \operatorname{rankcert}(P_GA)\le r,\qquad
  \|A-P_GA\|_F\le \tau+\kappa .
\]
If \(A_k\) is supplied with an exact best-rank-\(k\) Frobenius certificate and
\[
  \tau+\kappa \le \rho\,\|A-A_k\|_F,
\]
then
\[
  \operatorname{rankcert}(A_k)\le k,\qquad
  \operatorname{rankcert}(P_GA)\le r,\qquad
  \|A-P_GA\|_F\le \rho\,\|A-A_k\|_F .
\]

This is an exact-object residual bridge.  It does not construct \(H,T,W\) from
rectangular SVD data, prove Eckart--Young optimality, prove the source
\(V_k^TZ\) pseudoinverse residual inequality, or supply computed
SVD/projector/Gram/inverse/product error certificates.  If an implementation
computes \(U\), \(\Sigma\), \(V\), \(H\), \(T\), \(W\), \(G\), \(P_G\), or any
downstream product, those non-probability computations require explicit FP or
inexact-arithmetic certificates.  Sampling probabilities and laws remain exact
mathematical inputs.
\end{corollary}

\begin{formalized}
\leanfit{frobNormRect_neg}
\leanfit{frobNormRect_sub_le}
\leanfit{ColumnSketchHeadFactorization}
\leanfit{preconditionRows_reproduces_head_of_columnSketchHeadFactorization}
\leanfit{Equation9HeadTailSketchCertificate}
\leanfit{Equation9HeadTailSketchCertificate.to_residualCertificate}
\leanfit{equation9HeadTailSketchRankResidualSurface}
\leanfit{equation9HeadTailSketchRelativeResidualSurface}
\leanfit{columnSketchGramInverseProjector}
\leanfit{columnSketchGramInverseProjector_orthogonalProjectorSurface_of_sourceSVD_det_factors}
\leanfit{columnSketchGramInverseProjector_sourceSVD_headTailRankResidualSurface}
\leanfit{columnSketchGramInverseProjector_sourceSVD_headTailRelativeResidualSurface}
\end{formalized}

\begin{corollary}[Explicit-coefficient head-tail bridge for equation (9)]
Keep the exact source-SVD determinant hypotheses and the projector \(P_G\) from
the preceding corollary.  Let \(W\in\mathbb R^{r\times n}\) be any exact
coefficient table and define
\[
  H_W:=(AZ)W,\qquad T_W:=A-H_W .
\]
Then \(H_W\) factors through the sketch by definition and
\[
  A=H_W+T_W .
\]
Consequently, if exact radii \(\tau,\kappa\ge0\) satisfy
\[
  \|T_W\|_F\le \tau,\qquad \|P_GT_W\|_F\le \kappa ,
\]
then the head-tail certificate of the preceding corollary is instantiated with
\(H=H_W\), \(T=T_W\), and hence
\[
  \operatorname{rankcert}(P_GA)\le r,\qquad
  \|A-P_GA\|_F\le \tau+\kappa .
\]
If \(A_k\) carries an exact best-rank-\(k\) Frobenius certificate and
\[
  \tau+\kappa\le \rho\,\|A-A_k\|_F,
\]
then
\[
  \operatorname{rankcert}(A_k)\le k,\qquad
  \operatorname{rankcert}(P_GA)\le r,\qquad
  \|A-P_GA\|_F\le \rho\,\|A-A_k\|_F .
\]

This is still an exact-object adapter.  The remaining source proof must choose
the pseudoinverse-derived \(W\) and prove the two displayed norm bounds from
rectangular SVD, the \(V_k^TZ\) condition, and unitarily invariant norm
infrastructure.  If an implementation computes \(W\), \(T_W\), \(P_G\), or any
source SVD/projector/Gram/inverse/product data, those non-probability
computations require explicit FP or inexact-arithmetic certificates.
\end{corollary}

\begin{formalized}
\leanfit{columnSketchHead}
\leanfit{columnSketchTail}
\leanfit{columnSketchHead_headFactorization}
\leanfit{columnSketchHeadTail_split}
\leanfit{equation9HeadTailSketchCertificate_of_columnSketchHead}
\leanfit{columnSketchGramInverseProjector_sourceSVD_columnSketchHeadRankResidualSurface}
\leanfit{columnSketchGramInverseProjector_sourceSVD_columnSketchHeadRelativeResidualSurface}
\end{formalized}

\begin{corollary}[Source-coefficient residual-tail bridge for equation (9)]
Let \(U\in\mathbb R^{m\times r}\), \(\Sigma\in\mathbb R^{r\times r}\),
\(V\in\mathbb R^{n\times r}\), \(Z\in\mathbb R^{n\times r}\), and let
\[
  M:=V^TZ .
\]
Assume
\[
  \det(M)\ne0
\]
and define the exact coefficient table
\[
  W:=M^{-1}V^T .
\]
Then
\[
  MW=V^T,\qquad \Sigma MW=\Sigma V^T,
\]
and the source head is reproduced from its own sketch:
\[
  \bigl((U\Sigma V^T)Z\bigr)W=U\Sigma V^T .
\]
Now let \(A,T\in\mathbb R^{m\times n}\) satisfy the exact source split
\[
  A=U\Sigma V^T+T .
\]
Then the implemented coefficient-generated sketch head and sketch tail obey
\[
  (AZ)W=U\Sigma V^T+(TZ)W,
\]
and
\[
  A-(AZ)W=T-(TZ)W .
\]
Equivalently, with
\[
  T_{\mathrm{src}}:=T-(TZ)(V^TZ)^{-1}V^T,
\]
the canonical sketch-tail object \(T_W=A-(AZ)W\) is exactly
\[
  T_W=T_{\mathrm{src}} .
\]

Consequently, for any exact multiplier \(P_{AZ}\) that factors through the
columns of \(AZ\) and reproduces \(AZ\), if exact radii \(\tau,\kappa\ge0\)
satisfy
\[
  \|T_{\mathrm{src}}\|_F\le\tau,\qquad
  \|P_{AZ}T_{\mathrm{src}}\|_F\le\kappa,
\]
then
\[
  \operatorname{rankcert}(P_{AZ}A)\le r,\qquad
  \|A-P_{AZ}A\|_F\le\tau+\kappa .
\]
If \(A_k\) carries an exact best-rank-\(k\) Frobenius certificate and
\[
  \tau+\kappa\le \rho\,\|A-A_k\|_F,
\]
then
\[
  \operatorname{rankcert}(A_k)\le k,\qquad
  \operatorname{rankcert}(P_{AZ}A)\le r,\qquad
  \|A-P_{AZ}A\|_F\le\rho\,\|A-A_k\|_F .
\]

This is an exact-object source-coefficient theorem.  It chooses the paper's
pseudoinverse-derived coefficient table \(W=(V^TZ)^{-1}V^T\) and rewrites the
implemented sketch tail as the explicit source tail
\(T-(TZ)(V^TZ)^{-1}V^T\).  It still does not prove the two Frobenius norm
bounds, construct the exact orthogonal projector onto \(\operatorname{range}(AZ)\)
for the head-plus-tail sketch, prove Eckart--Young optimality, or supply
computed cross-product/inverse/coefficient/SVD/projector/product error
certificates.  Sampling probabilities and laws remain exact mathematical
inputs.
\end{corollary}

\begin{formalized}
\leanfit{sourceSketchCoefficient}
\leanfit{rightSketchCrossGram_sourceSketchCoefficient}
\leanfit{sourceSVDSketchRightFactor_sourceSketchCoefficient}
\leanfit{columnSketchHead_sourceSVDFactorMatrix_sourceSketchCoefficient}
\leanfit{columnSketch_eq_add_of_eq_add}
\leanfit{columnSketchHead_eq_add_of_columnSketch_eq_add}
\leanfit{columnSketchHead_sourceHeadTail_sourceSketchCoefficient}
\leanfit{sourceSketchResidualTail}
\leanfit{columnSketchTail_sourceHeadTail_sourceSketchCoefficient}
\leanfit{equation9HeadTailSketchCertificate_of_sourceHeadTail_sourceSketchCoefficient}
\leanfit{equation9RankResidualSurface_of_sourceHeadTail_sourceSketchCoefficient}
\leanfit{equation9RelativeResidualSurface_of_sourceHeadTail_sourceSketchCoefficient}
\end{formalized}

\begin{corollary}[Moore--Penrose source-tail projector bridge for equation (9)]
Keep the hypotheses and definitions of the preceding source-coefficient
corollary:
\[
  M=V^TZ,\qquad \det(M)\ne0,\qquad W=M^{-1}V^T,\qquad
  A=U\Sigma V^T+T,
\]
and
\[
  T_{\mathrm{src}}:=T-(TZ)(V^TZ)^{-1}V^T .
\]
Let \(C\in\mathbb R^{r\times m}\) be an exact coefficient table for the full
sketch \(B:=AZ\), and set
\[
  P_C:=BC=(AZ)C .
\]
Assume \(C\) satisfies the four exact Moore--Penrose certificate equations
\[
  BCB=B,\qquad CBC=C,\qquad (BC)^T=BC,\qquad (CB)^T=CB .
\]
Then \(P_C\) is symmetric, factors through \(AZ\), reproduces \(AZ\), and is
idempotent:
\[
  P_C^T=P_C,\qquad P_C(AZ)=AZ,\qquad P_C^2=P_C .
\]
If exact radii \(\tau,\kappa\ge0\) satisfy
\[
  \|T_{\mathrm{src}}\|_F\le\tau,\qquad
  \|P_C T_{\mathrm{src}}\|_F\le\kappa,
\]
then
\[
  \operatorname{rankcert}(P_CA)\le r,\qquad
  \|A-P_CA\|_F\le \tau+\kappa .
\]
If \(A_k\) carries an exact best-rank-\(k\) Frobenius certificate and
\[
  \tau+\kappa\le \rho\,\|A-A_k\|_F,
\]
then
\[
  \operatorname{rankcert}(A_k)\le k,\qquad
  \operatorname{rankcert}(P_CA)\le r,\qquad
  \|A-P_CA\|_F\le\rho\,\|A-A_k\|_F .
\]

This is an exact-object projector-certificate theorem.  It removes the generic
unnamed \(P_{AZ}\) hypothesis from the previous corollary and replaces it by the
four displayed Moore--Penrose equations for the full sketch \(AZ\).  It still
does not construct \(C=(AZ)^+\), prove that a concrete Gram or pseudoinverse
routine satisfies those equations for the head-plus-tail sketch, prove the two
source-tail Frobenius bounds, or supply computed cross-product/inverse/
coefficient/SVD/projector/product error certificates.  Sampling probabilities
and laws remain exact mathematical inputs.
\end{corollary}

\begin{formalized}
\leanfit{columnSketchLeftMultiplier_sourceHeadTail_sourceSketchCoefficientRankResidualSurface}
\leanfit{columnSketchLeftMultiplier_sourceHeadTail_sourceSketchCoefficientRelativeResidualSurface}
\end{formalized}

\begin{corollary}[Concrete Gram-inverse source-tail projector bridge for equation~(9)]
Keep the source-coefficient hypotheses
\[
  M=V^T Z,\qquad \det(M)\ne0,\qquad
  A=U\Sigma V^T+T,
\]
and
\[
  T_{\mathrm{src}}:=T-(TZ)(V^T Z)^{-1}V^T .
\]
Let \(B:=AZ\), \(G:=B^T B\), and assume the exact Gram determinant condition
\[
  \det(G)\ne0 .
\]
Define the exact Gram-inverse coefficient table and projector
\[
  C_G:=G^{-1}B^T,\qquad
  P_G:=B C_G=(AZ)\bigl((AZ)^T(AZ)\bigr)^{-1}(AZ)^T .
\]
Then the determinant-facing Gram route supplies the four exact
Moore--Penrose equations
\[
  BC_GB=B,\qquad C_GBC_G=C_G,\qquad
  (BC_G)^T=BC_G,\qquad (C_GB)^T=C_GB .
\]
Consequently \(P_G\) is symmetric, factors through \(AZ\), reproduces \(AZ\),
and is idempotent:
\[
  P_G^T=P_G,\qquad P_G(AZ)=AZ,\qquad P_G^2=P_G .
\]
If exact radii \(\tau,\kappa\ge0\) satisfy
\[
  \|T_{\mathrm{src}}\|_F\le\tau,\qquad
  \|P_G T_{\mathrm{src}}\|_F\le\kappa,
\]
then
\[
  \operatorname{rankcert}(P_GA)\le r,\qquad
  \|A-P_GA\|_F\le \tau+\kappa .
\]
If \(A_k\) carries an exact best-rank-\(k\) Frobenius certificate and
\[
  \tau+\kappa\le \rho\,\|A-A_k\|_F,
\]
then
\[
  \operatorname{rankcert}(A_k)\le k,\qquad
  \operatorname{rankcert}(P_GA)\le r,\qquad
  \|A-P_GA\|_F\le\rho\,\|A-A_k\|_F .
\]

This is an exact-object Gram-inverse theorem.  It instantiates the
Moore--Penrose certificate \(C\) of the previous corollary with the concrete
table \(C_G=\operatorname{nonsingInv}((AZ)^T(AZ))(AZ)^T\), under the displayed
exact determinant hypothesis.  It still does not prove \(\det((AZ)^T(AZ))\ne0\)
from source SVD/full-rank hypotheses for the head-plus-tail sketch, prove the
two source-tail Frobenius bounds, prove Eckart--Young optimality, build
unitarily invariant norm infrastructure, or supply computed
Gram/inverse/coefficient/SVD/projector/product error certificates.  Sampling
probabilities and laws remain exact mathematical inputs.
\end{corollary}

\begin{formalized}
\leanfit{columnSketchGramInverseProjector_sourceHeadTail_sourceSketchCoefficientRankResidualSurface_of_det_ne_zero}
\leanfit{columnSketchGramInverseProjector_sourceHeadTail_sourceSketchCoefficientRelativeResidualSurface_of_det_ne_zero}
\end{formalized}

\begin{corollary}[Source-tail orthogonality from exact tail factors for equation~(9)]
Let
\[
  T=U_\perp\Sigma_\perp V_\perp^T,
  \qquad
  T_{ij}=\sum_{\beta} (U_\perp)_{i\beta}
      \sum_{\gamma}(\Sigma_\perp)_{\beta\gamma}(V_\perp)_{j\gamma},
\]
and let \(U\in\mathbb R^{m\times r}\) be the exact source-head left basis.
If the exact head and tail left bases are cross-orthogonal,
\[
  \sum_i U_{i\alpha}(U_\perp)_{i\beta}=0
  \qquad\text{for every }\alpha,\beta,
\]
then the source-tail orthogonality predicate required by the head-tail
sketch-Gram split holds:
\[
  U^T T=0,
  \qquad
  \sum_i U_{i\alpha}T_{ij}=0
  \quad\text{for every }\alpha,j .
\]

This is an exact-object D3 theorem.  It derives the displayed \(U^T T=0\)
field from supplied tail factors and left-basis cross-orthogonality.  It does
not construct the rectangular SVD or the tail factors, prove their
cross-orthogonality, prove singular-value comparison or Eckart--Young
optimality, or certify computed SVD/singular-vector/projector/Gram data.
Sampling probabilities and laws remain exact mathematical inputs.
\end{corollary}

\begin{formalized}
\leanfit{sourceTailLeftOrthogonal_of_tail_factor_left_cross_zero}
\end{formalized}

\begin{corollary}[Orthogonal source head-tail sketch-Gram split for equation~(9)]
Let
\[
  H:=U\Sigma V^T,\qquad A=H+T,
\]
where \(U\in\mathbb R^{m\times r}\), \(\Sigma\in\mathbb R^{r\times r}\),
\(V\in\mathbb R^{n\times r}\), \(T\in\mathbb R^{m\times n}\), and
\(Z\in\mathbb R^{n\times r}\).  Assume the exact source-tail
left-orthogonality condition
\[
  U^T T=0,
  \qquad\text{i.e.}\qquad
  \sum_i U_{i\alpha}T_{ij}=0
  \quad\text{for all } \alpha,j .
\]
Define the sketched source-head and tail blocks
\[
  B_H:=HZ,\qquad B_T:=TZ.
\]
Then the sketched source head has the exact factorization
\[
  HZ=U\bigl(\Sigma(V^T Z)\bigr).
\]
Moreover the two cross-Gram terms vanish:
\[
  B_H^T B_T=0,\qquad B_T^T B_H=0.
\]
Consequently the full sketch Gram matrix splits as
\[
  (AZ)^T(AZ)
  =
  (HZ)^T(HZ)+(TZ)^T(TZ).
\]
Equivalently, for all sketch columns \(a,b\),
\[
  \sum_i (AZ)_{ia}(AZ)_{ib}
  =
  \sum_i (HZ)_{ia}(HZ)_{ib}
  +
  \sum_i (TZ)_{ia}(TZ)_{ib}.
\]

This is an exact-object determinant-route theorem.  It closes the first
algebraic component of the active equation~(9) determinant bottleneck: the
head-plus-tail sketch Gram is now reduced to a source-head sketch Gram plus
the tail sketch Gram under the displayed source-tail orthogonality hypothesis.
It does not yet prove that the tail sketch Gram is positive semidefinite, that
the source-head Gram is positive definite or has nonzero determinant from
\(\det(\Sigma)\ne0\) and \(\det(V^T Z)\ne0\), or that adding the tail Gram
preserves the nonzero determinant.  It also does not supply source-tail norm
bounds, SVD/Eckart--Young/unitarily invariant norm infrastructure, or computed
SVD/Gram/product certificates.  Sampling probabilities and laws remain exact
mathematical inputs.
\end{corollary}

\begin{formalized}
\leanfit{sourceTailLeftOrthogonal}
\leanfit{columnSketch_sourceSVDFactorMatrix}
\leanfit{columnSketch_sourceSVDFactorMatrix_tail_leftOrthogonal}
\leanfit{columnSketch_tail_sourceSVDFactorMatrix_leftOrthogonal}
\leanfit{columnSketchGram_sourceHeadTail_leftOrthogonal}
\end{formalized}

\begin{corollary}[Column-sketch Gram PSD and determinant preservation for equation~(9)]
For any exact matrix \(A\in\mathbb R^{m\times n}\) and sketch
\(Z\in\mathbb R^{n\times r}\), let
\[
  B:=AZ,\qquad G_A:=B^T B .
\]
Then for every vector \(x\in\mathbb R^r\),
\[
  x^T G_A x
  =
  \sum_i \left(\sum_a B_{ia}x_a\right)^2 .
\]
In particular \(G_A\succeq0\).

Now keep the orthogonal source split of the preceding corollary:
\[
  H=U\Sigma V^T,\qquad A=H+T,\qquad U^T T=0,
\]
and define
\[
  G_H:=(HZ)^T(HZ),\qquad G_T:=(TZ)^T(TZ).
\]
The previous corollary gives
\[
  G_A=G_H+G_T,
\]
and the first part gives \(G_T\succeq0\).  Therefore, if the exact source-head
sketch Gram carries a positive-definite certificate
\[
  G_H\succ0,
\]
then
\[
  G_A=G_H+G_T\succ0,\qquad \det(G_A)\ne0 .
\]

This is an exact-object determinant-preservation theorem.  It closes the PSD
tail-Gram and positive-definite-plus-PSD part of the active equation~(9)
determinant route.  The remaining determinant obligation is now narrower: prove
the displayed positive-definite certificate \(G_H\succ0\) from the source
thin-factor/full-rank data, in particular from \(\det(\Sigma)\ne0\) and
\(\det(V^T Z)\ne0\).  The result still does not prove source-tail norm bounds,
SVD/Eckart--Young/unitarily invariant norm infrastructure, or computed
SVD/Gram/product certificates.  Sampling probabilities and laws remain exact
mathematical inputs.
\end{corollary}

\begin{formalized}
\leanfit{finiteQuadraticForm_columnSketchGram_eq_sum_sq}
\leanfit{columnSketchGram_finitePSD}
\leanfit{matrix_det_ne_zero_of_posDef_add_posSemidef}
\leanfit{columnSketchGram_sourceHeadTail_det_ne_zero_of_head_posDef}
\end{formalized}

\begin{corollary}[Source-factor positive-definite sketch Gram and determinant
closure for equation~(9)]
Let \(R\in\mathbb R^{r\times r}\) be an exact square factor.  If
\[
  \det(R)\ne0,
\]
then
\[
  R^T R\succ0.
\]
Equivalently, for every nonzero \(x\in\mathbb R^r\),
\[
  x^T R^T R x=\|Rx\|_2^2>0.
\]
Consequently, if an exact sketch \(B=AZ\) has a thin factorization
\[
  B=UR,\qquad U^T U=I_r,\qquad \det(R)\ne0,
\]
then
\[
  (AZ)^T(AZ)=R^T R\succ0.
\]

Now let
\[
  H=U\Sigma V^T,\qquad A=H+T,\qquad U^T T=0,\qquad U^T U=I_r,
\]
and assume the exact source determinant hypotheses
\[
  \det(\Sigma)\ne0,\qquad \det(V^T Z)\ne0.
\]
Set
\[
  R:=\Sigma(V^T Z).
\]
Then \(\det(R)\ne0\), the source-head sketch satisfies
\[
  HZ=UR,
\]
and therefore
\[
  G_H:=(HZ)^T(HZ)\succ0.
\]
Combining this positive-definite source-head certificate with the preceding
PSD-tail determinant bridge gives
\[
  \det\!\big((AZ)^T(AZ)\big)\ne0 .
\]

This is an exact-object determinant theorem.  It closes the exact determinant
side of the active equation~(9) bottleneck from source factor hypotheses:
nonzero \(\det(\Sigma)\), nonzero \(\det(V^T Z)\), exact orthonormal source
columns, and exact source-tail left orthogonality.  The result still does not
prove the source-tail norm bounds, rectangular SVD/Eckart--Young/unitarily
invariant norm infrastructure, or implementation-facing certificates for
computed SVD/singular-vector/projector/Gram/inverse/product routines.  Sampling
probabilities and laws remain exact mathematical inputs.
\end{corollary}

\begin{formalized}
\leanfit{matrix_transpose_mul_self_posDef_of_det_ne_zero}
\leanfit{columnSketchGram_posDef_of_thinFactorCertificate}
\leanfit{columnSketchGram_posDef_of_sourceSVD_det_factors}
\leanfit{columnSketchGram_sourceHeadTail_det_ne_zero_of_sourceSVD_det_factors}
\end{formalized}

\begin{corollary}[Tail-factor determinant route for equation~(9)]
Let
\[
  A=U\Sigma V^T+T,
  \qquad
  T=U_\perp\Sigma_\perp V_\perp^T,
\]
where
\[
  T_{ij}=\sum_{\beta} (U_\perp)_{i\beta}
      \sum_{\gamma}(\Sigma_\perp)_{\beta\gamma}(V_\perp)_{j\gamma}.
\]
Assume the exact left-basis cross-orthogonality and source-head
orthonormality conditions
\[
  U^T U_\perp=0,
  \qquad
  U^T U=I_r .
\]
Equivalently,
\[
  \sum_i U_{i\alpha}(U_\perp)_{i\beta}=0,
  \qquad
  \sum_i U_{i\alpha}U_{i\beta}=\delta_{\alpha\beta}.
\]
If the exact source determinant hypotheses also hold,
\[
  \det(\Sigma)\ne0,\qquad \det(V^T Z)\ne0,
\]
then
\[
  \det\!\big((AZ)^T(AZ)\big)\ne0 .
\]

The proof first derives \(U^T T=0\) from the displayed tail factorization and
\(U^T U_\perp=0\), then invokes the source-factor determinant theorem above.
Thus the determinant-facing route can consume exact tail factors directly
instead of requiring a raw source-tail orthogonality hypothesis.

This is an exact-object D3-to-D1 theorem.  It still does not construct the
rectangular SVD or tail factors, prove left/right basis orthogonality from SVD
existence, prove the source-tail norm bounds, prove Eckart--Young optimality,
or certify computed SVD/singular-vector/projector/Gram/inverse/product
routines.  All of those are non-probability computed-object obligations in
implementation-facing versions.  Sampling probabilities and laws remain exact
mathematical inputs.
\end{corollary}

\begin{formalized}
\leanfit{columnSketchGram_sourceHeadTail_det_ne_zero_of_sourceSVD_det_factors_tail_factor_left_cross_zero}
\end{formalized}

\begin{corollary}[Diagonal singular-block determinant route for equation~(9)]
Let
\[
  A=U\Sigma V^T+T,
  \qquad
  T=U_\perp\Sigma_\perp V_\perp^T,
\]
with the same exact tail-factor interpretation as above.  Assume
\[
  U^T U_\perp=0,\qquad U^T U=I_r,\qquad \det(V^T Z)\ne0.
\]
Let \(\sigma:\{1,\ldots,r\}\to\mathbb R\) be the displayed singular-value
list for the exact source block, and assume the exact diagonal identity
\[
  \Sigma_{\alpha\beta}=
  \begin{cases}
    \sigma_\alpha, & \alpha=\beta,\\
    0, & \alpha\ne\beta .
  \end{cases}
\]
If
\[
  \sigma_\alpha\ne0\qquad\text{for every }\alpha,
\]
then
\[
  \det(\Sigma)=\prod_{\alpha=1}^r\sigma_\alpha\ne0,
\]
and therefore
\[
  \det\!\big((AZ)^T(AZ)\big)\ne0 .
\]
In particular, the same conclusion holds under the stronger displayed
singular-value hypothesis \(\sigma_\alpha>0\) for every \(\alpha\).

Thus the determinant theorem can consume the exact diagonal singular block
and nonzero displayed singular values instead of a raw hypothesis
\(\det(\Sigma)\ne0\).  This is still an exact-object D3-to-D1 result: it does
not construct the rectangular SVD, prove that an implementation computed the
diagonal block or singular values with perturbation guarantees, establish
singular-value ordering, prove source-tail norm bounds, prove Eckart--Young
optimality, or certify computed SVD, singular-vector, projector, Gram,
inverse, or product routines.  Sampling probabilities and laws remain exact
mathematical inputs.
\end{corollary}

\begin{formalized}
\leanfit{matrix_det_ne_zero_of_eq_diagonal_nonzero}
\leanfit{matrix_det_ne_zero_of_eq_diagonal_pos}
\leanfit{columnSketchGram_sourceHeadTail_det_ne_zero_of_sourceSVD_diagonal_tail_factor_left_cross_zero}
\end{formalized}

\begin{corollary}[Diagonal source-SVD sketch and projector route for equation~(9)]
Let
\[
  H=U\Sigma V^T,\qquad U^T U=I_r,\qquad \det(V^T Z)\ne0,
\]
and let the exact source singular block be diagonal,
\[
  \Sigma_{\alpha\beta}=
  \begin{cases}
    \sigma_\alpha, & \alpha=\beta,\\
    0, & \alpha\ne\beta ,
  \end{cases}
  \qquad
  \sigma_\alpha\ne0 \quad\text{for every }\alpha.
\]
Then the exact source sketch right factor is nonsingular:
\[
  \det\!\big(\Sigma(V^T Z)\big)\ne0.
\]
Consequently \(HZ\) has the exact thin factorization
\[
  HZ=U\,\Sigma(V^T Z),
\]
with orthonormal left factor and nonsingular square right factor.  Hence
\[
  \det\!\big((HZ)^T(HZ)\big)\ne0,
  \qquad
  (HZ)^T(HZ)\succ0 .
\]
If additionally \(A=H+T\) and \(U^T T=0\), then the head-plus-tail sketch Gram
satisfies
\[
  \det\!\big((AZ)^T(AZ)\big)\ne0 .
\]

For the source-head matrix \(H\), let
\[
  P_H := HZ\big((HZ)^T(HZ)\big)^{-1}(HZ)^T
\]
be the exact Gram-inverse column-sketch projector.  The same diagonal
singular-block hypotheses give the exact Gram-inverse certificate, the
Moore--Penrose coefficient certificate, and the projector surface
\[
  P_H^T=P_H,\qquad P_H^2=P_H,\qquad P_H(HZ)=HZ,\qquad
  \operatorname{rank}(P_HH)\le r .
\]

This exact-object result propagates the displayed diagonal singular-value
hypotheses through the source-SVD sketch/projector algebra.  It does not
construct the rectangular SVD, prove the diagonal block or singular-value
positivity/order from SVD existence, prove source-tail norm bounds, prove
Eckart--Young optimality, or certify computed SVD, singular-vector, projector,
Gram, inverse, or product routines.  Sampling probabilities and laws remain
exact mathematical inputs.
\end{corollary}

\begin{formalized}
\leanfit{sourceSVDSketchRightFactor_det_ne_zero_of_diagonal_nonzero}
\leanfit{sourceSVDSketchRightFactor_det_ne_zero_of_diagonal_pos}
\leanfit{columnSketchThinFactorCertificate_of_sourceSVD_diagonal_det_factors}
\leanfit{columnSketchGram_det_ne_zero_of_sourceSVD_diagonal_det_factors}
\leanfit{columnSketchGram_posDef_of_sourceSVD_diagonal_det_factors}
\leanfit{columnSketchGram_sourceHeadTail_det_ne_zero_of_sourceSVD_diagonal_det_factors}
\leanfit{columnSketchGramInverseCertificate_of_sourceSVD_diagonal_det_factors}
\leanfit{columnSketchGramInverseCoefficient_moorePenroseCertificate_of_sourceSVD_diagonal_det_factors}
\leanfit{columnSketchLeftMultiplier_orthogonalProjectorSurface_of_sourceSVD_diagonal_det_factors}
\leanfit{columnSketchGramInverseProjector_orthogonalProjectorSurface_of_sourceSVD_diagonal_det_factors}
\end{formalized}

\begin{corollary}[Diagonal source-SVD residual surfaces for equation~(9)]
Let
\[
  H=U\Sigma V^T,\qquad U^T U=I_r,\qquad \det(V^T Z)\ne0,
\]
and let the exact source singular block be diagonal with nonzero displayed
entries,
\[
  \Sigma_{\alpha\beta}=
  \begin{cases}
    \sigma_\alpha, & \alpha=\beta,\\
    0, & \alpha\ne\beta ,
  \end{cases}
  \qquad
  \sigma_\alpha\ne0 \quad\text{for every }\alpha.
\]
Let
\[
  P_H := HZ\big((HZ)^T(HZ)\big)^{-1}(HZ)^T
\]
be the exact Gram-inverse column-sketch projector generated by the preceding
corollary.  If an exact equation-(9) head/tail certificate supplies matrices
\(H_0,T_0\) and radii \(\tau,\kappa\) such that
\[
  \|H-(P_H H)\|_F
  \le
  \tau+\kappa
\]
through the repository head/tail certificate, then the concrete residual
surface holds with the displayed diagonal singular-block hypotheses, not with
a raw \(\det(\Sigma)\ne0\) assumption:
\[
  \rho_F(H,P_HH) \le \tau+\kappa .
\]
If \(A_k\) is a supplied best rank-\(k\) approximation certificate and
\[
  \tau+\kappa \le \rho\,\rho_F(H,A_k),
\]
then
\[
  \rho_F(H,P_HH)\le \rho\,\rho_F(H,A_k).
\]
The same diagonal replacement holds for the explicit-coefficient head/tail
certificate obtained from a supplied exact coefficient table \(W\).

For the source-head plus source-tail form, let
\[
  A=H+T,\qquad U^T T=0,
\]
and let
\[
  P_A := AZ\big((AZ)^T(AZ)\big)^{-1}(AZ)^T,\qquad
  W=(V^T Z)^{-1}V^T.
\]
The same exact diagonal singular-block hypotheses imply
\[
  \det\!\big((AZ)^T(AZ)\big)\ne0,
\]
so the concrete Gram-inverse projector is available without a raw
sketch-Gram determinant hypothesis.  If the exact source-tail residual and
its projected image obey
\[
  \bigl\|T-(T Z)W\bigr\|_F \le \tau,\qquad
  \bigl\|P_A\bigl(T-(T Z)W\bigr)\bigr\|_F \le \kappa ,
\]
then
\[
  \rho_F(A,P_AA)\le \tau+\kappa .
\]
If in addition \(A_k\) is a supplied best rank-\(k\) approximation certificate
and
\[
  \tau+\kappa \le \rho\,\rho_F(A,A_k),
\]
then
\[
  \rho_F(A,P_AA)\le \rho\,\rho_F(A,A_k).
\]

This is an exact-object residual-surface theorem.  It propagates the displayed
diagonal singular-value hypotheses into the concrete equation-(9) residual
and relative-residual wrappers.  It still assumes the tail and coupling
radii, and the relative form still assumes a supplied best-rank certificate.
It does not construct the rectangular SVD, prove source-tail norm bounds,
prove Eckart--Young optimality, derive randomness certificates, or certify
computed SVD, singular-vector, projector, Gram, inverse, coefficient, or
product routines.  Sampling probabilities and laws remain exact mathematical
inputs.
\end{corollary}

\begin{formalized}
\leanfit{columnSketchGramInverseProjector_sourceSVD_diagonal_headTailRankResidualSurface}
\leanfit{columnSketchGramInverseProjector_sourceSVD_diagonal_headTailRelativeResidualSurface}
\leanfit{columnSketchGramInverseProjector_sourceSVD_diagonal_columnSketchHeadRankResidualSurface}
\leanfit{columnSketchGramInverseProjector_sourceSVD_diagonal_columnSketchHeadRelativeResidualSurface}
\leanfit{columnSketchGramInverseProjector_sourceHeadTail_sourceSketchCoefficientRankResidualSurface_of_sourceSVD_diagonal_det_factors}
\leanfit{columnSketchGramInverseProjector_sourceHeadTail_sourceSketchCoefficientRelativeResidualSurface_of_sourceSVD_diagonal_det_factors}
\end{formalized}

\begin{corollary}[Diagonal source-SVD scalar tail rate for equation~(9)]
Let
\[
  A=H+T,\qquad H=U\Sigma_h V^T,\qquad U^T U=I_r,
  \qquad \det(V^T Z)\ne0,
\]
and let the exact source-head singular block be diagonal with nonzero displayed
entries,
\[
  (\Sigma_h)_{\alpha\beta}=
  \begin{cases}
    \sigma_\alpha, & \alpha=\beta,\\
    0, & \alpha\ne\beta ,
  \end{cases}
  \qquad
  \sigma_\alpha\ne0 \quad\text{for every }\alpha .
\]
Assume the exact source tail has the displayed SVD-style factorization
\[
  T=U_\perp\Sigma_\perp V_\perp^T,
  \qquad U_\perp^T U_\perp=I_q,
  \qquad U^T U_\perp=0 .
\]
Assume also the exact right-basis fields
\[
  V_\perp^T V_\perp=I_q,\qquad
  V^T V_\perp=0,\qquad
  V_\perp^T V=0,\qquad
  V^T V=I_r,
\]
together with row completeness of the concatenated right-basis block
\([V_\perp,V]\).  If \(\varepsilon\ge0\) and the exact CACM cross-term
certificate
\[
  \bigl\|\Sigma_\perp (V_\perp^T Z)(V^T Z)^{-1}\bigr\|_F
  \le
  \varepsilon\,\|\Sigma_\perp\|_F
\]
is supplied, then for the exact Gram-inverse column-sketch projector
\[
  P_A:=AZ\big((AZ)^T(AZ)\big)^{-1}(AZ)^T
\]
the diagonal source-head determinant route and the source-tail/coupling
certificates give
\[
  \operatorname{rank}(P_AA)\le r,
  \qquad
  \rho_F(A,P_AA)
  \le
  2\sqrt{1+\varepsilon^2}\,\|\Sigma_\perp\|_F .
\]
If \(A_k\) is a supplied best rank-\(k\) approximation certificate and
\[
  2\sqrt{1+\varepsilon^2}\,\|\Sigma_\perp\|_F
  \le
  \rho\,\rho_F(A,A_k),
\]
then
\[
  \rho_F(A,P_AA)\le \rho\,\rho_F(A,A_k).
\]

This is an exact-object scalar-rate assembly theorem.  It removes the two
opaque tail and coupling radii from the preceding source-head/tail residual
surface by instantiating both of them from the exact source-tail factorization
and exact cross-term certificate.  It still does not construct the rectangular
SVD/source-tail factors, prove singular-value ordering or Eckart--Young
optimality, derive the cross-term certificate from randomness, or certify
computed SVD, singular-vector, projector, Gram, inverse, coefficient, or
product routines.  Sampling probabilities and laws remain exact mathematical
inputs.
\end{corollary}

\begin{formalized}
\leanfit{columnSketchGramInverseProjector_sourceHeadTail_sourceSVDTailRankResidualSurface_of_sourceSVD_diagonal_crossTerm}
\leanfit{columnSketchGramInverseProjector_sourceHeadTail_sourceSVDTailRelativeResidualSurface_of_sourceSVD_diagonal_crossTerm}
\end{formalized}

\begin{corollary}[Source-tail-optimal scalar relative rate for equation~(9)]
Under the hypotheses of the preceding corollary, write
\[
  H=U\Sigma_h V^T,\qquad
  T=U_\perp\Sigma_\perp V_\perp^T,
\]
and assume in addition the supplied Frobenius tail-optimality inequality
\[
  \|T\|_F \le \rho_F(A,B)
  \qquad\text{for every } B \text{ with } \operatorname{rank}(B)\le r .
\]
Assume also the scalar comparison
\[
  2\sqrt{1+\varepsilon^2}\,\|\Sigma_\perp\|_F
  \le
  \rho\,\|T\|_F .
\]
Then the exact source head is a best rank-\(r\) Frobenius approximation
certificate,
\[
  H \in \operatorname*{argmin}_{\operatorname{rank}(B)\le r}\rho_F(A,B),
  \qquad
  \rho_F(A,H)=\|T\|_F .
\]
For the exact Gram-inverse column-sketch projector
\[
  P_A:=AZ\big((AZ)^T(AZ)\big)^{-1}(AZ)^T,
\]
the same diagonal source-head determinant route, source-tail/coupling
certificates, and tail-optimality handoff give
\[
  \operatorname{rank}(P_AA)\le r,
  \qquad
  \rho_F(A,P_AA)\le \rho\,\rho_F(A,H).
\]

This is an exact-object D4 handoff theorem.  It composes the scalar tail-rate
surface with a supplied tail-optimality inequality, so the old separate
\(A_k\)/best-rank hypothesis is replaced by the exact source head \(H\).  It
still does not prove the tail-optimality inequality from singular values or
Eckart--Young, construct the rectangular SVD/source-tail factors, derive the
cross-term certificate from randomness, or certify computed SVD,
singular-vector, projector, Gram, inverse, coefficient, or product routines.
Sampling probabilities and laws remain exact mathematical inputs.
\end{corollary}

\begin{formalized}
\leanfit{columnSketchGramInverseProjector_sourceHeadTail_sourceSVDTailRelativeResidualSurface_of_sourceSVD_diagonal_crossTerm_tailOptimal}
\end{formalized}

\begin{corollary}[Diagonal source-SVD certificate handoff for equation~(9)]
Define a diagonal source-SVD tail certificate
\[
  \mathsf{DiagTail}(A,T,U,\Sigma_h,\sigma,V,
    U_\perp,\Sigma_\perp,V_\perp)
\]
to consist of the following exact identities and orthogonality fields:
\[
\begin{aligned}
  &A=U\Sigma_hV^T+T,\qquad
    T=U_\perp\Sigma_\perp V_\perp^T,\\
  &U_\perp^TU_\perp=I,\qquad U^TU_\perp=0,\qquad U^TU=I,\\
  &(\Sigma_h)_{ab}=\begin{cases}\sigma_a,&a=b,\\0,&a\ne b,\end{cases}
    \qquad \sigma_a\ne0\quad\text{for all }a,\\
  &V_\perp^TV_\perp=I,\qquad V^TV_\perp=0,\qquad
    V_\perp^TV=0,\qquad V^TV=I,\\
  &V_\perp V_\perp^T+VV^T=I_n .
\end{aligned}
\]
This certificate implies the source-tail left-orthogonality field \(U^TT=0\).
If the supplied Frobenius tail-optimality inequality
\[
  \|T\|_F\le \rho_F(A,B)
  \qquad\text{for every }B\text{ with }\operatorname{rank}(B)\le r
\]
is also available, then it supplies the best-rank certificate for
\(H=U\Sigma_hV^T\).

Assume \(\varepsilon\ge0\),
\[
  \det(V^TZ)\ne0,\qquad
  \|\Sigma_\perp(V_\perp^TZ)(V^TZ)^{-1}\|_F
  \le \varepsilon\,\|\Sigma_\perp\|_F .
\]
For the exact Gram-inverse column-sketch projector
\[
  P_A:=AZ\big((AZ)^T(AZ)\big)^{-1}(AZ)^T,
\]
the certificate gives
\[
  \operatorname{rank}(P_AA)\le r,\qquad
  \rho_F(A,P_AA)\le
    2\sqrt{1+\varepsilon^2}\,\|\Sigma_\perp\|_F .
\]
If, in addition,
\[
  2\sqrt{1+\varepsilon^2}\,\|\Sigma_\perp\|_F
  \le \rho\,\|T\|_F,
\]
then
\[
  H\in\operatorname*{argmin}_{\operatorname{rank}(B)\le r}\rho_F(A,B),
  \qquad
  \rho_F(A,H)=\|T\|_F,\qquad
  \rho_F(A,P_AA)\le \rho\,\rho_F(A,H).
\]

This is an exact-object D3/D4 certificate handoff.  It packages the source-SVD
data used by the previous two scalar-rate corollaries into one mathematical
certificate, but it still does not construct that certificate from a
rectangular SVD, prove singular-value positivity/order or Eckart--Young
tail-optimality, derive the cross-term certificate from randomness, or certify
computed SVD, singular-vector, projector, Gram, inverse, coefficient, or
product routines.  Sampling probabilities and laws remain exact mathematical
inputs.
\end{corollary}

\begin{formalized}
\leanfit{DiagonalSourceSVDTailCertificate}
\leanfit{DiagonalSourceSVDTailCertificate.sourceTailLeftOrthogonal}
\leanfit{DiagonalSourceSVDTailCertificate.isBestRankApproxFrob_of_tail_optimal}
\leanfit{columnSketchGramInverseProjector_sourceHeadTail_sourceSVDTailRankResidualSurface_of_diagonalSourceSVDTailCertificate}
\leanfit{columnSketchGramInverseProjector_sourceHeadTail_sourceSVDTailRelativeResidualSurface_of_diagonalSourceSVDTailCertificate}
\end{formalized}

\begin{corollary}[Block source-SVD certificate constructor for equation~(9)]
Let
\[
  L=[U,U_\perp],\qquad R=[V_\perp,V]
\]
be the concatenated left and right source blocks.  Assume the exact block
decomposition
\[
  A=U\Sigma_hV^T+U_\perp\Sigma_\perp V_\perp^T,
\]
the exact block orthonormality and completeness conditions
\[
  L^TL=I_{r+q},\qquad R^TR=I_{q+r},\qquad RR^T=I_n,
\]
and the exact diagonal nonsingular head block
\[
  (\Sigma_h)_{ab}=\begin{cases}\sigma_a,&a=b,\\0,&a\ne b,\end{cases}
  \qquad \sigma_a\ne0\quad\text{for all }a.
\]
Then, with
\[
  T:=U_\perp\Sigma_\perp V_\perp^T,
\]
these block hypotheses construct the diagonal source-SVD tail certificate
\[
  \mathsf{DiagTail}(A,T,U,\Sigma_h,\sigma,V,
    U_\perp,\Sigma_\perp,V_\perp).
\]
Equivalently, the left block certificate gives
\[
  U^TU=I,\qquad U^TU_\perp=0,\qquad U_\perp^TU_\perp=I,
\]
and the right block certificate gives
\[
  V_\perp^TV_\perp=I,\qquad V^TV_\perp=0,\qquad
  V_\perp^TV=0,\qquad V^TV=I,\qquad
  V_\perp V_\perp^T+VV^T=I_n.
\]
Thus the source-tail left-orthogonality field \(U^TT=0\) follows from the
block data.

Assume further that \(\varepsilon\ge0\),
\[
  \det(V^TZ)\ne0,\qquad
  \|\Sigma_\perp(V_\perp^TZ)(V^TZ)^{-1}\|_F
  \le \varepsilon\,\|\Sigma_\perp\|_F .
\]
For the exact Gram-inverse column-sketch projector
\[
  P_A:=AZ\big((AZ)^T(AZ)\big)^{-1}(AZ)^T,
\]
the constructed certificate gives
\[
  \operatorname{rank}(P_AA)\le r,\qquad
  \rho_F(A,P_AA)\le
    2\sqrt{1+\varepsilon^2}\,\|\Sigma_\perp\|_F .
\]
If the supplied Frobenius tail-optimality inequality
\[
  \|U_\perp\Sigma_\perp V_\perp^T\|_F\le \rho_F(A,B)
  \qquad\text{for every }B\text{ with }\operatorname{rank}(B)\le r
\]
and the scalar comparison
\[
  2\sqrt{1+\varepsilon^2}\,\|\Sigma_\perp\|_F
  \le
  \rho\,\|U_\perp\Sigma_\perp V_\perp^T\|_F
\]
are also available, then
\[
  U\Sigma_hV^T\in
  \operatorname*{argmin}_{\operatorname{rank}(B)\le r}\rho_F(A,B),
  \qquad
  \rho_F(A,U\Sigma_hV^T)=\|U_\perp\Sigma_\perp V_\perp^T\|_F,
\]
and
\[
  \rho_F(A,P_AA)\le
  \rho\,\rho_F(A,U\Sigma_hV^T).
\]

This is an exact-object D3 constructor for the preceding certificate theorem.
It derives the certificate from primitive block decomposition and block
orthonormality data, but it still does not prove rectangular SVD existence,
singular-value ordering or positivity from an SVD construction, Eckart--Young
tail optimality, randomness-derived cross-term certificates, or computed
non-probability SVD, singular-vector, projector, Gram, inverse, coefficient, or
product routines.  Sampling probabilities and laws remain exact mathematical
inputs.
\end{corollary}

\begin{formalized}
\leanfit{leftBasisBlock}
\leanfit{leftBasisBlock_component_orthonormal_fields_of_col_orthonormal}
\leanfit{BlockDiagonalSourceSVDTailCertificate}
\leanfit{BlockDiagonalSourceSVDTailCertificate.to_diagonalSourceSVDTailCertificate}
\leanfit{BlockDiagonalSourceSVDTailCertificate.sourceTailLeftOrthogonal}
\leanfit{BlockDiagonalSourceSVDTailCertificate.isBestRankApproxFrob_of_tail_optimal}
\leanfit{columnSketchGramInverseProjector_sourceHeadTail_sourceSVDTailRankResidualSurface_of_blockDiagonalSourceSVDTailCertificate}
\leanfit{columnSketchGramInverseProjector_sourceHeadTail_sourceSVDTailRelativeResidualSurface_of_blockDiagonalSourceSVDTailCertificate}
\end{formalized}

\begin{corollary}[Square SVD split constructor for equation~(9)]
Let \(A\in\mathbb R^{(r+q)\times(r+q)}\).  Suppose exact square SVD-style data
are supplied:
\[
  U_{\rm full}^TU_{\rm full}=I,\qquad
  V_{\rm full}^TV_{\rm full}=I,\qquad
  U_{\rm full}U_{\rm full}^T=I,\qquad
  V_{\rm full}V_{\rm full}^T=I,
\]
and
\[
  A_{ij}
  =
  \sum_{\ell=1}^{r+q}
    (U_{\rm full})_{i\ell}\,
    \sigma_\ell\,
    (V_{\rm full})_{j\ell}.
\]
Let \(h(a)\) denote the head index \(a\in\{1,\ldots,r\}\) embedded in
\(\{1,\ldots,r+q\}\), and let \(t(c)\) denote the tail index
\(r+c\).  Define
\[
  U_a=(U_{\rm full})_{\cdot,h(a)},\qquad
  U_{\perp,c}=(U_{\rm full})_{\cdot,t(c)},\qquad
  V_a=(V_{\rm full})_{\cdot,h(a)},\qquad
  V_{\perp,c}=(V_{\rm full})_{\cdot,t(c)},
\]
and diagonal blocks
\[
  (\Sigma_h)_{ab}=
  \begin{cases}\sigma_{h(a)},&a=b,\\0,&a\ne b,\end{cases}
  \qquad
  (\Sigma_\perp)_{cd}=
  \begin{cases}\sigma_{t(c)},&c=d,\\0,&c\ne d.\end{cases}
\]
If the displayed head singular entries are nonzero,
\[
  \sigma_{h(a)}\ne0\qquad\text{for all }a,
\]
then the split exact SVD data construct
\[
  \mathsf{BlockTail}
    (A,U,\Sigma_h,\sigma_h,V,U_\perp,\Sigma_\perp,V_\perp),
\]
where
\[
  A=U\Sigma_hV^T+U_\perp\Sigma_\perp V_\perp^T.
\]
In particular, the square SVD representation proves the head and tail
expansions
\[
  U\Sigma_hV^T
  =
  \sum_{a=1}^r
    (U_{\rm full})_{\cdot,h(a)}\,
    \sigma_{h(a)}\,
    (V_{\rm full})_{\cdot,h(a)}^T,
\]
and
\[
  U_\perp\Sigma_\perp V_\perp^T
  =
  \sum_{c=1}^q
    (U_{\rm full})_{\cdot,t(c)}\,
    \sigma_{t(c)}\,
    (V_{\rm full})_{\cdot,t(c)}^T.
\]

Consequently, if \(\varepsilon\ge0\),
\[
  \det(V^TZ)\ne0,\qquad
  \|\Sigma_\perp(V_\perp^TZ)(V^TZ)^{-1}\|_F
  \le \varepsilon\,\|\Sigma_\perp\|_F ,
\]
then the exact Gram-inverse column-sketch projector
\[
  P_A=AZ\big((AZ)^T(AZ)\big)^{-1}(AZ)^T
\]
satisfies
\[
  \operatorname{rank}(P_AA)\le r,\qquad
  \rho_F(A,P_AA)\le
    2\sqrt{1+\varepsilon^2}\,\|\Sigma_\perp\|_F .
\]
If, in addition, the supplied Frobenius tail-optimality inequality
\[
  \|U_\perp\Sigma_\perp V_\perp^T\|_F
  \le \rho_F(A,B)
  \qquad\text{for every }B\text{ with }\operatorname{rank}(B)\le r
\]
and the scalar comparison
\[
  2\sqrt{1+\varepsilon^2}\,\|\Sigma_\perp\|_F
  \le
  \rho\,\|U_\perp\Sigma_\perp V_\perp^T\|_F
\]
hold, then
\[
  U\Sigma_hV^T\in
  \operatorname*{argmin}_{\operatorname{rank}(B)\le r}\rho_F(A,B),
  \qquad
  \rho_F(A,U\Sigma_hV^T)
  =\|U_\perp\Sigma_\perp V_\perp^T\|_F,
\]
and
\[
  \rho_F(A,P_AA)
  \le
  \rho\,\rho_F(A,U\Sigma_hV^T).
\]

This is an exact-object square-specialized D3 constructor.  It does not prove
existence of the SVD, singular-value order or positivity beyond the visible
head-nonzero hypothesis, Eckart--Young tail optimality, randomness-derived
cross-term certificates, or computed non-probability SVD, singular-vector,
projector, Gram, inverse, coefficient, or product routines.  Sampling
probabilities and laws remain exact mathematical inputs.
\end{corollary}

\begin{formalized}
\leanfit{squareSVDHeadLeft}
\leanfit{squareSVDTailLeft}
\leanfit{squareSVDHeadRight}
\leanfit{squareSVDTailRight}
\leanfit{squareSVDHeadDiagonal}
\leanfit{squareSVDTailDiagonal}
\leanfit{squareSVDHeadValues}
\leanfit{sourceSVDFactorMatrix_squareSVDHeadDiagonal}
\leanfit{sourceSVDFactorMatrix_squareSVDTailDiagonal}
\leanfit{BlockDiagonalSourceSVDTailCertificate.of_squareSVD}
\leanfit{columnSketchGramInverseProjector_sourceHeadTail_sourceSVDTailRankResidualSurface_of_squareSVD}
\leanfit{columnSketchGramInverseProjector_sourceHeadTail_sourceSVDTailRelativeResidualSurface_of_squareSVD}
\end{formalized}

\begin{corollary}[Rectangular-thin SVD split constructor for equation~(9)]
Let \(A\in\mathbb R^{m\times(r+q)}\).  Suppose exact thin rectangular
SVD-style data are supplied:
\[
  \sum_{i=1}^m (U_{\rm full})_{ia}(U_{\rm full})_{ib}
  =
  \begin{cases}1,&a=b,\\0,&a\ne b,\end{cases}
\]
for all \(a,b\in\{1,\ldots,r+q\}\), and suppose the right table is exactly
orthogonal,
\[
  V_{\rm full}^TV_{\rm full}=I,\qquad
  V_{\rm full}V_{\rm full}^T=I.
\]
Assume the exact representation
\[
  A_{ij}
  =
  \sum_{\ell=1}^{r+q}
    (U_{\rm full})_{i\ell}\,
    \sigma_\ell\,
    (V_{\rm full})_{j\ell}.
\]
Using the same head indices \(h(a)\) and tail indices \(t(c)\) as in the
square-SVD split, define
\[
  U_a=(U_{\rm full})_{\cdot,h(a)},\qquad
  U_{\perp,c}=(U_{\rm full})_{\cdot,t(c)},\qquad
  V_a=(V_{\rm full})_{\cdot,h(a)},\qquad
  V_{\perp,c}=(V_{\rm full})_{\cdot,t(c)},
\]
and the same diagonal blocks
\[
  (\Sigma_h)_{ab}=
  \begin{cases}\sigma_{h(a)},&a=b,\\0,&a\ne b,\end{cases}
  \qquad
  (\Sigma_\perp)_{cd}=
  \begin{cases}\sigma_{t(c)},&c=d,\\0,&c\ne d.\end{cases}
\]
If the displayed head singular entries are nonzero,
\[
  \sigma_{h(a)}\ne0\qquad\text{for all }a,
\]
then these thin rectangular data construct
\[
  \mathsf{BlockTail}
    (A,U,\Sigma_h,\sigma_h,V,U_\perp,\Sigma_\perp,V_\perp),
\]
where
\[
  A=U\Sigma_hV^T+U_\perp\Sigma_\perp V_\perp^T.
\]
The head and tail expansions are
\[
  U\Sigma_hV^T
  =
  \sum_{a=1}^r
    (U_{\rm full})_{\cdot,h(a)}\,
    \sigma_{h(a)}\,
    (V_{\rm full})_{\cdot,h(a)}^T,
\]
and
\[
  U_\perp\Sigma_\perp V_\perp^T
  =
  \sum_{c=1}^q
    (U_{\rm full})_{\cdot,t(c)}\,
    \sigma_{t(c)}\,
    (V_{\rm full})_{\cdot,t(c)}^T.
\]

Consequently, if \(\varepsilon\ge0\),
\[
  \det(V^TZ)\ne0,\qquad
  \|\Sigma_\perp(V_\perp^TZ)(V^TZ)^{-1}\|_F
  \le \varepsilon\,\|\Sigma_\perp\|_F ,
\]
then the exact Gram-inverse column-sketch projector
\[
  P_A=AZ\big((AZ)^T(AZ)\big)^{-1}(AZ)^T
\]
satisfies
\[
  \operatorname{rank}(P_AA)\le r,\qquad
  \rho_F(A,P_AA)\le
    2\sqrt{1+\varepsilon^2}\,\|\Sigma_\perp\|_F .
\]
If the supplied Frobenius tail-optimality inequality
\[
  \|U_\perp\Sigma_\perp V_\perp^T\|_F
  \le \rho_F(A,B)
  \qquad\text{for every }B\text{ with }\operatorname{rank}(B)\le r
\]
and the scalar comparison
\[
  2\sqrt{1+\varepsilon^2}\,\|\Sigma_\perp\|_F
  \le
  \rho\,\|U_\perp\Sigma_\perp V_\perp^T\|_F
\]
hold, then
\[
  U\Sigma_hV^T\in
  \operatorname*{argmin}_{\operatorname{rank}(B)\le r}\rho_F(A,B),
  \qquad
  \rho_F(A,U\Sigma_hV^T)
  =\|U_\perp\Sigma_\perp V_\perp^T\|_F,
\]
and
\[
  \rho_F(A,P_AA)
  \le
  \rho\,\rho_F(A,U\Sigma_hV^T).
\]

This is an exact-object thin-rectangular D3 constructor.  It removes the
square-left restriction from the preceding corollary, but it still does not
prove rectangular SVD existence, singular-value order or positivity beyond the
visible head-nonzero hypothesis, Eckart--Young tail optimality,
randomness-derived cross-term certificates, or computed non-probability SVD,
singular-vector, projector, Gram, inverse, coefficient, or product routines.
Sampling probabilities and laws remain exact mathematical inputs.
\end{corollary}

\begin{formalized}
\leanfit{rectangularThinSVDHeadLeft}
\leanfit{rectangularThinSVDTailLeft}
\leanfit{sourceSVDFactorMatrix_rectangularThinSVDHeadDiagonal}
\leanfit{sourceSVDFactorMatrix_rectangularThinSVDTailDiagonal}
\leanfit{BlockDiagonalSourceSVDTailCertificate.of_rectangularThinSVD}
\leanfit{columnSketchGramInverseProjector_sourceHeadTail_sourceSVDTailRankResidualSurface_of_rectangularThinSVD}
\leanfit{columnSketchGramInverseProjector_sourceHeadTail_sourceSVDTailRelativeResidualSurface_of_rectangularThinSVD}
\end{formalized}

\begin{corollary}[Head-positive SVD split constructors for equation~(9)]
In either of the two preceding supplied SVD-style settings, let
\[
  h(a)\in\{1,\ldots,r+q\},\qquad a=1,\ldots,r,
\]
denote the displayed head singular-value indices.  Assume the source-style
strict positivity condition
\[
  \sigma_{h(a)}>0\qquad\text{for every }a=1,\ldots,r.
\]
Then the split head-value table is strictly positive,
\[
  (\sigma_h)_a>0\qquad\text{for every }a,
\]
and in particular the nonzero-head field required by the diagonal source-SVD
constructors holds:
\[
  \sigma_{h(a)}\ne0\qquad\text{for every }a.
\]

Consequently, in the square SVD setting, the exact hypotheses
\[
  U_{\rm full}^TU_{\rm full}=I,\quad
  U_{\rm full}U_{\rm full}^T=I,\quad
  V_{\rm full}^TV_{\rm full}=I,\quad
  V_{\rm full}V_{\rm full}^T=I,\quad
  A=U_{\rm full}\operatorname{diag}(\sigma)V_{\rm full}^T
\]
together with \(\sigma_{h(a)}>0\) construct
\[
  \mathsf{BlockTail}
    (A,U,\Sigma_h,\sigma_h,V,U_\perp,\Sigma_\perp,V_\perp).
\]
The same conclusion holds in the thin rectangular setting from
\[
  U_{\rm full}^TU_{\rm full}=I,\qquad
  V_{\rm full}^TV_{\rm full}=I,\qquad
  V_{\rm full}V_{\rm full}^T=I,\qquad
  A=U_{\rm full}\operatorname{diag}(\sigma)V_{\rm full}^T,
\]
where \(U_{\rm full}\in\mathbb R^{m\times(r+q)}\) has exact orthonormal
columns.

In both settings, if \(\varepsilon\ge0\),
\[
  \det(V^TZ)\ne0,\qquad
  \|\Sigma_\perp(V_\perp^TZ)(V^TZ)^{-1}\|_F
  \le \varepsilon\,\|\Sigma_\perp\|_F,
\]
then the exact Gram-inverse column-sketch projector satisfies
\[
  \operatorname{rank}(P_AA)\le r,\qquad
  \rho_F(A,P_AA)\le
  2\sqrt{1+\varepsilon^2}\,\|\Sigma_\perp\|_F .
\]
If the supplied Frobenius tail-optimality inequality and scalar comparison
from the preceding two corollaries are also available, then the same
tail-optimal relative conclusion follows:
\[
  U\Sigma_hV^T\in
  \operatorname*{argmin}_{\operatorname{rank}(B)\le r}\rho_F(A,B),
  \qquad
  \rho_F(A,P_AA)\le
  \rho\,\rho_F(A,U\Sigma_hV^T).
\]

This is an exact-object D3/D4 positivity handoff.  It closes only the passage
from positive displayed leading singular values to the nonzero singular block
used by the determinant and split-constructor theorems.  It still does not
prove rectangular SVD existence, singular-value ordering, Eckart--Young tail
optimality, randomness-derived cross-term certificates, or computed
non-probability SVD, singular-vector, projector, Gram, inverse, coefficient, or
product routines.  Sampling probabilities and laws remain exact mathematical
inputs.
\end{corollary}

\begin{formalized}
\leanfit{squareSVDHeadValues_pos}
\leanfit{squareSVDHeadValues_nonzero_of_pos}
\leanfit{BlockDiagonalSourceSVDTailCertificate.of_squareSVD_head_pos}
\leanfit{columnSketchGramInverseProjector_sourceHeadTail_sourceSVDTailRankResidualSurface_of_squareSVD_head_pos}
\leanfit{columnSketchGramInverseProjector_sourceHeadTail_sourceSVDTailRelativeResidualSurface_of_squareSVD_head_pos}
\leanfit{BlockDiagonalSourceSVDTailCertificate.of_rectangularThinSVD_head_pos}
\leanfit{columnSketchGramInverseProjector_sourceHeadTail_sourceSVDTailRankResidualSurface_of_rectangularThinSVD_head_pos}
\leanfit{columnSketchGramInverseProjector_sourceHeadTail_sourceSVDTailRelativeResidualSurface_of_rectangularThinSVD_head_pos}
\end{formalized}

\begin{corollary}[Source-SVD tail Frobenius norm identity for equation~(9)]
Let \(U\in\mathbb R^{m\times q}\), \(V\in\mathbb R^{n\times q}\), and
\(\Sigma\in\mathbb R^{q\times q}\) be exact analysis objects satisfying
\[
  U^T U=I_q,\qquad V^T V=I_q .
\]
Then the exact source-SVD factor table satisfies
\[
  \|U\Sigma V^T\|_F^2=\|\Sigma\|_F^2,
  \qquad
  \|U\Sigma V^T\|_F=\|\Sigma\|_F .
\]

Consequently, suppose the block source-SVD certificate of the previous
equation-(9) route supplies
\[
  A=U\Sigma_hV^T+U_\perp\Sigma_\perp V_\perp^T
\]
with exact block orthonormality for the displayed left and right factors.
Then the source-tail matrix \(T=U_\perp\Sigma_\perp V_\perp^T\) satisfies
\[
  \|T\|_F=\|\Sigma_\perp\|_F,\qquad
  \rho_F(A,U\Sigma_hV^T)=\|\Sigma_\perp\|_F .
\]
Moreover, if \(\varepsilon\ge0\), \(\det(V^TZ)\ne0\), and
\[
  \|\Sigma_\perp(V_\perp^TZ)(V^TZ)^{-1}\|_F
  \le \varepsilon\,\|\Sigma_\perp\|_F,
\]
and if the supplied tail-optimality and scalar-comparison hypotheses are
already stated with the singular-tail norm,
\[
  \|\Sigma_\perp\|_F\le \rho_F(A,B)
  \quad\text{for all } \operatorname{rank}(B)\le r,
\]
\[
  2\sqrt{1+\varepsilon^2}\,\|\Sigma_\perp\|_F
  \le \rho\,\|\Sigma_\perp\|_F,
\]
then the exact Gram-inverse column-sketch projector satisfies
\[
  U\Sigma_hV^T\in
  \operatorname*{argmin}_{\operatorname{rank}(B)\le r}\rho_F(A,B),
  \qquad
  \operatorname{rank}(P_AA)\le r,
\]
\[
  \rho_F(A,P_AA)\le
  \rho\,\rho_F(A,U\Sigma_hV^T).
\]
This is still an exact-object theorem: it does not construct a rectangular
SVD, prove singular-value ordering, prove Eckart--Young tail optimality from
singular values, derive the cross-term radius from randomness, or certify
computed non-probability SVD, singular-vector, projector, Gram, inverse, or
product routines.  Sampling probabilities and laws remain exact mathematical
inputs.
\end{corollary}

\begin{formalized}
\leanfit{frobNormSqRect_sourceSVDFactorMatrix_orthonormal}
\leanfit{frobNormRect_sourceSVDFactorMatrix_orthonormal}
\leanfit{BlockDiagonalSourceSVDTailCertificate.tail_frobNorm_eq_sigma}
\leanfit{BlockDiagonalSourceSVDTailCertificate.tail_lowRankResidual_eq_sigma}
\leanfit{columnSketchGramInverseProjector_sourceHeadTail_sourceSVDTailRelativeResidualSurface_of_blockDiagonalSourceSVDTailCertificate_sigmaTail}
\end{formalized}

\begin{corollary}[Sigma-tail square and thin-rectangular SVD wrappers for equation~(9)]
Let \(r,q\in\mathbb N\).  In the square case, suppose
\[
  A=U_{\mathrm{full}}\operatorname{diag}(\sigma)V_{\mathrm{full}}^T,
  \qquad
  U_{\mathrm{full}}^TU_{\mathrm{full}}=
  V_{\mathrm{full}}^TV_{\mathrm{full}}=I_{r+q},
\]
and every displayed head singular entry is nonzero.  Let
\(\Sigma_{\mathrm{tail}}\) be the diagonal block
\(\operatorname{diag}(\sigma_{r+1},\ldots,\sigma_{r+q})\).  Then the split
tail and source-head residual satisfy
\[
  \|U_{\mathrm{tail}}\Sigma_{\mathrm{tail}}V_{\mathrm{tail}}^T\|_F
  =\|\Sigma_{\mathrm{tail}}\|_F,
  \qquad
  \rho_F(A,U_{\mathrm{head}}\Sigma_{\mathrm{head}}V_{\mathrm{head}}^T)
  =\|\Sigma_{\mathrm{tail}}\|_F .
\]
Consequently, if \(\varepsilon\ge0\), \(\det(V_{\mathrm{head}}^TZ)\ne0\),
\[
  \|\Sigma_{\mathrm{tail}}(V_{\mathrm{tail}}^TZ)
     (V_{\mathrm{head}}^TZ)^{-1}\|_F
  \le \varepsilon\|\Sigma_{\mathrm{tail}}\|_F,
\]
and
\[
  \|\Sigma_{\mathrm{tail}}\|_F\le \rho_F(A,B)
  \quad\text{for all }\operatorname{rank}(B)\le r,\qquad
  2\sqrt{1+\varepsilon^2}\|\Sigma_{\mathrm{tail}}\|_F
  \le \rho\|\Sigma_{\mathrm{tail}}\|_F,
\]
then the exact Gram-inverse column-sketch projector obeys
\[
  U_{\mathrm{head}}\Sigma_{\mathrm{head}}V_{\mathrm{head}}^T
  \in\operatorname*{argmin}_{\operatorname{rank}(B)\le r}\rho_F(A,B),
  \qquad
  \rho_F(A,P_AA)\le
  \rho\,\rho_F(A,U_{\mathrm{head}}\Sigma_{\mathrm{head}}V_{\mathrm{head}}^T).
\]
The same conclusion holds for the thin rectangular case
\[
  A=U_{\mathrm{full}}\operatorname{diag}(\sigma)V_{\mathrm{full}}^T,
  \qquad
  U_{\mathrm{full}}^TU_{\mathrm{full}}=I_{r+q},\quad
  V_{\mathrm{full}}^TV_{\mathrm{full}}=I_{r+q},
\]
with \(A\in\mathbb R^{m\times(r+q)}\).  In both cases, a strict-positive head
condition \(0<\sigma_a\) may replace the nonzero-head hypothesis.

This corollary is exact-object algebra for supplied square or thin-rectangular
SVD-style data.  It does not construct the SVD, prove singular-value ordering,
prove Eckart--Young tail optimality, derive the cross-term radius from
randomness, or certify computed non-probability SVD, singular-vector,
projector, Gram, inverse, or product routines.  Sampling probabilities and laws
remain exact mathematical inputs.
\end{corollary}

\begin{formalized}
\leanfit{frobNormRect_squareSVDTail_eq_sigmaTail}
\leanfit{lowRankResidualFrob_squareSVDHead_eq_sigmaTail}
\leanfit{columnSketchGramInverseProjector_sourceHeadTail_sourceSVDTailRelativeResidualSurface_of_squareSVD_sigmaTail}
\leanfit{columnSketchGramInverseProjector_sourceHeadTail_sourceSVDTailRelativeResidualSurface_of_squareSVD_head_pos_sigmaTail}
\leanfit{frobNormRect_rectangularThinSVDTail_eq_sigmaTail}
\leanfit{lowRankResidualFrob_rectangularThinSVDHead_eq_sigmaTail}
\leanfit{columnSketchGramInverseProjector_sourceHeadTail_sourceSVDTailRelativeResidualSurface_of_rectangularThinSVD_sigmaTail}
\leanfit{columnSketchGramInverseProjector_sourceHeadTail_sourceSVDTailRelativeResidualSurface_of_rectangularThinSVD_head_pos_sigmaTail}
\end{formalized}

\begin{corollary}[Sigma-tail best-rank handoff for square and thin SVD data]
Under the square-SVD hypotheses of the previous corollary, assume the visible
tail-optimality inequality is stated directly with the displayed tail singular
block:
\[
  \|\Sigma_{\mathrm{tail}}\|_F\le \rho_F(A,B)
  \quad\text{for every }B\text{ with }\operatorname{rank}(B)\le r .
\]
Then the exact source head is a Frobenius best rank-\(r\) approximation:
\[
  U_{\mathrm{head}}\Sigma_{\mathrm{head}}V_{\mathrm{head}}^T
  \in \operatorname*{argmin}_{\operatorname{rank}(B)\le r}\rho_F(A,B).
\]
The same implication holds in the thin rectangular setting, and in both cases
the nonzero-head hypothesis may be replaced by the strict-positive head
hypothesis \(0<\sigma_a\).

This is only a D4 handoff from a visible Eckart--Young-style inequality to the
repository certificate \(\mathrm{IsBestRankApproxFrob}\).  It does not prove
SVD existence, singular-value ordering, the Eckart--Young theorem that supplies
the displayed inequality, randomness, or computed non-probability SVD,
singular-vector, projector, Gram, inverse, or product routines.  Sampling
probabilities and laws remain exact mathematical inputs.
\end{corollary}

\begin{formalized}
\leanfit{isBestRankApproxFrob_of_squareSVD_sigmaTail_optimal}
\leanfit{isBestRankApproxFrob_of_squareSVD_head_pos_sigmaTail_optimal}
\leanfit{isBestRankApproxFrob_of_rectangularThinSVD_sigmaTail_optimal}
\leanfit{isBestRankApproxFrob_of_rectangularThinSVD_head_pos_sigmaTail_optimal}
\end{formalized}

\begin{corollary}[Right-Gram singular-value order and nonnegativity for equation~(9)]
Let \(A\in\mathbb R^{m\times n}\) be an exact analysis matrix and define the
right Gram matrix
\[
  G=A^T A,\qquad
  G_{jk}=\sum_{i=1}^m A_{ij}A_{ik}.
\]
Then \(G\) is symmetric and positive semidefinite.  Equivalently, for every
\(x\in\mathbb R^n\),
\[
  x^TGx=\sum_{i=1}^m\left(\sum_{j=1}^n A_{ij}x_j\right)^2\ge0 .
\]
Let \(\lambda_j(A)\) be the ordered Hermitian eigenvalues of \(G\), indexed in
the repository's zero-based finite order, and define
\[
  \sigma_j(A)=\sqrt{\lambda_j(A)} .
\]
Then
\[
  0\le\lambda_j(A),\qquad
  j\le k\Longrightarrow \lambda_k(A)\le\lambda_j(A),
\]
\[
  0\le\sigma_j(A),\qquad
  j\le k\Longrightarrow \sigma_k(A)\le\sigma_j(A),
  \qquad
  \sigma_j(A)^2=\lambda_j(A).
\]

This is an exact spectral adapter for the equation-(9) SVD/Eckart--Young
route.  It supplies ordered nonnegative singular-value magnitudes from the
right Gram \(A^TA\), but it does not construct singular vectors, a rectangular
SVD source split, or the Eckart--Young tail-optimality theorem.  If an
implementation computes \(A^TA\), singular vectors, projectors, bases, or SVD
data, those are non-probability computed quantities and still require explicit
floating-point/inexact-arithmetic certificates.  Sampling probabilities and
sampling laws remain exact mathematical inputs by the current convention.
\end{corollary}

\begin{formalized}
\leanfit{rectRightGram}
\leanfit{rectRightGram_symmetric}
\leanfit{finiteQuadraticForm_rectRightGram_eq_sum_sq}
\leanfit{rectRightGram_finitePSD}
\leanfit{rectRightGram_matrix_posSemidef}
\leanfit{rectSingularValueSq}
\leanfit{rectSingularValue}
\leanfit{rectSingularValueSq_nonneg}
\leanfit{rectSingularValueSq_antitone}
\leanfit{rectSingularValue_nonneg}
\leanfit{rectSingularValue_antitone}
\leanfit{rectSingularValue_sq_eq}
\end{formalized}

\begin{corollary}[Right-Gram eigenvector table and diagonalization for equation~(9)]
Let \(A\in\mathbb R^{m\times n}\) be an exact analysis matrix and let
\[
  G=A^TA,\qquad G_{jk}=\sum_{i=1}^m A_{ij}A_{ik}.
\]
Let \(\alpha_\ell(A)\) denote the basis-indexed Hermitian eigenvalues attached
to a mathlib Hermitian eigenvector basis of \(G\), and let
\[
  V_{j\ell}(A)
\]
be the corresponding exact real eigenvector table.  Define the basis-indexed
singular values
\[
  \tau_\ell(A)=\sqrt{\alpha_\ell(A)}.
\]
Then \(V(A)\) is orthogonal:
\[
  \sum_{j=1}^n V_{ja}(A)V_{jb}(A)=\delta_{ab},
  \qquad
  \sum_{a=1}^n V_{ia}(A)V_{ja}(A)=\delta_{ij}.
\]
Moreover
\[
  0\le \alpha_a(A),\qquad 0\le \tau_a(A),
  \qquad \tau_a(A)^2=\alpha_a(A),
\]
and each column of \(V(A)\) is a right-Gram eigenvector:
\[
  \sum_{k=1}^n G_{jk}V_{ka}(A)
  =
  \alpha_a(A)V_{ja}(A).
\]
Consequently,
\[
  \sum_{j=1}^n V_{ja}(A)
    \left(\sum_{k=1}^nG_{jk}V_{kb}(A)\right)
  =
  \begin{cases}
    \alpha_a(A),& a=b,\\
    0,& a\ne b,
  \end{cases}
\]
equivalently the diagonal entries may be written as \(\tau_a(A)^2\).

This is an exact-object eigenvector adapter for the equation-(9)
SVD/Eckart--Young route.  It is basis-indexed: it uses the index order of
mathlib's eigenvector basis, not the ordered zero-indexed singular-value
sequence of the previous corollary.  Thus it still does not construct the
ordered head/tail rectangular SVD source split or prove the Eckart--Young
tail-optimality theorem.  If an implementation computes \(A^TA\), the
eigenvector table, singular vectors, projectors, bases, or SVD data, those are
non-probability computed quantities and require explicit
floating-point/inexact-arithmetic certificates.  Sampling probabilities and
sampling laws remain exact mathematical inputs by the current convention.
\end{corollary}

\begin{formalized}
\leanfit{rectRightGramEigenvalue}
\leanfit{rectRightGramEigenbasis}
\leanfit{rectRightGramBasisSingularValue}
\leanfit{rectRightGramEigenbasis_isOrthogonal}
\leanfit{rectRightGramEigenbasis_col_orthonormal}
\leanfit{rectRightGramEigenbasis_row_orthonormal}
\leanfit{rectRightGramEigenvalue_nonneg}
\leanfit{rectRightGramBasisSingularValue_nonneg}
\leanfit{rectRightGramBasisSingularValue_sq_eq}
\leanfit{rectRightGramEigenbasis_eigenvector}
\leanfit{rectRightGramEigenbasis_diagonalizes}
\leanfit{rectRightGramEigenbasis_diagonalizes_singularValueSq}
\end{formalized}

\begin{corollary}[Full-positive right-Gram SVD reconstruction for equation~(9)]
Let \(A\in\mathbb R^{m\times n}\), let \(G=A^TA\), and let
\((\alpha_a,V_{\cdot a},\tau_a)\) be the basis-indexed eigenvalue,
right-eigenvector, and singular-value data from the previous corollary:
\[
  \tau_a=\sqrt{\alpha_a},\qquad
  GV_{\cdot a}=\alpha_aV_{\cdot a},\qquad
  V^TV=I_n,\qquad VV^T=I_n.
\]
Assume the full-positive hypothesis
\[
  0<\tau_a\qquad\text{for every }a\in\{1,\ldots,n\}.
\]
Define
\[
  y_{ia}=(AV_{\cdot a})_i
       =\sum_{j=1}^n A_{ij}V_{ja},
  \qquad
  u_{ia}=\frac{y_{ia}}{\tau_a}.
\]
Then the projected columns have diagonal Gram matrix
\[
  \sum_{i=1}^m y_{ia}y_{ib}
  =
  \begin{cases}
    \tau_a^2,& a=b,\\
    0,& a\ne b,
  \end{cases}
\]
and the normalized left candidates are exactly column-orthonormal:
\[
  \sum_{i=1}^m u_{ia}u_{ib}=\delta_{ab}.
\]
Moreover, the exact row-completeness of \(V\) gives
\[
  A_{ij}
  =
  \sum_{a=1}^n y_{ia}V_{ja}
  =
  \sum_{a=1}^n u_{ia}\tau_a V_{ja}.
\]

This is an exact-object reconstruction theorem for the full-positive,
basis-indexed case.  It is not yet the ordered rank-deficient rectangular SVD
source split used by the paper's final Eckart--Young route: zero singular
values, ordered head/tail selection, tail optimality, and the rank-deficient
extension remain open.  If an implementation computes \(A^TA\), \(V\), \(U\),
singular values, projectors, bases, or SVD data, those are non-probability
computed quantities and require explicit floating-point/inexact-arithmetic
certificates.  Sampling probabilities and sampling laws remain exact
mathematical inputs by the current convention.
\end{corollary}

\begin{formalized}
\leanfit{rectRightGramProjectedColumn}
\leanfit{rectRightGramLeftSingularFromEigenbasis}
\leanfit{rectRightGramBasisSingularDiagonal}
\leanfit{rectRightGramProjectedColumn_dot}
\leanfit{rectRightGramProjectedColumn_dot_diagonal}
\leanfit{rectRightGramLeftSingularFromEigenbasis_col_orthonormal_of_pos}
\leanfit{rectRightGramProjectedColumn_reconstruct}
\leanfit{rectRightGramLeftSingularFromEigenbasis_factor_column_of_pos}
\leanfit{rectRightGram_fullPositive_basisSVD_representation}
\end{formalized}

\begin{corollary}[Zero right-Gram projected columns for equation~(9)]
Let \(A\in\mathbb R^{m\times n}\), let \(G=A^TA\), and let
\((\alpha_a,V_{\cdot a},\tau_a)\) be the basis-indexed right-Gram data above.
For
\[
  y_{ia}=(AV_{\cdot a})_i
       =\sum_{j=1}^n A_{ij}V_{ja},
\]
one has the exact norm-square identities
\[
  \sum_{i=1}^m y_{ia}^2=\tau_a^2=\alpha_a.
\]
Consequently,
\[
  \tau_a=0 \quad\Longrightarrow\quad y_{ia}=0
  \quad\text{for every }i,
\]
and equivalently
\[
  \alpha_a=0 \quad\Longrightarrow\quad y_{ia}=0
  \quad\text{for every }i.
\]

This is an exact-object rank-deficient adapter for the basis-indexed
right-Gram route.  It removes the obstruction that zero singular values would
otherwise create in the full-positive normalization \(u_a=y_a/\tau_a\), but it
does not construct the ordered head/tail source split or prove
Eckart--Young tail optimality.  If an implementation computes \(A^TA\),
right singular vectors, singular values, left candidates, projectors, or SVD
data, those are non-probability computed quantities and require explicit
floating-point/inexact-arithmetic certificates.  Sampling probabilities and
sampling laws remain exact mathematical inputs by the current convention.
\end{corollary}

\begin{formalized}
\leanfit{rectRightGramProjectedColumn_normSq_eq_singularValue_sq}
\leanfit{rectRightGramProjectedColumn_normSq_eq_eigenvalue}
\leanfit{rectRightGramProjectedColumn_eq_zero_of_singularValue_eq_zero}
\leanfit{rectRightGramProjectedColumn_eq_zero_of_eigenvalue_eq_zero}
\end{formalized}

\begin{corollary}[Zero-safe right-Gram reconstruction for equation~(9)]
Let \(A\in\mathbb R^{m\times n}\), let \(G=A^TA\), and keep the
basis-indexed right-Gram data \((\alpha_a,V_{\cdot a},\tau_a)\).  With
\[
  y_{ia}=(AV_{\cdot a})_i,
\]
define the zero-safe left candidates by
\[
  \widetilde u_{ia}=
  \begin{cases}
    0,& \tau_a=0,\\
    y_{ia}/\tau_a,& \tau_a\ne 0.
  \end{cases}
\]
Then every basis index satisfies
\[
  \tau_a\,\widetilde u_{ia}=y_{ia}.
\]
Consequently the exact right-Gram row-completeness identity gives the
basis-indexed reconstruction
\[
  A_{ij}
  =
  \sum_{a=1}^n \widetilde u_{ia}\tau_aV_{ja}.
\]

This is an exact-object rank-deficient reconstruction theorem.  It removes the
full-positive hypothesis from the preceding reconstruction by using the zero
projected-column result when \(\tau_a=0\).  It is still basis-indexed rather
than the ordered source head/tail split needed for the paper's final
Eckart--Young route.  If an implementation computes \(A^TA\), singular
vectors, singular values, left candidates, projectors, or SVD data, those are
non-probability computed quantities and require explicit
floating-point/inexact-arithmetic certificates.  Sampling probabilities and
sampling laws remain exact mathematical inputs by the current convention.
\end{corollary}

\begin{formalized}
\leanfit{rectRightGramLeftSingularZeroSafe}
\leanfit{rectRightGramLeftSingularZeroSafe_eq_zero_of_singularValue_eq_zero}
\leanfit{rectRightGramLeftSingularZeroSafe_eq_inv_mul_of_singularValue_ne_zero}
\leanfit{rectRightGramLeftSingularZeroSafe_factor_column}
\leanfit{rectRightGram_basisSVD_representation}
\end{formalized}

\begin{corollary}[Selected right-Gram head/tail split for equation~(9)]
Let \(A\in\mathbb R^{m\times n}\), let \(G=A^TA\), and keep the
basis-indexed right-Gram data \((\alpha_a,V_{\cdot a},\tau_a)\) and the
zero-safe left candidates \(\widetilde u_{\cdot a}\) from the preceding
corollary.  Let \(s\subseteq\{1,\ldots,n\}\) be any finite selected set of
basis indices.  Define the exact selected head and complementary tail by
\[
  H^{(s)}_{ij}
  =
  \sum_{a\in s}\widetilde u_{ia}\tau_aV_{ja},
  \qquad
  T^{(s)}_{ij}
  =
  \sum_{a\notin s}\widetilde u_{ia}\tau_aV_{ja}.
\]
Then the zero-safe reconstruction partitions entrywise as
\[
  A_{ij}=H^{(s)}_{ij}+T^{(s)}_{ij}
  \qquad\text{for every } i,j.
\]
Moreover, if \(\phi:\{1,\ldots,|s|\}\to s\) is the increasing enumeration of
the selected finite set, then the selected head has the explicit rank
factorization
\[
  H^{(s)}_{ij}
  =
  \sum_{b=1}^{|s|}
    \widetilde u_{i,\phi(b)}
    \bigl(\tau_{\phi(b)}V_{j,\phi(b)}\bigr),
\]
and hence
\[
  \operatorname{rank}(H^{(s)})\le |s|.
\]

This is an exact-object selected-index source split.  It proves the algebra
needed once a later theorem identifies \(s\) with the ordered top singular
directions, but it does not yet prove that ordered handoff, rectangular SVD
existence, Eckart--Young tail optimality, randomness-derived certificates, or
computed non-probability SVD/singular-vector/projector/Gram routines.  Sampling
probabilities and sampling laws remain exact mathematical inputs by the
current convention.
\end{corollary}

\begin{formalized}
\leanfit{rectRightGramBasisSVDHead}
\leanfit{rectRightGramBasisSVDTail}
\leanfit{rectRightGramBasisSVD_head_add_tail}
\leanfit{rectRightGramBasisSVD_head_tail_entry}
\leanfit{rectRightGramBasisSVDHeadRankFactorization}
\leanfit{rectRightGramBasisSVDHead_rankAtMost}
\end{formalized}

\begin{corollary}[Selected right-Gram sketch bridge for equation~(9)]
Under the hypotheses of the preceding corollary, let
\(\phi:\{1,\ldots,|s|\}\to s\) be the selected-index enumeration.  Define the
exact selected eigenvector sketch and coefficient table by
\[
  (Z_s)_{jb}=V_{j,\phi(b)},
  \qquad
  (W_s)_{bj}=V_{j,\phi(b)} .
\]
Then the exact column sketch \(A Z_s\) is the selected projected-column table:
\[
  (A Z_s)_{ib}
  =
  (AV_{\cdot,\phi(b)})_i .
\]
Moreover the selected head from Corollary~549 is exactly the corresponding
column-sketch head,
\[
  H^{(s)}_{ij}
  =
  \bigl((A Z_s)W_s\bigr)_{ij}
  =
  \sum_{b=1}^{|s|}
    (A Z_s)_{ib}(W_s)_{bj}.
\]
Equivalently, the selected head has a formal
\(\texttt{ColumnSketchHeadFactorization}\) through the displayed sketch
columns \(A Z_s\).

This is an exact-object selected-head sketch-space theorem.  It closes the
algebraic statement that the arbitrary selected right-Gram head lies in the
column space of the matching selected eigenvector sketch.  It does not identify
the selected set with the ordered top singular directions, prove rectangular
SVD existence, prove Eckart--Young tail optimality, derive randomness
certificates, or certify computed non-probability SVD/singular-vector/projector/
Gram/sketch routines.  Sampling probabilities and sampling laws remain exact
mathematical inputs by the current convention.
\end{corollary}

\begin{formalized}
\leanfit{rectRightGramBasisSketchMatrix}
\leanfit{rectRightGramBasisSketchCoeff}
\leanfit{columnSketch_rectRightGramBasisSketchMatrix}
\leanfit{rectRightGramBasisSketch_head_eq}
\leanfit{rectRightGramBasisSketchHeadFactorization}
\leanfit{rectRightGramBasisSVDHead_columnSketchHeadFactorization}
\end{formalized}

\begin{corollary}[Selected right-Gram equation-(9) certificate adapter]
Let \(s\) be a finite selected set of basis-indexed right-Gram directions, and
let
\[
  Z_s,\qquad H^{(s)},\qquad T^{(s)}
\]
denote respectively the selected eigenvector sketch, selected head, and
selected tail from Corollaries~549 and~550.  Let
\(P_{A Z_s}\in\mathbb R^{m\times m}\) be an exact left multiplier.  Suppose
that the displayed radii \(\tau,\kappa\) are nonnegative and satisfy
\[
  \|T^{(s)}\|_F \le \tau,
  \qquad
  \|P_{A Z_s}T^{(s)}\|_F \le \kappa .
\]
Then the selected data instantiate the exact equation-(9) head/tail sketch
certificate:
\[
  A = H^{(s)} + T^{(s)},\qquad
  H^{(s)}\in \operatorname{range}(A Z_s),
\]
with tail radius \(\tau\) and projected-tail coupling radius \(\kappa\).

If, in addition, \(P_{A Z_s}\) factors through the selected sketch columns,
meaning that for some coefficient table \(C\),
\[
  (P_{A Z_s})_{ij}
  =
  \sum_{b=1}^{|s|}
    (A Z_s)_{ib} C_{bj},
\]
and \(P_{A Z_s}\) reproduces those columns,
\[
  P_{A Z_s}(A Z_s)=A Z_s,
\]
then the exact projector candidate satisfies the equation-(9) rank and
residual surface
\[
  \operatorname{rank}(P_{A Z_s}A)\le |s|,
  \qquad
  \|A-P_{A Z_s}A\|_F\le \tau+\kappa .
\]

This is an exact-object adapter theorem.  It closes the formal handoff from
the arbitrary selected right-Gram split to the existing equation-(9)
head/tail residual API.  It does not prove the tail or projected-tail coupling
bounds, does not prove that a concrete projector reproduces \(A Z_s\), does not
identify \(s\) with ordered top singular directions, does not prove
rectangular SVD existence or Eckart--Young tail optimality, and does not
certify computed non-probability SVD/singular-vector/projector/Gram/sketch
routines.  Sampling probabilities and sampling laws remain exact mathematical
inputs by the current convention.
\end{corollary}

\begin{formalized}
\leanfit{equation9HeadTailSketchCertificate_of_rectRightGramBasisSVDHead}
\leanfit{equation9HeadTailSketchRankResidualSurface_of_rectRightGramBasisSVDHead}
\end{formalized}

\begin{corollary}[Selected cardinality rank handoff for equation~(9)]
Let \(B\in\mathbb R^{m\times n}\) have a displayed exact rank-factorization
through \(r\) columns.  If \(r=k\), then the same factorization is a formal
rank-at-most-\(k\) certificate:
\[
  \operatorname{rank}(B)\le r,\qquad r=k
  \quad\Longrightarrow\quad
  \operatorname{rank}(B)\le k .
\]

Consequently, for the selected right-Gram head from Corollary~549, if the
selected set satisfies
\[
  |s|=k,
\]
then
\[
  \operatorname{rank}(H^{(s)})\le k .
\]

Under the hypotheses of Corollary~551, and under the same cardinality
certificate \(|s|=k\), the selected equation-(9) projector candidate satisfies
the paper-facing rank/residual surface
\[
  \operatorname{rank}(P_{A Z_s}A)\le k,
  \qquad
  \|A-P_{A Z_s}A\|_F\le \tau+\kappa .
\]

This is an exact rank-index handoff.  It closes the bookkeeping step that
turns the arbitrary selected-set rank \(|s|\) into a displayed paper rank
\(k\).  It does not identify \(s\) with ordered top singular directions, does
not prove that a concrete source-faithful top-\(k\) selection has cardinality
\(k\), does not prove rectangular SVD existence, tail/coupling bounds,
projector reproduction, Eckart--Young, randomness certificates, or computed
non-probability SVD/singular-vector/projector/Gram/sketch routines.  Sampling
probabilities and sampling laws remain exact mathematical inputs by the current
convention.
\end{corollary}

\begin{formalized}
\leanfit{rectRankAtMost_of_eq_rank}
\leanfit{rectRightGramBasisSVDHead_rankAtMost_of_card_eq}
\leanfit{equation9HeadTailSketchRankResidualSurface_of_rectRightGramBasisSVDHead_card_eq}
\end{formalized}

\begin{corollary}[Injective selected-index handoff for equation~(9)]
Let \(e:\{1,\ldots,k\}\hookrightarrow \{1,\ldots,n\}\) be an injective map of
displayed selected indices into the right-Gram basis indices.  Define the
selected set
\[
  s_e=\operatorname{image}(e).
\]
Then
\[
  |s_e|=k.
\]
Consequently, the selected right-Gram head induced by \(s_e\) satisfies
\[
  \operatorname{rank}(H^{(s_e)})\le k.
\]

Under the hypotheses of Corollary~551 with \(s=s_e\), the exact projector
candidate satisfies the embedding-selected equation-(9) rank/residual surface
\[
  \operatorname{rank}(P_{A Z_{s_e}}A)\le k,
  \qquad
  \|A-P_{A Z_{s_e}}A\|_F\le \tau+\kappa .
\]

This is an exact injective-selection cardinality theorem.  It removes the
separate \(|s|=k\) hypothesis when the selected directions are supplied by an
injective displayed index map.  It does not prove that \(e\) enumerates the
ordered top singular directions, does not prove rectangular SVD existence,
tail/coupling bounds, projector reproduction, Eckart--Young, randomness
certificates, or computed non-probability SVD/singular-vector/projector/Gram/
sketch routines.  Sampling probabilities and sampling laws remain exact
mathematical inputs by the current convention.
\end{corollary}

\begin{formalized}
\leanfit{rectRightGramSelectedIndexSet}
\leanfit{rectRightGramSelectedIndexSet_card}
\leanfit{rectRightGramBasisSVDHead_rankAtMost_of_embedding}
\leanfit{equation9HeadTailSketchRankResidualSurface_of_rectRightGramBasisSVDHead_embedding}
\end{formalized}

\begin{corollary}[Semantic ordered-top embedding handoff for equation~(9)]
Assume \(k\le n\), and let
\[
  \iota_k^n:\{1,\ldots,k\}\to\{1,\ldots,n\}
\]
be the displayed top-index inclusion.  Let
\[
  e:\{1,\ldots,k\}\hookrightarrow \{1,\ldots,n\}
\]
be an injective selected-index map into the right-Gram basis indices.  Suppose
the semantic ordered-top certificate holds:
\[
  \widehat\sigma_{\mathrm{basis}}(e(a))
    = \sigma_{\mathrm{ord}}(\iota_k^n(a)),
  \qquad a=1,\ldots,k,
\]
where \(\widehat\sigma_{\mathrm{basis}}\) is the basis-indexed right-Gram
singular value attached to mathlib's exact right-Gram eigenbasis and
\(\sigma_{\mathrm{ord}}\) is the ordered zero-indexed right-Gram singular-value
sequence.  Then
\[
  \widehat\sigma_{\mathrm{basis}}(e(a))^2
    = \lambda_{\mathrm{ord}}(\iota_k^n(a)),
  \qquad a=1,\ldots,k,
\]
and the selected basis singular values are ordered:
\[
  a\le b
  \quad\Longrightarrow\quad
  \widehat\sigma_{\mathrm{basis}}(e(b))
    \le \widehat\sigma_{\mathrm{basis}}(e(a)).
\]

Moreover, under the hypotheses of Corollary~553 for the same embedding \(e\),
the exact projector candidate satisfies
\[
  \operatorname{rank}(P_{A Z_{s_e}}A)\le k,
  \qquad
  \|A-P_{A Z_{s_e}}A\|_F\le \tau+\kappa .
\]

This is an exact semantic-order handoff.  It exposes the precise certificate
needed to say that the embedding-selected basis directions are the ordered top
right-Gram singular directions, and then composes that certificate with the
embedding rank/residual surface.  It does not construct the certificate from
mathlib's arbitrary right-Gram eigenbasis order, and it does not prove
rectangular SVD existence, tail/coupling bounds, projector reproduction,
Eckart--Young, randomness certificates, or computed non-probability
SVD/singular-vector/projector/Gram/sketch routines.  Sampling probabilities and
sampling laws remain exact mathematical inputs by the current convention.
\end{corollary}

\begin{formalized}
\leanfit{rectTopIndex}
\leanfit{rectTopIndex_le}
\leanfit{RectRightGramOrderedTopEmbeddingCertificate}
\leanfit{rectRightGramOrderedTopEmbeddingCertificate_selected_sq_eq}
\leanfit{rectRightGramOrderedTopEmbeddingCertificate_selected_antitone}
\leanfit{equation9HeadTailSketchRankResidualSurface_of_rectRightGramBasisSVDHead_orderedTopEmbedding}
\end{formalized}

\begin{corollary}[Constructed ordered-top embedding for equation~(9)]
For the exact right-Gram matrix \(A^T A\), let
\[
  \phi_n:\operatorname{Fin}(\#\operatorname{Fin}(n))\simeq \operatorname{Fin}(n)
\]
be the same finite equivalence used by mathlib to reindex the ordered Hermitian
eigenvalue sequence \(\lambda_{\mathrm{ord}}\) into the basis-indexed
right-Gram eigenvector table.  Given \(k\le n\), define
\[
  e_{\mathrm{ord}}(a)=\phi_n(\iota_k^n(a)),\qquad a=1,\ldots,k.
\]
Then \(e_{\mathrm{ord}}:\{1,\ldots,k\}\hookrightarrow\{1,\ldots,n\}\) is
injective and satisfies the ordered-top certificate
\[
  \widehat\sigma_{\mathrm{basis}}(e_{\mathrm{ord}}(a))
    = \sigma_{\mathrm{ord}}(\iota_k^n(a)),
  \qquad a=1,\ldots,k.
\]
Consequently,
\[
  \widehat\sigma_{\mathrm{basis}}(e_{\mathrm{ord}}(a))^2
    = \lambda_{\mathrm{ord}}(\iota_k^n(a)),
  \qquad
  a\le b\Longrightarrow
  \widehat\sigma_{\mathrm{basis}}(e_{\mathrm{ord}}(b))
    \le \widehat\sigma_{\mathrm{basis}}(e_{\mathrm{ord}}(a)).
\]
Under the hypotheses of Corollary~554 with \(e=e_{\mathrm{ord}}\), the exact
projector candidate satisfies
\[
  \operatorname{rank}(P_{A Z_{s_{e_{\mathrm{ord}}}}}A)\le k,
  \qquad
  \|A-P_{A Z_{s_{e_{\mathrm{ord}}}}}A\|_F\le \tau+\kappa .
\]

This closes the exact ordered-top embedding construction for the right-Gram
eigenbasis chosen by mathlib.  It does not prove rectangular SVD existence,
tail/coupling bounds, projector reproduction, Eckart--Young, randomness
certificates, or computed non-probability SVD/singular-vector/projector/Gram/
sketch routines.  Sampling probabilities and sampling laws remain exact
mathematical inputs by the current convention.
\end{corollary}

\begin{formalized}
\leanfit{rectRightGramOrderedEigenbasisEquiv}
\leanfit{rectRightGramOrderedTopEmbedding}
\leanfit{rectRightGramOrderedTopEmbedding_certificate}
\leanfit{equation9HeadTailSketchRankResidualSurface_of_rectRightGramBasisSVDHead_constructedOrderedTopEmbedding}
\end{formalized}

\begin{corollary}[Constructed ordered top/complement singular-value dominance]
For the exact right-Gram matrix \(A^T A\), let \(\phi_n\) and
\(e_{\mathrm{ord}}\) be as in the preceding corollary.  For a basis-indexed
right-Gram direction \(b\in\{1,\ldots,n\}\), define its ordered coordinate
\[
  j_{\mathrm{ord}}(b)=\phi_n^{-1}(b),
\]
viewed through the canonical cast back to the displayed ordered index set
\(\{1,\ldots,n\}\).  The basis-indexed singular value is exactly the ordered
singular value at this coordinate:
\[
  \widehat\sigma_{\mathrm{basis}}(b)
    = \sigma_{\mathrm{ord}}(j_{\mathrm{ord}}(b)).
\]
If \(b\notin s_{e_{\mathrm{ord}}}\), then
\[
  k\le j_{\mathrm{ord}}(b).
\]
Consequently, for every displayed selected index \(a=1,\ldots,k\),
\[
  \widehat\sigma_{\mathrm{basis}}(b)
    \le
  \widehat\sigma_{\mathrm{basis}}(e_{\mathrm{ord}}(a)).
\]

This is an exact spectral-index dominance theorem for the constructed
top-\(k\) right-Gram selection.  It supplies the ordered head/complement
comparison needed before a future Eckart--Young or tail-optimality proof, but
it is not itself an Eckart--Young theorem, a rectangular SVD existence theorem,
a tail/coupling bound, a projector-reproduction theorem, a randomness-derived
certificate, or a computed non-probability SVD/singular-vector/projector/Gram/
sketch routine certificate.  Sampling probabilities and sampling laws remain
exact mathematical inputs by the current convention.
\end{corollary}

\begin{formalized}
\leanfit{rectRightGramBasisOrderedIndex}
\leanfit{finCardIndex_rectRightGramBasisOrderedIndex}
\leanfit{rectRightGramBasisSingularValue_eq_orderedIndex}
\leanfit{rectRightGramOrderedTopEmbedding_not_mem_index_ge}
\leanfit{rectTopIndex_le_rectRightGramBasisOrderedIndex_of_not_mem_orderedTopEmbedding}
\leanfit{rectRightGramOrderedTopEmbedding_complement_singularValue_le_selected}
\end{formalized}

\begin{corollary}[Constructed top-\(k\) head positivity from the kth singular value]
Assume \(0<k\le n\).  Let
\[
  a_\star=k-1
\]
be the last displayed index in the top-\(k\) block, and suppose the kth ordered
right-Gram singular value is positive:
\[
  0<\sigma_{\mathrm{ord}}(\iota_k^n(a_\star)).
\]
Then every displayed top-\(k\) ordered singular value is positive:
\[
  0<\sigma_{\mathrm{ord}}(\iota_k^n(a)),
  \qquad a=1,\ldots,k.
\]
Consequently, for the constructed ordered-top embedding \(e_{\mathrm{ord}}\),
every selected basis-indexed head singular value is positive and hence nonzero:
\[
  0<\widehat\sigma_{\mathrm{basis}}(e_{\mathrm{ord}}(a)),
  \qquad
  \widehat\sigma_{\mathrm{basis}}(e_{\mathrm{ord}}(a))\ne 0,
  \qquad a=1,\ldots,k.
\]

This is an exact selected-head positivity handoff.  It turns the visible
source-rank-style hypothesis that the kth ordered singular value is positive
into the nonzero selected-head entries needed by source-SVD determinant and
diagonal routes.  It does not prove rectangular SVD existence, source-tail
factor construction, tail/coupling bounds, projector reproduction,
Eckart--Young, randomness certificates, or computed non-probability
SVD/singular-vector/projector/Gram/sketch routines.  Sampling probabilities and
sampling laws remain exact mathematical inputs by the current convention.
\end{corollary}

\begin{formalized}
\leanfit{rectTopLastIndex}
\leanfit{le_rectTopLastIndex}
\leanfit{rectTopIndex_le_last}
\leanfit{rectSingularValue_top_pos_of_last_pos}
\leanfit{rectRightGramOrderedTopEmbedding_selected_pos_of_last_pos}
\leanfit{rectRightGramOrderedTopEmbedding_selected_nonzero_of_last_pos}
\end{formalized}

\begin{corollary}[Constructed top-\(k\) zero-safe left-basis orthonormality]
For a basis-indexed right-Gram direction \(a\), let
\[
  u^{0}_{a} =
  \begin{cases}
    0, & \widehat\sigma_{\mathrm{basis}}(a)=0,\\
    A v_a/\widehat\sigma_{\mathrm{basis}}(a),
      & \widehat\sigma_{\mathrm{basis}}(a)\ne 0,
  \end{cases}
\]
be the zero-safe exact left candidate.  If two displayed basis-indexed singular
values are positive, then their zero-safe left candidates have the exact
orthonormal dot product
\[
  \sum_{i=1}^m u^{0}_{a}(i)u^{0}_{b}(i)=\delta_{ab}.
\]

Consequently, assume \(0<k\le n\) and suppose, as in the preceding corollary,
that the kth ordered right-Gram singular value is positive:
\[
  0<\sigma_{\mathrm{ord}}(\iota_k^n(a_\star)),
  \qquad a_\star=k-1.
\]
For the constructed ordered-top embedding \(e_{\mathrm{ord}}\), define the
selected zero-safe left candidate block
\[
  U_{\mathrm{ord}}(i,a)=u^{0}_{e_{\mathrm{ord}}(a)}(i),
  \qquad a=1,\ldots,k.
\]
Then
\[
  U_{\mathrm{ord}}^T U_{\mathrm{ord}}=I_k,
\]
entrywise:
\[
  \sum_{i=1}^m
    u^{0}_{e_{\mathrm{ord}}(a)}(i)\,
    u^{0}_{e_{\mathrm{ord}}(b)}(i)
  = \delta_{ab},
  \qquad a,b=1,\ldots,k.
\]

This is an exact source-SVD left-basis ingredient for the ordered top-\(k\)
route.  It does not construct the complementary tail factor, does not prove
right-basis completeness for the source split, does not prove rectangular SVD
existence or Eckart--Young tail optimality, and does not certify computed
non-probability SVD/singular-vector/projector/Gram/sketch routines.  Sampling
probabilities and sampling laws remain exact mathematical inputs by the current
convention.
\end{corollary}

\begin{formalized}
\leanfit{rectRightGramLeftSingularZeroSafe_col_orthonormal_of_pos}
\leanfit{rectRightGramOrderedTopEmbedding_leftZeroSafe_col_orthonormal_of_last_pos}
\end{formalized}

\begin{corollary}[Constructed ordered source-head factorization]
Assume \(k\le n\), and let \(e_{\mathrm{ord}}\) be the constructed ordered
top-\(k\) right-Gram embedding.  Define the exact ordered head factors
\[
  U_{\mathrm{ord}}(i,a)=u^0_{e_{\mathrm{ord}}(a)}(i),\qquad
  V_{\mathrm{ord}}(j,a)=v_{e_{\mathrm{ord}}(a)}(j),
\]
and the diagonal head singular-value block
\[
  \Sigma_{\mathrm{ord}}(a,b)=
  \begin{cases}
    \widehat\sigma_{\mathrm{basis}}(e_{\mathrm{ord}}(a)), & a=b,\\
    0, & a\ne b.
  \end{cases}
\]
Let \(s_{\mathrm{ord}}\) be the selected finite set induced by
\(e_{\mathrm{ord}}\), and let \(H_{s_{\mathrm{ord}}}\) be the selected
right-Gram head.  Then, entrywise,
\[
  H_{s_{\mathrm{ord}}}
  =
  U_{\mathrm{ord}}\Sigma_{\mathrm{ord}}V_{\mathrm{ord}}^T.
\]
Equivalently,
\[
  H_{s_{\mathrm{ord}}}(i,j)
  =
  \sum_{a=1}^k
    u^0_{e_{\mathrm{ord}}(a)}(i)\,
    \widehat\sigma_{\mathrm{basis}}(e_{\mathrm{ord}}(a))\,
    v_{e_{\mathrm{ord}}(a)}(j).
\]

If \(0<k\le n\) and the kth ordered singular value is positive, then the head
left factor has exact orthonormal columns:
\[
  U_{\mathrm{ord}}^T U_{\mathrm{ord}}=I_k.
\]
The head right factor has exact orthonormal columns without the positivity
hypothesis:
\[
  V_{\mathrm{ord}}^T V_{\mathrm{ord}}=I_k.
\]

This is an exact ordered source-head factorization theorem.  It closes the
head side of the constructed top-\(k\) source split, but it does not construct
the complementary tail factor, prove row completeness for the full right-basis
split, prove Eckart--Young tail optimality, derive randomness certificates, or
certify computed non-probability SVD/singular-vector/projector/Gram/sketch
routines.  Sampling probabilities and sampling laws remain exact mathematical
inputs by the current convention.
\end{corollary}

\begin{formalized}
\leanfit{rectRightGramOrderedHeadLeft}
\leanfit{rectRightGramOrderedHeadRight}
\leanfit{rectRightGramOrderedHeadSingularDiagonal}
\leanfit{sourceSVDFactorMatrix_rectRightGramOrderedHead_entry}
\leanfit{rectRightGramBasisSVDHead_orderedTopEmbedding_eq_sourceSVDFactorMatrix}
\leanfit{rectRightGramOrderedHeadLeft_col_orthonormal_of_last_pos}
\leanfit{rectRightGramOrderedHeadRight_col_orthonormal}
\end{formalized}

\begin{corollary}[Constructed ordered source-tail factorization]
Assume \(k\le n\), let \(e_{\mathrm{ord}}\) be the constructed ordered
top-\(k\) right-Gram embedding, and let \(s_{\mathrm{ord}}\) be its selected
finite set.  Put \(q=|s_{\mathrm{ord}}^c|\), and enumerate the complement by
\[
  c \longmapsto e_{\mathrm{tail}}(c),\qquad c=1,\ldots,q.
\]
Define the exact ordered complement-tail factors
\[
  U_{\mathrm{tail}}(i,c)=u^0_{e_{\mathrm{tail}}(c)}(i),\qquad
  V_{\mathrm{tail}}(j,c)=v_{e_{\mathrm{tail}}(c)}(j),
\]
and
\[
  \Sigma_{\mathrm{tail}}(c,d)=
  \begin{cases}
    \widehat\sigma_{\mathrm{basis}}(e_{\mathrm{tail}}(c)), & c=d,\\
    0, & c\ne d.
  \end{cases}
\]
Let \(T_{s_{\mathrm{ord}}}\) be the complementary right-Gram tail obtained by
summing the zero-safe reconstruction terms over \(s_{\mathrm{ord}}^c\).  Then,
entrywise,
\[
  T_{s_{\mathrm{ord}}}
  =
  U_{\mathrm{tail}}\Sigma_{\mathrm{tail}}V_{\mathrm{tail}}^T,
\]
or equivalently
\[
  T_{s_{\mathrm{ord}}}(i,j)
  =
  \sum_{c=1}^q
    u^0_{e_{\mathrm{tail}}(c)}(i)\,
    \widehat\sigma_{\mathrm{basis}}(e_{\mathrm{tail}}(c))\,
    v_{e_{\mathrm{tail}}(c)}(j).
\]
The exact right-tail table has orthonormal columns:
\[
  V_{\mathrm{tail}}^T V_{\mathrm{tail}}=I_q.
\]
Together with the ordered source-head factorization, the exact ordered source
split reconstructs \(A\):
\[
  U_{\mathrm{ord}}\Sigma_{\mathrm{ord}}V_{\mathrm{ord}}^T
  +
  U_{\mathrm{tail}}\Sigma_{\mathrm{tail}}V_{\mathrm{tail}}^T
  =
  A.
\]

This is an exact source-split algebra theorem.  It does not assert that
\(U_{\mathrm{tail}}\) is an orthonormal rectangular-SVD tail basis when some
tail singular values vanish; a nullspace completion is still required for that
stronger certificate.  It also does not prove full right-basis row
completeness, Eckart--Young tail optimality, randomness-derived certificates,
or computed non-probability SVD/singular-vector/projector/Gram/sketch routines.
Sampling probabilities and sampling laws remain exact mathematical inputs by
the current convention.
\end{corollary}

\begin{formalized}
\leanfit{rectRightGramBasisSVDTailLeft}
\leanfit{rectRightGramBasisSVDTailRight}
\leanfit{rectRightGramBasisSVDTailSingularDiagonal}
\leanfit{rectRightGramOrderedTailLeft}
\leanfit{rectRightGramOrderedTailRight}
\leanfit{rectRightGramOrderedTailSingularDiagonal}
\leanfit{rectRightGramBasisSVDTailRight_col_orthonormal}
\leanfit{rectRightGramOrderedTailRight_col_orthonormal}
\leanfit{sourceSVDFactorMatrix_rectRightGramBasisSVDTail_entry}
\leanfit{rectRightGramBasisSVDTail_eq_sourceSVDFactorMatrix}
\leanfit{rectRightGramBasisSVDTail_orderedTopEmbedding_eq_sourceSVDFactorMatrix}
\leanfit{rectRightGramOrdered_source_head_add_tail}
\end{formalized}

\begin{corollary}[Constructed ordered right-basis block completeness]
Assume \(k\le n\), let \(s_{\mathrm{ord}}\) be the selected finite set induced
by the constructed ordered top-\(k\) right-Gram embedding, and put
\(q=|s_{\mathrm{ord}}^c|\).  Let \(V_{\mathrm{ord}}\) and
\(V_{\mathrm{tail}}\) be the exact right factors from the ordered source-head
and ordered source-tail constructions.  Define the block right table
\[
  Q_{\mathrm{ord}}(j,\beta)=
  \begin{cases}
    V_{\mathrm{tail}}(j,c), & \beta=\mathrm{inl}(c),\\
    V_{\mathrm{ord}}(j,a), & \beta=\mathrm{inr}(a).
  \end{cases}
\]
Then \(Q_{\mathrm{ord}}\) has exact orthonormal columns:
\[
  \sum_{j=1}^n Q_{\mathrm{ord}}(j,\beta)\,
        Q_{\mathrm{ord}}(j,\beta')
  =
  \delta_{\beta\beta'},
  \qquad \beta,\beta'\in\{1,\ldots,q\}\sqcup\{1,\ldots,k\}.
\]
It is also row-complete:
\[
  \sum_{\beta}
    Q_{\mathrm{ord}}(j,\beta)\,Q_{\mathrm{ord}}(\ell,\beta)
  =
  \delta_{j\ell},
  \qquad j,\ell=1,\ldots,n.
\]
Equivalently,
\[
  Q_{\mathrm{ord}}^TQ_{\mathrm{ord}}=I_{q+k},
  \qquad
  Q_{\mathrm{ord}}Q_{\mathrm{ord}}^T=I_n.
\]

The proof uses only exact selected-set and complement enumeration, disjointness
of a finite set and its complement, and exact column/row orthonormality of the
right-Gram eigenbasis.  This closes the constructed ordered right-basis block
fields needed by the source-tail Frobenius route.  It does not construct a
nullspace-completed orthonormal tail-left basis, prove Eckart--Young tail
optimality, derive randomness certificates, or certify computed non-probability
SVD/singular-vector/projector/Gram/sketch routines.  Sampling probabilities and
sampling laws remain exact mathematical inputs by the current convention.
\end{corollary}

\begin{formalized}
\leanfit{rectRightGramSelectedIndexSet_sum}
\leanfit{rectRightGramComplement_sum_orderEmbOfFin}
\leanfit{rectRightGramSelectedIndexSet_head_tail_cross_zero}
\leanfit{rectRightGramSelectedIndexSet_tail_head_cross_zero}
\leanfit{rectRightGramSelectedIndexSet_tail_head_row_complete}
\leanfit{rectRightGramOrderedRightBasisBlock_col_orthonormal}
\leanfit{rectRightGramOrderedRightBasisBlock_row_orthonormal}
\leanfit{rectRightGramOrderedRightBasisBlock_col_row_orthonormal}
\end{formalized}

\begin{corollary}[Constructed ordered block source-SVD certificate]
Assume \(k\le n\).  Let
\[
  q=\left|s_{\rm ord}^c\right|,
\]
where \(s_{\rm ord}\) is the selected finite set induced by the constructed
ordered top-\(k\) right-Gram embedding.  Let
\[
  H_{\rm ord}=U_{\rm ord}\Sigma_{\rm ord}V_{\rm ord}^T,
  \qquad
  T_{\rm ord}=U_{\rm tail}\Sigma_{\rm tail}V_{\rm tail}^T
\]
be the constructed ordered source-head and source-tail factors from the exact
right-Gram eigenbasis.  Suppose the remaining left block is certified:
\[
  L_{\rm ord}(i,\alpha)=
  \begin{cases}
    U_{\rm ord}(i,a), & \alpha=\operatorname{inl}(a),\\
    U_{\rm tail}(i,c), & \alpha=\operatorname{inr}(c),
  \end{cases}
  \qquad
  L_{\rm ord}^TL_{\rm ord}=I_{k+q},
\]
and suppose the displayed ordered head singular values are nonzero,
\[
  \sigma_a
  =
  \sigma_{\rm right}\!\left(\operatorname{ordTop}(a)\right)\ne0,
  \qquad a=1,\ldots,k.
\]
Then the constructed ordered data instantiate the block diagonal source-SVD
certificate:
\[
  A=
  U_{\rm ord}\Sigma_{\rm ord}V_{\rm ord}^T+
  U_{\rm tail}\Sigma_{\rm tail}V_{\rm tail}^T,
\]
\[
  L_{\rm ord}^TL_{\rm ord}=I_{k+q},
  \qquad
  Q_{\rm ord}^TQ_{\rm ord}=I_{q+k},
  \qquad
  Q_{\rm ord}Q_{\rm ord}^T=I_n,
\]
where \(Q_{\rm ord}=[V_{\rm tail},V_{\rm ord}]\), and
\[
  (\Sigma_{\rm ord})_{ab}=\delta_{ab}\sigma_a,
  \qquad
  \sigma_a\ne0.
\]
Equivalently, the exact constructed ordered split supplies
\[
  \mathrm{BlockDiagonalSourceSVDTailCertificate}
  \bigl(A,U_{\rm ord},\Sigma_{\rm ord},\sigma,V_{\rm ord},
        U_{\rm tail},\Sigma_{\rm tail},V_{\rm tail}\bigr).
\]

There is also a component-left form.  It is enough to prove
\[
  U_{\rm ord}^TU_{\rm ord}=I_k,\qquad
  U_{\rm ord}^TU_{\rm tail}=0,\qquad
  U_{\rm tail}^TU_{\rm tail}=I_q,
\]
together with the same nonzero head singular values.  If \(k>0\) and the kth
ordered right-Gram singular value is positive, then the nonzero head field and
\(U_{\rm ord}^TU_{\rm ord}=I_k\) are supplied by the constructed ordered
head-positivity theorem.  In that case the only remaining left-side hypotheses
are
\[
  U_{\rm ord}^TU_{\rm tail}=0,
  \qquad
  U_{\rm tail}^TU_{\rm tail}=I_q.
\]

This is an exact-object certificate adapter.  It closes the route from the
constructed ordered right-Gram source split and right-basis block to the
existing block source-SVD certificate surface.  It does not construct the
nullspace-completed orthonormal tail-left basis, prove Eckart--Young tail
optimality, derive randomness certificates, or certify computed
non-probability SVD/singular-vector/projector/Gram/sketch routines.  Sampling
probabilities and sampling laws remain exact mathematical inputs by the current
convention.
\end{corollary}

\begin{formalized}
\leanfit{leftBasisBlock_col_orthonormal_of_component_orthonormal_fields}
\leanfit{BlockDiagonalSourceSVDTailCertificate.of_rectRightGramOrderedSourceSplit}
\leanfit{BlockDiagonalSourceSVDTailCertificate.of_rectRightGramOrderedSourceSplit_component_left}
\leanfit{BlockDiagonalSourceSVDTailCertificate.of_rectRightGramOrderedSourceSplit_component_left_of_last_pos}
\end{formalized}

\begin{corollary}[Constructed ordered head-tail left cross field]
Assume \(k>0\), \(k\le n\), and the kth ordered right-Gram singular value is
positive:
\[
  0<\sigma_k(A).
\]
Let \(U_{\rm ord}\) be the constructed ordered head-left zero-safe table and
let \(U_{\rm tail}\) be the constructed complement-tail zero-safe table.  Then
the exact head/tail left cross field vanishes:
\[
  U_{\rm ord}^T U_{\rm tail}=0,
  \qquad
  \sum_{i=1}^m U_{\rm ord}(i,a)U_{\rm tail}(i,c)=0
\]
for every displayed head index \(a=1,\ldots,k\) and every complement-tail
index \(c=1,\ldots,q\).

More generally, if \(u_a\) and \(u_b\) are zero-safe left candidates built from
right-Gram basis directions \(a\ne b\), and the singular value for \(a\) is
strictly positive, then
\[
  \sum_{i=1}^m u_a(i)u_b(i)=0.
\]
If the singular value for \(b\) is zero, this follows because \(u_b=0\); if it
is nonzero, the existing positive-pair orthonormality theorem gives the
off-diagonal zero.

Consequently, under the same kth-positivity hypothesis, the constructed
ordered block source-SVD certificate from the preceding corollary only needs
the remaining tail-left orthonormality field
\[
  U_{\rm tail}^T U_{\rm tail}=I_q.
\]
The head-left orthonormality, head nonzero field, and head-tail left
cross-orthogonality are all supplied by the constructed ordered right-Gram
theorems.

This is still an exact-object D3 theorem.  It does not construct a
nullspace-completed orthonormal tail-left basis, prove Eckart--Young tail
optimality, derive randomness certificates, or certify computed
non-probability SVD/singular-vector/projector/Gram/sketch routines.  Sampling
probabilities and sampling laws remain exact mathematical inputs by the current
convention.
\end{corollary}

\begin{formalized}
\leanfit{rectRightGramLeftSingularZeroSafe_cross_zero_of_pos_ne}
\leanfit{rectRightGramOrderedHeadTailLeft_cross_zero_of_last_pos}
\leanfit{BlockDiagonalSourceSVDTailCertificate.of_rectRightGramOrderedSourceSplit_tail_left_of_last_pos}
\end{formalized}

\begin{corollary}[Positive-complement ordered tail-left certificate for equation~(9)]
Let \(s\) be any selected right-Gram basis-index set and enumerate its
complement by \(e_c\), \(c=1,\ldots,q\).  Assume every complement singular
value is strictly positive:
\[
  0<\tau_{e_c}(A),\qquad c=1,\ldots,q.
\]
Let \(U_{\rm tail}\) be the complement-tail zero-safe left candidate table,
\[
  U_{\rm tail}(i,c)=u_{e_c}(i).
\]
Then the complement-tail left table has exact orthonormal columns:
\[
  U_{\rm tail}^T U_{\rm tail}=I_q,
  \qquad
  \sum_{i=1}^m U_{\rm tail}(i,c)U_{\rm tail}(i,d)=\delta_{cd}.
\]
For the constructed ordered top-\(k\) selected set, if additionally
\[
  k>0,\qquad k\le n,\qquad 0<\sigma_k(A),
\]
and every constructed complement singular value is strictly positive, then the
constructed ordered block source-SVD certificate is instantiated entirely from
the constructed exact right-Gram objects:
\[
  \operatorname{BlockDiagonalSourceSVDTailCertificate}
  \bigl(A,U_{\rm ord},\Sigma_{\rm ord},\sigma_{\rm ord},V_{\rm ord},
        U_{\rm tail},\Sigma_{\rm tail},V_{\rm tail}\bigr).
\]
Here kth head positivity supplies the head nonzero and head-left fields,
Corollary~563 supplies \(U_{\rm ord}^T U_{\rm tail}=0\), and complement
positivity supplies \(U_{\rm tail}^T U_{\rm tail}=I_q\).

This is an exact-object positive-complement branch.  It does not construct a
nullspace-completed tail-left basis when a complement singular value is zero,
prove Eckart--Young tail optimality, derive randomness certificates, or
certify computed non-probability SVD/singular-vector/projector/Gram/sketch
routines.  Sampling probabilities and sampling laws remain exact mathematical
inputs by the current convention.
\end{corollary}

\begin{formalized}
\leanfit{rectRightGramBasisSVDTailLeft_col_orthonormal_of_pos}
\leanfit{rectRightGramOrderedTailLeft_col_orthonormal_of_complement_pos}
\leanfit{BlockDiagonalSourceSVDTailCertificate.of_rectRightGramOrderedSourceSplit_all_tail_pos_of_last_pos}
\end{formalized}

\begin{corollary}[Zero-tail obstruction for the raw zero-safe tail-left table]
Let \(s\) be any selected right-Gram basis-index set and enumerate its
complement by \(e_c\), \(c=1,\ldots,q\).  If one complement singular value
vanishes,
\[
  \tau_{e_c}(A)=0,
\]
then the corresponding zero-safe tail-left column is the zero column:
\[
  U_{\rm tail}(i,c)=0,\qquad i=1,\ldots,m.
\]
Consequently its self-dot is zero,
\[
  \sum_{i=1}^m U_{\rm tail}(i,c)^2=0,
\]
and the raw zero-safe complement-tail table cannot be column orthonormal,
because column orthonormality would require the same diagonal self-dot to be
\[
  (I_q)_{cc}=1.
\]
For the constructed ordered top-\(k\) selected set, the same obstruction rules
out any block source-SVD certificate whose tail-left block is the raw
constructed zero-safe table:
\[
  \neg\operatorname{BlockDiagonalSourceSVDTailCertificate}
  \bigl(A,U_{\rm ord},\Sigma_{\rm ord},\sigma_{\rm ord},V_{\rm ord},
        U_{\rm tail},\Sigma_{\rm tail},V_{\rm tail}\bigr)
\]
whenever a constructed complement singular value is zero.

This is a route-elimination result, not a replacement construction.  It proves
that the remaining zero-tail case genuinely needs a nullspace-completed
orthonormal tail-left basis.  It does not construct that basis, prove
Eckart--Young tail optimality, derive randomness certificates, or certify
computed non-probability SVD/singular-vector/projector/Gram/sketch routines.
Sampling probabilities and sampling laws remain exact mathematical inputs by
the current convention.
\end{corollary}

\begin{formalized}
\leanfit{rectRightGramBasisSVDTailLeft_self_dot_eq_zero_of_singularValue_eq_zero}
\leanfit{not_rectRightGramBasisSVDTailLeft_col_orthonormal_of_zero_singularValue}
\leanfit{not_rectRightGramOrderedTailLeft_col_orthonormal_of_zero_complement_singularValue}
\leanfit{not_BlockDiagonalSourceSVDTailCertificate_rectRightGramOrdered_zero_safe_tail_of_zero_complement_singularValue}
\end{formalized}

\begin{corollary}[Replacement tail-left adapter for the zero-tail source split]
Let \(s\) be a selected right-Gram basis-index set with complement enumeration
\(e_c\), \(c=1,\ldots,q\).  Let \(\widetilde U_{\rm tail}\in\mathbb R^{m\times q}\)
be any replacement tail-left table satisfying the nonzero-direction agreement
condition
\[
  \tau_{e_c}(A)\ne 0
  \quad\Longrightarrow\quad
  \widetilde U_{\rm tail}(i,c)=U_{\rm tail}(i,c)
  \qquad\text{for all }i,c.
\]
Then the exact complement-tail source factor is unchanged:
\[
  T_s
  =
  \widetilde U_{\rm tail}\Sigma_{\rm tail}V_{\rm tail}^T.
\]
Indeed, if \(\tau_{e_c}(A)=0\), the diagonal entry of
\(\Sigma_{\rm tail}\) erases the \(c\)-th replacement column; if
\(\tau_{e_c}(A)\ne0\), the replacement column agrees with the zero-safe
column.

For the constructed ordered top-\(k\) selected set, assume
\[
  k>0,\qquad k\le n,\qquad 0<\sigma_k(A),
\]
and suppose \(\widetilde U_{\rm tail}\) satisfies the same nonzero-direction
agreement, the replacement orthonormality field
\[
  \widetilde U_{\rm tail}^T\widetilde U_{\rm tail}=I_q,
\]
and the replacement head-tail cross field
\[
  U_{\rm ord}^T\widetilde U_{\rm tail}=0.
\]
Then the constructed ordered block source-SVD certificate holds with
\(\widetilde U_{\rm tail}\) in place of the raw zero-safe tail-left table:
\[
  \operatorname{BlockDiagonalSourceSVDTailCertificate}
  \bigl(A,U_{\rm ord},\Sigma_{\rm ord},\sigma_{\rm ord},V_{\rm ord},
        \widetilde U_{\rm tail},\Sigma_{\rm tail},V_{\rm tail}\bigr).
\]

This is the exact adapter for a future nullspace-completed tail-left
construction.  It does not construct \(\widetilde U_{\rm tail}\), prove
Eckart--Young tail optimality, derive randomness certificates, or certify
computed non-probability SVD/singular-vector/projector/Gram/sketch routines.
Sampling probabilities and sampling laws remain exact mathematical inputs by
the current convention.
\end{corollary}

\begin{formalized}
\leanfit{sourceSVDFactorMatrix_rectRightGramBasisSVDTail_replacement_left_entry}
\leanfit{rectRightGramBasisSVDTail_eq_sourceSVDFactorMatrix_replacement_left}
\leanfit{rectRightGramOrdered_source_head_add_tail_replacement_left}
\leanfit{BlockDiagonalSourceSVDTailCertificate.of_rectRightGramOrderedSourceSplit_replacement_tail_left_of_last_pos}
\end{formalized}

\begin{corollary}[Left-block dimension guard for equation~(9)]
Let \(\mathcal I\) be a finite index set and let
\(X\in\mathbb R^{m\times \mathcal I}\) satisfy the exact column
orthonormality equations
\[
  \sum_{i=1}^m X(i,\alpha)X(i,\beta)
  =
  \begin{cases}
  1,&\alpha=\beta,\\
  0,&\alpha\ne\beta,
  \end{cases}
  \qquad \alpha,\beta\in\mathcal I.
\]
Then
\[
  |\mathcal I|\le m.
\]
Consequently, if the concatenated left source-SVD block
\[
  [\,U\;U_{\rm tail}\,]\in\mathbb R^{m\times(r+q)}
\]
is column orthonormal, then
\[
  r+q\le m.
\]
In particular, every exact block source-SVD certificate of the form
\[
  \operatorname{BlockDiagonalSourceSVDTailCertificate}
  \bigl(A,U,\Sigma,\sigma,V,U_{\rm tail},\Sigma_{\rm tail},V_{\perp}\bigr)
\]
exposes the necessary dimension side condition \(r+q\le m\).

For the full right-Gram split used in the current ordered equation~(9) route,
whenever the selected and complement blocks partition \(n\) right directions so
that \(r+q=n\), this guard reads \(n\le m\).  Thus the nullspace-completed
replacement-tail construction must either carry this tall/thin condition
explicitly or move to a rectangular SVD surface with only the appropriate
number of left columns.  This is an exact-object dimension guard.  It does not
construct the replacement basis, prove Eckart--Young tail optimality, derive
randomness certificates, or certify computed non-probability
SVD/singular-vector/projector/Gram/sketch routines.  Sampling probabilities and
sampling laws remain exact mathematical inputs by the current convention.
\end{corollary}

\begin{formalized}
\leanfit{colOrthonormal_fintype_card_le_rows}
\leanfit{leftBasisBlock_col_orthonormal_card_le_rows}
\leanfit{BlockDiagonalSourceSVDTailCertificate.left_column_count_le_row_dim}
\end{formalized}

\begin{corollary}[Partial left-block orthonormal completion for equation~(9)]
Let \(S\subseteq\{1,\ldots,m\}\), and let \(X\in\mathbb R^{m\times m}\)
satisfy the exact partial column-orthonormality equations
\[
  \sum_{i=1}^m X(i,a)X(i,b)
  =
  \begin{cases}
  1,&a=b,\\
  0,&a\ne b,
  \end{cases}
  \qquad a,b\in S.
\]
Then there exists \(Y\in\mathbb R^{m\times m}\) such that
\[
  Y(:,a)=X(:,a)\qquad (a\in S)
\]
and
\[
  \sum_{i=1}^m Y(i,a)Y(i,b)
  =
  \begin{cases}
  1,&a=b,\\
  0,&a\ne b,
  \end{cases}
  \qquad a,b\in\{1,\ldots,m\}.
\]

More generally, let
\[
  e:\{1,\ldots,r\}\sqcup\{1,\ldots,q\}\hookrightarrow \{1,\ldots,m\}
\]
embed the displayed head/tail block columns into the ambient left-column
index set.  Let \(S\subseteq \{1,\ldots,r\}\sqcup\{1,\ldots,q\}\) contain
every head index, and suppose the columns of the partially specified block
\[
  [\,U\;U_{\rm tail}^{(0)}\,]
\]
indexed by \(S\) are exactly orthonormal.  Then there exists a replacement
tail-left table \(U_{\rm tail}\in\mathbb R^{m\times q}\) such that
\[
  U_{\rm tail}(:,a)=U_{\rm tail}^{(0)}(:,a)
  \quad\text{for every tail index }a\text{ with }{\rm tail}(a)\in S,
\]
and
\[
  [\,U\;U_{\rm tail}\,]^T[\,U\;U_{\rm tail}\,]=I_{r+q}.
\]

This is exact-object nullspace-completion infrastructure for the
equation~(9) route.  It does not yet instantiate \(S\) with the ordered
nonzero-tail directions, compose the completion with the replacement
source-factor adapter, prove Eckart--Young tail optimality, derive randomness
certificates, or certify computed non-probability
SVD/singular-vector/projector/Gram/sketch routines.  Sampling probabilities and
sampling laws remain exact mathematical inputs by the current convention.
\end{corollary}

\begin{formalized}
\leanfit{partialColOrthonormal_exists_fullColOrthonormal}
\leanfit{partialLeftBasisBlock_exists_replacement_tail}
\end{formalized}

\begin{corollary}[Ordered nonzero-tail completion for equation~(9)]
Let \(k\le n\), \(k>0\), and suppose the \(k\)th ordered singular value of
the exact right Gram route is positive:
\[
  0<\sigma_k.
\]
Let \(S_{\rm tail}\) be the complement of the constructed ordered top-\(k\)
right-index set, enumerated by \(c=1,\ldots,q\), and write
\[
  \tau_c
  =
  \operatorname{rectRightGramBasisSingularValue}
    \bigl(A,S_{\rm tail}(c)\bigr).
\]
Define the partial left-block set
\[
  S_{\ne 0}
  =
  \{\operatorname{head}(a):1\le a\le k\}
  \;\cup\;
  \{\operatorname{tail}(c):1\le c\le q,\ \tau_c\ne 0\}.
\]
For the exact zero-safe ordered left candidates
\[
  U_{\rm ord}=\operatorname{rectRightGramOrderedHeadLeft}(A),
  \qquad
  U_{\rm tail}^{(0)}
    =\operatorname{rectRightGramOrderedTailLeft}(A),
\]
the columns of \([\,U_{\rm ord}\;U_{\rm tail}^{(0)}\,]\) indexed by
\(S_{\ne 0}\) are exactly orthonormal:
\[
  \sum_{i=1}^m
    [\,U_{\rm ord}\;U_{\rm tail}^{(0)}\,](i,\alpha)
    [\,U_{\rm ord}\;U_{\rm tail}^{(0)}\,](i,\beta)
  =
  \begin{cases}
  1,&\alpha=\beta,\\
  0,&\alpha\ne\beta,
  \end{cases}
  \qquad \alpha,\beta\in S_{\ne 0}.
\]

Assume in addition that the full displayed head/tail left-block index set
embeds into the ambient row-coordinate set,
\[
  e:\{1,\ldots,k\}\sqcup\{1,\ldots,q\}\hookrightarrow \{1,\ldots,m\}.
\]
Then there exists a replacement tail-left table
\(U_{\rm tail}\in\mathbb R^{m\times q}\) such that
\[
  U_{\rm tail}(:,c)=U_{\rm tail}^{(0)}(:,c)
  \qquad \text{whenever }\tau_c\ne 0,
\]
\[
  [\,U_{\rm ord}\;U_{\rm tail}\,]^T
  [\,U_{\rm ord}\;U_{\rm tail}\,]=I_{k+q},
\]
and the ordered block source-SVD certificate
\[
\begin{aligned}
&\operatorname{BlockDiagonalSourceSVDTailCertificate}
 \bigl(A,U_{\rm ord},\Sigma_{\rm ord},\sigma_{\rm ord},V_{\rm ord},
       U_{\rm tail},\Sigma_{\rm tail},V_{\rm tail}\bigr)
\end{aligned}
\]
holds.  Thus the replacement-tail source factor agrees with the ordered
zero-safe tail source factor on all nonzero tail singular directions, while
zero tail singular directions are supplied by orthonormal completion.

This is exact-object nullspace-completion infrastructure.  It does not prove
Eckart--Young tail optimality, randomness-derived cross-term certificates, or
computed non-probability SVD/singular-vector/projector/Gram/sketch routines.
Sampling probabilities and sampling laws remain exact mathematical inputs by
the current convention.
\end{corollary}

\begin{formalized}
\leanfit{rectRightGramOrderedTailIndex}
\leanfit{rectRightGramOrderedNonzeroTailPartialSet}
\leanfit{rectRightGramOrderedNonzeroTailPartialSet_head}
\leanfit{rectRightGramOrderedNonzeroTailPartialSet_tail_iff}
\leanfit{rectRightGramOrderedNonzeroTailPartialSet_leftBasisBlock_col_orthonormal_of_last_pos}
\leanfit{exists_rectRightGramOrdered_replacement_tail_left_block_certificate_of_last_pos}
\end{formalized}

\begin{corollary}[Ordered replacement-tail rank surface for equation~(9)]
Keep the hypotheses of the ordered nonzero-tail completion corollary.  Thus
\(k\le n\), \(k>0\), the \(k\)th ordered right-Gram singular value is positive,
and the full ordered head/tail left-block index set embeds in the ambient row
coordinates.  Let
\[
  U_{\rm ord},\quad
  \Sigma_{\rm ord},\quad
  V_{\rm ord},\quad
  \Sigma_{\rm tail},\quad
  V_{\rm tail}
\]
be the exact constructed ordered right-Gram source factors, and let
\(Z\in\mathbb R^{n\times k}\) be the exact sketching matrix.  Assume the exact
source full-rank condition
\[
  \det(V_{\rm ord}^T Z)\ne 0,
\]
and assume the exact cross-term certificate
\[
  \bigl\|
    \Sigma_{\rm tail}(V_{\rm tail}^T Z)(V_{\rm ord}^T Z)^{-1}
  \bigr\|_F
  \le
  \varepsilon\,\|\Sigma_{\rm tail}\|_F,
  \qquad \varepsilon\ge 0.
\]
Then there exists a replacement tail-left table \(U_{\rm tail}\) such that
\[
  U_{\rm tail}(:,c)=U_{\rm tail}^{(0)}(:,c)
  \quad\text{for every nonzero complement-tail singular direction }c,
\]
\[
  [\,U_{\rm ord}\;U_{\rm tail}\,]^T
  [\,U_{\rm ord}\;U_{\rm tail}\,]=I,
\]
and the exact Gram-inverse column-sketch projector
\[
  P_{AZ}=AZ\bigl((AZ)^T(AZ)\bigr)^{-1}(AZ)^T
\]
satisfies the equation-(9) rank/residual surface
\[
  \operatorname{rank}(P_{AZ}A)\le k,
  \qquad
  \rho_F(A,P_{AZ}A)
  \le
  2\sqrt{1+\varepsilon^2}\,\|\Sigma_{\rm tail}\|_F.
\]
The theorem also returns the exact projector reproduction fields
\[
  P_{AZ}(AZ)=AZ,
  \qquad
  P_{AZ}P_{AZ}=P_{AZ},
\]
and the repository symmetric-projector certificate for \(P_{AZ}\).

This is a new exact-object source-split-to-rank-surface result.  It does not
prove the relative Eckart--Young conclusion
\(\rho_F(A,P_{AZ}A)\le\rho\,\rho_F(A,A_k)\), does not derive the displayed
cross-term certificate from randomness, and does not certify computed
non-probability SVD/singular-vector/projector/Gram/inverse/sketch/product
routines.  Sampling probabilities and sampling laws remain exact mathematical
inputs by the current convention.
\end{corollary}

\begin{formalized}
\leanfit{exists_columnSketchGramInverseProjector_sourceHeadTail_sourceSVDTailRankResidualSurface_of_rectRightGramOrdered_replacement_tail_left_of_last_pos}
\end{formalized}

\begin{corollary}[Ordered replacement-tail relative surface for equation~(9)]
Keep the hypotheses and notation of the ordered replacement-tail rank-surface
corollary.  In addition, assume the exact tail-optimality inequality
\[
  \|\Sigma_{\rm tail}\|_F
  \le
  \rho_F(A,B)
  \qquad
  \text{for every exact }B\text{ with }\operatorname{rank}(B)\le k,
\]
and assume the scalar comparison
\[
  2\sqrt{1+\varepsilon^2}\,\|\Sigma_{\rm tail}\|_F
  \le
  \rho\,\|\Sigma_{\rm tail}\|_F.
\]
Then there exists a replacement tail-left table \(U_{\rm tail}\) satisfying
the same nonzero-direction agreement and block-orthonormality conditions as in
the preceding corollary, and the constructed ordered source head
\[
  A_k^{\rm ord}
  =
  U_{\rm ord}\Sigma_{\rm ord}V_{\rm ord}^T
\]
is an exact best rank-\(k\) approximation in the repository Frobenius sense:
\[
  A_k^{\rm ord}\in
  \arg\min_{\operatorname{rank}(B)\le k}\rho_F(A,B),
  \qquad
  \rho_F(A,A_k^{\rm ord})=\|\Sigma_{\rm tail}\|_F.
\]
The same exact Gram-inverse column-sketch projector
\[
  P_{AZ}=AZ\bigl((AZ)^T(AZ)\bigr)^{-1}(AZ)^T
\]
satisfies
\[
  \operatorname{rank}(P_{AZ}A)\le k,
  \qquad
  \rho_F(A,P_{AZ}A)
  \le
  \rho\,\rho_F(A,A_k^{\rm ord}).
\]

This is an exact-object relative theorem surface.  It removes the constructed
source-split handoff from the relative equation-(9) proof, but it still assumes
the exact Eckart--Young tail-optimality inequality displayed above.  It also
does not derive the cross-term certificate from randomness and does not certify
computed non-probability SVD/singular-vector/projector/Gram/inverse/sketch/
product routines.  Sampling probabilities and sampling laws remain exact
mathematical inputs by the current convention.
\end{corollary}

\begin{formalized}
\leanfit{exists_columnSketchGramInverseProjector_sourceHeadTail_sourceSVDTailRelativeResidualSurface_of_rectRightGramOrdered_replacement_tail_left_of_last_pos}
\end{formalized}

\begin{corollary}[Constructed ordered tail-diagonal Frobenius expansion]
Let \(S_{\rm ord}\) be the constructed ordered top-\(k\) right-Gram index set,
let \(S_{\rm ord}^c\) be its complement, and enumerate the complement by
\(\tau(c)\), \(c=1,\ldots,q\).  Define the exact ordered tail diagonal by
\[
  (\Sigma_{\rm tail})_{cd}
  =
  \begin{cases}
    \sigma_{\tau(c)}, & c=d,\\
    0, & c\ne d,
  \end{cases}
\]
where \(\sigma_{\tau(c)}\) is the exact basis-indexed right-Gram singular
value in the complement direction.  Then
\[
  \|\Sigma_{\rm tail}\|_F^2
  =
  \sum_{c=1}^q \sigma_{\tau(c)}^2,
  \qquad
  \|\Sigma_{\rm tail}\|_F
  =
  \left(\sum_{c=1}^q \sigma_{\tau(c)}^2\right)^{1/2}.
\]

More generally, for any exact diagonal table \(D=\operatorname{diag}(\alpha)\),
\[
  \|D\|_F^2=\sum_c \alpha_c^2,
  \qquad
  \|D\|_F=\left(\sum_c \alpha_c^2\right)^{1/2}.
\]

This is the exact-object diagonal-norm expansion needed before the remaining
LR.1dt tail-optimality hypothesis can be discharged by the q-dimensional
Eckart--Young lower-bound route.  It does not prove the residual lower bound
\(\|\Sigma_{\rm tail}\|_F\le \rho_F(A,B)\), does not derive any randomized
cross-term certificate, and does not certify computed non-probability
SVD/singular-vector/projector/Gram/inverse/sketch/product routines.  Sampling
probabilities and sampling laws remain exact mathematical inputs by the current
convention.
\end{corollary}

\begin{formalized}
\leanfit{frobNormSq_diagonal_eq_sum}
\leanfit{frobNorm_diagonal_eq_sqrt_sum}
\leanfit{frobNormSq_rectRightGramOrderedTailSingularDiagonal_eq_sum}
\leanfit{frobNorm_rectRightGramOrderedTailSingularDiagonal_eq_sqrt_sum}
\end{formalized}

\begin{corollary}[Constructed ordered head--tail cardinality bridge]
Let \(e:\{1,\ldots,k\}\hookrightarrow\{1,\ldots,n\}\) be any injective
selected-index map, and let
\[
  S=\{e(a):a=1,\ldots,k\}.
\]
Then the exact selected set and its complement exhaust the ambient right
coordinates:
\[
  |S|=k,
  \qquad
  k+|S^c|=n.
\]
In particular, for the constructed ordered top-\(k\) right-Gram embedding
\(e_{\rm ord}\), if \(q\) denotes the cardinality of the ordered complement
tail index type, then
\[
  k+q=n.
\]

This is an exact-object reindexing dependency for the remaining
Eckart--Young transport.  It supplies only the head-plus-tail cardinality
bridge needed before the q-dimensional lower-bound theorem can be moved to the
original \(n\)-column ordered source split.  It does not construct the column
permutation/equivalence transport, does not prove the residual lower bound,
does not derive randomness, and does not certify computed non-probability
SVD/singular-vector/projector/Gram/inverse/sketch/product routines.  Sampling
probabilities and sampling laws remain exact mathematical inputs by the current
convention.
\end{corollary}

\begin{formalized}
\leanfit{rectRightGramSelectedIndexSet_card_add_compl_card}
\leanfit{rectRightGramOrderedTailIndex_card_add}
\end{formalized}

\begin{corollary}[Constructed ordered head--tail square gap]
Let \(A\in\mathbb R^{m\times n}\), let \(k\le n\), and assume \(0<k\).  Let
\[
  e_{\rm ord}:\{1,\ldots,k\}\hookrightarrow\{1,\ldots,n\}
\]
be the constructed ordered top-\(k\) right-Gram embedding, and let
\[
  \tau:\{1,\ldots,q\}\to \{1,\ldots,n\}
\]
enumerate the complement-tail right-Gram indices
\(\{e_{\rm ord}(1),\ldots,e_{\rm ord}(k)\}^{c}\).  If
\(\sigma_\ell\) denotes the exact right-Gram basis singular value associated
with right coordinate \(\ell\), and \(a_\ast\) is the last selected ordered
head index, define
\[
  \eta = \sigma_{e_{\rm ord}(a_\ast)}^2.
\]
Then \(\eta\) separates the constructed ordered head squares from all
complement-tail squares:
\[
  \eta \le \sigma_{e_{\rm ord}(a)}^2
  \quad\text{for every } a=1,\ldots,k,
  \qquad
  \sigma_{\tau(c)}^2 \le \eta
  \quad\text{for every } c=1,\ldots,q.
\]

The head inequality is proved by the exact selected-square equality for the
constructed ordered embedding and antitonicity of ordered right-Gram
singular-value squares.  The tail inequality is proved by the
complement-versus-selected singular-value comparison, nonnegativity of the
basis singular values, and monotonicity of squaring on nonnegative reals.  Thus
the LR.1dw visible gap hypothesis is instantiated for the constructed ordered
split without assuming that the complement-tail enumeration is internally
sorted.

This is an exact-object gap theorem.  It does not construct the
original-column reindexing/equivalence transport, does not prove the LR.1dt
residual lower-bound/tail-optimality discharge, does not derive randomness, and
does not certify computed non-probability SVD/singular-vector/projector/Gram/
inverse/sketch/product routines.  Sampling probabilities and sampling laws
remain exact mathematical inputs by the current convention.
\end{corollary}

\begin{formalized}
\leanfit{rectRightGramOrdered_head_tail_square_gap}
\end{formalized}

\begin{corollary}[Exact column-permutation transport for low-rank residuals]
Let \(A,B\in\mathbb R^{m\times n}\), let \(r\) be a target rank, and let
\[
  \pi:\{1,\ldots,n\}\simeq\{1,\ldots,n\}
\]
be an exact column equivalence.  Define the column-reindexed matrix
\[
  (A^\pi)_{ij}=A_{i,\pi(j)}.
\]
If \(B\) has an exact rank factorization
\[
  B_{ij}=\sum_{a=1}^{r} L_{ia}R_{aj},
\]
then \(B^\pi\) has the exact rank factorization
\[
  B^\pi_{ij}=\sum_{a=1}^{r} L_{ia}R_{a,\pi(j)}.
\]
Consequently,
\[
  \operatorname{rank}(B)\le r
  \quad\Longrightarrow\quad
  \operatorname{rank}(B^\pi)\le r,
  \qquad
  \operatorname{rank}(B^\pi)\le r
  \quad\Longrightarrow\quad
  \operatorname{rank}(B)\le r.
\]
Moreover, if the same exact column equivalence is applied to the source and
competitor, the Frobenius residual is unchanged:
\[
  \rho_F(A^\pi,B^\pi)=\rho_F(A,B).
\]

This is exact-object column-reindexing infrastructure.  It supplies the
generic transport needed after the \(k+q=n\) cardinality bridge, but this
theorem itself does not construct the specific head-plus-complement-tail
equivalence \(\{1,\ldots,k+q\}\simeq\{1,\ldots,n\}\), does not discharge the LR.1dt
tail-optimality inequality, does not derive randomness, and does not certify
computed non-probability SVD/singular-vector/projector/Gram/inverse/sketch/
product routines.  Sampling probabilities and sampling laws remain exact
mathematical inputs by the current convention.
\end{corollary}

\begin{formalized}
\leanfit{RectRankFactorization.permuteCols}
\leanfit{RectRankAtMost.permuteCols}
\leanfit{RectRankAtMost.of_permuteCols}
\leanfit{lowRankResidualFrob_permuteCols}
\end{formalized}

\begin{corollary}[Constructed ordered head--tail column equivalence]
Let \(k\le n\), let
\[
  e_{\rm ord}:\{1,\ldots,k\}\hookrightarrow\{1,\ldots,n\}
\]
be the constructed ordered top-\(k\) embedding, and define
\[
  S=\{e_{\rm ord}(1),\ldots,e_{\rm ord}(k)\}.
\]
Let \(q=|S^c|\), and let
\[
  \tau:\{1,\ldots,q\}\to\{1,\ldots,n\}
\]
be the canonical increasing enumeration of the complement \(S^c\).  Define
\[
  \Phi:\{1,\ldots,k\}\sqcup\{1,\ldots,q\}\to\{1,\ldots,n\}
\]
by
\[
  \Phi(\operatorname{inl} a)=e_{\rm ord}(a),
  \qquad
  \Phi(\operatorname{inr} c)=\tau(c).
\]
Then \(\Phi\) is injective and surjective.  Hence it induces an exact
equivalence
\[
  \Phi_{\Sigma}:\{1,\ldots,k\}\sqcup\{1,\ldots,q\}
    \simeq \{1,\ldots,n\}.
\]
Composing with the standard finite sum equivalence gives the paper-facing
column equivalence
\[
  \Phi_{\rm fin}:\{1,\ldots,k+q\}\simeq\{1,\ldots,n\},
\]
with the displayed action
\[
  \Phi_{\rm fin}(\operatorname{sumFin}(\operatorname{inl} a))=e_{\rm ord}(a),
  \qquad
  \Phi_{\rm fin}(\operatorname{sumFin}(\operatorname{inr} c))=\tau(c).
\]

The injectivity proof uses disjointness of \(S\) and \(S^c\), injectivity of
the constructed ordered top embedding on the head block, and injectivity of
the canonical complement enumeration on the tail block.  The surjectivity
proof splits every original column coordinate into the selected set \(S\) or
the complement \(S^c\).

This is exact-object column-equivalence infrastructure.  It constructs the
specific head-plus-complement-tail equivalence needed after LR.1dv/LR.1dy, but
it does not discharge the LR.1dt tail-optimality inequality, does not derive
randomness, and does not certify computed non-probability
SVD/singular-vector/projector/Gram/inverse/sketch/product routines.  Sampling
probabilities and sampling laws remain exact mathematical inputs by the
current convention.
\end{corollary}

\begin{formalized}
\leanfit{rectRightGramOrderedHeadTailColumnMap}
\leanfit{rectRightGramOrderedHeadTailColumnMap_injective}
\leanfit{rectRightGramOrderedHeadTailColumnMap_surjective}
\leanfit{rectRightGramOrderedHeadTailColumnSumEquiv}
\leanfit{rectRightGramOrderedHeadTailColumnEquiv}
\leanfit{rectRightGramOrderedHeadTailColumnEquiv_head}
\leanfit{rectRightGramOrderedHeadTailColumnEquiv_tail}
\end{formalized}

\begin{corollary}[Cross-domain exact column-equivalence transport]
Let
\[
  \pi:\{1,\ldots,p\}\simeq\{1,\ldots,n\}
\]
be an exact column equivalence.  For \(A\in\mathbb R^{m\times n}\), define
the pulled-back matrix
\[
  A^\pi_{ij}=A_{i,\pi(j)},\qquad i=1,\ldots,m,\quad j=1,\ldots,p.
\]
If \(B\in\mathbb R^{m\times n}\) has an exact rank factorization
\[
  B_{ij}=\sum_{a=1}^r L_{ia}R_{aj},
\]
then \(B^\pi\in\mathbb R^{m\times p}\) has the exact rank factorization
\[
  B^\pi_{ij}=\sum_{a=1}^r L_{ia}R_{a,\pi(j)}.
\]
Consequently, exact rank-at-most certificates transport forward and backward:
\[
  \operatorname{rank}(B)\le r
  \quad\Longrightarrow\quad
  \operatorname{rank}(B^\pi)\le r,
  \qquad
  \operatorname{rank}(B^\pi)\le r
  \quad\Longrightarrow\quad
  \operatorname{rank}(B)\le r.
\]
Moreover,
\[
  \|A^\pi\|_F=\|A\|_F,
  \qquad
  \rho_F(A^\pi,B^\pi)=\rho_F(A,B).
\]
Specializing \(\pi\) to the constructed ordered head--tail equivalence
\(\Phi_{\rm fin}\) from Corollary~590 gives the same rank and residual
transport for the paper-facing reindexing
\[
  \{1,\ldots,k+q\}\simeq\{1,\ldots,n\}.
\]

This is exact-object cross-domain reindexing infrastructure.  It supplies the
rank/residual transport needed before applying \(q\)-dimensional lower-bound
theorems to the original \(n\)-column ordered source split, but it does not
discharge the LR.1dt tail-optimality inequality, does not derive randomness,
and does not certify computed non-probability
SVD/singular-vector/projector/Gram/inverse/sketch/product routines.  Sampling
probabilities and sampling laws remain exact mathematical inputs by the
current convention.
\end{corollary}

\begin{formalized}
\leanfit{rectReindexCols}
\leanfit{RectRankFactorization.reindexCols}
\leanfit{RectRankAtMost.reindexCols}
\leanfit{RectRankAtMost.of_reindexCols}
\leanfit{frobNormSqRect_reindexCols}
\leanfit{frobNormRect_reindexCols}
\leanfit{lowRankResidualFrob_reindexCols}
\leanfit{RectRankAtMost.reindexCols_rectRightGramOrderedHeadTailColumnEquiv}
\leanfit{RectRankAtMost.of_reindexCols_rectRightGramOrderedHeadTailColumnEquiv}
\leanfit{lowRankResidualFrob_rectRightGramOrderedHeadTailColumnEquiv}
\end{formalized}

\begin{corollary}[Constructed ordered tail optimality for the LR.1dt surface]
Let \(k\le n\), \(0<k\), and suppose the last displayed ordered head singular
value is positive:
\[
  0<\sigma_k(A).
\]
Let \(S\) be the constructed ordered top-\(k\) right-Gram index set, let
\(q=|S^c|\), and let
\[
  e:\{1,\ldots,k\}\sqcup\{1,\ldots,q\}\hookrightarrow\{1,\ldots,m\}
\]
be the exact embedding used by the replacement-tail left-block completion.
Let
\[
  \Sigma_{\rm tail}
    =\operatorname{rectRightGramOrderedTailSingularDiagonal}(A,k)
\]
be the constructed complement-tail diagonal singular-value block.  Then every
exact rank-at-most-\(k\) competitor \(B\in\mathbb R^{m\times n}\) satisfies
\[
  \|\Sigma_{\rm tail}\|_F \le \rho_F(A,B).
\]
Equivalently, the visible LR.1dt tail-optimality hypothesis
\[
  \forall B,\quad \operatorname{rank}(B)\le k
  \Longrightarrow \|\Sigma_{\rm tail}\|_F\le \rho_F(A,B)
\]
is discharged for the constructed ordered replacement-tail source split.

The proof builds the head-first exact source factor on
\(\{1,\ldots,k+q\}\).  Its left block is orthonormal by the replacement-tail
completion certificate, its right block is orthogonal after pulling rows back
along \(\Phi_{\rm fin}\), and the exact source factor equals
\[
  A^{\Phi_{\rm fin}}_{ij}=A_{i,\Phi_{\rm fin}(j)}.
\]
The constructed head-tail square gap then feeds the \(q\)-dimensional
source-factor lower-bound theorem, the tail diagonal norm is rewritten as the
square root of the complement singular-square sum, and the Frobenius residual
is transported back along \(\Phi_{\rm fin}\).

Consequently, under the exact determinant hypothesis
\[
  \det(V_{\rm head}^{T}Z)\ne 0,
\]
the exact cross-term bound
\[
  \|\Sigma_{\rm tail}
      (V_{\rm tail}^{T}Z)(V_{\rm head}^{T}Z)^{-1}\|_F
  \le \varepsilon\|\Sigma_{\rm tail}\|_F,
\]
and the scalar comparison
\[
  2\sqrt{1+\varepsilon^2}\,\|\Sigma_{\rm tail}\|_F
  \le \rho\,\|\Sigma_{\rm tail}\|_F,
\]
the exact Gram-inverse column-sketch projector satisfies the relative surface
\[
  \rho_F\!\left(A,P_{AZ}A\right)
  \le
  \rho\,
  \rho_F\!\left(
    A,
    U_{\rm head}\Sigma_{\rm head}V_{\rm head}^{T}
  \right),
\]
without assuming a separate supplied tail-optimality hypothesis.

This is exact-object D4 infrastructure for equation~(9).  It does not derive
the cross-term radius from randomness, does not prove the scalar comparison,
and does not certify computed non-probability
SVD/singular-vector/projector/Gram/inverse/sketch/product routines.  Sampling
probabilities and sampling laws remain exact mathematical inputs by the
current convention.
\end{corollary}

\begin{formalized}
\leanfit{rectRightGramOrderedHeadTailLeftFinBlock}
\leanfit{rectRightGramOrderedHeadTailSigmaFin}
\leanfit{rectRightGramOrderedHeadTailRightFinBlock}
\leanfit{rectRightGramOrderedHeadTailRightOriginalFinBlock}
\leanfit{rectRightGramOrderedHeadTailLeftFinBlock_col_orthonormal}
\leanfit{rectRightGramOrderedHeadTailRightOriginalFinBlock_col_orthonormal}
\leanfit{rectRightGramOrderedHeadTailRightOriginalFinBlock_row_orthonormal}
\leanfit{rectRightGramOrderedHeadTailRightFinBlock_isOrthogonal}
\leanfit{sourceSVDFactorMatrix_rectRightGramOrderedHeadTailFinBlock_eq_reindexCols}
\leanfit{frobNorm_rectRightGramOrderedTailSingularDiagonal_le_lowRankResidualFrob}
\leanfit{exists_columnSketchGramInverseProjector_sourceHeadTail_sourceSVDTailRelativeResidualSurface_of_rectRightGramOrdered_replacement_tail_left_of_last_pos_tailOptimal}
\end{formalized}

\begin{corollary}[Scalar-relative ordered tail-optimal surface]
Under the hypotheses of Corollary~592, the product-form scalar comparison
\[
  2\sqrt{1+\varepsilon^2}\,\|\Sigma_{\rm tail}\|_F
  \le \rho\,\|\Sigma_{\rm tail}\|_F
\]
may be replaced by the coefficient condition
\[
  2\sqrt{1+\varepsilon^2}\le \rho.
\]
Indeed, for any nonnegative tail radius \(t\),
\[
  2\bigl(\sqrt{1+\varepsilon^2}\,t\bigr)
  =
  \bigl(2\sqrt{1+\varepsilon^2}\bigr)t
  \le \rho t .
\]
Taking \(t=\|\Sigma_{\rm tail}\|_F\), which is nonnegative, gives the
product-form comparison required by Corollary~592.  Consequently, under
\[
  \det(V_{\rm head}^{T}Z)\ne 0,\qquad
  \|\Sigma_{\rm tail}(V_{\rm tail}^{T}Z)
      (V_{\rm head}^{T}Z)^{-1}\|_F
  \le \varepsilon\|\Sigma_{\rm tail}\|_F,
  \qquad
  2\sqrt{1+\varepsilon^2}\le \rho,
\]
the exact Gram-inverse column-sketch projector satisfies
\[
  \rho_F(A,P_{AZ}A)
  \le
  \rho\,
  \rho_F\!\left(
    A,
    U_{\rm head}\Sigma_{\rm head}V_{\rm head}^{T}
  \right),
\]
with no supplied tail-optimality hypothesis and no raw product-form scalar
comparison hypothesis.

This is exact-object D4 theorem-surface cleanup for equation~(9).  It still
does not derive the cross-term radius from randomness and does not certify
computed non-probability SVD/singular-vector/projector/Gram/inverse/sketch/
product routines.  Sampling probabilities and sampling laws remain exact
mathematical inputs by the current convention.
\end{corollary}

\begin{formalized}
\leanfit{two_sqrt_one_add_sq_mul_tail_le_of_scalar}
\leanfit{exists_columnSketchGramInverseProjector_sourceHeadTail_sourceSVDTailRelativeResidualSurface_of_rectRightGramOrdered_replacement_tail_left_of_last_pos_tailOptimal_of_scalarRelative}
\end{formalized}

\begin{corollary}[Rank-nullity kernel vector for the equation~(9) Eckart--Young route]
Let \(B\in\mathbb R^{m\times (r+1)}\) have repository rank at most \(r\),
meaning that there are exact factors
\[
  L\in\mathbb R^{m\times r},\qquad R\in\mathbb R^{r\times (r+1)}
\]
with
\[
  B_{ij}=\sum_{a=1}^r L_{ia}R_{aj}.
\]
Then there exists a nonzero vector \(x\in\mathbb R^{r+1}\) such that
\[
  x\ne 0,
  \qquad
  \sum_{j=1}^{r+1} B_{ij}x_j=0
  \quad\text{for every }i=1,\ldots,m.
\]

Equivalently, every exact `RectRankAtMost \(m\,(r+1)\,r\)' competitor has a
nontrivial right-kernel vector on any displayed \(r+1\)-coordinate right
domain.  The proof forms the exact linear map
\[
  x\longmapsto Rx:\mathbb R^{r+1}\to\mathbb R^r,
\]
uses finite-dimensional rank-nullity to choose \(x\ne 0\) with \(Rx=0\), and
then pushes \(Rx=0\) through the factorization \(B=LR\).

This is the first min--max/rank-nullity foundation for the
Eckart--Young lower-bound route.  It does not prove the singular-value lower
bound, the Frobenius tail-optimality inequality, rectangular SVD existence, or
any computed SVD/singular-vector/projector/Gram/sketch routine.  Sampling
probabilities and sampling laws remain exact mathematical inputs by the current
convention.
\end{corollary}

\begin{formalized}
\leanfit{rectRankFactorization_exists_rightKernelVector_succ}
\leanfit{rectRankAtMost_exists_rightKernelVector_succ}
\end{formalized}

\begin{corollary}[Kernel-to-residual lower bound for the equation~(9) Eckart--Young route]
Let \(A\in\mathbb R^{m\times(r+1)}\), let \(\sigma\in\mathbb R\), and assume
the exact vector-action lower bound
\[
  \sigma\|x\|_2 \le \|Ax\|_2
  \qquad\text{for every nonzero }x\in\mathbb R^{r+1}.
\]
Let \(B\in\mathbb R^{m\times(r+1)}\) have repository rank at most \(r\).
Then
\[
  \sigma \le \|A-B\|_F.
\]
Equivalently, the displayed source block has Frobenius distance at least
\(\sigma\) from every exact rank-\(r\) competitor on the same \(r+1\) right
coordinates.

Indeed, the rank-nullity corollary gives a nonzero vector \(x\) with
\[
  Bx=0.
\]
Therefore
\[
  (A-B)x=Ax.
\]
The exact Frobenius matrix-vector domination theorem gives
\[
  \|Ax\|_2
  =\|(A-B)x\|_2
  \le \|A-B\|_F\|x\|_2.
\]
Combining this with the lower-action hypothesis yields
\[
  \sigma\|x\|_2 \le \|A-B\|_F\|x\|_2.
\]
Since \(x\ne0\), \(\|x\|_2>0\), and cancellation gives
\[
  \sigma \le \|A-B\|_F.
\]

This is an exact-object D4 min--max adapter.  It closes the
kernel-to-residual step after rank-nullity, but it does not prove the
vector-action lower bound from an ordered singular-value block, does not prove
the Frobenius Eckart--Young tail-optimality theorem, does not construct a
rectangular SVD/source split, and does not certify computed
SVD/singular-vector/projector/Gram/sketch routines.  Sampling probabilities and
sampling laws remain exact mathematical inputs by the current convention.
\end{corollary}

\begin{formalized}
\leanfit{rectMatMulVec_sub_eq_left_of_rightKernel}
\leanfit{rectRankAtMost_lowRankResidualFrob_ge_of_vector_lower_bound_succ}
\end{formalized}

\begin{corollary}[Diagonal source-block vector-action lower bound for equation~(9)]
Let \(U\in\mathbb R^{m\times q}\) have exact orthonormal columns,
\[
  U^TU=I_q,
\]
let \(V\in\mathbb R^{q\times q}\) be exact orthogonal, and let
\(\Sigma\in\mathbb R^{q\times q}\) be diagonal:
\[
  \Sigma_{ab}=
  \begin{cases}
    \sigma_a, & a=b,\\
    0, & a\ne b.
  \end{cases}
\]
Assume that \(\sigma\ge 0\) and that every displayed diagonal entry satisfies
\[
  \sigma\le \sigma_a.
\]
Define the exact source block
\[
  A = U\Sigma V^T.
\]
Then, for every \(x\in\mathbb R^q\),
\[
  \sigma\|x\|_2 \le \|Ax\|_2.
\]

Indeed, set \(y=V^Tx\) and \(z=\Sigma y\).  Since \(V\) is orthogonal,
\[
  \|y\|_2=\|x\|_2.
\]
The diagonal hypothesis gives the squared lower bound
\[
  \|z\|_2^2
  =\sum_a \sigma_a^2 y_a^2
  \ge \sum_a \sigma^2 y_a^2
  =\sigma^2\|y\|_2^2.
\]
Because \(\sigma\ge0\), this implies
\[
  \sigma\|y\|_2\le \|z\|_2.
\]
Finally, exact left-column orthonormality gives
\[
  \|Uz\|_2=\|z\|_2,
\]
and \(Ax=U\Sigma V^Tx=Uz\).  Combining the displayed identities yields
\[
  \sigma\|x\|_2 \le \|Ax\|_2.
\]

Thus an exact diagonal source block supplies the vector-action hypothesis used
by the preceding min--max adapter.  This corollary is exact-object spectral
infrastructure: it does not instantiate the ordered singular-value monotonicity
needed for the final Eckart--Young source block, does not prove the Frobenius
tail-optimality theorem, does not construct a rectangular SVD/source split, and
does not certify computed SVD/singular-vector/projector/Gram/sketch/product
routines.  Sampling probabilities and sampling laws remain exact mathematical
inputs by the current convention.
\end{corollary}

\begin{formalized}
\leanfit{vecNorm2Sq_leftOrthonormalFactor}
\leanfit{vecNorm2Sq_diagonal_lower_bound}
\leanfit{rectMatMulVec_sourceSVDFactorMatrix}
\leanfit{sourceSVDFactorMatrix_diagonal_vector_action_lower_bound}
\end{formalized}

\begin{corollary}[Supplied square and thin SVD lower-action specializations for equation~(9)]
Let \(\sigma\ge0\).

\emph{Square case.}  Let \(U,V\in\mathbb R^{q\times q}\) be exact
orthogonal tables, and let \(\sigma_a\) be displayed diagonal entries with
\[
  \sigma\le \sigma_a\qquad\text{for every }a.
\]
Define
\[
  \Sigma_{ab}=
  \begin{cases}
    \sigma_a, & a=b,\\
    0, & a\ne b,
  \end{cases}
  \qquad
  A_{\rm sq}=U\Sigma V^T .
\]
Then for every \(x\in\mathbb R^q\),
\[
  \sigma\|x\|_2\le \|A_{\rm sq}x\|_2 .
\]

\emph{Thin-rectangular case.}  Let \(U\in\mathbb R^{m\times q}\) have exact
orthonormal columns, let \(V\in\mathbb R^{q\times q}\) be exact orthogonal,
and use the same diagonal \(\Sigma\).  Define
\[
  A_{\rm thin}=U\Sigma V^T .
\]
Then for every \(x\in\mathbb R^q\),
\[
  \sigma\|x\|_2\le \|A_{\rm thin}x\|_2 .
\]

Both conclusions are direct mathematical specializations of the previous
diagonal source-block result.  In the square case, the exact orthogonality of
\(U\) supplies \(U^TU=I_q\); in the thin case, this field is supplied as the
explicit column-orthonormality hypothesis.  The diagonal equality is the
literal displayed matrix formula for \(\Sigma\), and the lower-entry hypothesis
is exactly \(\sigma\le\sigma_a\) for every displayed source coordinate.

This is an exact-object specialization for supplied SVD-style data.  It does
not prove existence of an SVD, does not prove singular-value ordering, does not
prove Frobenius Eckart--Young tail optimality, and does not certify computed
non-probability SVD/singular-vector/projector/Gram/sketch/product routines.
Sampling probabilities and sampling laws remain exact mathematical inputs by
the current convention.
\end{corollary}

\begin{formalized}
\leanfit{squareSVD_diagonal_vector_action_lower_bound}
\leanfit{rectangularThinSVD_diagonal_vector_action_lower_bound}
\end{formalized}

\begin{corollary}[Supplied diagonal SVD residual lower bound for equation~(9)]
Let \(\sigma\ge0\).

\emph{Square case.}  Let \(U,V\in\mathbb R^{(r+1)\times(r+1)}\) be exact
orthogonal tables, and let the displayed diagonal entries \(\sigma_a\) satisfy
\[
  \sigma\le \sigma_a\qquad\text{for every }a\in\{1,\ldots,r+1\}.
\]
With
\[
  \Sigma_{ab}=
  \begin{cases}
    \sigma_a, & a=b,\\
    0, & a\ne b,
  \end{cases}
  \qquad
  A_{\rm sq}=U\Sigma V^T,
\]
every rank-at-most-\(r\) competitor \(B\in
\mathbb R^{(r+1)\times(r+1)}\) satisfies
\[
  \sigma\le \|A_{\rm sq}-B\|_F.
\]

\emph{Thin-rectangular case.}  Let \(U\in\mathbb R^{m\times(r+1)}\) have
exact orthonormal columns, let \(V\in\mathbb R^{(r+1)\times(r+1)}\) be exact
orthogonal, and define \(A_{\rm thin}=U\Sigma V^T\).  Then every
rank-at-most-\(r\) competitor \(B\in\mathbb R^{m\times(r+1)}\) satisfies
\[
  \sigma\le \|A_{\rm thin}-B\|_F.
\]

The proof is the exact min--max composition.  The preceding supplied-SVD
lower-action result gives
\[
  \sigma\|x\|_2\le \|A_{\rm sq}x\|_2
  \quad\text{or}\quad
  \sigma\|x\|_2\le \|A_{\rm thin}x\|_2
\]
for every vector \(x\).  The rank-nullity residual adapter then chooses a
nonzero right-kernel vector of any rank-at-most-\(r\) competitor \(B\), rewrites
\((A-B)x=Ax\), and uses Frobenius matrix-vector domination to conclude the
displayed residual lower bound.

This is still an exact-object one-block lower-bound theorem.  It does not prove
the full Frobenius Eckart--Young tail-optimality theorem, singular-value
ordering, SVD/source-split existence, randomness-derived certificates, or
computed non-probability SVD/singular-vector/projector/Gram/sketch/product
routines.  Sampling probabilities and sampling laws remain exact mathematical
inputs by the current convention.
\end{corollary}

\begin{formalized}
\leanfit{rectRankAtMost_lowRankResidualFrob_ge_of_squareSVD_diagonal_succ}
\leanfit{rectRankAtMost_lowRankResidualFrob_ge_of_rectangularThinSVD_diagonal_succ}
\end{formalized}

\begin{corollary}[Ordered right-Gram head residual lower bound for equation~(9)]
Let \(A\in\mathbb R^{m\times n}\) and let \(r+1\le n\).  Let
\(\sigma_{r+1}^{\rm ord}\) denote the last selected ordered right-Gram singular
value in the constructed top-\(r+1\) list, and assume
\[
  \sigma_{r+1}^{\rm ord}>0.
\]
Let \(U_{\rm ord}\in\mathbb R^{m\times(r+1)}\) be the constructed zero-safe
ordered head-left table, and let
\[
  \Sigma_{\rm ord,ab}=
  \begin{cases}
    \tau_a, & a=b,\\
    0, & a\ne b,
  \end{cases}
\]
where \(\tau_a\) is the exact basis-indexed right-Gram singular value selected
by the \(a\)-th ordered top embedding.  Define the displayed coefficient block
\[
  H_{\rm coeff}=U_{\rm ord}\Sigma_{\rm ord}I_{r+1}.
\]
Then every repository rank-at-most-\(r\) competitor
\(B\in\mathbb R^{m\times(r+1)}\) satisfies
\[
  \sigma_{r+1}^{\rm ord}
  \le
  \|H_{\rm coeff}-B\|_F.
\]

The proof is an ordered instantiation of the previous one-block residual
theorem.  Positivity of \(\sigma_{r+1}^{\rm ord}\) gives exact column
orthonormality of \(U_{\rm ord}\).  The right block is the identity, hence
orthogonal.  Ordered singular-value monotonicity gives
\[
  \sigma_{r+1}^{\rm ord}\le \tau_a
  \qquad\text{for every displayed head coordinate }a.
\]
The thin-rectangular supplied-SVD residual lower-bound theorem therefore
applies with \(\sigma=\sigma_{r+1}^{\rm ord}\).

This is an exact-object coefficient-block lower bound.  It is not the full
Frobenius Eckart--Young tail-optimality theorem, does not construct a full
rectangular SVD/source split, does not derive randomness certificates, and does
not certify computed non-probability SVD/singular-vector/projector/Gram/sketch
or downstream product routines.  Sampling probabilities and sampling laws
remain exact mathematical inputs by the current convention.
\end{corollary}

\begin{formalized}
\leanfit{rectRankAtMost_lowRankResidualFrob_ge_of_rectRightGramOrderedHeadDiagonal_succ}
\end{formalized}

\begin{corollary}[Ordered one-step best-rank coefficient block for equation~(9)]
Let \(A\in\mathbb R^{m\times n}\), let \(r+1\le n\), and keep the notation of
the previous corollary.  Assume again that the last selected ordered
right-Gram singular value is positive:
\[
  \sigma_{r+1}^{\rm ord}>0.
\]
Let
\[
  \tau_a =
  \text{the exact basis-indexed right-Gram singular value selected by }a,
  \qquad a\in\{1,\ldots,r+1\},
\]
and set
\[
  H_{\rm coeff}=U_{\rm ord}\operatorname{diag}(\tau_1,\ldots,\tau_{r+1})
  I_{r+1}.
\]
Let \(H_r\) be the supplied first-\(r\) truncation of this coefficient block,
\[
  H_r
  =
  U_{\rm ord}(:,1{:}r)\,
  \operatorname{diag}(\tau_1,\ldots,\tau_r)\,
  I_{r+1}(:,1{:}r)^T .
\]
Then \(H_r\) has rank at most \(r\), and for every repository
rank-at-most-\(r\) competitor \(B\in\mathbb R^{m\times(r+1)}\),
\[
  \|H_{\rm coeff}-H_r\|_F
  \le
  \|H_{\rm coeff}-B\|_F.
\]
Equivalently, \(H_r\) is a Frobenius best rank-\(r\) approximant to the exact
ordered top-\(r+1\) coefficient block.

The one-entry tail identity used in the proof is
\[
  \left\|
    \operatorname{diag}_{q=1}(\tau_{r+1})
  \right\|_F
  =
  \tau_{r+1}
  =
  \sigma_{r+1}^{\rm ord}.
\]
The first equality is the local diagonal Frobenius identity for a one-column
tail, and the second equality is the ordered right-Gram embedding certificate.
The preceding corollary gives the required tail-optimality inequality:
every rank-at-most-\(r\) competitor has residual at least
\(\sigma_{r+1}^{\rm ord}\).  The local thin-rectangular supplied-SVD
best-rank handoff then yields the displayed optimality statement.

This is an exact-object one-step Eckart--Young coefficient-block theorem.  It
does not prove the full multi-tail Frobenius theorem, rectangular
SVD/source-split existence, randomness-derived certificates, or computed
non-probability SVD/singular-vector/projector/Gram/sketch/product routines.
Sampling probabilities and sampling laws remain exact mathematical inputs by
the current convention.
\end{corollary}

\begin{formalized}
\leanfit{frobNorm_squareSVDTailDiagonal_one}
\leanfit{isBestRankApproxFrob_of_rectRightGramOrderedHeadDiagonal_succ}
\end{formalized}

\begin{corollary}[Multi-tail diagonal Frobenius identity for equation~(9)]
Let \(r,q\ge 0\) and let
\[
  \tau : \{1,\ldots,r+q\}\to\mathbb R
\]
be the displayed exact singular-value table for a supplied square or
thin-rectangular SVD-style split.  Let
\[
  \Sigma_{\rm tail}
  =
  \operatorname{diag}(\tau_{r+1},\ldots,\tau_{r+q})
  \in\mathbb R^{q\times q}.
\]
Then
\[
  \|\Sigma_{\rm tail}\|_F^2
  =
  \sum_{c=1}^{q}\tau_{r+c}^2
\]
and therefore
\[
  \|\Sigma_{\rm tail}\|_F
  =
  \left(\sum_{c=1}^{q}\tau_{r+c}^2\right)^{1/2}.
\]

The proof is exact diagonal algebra: all off-diagonal entries of
\(\Sigma_{\rm tail}\) are zero, and the diagonal entries are precisely the
tail singular values selected by the `Fin.natAdd` tail embedding.  This result
supplies the displayed multi-tail norm used by the full Frobenius
Eckart--Young route, but it does not prove the q-dimensional
min-max/tail-optimality inequality itself.  It also does not construct a
rectangular SVD/source split, derive randomness certificates, or certify
computed non-probability SVD/singular-vector/projector/Gram/sketch/product
routines.  Sampling probabilities and sampling laws remain exact mathematical
inputs by the current convention.
\end{corollary}

\begin{formalized}
\leanfit{frobNormSq_squareSVDTailDiagonal_eq_sum}
\leanfit{frobNorm_squareSVDTailDiagonal_eq_sqrt_sum}
\end{formalized}

\begin{corollary}[Q-dimensional right-kernel foundation for equation~(9)]
Let \(r,q\ge 0\), and let
\[
  R:\{1,\ldots,r\}\times\{1,\ldots,r+q\}\to\mathbb R
\]
be an exact right factor.  Define the associated exact linear map
\[
  \mathcal R_R:\mathbb R^{r+q}\to\mathbb R^r,\qquad
  (\mathcal R_R x)_a=\sum_{j=1}^{r+q} R_{aj}x_j .
\]
Then the right kernel has dimension at least \(q\):
\[
  q\le \dim\ker(\mathcal R_R).
\]
Indeed, rank--nullity gives
\[
  \dim\operatorname{range}(\mathcal R_R)+\dim\ker(\mathcal R_R)=r+q,
\]
while \(\dim\operatorname{range}(\mathcal R_R)\le r\).

Now suppose an exact repository rank factorization of
\(B\in\mathbb R^{m\times(r+q)}\) is supplied:
\[
  B_{ij}=\sum_{a=1}^{r} L_{ia}R_{aj}.
\]
If \(x\in\ker(\mathcal R_R)\), then \(x\) is a right-kernel vector of the
represented matrix:
\[
  \forall i,\qquad
  \sum_{j=1}^{r+q}B_{ij}x_j=0.
\]
This follows by substituting the factorization and exchanging the two finite
sums:
\[
  \sum_j B_{ij}x_j
  =
  \sum_a L_{ia}\left(\sum_j R_{aj}x_j\right)
  =
  0.
\]

This is exact-object \(q\)-dimensional rank-nullity infrastructure for the
multi-tail min--max route.  It does not select a \(q\)-dimensional witness
family inside the kernel, prove the tail Frobenius lower bound, prove the full
Eckart--Young theorem, construct rectangular SVD/source-split data, derive
randomness certificates, or certify computed non-probability
SVD/singular-vector/projector/Gram/sketch/product routines.  Sampling
probabilities and sampling laws remain exact mathematical inputs by the
current convention.
\end{corollary}

\begin{formalized}
\leanfit{rectRankRightFactorMap}
\leanfit{rectRankRightFactorMap_ker_finrank_ge}
\leanfit{rectRankFactorization_rightKernel_finrank_ge}
\leanfit{rectRankFactorization_matrix_rightKernel_of_rightFactor_ker}
\end{formalized}

\begin{corollary}[Selected right-kernel family for equation~(9)]
Let \(B\in\mathbb R^{m\times(r+q)}\) be supplied with an exact repository
rank factorization
\[
  B_{ij}=\sum_{a=1}^{r} L_{ia}R_{aj}.
\]
Let \(\mathcal R_R:\mathbb R^{r+q}\to\mathbb R^r\) be the right-factor map
from the preceding corollary.  Then there exists a family
\[
  x_1,\ldots,x_q \in \ker(\mathcal R_R)
\]
that is linearly independent as a family in the kernel subtype, and every
selected vector annihilates the represented competitor:
\[
  \forall c\in\{1,\ldots,q\},\ \forall i\in\{1,\ldots,m\},\qquad
  \sum_{j=1}^{r+q}B_{ij}(x_c)_j=0.
\]

The proof uses the dimension lower bound
\[
  q\le \dim\ker(\mathcal R_R)
\]
from the preceding corollary and applies the finite-dimensional selection
lemma for linearly independent families.  Since each selected vector is a
member of the right-factor kernel, the entrywise annihilation equation follows
from the factorization bridge in the preceding corollary.

This is exact-object vector-selection infrastructure for the multi-tail
min--max route.  It does not prove the \(q\)-dimensional tail lower-bound
argument, the full Frobenius Eckart--Young theorem, rectangular
SVD/source-split construction, randomness-derived certificates, or computed
non-probability SVD/singular-vector/projector/Gram/sketch/product routines.
Sampling probabilities and sampling laws remain exact mathematical inputs by
the current convention.
\end{corollary}

\begin{formalized}
\leanfit{rectRankFactorization_exists_rightKernelFamily}
\end{formalized}

\begin{corollary}[Orthonormal right-kernel family for equation~(9)]
Let \(B\in\mathbb R^{m\times(r+q)}\) be supplied with an exact repository
rank factorization
\[
  B_{ij}=\sum_{a=1}^{r} L_{ia}R_{aj}.
\]
Define the Euclidean-coordinate right-factor map
\[
  \mathcal R_R^{\mathrm E}:\mathbb R^{r+q}_{2}\to\mathbb R^r,
  \qquad
  (\mathcal R_R^{\mathrm E}x)_a=\sum_{j=1}^{r+q}R_{aj}x_j.
\]
Here \(\mathbb R^{r+q}_{2}\) denotes the same coordinate vector space with the
finite-dimensional inner-product structure used for orthonormal-basis
selection.  The coordinate formula is identical to the right-factor map in the
preceding corollaries.  Then
\[
  q\le \dim\ker(\mathcal R_R^{\mathrm E}),
\]
and there exists an orthonormal family
\[
  x_1,\ldots,x_q\in \ker(\mathcal R_R^{\mathrm E})
\]
such that every selected vector annihilates the represented competitor:
\[
  \forall c\in\{1,\ldots,q\},\ \forall i\in\{1,\ldots,m\},\qquad
  \sum_{j=1}^{r+q}B_{ij}(x_c)_j=0.
\]

The proof repeats the rank--nullity lower bound for
\(\mathcal R_R^{\mathrm E}\), chooses the standard orthonormal basis of the
kernel subtype, restricts it to the first \(q\) basis vectors using the
dimension inequality, and uses injectivity of the restriction map to preserve
orthonormality.  Since each selected vector is in the Euclidean-coordinate
right-factor kernel, substituting the exact factorization gives the same
finite-sum annihilation identity as before:
\[
  \sum_j B_{ij}(x_c)_j
  =
  \sum_a L_{ia}\left(\sum_j R_{aj}(x_c)_j\right)
  =0.
\]

This is exact-object min--max-ready vector-selection infrastructure.  It does
not prove the Ky Fan/Courant--Fischer trace/Rayleigh lower bound, the full
Frobenius Eckart--Young theorem, rectangular SVD/source-split construction,
randomness-derived certificates, or computed non-probability
SVD/singular-vector/projector/Gram/sketch/product routines.  Sampling
probabilities and sampling laws remain exact mathematical inputs by the current
convention.
\end{corollary}

\begin{formalized}
\leanfit{rectRankRightFactorEuclideanMap}
\leanfit{rectRankRightFactorEuclideanMap_ker_finrank_ge}
\leanfit{rectRankFactorization_euclideanRightKernel_finrank_ge}
\leanfit{rectRankFactorization_matrix_rightKernel_of_euclideanRightFactor_ker}
\leanfit{rectRankFactorization_exists_orthonormalRightKernelFamily}
\end{formalized}

\begin{corollary}[Orthonormal right-kernel residual energy for equation~(9)]
Let \(M\in\mathbb R^{m\times n}\), and let
\[
  x_1,\ldots,x_q\in \mathbb R^n_2
\]
be an exact orthonormal family.  Then finite Bessel inequality, applied row by
row, gives the Frobenius domination
\[
  \sum_{c=1}^{q}\|Mx_c\|_2^2 \le \|M\|_F^2.
\]
Equivalently, in the repository's coordinate notation,
\[
  \sum_{c=1}^{q}
    \sum_{i=1}^{m}
      \left(\sum_{j=1}^{n}M_{ij}(x_c)_j\right)^2
  \le
  \sum_{i=1}^{m}\sum_{j=1}^{n}M_{ij}^2.
\]

Consequently, let \(A,B\in\mathbb R^{m\times n}\) and suppose the same
orthonormal family lies in the exact right kernel of \(B\):
\[
  \forall c,\ \forall i,\qquad
  \sum_{j=1}^{n}B_{ij}(x_c)_j=0.
\]
Then \((A-B)x_c=Ax_c\) for every \(c\), and hence
\[
  \sum_{c=1}^{q}\|Ax_c\|_2^2
  \le
  \|A-B\|_F^2.
\]
In particular, if \(B\in\mathbb R^{m\times(r+q)}\) is supplied with an exact
repository rank factorization of inner dimension \(r\), the orthonormal
right-kernel family selected in the previous corollary satisfies
\[
  \sum_{c=1}^{q}\|Ax_c\|_2^2
  \le
  \|A-B\|_F^2.
\]

This is the residual side of the \(q\)-dimensional trace/Rayleigh route toward
the equation~(9) tail lower bound.  It does not prove the matching source-side
tail-energy lower bound, the Ky Fan/Courant--Fischer comparison, the full
Frobenius Eckart--Young theorem, rectangular SVD/source-split construction,
randomness-derived certificates, or computed non-probability
SVD/singular-vector/projector/Gram/sketch/product routines.  Sampling
probabilities and sampling laws remain exact mathematical inputs by the current
convention.
\end{corollary}

\begin{formalized}
\leanfit{sum_vecNorm2Sq_rectMatMulVec_le_frobNormSqRect_of_orthonormal}
\leanfit{sum_vecNorm2Sq_rectMatMulVec_lowRankResidual_le_of_orthonormal_rightKernel}
\leanfit{rectRankFactorization_exists_orthonormalRightKernelFamily_energy_le}
\end{formalized}

\begin{corollary}[Diagonal source-tail energy lower bound for equation~(9)]
Let \(r,q\ge 0\), let
\[
  \tau:\{1,\ldots,r+q\}\to\mathbb R,
\qquad
  D=\operatorname{diag}(\tau_1,\ldots,\tau_{r+q}),
\]
and let \(x_1,\ldots,x_q\in\mathbb R^{r+q}_2\) be an exact orthonormal
family.  Suppose there is a visible gap parameter \(\eta\in\mathbb R\) such
that the displayed head and tail diagonal squares satisfy
\[
  \eta \le \tau_a^2 \quad (1\le a\le r),
  \qquad
  \tau_{r+c}^2 \le \eta \quad (1\le c\le q).
\]
Then the diagonal source action on the frame dominates the displayed tail
energy:
\[
  \sum_{c=1}^{q}\tau_{r+c}^2
  \le
  \sum_{d=1}^{q}\|D x_d\|_2^2.
\]
Equivalently, writing
\[
  \alpha_j=\sum_{d=1}^{q}(x_d)_j^2,
\]
the repository proves the expansion
\[
  \sum_{d=1}^{q}\|D x_d\|_2^2
  =
  \sum_{j=1}^{r+q}\tau_j^2\alpha_j,
\]
with coordinate-frame facts
\[
  0\le \alpha_j\le 1,\qquad
  \sum_{j=1}^{r+q}\alpha_j=q.
\]
The scalar mass-transfer step is explicit: the nonnegative head coordinate
mass equals the missing tail coordinate mass, and moving this mass from tail
coordinates with square at most \(\eta\) to head coordinates with square at
least \(\eta\) cannot decrease the weighted sum.

This is the source-side diagonal half of the \(q\)-dimensional
trace/Rayleigh route.  It does not yet instantiate the gap certificate from
ordered singular values, transport the theorem through a nontrivial right
singular-vector table, prove the full Ky Fan/Courant--Fischer theorem, close
Frobenius Eckart--Young optimality, construct rectangular SVD/source-split
data, derive randomness certificates, or certify computed non-probability
SVD/singular-vector/projector/Gram/sketch/product routines.  Sampling
probabilities and sampling laws remain exact mathematical inputs by the current
convention.
\end{corollary}

\begin{formalized}
\leanfit{headTail_weighted_tail_sum_le_of_gap}
\leanfit{orthonormal_sum_coord_sq_le_one}
\leanfit{orthonormal_sum_coord_sq_eq_card}
\leanfit{sum_vecNorm2Sq_diagonal_rectMatMulVec_eq_weighted_coord_sq}
\leanfit{sum_vecNorm2Sq_diagonal_rectMatMulVec_ge_tail_sq_of_orthonormal_gap}
\end{formalized}

\begin{corollary}[Ordered diagonal source-tail energy, positive-head case]
Let \(r,q\ge 0\), assume \(0<r\), and let
\[
  \tau:\{1,\ldots,r+q\}\to\mathbb R,
\qquad
  D=\operatorname{diag}(\tau_1,\ldots,\tau_{r+q}).
\]
Assume the displayed diagonal/singular-value squares are antitone in coordinate
order:
\[
  i\le j\quad\Longrightarrow\quad \tau_j^2\le \tau_i^2
  \qquad (1\le i,j\le r+q).
\]
Let \(x_1,\ldots,x_q\in\mathbb R^{r+q}_2\) be an exact orthonormal family.
Then the diagonal source action satisfies
\[
  \sum_{c=1}^{q}\tau_{r+c}^2
  \le
  \sum_{d=1}^{q}\|D x_d\|_2^2.
\]
The proof first constructs the visible gap parameter
\[
  \eta=\tau_r^2 .
\]
For every head coordinate \(a\le r\), antitonicity gives
\(\eta=\tau_r^2\le \tau_a^2\).  For every tail coordinate \(r+c\),
antitonicity gives \(\tau_{r+c}^2\le \tau_r^2=\eta\).  Thus the hypotheses of
the preceding diagonal mass-transfer theorem hold, and the displayed
tail-energy inequality follows.

This is an exact-object positive-head ordered-diagonal adapter.  It removes
the abstract gap parameter from the LR.1dk statement when the source supplies
an ordered displayed diagonal table, but it does not cover the zero-head edge
case, transport the result through a nontrivial right singular-vector table,
prove the full Rayleigh/Ky Fan or Eckart--Young theorem, construct
rectangular SVD/source-split data, derive randomness certificates, or certify
computed non-probability SVD/singular-vector/projector/Gram/sketch/product
routines.  Sampling probabilities and sampling laws remain exact mathematical
inputs by the current convention.
\end{corollary}

\begin{formalized}
\leanfit{diagonal_headTail_square_gap_of_antitone_head_pos}
\leanfit{sum_vecNorm2Sq_diagonal_rectMatMulVec_ge_tail_sq_of_orthonormal_antitone_head_pos}
\end{formalized}

\begin{corollary}[Zero-head diagonal source-tail energy]
Let \(q\ge 0\), let
\[
  \tau:\{1,\ldots,q\}\to\mathbb R,
\qquad
  D=\operatorname{diag}(\tau_1,\ldots,\tau_q),
\]
and let \(x_1,\ldots,x_q\in\mathbb R^q_2\) be an exact orthonormal family.
Then every coordinate has full squared mass across the frame:
\[
  \sum_{d=1}^{q}(x_d)_j^2 = 1
  \qquad (1\le j\le q).
\]
Consequently the diagonal source action is exactly the displayed zero-head
tail energy:
\[
  \sum_{d=1}^{q}\|D x_d\|_2^2
  =
  \sum_{j=1}^{q}\tau_j^2,
\]
and therefore
\[
  \sum_{j=1}^{q}\tau_j^2
  \le
  \sum_{d=1}^{q}\|D x_d\|_2^2.
\]
The proof combines the Bessel coordinate bound
\(\sum_d(x_d)_j^2\le 1\) with the total-mass identity
\(\sum_j\sum_d(x_d)_j^2=q\).  The nonnegative deficits
\(1-\sum_d(x_d)_j^2\) sum to zero, so each deficit is zero.  The diagonal
energy expansion from the preceding source-tail theorem then gives the
displayed equality.

This is the exact-object \(r=0\) companion to the positive-head ordered-gap
adapter.  It removes the zero-head exception from the diagonal source-side
layer, but it does not transport the result through a nontrivial right
singular-vector table, prove the full Rayleigh/Ky Fan or Eckart--Young
theorem, construct rectangular SVD/source-split data, derive randomness
certificates, or certify computed non-probability
SVD/singular-vector/projector/Gram/sketch/product routines.  Sampling
probabilities and sampling laws remain exact mathematical inputs by the current
convention.
\end{corollary}

\begin{formalized}
\leanfit{orthonormal_sum_coord_sq_eq_one_of_card_eq}
\leanfit{orthonormal_sum_coord_sq_eq_one_of_full}
\leanfit{sum_vecNorm2Sq_diagonal_rectMatMulVec_eq_sum_sq_of_orthonormal_full}
\leanfit{sum_vecNorm2Sq_diagonal_rectMatMulVec_ge_tail_sq_of_orthonormal_zero_head}
\end{formalized}

\begin{corollary}[Combined ordered diagonal source-tail energy]
Let \(r,q\ge 0\), let
\[
  \tau:\{1,\ldots,r+q\}\to\mathbb R,
\qquad
  D=\operatorname{diag}(\tau_1,\ldots,\tau_{r+q}),
\]
and assume the displayed diagonal/singular-value squares are antitone:
\[
  i\le j\quad\Longrightarrow\quad \tau_j^2\le \tau_i^2
  \qquad (1\le i,j\le r+q).
\]
For every exact orthonormal family
\[
  x_1,\ldots,x_q\in\mathbb R^{r+q}_2,
\]
the source-side diagonal lower bound holds:
\[
  \sum_{c=1}^{q}\tau_{r+c}^2
  \le
  \sum_{d=1}^{q}\|D x_d\|_2^2 .
\]
The proof splits on the head count.  If \(0<r\), the last head square
\(\tau_r^2\) supplies the visible gap used in the positive-head ordered
corollary.  If \(r=0\), the frame is a full orthonormal frame in the whole
\(q\)-dimensional coordinate space, so the zero-head equality above applies.

This is the closed exact-object ordered-diagonal source-side layer.  It still
does not transport the result through a nontrivial right singular-vector table,
prove the full Rayleigh/Ky Fan or Eckart--Young theorem, construct rectangular
SVD/source-split data, derive randomness certificates, or certify computed
non-probability SVD/singular-vector/projector/Gram/sketch/product routines.
Sampling probabilities and sampling laws remain exact mathematical inputs by
the current convention.
\end{corollary}

\begin{formalized}
\leanfit{sum_vecNorm2Sq_diagonal_rectMatMulVec_ge_tail_sq_of_orthonormal_antitone}
\end{formalized}

\begin{corollary}[Source-factor ordered source-tail energy]
Let \(r,q\ge 0\), let
\[
  U\in\mathbb R^{m\times(r+q)},\qquad
  V\in\mathbb R^{(r+q)\times(r+q)},\qquad
  D=\operatorname{diag}(\tau_1,\ldots,\tau_{r+q}),
\]
and assume
\[
  U^T U=I_{r+q},\qquad V^T V=VV^T=I_{r+q}.
\]
Assume the displayed diagonal/singular-value squares are antitone:
\[
  i\le j \quad\Longrightarrow\quad \tau_j^2\le \tau_i^2
  \qquad (1\le i,j\le r+q).
\]
Then for every exact orthonormal family
\[
  x_1,\ldots,x_q\in\mathbb R^{r+q}_2,
\]
the transported source-factor lower bound holds:
\[
  \sum_{c=1}^{q}\tau_{r+c}^2
  \le
  \sum_{d=1}^{q}\|UDV^T x_d\|_2^2 .
\]

Indeed, set \(y_d=V^T x_d\).  Exact right orthogonality of \(V\) gives
\[
  \langle y_c,y_d\rangle=\langle x_c,x_d\rangle,
\]
so \(y_1,\ldots,y_q\) is again an exact orthonormal family.  Exact left
column-orthonormality of \(U\) gives, for each \(d\),
\[
  \|UDV^T x_d\|_2^2=\|D y_d\|_2^2.
\]
The result is therefore exactly the ordered diagonal source-tail theorem above
applied to the transported frame \(y_d\).

This is the closed exact-object right-basis transport layer.  It still does
not construct a rectangular SVD/source split, prove the full Rayleigh/Ky Fan or
Eckart--Young theorem, derive randomness certificates, or certify computed
non-probability SVD/singular-vector/projector/Gram/sketch/product routines.
Sampling probabilities and sampling laws remain exact mathematical inputs by
the current convention.
\end{corollary}

\begin{formalized}
\leanfit{inner_matTranspose_mulVec_eq_of_isOrthogonal}
\leanfit{orthonormal_matTranspose_mulVec_of_isOrthogonal}
\leanfit{vecNorm2Sq_sourceSVDFactorMatrix_eq_diagonal_transpose_action}
\leanfit{sum_vecNorm2Sq_sourceSVDFactorMatrix_ge_tail_sq_of_orthonormal_antitone}
\end{formalized}

\begin{corollary}[Source-factor gap lower-bound bridge for equation~(9)]
Let \(r,q\ge 0\), let
\[
  U\in\mathbb R^{m\times(r+q)},\qquad
  V\in\mathbb R^{(r+q)\times(r+q)},\qquad
  D=\operatorname{diag}(\tau_1,\ldots,\tau_{r+q}),
\]
and assume
\[
  U^T U=I_{r+q},\qquad V^T V=VV^T=I_{r+q}.
\]
Suppose there exists a visible separator \(\eta\in\mathbb R\) such that
\[
  \eta\le \tau_a^2 \quad (1\le a\le r),
  \qquad
  \tau_{r+c}^2\le \eta \quad (1\le c\le q).
\]
Then every exact orthonormal family
\[
  x_1,\ldots,x_q\in\mathbb R^{r+q}_2
\]
satisfies the source-factor gap lower bound
\[
  \sum_{c=1}^{q}\tau_{r+c}^2
  \le
  \sum_{d=1}^{q}\|UDV^T x_d\|_2^2 .
\]

Moreover, set \(A=UDV^T\).  For every exact rank-at-most-\(r\) competitor
\(B\in\mathbb R^{m\times(r+q)}\), the squared residual lower bound holds:
\[
  \sum_{c=1}^{q}\tau_{r+c}^2
  \le
  \|A-B\|_F^2,
\]
and therefore the norm form holds:
\[
  \sqrt{\sum_{c=1}^{q}\tau_{r+c}^2}
  \le
  \|A-B\|_F .
\]

The proof of the first inequality sets \(y_d=V^T x_d\).  Exact right
orthogonality of \(V\) makes \(y_1,\ldots,y_q\) an exact orthonormal family,
and exact left column-orthonormality of \(U\) gives
\[
  \|UDV^T x_d\|_2^2=\|D y_d\|_2^2 .
\]
The diagonal gap theorem then applies with the displayed separator \(\eta\).
For the residual statement, the exact rank factorization of \(B\) supplies a
\(q\)-element orthonormal family in the right kernel of \(B\).  The residual
energy theorem gives
\[
  \sum_{d=1}^{q}\|A x_d\|_2^2\le \|A-B\|_F^2,
\]
and the source-factor gap lower bound gives the matching lower inequality for
the same family.  Chaining the inequalities proves the squared statement, and
monotonicity of the square root gives the norm statement.

This is the exact-object form needed for constructed ordered source splits
whose complement-tail enumeration is not itself sorted by singular value.  It
does not instantiate the separator \(\eta\) from the constructed right-Gram
top set, build the original-column reindexing/equivalence transport, prove the
remaining LR.1dt tail-optimality discharge, derive randomness certificates,
or certify computed non-probability SVD/singular-vector/projector/Gram/sketch
or product routines.  Sampling probabilities and sampling laws remain exact
mathematical inputs by the current convention.
\end{corollary}

\begin{formalized}
\leanfit{sum_vecNorm2Sq_sourceSVDFactorMatrix_ge_tail_sq_of_orthonormal_gap}
\leanfit{rectRankAtMost_lowRankResidualFrob_sq_ge_tail_sum_of_sourceSVDFactorMatrix_gap}
\leanfit{sqrt_tail_sum_le_lowRankResidualFrob_of_sourceSVDFactorMatrix_gap}
\end{formalized}

\begin{corollary}[Q-dimensional source-factor lower bound for Eckart--Young]
Use the hypotheses and notation of the source-factor ordered source-tail
corollary above, including exact singular-square antitonicity, and
set
\[
  A=UDV^T .
\]
Let \(B\in\mathbb R^{m\times(r+q)}\) be any exact rank-at-most-\(r\)
competitor.  Then the squared residual lower bound holds:
\[
  \sum_{c=1}^{q}\tau_{r+c}^2
  \le
  \|A-B\|_F^2 .
\]
Equivalently,
\[
  \sqrt{\sum_{c=1}^{q}\tau_{r+c}^2}
  \le
  \|A-B\|_F .
\]

The proof selects, from the exact rank factorization of \(B\), a \(q\)-element
orthonormal family \(x_1,\ldots,x_q\) in the right kernel of \(B\).  The
residual-energy theorem gives
\[
  \sum_{d=1}^{q}\|A x_d\|_2^2
  \le
  \|A-B\|_F^2,
\]
because \(B x_d=0\).  The source-factor ordered source-tail theorem gives the
matching lower bound
\[
  \sum_{c=1}^{q}\tau_{r+c}^2
  \le
  \sum_{d=1}^{q}\|A x_d\|_2^2.
\]
Chaining the two inequalities proves the squared statement; monotonicity of
the square root and nonnegativity of the Frobenius norm give the norm form.

This is an exact-object lower-bound bridge for supplied ordered source-factor
data.  It still does not prove that a source head realizes the matching tail
norm, construct a rectangular SVD/source split, derive randomness
certificates, or certify computed non-probability
SVD/singular-vector/projector/Gram/sketch/product routines.  Sampling
probabilities and sampling laws remain exact mathematical inputs by the
current convention.
\end{corollary}

\begin{formalized}
\leanfit{rectRankAtMost_lowRankResidualFrob_sq_ge_tail_sum_of_sourceSVDFactorMatrix_antitone}
\leanfit{sqrt_tail_sum_le_lowRankResidualFrob_of_sourceSVDFactorMatrix_antitone}
\end{formalized}

\begin{corollary}[Ordered supplied-SVD best-rank adapter for equation~(9)]
Let \(r,q\in\mathbb N\), let
\(\tau:\{1,\ldots,r+q\}\to\mathbb R\), and assume the displayed singular-square
table is antitone:
\[
  i\le j\quad\Longrightarrow\quad \tau_j^2\le \tau_i^2.
\]

For the square case, suppose \(U,V\in\mathbb R^{(r+q)\times(r+q)}\) are exact
orthogonal tables and
\[
  A_{ij}=\sum_{k=1}^{r+q} U_{ik}\tau_k V_{jk}.
\]
Then every exact rank-at-most-\(r\) competitor \(B\) satisfies
\[
  \|D_{\rm tail}\|_F
  =
  \sqrt{\sum_{c=1}^{q}\tau_{r+c}^2}
  \le
  \|A-B\|_F,
\]
where \(D_{\rm tail}=\operatorname{diag}(\tau_{r+1},\ldots,\tau_{r+q})\).
Consequently, if the displayed head singular entries are nonzero, the exact
source head \(U_{\rm head}D_{\rm head}V_{\rm head}^T\) is a Frobenius best
rank-\(r\) approximation to \(A\).  A strict-positive head version follows by
deriving the nonzero-head field from positivity.

For the thin rectangular case, suppose
\(U\in\mathbb R^{m\times(r+q)}\) has exact orthonormal columns, \(V\) is exact
orthogonal, and the same representation holds with \(i=1,\ldots,m\).  Then
the same tail optimality inequality holds for every exact rank-at-most-\(r\)
competitor \(B\in\mathbb R^{m\times(r+q)}\), and the displayed thin source
head is a Frobenius best rank-\(r\) approximation whenever the head entries are
nonzero.

The proof identifies the paper-style diagonal SVD sum with the repository
source-factor object \(U\operatorname{diag}(\tau)V^T\).  The preceding
q-dimensional lower-bound theorem supplies the residual lower bound, and the
tail-diagonal identity rewrites
\(\|D_{\rm tail}\|_F=\sqrt{\sum_c\tau_{r+c}^2}\).  The existing supplied
square/thin SVD best-rank constructors then consume this tail-optimality
inequality.

This is an exact-object, exact-law result for supplied ordered SVD-style data.
It does not prove existence of such an SVD, does not derive the ordering from a
computed spectral routine, does not derive randomness certificates, and does
not certify computed non-probability SVD/singular-vector/projector/Gram/sketch
or product routines.  Sampling probabilities and sampling laws remain exact
mathematical inputs by the current convention.
\end{corollary}

\begin{formalized}
\leanfit{sourceSVDFactorMatrix_diagonal_eq_sum}
\leanfit{squareSVD_sigmaTail_le_lowRankResidualFrob_of_antitone}
\leanfit{isBestRankApproxFrob_of_squareSVD_antitone}
\leanfit{isBestRankApproxFrob_of_squareSVD_head_pos_antitone}
\leanfit{rectangularThinSVD_sigmaTail_le_lowRankResidualFrob_of_antitone}
\leanfit{isBestRankApproxFrob_of_rectangularThinSVD_antitone}
\leanfit{isBestRankApproxFrob_of_rectangularThinSVD_head_pos_antitone}
\end{formalized}

\begin{corollary}[Ordered supplied-SVD relative surfaces for equation~(9)]
Use the square or thin rectangular hypotheses of the preceding supplied-SVD
corollary, including exact singular-square antitonicity.  Let \(Z\) be the
exact sketching matrix and assume the exact full-rank source condition
\[
  \det(V_{\rm head}^T Z)\ne 0.
\]
Assume also the exact cross-term certificate
\[
  \|\Sigma_{\rm tail}(V_{\rm tail}^T Z)(V_{\rm head}^T Z)^{-1}\|_F
  \le
  \varepsilon \|\Sigma_{\rm tail}\|_F,
  \qquad \varepsilon\ge 0,
\]
and the visible scalar comparison
\[
  2\sqrt{1+\varepsilon^2}\,\|\Sigma_{\rm tail}\|_F
  \le
  \rho\,\|\Sigma_{\rm tail}\|_F.
\]
Then the exact Gram-inverse column-sketch projector
\[
  P_{AZ}=AZ\bigl((AZ)^T(AZ)\bigr)^{-1}(AZ)^T
\]
satisfies the relative equation-(9) surface
\[
  \|A-P_{AZ}A\|_F
  \le
  \rho\,\|A-A_r\|_F,
\]
where \(A_r\) is the displayed square or thin supplied-SVD source head.
The theorem also returns the exact best-rank certificate for \(A_r\), the
source-head residual identity
\[
  \|A-A_r\|_F=\|\Sigma_{\rm tail}\|_F,
\]
the exact projector reproduction fields, and the rank-at-most-\(r\) conclusion
for \(P_{AZ}A\).

The only change from the earlier sigma-tail relative surfaces is that the raw
tail-optimality hypothesis is no longer assumed.  It is supplied by the
preceding ordered supplied-SVD best-rank adapter.  The determinant, cross-term,
and scalar comparison assumptions remain visible exact-object hypotheses.

This is an exact-law/exact-object theorem surface.  It does not prove SVD
existence, derive the ordering or cross-term certificate from randomness, or
certify computed non-probability SVD/singular-vector/projector/Gram/inverse/
sketch/product routines.  Sampling probabilities and sampling laws remain exact
mathematical inputs by the current convention.
\end{corollary}

\begin{formalized}
\leanfit{columnSketchGramInverseProjector_sourceHeadTail_sourceSVDTailRelativeResidualSurface_of_squareSVD_sigmaTail_antitone}
\leanfit{columnSketchGramInverseProjector_sourceHeadTail_sourceSVDTailRelativeResidualSurface_of_squareSVD_head_pos_sigmaTail_antitone}
\leanfit{columnSketchGramInverseProjector_sourceHeadTail_sourceSVDTailRelativeResidualSurface_of_rectangularThinSVD_sigmaTail_antitone}
\leanfit{columnSketchGramInverseProjector_sourceHeadTail_sourceSVDTailRelativeResidualSurface_of_rectangularThinSVD_head_pos_sigmaTail_antitone}
\end{formalized}

\begin{corollary}[Source-tail Frobenius norm reduction for equation~(9)]
Let \(U_\perp\in\mathbb R^{m\times q}\) have exact orthonormal columns,
\[
  U_\perp^T U_\perp=I_q,
\]
and let \(C\in\mathbb R^{q\times n}\).  Then
\[
  \|U_\perp C\|_F=\|C\|_F.
\]
Equivalently, the squared Frobenius norms agree:
\[
  \|U_\perp C\|_F^2=\|C\|_F^2.
\]

Now suppose the exact source tail factors as
\[
  T = U_\perp T_c,
\]
where \(T_c\in\mathbb R^{q\times n}\).  For the exact source coefficient table
\[
  W=(V^T Z)^{-1}V^T,
\]
the source residual tail satisfies the exact factorization
\[
  T-(T Z)W
  =
  U_\perp\bigl(T_c-(T_c Z)W\bigr).
\]
Consequently,
\[
  \bigl\|T-(T Z)(V^T Z)^{-1}V^T\bigr\|_F
  =
  \bigl\|T_c-(T_c Z)(V^T Z)^{-1}V^T\bigr\|_F.
\]

This is an exact-object D2 reduction theorem.  It moves the source-tail norm
obligation from the ambient residual tail to the coordinate residual tail.  It
does not yet prove the coordinate residual bound from the \(\Sigma_\perp\),
\(V_\perp^T Z\), and \(V_k^T Z\) factors, does not prove unitarily invariant or
spectral-norm comparison infrastructure, and does not certify computed
SVD/cross-product/inverse/coefficient/projector/product routines.  Sampling
probabilities and laws remain exact mathematical inputs.
\end{corollary}

\begin{formalized}
\leanfit{frobNormSqRect_leftOrthonormalFactor}
\leanfit{frobNormRect_leftOrthonormalFactor}
\leanfit{sourceSketchResidualTail_leftFactor}
\leanfit{frobNormSqRect_sourceSketchResidualTail_leftOrthonormalFactor}
\leanfit{frobNormRect_sourceSketchResidualTail_leftOrthonormalFactor}
\end{formalized}

\begin{corollary}[Coordinate-tail residual factorization for equation~(9)]
Let \(V_k\in\mathbb R^{n\times r}\), \(V_\perp\in\mathbb R^{n\times q}\), and
\(Z\in\mathbb R^{n\times r}\).  Set the exact source coefficient table
\[
  W=(V_k^T Z)^{-1}V_k^T
\]
and the exact rectangular cross factor
\[
  M_\perp=V_\perp^T Z .
\]
For the exact coordinate right-tail residual,
\[
  R_\perp
  =
  V_\perp^T - (V_\perp^T Z)(V_k^T Z)^{-1}V_k^T
  =
  V_\perp^T - M_\perp W,
\]
the formal entrywise statement is
\[
  R_\perp(a,j)
  =
  (V_\perp)_{j,a}
  -
  \sum_{b=1}^r (M_\perp)_{a,b}\, W_{b,j}.
\]

Now let \(\Sigma_\perp\in\mathbb R^{q\times q}\) and define the coordinate tail
\[
  T_c=\Sigma_\perp V_\perp^T.
\]
Then the exact residual induced by the paper's coefficient table factors as
\[
  T_c-(T_c Z)(V_k^T Z)^{-1}V_k^T
  =
  \Sigma_\perp
  \bigl(V_\perp^T-(V_\perp^T Z)(V_k^T Z)^{-1}V_k^T\bigr).
\]
Consequently the repository proves the visible Frobenius bound
\[
  \bigl\|T_c-(T_c Z)(V_k^T Z)^{-1}V_k^T\bigr\|_F
  \le
  \|\Sigma_\perp\|_F\,
  \bigl\|V_\perp^T-(V_\perp^T Z)(V_k^T Z)^{-1}V_k^T\bigr\|_F .
\]

This is an exact-object D2 factorization theorem.  It exposes the
\(\Sigma_\perp\), \(V_\perp^T Z\), and \(V_k^T Z\) factors in the CACM
equation~(9) residual.  The displayed Frobenius bound is deliberately
non-sharp: it uses rectangular Frobenius submultiplicativity through
\(\|\Sigma_\perp\|_F\).  The sharp spectral/unitarily invariant source-tail
bound, the right-orthogonal block invariance step, Eckart--Young, and
computed SVD/cross-product/inverse/coefficient/projector/product certificates
remain open.  Sampling probabilities and laws remain exact mathematical
inputs.
\end{corollary}

\begin{formalized}
\leanfit{sourceRightBasisTranspose}
\leanfit{rightSketchCrossGramRect}
\leanfit{columnSketch_sourceRightBasisTranspose}
\leanfit{columnSketchHead_sourceRightBasisTranspose}
\leanfit{sourceSketchResidualTail_sourceRightBasisTranspose}
\leanfit{sourceSketchResidualTail_leftSquareFactor}
\leanfit{frobNormRect_sourceSketchResidualTail_leftSquareFactor_le}
\leanfit{sourceSketchResidualTail_sigmaRightBasisTranspose_explicit}
\leanfit{frobNormRect_sourceSketchResidualTail_sigmaRightBasisTranspose_le}
\end{formalized}

\begin{corollary}[Right-tail residual block products for equation~(9)]
Let \(V_k\in\mathbb R^{n\times r}\), \(V_\perp\in\mathbb R^{n\times q}\), and
\(Z\in\mathbb R^{n\times r}\).  Set the exact source coefficient table
\[
  W=(V_k^T Z)^{-1}V_k^T
\]
and the exact inverse cross factor
\[
  M_\perp^\#=(V_\perp^T Z)(V_k^T Z)^{-1}.
\]
Assume the exact right-basis hypotheses
\[
  V_\perp^T V_\perp=I_q,\qquad
  V_k^T V_\perp=0,\qquad
  V_\perp^T V_k=0,\qquad
  V_k^T V_k=I_r .
\]
Then the coefficient table has the two exact block actions
\[
  W V_\perp=0,
  \qquad
  W V_k=(V_k^T Z)^{-1}.
\]
For the exact right-tail residual
\[
  R_\perp=V_\perp^T-(V_\perp^T Z)W,
\]
the repository proves the block products
\[
  R_\perp V_\perp=I_q,
  \qquad
  R_\perp V_k=-M_\perp^\#
  =-(V_\perp^T Z)(V_k^T Z)^{-1}.
\]
Equivalently, entrywise,
\[
  \sum_j R_\perp(a,j)(V_\perp)_{j,c}=\delta_{ac},
  \qquad
  \sum_j R_\perp(a,j)(V_k)_{j,c}=-(M_\perp^\#)_{ac}.
\]

This is an exact-object D2 block-product theorem.  It closes the displayed
right-orthogonal algebra needed after the coordinate-tail factorization:
\(R_\perp\) is identity on the tail right-basis block and the negative inverse
cross factor on the head right-basis block.  It still does not prove the sharp
Frobenius/unitarily invariant norm identity or bound that turns these block
products into the CACM equation~(9) source-tail rate.  It also does not prove
rectangular SVD/Eckart--Young foundations or implementation-facing
certificates for computed SVD/singular vectors/cross-products/inverses/
coefficient tables/products.  Sampling probabilities and laws remain exact
mathematical inputs.
\end{corollary}

\begin{formalized}
\leanfit{rightSketchCrossGramRectInvFactor}
\leanfit{sourceSketchCoefficient_mul_rightTailBasis_of_cross_zero}
\leanfit{sourceSketchCoefficient_mul_headRightBasis_of_orthonormal}
\leanfit{sourceRightResidual_mul_rightTailBasis_eq_id}
\leanfit{sourceRightResidual_mul_headRightBasis_eq_neg_invFactor}
\end{formalized}

\begin{corollary}[Right-tail Frobenius block identity for equation~(9)]
Let \(V_k\in\mathbb R^{n\times r}\), \(V_\perp\in\mathbb R^{n\times q}\), and
\(Z\in\mathbb R^{n\times r}\).  Let
\[
  R_\perp=V_\perp^T-(V_\perp^T Z)(V_k^T Z)^{-1}V_k^T,\qquad
  M_\perp^\#=(V_\perp^T Z)(V_k^T Z)^{-1}.
\]
Assume the exact right-basis hypotheses of the preceding corollary and the
row-completeness identity for the concatenated right-basis block
\[
  [V_\perp,V_k][V_\perp,V_k]^T=I_n,
\]
equivalently
\[
  \sum_c (V_\perp)_{j,c}(V_\perp)_{\ell,c}
  +
  \sum_b (V_k)_{j,b}(V_k)_{\ell,b}
  =
  \delta_{j\ell}
  \qquad\text{for all }j,\ell .
\]
Then for every exact square tail singular-value block
\(\Sigma_\perp\in\mathbb R^{q\times q}\),
\[
  \|\Sigma_\perp R_\perp\|_F^2
  =
  \|\Sigma_\perp\|_F^2
  +
  \|\Sigma_\perp M_\perp^\#\|_F^2 .
\]
Equivalently,
\[
  \|\Sigma_\perp R_\perp\|_F
  =
  \sqrt{
    \|\Sigma_\perp\|_F^2
    +
    \|\Sigma_\perp (V_\perp^T Z)(V_k^T Z)^{-1}\|_F^2}.
\]

The formal proof introduces the sum-indexed right-basis block
\([V_\perp,V_k]\).  Right multiplication by this exact row-complete block
preserves squared Frobenius norm, and the preceding block-product theorem gives
\[
  R_\perp[V_\perp,V_k]=[I_q,-M_\perp^\#].
\]
After left multiplication by \(\Sigma_\perp\), the squared Frobenius norm of
\([\Sigma_\perp,-\Sigma_\perp M_\perp^\#]\) splits into the two displayed
summands.

This is an exact-object D2 Frobenius identity theorem.  It closes the
right-orthogonal Frobenius invariance step after the block-product corollary.
It still does not prove the sharp spectral/unitarily invariant bound on
\(\Sigma_\perp(V_\perp^T Z)(V_k^T Z)^{-1}\), rectangular
SVD/Eckart--Young foundations, or implementation-facing certificates for
computed SVD/singular vectors/cross-products/inverses/coefficient tables/
products.  Sampling probabilities and laws remain exact mathematical inputs.
\end{corollary}

\begin{formalized}
\leanfit{finiteFrobNormSq_fin}
\leanfit{finiteFrobNormSq_neg}
\leanfit{finiteFrobNormSq_sumColumns}
\leanfit{finiteFrobNormSq_rectRightOrthonormal}
\leanfit{rightBasisBlock}
\leanfit{rightBasisBlock_row_orthonormal_of_sum}
\leanfit{sourceRightResidualBlock}
\leanfit{sigmaRightResidualBlock}
\leanfit{finiteFrobNormSq_sigmaRightResidualBlock}
\leanfit{sourceRightResidual_rightBasisBlock_eq_block}
\leanfit{sourceRightResidual_sigma_rightBasisBlock_eq_block}
\leanfit{frobNormSqRect_sigma_sourceRightResidual_eq_block}
\leanfit{frobNormRect_sigma_sourceRightResidual_eq_sqrt_block}
\end{formalized}

\begin{corollary}[Right-basis block orthonormality certificate for equation~(9)]
Let the exact concatenated right-basis block be
\[
  Q=[V_\perp,V_k],
  \qquad
  Q_{j,\operatorname{inl}(c)}=(V_\perp)_{j,c},
  \qquad
  Q_{j,\operatorname{inr}(b)}=(V_k)_{j,b}.
\]
Assume exact column orthonormality of this block:
\[
  \sum_j Q_{j,\mu}Q_{j,\nu}
  =
  \begin{cases}
    1,&\mu=\nu,\\
    0,&\mu\ne\nu.
  \end{cases}
\]
Then the component right-basis fields used by the source-tail residual proof
follow:
\[
  V_\perp^T V_\perp=I_q,\qquad
  V_k^T V_\perp=0,\qquad
  V_\perp^T V_k=0,\qquad
  V_k^T V_k=I_r .
\]
Assume also exact row orthonormality of the same block,
\[
  \sum_\mu Q_{j,\mu}Q_{\ell,\mu}=\delta_{j\ell}.
\]
Equivalently,
\[
  \sum_c (V_\perp)_{j,c}(V_\perp)_{\ell,c}
  +
  \sum_b (V_k)_{j,b}(V_k)_{\ell,b}
  =
  \delta_{j\ell}.
\]

Consequently, if \(\varepsilon\ge0\) and the exact source cross-term radius
satisfies
\[
  \|\Sigma_\perp (V_\perp^T Z)(V_k^T Z)^{-1}\|_F
  \le
  \varepsilon\|\Sigma_\perp\|_F,
\]
then
\[
  \|\Sigma_\perp
    (V_\perp^T-(V_\perp^T Z)(V_k^T Z)^{-1}V_k^T)\|_F
  \le
  \sqrt{1+\varepsilon^2}\,\|\Sigma_\perp\|_F .
\]

This is an exact-object D3-to-D2 theorem.  It reduces the several
right-basis component hypotheses in the source cross-term proof to one
column/row orthonormality certificate for the exact block \([V_\perp,V_k]\).
It still does not construct that block from a rectangular SVD, prove
singular-value ordering or Eckart--Young optimality, derive the cross-term
radius from randomness, or certify computed SVD/singular-vector/cross-product/
inverse/product routines.  Sampling probabilities and laws remain exact
mathematical inputs.
\end{corollary}

\begin{formalized}
\leanfit{rightBasisBlock_component_orthonormal_fields_of_col_orthonormal}
\leanfit{rightBasisBlock_complete_sum_of_row_orthonormal}
\leanfit{frobNormRect_sigma_sourceRightResidual_le_sqrt_one_add_eps_sq_of_rightBasisBlock_orthonormal}
\end{formalized}

\begin{corollary}[Right-basis block assembly from component fields for equation~(9)]
Let \(V_k\in\mathbb R^{n\times r}\) and
\(V_\perp\in\mathbb R^{n\times q}\), and define the exact concatenated block
\[
  Q=[V_\perp,V_k],
  \qquad
  Q_{j,\operatorname{inl}(c)}=(V_\perp)_{j,c},
  \qquad
  Q_{j,\operatorname{inr}(b)}=(V_k)_{j,b}.
\]
Assume the exact component right-basis fields
\[
  V_\perp^T V_\perp=I_q,\qquad
  V_k^T V_\perp=0,\qquad
  V_\perp^T V_k=0,\qquad
  V_k^T V_k=I_r .
\]
Then \(Q\) is column orthonormal:
\[
  \sum_j Q_{j,\mu}Q_{j,\nu}
  =
  \begin{cases}
    1,&\mu=\nu,\\
    0,&\mu\ne\nu.
  \end{cases}
\]
If, in addition, the component row-completeness identity holds,
\[
  \sum_c (V_\perp)_{j,c}(V_\perp)_{\ell,c}
  +
  \sum_b (V_k)_{j,b}(V_k)_{\ell,b}
  =
  \delta_{j\ell}
  \qquad\text{for all }j,\ell ,
\]
then \(Q\) is row orthonormal:
\[
  \sum_\mu Q_{j,\mu}Q_{\ell,\mu}=\delta_{j\ell}.
\]

Consequently, if \(Z\in\mathbb R^{n\times r}\), \(\varepsilon\ge0\), and the
exact source cross-term radius satisfies
\[
  \|\Sigma_\perp (V_\perp^T Z)(V_k^T Z)^{-1}\|_F
  \le
  \varepsilon\|\Sigma_\perp\|_F,
\]
then the assembled block certificate gives
\[
  \|\Sigma_\perp
    (V_\perp^T-(V_\perp^T Z)(V_k^T Z)^{-1}V_k^T)\|_F
  \le
  \sqrt{1+\varepsilon^2}\,\|\Sigma_\perp\|_F .
\]

This is an exact-object D3 source-SVD assembly theorem.  It proves the
direction a rectangular SVD construction would naturally supply: separate
component orthogonality and row-completeness fields assemble the block
orthonormality certificate used by the preceding corollary.  It still does not
construct \(V_\perp,V_k\) or \(\Sigma_\perp\) from a rectangular SVD, prove
singular-value ordering or Eckart--Young optimality, derive the cross-term
radius from randomness, or certify computed SVD/singular-vector/cross-product/
inverse/product routines.  Sampling probabilities and laws remain exact
mathematical inputs.
\end{corollary}

\begin{formalized}
\leanfit{rightBasisBlock_col_orthonormal_of_component_orthonormal_fields}
\leanfit{rightBasisBlock_col_row_orthonormal_of_component_fields}
\leanfit{frobNormRect_sigma_sourceRightResidual_le_sqrt_one_add_eps_sq_of_component_block_assembly}
\end{formalized}

\begin{corollary}[Source cross-term certificate for equation~(9)]
Keep the exact hypotheses and notation of the preceding corollary.  Let
\(\varepsilon\ge 0\), and assume the Frobenius cross-term condition appearing
in the CACM structural theorem:
\[
  \|\Sigma_\perp (V_\perp^T Z)(V_k^T Z)^{-1}\|_F
  \le
  \varepsilon \|\Sigma_\perp\|_F .
\]
Then the exact right-tail residual satisfies the source-tail rate
\[
  \|\Sigma_\perp
    (V_\perp^T-(V_\perp^T Z)(V_k^T Z)^{-1}V_k^T)\|_F
  \le
  \sqrt{1+\varepsilon^2}\,\|\Sigma_\perp\|_F .
\]

Equivalently, with
\[
  R_\perp=V_\perp^T-(V_\perp^T Z)(V_k^T Z)^{-1}V_k^T,
  \qquad
  M_\perp^\#=(V_\perp^T Z)(V_k^T Z)^{-1},
\]
the exact implication is
\[
  \|\Sigma_\perp M_\perp^\#\|_F
  \le
  \varepsilon\|\Sigma_\perp\|_F
  \quad\Longrightarrow\quad
  \|\Sigma_\perp R_\perp\|_F
  \le
  \sqrt{1+\varepsilon^2}\,\|\Sigma_\perp\|_F .
\]

The proof combines the squared block identity above
\[
  \|\Sigma_\perp R_\perp\|_F^2
  =
  \|\Sigma_\perp\|_F^2+\|\Sigma_\perp M_\perp^\#\|_F^2
\]
with the squared cross-term bound and monotonicity of the square root.  This is
an exact-object structural-certificate theorem: the cross-term inequality is a
visible hypothesis from the source equation~(9), not a consequence of the
full-rank condition by itself.  It does not prove the spectral-norm or general
unitarily invariant norm variants, rectangular SVD/Eckart--Young foundations,
or implementation-facing certificates for computed SVD/singular vectors/
cross-products/inverses/coefficient tables/products.  Sampling probabilities
and laws remain exact mathematical inputs.
\end{corollary}

\begin{formalized}
\leanfit{frobNormRect_sigma_sourceRightResidual_le_sqrt_one_add_eps_sq}
\end{formalized}

\begin{corollary}[Ambient source-tail certificate for equation~(9)]
Let \(T\in\mathbb R^{m\times n}\) be an exact source-tail matrix, and assume
that it has the exact source-tail factorization
\[
  T=U_\perp\Sigma_\perp V_\perp^T,
  \qquad
  U_\perp^T U_\perp=I_q .
\]
Keep the exact right-basis hypotheses, row-completeness identity, and
cross-term certificate of the preceding corollary:
\[
  \varepsilon\ge 0,\qquad
  \|\Sigma_\perp (V_\perp^T Z)(V_k^T Z)^{-1}\|_F
  \le
  \varepsilon\|\Sigma_\perp\|_F .
\]
Let \(W=(V_k^T Z)^{-1}V_k^T\).  Then the exact ambient source-tail residual
satisfies
\[
  \|T-(TZ)W\|_F
  =
  \|T-(TZ)(V_k^T Z)^{-1}V_k^T\|_F
  \le
  \sqrt{1+\varepsilon^2}\,\|\Sigma_\perp\|_F .
\]

Equivalently, this result composes the three exact reductions
\[
  \|T-(TZ)W\|_F
  =
  \|\Sigma_\perp(V_\perp^T-(V_\perp^T Z)W)\|_F,
  \qquad
  W=(V_k^T Z)^{-1}V_k^T,
\]
with the displayed CACM cross-term certificate.  It is an exact-object
source-tail theorem: the left and right singular-vector blocks, the
singular-value block, the inverse factor, and the products are supplied as
exact mathematical objects.  It does not construct a rectangular SVD/source
split, prove Eckart--Young, prove the spectral-norm or general unitarily
invariant variants, bound the projector-applied coupling tail, or certify any
computed non-probability SVD/projector/cross-product/Gram/inverse/coefficient/
product routine.  Sampling probabilities and laws remain exact mathematical
inputs.
\end{corollary}

\begin{formalized}
\leanfit{frobNormRect_sourceSketchResidualTail_sourceSVDTail_le_sqrt_one_add_eps_sq}
\end{formalized}

\begin{corollary}[Projected source-tail coupling certificate for equation~(9)]
Let \(P\in\mathbb R^{m\times m}\) be an exact symmetric idempotent row
multiplier:
\[
  P^T=P,
  \qquad
  P^2=P .
\]
Then for every \(R\in\mathbb R^{m\times n}\),
\[
  \|PR\|_F\le \|R\|_F .
\]
In particular, let \(B=AZ\) and \(P=(AZ)C\), and suppose that \(C\) carries an
exact orthogonal column-sketch projector certificate for \(B\), equivalently the
multiplier \(P\) is symmetric, idempotent, and reproduces the sketch.  The same
contraction gives
\[
  \|P R\|_F\le \|R\|_F
  \qquad\text{for all }R\in\mathbb R^{m\times n}.
\]

Now keep all exact source-tail hypotheses of the preceding corollary:
\[
  T=U_\perp\Sigma_\perp V_\perp^T,
  \qquad
  U_\perp^T U_\perp=I_q,
\]
the exact right-basis and row-completeness hypotheses, and the exact
cross-term certificate
\[
  \varepsilon\ge0,
  \qquad
  \|\Sigma_\perp (V_\perp^T Z)(V_k^T Z)^{-1}\|_F
  \le
  \varepsilon\|\Sigma_\perp\|_F .
\]
With \(W=(V_k^T Z)^{-1}V_k^T\), if \(C\) carries either the orthogonal-projector
certificate for \(P=(AZ)C\) or the four-equation Moore--Penrose certificate for
the sketch \(B=AZ\), then the projected coupling tail satisfies
\[
  \|P(T-(TZ)W)\|_F
  =
  \|P(T-(TZ)(V_k^T Z)^{-1}V_k^T)\|_F
  \le
  \sqrt{1+\varepsilon^2}\,\|\Sigma_\perp\|_F .
\]

This is an exact-object coupling theorem.  It closes the Frobenius
projector-applied tail handoff for supplied exact source-tail and projector
certificates.  It does not construct a rectangular SVD/source split, prove
Eckart--Young, derive the cross-term certificate from randomness, prove the
spectral-norm or general unitarily invariant variants, or certify computed
non-probability SVD/projector/Gram/cross-product/inverse/coefficient/product
routines.  Sampling probabilities and laws remain exact mathematical inputs.
\end{corollary}

\begin{formalized}
\leanfit{finiteVecNorm2_finiteMatVec_le_of_symmetric_idempotent}
\leanfit{frobNormRect_preconditionRows_le_of_symmetric_idempotent}
\leanfit{frobNormRect_preconditionRows_columnSketchLeftMultiplier_le_of_orthogonalProjectorCertificate}
\leanfit{frobNormRect_preconditionRows_sourceSketchResidualTail_sourceSVDTail_orthogonalProjector_le_sqrt_one_add_eps_sq}
\leanfit{frobNormRect_preconditionRows_sourceSketchResidualTail_sourceSVDTail_moorePenrose_le_sqrt_one_add_eps_sq}
\end{formalized}

\begin{corollary}[Transpose-action spectral certificate for equation~(9)]
Let \(M\in\mathbb R^{q\times r}\), let
\(\Sigma\in\mathbb R^{q\times q}\), and assume
\(\varepsilon\ge0\).  Suppose the exact transpose action of \(M\) has
operator-2 radius \(\varepsilon\):
\[
  \|M^T x\|_2\le \varepsilon\|x\|_2
  \qquad\text{for every }x\in\mathbb R^q .
\]
Then
\[
  \|\Sigma M\|_F
  \le
  \varepsilon\|\Sigma\|_F .
\]
Equivalently, the squared form is
\[
  \|\Sigma M\|_F^2
  \le
  \varepsilon^2\|\Sigma\|_F^2 .
\]

In the CACM equation~(9) notation, put
\[
  M=(V_\perp^T Z)(V_k^T Z)^{-1}.
\]
If the exact transpose-action certificate
\[
  \|M^T x\|_2\le \varepsilon\|x\|_2
  \qquad(x\in\mathbb R^q)
\]
holds, then the source cross-term hypothesis used above follows:
\[
  \|\Sigma_\perp (V_\perp^T Z)(V_k^T Z)^{-1}\|_F
  \le
  \varepsilon\|\Sigma_\perp\|_F .
\]
Consequently, under the exact source-tail hypotheses
\[
  T=U_\perp\Sigma_\perp V_\perp^T,\qquad
  U_\perp^T U_\perp=I_q,
\]
the exact right-basis and row-completeness hypotheses, and a Moore--Penrose
certificate for \(C\) relative to \(B=AZ\), the projected source-tail coupling
bound can be stated directly from the transpose-action operator certificate:
\[
  \|(AZ)C\bigl(T-(TZ)(V_k^T Z)^{-1}V_k^T\bigr)\|_F
  \le
  \sqrt{1+\varepsilon^2}\,\|\Sigma_\perp\|_F .
\]

This is an exact-object-plus-certificate theorem.  The singular-vector blocks,
cross products, inverse factor, projector coefficient table, and matrix
products are exact mathematical objects with exact certificates.  The result
does not prove equivalence with an ordinary spectral-norm certificate for the
non-transposed factor \(M\), does not prove general unitarily invariant
variants, does not derive the operator certificate from the randomness of
\(Z\), and does not certify computed non-probability
SVD/projector/Gram/cross-product/inverse/coefficient/product routines.
Sampling probabilities and laws remain exact mathematical inputs.
\end{corollary}

\begin{formalized}
\leanfit{frobNormSqRect_matMulRectLeft_le_sq_mul_of_transpose_rectOpNorm2Le}
\leanfit{frobNormRect_matMulRectLeft_le_of_transpose_rectOpNorm2Le}
\leanfit{frobNormRect_sigma_rightSketchCrossGramRectInvFactor_le_of_transpose_rectOpNorm2Le}
\leanfit{frobNormRect_preconditionRows_sourceSketchResidualTail_sourceSVDTail_moorePenrose_le_sqrt_one_add_eps_sq_of_transpose_rectOpNorm2Le}
\end{formalized}

\begin{corollary}[Ordinary operator certificate for equation~(9)]
Let \(M\in\mathbb R^{m\times n}\), let \(\varepsilon\ge0\), and suppose the
ordinary rectangular operator certificate holds:
\[
  \|Mx\|_2\le \varepsilon\|x\|_2
  \qquad\text{for every }x\in\mathbb R^n .
\]
Then the transpose action has the same radius:
\[
  \|M^T y\|_2\le \varepsilon\|y\|_2
  \qquad\text{for every }y\in\mathbb R^m .
\]
Consequently, for every \(\Sigma\in\mathbb R^{m\times m}\),
\[
  \|\Sigma M\|_F\le \varepsilon\|\Sigma\|_F .
\]

For the CACM equation~(9) cross factor
\[
  M=(V_\perp^T Z)(V_k^T Z)^{-1},
\]
the non-transposed certificate
\[
  \|(V_\perp^T Z)(V_k^T Z)^{-1}x\|_2
  \le
  \varepsilon\|x\|_2
  \qquad(x\in\mathbb R^r)
\]
therefore implies
\[
  \|\Sigma_\perp (V_\perp^T Z)(V_k^T Z)^{-1}\|_F
  \le
  \varepsilon\|\Sigma_\perp\|_F .
\]
Under the exact source-tail, right-basis, row-completeness, and
Moore--Penrose sketch certificates of the preceding corollaries, it follows
directly that
\[
  \|(AZ)C\bigl(T-(TZ)(V_k^T Z)^{-1}V_k^T\bigr)\|_F
  \le
  \sqrt{1+\varepsilon^2}\,\|\Sigma_\perp\|_F .
\]

This is an exact-object operator-certificate theorem.  It removes the
transpose-action assumption from the previous corollary when the supplied
certificate is the ordinary non-transposed rectangular operator predicate.
It still does not prove general unitarily invariant variants, derive the
operator certificate from the randomness of \(Z\), construct the rectangular
SVD/source split or Eckart--Young certificate, or certify computed
non-probability SVD/projector/Gram/cross-product/inverse/coefficient/product
routines.  Sampling probabilities and laws remain exact mathematical inputs.
\end{corollary}

\begin{formalized}
\leanfit{rectOpNorm2Le_finiteTranspose_of_rectOpNorm2Le}
\leanfit{frobNormRect_sigma_rightSketchCrossGramRectInvFactor_le_of_rectOpNorm2Le}
\leanfit{frobNormRect_preconditionRows_sourceSketchResidualTail_sourceSVDTail_moorePenrose_le_sqrt_one_add_eps_sq_of_rectOpNorm2Le}
\end{formalized}

\begin{corollary}[Computed cross-factor perturbation for equation~(9)]
Let \(M\in\mathbb R^{q\times r}\) be the exact analysis cross factor and let
\(\widehat M\in\mathbb R^{q\times r}\) be the computed non-probability factor
actually supplied by an implementation.  Let \(\varepsilon\ge0\), and suppose
that the computed factor has the ordinary operator certificate
\[
  \|\widehat Mx\|_2\le\varepsilon\|x\|_2
  \qquad\text{for every }x\in\mathbb R^r .
\]
Suppose also that the computed factor approximates the exact analysis factor
with Frobenius radius \(\tau\):
\[
  \|M-\widehat M\|_F\le\tau .
\]
Then for every \(\Sigma\in\mathbb R^{q\times q}\),
\[
  \|\Sigma M\|_F
  \le
  (\varepsilon+\tau)\|\Sigma\|_F .
\]

In the CACM equation~(9) notation, take
\[
  M=(V_\perp^T Z)(V_k^T Z)^{-1}.
\]
If a concrete routine returns a computed factor \(\widehat M\) satisfying
\[
  \|\widehat Mx\|_2\le\varepsilon\|x\|_2
  \quad(x\in\mathbb R^r),
  \qquad
  \|(V_\perp^T Z)(V_k^T Z)^{-1}-\widehat M\|_F\le\tau ,
\]
then the exact source cross-term certificate holds with enlarged radius:
\[
  \|\Sigma_\perp (V_\perp^T Z)(V_k^T Z)^{-1}\|_F
  \le
  (\varepsilon+\tau)\|\Sigma_\perp\|_F .
\]
Consequently, under the exact source-tail, right-basis, row-completeness, and
Moore--Penrose sketch certificates of the preceding corollaries,
\[
  \|(AZ)C\bigl(T-(TZ)(V_k^T Z)^{-1}V_k^T\bigr)\|_F
  \le
  \sqrt{1+(\varepsilon+\tau)^2}\,\|\Sigma_\perp\|_F .
\]

This is an implementation-facing certificate transfer for computed
non-probability cross-factor data.  The sampling probabilities and sampling
law for \(Z\) remain exact mathematical inputs.  The theorem does not claim
that a particular floating-point routine has produced \(\widehat M\), nor does
it certify the computed singular vectors, projector, Gram matrix, inverse, or
coefficient table; those concrete routine instantiations remain not-proved
ledger obligations.
\end{corollary}

\begin{formalized}
\leanfit{frobNormRect_sigma_exactFactor_le_of_computed_rectOpNorm2Le_of_frobNormRect_error}
\leanfit{frobNormRect_sigma_rightSketchCrossGramRectInvFactor_le_of_computed_rectOpNorm2Le_of_frobNormRect_error}
\leanfit{frobNormRect_preconditionRows_sourceSketchResidualTail_sourceSVDTail_moorePenrose_le_sqrt_one_add_eps_plus_tau_sq_of_computed_rectOpNorm2Le}
\end{formalized}

\begin{corollary}[Entrywise computed cross-factor radius for equation~(9)]
Let \(M,\widehat M\in\mathbb R^{q\times r}\), let
\(\eta\ge0\), and suppose a concrete computed non-probability routine supplies
the uniform entrywise certificate
\[
  |M_{ab}-\widehat M_{ab}|\le\eta
  \qquad
  (1\le a\le q,\;1\le b\le r).
\]
Then
\[
  \|M-\widehat M\|_F
  \le
  \sqrt{qr}\,\eta .
\]
Consequently, if the computed factor also satisfies the ordinary operator
certificate
\[
  \|\widehat Mx\|_2\le\varepsilon\|x\|_2
  \qquad(x\in\mathbb R^r),
  \qquad \varepsilon\ge0,
\]
then every \(\Sigma\in\mathbb R^{q\times q}\) satisfies
\[
  \|\Sigma M\|_F
  \le
  \bigl(\varepsilon+\sqrt{qr}\,\eta\bigr)\|\Sigma\|_F .
\]

For the exact CACM cross factor
\[
  M=(V_\perp^T Z)(V_k^T Z)^{-1},
\]
the same hypotheses on the computed factor \(\widehat M\) give
\[
  \|\Sigma_\perp (V_\perp^T Z)(V_k^T Z)^{-1}\|_F
  \le
  \bigl(\varepsilon+\sqrt{qr}\,\eta\bigr)\|\Sigma_\perp\|_F .
\]
Under the exact source-tail, right-basis, row-completeness, and Moore--Penrose
sketch certificates used above, this yields the projected equation~(9) bound
\[
  \|(AZ)C\bigl(T-(TZ)(V_k^T Z)^{-1}V_k^T\bigr)\|_F
  \le
  \sqrt{1+\bigl(\varepsilon+\sqrt{qr}\,\eta\bigr)^2}\,
  \|\Sigma_\perp\|_F .
\]

This is an implementation-facing entrywise-error transfer theorem for computed
non-probability cross-factor data.  Sampling probabilities and sampling laws
remain exact mathematical inputs.  The theorem does not instantiate a concrete
floating-point routine for forming \(\widehat M\); such a routine must still
prove the entrywise budget, the computed operator certificate, and the
computed SVD/singular-vector/projector/Gram/inverse/coefficient-table
obligations recorded in the not-proved ledger.
\end{corollary}

\begin{formalized}
\leanfit{frobNormRect_le_sqrt_mul_nat_of_entry_abs_le}
\leanfit{frobNormRect_sigma_rightSketchCrossGramRectInvFactor_le_of_computed_rectOpNorm2Le_of_entry_abs_error}
\leanfit{frobNormRect_preconditionRows_sourceSketchResidualTail_sourceSVDTail_moorePenrose_le_sqrt_one_add_eps_plus_entry_sq_of_computed_rectOpNorm2Le}
\end{formalized}

\begin{corollary}[Component-certified computed cross factor for equation~(9)]
Let
\[
  X=V_\perp^T Z,\qquad
  Y=(V_k^T Z)^{-1},
  \qquad M=XY .
\]
Let \(\widehat X\), \(\widehat Y\), and \(\widehat M\) be the computed
non-probability cross-gram, inverse factor, and rounded product supplied by an
implementation.  Suppose the component certificates
\[
  \sum_b |X_{ab}-\widehat X_{ab}|\,|Y_{bc}|\le\alpha ,
\]
\[
  \sum_b |\widehat X_{ab}|\,|Y_{bc}-\widehat Y_{bc}|\le\beta ,
\]
and
\[
  \left|\sum_b \widehat X_{ab}\widehat Y_{bc}
    -\widehat M_{ac}\right|\le\rho
\]
hold for every \(a,c\), with
\(\alpha,\beta,\rho\ge0\).  Then
\[
  |M_{ac}-\widehat M_{ac}|
  \le
  \alpha+\beta+\rho
  \qquad (a,c).
\]
Consequently, if \(\widehat M\) also satisfies
\[
  \|\widehat Mx\|_2\le\varepsilon\|x\|_2
  \qquad(x\in\mathbb R^r),
  \qquad \varepsilon\ge0,
\]
then
\[
  \|\Sigma_\perp (V_\perp^T Z)(V_k^T Z)^{-1}\|_F
  \le
  \bigl(\varepsilon+\sqrt{qr}\,(\alpha+\beta+\rho)\bigr)
  \|\Sigma_\perp\|_F .
\]
Under the exact source-tail, right-basis, row-completeness, and Moore--Penrose
sketch certificates used above,
\[
  \|(AZ)C\bigl(T-(TZ)(V_k^T Z)^{-1}V_k^T\bigr)\|_F
  \le
  \sqrt{1+
    \bigl(\varepsilon+\sqrt{qr}\,(\alpha+\beta+\rho)\bigr)^2}\,
  \|\Sigma_\perp\|_F .
\]

This is an implementation-facing component-certificate theorem for the
computed cross factor.  It charges the computed cross-gram, inverse factor, and
rounded product, but it still does not instantiate a particular floating-point
dot-product routine for \(V_\perp^T Z\), a concrete inverse routine for
\((V_k^T Z)^{-1}\), a final matrix-product routine, or the computed operator
certificate for \(\widehat M\).  Sampling probabilities and sampling laws
remain exact mathematical inputs.
\end{corollary}

\begin{formalized}
\leanfit{rectMatMul_entry_abs_sub_computed_le_of_component_sums}
\leanfit{rightSketchCrossGramRectInvFactor_entry_abs_error_le_of_component_sums}
\leanfit{frobNormRect_sigma_rightSketchCrossGramRectInvFactor_le_of_computed_rectOpNorm2Le_of_component_error}
\leanfit{frobNormRect_preconditionRows_sourceSketchResidualTail_sourceSVDTail_moorePenrose_le_sqrt_one_add_eps_plus_component_sq_of_computed_rectOpNorm2Le}
\end{formalized}

\begin{corollary}[Concrete floating-point cross Gram for equation~(9)]
Let \(fp\) be a floating-point model, assume the Higham validity condition
\(\gamma(fp,n)\) for dot products of length \(n\), and define the computed
rectangular cross Gram
\[
  \widehat X
  =
  \operatorname{fl\_matMul}_{fp}\bigl(V_\perp^T,Z\bigr).
\]
For each entry set
\[
  B_{ab}
  =
  \gamma(fp,n)\sum_j |(V_\perp)_{ja}|\,|Z_{jb}|.
\]
Then
\[
  |(V_\perp^T Z)_{ab}-\widehat X_{ab}|
  \le B_{ab}.
\]
Consequently, with
\[
  Y=(V_k^T Z)^{-1},
  \qquad M=(V_\perp^T Z)Y,
\]
suppose a computed inverse factor \(\widehat Y\), rounded product
\(\widehat M\), and radii \(\alpha,\beta,\rho\ge0\) satisfy
\[
  \sum_b B_{ab}|Y_{bc}|\le\alpha ,
\]
\[
  \sum_b |\widehat X_{ab}|\,|Y_{bc}-\widehat Y_{bc}|\le\beta ,
\]
and
\[
  \left|\sum_b \widehat X_{ab}\widehat Y_{bc}
    -\widehat M_{ac}\right|\le\rho
\]
for every \(a,c\).  Then
\[
  |M_{ac}-\widehat M_{ac}|
  \le \alpha+\beta+\rho .
\]
If, in addition,
\[
  \|\widehat Mx\|_2\le\varepsilon\|x\|_2
  \qquad(x\in\mathbb R^r),\qquad \varepsilon\ge0,
\]
then
\[
  \|\Sigma_\perp (V_\perp^T Z)(V_k^T Z)^{-1}\|_F
  \le
  \bigl(\varepsilon+\sqrt{qr}\,(\alpha+\beta+\rho)\bigr)
  \|\Sigma_\perp\|_F .
\]
Under the exact source-tail, right-basis, row-completeness, and Moore--Penrose
sketch certificates used above,
\[
  \|(AZ)C\bigl(T-(TZ)(V_k^T Z)^{-1}V_k^T\bigr)\|_F
  \le
  \sqrt{1+
    \bigl(\varepsilon+\sqrt{qr}\,(\alpha+\beta+\rho)\bigr)^2}\,
  \|\Sigma_\perp\|_F .
\]

This is an implementation-facing theorem for the concrete floating-point
computation of the rectangular cross Gram \(V_\perp^T Z\).  It still keeps the
inverse routine for \((V_k^T Z)^{-1}\), the final rounded product
\(\widehat X\widehat Y\), the operator certificate for \(\widehat M\), the
computed square cross Gram used as inverse input, and computed
SVD/projector/Gram obligations as explicit remaining certificates.  Sampling
probabilities and sampling laws remain exact mathematical inputs.
\end{corollary}

\begin{formalized}
\leanfit{flRightSketchCrossGramRect}
\leanfit{rightSketchCrossGramRectDotBudget}
\leanfit{rightSketchCrossGramRect_flMatMul_entry_abs_error_le}
\leanfit{rightSketchCrossGramRect_flMatMul_component_left_error_le}
\leanfit{rightSketchCrossGramRectInvFactor_entry_abs_error_le_of_flMatMul_crossGram_component_sums}
\leanfit{frobNormRect_sigma_rightSketchCrossGramRectInvFactor_le_of_computed_rectOpNorm2Le_of_flMatMul_crossGram_component_error}
\leanfit{frobNormRect_preconditionRows_sourceSketchResidualTail_sourceSVDTail_moorePenrose_le_sqrt_one_add_eps_plus_flMatMul_crossGram_component_sq_of_computed_rectOpNorm2Le}
\end{formalized}

\begin{corollary}[Concrete floating-point square cross Gram for equation~(9)]
Let \(fp\) be a floating-point model satisfying the Higham validity condition
\(\gamma(fp,n)\) for length-\(n\) dot products, and define
\[
  \widehat G
  =
  \operatorname{fl\_matMul}_{fp}\bigl(V_k^T,Z\bigr).
\]
For each entry set
\[
  D_{ab}
  =
  \gamma(fp,n)\sum_j |(V_k)_{ja}|\,|Z_{jb}|.
\]
Then the exact square cross Gram \(G=V_k^T Z\) satisfies
\[
  |G_{ab}-\widehat G_{ab}|\le D_{ab}.
\]
If, for some \(\omega\ge0\),
\[
  D_{ab}\le\omega
  \qquad\text{for every }a,b,
\]
then
\[
  \|G-\widehat G\|_F
  \le
  \sqrt{r\,r}\,\omega .
\]

This is an implementation-facing inverse-input theorem: it certifies the
computed square cross Gram that an inverse routine would consume when forming
an approximation to \((V_k^T Z)^{-1}\).  It does not by itself prove inverse
stability, a computed inverse-factor budget for \(\widehat Y\), or the final
rounded-product/operator certificates.  Sampling probabilities and sampling
laws remain exact mathematical inputs.
\end{corollary}

\begin{formalized}
\leanfit{flRightSketchCrossGram}
\leanfit{rightSketchCrossGramDotBudget}
\leanfit{rightSketchCrossGram_flMatMul_entry_abs_error_le}
\leanfit{frobNorm_rightSketchCrossGram_sub_flMatMul_le_of_dotBudget_le}
\end{formalized}

\begin{corollary}[Entrywise computed inverse factor for equation~(9)]
Let
\[
  \widehat X
  =
  \operatorname{fl\_matMul}_{fp}\bigl(V_\perp^T,Z\bigr),
  \qquad
  Y=(V_k^T Z)^{-1},
\]
and let \(\widehat Y\) be the computed inverse-factor output supplied by an
implementation.  Suppose that the implementation supplies the entrywise
inverse certificate
\[
  |Y_{bc}-\widehat Y_{bc}|\le \eta
  \qquad\text{for every }b,c,
  \qquad \eta\ge0,
\]
and the computed left-factor row certificate
\[
  \sum_b |\widehat X_{ab}| \le \chi
  \qquad\text{for every }a.
\]
Then the right component sum used by the computed cross-factor theorem is
\[
  \sum_b |\widehat X_{ab}|\,|Y_{bc}-\widehat Y_{bc}|
  \le
  \chi\eta
  \qquad\text{for every }a,c.
\]
Consequently, if the concrete cross-Gram left budget is bounded by
\(\alpha\),
\[
  \sum_b
  B_{ab}\,|Y_{bc}|\le \alpha,
  \qquad
  B_{ab}=\gamma(fp,n)\sum_j |(V_\perp)_{ja}|\,|Z_{jb}|,
\]
and the final rounded product
\[
  \widehat M_{ac}
  \approx
  \sum_b \widehat X_{ab}\widehat Y_{bc}
\]
has entrywise rounding certificate
\[
  \left|
    \sum_b \widehat X_{ab}\widehat Y_{bc}-\widehat M_{ac}
  \right|
  \le \rho,
\]
then
\[
  \left|
    ((V_\perp^T Z)(V_k^T Z)^{-1})_{ac}-\widehat M_{ac}
  \right|
  \le
  \alpha+\chi\eta+\rho .
\]
If, moreover, \(\widehat M\) has operator certificate
\(\|\widehat M x\|_2\le\varepsilon\|x\|_2\), with
\(\varepsilon,\alpha,\chi,\eta,\rho\ge0\), then
\[
  \|\Sigma_\perp(V_\perp^T Z)(V_k^T Z)^{-1}\|_F
  \le
  \bigl(\varepsilon+
    \sqrt{qr}\,(\alpha+\chi\eta+\rho)\bigr)\|\Sigma_\perp\|_F .
\]
Under the same exact source-tail, right-basis, row-completeness, and
Moore--Penrose sketch certificates used above,
\[
  \|(AZ)C(T-(T Z)(V_k^T Z)^{-1}V_k^T)\|_F
  \le
  \sqrt{1+
    \bigl(\varepsilon+
      \sqrt{qr}\,(\alpha+\chi\eta+\rho)\bigr)^2}\,
  \|\Sigma_\perp\|_F .
\]

This is an implementation-facing inverse-factor adapter.  It does not prove
that a concrete inverse routine supplies the entrywise radius \(\eta\); that
routine remains a not-proved obligation.  It does prove that, once such a
non-probability certificate is available, the previous abstract inverse
component radius \(\beta\) may be chosen as \(\chi\eta\).  Sampling
probabilities and sampling laws remain exact mathematical inputs.
\end{corollary}

\begin{formalized}
\leanfit{rightSketchCrossGramRectInvFactor_inverse_component_sum_le_of_entry_abs_error}
\leanfit{rightSketchCrossGramRectInvFactor_flMatMul_inverse_component_sum_le_of_entry_abs_error}
\leanfit{rightSketchCrossGramRectInvFactor_entry_abs_error_le_of_flMatMul_crossGram_inverse_entry_abs_error}
\leanfit{frobNormRect_sigma_rightSketchCrossGramRectInvFactor_le_of_computed_rectOpNorm2Le_of_flMatMul_crossGram_inverse_entry_error}
\leanfit{frobNormRect_preconditionRows_sourceSketchResidualTail_sourceSVDTail_moorePenrose_le_sqrt_one_add_eps_plus_flMatMul_crossGram_inverse_entry_sq_of_computed_rectOpNorm2Le}
\end{formalized}

\begin{corollary}[Concrete rounded product for equation~(9)]
Let
\[
  \widehat X
  =
  \operatorname{fl\_matMul}_{fp}\bigl(V_\perp^T,Z\bigr),
  \qquad
  Y=(V_k^T Z)^{-1},
\]
let \(\widehat Y\) be a computed inverse-factor output, and define the
computed cross factor by the concrete rounded matrix product
\[
  \widehat M
  =
  \operatorname{fl\_matMul}_{fp}(\widehat X,\widehat Y).
\]
For each \(a,c\), set the visible product dot-budget
\[
  R_{ac}
  =
  \gamma(fp,r)
  \sum_b |\widehat X_{ab}|\,|\widehat Y_{bc}|.
\]
Assume the non-probability inverse routine and computed left factor satisfy
\[
  |Y_{bc}-\widehat Y_{bc}|\le \eta
  \quad\text{for every }b,c,
  \qquad
  \sum_b |\widehat X_{ab}|\le \chi
  \quad\text{for every }a,
\]
and that the concrete cross-Gram left budget satisfies
\[
  \sum_b
  B_{ab}\,|Y_{bc}|\le \alpha,
  \qquad
  B_{ab}=\gamma(fp,n)\sum_j |(V_\perp)_{ja}|\,|Z_{jb}|.
\]
If
\[
  R_{ac}\le \rho
  \qquad\text{for every }a,c,
\]
then the exact analysis factor
\[
  M=(V_\perp^T Z)(V_k^T Z)^{-1}
\]
and the implemented rounded product \(\widehat M\) satisfy
\[
  |M_{ac}-\widehat M_{ac}|
  \le
  \alpha+\chi\eta+\rho
  \qquad\text{for every }a,c.
\]
If, moreover, the computed product has operator certificate
\[
  \|\widehat M x\|_2\le\varepsilon\|x\|_2
  \qquad\text{for every }x,
\]
with \(\varepsilon,\alpha,\chi,\eta,\rho\ge0\), then
\[
  \|\Sigma_\perp(V_\perp^T Z)(V_k^T Z)^{-1}\|_F
  \le
  \bigl(\varepsilon+
    \sqrt{qr}\,(\alpha+\chi\eta+\rho)\bigr)\|\Sigma_\perp\|_F .
\]
Under the same exact source-tail, right-basis, row-completeness, and
Moore--Penrose sketch certificates used above,
\[
  \|(AZ)C(T-(T Z)(V_k^T Z)^{-1}V_k^T)\|_F
  \le
  \sqrt{1+
    \bigl(\varepsilon+
      \sqrt{qr}\,(\alpha+\chi\eta+\rho)\bigr)^2}\,
  \|\Sigma_\perp\|_F .
\]

This is an implementation-facing rounded-product theorem.  It instantiates
the final product radius \(\rho\) by the concrete repository
\(\operatorname{fl\_matMul}\) dot-budget \(R_{ac}\).  It does not prove that a
concrete inverse routine supplies the entrywise radius \(\eta\), and it does
not prove the computed operator certificate for \(\widehat M\).  Sampling
probabilities and sampling laws remain exact mathematical inputs.
\end{corollary}

\begin{formalized}
\leanfit{flRightSketchCrossGramRectInvFactorProduct}
\leanfit{rightSketchCrossGramRectInvFactorProductDotBudget}
\leanfit{rightSketchCrossGramRectInvFactorProduct_flMatMul_entry_abs_error_le}
\leanfit{rightSketchCrossGramRectInvFactor_entry_abs_error_le_of_flMatMul_crossGram_inverse_entry_abs_error_flMatMul_product}
\leanfit{frobNormRect_sigma_rightSketchCrossGramRectInvFactor_le_of_computed_rectOpNorm2Le_of_flMatMul_crossGram_inverse_entry_flMatMul_product}
\leanfit{frobNormRect_preconditionRows_sourceSketchResidualTail_sourceSVDTail_moorePenrose_le_sqrt_one_add_eps_plus_flMatMul_crossGram_inverse_entry_flMatMul_product_sq_of_computed_rectOpNorm2Le}
\end{formalized}

\begin{corollary}[Computed-product Frobenius certificate for equation~(9)]
With the notation and hypotheses of the preceding corollary, suppose the
computed product itself has the visible Frobenius certificate
\[
  \|\widehat M\|_F\le\varepsilon,
  \qquad
  \widehat M=
  \operatorname{fl\_matMul}_{fp}(\widehat X,\widehat Y).
\]
Then \(\varepsilon\ge0\), and the ordinary rectangular operator certificate
\[
  \|\widehat Mx\|_2\le\varepsilon\|x\|_2
  \qquad\text{for every }x
\]
holds.  Consequently, if the same concrete cross-Gram left budget, inverse
entrywise certificate, computed-left row budget, and final product budget
hold,
\[
  \sum_b B_{ab}|Y_{bc}|\le\alpha,\qquad
  |Y_{bc}-\widehat Y_{bc}|\le\eta,\qquad
  \sum_b|\widehat X_{ab}|\le\chi,\qquad
  R_{ac}\le\rho,
\]
with \(\alpha,\chi,\eta,\rho\ge0\), then
\[
  \|\Sigma_\perp(V_\perp^T Z)(V_k^T Z)^{-1}\|_F
  \le
  \bigl(\varepsilon+
    \sqrt{qr}\,(\alpha+\chi\eta+\rho)\bigr)\|\Sigma_\perp\|_F .
\]
Under the same exact source-tail, right-basis, row-completeness, and
Moore--Penrose sketch certificates used above,
\[
  \|(AZ)C(T-(T Z)(V_k^T Z)^{-1}V_k^T)\|_F
  \le
  \sqrt{1+
    \bigl(\varepsilon+
      \sqrt{qr}\,(\alpha+\chi\eta+\rho)\bigr)^2}\,
  \|\Sigma_\perp\|_F .
\]

This is an implementation-facing computed-operator handoff.  It replaces the
abstract operator certificate for the computed product by the concrete,
checkable Frobenius certificate \(\|\widehat M\|_F\le\varepsilon\).  It does
not prove that an inverse routine supplies \(\eta\), nor that a randomized
argument or a particular computation routine supplies the displayed Frobenius
certificate.  Sampling probabilities and sampling laws remain exact
mathematical inputs.
\end{corollary}

\begin{formalized}
\leanfit{rectOpNorm2Le_flRightSketchCrossGramRectInvFactorProduct_of_frobNormRect_le}
\leanfit{frobNormRect_sigma_rightSketchCrossGramRectInvFactor_le_of_frobNormRect_flMatMul_crossGram_inverse_entry_flMatMul_product}
\leanfit{frobNormRect_preconditionRows_sourceSketchResidualTail_sourceSVDTail_moorePenrose_le_sqrt_one_add_frobNormRect_flMatMul_crossGram_inverse_entry_flMatMul_product_sq}
\end{formalized}

\begin{corollary}[Product absolute-sum Frobenius certificate for equation~(9)]
Let
\[
  \widehat M=
  \operatorname{fl\_matMul}_{fp}(\widehat X,\widehat Y),
  \qquad
  R_{ac}=\gamma(fp,r)
  \sum_b |\widehat X_{ab}|\,|\widehat Y_{bc}|.
\]
Assume that the computed inputs to the final product satisfy the visible
absolute-product budget
\[
  \sum_b |\widehat X_{ab}|\,|\widehat Y_{bc}|
  \le \kappa
  \qquad\text{for every }a,c,
\]
and that the final rounded-product budget satisfies
\[
  R_{ac}\le\rho
  \qquad\text{for every }a,c.
\]
Then
\[
  |\widehat M_{ac}|\le \kappa+\rho
  \qquad\text{for every }a,c.
\]
If \(\kappa,\rho\ge0\), then
\[
  \|\widehat M\|_F
  \le
  \sqrt{qr}\,(\kappa+\rho).
\]
Consequently, any displayed radius \(\varepsilon\) satisfying
\[
  \sqrt{qr}\,(\kappa+\rho)\le\varepsilon
\]
supplies the ordinary rectangular operator certificate
\[
  \|\widehat Mx\|_2\le\varepsilon\|x\|_2
  \qquad\text{for every }x.
\]
With \(\widehat X=\operatorname{fl\_matMul}_{fp}(V_\perp^T,Z)\), if the same
concrete cross-Gram left budget, inverse entrywise certificate, computed-left
row budget, and final product budget from the preceding corollaries hold,
\[
  \sum_b B_{ab}|Y_{bc}|\le\alpha,\qquad
  |Y_{bc}-\widehat Y_{bc}|\le\eta,\qquad
  \sum_b|\widehat X_{ab}|\le\chi,\qquad
  R_{ac}\le\rho,
\]
with \(\alpha,\chi,\eta,\rho\ge0\), then
\[
  \|\Sigma_\perp(V_\perp^T Z)(V_k^T Z)^{-1}\|_F
  \le
  \bigl(\varepsilon+
    \sqrt{qr}\,(\alpha+\chi\eta+\rho)\bigr)\|\Sigma_\perp\|_F .
\]
Under the same exact source-tail, right-basis, row-completeness, and
Moore--Penrose sketch certificates used above,
\[
  \|(AZ)C(T-(T Z)(V_k^T Z)^{-1}V_k^T)\|_F
  \le
  \sqrt{1+
    \bigl(\varepsilon+
      \sqrt{qr}\,(\alpha+\chi\eta+\rho)\bigr)^2}\,
  \|\Sigma_\perp\|_F .
\]

This is an implementation-facing source for the product Frobenius/operator
certificate.  It charges only computed non-probability arithmetic: the final
rounded product and the magnitudes of its computed inputs.  It still does not
prove that an inverse routine supplies \(\eta\), nor does it derive
randomness-generated certificates or computed SVD/projector/Gram certificates.
Sampling probabilities and sampling laws remain exact mathematical inputs.
\end{corollary}

\begin{formalized}
\leanfit{rightSketchCrossGramRectInvFactorProduct_entry_abs_le_of_product_sum_budget}
\leanfit{frobNormRect_flRightSketchCrossGramRectInvFactorProduct_le_sqrt_mul_product_sum_budget}
\leanfit{rectOpNorm2Le_flRightSketchCrossGramRectInvFactorProduct_of_product_sum_budget}
\leanfit{frobNormRect_sigma_rightSketchCrossGramRectInvFactor_le_of_product_sum_budget_flMatMul_crossGram_inverse_entry_flMatMul_product}
\leanfit{frobNormRect_preconditionRows_sourceSketchResidualTail_sourceSVDTail_moorePenrose_le_sqrt_one_add_product_sum_budget_flMatMul_crossGram_inverse_entry_flMatMul_product_sq}
\end{formalized}

\begin{corollary}[Perturbed-inverse certificate for equation~(9)]
Let \(A_0\in\mathbb R^{r\times r}\) have nonzero determinant, and let
\(\widehat Y\) and \(\Delta A\) satisfy the perturbed-inverse equation
\[
  (A_0+\Delta A)\widehat Y=I.
\]
Assume the componentwise perturbation budget
\[
  |\Delta A_{ij}|\le\varepsilon_{\rm inv}|(A_0)_{ij}|
  \qquad\text{for every }i,j,
\]
and the visible inverse-sensitivity budget
\[
  \varepsilon_{\rm inv}
  \sum_{k_1}|(A_0^{-1})_{ik_1}|
    \sum_{k_2}|(A_0)_{k_1k_2}|\,|\widehat Y_{k_2j}|
  \le \eta
  \qquad\text{for every }i,j.
\]
Then
\[
  |(A_0^{-1})_{ij}-\widehat Y_{ij}|\le\eta
  \qquad\text{for every }i,j.
\]
Specializing \(A_0=V_k^T Z\) gives the inverse-entry certificate
\[
  |((V_k^T Z)^{-1})_{bc}-\widehat Y_{bc}|\le\eta
  \qquad\text{for every }b,c
\]
consumed by the computed cross-factor results above.

Consequently, with
\(\widehat X=\operatorname{fl\_matMul}_{fp}(V_\perp^T,Z)\), if the product
absolute-sum certificate of the preceding corollary supplies
\(\|\widehat Mx\|_2\le\varepsilon_M\|x\|_2\), and if the same concrete
cross-Gram left budget, computed-left row budget, and final product budget
hold,
\[
  \sum_b B_{ab}|Y_{bc}|\le\alpha,\qquad
  \sum_b|\widehat X_{ab}|\le\chi,\qquad
  R_{ac}\le\rho,
\]
with \(\alpha,\chi,\eta,\rho\ge0\), then
\[
  \|\Sigma_\perp(V_\perp^T Z)(V_k^T Z)^{-1}\|_F
  \le
  \bigl(\varepsilon_M+
    \sqrt{qr}\,(\alpha+\chi\eta+\rho)\bigr)\|\Sigma_\perp\|_F .
\]
Under the same exact source-tail, right-basis, row-completeness, and
Moore--Penrose sketch certificates used above,
\[
  \|(AZ)C(T-(T Z)(V_k^T Z)^{-1}V_k^T)\|_F
  \le
  \sqrt{1+
    \bigl(\varepsilon_M+
      \sqrt{qr}\,(\alpha+\chi\eta+\rho)\bigr)^2}\,
  \|\Sigma_\perp\|_F .
\]

This is an implementation-facing inverse-radius source theorem.  It uses the
local Higham~\S13.1 forward-error theorem for computed inverses to replace the
raw assumption \(|Y-\widehat Y|\le\eta\) by a perturbed-inverse contract.  It
does not prove that a particular inversion loop supplies \(\Delta A\) and
\(\widehat Y\).  Sampling probabilities and sampling laws remain exact
mathematical inputs.
\end{corollary}

\begin{formalized}
\leanfit{nonsingInv_entry_abs_sub_computed_inverse_le_of_perturbed_inverse_component_budget}
\leanfit{rightSketchCrossGram_inverse_entry_abs_error_le_of_perturbed_inverse_component_budget}
\leanfit{frobNormRect_sigma_rightSketchCrossGramRectInvFactor_le_of_perturbed_inverse_product_sum_budget}
\leanfit{frobNormRect_preconditionRows_sourceSketchResidualTail_sourceSVDTail_moorePenrose_le_sqrt_one_add_perturbed_inverse_product_sum_budget_sq}
\end{formalized}

\begin{corollary}[Method-A LU inverse certificate for equation~(9)]
Let \(A_0=V_k^T Z\in\mathbb R^{r\times r}\) have nonzero determinant, and
let \(\widehat L,\widehat U\in\mathbb R^{r\times r}\) have nonzero diagonals.
Assume that \(\widehat L,\widehat U\) satisfy the componentwise
LU backward-error certificate
\[
  \operatorname{LUBackwardError}_r(A_0,\widehat L,\widehat U,\gamma(fp,r)).
\]
Define the Method-A computed inverse
\[
  \widehat Y
  =\operatorname{methodAComputedInverse}_{fp,r}(\widehat L,\widehat U),
\]
meaning that the \(c\)-th column is computed by solving
\[
  \widehat L\,\widehat y_c=e_c,\qquad
  \widehat U\,\widehat x_c=\widehat y_c
\]
with the repository floating-point forward- and back-substitution routines,
and setting \(\widehat Y_{\cdot c}=\widehat x_c\).  Put
\[
  c_r=3\gamma(fp,r)+\gamma(fp,r)^2.
\]
Assume the visible Method-A inverse-sensitivity budget
\[
  c_r
  \sum_{k_1}|(A_0^{-1})_{bk_1}|
    \sum_{k_2}
      \left(\sum_{\ell}|\widehat L_{k_1\ell}|\,|\widehat U_{\ell k_2}|\right)
      |\widehat Y_{k_2c}|
  \le \eta
  \qquad\text{for every }b,c .
\]
Then
\[
  |(A_0^{-1})_{bc}-\widehat Y_{bc}|\le\eta
  \qquad\text{for every }b,c .
\]
Equivalently, specializing \(A_0=V_k^T Z\),
\[
  |((V_k^T Z)^{-1})_{bc}
     -\operatorname{methodAComputedInverse}_{fp,r}
        (\widehat L,\widehat U)_{bc}|
  \le \eta .
\]

Consequently, with
\(\widehat X=\operatorname{fl\_matMul}_{fp}(V_\perp^T,Z)\) and
\(\widehat Y=\operatorname{methodAComputedInverse}_{fp,r}
(\widehat L,\widehat U)\), if the same product absolute-sum certificate
supplies \(\|\widehat Mx\|_2\le\varepsilon_M\|x\|_2\), and if
\[
  \sum_b B_{ab}|Y_{bc}|\le\alpha,\qquad
  \sum_b|\widehat X_{ab}|\le\chi,\qquad
  R_{ac}\le\rho,
\]
with \(\alpha,\chi,\eta,\rho\ge0\), then
\[
  \|\Sigma_\perp(V_\perp^T Z)(V_k^T Z)^{-1}\|_F
  \le
  \bigl(\varepsilon_M+
    \sqrt{qr}\,(\alpha+\chi\eta+\rho)\bigr)\|\Sigma_\perp\|_F .
\]
Under the same exact source-tail, right-basis, row-completeness, and
Moore--Penrose sketch certificates used above,
\[
  \|(AZ)C(T-(T Z)(V_k^T Z)^{-1}V_k^T)\|_F
  \le
  \sqrt{1+
    \bigl(\varepsilon_M+
      \sqrt{qr}\,(\alpha+\chi\eta+\rho)\bigr)^2}\,
  \|\Sigma_\perp\|_F .
\]

This is an implementation-facing inverse-solve source theorem for Method A.
It charges the computed non-probability inverse columns through LU,
forward-substitution, and back-substitution certificates.  It is still
conditional on the displayed LU backward-error certificate for the square
cross Gram, and it does not prove computed SVD/projector/Gram or
randomness-derived certificates.  Sampling probabilities and sampling laws
remain exact mathematical inputs.
\end{corollary}

\begin{formalized}
\leanfit{methodAComputedInverse}
\leanfit{methodA_column_backward_error_computed_inverse}
\leanfit{methodA_computed_inverse_entry_abs_sub_nonsingInv_le_of_lu_budget}
\leanfit{rightSketchCrossGram_inverse_entry_abs_error_le_of_methodA_lu_budget}
\leanfit{frobNormRect_sigma_rightSketchCrossGramRectInvFactor_le_of_methodA_inverse_product_sum_budget}
\leanfit{frobNormRect_preconditionRows_sourceSketchResidualTail_sourceSVDTail_moorePenrose_le_sqrt_one_add_methodA_inverse_product_sum_budget_sq}
\end{formalized}

\begin{corollary}[Computed-input Method-A LU certificate for equation~(9)]
Let \(A_0=V_k^T Z\in\mathbb R^{r\times r}\) have nonzero determinant, and
let
\[
  \widehat A_0
  =\operatorname{flRightSketchCrossGram}_{fp}(V_k,Z)
  =\operatorname{fl\_matMul}_{fp}(V_k^T,Z)
\]
be the rounded square cross Gram actually passed to the LU factorization.
Assume that the computed factors \(\widehat L,\widehat U\) satisfy
\[
  \operatorname{LUBackwardError}_r
    (\widehat A_0,\widehat L,\widehat U,\varepsilon_{\rm LU}),
\]
and that the rounded-input error is controlled in the same componentwise
\(|\widehat L|\,|\widehat U|\) weights:
\[
  \operatorname{rightSketchCrossGramDotBudget}_{fp}(V_k,Z)_{bc}
  \le
  \mu\sum_{\ell}|\widehat L_{b\ell}|\,|\widehat U_{\ell c}|
  \qquad\text{for every }b,c .
\]
Then the exact analysis matrix inherits the LU certificate
\[
  \operatorname{LUBackwardError}_r
    (A_0,\widehat L,\widehat U,\varepsilon_{\rm LU}+\mu).
\]
Assume \(\varepsilon_{\rm LU}\ge0\), \(\mu\ge0\), the triangular diagonals of
\(\widehat L,\widehat U\) are nonzero, and put
\[
  c_r^{\rm in}
  =(\varepsilon_{\rm LU}+\mu)+2\gamma(fp,r)+\gamma(fp,r)^2 .
\]
With
\[
  \widehat Y
  =\operatorname{methodAComputedInverse}_{fp,r}(\widehat L,\widehat U),
\]
assume the visible inverse-sensitivity budget
\[
  c_r^{\rm in}
  \sum_{k_1}|(A_0^{-1})_{bk_1}|
    \sum_{k_2}
      \left(\sum_{\ell}|\widehat L_{k_1\ell}|\,|\widehat U_{\ell k_2}|\right)
      |\widehat Y_{k_2c}|
  \le \eta
  \qquad\text{for every }b,c .
\]
Then
\[
  |(A_0^{-1})_{bc}-\widehat Y_{bc}|\le\eta
  \qquad\text{for every }b,c .
\]
Consequently, with
\(\widehat X=\operatorname{fl\_matMul}_{fp}(V_\perp^T,Z)\), if the same
product absolute-sum certificate supplies
\(\|\widehat Mx\|_2\le\varepsilon_M\|x\|_2\), and if
\[
  \sum_b B_{ab}|Y_{bc}|\le\alpha,\qquad
  \sum_b|\widehat X_{ab}|\le\chi,\qquad
  R_{ac}\le\rho,
\]
with \(\alpha,\chi,\eta,\rho\ge0\), then
\[
  \|\Sigma_\perp(V_\perp^T Z)(V_k^T Z)^{-1}\|_F
  \le
  \bigl(\varepsilon_M+
    \sqrt{qr}\,(\alpha+\chi\eta+\rho)\bigr)\|\Sigma_\perp\|_F ,
\]
and under the same exact source-tail, right-basis, row-completeness, and
Moore--Penrose sketch certificates used above,
\[
  \|(AZ)C(T-(T Z)(V_k^T Z)^{-1}V_k^T)\|_F
  \le
  \sqrt{1+
    \bigl(\varepsilon_M+
      \sqrt{qr}\,(\alpha+\chi\eta+\rho)\bigr)^2}\,
  \|\Sigma_\perp\|_F .
\]

This is an implementation-facing non-probability input-transfer theorem.  It
does not charge probability construction or the sampling law; those remain
exact mathematical inputs.  It does charge the rounded square cross Gram
actually consumed by the LU factorization.  It is still conditional on a
concrete routine proving the displayed `LUBackwardError' certificate for
\(\widehat A_0=\operatorname{flRightSketchCrossGram}_{fp}(V_k,Z)\).
\end{corollary}

\begin{formalized}
\leanfit{LUBackwardError.of_input_abs_error_le_absLUProduct}
\leanfit{lu_solve_backward_error_factor_gamma}
\leanfit{methodA_computed_inverse_entry_abs_sub_nonsingInv_le_of_lu_factor_budget}
\leanfit{rightSketchCrossGram_LUBackwardError_of_flRightSketchCrossGram_input_budget}
\leanfit{rightSketchCrossGram_inverse_entry_abs_error_le_of_methodA_fl_lu_input_budget}
\leanfit{frobNormRect_sigma_rightSketchCrossGramRectInvFactor_le_of_methodA_lu_factor_product_sum_budget}
\leanfit{frobNormRect_preconditionRows_sourceSketchResidualTail_sourceSVDTail_moorePenrose_le_sqrt_one_add_methodA_lu_factor_product_sum_budget_sq}
\end{formalized}

\begin{corollary}[Doolittle-generated computed-input Method-A LU certificate for equation~(9)]
Let \(A_0=V_k^T Z\in\mathbb R^{r\times r}\) have nonzero determinant, and
let
\[
  \widehat A_0
  =\operatorname{flRightSketchCrossGram}_{fp}(V_k,Z)
  =\operatorname{fl\_matMul}_{fp}(V_k^T,Z)
\]
be the rounded square cross Gram consumed by the dense LU factorization.
Assume that \(\widehat L,\widehat U\) satisfy the Doolittle recurrence
certificate
\[
  \operatorname{DoolittleLU}_r
    (\widehat A_0,\widehat L,\widehat U,fp).
\]
Then the repository recurrence-to-backward-error theorem gives
\[
  \operatorname{LUBackwardError}_r
    (\widehat A_0,\widehat L,\widehat U,\gamma(fp,r)).
\]
Assume moreover that the rounded square-cross-Gram input error is controlled
in the same componentwise \(|\widehat L|\,|\widehat U|\) weights:
\[
  \operatorname{rightSketchCrossGramDotBudget}_{fp}(V_k,Z)_{bc}
  \le
  \mu\sum_{\ell}|\widehat L_{b\ell}|\,|\widehat U_{\ell c}|
  \qquad\text{for every }b,c .
\]
Then the exact analysis matrix inherits
\[
  \operatorname{LUBackwardError}_r
    (A_0,\widehat L,\widehat U,\gamma(fp,r)+\mu).
\]
Assume \(\mu\ge0\), the triangular diagonals of \(\widehat L,\widehat U\)
are nonzero, and put
\[
  c_r^{\rm Doolittle}
  =
  (\gamma(fp,r)+\mu)+2\gamma(fp,r)+\gamma(fp,r)^2 .
\]
With
\[
  \widehat Y
  =\operatorname{methodAComputedInverse}_{fp,r}(\widehat L,\widehat U),
\]
assume the visible inverse-sensitivity budget
\[
  c_r^{\rm Doolittle}
  \sum_{k_1}|(A_0^{-1})_{bk_1}|
    \sum_{k_2}
      \left(\sum_{\ell}|\widehat L_{k_1\ell}|\,|\widehat U_{\ell k_2}|\right)
      |\widehat Y_{k_2c}|
  \le \eta
  \qquad\text{for every }b,c .
\]
Then
\[
  |(A_0^{-1})_{bc}-\widehat Y_{bc}|\le\eta
  \qquad\text{for every }b,c .
\]
Consequently, with
\(\widehat X=\operatorname{fl\_matMul}_{fp}(V_\perp^T,Z)\), if the same
product absolute-sum certificate supplies
\(\|\widehat Mx\|_2\le\varepsilon_M\|x\|_2\), and if
\[
  \sum_b B_{ab}|Y_{bc}|\le\alpha,\qquad
  \sum_b|\widehat X_{ab}|\le\chi,\qquad
  R_{ac}\le\rho,
\]
with \(\alpha,\chi,\eta,\rho\ge0\), then
\[
  \|\Sigma_\perp(V_\perp^T Z)(V_k^T Z)^{-1}\|_F
  \le
  \bigl(\varepsilon_M+
    \sqrt{qr}\,(\alpha+\chi\eta+\rho)\bigr)\|\Sigma_\perp\|_F ,
\]
and under the same exact source-tail, right-basis, row-completeness, and
Moore--Penrose sketch certificates used above,
\[
  \|(AZ)C(T-(T Z)(V_k^T Z)^{-1}V_k^T)\|_F
  \le
  \sqrt{1+
    \bigl(\varepsilon_M+
      \sqrt{qr}\,(\alpha+\chi\eta+\rho)\bigr)^2}\,
  \|\Sigma_\perp\|_F .
\]

This is an implementation-facing, certificate-level dense-LU result.  It
charges the rounded square cross Gram, the Doolittle recurrence
factorization, and the Method-A triangular solves.  It does not charge
probability construction or the sampling law; those remain exact mathematical
inputs.  It still requires a concrete executable dense-LU implementation to
produce the displayed \(\operatorname{DoolittleLU}\) witness.
\end{corollary}

\begin{formalized}
\leanfit{DoolittleLU.to_LUBackwardError}
\leanfit{doolittle_backward_error}
\leanfit{rightSketchCrossGram_LUBackwardError_of_DoolittleLU_flRightSketchCrossGram}
\leanfit{rightSketchCrossGram_inverse_entry_abs_error_le_of_methodA_doolittle_fl_input_budget}
\leanfit{frobNormRect_sigma_rightSketchCrossGramRectInvFactor_le_of_methodA_doolittle_fl_input_product_sum_budget}
\leanfit{frobNormRect_preconditionRows_sourceSketchResidualTail_sourceSVDTail_moorePenrose_le_sqrt_one_add_methodA_doolittle_fl_input_product_sum_budget_sq}
\end{formalized}

\begin{corollary}[Dense-loop Doolittle residual-compression certificate for equation~(9)]
Let
\[
  \widehat A_0
  =\operatorname{flRightSketchCrossGram}_{fp}(V_k,Z)
  =\operatorname{fl\_matMul}_{fp}(V_k^T,Z)
\]
be the rounded square cross Gram used by Method~A.  For \(k\le j\) define the
literal rounded Doolittle upper-entry fold
\[
  \widehat u_{kj}
  =
  \operatorname{flDoolittleUEntry}_{fp}
    (\widehat A_0,\widehat L,\widehat U)_{kj},
\]
and for \(i>k\) define the rounded lower-entry numerator and division
\[
  \widehat n_{ik}
  =
  \operatorname{flDoolittleLNumerator}_{fp}
    (\widehat A_0,\widehat L,\widehat U)_{ik},
  \qquad
  \widehat \ell_{ik}
  =
  \operatorname{flDoolittleLEntry}_{fp}
    (\widehat A_0,\widehat L,\widehat U)_{ik}.
\]
Assume the stored factors satisfy
\[
  \widehat U_{kj}=\widehat u_{kj}\quad(k\le j),\qquad
  \widehat L_{ik}=\widehat \ell_{ik}\quad(i>k),
\]
with \(\widehat L\) unit lower triangular and \(\widehat U\) upper triangular.
Assume also the two visible residual-compression budgets
\[
  \left|
    \widehat A_{0,kj}
    -\sum_{s<k}\widehat L_{ks}\widehat U_{sj}
    -\widehat U_{kj}
  \right|
  \le
  \gamma(fp,r)|\widehat U_{kj}|
  \qquad(k\le j),
\]
and
\[
  \left|
    \widehat A_{0,ik}
    -\sum_{s<k}\widehat L_{is}\widehat U_{sk}
    -\widehat L_{ik}\widehat U_{kk}
  \right|
  \le
  \gamma(fp,r)|\widehat L_{ik}\widehat U_{kk}|
  \qquad(i>k).
\]
Then these data form a
\[
  \operatorname{DoolittleDenseLoopCertificate}_r
    (\widehat A_0,\widehat L,\widehat U,fp),
\]
and the residual-compression theorem constructs the compact recurrence
certificate
\[
  \operatorname{DoolittleLU}_r
    (\widehat A_0,\widehat L,\widehat U,fp).
\]
If \(\operatorname{gammaValid}(fp,r)\), then
\[
  \operatorname{LUBackwardError}_r
    (\widehat A_0,\widehat L,\widehat U,\gamma(fp,r)).
\]
Consequently, under the same determinant, triangular-diagonal, rounded-input
budget
\[
  \operatorname{rightSketchCrossGramDotBudget}_{fp}(V_k,Z)_{bc}
  \le
  \mu\sum_{\ell}|\widehat L_{b\ell}|\,|\widehat U_{\ell c}|,
\]
Method-A inverse-sensitivity budget with
\[
  c_r^{\rm loop}
  =
  (\gamma(fp,r)+\mu)+2\gamma(fp,r)+\gamma(fp,r)^2,
\]
and product absolute-sum hypotheses from the previous Doolittle corollary, the
computed inverse
\[
  \widehat Y
  =
  \operatorname{methodAComputedInverse}_{fp,r}(\widehat L,\widehat U)
\]
satisfies
\[
  |(V_k^T Z)^{-1}_{bc}-\widehat Y_{bc}|\le\eta
  \qquad\text{for every }b,c,
\]
and the equation-(9) cross-term and projected-tail bounds are
\[
  \|\Sigma_\perp(V_\perp^T Z)(V_k^T Z)^{-1}\|_F
  \le
  \bigl(\varepsilon_M+
    \sqrt{qr}\,(\alpha+\chi\eta+\rho)\bigr)\|\Sigma_\perp\|_F ,
\]
\[
  \|(AZ)C(T-(T Z)(V_k^T Z)^{-1}V_k^T)\|_F
  \le
  \sqrt{1+
    \bigl(\varepsilon_M+
      \sqrt{qr}\,(\alpha+\chi\eta+\rho)\bigr)^2}\,
  \|\Sigma_\perp\|_F .
\]

This is an implementation-facing non-probability theorem surface.  Sampling
probabilities and sampling laws remain exact mathematical inputs.  The theorem
does not yet prove the residual-compression budgets from a concrete pivoting or
no-cancellation regime; it states exactly the certificate that such an
implementation proof must instantiate.
\end{corollary}

\begin{formalized}
\leanfit{flDoolittleUEntry}
\leanfit{flDoolittleLNumerator}
\leanfit{flDoolittleLEntry}
\leanfit{DoolittleDenseLoopCertificate}
\leanfit{DoolittleDenseLoopCertificate.to_DoolittleLU}
\leanfit{DoolittleDenseLoopCertificate.to_LUBackwardError}
\leanfit{rightSketchCrossGram_LUBackwardError_of_DoolittleDenseLoopCertificate_flRightSketchCrossGram}
\leanfit{rightSketchCrossGram_inverse_entry_abs_error_le_of_methodA_doolittleDenseLoop_fl_input_budget}
\leanfit{frobNormRect_sigma_rightSketchCrossGramRectInvFactor_le_of_methodA_doolittleDenseLoop_fl_input_product_sum_budget}
\leanfit{frobNormRect_preconditionRows_sourceSketchResidualTail_sourceSVDTail_moorePenrose_le_sqrt_one_add_methodA_doolittleDenseLoop_fl_input_product_sum_budget_sq}
\end{formalized}

\begin{corollary}[Absolute-budget dense-loop Doolittle handoff for equation~(9)]
Let \(\widehat A_0=\operatorname{flRightSketchCrossGram}_{fp}(V_k,Z)\) be the
rounded square cross Gram used by Method~A, and let
\(\widehat L,\widehat U\in\mathbb R^{r\times r}\) be stored dense-Doolittle
factors with \(\widehat L\) unit lower triangular and \(\widehat U\) upper
triangular.  Suppose the stored entries agree with the literal rounded
Doolittle folds:
\[
  \widehat U_{kj}
  =
  \operatorname{flDoolittleUEntry}_{fp}
    (\widehat A_0,\widehat L,\widehat U)_{kj}
  \quad(k\le j),
  \qquad
  \widehat L_{ik}
  =
  \operatorname{flDoolittleLEntry}_{fp}
    (\widehat A_0,\widehat L,\widehat U)_{ik}
  \quad(i>k).
\]
Let \(B^U,B^L\in\mathbb R^{r\times r}\) be absolute residual budgets satisfying
\[
  \left|
    \widehat A_{0,kj}
    -\sum_{s<k}\widehat L_{ks}\widehat U_{sj}
    -\widehat U_{kj}
  \right|
  \le B^U_{kj}
  \qquad(k\le j),
\]
\[
  \left|
    \widehat A_{0,ik}
    -\sum_{s<k}\widehat L_{is}\widehat U_{sk}
    -\widehat L_{ik}\widehat U_{kk}
  \right|
  \le B^L_{ik}
  \qquad(i>k).
\]
Assume further that these absolute budgets are dominated by the relative
compression radii
\[
  B^U_{kj}\le \gamma(fp,r)|\widehat U_{kj}|
  \qquad(k\le j),
  \qquad
  B^L_{ik}\le \gamma(fp,r)|\widehat L_{ik}\widehat U_{kk}|
  \qquad(i>k).
\]
Then the data form an implementation-facing certificate
\[
  \operatorname{DoolittleDenseLoopAbsBudgetCertificate}_r
    (\widehat A_0,\widehat L,\widehat U,fp,B^U,B^L),
\]
which yields
\[
  \operatorname{DoolittleDenseLoopCertificate}_r
    (\widehat A_0,\widehat L,\widehat U,fp),
  \qquad
  \operatorname{DoolittleLU}_r
    (\widehat A_0,\widehat L,\widehat U,fp).
\]
If \(\operatorname{gammaValid}(fp,r)\), then
\[
  \operatorname{LUBackwardError}_r
    (\widehat A_0,\widehat L,\widehat U,\gamma(fp,r)).
\]

Consequently, under the determinant and nonzero triangular-diagonal hypotheses
from the previous corollary, the rounded-input budget
\[
  \operatorname{rightSketchCrossGramDotBudget}_{fp}(V_k,Z)_{bc}
  \le
  \mu\sum_{\ell}|\widehat L_{b\ell}|\,|\widehat U_{\ell c}|,
\]
the Method-A inverse-sensitivity budget with
\[
  c_r^{\rm abs}
  =
  (\gamma(fp,r)+\mu)+2\gamma(fp,r)+\gamma(fp,r)^2,
\]
and the same product absolute-sum, final-product, and left-cross-Gram budgets,
the computed inverse
\[
  \widehat Y
  =
  \operatorname{methodAComputedInverse}_{fp,r}(\widehat L,\widehat U)
\]
satisfies
\[
  |(V_k^T Z)^{-1}_{bc}-\widehat Y_{bc}|\le\eta
  \qquad\text{for all }b,c,
\]
and the equation-(9) cross-term and projected-tail bounds are
\[
  \|\Sigma_\perp(V_\perp^T Z)(V_k^T Z)^{-1}\|_F
  \le
  \bigl(\varepsilon_M+
    \sqrt{qr}\,(\alpha+\chi\eta+\rho)\bigr)\|\Sigma_\perp\|_F ,
\]
\[
  \|(AZ)C(T-(T Z)(V_k^T Z)^{-1}V_k^T)\|_F
  \le
  \sqrt{1+
    \bigl(\varepsilon_M+
      \sqrt{qr}\,(\alpha+\chi\eta+\rho)\bigr)^2}\,
  \|\Sigma_\perp\|_F .
\]

This is a fully computed-object theorem surface for the non-probability
dense-Doolittle arithmetic once the absolute residual budgets and dominance
inequalities are supplied.  Sampling probabilities and sampling laws remain
exact mathematical inputs.  The theorem does not yet prove \(B^U,B^L\) or the
dominance inequalities from a concrete pivoting or no-cancellation regime; that
routine instantiation remains in the not-proved ledger.
\end{corollary}

\begin{formalized}
\leanfit{DoolittleDenseLoopAbsBudgetCertificate}
\leanfit{DoolittleDenseLoopAbsBudgetCertificate.to_denseLoopCertificate}
\leanfit{DoolittleDenseLoopAbsBudgetCertificate.to_DoolittleLU}
\leanfit{DoolittleDenseLoopAbsBudgetCertificate.to_LUBackwardError}
\leanfit{rightSketchCrossGram_LUBackwardError_of_DoolittleDenseLoopAbsBudgetCertificate_flRightSketchCrossGram}
\leanfit{rightSketchCrossGram_inverse_entry_abs_error_le_of_methodA_doolittleDenseLoopAbsBudget_fl_input_budget}
\leanfit{frobNormRect_sigma_rightSketchCrossGramRectInvFactor_le_of_methodA_doolittleDenseLoopAbsBudget_fl_input_product_sum_budget}
\leanfit{frobNormRect_preconditionRows_sourceSketchResidualTail_sourceSVDTail_moorePenrose_le_sqrt_one_add_methodA_doolittleDenseLoopAbsBudget_fl_input_product_sum_budget_sq}
\end{formalized}

\begin{corollary}[Rounded-product Doolittle fold budgets for equation~(9)]
Let \(m\in\mathbb N\), let \(s\in\mathbb R\), and let
\(t_0,\ldots,t_{m-1}\in\mathbb R\) be the rounded terms supplied to a
sequential floating-point subtraction fold.  If
\(\operatorname{gammaValid}(fp,m)\), then
\[
  \left|
    \left(s-\sum_{a=0}^{m-1}t_a\right)
    -\operatorname{foldl}_{a<m}\bigl(x\mapsto
        \operatorname{fl\_sub}_{fp}(x,t_a)\bigr)(s)
  \right|
  \le
  \gamma(fp,m)\left(|s|+\sum_{a=0}^{m-1}|t_a|\right).
\]
Specializing this to the dense-Doolittle upper-entry fold, with
\[
  \widehat p^U_{ksj}
  =
  \operatorname{fl\_mul}_{fp}(\widehat L_{ks},\widehat U_{sj}),
\]
gives
\[
  \left|
    \widehat A_{0,kj}
    -\sum_{s<k}\widehat p^U_{ksj}
    -\operatorname{flDoolittleUEntry}_{fp}
      (\widehat A_0,\widehat L,\widehat U)_{kj}
  \right|
  \le
  \gamma(fp,k)
  \left(|\widehat A_{0,kj}|+\sum_{s<k}|\widehat p^U_{ksj}|\right).
\]
Specializing it to the lower-entry numerator fold, with
\[
  \widehat p^L_{isk}
  =
  \operatorname{fl\_mul}_{fp}(\widehat L_{is},\widehat U_{sk}),
\]
gives
\[
  \left|
    \widehat A_{0,ik}
    -\sum_{s<k}\widehat p^L_{isk}
    -\operatorname{flDoolittleLNumerator}_{fp}
      (\widehat A_0,\widehat L,\widehat U)_{ik}
  \right|
  \le
  \gamma(fp,k)
  \left(|\widehat A_{0,ik}|+\sum_{s<k}|\widehat p^L_{isk}|\right).
\]

This is an implementation-facing non-probability theorem for the subtraction
accumulation that the dense Doolittle loop actually performs.  It does not yet
turn rounded products \(\widehat p^U,\widehat p^L\) into exact products
\(\widehat L\widehat U\), does not charge the lower rounded division and
subsequent multiplication by \(\widehat U_{kk}\), and does not prove the
dominance inequalities required by the absolute-budget certificate above.
Sampling probabilities and sampling laws remain exact mathematical inputs.
\end{corollary}

\begin{formalized}
\leanfit{fl_sub_sum_error_init_abs_residual_le}
\leanfit{flDoolittleUEntry_rounded_residual_abs_le}
\leanfit{flDoolittleLNumerator_rounded_residual_abs_le}
\end{formalized}

\begin{corollary}[Exact-product Doolittle fold budgets for equation~(9)]
For the dense-Doolittle upper-entry fold, define the exact and rounded
non-probability products
\[
  p^U_{ksj}=\widehat L_{ks}\widehat U_{sj},
  \qquad
  \widehat p^U_{ksj}
  =
  \operatorname{fl\_mul}_{fp}(\widehat L_{ks},\widehat U_{sj}).
\]
The primitive multiplication model gives
\[
  \left|
    \operatorname{fl\_mul}_{fp}(x,y)-xy
  \right|
  \le u_{fp}|xy|.
\]
Consequently, if \(\operatorname{gammaValid}(fp,k)\), then
\[
  \begin{aligned}
  &\left|
    \widehat A_{0,kj}
    -\sum_{s<k}p^U_{ksj}
    -\operatorname{flDoolittleUEntry}_{fp}
      (\widehat A_0,\widehat L,\widehat U)_{kj}
  \right|\\
  &\qquad\le
  \gamma(fp,k)
  \left(|\widehat A_{0,kj}|+\sum_{s<k}|\widehat p^U_{ksj}|\right)\\
  &\qquad\quad
  +u_{fp}\sum_{s<k}|p^U_{ksj}|.
  \end{aligned}
\]
For the lower-entry numerator fold, with
\[
  p^L_{isk}=\widehat L_{is}\widehat U_{sk},
  \qquad
  \widehat p^L_{isk}
  =
  \operatorname{fl\_mul}_{fp}(\widehat L_{is},\widehat U_{sk}),
\]
the same hypotheses give
\[
  \begin{aligned}
  &\left|
    \widehat A_{0,ik}
    -\sum_{s<k}p^L_{isk}
    -\operatorname{flDoolittleLNumerator}_{fp}
      (\widehat A_0,\widehat L,\widehat U)_{ik}
  \right|\\
  &\qquad\le
  \gamma(fp,k)
  \left(|\widehat A_{0,ik}|+\sum_{s<k}|\widehat p^L_{isk}|\right)\\
  &\qquad\quad
  +u_{fp}\sum_{s<k}|p^L_{isk}|.
  \end{aligned}
\]

This closes the product-transfer part of the dense-Doolittle absolute-budget
source: the residuals are now measured against the exact products
\(\widehat L\widehat U\), not only against the rounded products supplied to
the subtraction fold.  It still does not package the \(s<k\) sums into the
masked \(n\)-term sums used by \(\operatorname{DoolittleDenseLoopAbsBudgetCertificate}\),
does not charge the lower rounded division and multiplication by
\(\widehat U_{kk}\), and does not prove the dominance inequalities required by
that certificate.  Sampling probabilities and sampling laws remain exact
mathematical inputs.
\end{corollary}

\begin{formalized}
\leanfit{fl_mul_abs_sub_mul_le}
\leanfit{flDoolittleUEntry_exact_product_residual_abs_le}
\leanfit{flDoolittleLNumerator_exact_product_residual_abs_le}
\end{formalized}

\begin{corollary}[Masked-prefix Doolittle fold budgets for equation~(9)]
For every \(k<n\) and every table \(f_s\),
\[
  \sum_{s=0}^{n-1}{\bf 1}_{\{s<k\}}f_s
  =
  \sum_{s<k} f_s .
\]
Thus the exact-product Doolittle fold budgets above can be stated in the
masked \(n\)-term recurrence shape used by the dense-loop certificate.
Let
\[
  p^U_{ksj}=\widehat L_{ks}\widehat U_{sj},
  \qquad
  \widehat p^U_{ksj}
  =
  \operatorname{fl\_mul}_{fp}(\widehat L_{ks},\widehat U_{sj}).
\]
If \(\operatorname{gammaValid}(fp,k)\), then
\[
  \begin{aligned}
  &\left|
    \widehat A_{0,kj}
    -\sum_{s=0}^{n-1}{\bf 1}_{\{s<k\}}p^U_{ksj}
    -\operatorname{flDoolittleUEntry}_{fp}
      (\widehat A_0,\widehat L,\widehat U)_{kj}
  \right|\\
  &\qquad\le
  \gamma(fp,k)
  \left(|\widehat A_{0,kj}|+\sum_{s<k}|\widehat p^U_{ksj}|\right)\\
  &\qquad\quad
  +u_{fp}\sum_{s<k}|p^U_{ksj}|.
  \end{aligned}
\]
Likewise, with
\[
  p^L_{isk}=\widehat L_{is}\widehat U_{sk},
  \qquad
  \widehat p^L_{isk}
  =
  \operatorname{fl\_mul}_{fp}(\widehat L_{is},\widehat U_{sk}),
\]
the lower-entry numerator satisfies
\[
  \begin{aligned}
  &\left|
    \widehat A_{0,ik}
    -\sum_{s=0}^{n-1}{\bf 1}_{\{s<k\}}p^L_{isk}
    -\operatorname{flDoolittleLNumerator}_{fp}
      (\widehat A_0,\widehat L,\widehat U)_{ik}
  \right|\\
  &\qquad\le
  \gamma(fp,k)
  \left(|\widehat A_{0,ik}|+\sum_{s<k}|\widehat p^L_{isk}|\right)\\
  &\qquad\quad
  +u_{fp}\sum_{s<k}|p^L_{isk}|.
  \end{aligned}
\]

This is an implementation-facing non-probability theorem for the recurrence
shape: the computed multiplication and subtraction-fold errors are charged,
and the prefix sums are reindexed into the masked \(n\)-term form.  It still
does not charge the lower rounded division and multiplication by the computed
pivot \(\widehat U_{kk}\), and it does not prove the dominance inequalities
needed to instantiate the full dense-Doolittle absolute-budget certificate.
Sampling probabilities and sampling laws remain exact mathematical inputs.
\end{corollary}

\begin{formalized}
\leanfit{finMaskedPrefixSum_eq_finSum}
\leanfit{flDoolittleUEntry_masked_exact_product_residual_abs_le}
\leanfit{flDoolittleLNumerator_masked_exact_product_residual_abs_le}
\end{formalized}

\begin{corollary}[Lower rounded-division Doolittle budget for equation~(9)]
Let
\[
  \widehat n_{ik}
  =
  \operatorname{flDoolittleLNumerator}_{fp}
    (\widehat A_0,\widehat L,\widehat U)_{ik},
  \qquad
  \widehat \ell_{ik}
  =
  \operatorname{flDoolittleLEntry}_{fp}
    (\widehat A_0,\widehat L,\widehat U)_{ik}.
\]
If the computed pivot is nonzero, \(\widehat U_{kk}\ne0\), then the
primitive rounded-division model gives
\[
  |\widehat n_{ik}-\widehat \ell_{ik}\widehat U_{kk}|
  \le u_{fp}|\widehat n_{ik}|.
\]
Consequently, if also \(\operatorname{gammaValid}(fp,k)\) and the stored lower
entry is the literal rounded Doolittle update,
\[
  \widehat L_{ik}
  =
  \operatorname{flDoolittleLEntry}_{fp}
    (\widehat A_0,\widehat L,\widehat U)_{ik},
\]
then, with
\[
  p^L_{isk}=\widehat L_{is}\widehat U_{sk},
  \qquad
  \widehat p^L_{isk}
  =
  \operatorname{fl\_mul}_{fp}(\widehat L_{is},\widehat U_{sk}),
\]
the lower-entry residual after multiplication by the computed pivot satisfies
\[
  \begin{aligned}
  &\left|
    \widehat A_{0,ik}
    -\sum_{s=0}^{n-1}{\bf 1}_{\{s<k\}}p^L_{isk}
    -\widehat L_{ik}\widehat U_{kk}
  \right|\\
  &\qquad\le
  \left[
    \gamma(fp,k)
    \left(|\widehat A_{0,ik}|+\sum_{s<k}|\widehat p^L_{isk}|\right)
    +u_{fp}\sum_{s<k}|p^L_{isk}|
  \right]\\
  &\qquad\quad
  +u_{fp}|\widehat n_{ik}|.
  \end{aligned}
\]

This is an implementation-facing non-probability theorem for the lower-entry
division and computed-pivot multiplication in dense Doolittle.  It closes the
lower rounded-division/pivot source for the absolute-budget certificate.  It
still does not prove the dominance inequality
\[
  B^L_{ik}\le \gamma(fp,n)|\widehat L_{ik}\widehat U_{kk}|
\]
or any concrete pivot/no-cancellation regime.  Sampling probabilities and
sampling laws remain exact mathematical inputs.
\end{corollary}

\begin{formalized}
\leanfit{flDoolittleLEntry_mul_pivot_sub_numerator_abs_le}
\leanfit{flDoolittleLEntry_masked_exact_product_residual_abs_le}
\end{formalized}

\begin{corollary}[Literal Doolittle source-budget certificate for equation~(9)]
Define the concrete upper and lower dense-Doolittle absolute budgets
\[
  \begin{aligned}
  B^U_{kj}
  &=
  \gamma(fp,k)
  \left(|\widehat A_{0,kj}|+\sum_{s<k}|\widehat p^U_{ksj}|\right)
  +u_{fp}\sum_{s<k}|p^U_{ksj}|,\\
  B^L_{ik}
  &=
  \gamma(fp,k)
  \left(|\widehat A_{0,ik}|+\sum_{s<k}|\widehat p^L_{isk}|\right)
  +u_{fp}\sum_{s<k}|p^L_{isk}|
  +u_{fp}|\widehat n_{ik}|.
  \end{aligned}
\]
Here
\[
  p^U_{ksj}=\widehat L_{ks}\widehat U_{sj},\quad
  p^L_{isk}=\widehat L_{is}\widehat U_{sk},\quad
  \widehat p=\operatorname{fl\_mul}_{fp}(p\text{'s factors}),
\]
and
\(\widehat n_{ik}=
\operatorname{flDoolittleLNumerator}_{fp}
(\widehat A_0,\widehat L,\widehat U)_{ik}\).
Assume the stored factors have the required triangular shape, the upper and
lower stored entries agree with the literal rounded Doolittle updates, all
computed pivots are nonzero, and \(\operatorname{gammaValid}(fp,n)\).  If the
visible dominance inequalities
\[
  B^U_{kj}\le \gamma(fp,n)|\widehat U_{kj}|
  \quad(k\le j),
  \qquad
  B^L_{ik}\le \gamma(fp,n)|\widehat L_{ik}\widehat U_{kk}|
  \quad(k<i)
\]
hold, then the concrete budgets \(B^U,B^L\) instantiate
\[
  \operatorname{DoolittleDenseLoopAbsBudgetCertificate}
    (\widehat A_0,\widehat L,\widehat U,fp,B^U,B^L).
\]
Consequently all downstream dense-loop Method-A handoffs from the absolute
budget certificate apply with these explicit non-probability arithmetic
budgets.

This closes the source-budget assembly for literal dense Doolittle.  The only
remaining dense-Doolittle arithmetic bottleneck is proving the displayed
dominance inequalities from a concrete pivot/no-cancellation regime or
recording that routine-specific regime as an implementation obligation.
Sampling probabilities and sampling laws remain exact mathematical inputs.
\end{corollary}

\begin{formalized}
\leanfit{doolittleUAbsBudget}
\leanfit{doolittleLAbsBudget}
\leanfit{DoolittleDenseLoopAbsBudgetCertificate.of_literal_doolittle_source_budgets}
\end{formalized}

\begin{corollary}[Componentwise Doolittle dominance for equation~(9)]
For the literal dense-Doolittle budgets above define the upper component
quantities
\[
  W^U_{kj}=|\widehat A_{0,kj}|+
    \sum_{s<k}\left|
      \operatorname{fl\_mul}_{fp}(\widehat L_{ks},\widehat U_{sj})
    \right|,
  \qquad
  E^U_{kj}=\sum_{s<k}
    |\widehat L_{ks}\widehat U_{sj}|,
\]
and the lower component quantities
\[
  W^L_{ik}=|\widehat A_{0,ik}|+
    \sum_{s<k}\left|
      \operatorname{fl\_mul}_{fp}(\widehat L_{is},\widehat U_{sk})
    \right|,
  \quad
  E^L_{ik}=\sum_{s<k}
    |\widehat L_{is}\widehat U_{sk}|,
  \quad
  N^L_{ik}=|\widehat n_{ik}|.
\]
Assume \(\operatorname{gammaValid}(fp,n)\).  If \(k\le j\) and
\[
  W^U_{kj}\le |\widehat U_{kj}|,
  \qquad
  E^U_{kj}\le |\widehat U_{kj}|,
\]
then
\[
  B^U_{kj}\le
  \bigl(\gamma(fp,k)+u_{fp}\bigr)|\widehat U_{kj}|
  \le \gamma(fp,n)|\widehat U_{kj}|.
\]
If \(k<i\) and
\[
  W^L_{ik}\le |\widehat L_{ik}\widehat U_{kk}|,
  \qquad
  E^L_{ik}\le |\widehat L_{ik}\widehat U_{kk}|,
  \qquad
  N^L_{ik}\le |\widehat L_{ik}\widehat U_{kk}|,
\]
then
\[
  B^L_{ik}\le
  \bigl(\gamma(fp,k)+u_{fp}+u_{fp}\bigr)
    |\widehat L_{ik}\widehat U_{kk}|
  \le \gamma(fp,n)|\widehat L_{ik}\widehat U_{kk}|.
\]
Consequently, under the same triangular-shape, literal-update, nonzero-pivot,
and \(\operatorname{gammaValid}(fp,n)\) hypotheses as the previous corollary,
these componentwise no-cancellation inequalities instantiate
\[
  \operatorname{DoolittleDenseLoopAbsBudgetCertificate}
    (\widehat A_0,\widehat L,\widehat U,fp,B^U,B^L).
\]

This closes the dense-Doolittle dominance route from visible componentwise
no-cancellation hypotheses.  It still does not derive those component
inequalities from a concrete pivoting or no-cancellation implementation
regime; that routine-specific source remains an implementation obligation.
Sampling probabilities and sampling laws remain exact mathematical inputs.
\end{corollary}

\begin{formalized}
\leanfit{doolittleUWorkAbs}
\leanfit{doolittleUProductAbs}
\leanfit{doolittleLWorkAbs}
\leanfit{doolittleLProductAbs}
\leanfit{doolittleLNumeratorAbs}
\leanfit{doolittleUAbsBudget_le_compression_of_component_dominance}
\leanfit{doolittleLAbsBudget_le_compression_of_component_dominance}
\leanfit{DoolittleDenseLoopAbsBudgetCertificate.of_literal_doolittle_component_dominance}
\end{formalized}

\begin{corollary}[Exact-product Doolittle margins for equation~(9)]
For every real pair \(x,y\), the primitive multiplication model gives the
rounded-product growth estimate
\[
  |\operatorname{fl\_mul}_{fp}(x,y)|
  \le (1+u_{fp})\,|xy|.
\]
Hence the componentwise dominance hypotheses in the previous corollary may be
checked from exact-product no-cancellation margins.  For an upper entry define
\[
  E^U_{kj}=\sum_{s<k}|\widehat L_{ks}\widehat U_{sj}|.
\]
If \(k\le j\) and
\[
  |\widehat A_{0,kj}|+(1+u_{fp})E^U_{kj}
  \le |\widehat U_{kj}|,
\]
then
\[
  W^U_{kj}\le |\widehat U_{kj}|,
  \qquad
  E^U_{kj}\le |\widehat U_{kj}|.
\]
For a lower entry define
\[
  E^L_{ik}=\sum_{s<k}|\widehat L_{is}\widehat U_{sk}|.
\]
If \(k<i\) and
\[
  |\widehat A_{0,ik}|+(1+u_{fp})E^L_{ik}
  \le |\widehat L_{ik}\widehat U_{kk}|,
\]
then
\[
  W^L_{ik}\le |\widehat L_{ik}\widehat U_{kk}|,
  \qquad
  E^L_{ik}\le |\widehat L_{ik}\widehat U_{kk}|.
\]
Consequently, under the same triangular-shape, literal-update, nonzero-pivot,
and \(\operatorname{gammaValid}(fp,n)\) hypotheses as above, if the exact-product
margins hold for all upper and lower entries and the rounded lower numerator
also satisfies
\[
  N^L_{ik}\le |\widehat L_{ik}\widehat U_{kk}|
  \qquad (k<i),
\]
then the literal dense-Doolittle absolute-budget certificate is instantiated:
\[
  \operatorname{DoolittleDenseLoopAbsBudgetCertificate}
    (\widehat A_0,\widehat L,\widehat U,fp,B^U,B^L).
\]

This closes the product-growth handoff from exact-product no-cancellation
margins to the componentwise dense-Doolittle dominance route.  It is an
implementation-facing non-probability arithmetic theorem: products computed
inside the fold pay the explicit \((1+u_{fp})\) growth factor, while sampling
probabilities and sampling laws remain exact mathematical inputs.  Proving the
displayed exact-product margins and the lower rounded-numerator dominance from
a concrete pivoting/no-cancellation implementation regime remains a separate
routine obligation.
\end{corollary}

\begin{formalized}
\leanfit{fl_mul_abs_le_one_add_u_mul_abs_mul}
\leanfit{doolittleUWorkAbs_le_of_exact_product_margin}
\leanfit{doolittleUProductAbs_le_of_exact_product_margin}
\leanfit{doolittleLWorkAbs_le_of_exact_product_margin}
\leanfit{doolittleLProductAbs_le_of_exact_product_margin}
\leanfit{DoolittleDenseLoopAbsBudgetCertificate.of_literal_doolittle_exact_product_margins}
\end{formalized}

\begin{corollary}[Exact-product lower-numerator margins for equation~(9)]
Let \(k<i\), assume \(\operatorname{gammaValid}(fp,k)\), and define the
exact-product lower prefix magnitude
\[
  E^L_{ik}=\sum_{s<k}|\widehat L_{is}\widehat U_{sk}|,
\]
together with the rounded lower numerator magnitude
\[
  N^L_{ik}
  =
  \left|
    \operatorname{flDoolittleLNumerator}_{fp}
      (\widehat A_0,\widehat L,\widehat U,i,k)
  \right|.
\]
If the stronger exact-product numerator margin
\[
  \bigl(|\widehat A_{0,ik}|+E^L_{ik}\bigr)
  +
  \left[
    \gamma(fp,k)\bigl(|\widehat A_{0,ik}|+(1+u_{fp})E^L_{ik}\bigr)
    +u_{fp}E^L_{ik}
  \right]
  \le |\widehat L_{ik}\widehat U_{kk}|,
\]
then
\[
  N^L_{ik}\le |\widehat L_{ik}\widehat U_{kk}|.
\]
Consequently, under triangular-shape hypotheses, literal
\(\operatorname{flDoolittleUEntry}\) and \(\operatorname{flDoolittleLEntry}\)
update equalities, nonzero computed pivots \(\widehat U_{kk}\), and
\(\operatorname{gammaValid}(fp,n)\), if the upper exact-product margins
\[
  |\widehat A_{0,kj}|+(1+u_{fp})E^U_{kj}
  \le |\widehat U_{kj}| \qquad (k\le j),
\]
the lower exact-product margins
\[
  |\widehat A_{0,ik}|+(1+u_{fp})E^L_{ik}
  \le |\widehat L_{ik}\widehat U_{kk}| \qquad (k<i),
\]
and the lower exact-product numerator margins above hold for all required
entries, then
\[
  \operatorname{DoolittleDenseLoopAbsBudgetCertificate}
    (\widehat A_0,\widehat L,\widehat U,fp,B^U,B^L).
\]

Thus the lower rounded-numerator dominance hypothesis in the previous
corollary can be discharged by an explicit source-shaped margin.  The theorem
charges all non-probability arithmetic computed by the literal dense-Doolittle
fold: rounded products, the lower subtraction fold, and the lower division
surface inherited from the source-budget certificate.  Sampling probabilities
and sampling laws remain exact mathematical inputs.  Proving these margins from
a concrete pivoting/no-cancellation implementation regime remains an open
routine obligation.
\end{corollary}

\begin{formalized}
\leanfit{doolittleLNumeratorAbs_le_of_exact_product_numerator_margin}
\leanfit{DoolittleDenseLoopAbsBudgetCertificate.of_literal_doolittle_exact_product_numerator_margins}
\end{formalized}

\begin{corollary}[Exact-target Doolittle gaps for equation~(9)]
For the literal dense-Doolittle routine define the exact pre-rounded upper and
lower targets
\[
  T^U_{kj}
  =
  \widehat A_{0,kj}
  -
  \sum_{s<k}\widehat L_{ks}\widehat U_{sj},
  \qquad
  T^L_{ik}
  =
  \widehat A_{0,ik}
  -
  \sum_{s<k}\widehat L_{is}\widehat U_{sk}.
\]
Let
\[
  E^U_{kj}=\sum_{s<k}|\widehat L_{ks}\widehat U_{sj}|,
  \qquad
  E^L_{ik}=\sum_{s<k}|\widehat L_{is}\widehat U_{sk}|.
\]
The rounded-product growth estimates for the products actually used by the
subtraction folds are
\[
  \sum_{s<k}
    |\operatorname{fl\_mul}_{fp}(\widehat L_{ks},\widehat U_{sj})|
  \le (1+u_{fp})E^U_{kj},
  \qquad
  \sum_{s<k}
    |\operatorname{fl\_mul}_{fp}(\widehat L_{is},\widehat U_{sk})|
  \le (1+u_{fp})E^L_{ik}.
\]
Consequently the literal upper fold satisfies the exact-target residual bound
\[
  |T^U_{kj}-\widehat U_{kj}|
  \le
  R^U_{kj}
  :=
  \gamma(fp,k)\bigl(|\widehat A_{0,kj}|+(1+u_{fp})E^U_{kj}\bigr)
  +u_{fp}E^U_{kj},
\]
assuming \(\widehat U_{kj}=
\operatorname{flDoolittleUEntry}_{fp}(\widehat A_0,\widehat L,\widehat U,k,j)\).
For the lower numerator,
\[
  |T^L_{ik}-
    \operatorname{flDoolittleLNumerator}_{fp}
      (\widehat A_0,\widehat L,\widehat U,i,k)|
  \le
  R^N_{ik}
  :=
  \gamma(fp,k)\bigl(|\widehat A_{0,ik}|+(1+u_{fp})E^L_{ik}\bigr)
  +u_{fp}E^L_{ik}.
\]
If the computed pivot is nonzero and
\[
  \widehat L_{ik}=
  \operatorname{flDoolittleLEntry}_{fp}
    (\widehat A_0,\widehat L,\widehat U,i,k),
\]
then the lower entry multiplied back by the computed pivot satisfies
\[
  |T^L_{ik}-\widehat L_{ik}\widehat U_{kk}|
  \le
  R^L_{ik}
  :=
  R^N_{ik}
  +u_{fp}
   \left|
    \operatorname{flDoolittleLNumerator}_{fp}
      (\widehat A_0,\widehat L,\widehat U,i,k)
   \right|.
\]
Thus exact-target no-cancellation gaps imply the source-shaped margin
hypotheses used above.  If \(k\le j\) and
\[
  |\widehat A_{0,kj}|+(1+u_{fp})E^U_{kj}+R^U_{kj}
  \le |T^U_{kj}|,
\]
then
\[
  |\widehat A_{0,kj}|+(1+u_{fp})E^U_{kj}
  \le |\widehat U_{kj}|.
\]
If \(k<i\), \(\widehat U_{kk}\ne0\), and
\[
  |\widehat A_{0,ik}|+(1+u_{fp})E^L_{ik}+R^L_{ik}
  \le |T^L_{ik}|,
\]
then
\[
  |\widehat A_{0,ik}|+(1+u_{fp})E^L_{ik}
  \le |\widehat L_{ik}\widehat U_{kk}|.
\]
Finally, if the stronger exact-target lower gap
\[
  \bigl(|\widehat A_{0,ik}|+E^L_{ik}+R^N_{ik}\bigr)+R^L_{ik}
  \le |T^L_{ik}|
\]
holds, then
\[
  |\widehat A_{0,ik}|+E^L_{ik}+R^N_{ik}
  \le |\widehat L_{ik}\widehat U_{kk}|.
\]
Therefore, under triangular-shape hypotheses, literal
\(\operatorname{flDoolittleUEntry}\) and \(\operatorname{flDoolittleLEntry}\)
update equalities, nonzero computed pivots, and
\(\operatorname{gammaValid}(fp,n)\), if the displayed exact-target gaps hold
for all required entries, then
\[
  \operatorname{DoolittleDenseLoopAbsBudgetCertificate}
    (\widehat A_0,\widehat L,\widehat U,fp,B^U,B^L).
\]

This is an implementation-facing non-probability arithmetic result: it charges
rounded products, subtraction folds, lower rounded division, and multiplication
by the computed pivot.  It strictly reduces the remaining dense-Doolittle
obligation to proving the exact-target gaps from a concrete
pivoting/off-diagonal-control/no-cancellation invariant.  Sampling
probabilities and sampling laws remain exact mathematical inputs.
\end{corollary}

\begin{formalized}
\leanfit{doolittleUExactTarget}
\leanfit{doolittleLExactTarget}
\leanfit{doolittleUExactTargetResidualBudget}
\leanfit{doolittleLExactTargetNumeratorResidualBudget}
\leanfit{doolittleLExactTargetEntryResidualBudget}
\leanfit{doolittleURoundedProductAbsSum_le_one_add_u_productAbs}
\leanfit{doolittleLRoundedProductAbsSum_le_one_add_u_productAbs}
\leanfit{doolittleUExactTarget_residual_abs_le_of_literal}
\leanfit{doolittleLExactTarget_numerator_residual_abs_le_of_literal}
\leanfit{doolittleLExactTarget_entry_residual_abs_le_of_literal}
\leanfit{doolittleUExactProductMargin_of_exactTarget_gap}
\leanfit{doolittleLExactProductMargin_of_exactTarget_gap}
\leanfit{doolittleLExactProductNumeratorMargin_of_exactTarget_gap}
\leanfit{DoolittleDenseLoopAbsBudgetCertificate.of_literal_doolittle_exact_target_gaps}
\end{formalized}

\begin{corollary}[Triangle obstruction for exact-target Doolittle gaps]
The exact-target transfer above is a valid conditional theorem, but its
source-facing exact-target gaps are not an ordinary no-cancellation condition
in nondegenerate floating point.  The exact targets always obey the triangle
upper bounds
\[
  |T^U_{kj}|
  \le |\widehat A_{0,kj}|+E^U_{kj},
  \qquad
  |T^L_{ik}|
  \le |\widehat A_{0,ik}|+E^L_{ik}.
\]
Consequently, if the upper exact-target gap from the previous corollary holds,
then
\[
  u_{fp}E^U_{kj}+R^U_{kj}\le 0.
\]
If the lower entry exact-target gap holds, then
\[
  u_{fp}E^L_{ik}+R^L_{ik}\le 0.
\]
If the stronger lower numerator exact-target gap holds, then
\[
  R^N_{ik}+R^L_{ik}\le 0.
\]
Equivalently, the corresponding exact-target gap is contradictory whenever the
displayed excess term is positive.

Thus LR.1bp remains useful as a conditional transfer from already-supplied
exact-target gaps to stored-entry margins, but LR.1bq rules out using ordinary
pivoting/no-cancellation or off-diagonal-control estimates to supply those
gaps when the FP residual excess is positive.  The live dense-Doolittle D5
route must instead prove stored-entry/component dominance directly, revise the
compression surface, or move to sharper product/operator certificates.
Sampling probabilities and sampling laws remain exact mathematical inputs.
\end{corollary}

\begin{formalized}
\leanfit{doolittleUExactTarget_abs_le_source_plus_productAbs}
\leanfit{doolittleLExactTarget_abs_le_source_plus_productAbs}
\leanfit{doolittleUExactTarget_gap_excess_nonpos}
\leanfit{doolittleUExactTarget_gap_false_of_positive_excess}
\leanfit{doolittleLExactTarget_gap_excess_nonpos}
\leanfit{doolittleLExactTarget_gap_false_of_positive_excess}
\leanfit{doolittleLExactTarget_numerator_gap_excess_nonpos}
\leanfit{doolittleLExactTarget_numerator_gap_false_of_positive_excess}
\end{formalized}

\begin{corollary}[Norm-generic head-tail surface for equation~(9)]
Let \(\xi_{m,n}\) be a rectangular norm-like functional on
\(\mathbb R^{m\times n}\) satisfying
\[
  \xi_{m,n}(X)\ge0,\qquad
  \xi_{m,n}(X-Y)\le \xi_{m,n}(X)+\xi_{m,n}(Y)
  \quad\text{for all }X,Y\in\mathbb R^{m\times n}.
\]
Let \(Z\in\mathbb R^{n\times r}\), let \(P\in\mathbb R^{m\times m}\),
and suppose that \(P\) factors through the exact column sketch:
\[
  P_{ij}=\sum_{a=1}^r (AZ)_{ia}C_{aj}.
\]
Assume also that \(P\) exactly reproduces the sketch columns,
\[
  P(AZ)=AZ,
\]
and that \(A\) has an exact head-tail split
\[
  A=H+T,\qquad
  H=(AZ)W
\]
for some coefficient table \(W\).  If the two visible norm bounds
\[
  \xi_{m,n}(T)\le \tau,\qquad
  \xi_{m,n}(PT)\le \kappa
\]
hold with \(\tau,\kappa\ge0\), then
\[
  \operatorname{rank}(PA)\le r,\qquad
  \xi_{m,n}(A-PA)\le \tau+\kappa .
\]
Consequently, if \(A_k\) is a certified best rank-\(k\) approximation with
respect to \(\xi_{m,n}\) and
\[
  \tau+\kappa\le \rho\,\xi_{m,n}(A-A_k),
\]
then
\[
  \operatorname{rank}(A_k)\le k,\qquad
  \operatorname{rank}(PA)\le r,\qquad
  \xi_{m,n}(A-PA)\le \rho\,\xi_{m,n}(A-A_k).
\]

This is an exact-object, norm-abstract theorem surface.  It proves the
equation~(9) rank/residual algebra for any supplied norm-like functional and
head-tail certificate.  It does not instantiate a concrete unitarily invariant
norm, prove right/left orthogonal invariance, prove the singular-value
comparison for the cross term, or prove Eckart--Young for that norm.  Sampling
probabilities and laws remain exact mathematical inputs.
\end{corollary}

\begin{formalized}
\leanfit{RectNormLike}
\leanfit{lowRankResidualNorm}
\leanfit{IsBestRankApproxNorm}
\leanfit{Equation9ResidualNormCertificate}
\leanfit{Equation9HeadTailSketchNormCertificate}
\leanfit{Equation9HeadTailSketchNormCertificate.to_residualNormCertificate}
\leanfit{equation9RankResidualNormSurface}
\leanfit{equation9RelativeResidualNormSurface}
\leanfit{equation9HeadTailSketchNormRankResidualSurface}
\leanfit{equation9HeadTailSketchNormRelativeResidualSurface}
\end{formalized}

\begin{corollary}[Frobenius instance of the norm-generic equation~(9) surface]
Define
\[
  \xi^F_{m,n}(X)=\|X\|_F,\qquad X\in\mathbb R^{m\times n}.
\]
Then \(\xi^F_{m,n}\) is a rectangular norm-like functional:
\[
  \xi^F_{m,n}(X)\ge0,\qquad
  \xi^F_{m,n}(X-Y)\le \xi^F_{m,n}(X)+\xi^F_{m,n}(Y).
\]
Moreover, for every orthogonal \(U\in\mathbb R^{m\times m}\) and
orthogonal \(V\in\mathbb R^{n\times n}\),
\[
  \xi^F_{m,n}(UX)=\xi^F_{m,n}(X),\qquad
  \xi^F_{m,n}(XV)=\xi^F_{m,n}(X).
\]
Let \(Z\in\mathbb R^{n\times r}\), let \(P\in\mathbb R^{m\times m}\), and
suppose that \(P\) factors through the exact column sketch and reproduces it:
\[
  P_{ij}=\sum_{a=1}^r (AZ)_{ia}C_{aj},\qquad P(AZ)=AZ.
\]
Assume that
\[
  A=H+T,\qquad H=(AZ)W,
\]
and that the visible Frobenius bounds
\[
  \|T\|_F\le\tau,\qquad \|PT\|_F\le\kappa,\qquad \tau,\kappa\ge0
\]
hold.  Then
\[
  \operatorname{rank}(PA)\le r,\qquad
  \|A-PA\|_F\le \tau+\kappa .
\]
Consequently, if \(A_k\) is a certified best rank-\(k\) Frobenius
approximation and
\[
  \tau+\kappa\le \rho\,\|A-A_k\|_F,
\]
then
\[
  \operatorname{rank}(A_k)\le k,\qquad
  \operatorname{rank}(PA)\le r,\qquad
  \|A-PA\|_F\le \rho\,\|A-A_k\|_F .
\]

This is a concrete Frobenius instantiation of the norm-generic exact-object
surface above.  It closes the `RectNormLike' bridge for the repository
Frobenius norm and exposes the already-formalized left/right orthogonal
invariance facts.  It still does not prove the all-unitarily-invariant norm
theorem, singular-value comparison for the equation~(9) cross term, or
Eckart--Young.  Sampling probabilities and laws remain exact mathematical
inputs.
\end{corollary}

\begin{formalized}
\leanfit{frobRectNormLike}
\leanfit{frobRectNormLike_norm}
\leanfit{lowRankResidualNorm_frobRectNormLike}
\leanfit{IsBestRankApproxFrob.to_norm_frobRectNormLike}
\leanfit{Equation9ResidualCertificate.to_norm_frobRectNormLike}
\leanfit{Equation9HeadTailSketchCertificate.to_norm_frobRectNormLike}
\leanfit{frobRectNormLike_orthogonal_left}
\leanfit{frobRectNormLike_orthogonal_right}
\leanfit{equation9HeadTailSketchFrobNormRankResidualSurface}
\leanfit{equation9HeadTailSketchFrobNormRelativeResidualSurface}
\end{formalized}

\begin{corollary}[Unitary-invariant norm API for equation~(9)]
Let \(\xi_{m,n}\) be a rectangular norm-like functional equipped with exact
left and right orthogonal invariance:
\[
  \xi_{m,n}(UX)=\xi_{m,n}(X),\qquad
  \xi_{m,n}(XV)=\xi_{m,n}(X)
\]
for every orthogonal \(U\in\mathbb R^{m\times m}\), every orthogonal
\(V\in\mathbb R^{n\times n}\), and every \(X\in\mathbb R^{m\times n}\).
The repository packages this data as
\(\operatorname{UnitaryInvariantRectNormLike}_{m,n}\), extending
\(\operatorname{RectNormLike}_{m,n}\).  The Frobenius norm instance satisfies
\[
  \operatorname{frobUnitaryInvariantRectNormLike}_{m,n}
    \longmapsto \operatorname{frobRectNormLike}_{m,n},
  \qquad
  \xi_{m,n}(X)=\|X\|_F .
\]

For any supplied \(\operatorname{UnitaryInvariantRectNormLike}\) certificate,
the norm-generic equation~(9) rank/residual and relative-residual surfaces
apply verbatim to its underlying \(\operatorname{RectNormLike}\) part.  In
particular, if the exact head/tail sketch certificate and the visible norm
bounds of the previous norm-generic corollary are supplied, then
\[
  \operatorname{rank}(PA)\le r,\qquad
  \xi_{m,n}(A-PA)\le \tau+\kappa,
\]
and if \(A_k\) is a certified best rank-\(k\) approximation for \(\xi_{m,n}\)
and
\[
  \tau+\kappa\le \rho\,\xi_{m,n}(A-A_k),
\]
then
\[
  \operatorname{rank}(A_k)\le k,\qquad
  \operatorname{rank}(PA)\le r,\qquad
  \xi_{m,n}(A-PA)\le \rho\,\xi_{m,n}(A-A_k).
\]

This closes the supplied-norm API layer for unitarily invariant equation~(9)
surfaces.  It still does not prove singular-value comparison, Eckart--Young,
or the construction of best-rank certificates from rectangular SVD data.
Sampling probabilities and sampling laws remain exact mathematical inputs.
\end{corollary}

\begin{formalized}
\leanfit{UnitaryInvariantRectNormLike}
\leanfit{UnitaryInvariantRectNormLike.norm_matMulRectLeft}
\leanfit{UnitaryInvariantRectNormLike.norm_matMulRectRight}
\leanfit{frobUnitaryInvariantRectNormLike}
\leanfit{frobUnitaryInvariantRectNormLike_toRectNormLike}
\leanfit{frobUnitaryInvariantRectNormLike_norm}
\leanfit{equation9RankResidualUnitaryNormSurface}
\leanfit{equation9HeadTailSketchUnitaryNormRankResidualSurface}
\leanfit{equation9RelativeResidualUnitaryNormSurface}
\leanfit{equation9HeadTailSketchUnitaryNormRelativeResidualSurface}
\end{formalized}

\begin{remark}
The paper-level structural relative-error theorem behind equation~(9) is still
open.  The source-SVD and head-tail bridges above narrow the determinant and
residual-certificate sides to exact source factor data plus explicit
head/tail radii.  The explicit-coefficient and source-coefficient bridges now
make the pseudoinverse-derived choice of \(W\) concrete, and the Moore--Penrose
bridge replaces the generic projector hypothesis by four exact equations for a
coefficient table \(C\).  The concrete Gram-inverse bridge then instantiates
that table as \(C_G=\operatorname{nonsingInv}((AZ)^T(AZ))(AZ)^T\) under the
exact determinant condition on the sketch Gram matrix, and the orthogonal
source head-tail Gram split reduces the determinant side to a source-head
sketch Gram plus the tail sketch Gram.  The PSD/determinant bridge proves that
this split gives a nonzero full sketch-Gram determinant as soon as the
source-head sketch Gram is certified positive definite, and the source-factor
positive-definiteness bridge now derives that certificate from
\(\det(\Sigma)\ne0\) and \(\det(V^T Z)\ne0\).  Thus the exact determinant side
is closed for supplied source certificates.  The source-tail norm reduction now
also moves the ambient residual-tail Frobenius obligation to a coordinate
residual-tail obligation under a supplied orthonormal left tail basis, and the
coordinate-tail factorization exposes the exact
\(\Sigma_\perp(V_\perp^T-(V_\perp^T Z)(V_k^T Z)^{-1}V_k^T)\) residual with a
non-sharp Frobenius submultiplicative bound.  The right-tail block-product
theorem now proves \(R_\perp V_\perp=I\) and
\(R_\perp V_k=-(V_\perp^T Z)(V_k^T Z)^{-1}\).  The Frobenius block identity
then proves
\[
  \|\Sigma_\perp R_\perp\|_F^2
  =
  \|\Sigma_\perp\|_F^2+
  \|\Sigma_\perp(V_\perp^T Z)(V_k^T Z)^{-1}\|_F^2
\]
under exact row-completeness of \([V_\perp,V_k]\).  The cross-term
certificate corollary now converts the supplied Frobenius
condition
\(\|\Sigma_\perp(V_\perp^T Z)(V_k^T Z)^{-1}\|_F\le
\varepsilon\|\Sigma_\perp\|_F\) into the displayed
\(\sqrt{1+\varepsilon^2}\) coordinate source-tail rate, and the ambient
source-tail certificate now transports that rate through a supplied
left-orthonormal \(U_\perp\) factor for \(T=U_\perp\Sigma_\perp V_\perp^T\).
Finally, the projected source-tail coupling certificate uses exact
symmetric-idempotent Frobenius contractivity to pass the same bound through a
supplied exact orthogonal or Moore--Penrose column-sketch projector
\(P=(AZ)C\).  The transpose-action spectral certificate then proves the
Frobenius cross-term condition from the exact right-acting operator hypothesis
\(\|((V_\perp^T Z)(V_k^T Z)^{-1})^T x\|_2\le
\varepsilon\|x\|_2\), and composes that certificate through the projected
Moore--Penrose source-tail coupling theorem.  The ordinary operator
certificate corollary now derives that transpose-action hypothesis from the
non-transposed certificate
\(\|(V_\perp^T Z)(V_k^T Z)^{-1}x\|_2\le\varepsilon\|x\|_2\).  The computed
cross-factor perturbation corollary now allows a computed non-probability
factor \(\widehat M\) with operator radius \(\varepsilon\) and Frobenius
distance \(\tau\) from the exact
\((V_\perp^T Z)(V_k^T Z)^{-1}\) factor to feed the same projected
Moore--Penrose theorem with radius \(\sqrt{1+(\varepsilon+\tau)^2}\).  Its
entrywise specialization converts a uniform computed-entry budget
\(|M_{ab}-\widehat M_{ab}|\le\eta\) into the Frobenius radius
\(\sqrt{qr}\eta\), giving projected radius
\(\sqrt{1+(\varepsilon+\sqrt{qr}\eta)^2}\).  The component-certified computed
cross-factor corollary now derives that entrywise budget from separate
computed cross-gram, inverse-factor, and rounded-product radii
\(\alpha,\beta,\rho\), replacing \(\eta\) by
\(\alpha+\beta+\rho\).  The concrete floating-point cross-Gram corollary now
instantiates the \(V_\perp^T Z\) side of that component certificate with the
repository `fl\_matMul' dot-product budget
\(\gamma(fp,n)\sum_j |(V_\perp)_{ja}|\,|Z_{jb}|\).  The square cross-Gram
corollary similarly certifies the computed inverse input \(V_k^T Z\) and its
uniform Frobenius radius \(\sqrt{rr}\,\omega\).  The inverse-entry adapter now
turns a supplied entrywise computed-inverse certificate
\(|Y-\widehat Y|\le\eta\) and computed-left row budget \(\chi\) into the
component radius \(\chi\eta\), hence replaces
\(\alpha+\beta+\rho\) by \(\alpha+\chi\eta+\rho\) in the same projected
computed-factor bound.  The concrete rounded-product corollary now
instantiates the remaining \(\rho\) component by the repository
\(\operatorname{fl\_matMul}\) budget
\(\gamma(fp,r)\sum_b|\widehat X_{ab}|\,|\widehat Y_{bc}|\), so the final
product arithmetic is closed for supplied inverse and operator certificates.
The computed-product Frobenius certificate corollary then replaces that
abstract operator certificate by the visible condition
\(\|\widehat M\|_F\le\varepsilon\), using the deterministic
Frobenius-to-operator handoff.
The product absolute-sum Frobenius certificate corollary further supplies that
condition from
\(\sum_b|\widehat X_{ab}|\,|\widehat Y_{bc}|\le\kappa\) and the same final
product radius \(\rho\), giving
\(\|\widehat M\|_F\le\sqrt{qr}(\kappa+\rho)\).
The perturbed-inverse certificate corollary similarly replaces the raw
inverse-entry assumption \(|Y-\widehat Y|\le\eta\) by the Higham~\S13.1
contract \((A_0+\Delta A)\widehat Y=I\), a componentwise perturbation budget,
and a visible inverse-sensitivity bound.
The Method-A LU inverse certificate corollary goes one step further for a
specific computed inverse surface: the columns of
\(\widehat Y=\operatorname{methodAComputedInverse}_{fp,r}(\widehat L,\widehat U)\)
are generated by the repository forward/back substitution routines, and an
`LUBackwardError' certificate for \(V_k^T Z\) plus the visible
\((3\gamma+\gamma^2)\)-budget supplies the same \(\eta\) radius.
The computed-input Method-A corollary now removes the silent exact-input
assumption from that LU surface: if the LU routine is certified for the
rounded square cross Gram
\(\widehat A_0=\operatorname{flRightSketchCrossGram}_{fp}(V_k,Z)\), the
rounded-input dot-product budget transfers the certificate to \(V_k^T Z\) at
coefficient \(\varepsilon_{\rm LU}+\mu\), and the Method-A sensitivity budget
uses \((\varepsilon_{\rm LU}+\mu)+2\gamma+\gamma^2\).
The Doolittle-generated computed-input corollary now discharges that displayed
LU certificate from a dense Doolittle recurrence certificate for
\(\widehat A_0\), giving coefficient
\((\gamma+\mu)+2\gamma+\gamma^2\) for Method A.
The dense-loop Doolittle residual-compression corollary then exposes the
literal rounded upper-entry fold, lower-entry numerator fold, and rounded
division that a dense implementation computes, and proves that the two
displayed compression inequalities are sufficient to produce the same
`DoolittleLU' witness.
The absolute-budget dense-loop corollary lowers that certificate one step
further: a concrete implementation may prove ordinary absolute residual budgets
\(B^U,B^L\), then prove their dominance by the displayed relative
\(\gamma\)-radii, and the same Method-A equation-(9) bounds follow.
The rounded-product Doolittle fold-budget corollary now supplies the
subtraction-accumulation part of those absolute budgets for the literal upper
fold and lower numerator fold, measured against the rounded products actually
subtracted by the implementation.
The subsequent exact-product, masked-prefix, lower-division, product-margin,
and lower-numerator-margin corollaries pay the rounded-product transfer,
prefix reindexing, computed-pivot division surface, and source-shaped
stored-entry margins.  The exact-target gap corollary gives a valid
conditional transfer from exact pre-rounded target gaps to those stored-entry
margins, with explicit residual budgets for all non-probability arithmetic in
the literal routine.  The following triangle-obstruction proposition then
shows that this exact-target source route is too strong whenever the FP excess
is positive, so it should not be treated as the remaining ordinary
dense-Doolittle no-cancellation proof path.
The norm-generic head-tail corollary now closes the exact algebraic
rank/residual theorem surface for any supplied \(\xi\) satisfying the displayed
nonnegativity and triangle assumptions.  The Frobenius-instance corollary now
instantiates that surface for \(\xi(X)=\|X\|_F\), transports the Frobenius
best-rank and head-tail certificates into the norm-generic API, and exposes
the existing left/right orthogonal-invariance handoffs.  The concrete
all-unitarily-invariant norm theorem still requires the missing norm API,
singular-value comparison, and Eckart--Young instantiations.
The next missing foundations
are the rectangular SVD existence/Eckart--Young layer, general unitarily
invariant variants of the cross-term condition, randomness-derived
certificates for the displayed operator condition, a non-vacuous
stored-entry/component-dominance or revised-compression source for the
dense-Doolittle Method-A route, and
computed SVD/projector/Gram, coefficient, and downstream product routines.
\end{remark}

\section{What Is Not Proved}

The repository does not yet prove all RandNLA accuracy theorems mentioned in
the CACM survey.  The open paper-level items are tracked in
\lean{docs/RANDNLA_CACM_NOT_PROVED_LEDGER.md}.  The active bottleneck ledger is
\lean{docs/RANDNLA_CACM_BOTTLENECK_LEDGER.md}; it names the exact blocking
theorem family, closed dependencies, open dependencies, and rejected proof
routes for stalled paper-level rows.  The live theorem ledger is
\lean{docs/RANDNLA_CACM_THEOREM_LEDGER.md}; it records extracted claims,
events, hypotheses, Lean names, and next proof steps.  The proof-source ledger
is \lean{docs/RANDNLA_CACM_PROOF_SOURCE_LEDGER.md}; it records external
primary sources and chosen proof routes for incomplete paper proofs.  In
particular, the repository still does not prove:
\begin{itemize}
\item the untruncated/general-rectangular version of Algorithm~1
  equation~(2).  The source-aligned square theorem with truncation
  \(\tau=\varepsilon/(2n)\), source sample budget, deterministic truncation
  transfer, and derived floating-point \(\gamma\)-budget transfer is now
  formalized for the nonzero truncated squared-magnitude denominator case.
  Earlier weaker fixed-vector,
  finite-test-set, supplied finite-cover/Markov, Frobenius/Markov, and
  scalar-entry/union-bound routes remain available as fallback statements.
  The residual has been expressed
  as a sum of zero-mean increments with second-moment bounds, and a
  self-adjoint dilation of that residual is now formalized with trace-of-square
  identities plus quadratic-form, Loewner, PSD, and trace variance proxies.
  The one-step zero-mean and variance proxies have also been lifted through the
  independent product trace law to give a mean-zero full dilation residual and
  trace-law expected squared-increment variance matrix bounds. A generic
  scalar MGF factorization, exponential-Markov tail for sums of arbitrary
  one-step trace statistics, no-hidden-Lieb iid C\(^*\) trace-MGF domination,
  finite-real trace-exponential adapter, residual-dilation trace-MGF
  instantiation, trace-exponential Markov-to-eigenvalue step, scalar
  Bernstein parabola constants, and explicit scalar-to-operator CFC
  Bernstein lift are also formalized.  The generic one-sample
  matrix-CGF/log-MGF variance-proxy bound
  \(K\preceq g(\theta,R)V\), its support-aware variant, and its truncated
  Algorithm~1 dilation-increment instantiation are now formalized too.  The
  one-sample C\(^*\) dilation variance proxy, the negative-support lower-tail
  companion, the positive and negative scalar trace-MGF bounds, and the
  parameterized upper-tail and two-sided eigenvalue tails are also formalized,
  including the explicit \(1-\delta\) scaled-eigenvalue corollary with
  \(T=\log(2B/\delta)\).  The scaled event is converted to a rectangular
  spectral event at radius \(\log(2B/\delta)/\theta\), and the exact Bennett
  optimizer \(\theta=\log(1+Lr/W)/L\) gives a radius-\(r\) theorem under the
  corresponding Bennett budget.  The square source route now also has the
  sharp \(V=n\|\widehat A\|_F^2/s^2\) variance proxy and no-\(\sqrt2\)
  truncated support radius threaded through the two-sided \(1-\delta\)
  eigenvalue skeleton, plus matching source-sharp scaled-radius and
  Bennett-radius rectangular spectral-event corollaries.  Both the
  conservative denominator bridge \(q\le r^2/(2W+Lr)\) and the source
  denominator bridge \(q\le r^2/(2W+\frac23Lr)\) are closed.  The final source
  sample-budget theorem, deterministic truncation transfer, and support-aware
  floating-point gamma-budget transfer are also closed.  The
  product-grid cover reduction and cardinality
  \(|\alpha^n|=|\alpha|^n\) are formalized, but the covering route still lacks
  an explicit fine one-dimensional interval grid and sharp finite-net scalar
  tails;
  moreover, the primary Drineas--Zouzias theorem uses a truncated matrix
	  \(\hat A\); the deterministic truncation transfer, probability-one support,
	  bounded-increment, two-sided Loewner boundedness, bounded-square, and
	  simultaneous Bernstein-boundedness prerequisites for sampling \(\hat A\),
	  plus a weak accumulated \(sL\) spectral consequence and a support-aware finite-test MGF
	  theorem for truncated dilation quadratic forms, a one-sided dilation
	  Loewner event-transfer theorem, local-to-mathlib symmetry/PSD/Loewner bridges,
	  a mathlib-backed finite Hermitian eigenvalue/trace/PSD bridge, and an
	  eigenvalue-event plus per-eigenvalue union-bound transfer layer, as well as
	  the scalar matrix-exponential normalization \(\operatorname{tr}(\exp(LI))=d e^L\),
	  diagonal trace-exponential normalization, matrix-exponential symmetry preservation,
	  and the Hermitian CFC trace identity
	  \(\operatorname{tr}(E_{\mathrm{cfc}}(M))=\sum_i e^{\lambda_i(M)}\),
		  plus the repository-native matrix-exponential trace identity
		  \(\operatorname{tr}(\operatorname{Exp}(M))=\sum_i e^{\lambda_i(M)}\),
		  its lower-tail companion
		  \(\operatorname{tr}(\operatorname{Exp}(-M))=\sum_i e^{-\lambda_i(M)}\),
		  and the upper- and lower-tail MGF-to-eigenvalue tail/complement interfaces
		  \(\Prob\{\exists i:\lambda_i(M)\ge T\}
		    \le e^{-T}\E[\operatorname{tr}(\operatorname{Exp}(M))]\),
		  \(\Prob\{\exists i:\lambda_i(M)\le T\}
		    \le e^{T}\E[\operatorname{tr}(\operatorname{Exp}(-M))]\),
		  together with the two-sided
		  \(\Prob\{\forall i:|\lambda_i(M)|<T\}\) combiner and the scalar
		  Loewner-to-trace-exponential bound
		  \(M\preceq cI\Rightarrow
		  \operatorname{tr}(\operatorname{Exp}(M))\le d e^c\),
		  plus the complex C\(^*\)-matrix embedding of finite real matrices and
		  its subtraction, identity, scalar-identity, self-adjointness,
		  local PSD, local Loewner-order preservation, strict-positive
		  scalar-identity regularization, and regularized operator-log
		  preservation lemmas, together with the complex C\(^*\)-matrix trace
		  bridge back to the repository-native finite real trace, cyclic trace
		  algebra \(\operatorname{tr}_{\mathbb C}(XY)=\operatorname{tr}_{\mathbb C}(YX)\),
		  spectral-order trace positivity/monotonicity,
		  and scalar CFC-exponential trace normalization
		  \(\operatorname{tr}_{\mathbb C}(\exp_{\mathrm{cfc}}(aI))=e^a d\),
		  plus finite-probability expectation for complex scalars and complex
		  C\(^*\)-matrix-valued random variables, trace-expectation
		  interchange, compatibility with embedded finite real matrices,
		  complex and real expectation weighted-sum forms, C\(^*\)-order
		  nonnegativity, strict positivity, and monotonicity preservation,
		  positive-support existence for finite probability laws, strict-positive
		  regularization after adding \(\varepsilon I\), and the finite
		  Jensen adapter for concave real-valued functions of C\(^*\)-matrix
		  random variables, plus the self-adjoint
		  \(\log(\exp X)=X\) identities in normed-space, complex CFC, and
		  real CFC exponential forms, self-adjoint exponential strict
		  positivity in normed-space, complex CFC, and real CFC forms, and
		  the strictly positive
		  \(\exp(\log A)=A\) identities in normed-space, complex CFC, and
		  real CFC exponential forms, together with the normalization of the
		  local Lieb functional against the standard normed-algebra
		  exponential and the \(H=0\) trace case, plus the affine \(H=0\)
		  special case of the local Lieb concavity target, and the
		  conditional and unconditional one-step Tropp--Jensen trace-MGF
		  inequalities, together with scalar and finite-vector
		  relative-entropy nonnegativity, the finite log-sum inequality,
		  scalar and finite-vector relative-entropy joint convexity,
		  relative-entropy vocabulary
		  \(D(A;B)\), the diagonal normalization \(D(A;A)=0\), and
		  scalar-identity matrix relative-entropy nonnegativity
		  \(D(aI;bI)\ge 0\) for \(a,b>0\), as well as the real diagonal
		  matrix reduction
		  \(D(\Delta(a);\Delta(b))=D_{\mathrm{vec}}(a;b)\) and its
		  nonnegativity and joint convexity for positive diagonals, plus the
		  left/right multiplication endomorphism substrate
		  \(L_A,R_A\), their product/power algebra, and their
		  unit/strict-positive invertibility, plus the
		  explicit ratio endomorphism \(L_XR_A^{-1}\) and its base-point
		  normalization, plus the finite Kronecker lift substrate
		  \(A\otimes I\), \(I\otimes H\), its affine laws, commutation/product
		  identities, positive-definiteness preservation, trace
		  normalizations, product-index vectorization action, and
		  vectorized-identity trace pairing and polynomial powered
		  trace-pairing identity, plus its finite-polynomial packaging,
		  plus continuity of the finite quadratic form and polynomial evaluation
		  in the matrix argument,
		  plus the Hansen--Pedersen
		  source-route split into
		  ordinary positive-cone operator convexity, transfer, and assembled
		  Jensen targets, including identity-function sanity checks, the
		  assembly adapter from the concrete inputs to the assembled target,
		  and the direct unit-interval-kernel proof of the ordinary
		  \(x\log x\) operator-convexity input,
		  plus Bendat--Sherman-route derivative monotonicity of
		  \(1+\log x\) on the strictly positive cone, plus the
		  source-faithful first-divided-difference route names and
		  scalar normalizations, plus the unital fractional-kernel
		  monotonicity theorem for \(x\mapsto 1-(1+x)^{-1}=x/(1+x)\)
		  and the scaled theorem for \(x\mapsto x/(s+x)\), \(s>0\),
		  plus finite nonnegative-combination closure for such kernels
		  on the nonnegative C\(^*\)-matrix cone, plus the conditional adapter from
		  divided-difference monotonicity and the
		  Bendat--Sherman divided-difference bridge to operator convexity,
		  plus the base-point CFC normalization and the full
		  divided-difference monotonicity theorem,
		  plus finite matrix inverse convexity on the positive-definite cone
		  by Schur complements and its finite C\(^*\)-matrix CFC
		  inverse-kernel bridge, plus shifted inverse-kernel convexity
		  for \(x\mapsto(s+x)^{-1}\), \(s>0\),
		  plus the corrected scalar \(x(x-1)/(u+(1-u)x)\) kernel identity,
		  its CFC affine-plus-shifted-inverse decomposition, the CFC integral
		  representation of \(x\log x\), and the integrated direct proof of
		  ordinary positive-cone operator convexity of \(x\log x\),
		  plus the
		  conditional
		  relative-entropy route theorem showing that joint convexity and
		  the normalized variational formula imply the local Lieb
		  trace-concavity target, and the optimizer-candidate equality part
		  of that normalized formula, plus the reduction of variational
		  maximality to relative-entropy nonnegativity, plus the reduction
		  of relative-entropy nonnegativity to the generalized Klein
		  first-order trace inequality, the equivalence between that
		  first-order inequality and local relative-entropy nonnegativity,
		  the Hermitian spectral-overlap expansion, and the positive-spectrum
		  scalar entropy first-order specialization, plus the bridge from the
		  separated Hermitian entropy statement to the compact complex
		  C\(^*\)-matrix logarithm form of generalized Klein, local
		  relative-entropy nonnegativity, the normalized entropy variational
		  formula, and the reduction from joint convexity alone to the local
			  Lieb target, the generic CFC Bochner-integral order lift under
			  the \texttt{cfc\_integral} side conditions, the
			  scalar-integral-to-CFC equality for the normalized logarithmic kernel,
			  the scalar normalization back to \(x\log x\), the auxiliary
			  square-kernel affine-plus-shifted-inverse decomposition, and
			  normalized-kernel CFC monotonicity, the all-finite-size
			  packaging of ordinary \(x\log x\) operator convexity, and the
			  Hansen--Pedersen block-compression algebra plus range-projection
			  algebra plus block-diagonal star/order substrate plus
			  block-diagonal CFC decomposition plus unitary CFC conjugation,
				  strict-positive domain preservation, block-compression
				  domain preservation, rectangular compression order preservation,
				  the reflection-average CFC pinching inequality, and the
				  shifted-inverse kernel nonlinear corner identity, the
					  compression-integral assembly, the full \(x\log x\) corner
					  identity, and the concrete two-point Hansen--Pedersen Jensen
					  theorem, plus the affine-corrected normalized entropy-kernel
					  ordinary operator-convexity and Jensen theorems, plus the
					  CFC square-root and inverse-square-root perspective substrate,
					  plus the ordinary finite perspective theorem for the normalized
					  entropy kernel, plus the finite-polynomial Kronecker
					  trace-representation layer and the source-faithful
					  right-multiplication polynomial identity for
					  \(v_I^*p(L_XR_A^{-1})R_Av_I\), including finite-matrix
					  polynomial CFC, Weierstrass approximation wrappers, and the
					  explicit uniform-approximation transfer plus the
					  positive-spectrum construction of the approximating
					  polynomial sequence for the
					  entropy-kernel CFC trace term, the finite-polynomial
					  overlap formulas for \(\operatorname{tr}(X^kA(A^{-1})^k)\)
					  and real-polynomial sums, the common limiting overlap
					  theorem, and the source-faithful superoperator
					  trace-representation theorem, plus the product-index
					  \(L_X\)/\(R_A\) lift domain lemmas and the real
					  quadratic-form monotonicity needed to extract scalar
					  inequalities from product-index Loewner inequalities,
					  plus joint convexity of the ordinary product-index
					  perspective trace
					  \(v_I^*P_f(L_X,R_A)v_I\), and the first
					  CFC-commutation step showing that \(L_X\) commutes with
					  \(R_A^{-1/2}\) so the normalized perspective argument can
					  be reordered, plus the inverse-square-root square
					  identity putting that normalized argument in the
					  \(L_XR_A^{-1}\) shape and proving this ratio commutes
					  with \(R_A^{1/2}\), the final CFC transport and
					  product-index equality bridge
					  \(v_I^*P_f(L_X,R_A)v_I =
					  v_I^*f(L_XR_A^{-1})R_Av_I\), finite-dimensional
					  relative-entropy joint convexity, arbitrary-\(H\)
					  Lieb trace concavity, and the one-step Tropp trace-MGF
					  inequality, together with the one-coordinate Algorithm~1
					  product-law adapter from this one-step inequality to the
					  squared-magnitude trace sampler and the iid sampled-sum
					  trace-MGF iteration over the product trace law,
					  and the source-aligned scalar matrix-CGF/Bernstein
					  constants proving the half-budget event for sampling
					  \(\hat A\), are now formalized.  The advisory
			  Bendat--Sherman bridge from divided-difference monotonicity to
			  operator convexity remains advisory, and the exact untruncated
			  squared-magnitude law remains a separate theorem target.  A
			  locality check also found that mathlib currently supplies
	  operator-log monotonicity but explicitly lists operator-log concavity,
	  operator convexity of \(x\log x\), and operator concavity of \(t^p\) as TODOs
	  in its continuous-functional-calculus files.  The repository now supplies
	  the \(x\log x\) operator-convexity theorem and finite-dimensional Lieb
	  theorem locally, the iid trace-MGF iteration locally, and the
	  source-aligned Algorithm~1 Bernstein specialization locally; no generic
	  matrix Bernstein theorem is used as a hidden hypothesis;
\item a reusable general-purpose matrix Bernstein or Khintchine API beyond the
  specialized Algorithm~1 and Algorithm~2 source-aligned routes;
\item a cosmetic one-hypothesis restatement of the sharper Algorithm~2
  leverage-score equation~(7) sample budget.  The current Lean theorem already
  proves the source-aligned finite-Loewner high-probability event and the
  matching floating-point finite-Loewner transfer from two explicit Bennett
  sample-budget inequalities, one for the upper tail and one for the lower
  tail.  What remains here is only a presentation simplification if one wants
  to replace those two sufficient inequalities by a single more conservative
  displayed condition;
\item the rectangular QR/preconditioner FP theorem for
  equation~(8).  The sampled-row objective algebra itself is now formalized
  with canonical coordinates, the augmented-span basis theorem is formalized
  for \(d=\dim\operatorname{span}\{A_{*,j},b\}\), the exact augmented-span
  Bennett sample-budget least-squares theorem is formalized for \(d>1\), the
  abstract rounded-Gram FP transfer is formalized with the explicit Gram
  perturbation radius, the full identity-basis fallback is formalized for
  \(d=m\), and the literal rounded sampled/scaled \(A,b\) construction has
  rowwise residual, deterministic objective perturbation, high-probability
  rounded-minimizer, additive solver-gap, forward-error-certificate,
  perturbed-Gram-system transfer, local QR backward-error-spec transfer, and a
  concrete normal-equations/Cholesky solver transfer.  The deterministic
  rectangular perturbed-normal-equation adapter, the explicit norm-budget
  handoff into the local QR specification, and the induced Gram/RHS
  perturbation-budget lemmas are also formalized.  As route-A substrate,
  left and right multiplication of a rectangular matrix by a compatible square
  orthogonal matrix are now formalized to preserve the rectangular Frobenius
  norm, and compatible square left/right factors now satisfy rectangular
  Frobenius submultiplicativity; rectangular square-left action also has the
  needed identity, associativity, and additivity lemmas.  The rectangular
  one-step and multi-step orthogonal-transformation accumulation theorems are
  formalized as well, with the multi-step theorem using a geometric
  perturbation radius; the matching right-hand-side/vector accumulation theorem
  is also formalized with the same geometric radius; and the orthogonal
  handoff into rectangular normal equations is closed.  The exact
  zero-prefix and active-column Householder substrates are also closed.  The
  strong common panel/update contract
  \(\lean{HouseholderPanelAppError}\) is formalized, but the active
  source-faithful route is now the columnwise contract
  \(\lean{HouseholderColumnwisePanelAppError}\), matching Higham's
  columnwise Householder QR proof in which the perturbation matrix may depend
  on the panel column.  The columnwise accumulation theorem now supplies one
  final \(Q\), per-column perturbation bounds, and the right-hand-side
  perturbation bound.  The formed-reflector route via
  \(\lean{fl_householderApplyExplicit}\) is proved using the local
  matrix-vector product theorem.  The compact dot/scale/subtract route via
  \(\lean{fl_householderApplyCompact}\) is also proved with an explicit
  deterministic component budget and a visible budget-domination adapter into
  \(\lean{HouseholderAppError}\) and the columnwise panel contract.  The
  compact supplied-sequence adapter is also closed, so compact panel/RHS
  updates now feed the columnwise geometric accumulation theorem.  The stored
  trailing loop also has formalized \([R;0]\)/top-RHS shape facts, top-block
  triangular-solve handoffs into \(\lean{LSQRSolveBackwardError}\), and several
  visible-domain nonbreakdown routes: active-entry, dimensioned trailing-norm,
  leading-dual, local inverse row/Frobenius/\(\infty\)-norm, nonsingular
  leading blocks, triangular principal minors, and diagonal-dominant triangular
  inverse estimates.  The final Higham-style columnwise stored QR
  factorization assembly is closed: the Lean theorem now packages the common
  orthogonal factor, perturbation bounds, final \([R;0]\) block form, top RHS,
  and upper triangularity of \(R\).  Still open is the solver/preconditioner
  theorem that discharges the remaining visible nonzero-diagonal,
  determinant/rank, conditioning or inverse-budget, and compact-smallness
  hypotheses from the actual computed loop, rather than keeping them as domain
  assumptions.  The proof-source route is now a genuine mathematical choice
  beyond this columnwise factorization theorem: Higham's standard Householder
  QR theorem is columnwise/normwise, while the stronger row-wise weighted
  least-squares theorem of Cox--Higham requires pivoting or row sorting plus
  the correct sign convention.  The user-selected off-diagonal-control route
  now has a named Lean invariant
  \(\lean{StoredQROffDiagonalControlInvariant}\) and a solver-facing theorem
  showing that this invariant yields the local
  \(\lean{LSQRSolveBackwardError}\) certificate with repository final
  Gram/RHS budgets.  That invariant has also been reduced to the source-shaped
  data \(\lean{StoredQRSourceOffDiagonalControl}\): local triangular leading
  blocks, nonzero local diagonals, row-wise off-diagonal domination, and the
  stored-sequence product bound.  The triangular leading-block field is now
  derived from the stored recurrence using the repository's
  prefix-lower-zero theorem, so the remaining open source-shaped obligations
  are nonzero displayed diagonals, row-wise off-diagonal domination, and the
  compact-product bound.  The nonzero-diagonal obligation has been reduced
  further: previous displayed diagonals follow from the stored signed-alpha
  prefix-diagonal theorem and a square-root nonbreakdown budget, leaving the
  current pivot nonzero condition visible.  The current-pivot condition has
  now also been reduced to nonsingularity of the displayed local leading block
  \(S_k\) by the stored lower-zero determinant bridge, and the square-root
  budget has been replaced by the dimensioned norm-square margin using the
  repository's Householder norm-square bridge.  The row-wise off-diagonal
  domination field is now decomposed into a row-growth budget plus a diagonal
  lower-bound obligation, matching the Cox--Higham pivoted/sorted route.  The
  row-growth half now has a QR-layer propagation theorem from Cox--Higham
  stage budgets to displayed leading-block off-diagonal row budgets, and a
  solver-facing handoff now composes those stage budgets with the source-control
  and \(\lean{LSQRSolveBackwardError}\) certificates under explicit diagonal
  lower bounds.  The stage-entry assumptions have also been split into
  one-step active-block and prefix displayed-row recurrences, matching the
  zero-prefix/active-row split in the exact signed-stage reflector proof.
  A finite global compact-step budget now supplies the displayed off-diagonal,
  active-block, and prefix-row one-step fields from a single scalar recurrence
  \( \Gamma_mB_t+\Delta_t^{\max}\le B_{t+1}\).  The same recurrence, together
  with nonnegative budgets, now also supplies monotonicity on the QR stage
  horizon \(a\le b\le n\).  The completed-column preservation field is now
  derived from the stored-loop lower-zero invariant and the zero-prefix
  Householder support lemma.  The per-pivot
  compact-product hypotheses are now also packaged by a finite global product
  budget \(\Pi_{\max}<1\), and the equation~(8) objective theorem is assembled
  from this active/prefix global-product QR route under the remaining visible
  fields.  Still open is proving
  local leading-block nonsingularity, the norm-square nonbreakdown margin,
  diagonal lower bounds, row-max scalar defect and row-max comparison fields,
  and global compact-product smallness from genuine pivoting, ordering,
  or off-diagonal-growth assumptions, or explicitly keeping them visible as
  domain hypotheses.  The
  local route-elimination
  theorems already rule out
  deriving these visible fields from full rank, whole-matrix determinant
  nonzeroness, positive trailing norm, exact final QR shape, finite local
  conditioning, diagonal dominance alone, product smallness alone, active-pivot
  plus active/off-diagonal budgets alone, or the ordinary exact no-pivot
  recurrence.  For the current unpivoted theorem family, the repository keeps
  these fields as explicit source/domain assumptions.  This scopes the
  implemented conditional QR stability theorem, but it does not close the
  generic paper-level claim that the ordinary unpivoted loop supplies them.
  Adjacent no-pivot QR adapters remain frozen by the bottleneck ledger.  The
  positive QR route has therefore switched to the Cox--Higham pivoted/sorted
  weighted least-squares theorem family: the signed Householder denominator
  lower bound \(2\|x_{\mathrm{tail}}\|_2^2\le v^Tv\) and the column-pivoting
  active-column norm comparison are now formalized.  The scalar endpoint
  \(|\beta v^Ty_{\mathrm{tail}}|\le\sqrt2\) and the one-step row-growth
  estimate \(|a_{ij}-\phi a_{ik}|\le(1+\sqrt2)B\) are also formalized.  The
  row-sorting stage-accumulation step is now formalized as a scalar induction
  and sorted-active-row transfer, and the pivot-row active-tail norm step from
  equations (4.4)--(4.5) is formalized with the explicit ambient factor
  \(\sqrt m\).  The scalar row-wise accumulated perturbation
  bound is also formalized as an affine recurrence plus an exact/computed
  triangle-inequality transfer, and the concrete stored rounded panel step now
  supplies the per-step compact Householder FP budgets for that recurrence.
  A direct stored-panel row-magnitude adapter is also formalized, so the
  source-shaped route can combine exact same-reflector row growth with the
  concrete compact FP budget without routing through a separate exact-sequence
  error field.  The non-pivot active-row exact same-reflector bridge is also
  formalized: the signed Householder matrix-vector update is rewritten into
  Cox--Higham's scalar row update before applying the row-growth bound.  The
  exact signed pivot-row same-reflector bridge is now formalized as well, in
  the same ambient-dimension form as the local pivot-row norm step.  The
  unified active-row case split and the exact signed-pivot panel sequence
  propagation theorem are now formalized too, with the stage budgets and their
  geometric \(c_m^tB_0\) form kept explicit.  The active-block one-step
  adapter is also formalized: a single bound over the active rows and remaining
  columns supplies the separate row and column fields for one signed-pivot
  step.  The exact sequence now also has active-block wrappers, so the separate
  row and column stage fields can be replaced by one active-block budget family
  at every stage.  The active-block family itself is now propagated from an
  initial entrywise bound under monotone active row/column windows, and the
  positive-norm field is reduced to a nonzero-active-block witness plus
  pivot-maximality; the raw pivot-maximality inequality is in turn supplied by
  the finite active max-pivot selector.  The nonzero-active-block witness can
  now be supplied by positivity of the displayed active-block sum of squares,
  and the concrete column-swap bridge proves that swapping the selected active
  max column into the displayed active position supplies the post-swap
  pivot-maximality field.
  The stored-panel handoff is now also available in stage-budget form: a
  signed-pivot exact step bounded by
  \(c_mB_t\) transfers to the rounded stored step with the concrete compact
  Householder FP budget.  The rounded active-block recurrence is now
  formalized under visible signed-pivot nonbreakdown, pivot-maximality,
  storage, active-window, and compact-budget fields, and the prefix displayed
  rows are now propagated from a one-step prefix compact-budget recurrence
  using the active-block theorem and the exact same-reflector prefix/active
  split.  A finite global compact-step budget now packages the off-diagonal,
  active-block, and prefix-row compact-update recurrences into one scalar
  stage-budget condition.  The raw-stage positive
  active-block mass field is now supplied from local QR rank/determinant
  nonbreakdown data, including a stagewise determinant wrapper.  A checked
  route-elimination counterexample also shows that this positive active-block
  mass, even together with upper-triangular nonsingular leading blocks, does
  not imply the row-wise off-diagonal-control field.  A second checked
  counterexample shows that adding a valid strict-upper-entry row-growth
  budget still does not imply the diagonal lower-bound/nonbreakdown field:
  the local witness has row budget \(2\) and first diagonal magnitude \(1\).
  The Lean structure \(\lean{StoredQRDisplayedRowBudgetControl}\) now names
  this residual row-budget plus diagonal-lower-bound package explicitly, and
  the solver wrapper consumes it as a visible source/domain assumption; it is
  not advertised as a proof of the package.  A newer route-1 handoff shows
  that a scalar row-max contraction
  \(R^{\max}_{k,i}\le \rho |(S_k)_{ii}|\) with \(\rho\le 1\) is sufficient to
  build the package and feed the local solver theorem, but the contraction
  itself is still a source/domain invariant to prove for the chosen concrete
  loop.
  The active/prefix stage-diagonal lower-bound family is now similarly
  packaged by the finite scalar defect
  \(\lean{storedQRStageDiagLowerDefectBudget}\): nonpositivity of that scalar
  supplies \(B_k\le |(S_k)_{ii}|\) for every displayed off-diagonal row
  \(i<k\), and the converse theorem packages that same pointwise family back
  into scalar nonpositivity.  The local source-control and solver wrappers
  consume the scalar condition directly.  The active max-pivot sibling now also
  derives the raw pivot-maximality field from the finite selector policy on
  this scalar route.
  A row-max bridge now proves this scalar condition from nonpositive row-max
  defect, or from local diagonal dominance, provided the remaining comparison
  \(B_k\le R^{\max}_{k,i}\) is supplied for each displayed off-diagonal row.
  That comparison is now also packaged by the finite scalar defect
  \(\Delta_{\mathrm{cmp}}\): nonpositivity extracts
  \(B_k\le R^{\max}_{k,i}\), and the converse theorem packages a pointwise
  comparison proof back into scalar nonpositivity.  The scalar row-max bridge
  therefore has a two-scalar form.  The local source-control, local solver,
  and probability-level active-pivot wrappers now consume this two-scalar form
  directly, so the displayed comparison family no longer appears at those
  theorem surfaces; nevertheless both scalar row-max conditions still need
  concrete proof or visible hypotheses.
  The diagonal-dominant scalar-comparison local source-control and solver
  siblings sharpen that interface when local diagonal dominance is supplied:
  diagonal dominance now gives leading-block determinant nonzeroness, and
  diagonal dominance together with the scalar comparison defect gives the
  scalar stage-diagonal condition.  Thus the separate determinant and row-max
  scalar-defect fields no longer appear on that local branch, but local
  diagonal dominance itself, the scalar comparison defect, the conditioning/dual
  fields on the \(\kappa\)-budget branch, active-pivot policy,
  compact-product smallness, and the final generic QR/preconditioner theorem
  still remain visible or open.
  The probability-level equation~(8) sibling now exposes this same narrowed
  surface at the sampled objective theorem: samplewise diagonal dominance
  supplies the determinant and row-max scalar-defect fields, and
  \(\Delta_{\mathrm{cmp}}\le0\) supplies the displayed comparison internally.
  This is a lift of the local reduction only; it still leaves local diagonal
  dominance, the scalar comparison defect, conditioning/dual compact-budget
  data, active-pivot policy, raw compact-product smallness, and the final
  generic QR/preconditioner theorem visible or open.  The raw product field on
  this active diagonal-dominant scalar-comparison branch is now reducible to the
  canonical finite-max smallness scalar.  On the concrete-dual finite-max
  sibling, the auxiliary \(\kappa_k\), \(K_k\), and dual compact-budget fields
  are also removed by deriving norm-square nonbreakdown from diagonal dominance
  and the per-pivot product-sequence budget.  The finite-max scalar smallness
  inequality itself remains a visible numerical/domain obligation.  A further
  local source-denominator rational-gamma sibling reduces that active local
  source-control/solver product surface to source-denominator nonbreakdown, a
  unit-roundoff cap, \((m:\mathbb R)U_{\mathrm{cap}}<1\), and the canonical
  rational-gamma cap-smallness scalar, but those source-denominator and cap
  fields are still visible obligations.  The sampled equation~(8) sibling
  applies the same source-control certificate per trace, so the probability
  theorem has the same source-denominator/cap surface instead of the assembled
  finite-max compact-product scalar.  The source-control, local solver, and
  sampled actual-unit siblings specialize the cap to \(fp.u\) and derive the
  former \texttt{gammaValid} guards from
  \((m:\mathbb R)fp.u<1\) or \((s:\mathbb R)fp.u<1\), but the actual-unit
  scalar smallness, source-denominator, diagonal-dominance, scalar-comparison,
  active-pivot, and signed-stage obligations remain visible.  On the
  stored-lower actual-unit branch, the raw source-denominator field is now
  derived from the stored recurrence, signed-alpha definition, and local
  diagonal dominance, leaving actual-unit scalar smallness, local
  diagonal-dominance/off-diagonal control, scalar comparison, active-pivot, and
  signed-stage loop invariants as the remaining visible QR/preconditioner
  obligations.  The stored-recurrence scalar-smallness route is also now
  ruled out: those stored-loop, diagonal-dominance, signed-alpha, and
  actual-unit validity fields do not by themselves force the canonical
  rational-gamma smallness inequality.  The stored-recurrence
  diagonal-dominance route is likewise ruled out: stored recurrence,
  signed-alpha choice, nonzero source denominators, and actual-unit validity do
  not by themselves force local diagonal dominance.
  A new route-elimination theorem shows that local diagonal dominance alone
  cannot discharge the scalar comparison defect: the diagonally dominant
  witness with displayed strict-upper row maximum \(2\) and stage budget \(3\)
  has \(\Delta_{\mathrm{cmp}}>0\).
  The same witness satisfies the displayed compact-product inequality with
  budget \(B=1/16\), so product smallness does not rescue this diagonal-
  dominance shortcut.
  The scalar comparison defect is not forced by the current active-pivot
  active/off-diagonal budget surface, and product compact-smallness does not
  rescue that route: the finite-max extractor would otherwise imply the
  already-refuted displayed comparison \(B_k\le R^{\max}_{k,i}\).
  Nor is it forced after adding diagonal dominance and product compact-smallness
  to that same active surface: the diagonally dominant witness with
  \(B=1/16\) still satisfies active max-pivot and active/off-diagonal budgets
  while \(\Delta_{\mathrm{cmp}}>0\).
  Even granting the scalar row-max defect does not make the scalar comparison
  defect automatic; the checked diagonally safe witness has
  \(R^{\max}_{1,0}=2\), displayed diagonal magnitude \(3\), and stage budget
  \(3\), so \(\Delta_{\max}\le 0\) but
  \(\Delta_{\mathrm{cmp}}>0\).  Product compact-smallness still does not
  rescue this stronger route.  Nor does \(\Delta_{\max}\le0\) alone imply the
  downstream stage-diagonal scalar condition: with the same witness and stage
  budget \(4\), the displayed defect \(4-3\) makes
  \(\Delta_{\mathrm{diag}}>0\).  Adding active max-pivoting,
  active/off-diagonal budget control, and compact-product smallness with
  budget \(B=1/16\) still does not rescue this scalar-stage shortcut.
  The active-pivot row-max-defect counterexample rules out deriving
  nonpositive row-max defect from active max-pivoting plus active/off-diagonal
  budget bounds alone.  A separate compact-product counterexample rules out
  deriving the same nonpositive row-max defect merely from upper-triangular
  nonsingularity plus a satisfiable product compact-smallness inequality, and
  the product/active-budget strengthening shows that adding the active-pivot
  budget surface still does not rescue this row-max defect route.  The
  same product-smallness surface also fails to imply the scalar
  stage-diagonal lower-bound condition for the fixed stage budget \(B_k=2\),
  and adding the active-pivot budget surface still does not rescue that scalar
  stage-diagonal route.  The active-only implication is also refuted as a
  consequence: if active max-pivoting and active/off-diagonal budgets alone
  proved the scalar stage-diagonal condition, the product-plus-active
  implication would follow.
  Even adding product compact-smallness to the active-pivot budget surface does
  not rescue the comparison \(B_k\le R^{\max}_{k,i}\).  The active-pivot
  oversized-budget counterexample separately rules out deriving that comparison
  from the same active surface.
  The probability-level equation~(8) wrapper now exposes these row-max fields
  directly at the sampled least-squares objective theorem, rather than hiding
  them behind a diagonal lower-bound family, but it does not prove them.  Its
  horizon-clamped sibling applies the same row-max bridge and then reuses the
  active-pivot horizon theorem, so the visible row-max probability surface no
  longer exposes samplewise global budget-monotonicity.
  The local source-control and solver actual-unit siblings likewise remove
  only the abstract \(\lean{gammaValid}\) fields from the row-max active-pivot
  local route, deriving them from \((m:\mathbb R)u<1\).
  Its actual-unit-roundoff sibling removes the abstract sampled
  \(\lean{gammaValid}\) fields by deriving them from \((s:\mathbb R)u<1\);
  this is only a validity-surface reduction.  The actual-unit horizon sibling
  combines that scalar guard with the row-max horizon theorem, removing both
  the sampled validity fields and samplewise global budget-monotonicity from
  this row-max probability surface.  The scalar-comparison explicit-validity
  horizon sibling first extracts the displayed row-max comparison from
  \(\lean{storedQRStageRowMaxComparisonDefectBudget}\le0\), then reuses that
  horizon row-max wrapper, so the two-scalar probability surface no longer
  needs samplewise global budget-monotonicity.  Its actual-unit horizon sibling
  derives the sampled validity fields from \((s:\mathbb R)u<1\) and calls that
  explicit-validity wrapper, while still keeping both scalar defects visible.
  The exact no-pivot shortcut for this scalar condition is now ruled out:
  the existing two-stage exact Householder counterexample has positive
  \(\lean{storedQRStageDiagLowerDefectBudget}\) for the constant stage budget
  \(2\).  The exact no-pivot comparison-scalar shortcut is separately ruled
  out as well: the same exact witness has positive
  \(\lean{storedQRStageRowMaxComparisonDefectBudget}\) for the nonnegative
  constant stage budget \(3\), because the displayed strict-upper row maximum
  is \(2\).  Thus exact recurrence plus signed squared-norm identities,
  nonzero denominators, and nonnegative stage budgets still do not imply the
  comparison scalar.  The scalar conditions themselves remain open for the
  concrete pivoted/sorted/off-diagonal-controlled loop.  The stored-recurrence
  active-pivot shortcut is also ruled out: the same exact-with-\(u=0\)
  stored-panel witness satisfies the stored recurrence, signed-alpha equation,
  nonzero source denominators, and \((2:\mathbb R)u<1\), but its displayed
  first pivot is not the finite active max-pivot selector because the second
  column has larger active trailing norm.  Thus the active-pivot equation must
  be proved by a concrete pivoted/sorted loop invariant or kept visible.  The
  stored-recurrence scalar row-budget shortcut is also ruled out: the same
  exact-with-\(u=0\) stored-panel witness satisfies the stored recurrence,
  signed-alpha equation, nonzero source denominators, and \((2:\mathbb R)u<1\),
  while constant stage budgets \(2\) and \(3\) make the scalar
  stage-diagonal and stage-budget/row-max comparison defects positive.  Thus
  those scalar fields must also come from a concrete pivoted/sorted
  off-diagonal-control invariant or remain visible.  The diagonal-dominant
  stored-surface shortcut is ruled out as well: the exact stored sequence from
  \(\begin{bmatrix}3&2\\0&1\end{bmatrix}\) remains locally diagonally dominant
  and satisfies the stored, signed-alpha, source-denominator, actual-unit, and
  nonnegative-stage-budget fields, but stage budgets \(4\) and \(3\) still make
  the scalar diagonal-lower and comparison defects positive.
  Row sorting itself now
  preserves the rectangular least-squares residual objective and normal
  equations, by the row-permutation lemmas
  \(\lean{lsObjective\_permuteRows}\),
  \(\lean{rectLSGram\_permuteRows}\), and
  \(\lean{rectLSRhs\_permuteRows}\).  Column pivoting now also has the matching
  coefficient-relabeling foundation:
  \(\lean{rectMatMulVec\_permuteCols}\),
  \(\lean{RectLSNormalEquations.of\_permuteCols}\), and
  \(\lean{IsLeastSquaresMinimizer.of\_permuteCols}\).  The combined wrappers
  \(\lean{rectLSGram\_permuteRowsCols}\),
  \(\lean{rectLSRhs\_permuteRowsCols}\),
  \(\lean{RectLSNormalEquations.of\_permuteRowsCols}\),
  \(\lean{lsObjective\_permuteRowsCols}\), and
  \(\lean{IsLeastSquaresMinimizer.of\_permuteRowsCols}\) now package the
  row-sorted, column-pivoted least-squares pullback in one step.  On the
  unpivoted
  source-controlled branch, the stored-QR final triangular block and
  back-substitution solution are now connected directly to the equation~(8)
  high-probability rounded sampled-row objective theorem.  The remaining
  source dependency is therefore to supply the determinant/lower-zero,
  global compact-step budget, and row-budget/diagonal-lower-bound fields for
  the chosen concrete rounded loop,
  or to replace that branch by a pivoted/sorted theorem family with its own
  proved growth invariant;
\item distribution-specific random-projection/FJLT/Gaussian/Rademacher
  uniformization guarantees for Algorithm 3.  The weaker coordinate-Hoeffding
  leverage-probability event and the induced one-step uniform-row Loewner
  bound are now formalized.  The deterministic-after-preconditioning iid
  uniform sample-average concentration theorem is also formalized in
  tail-budget finite-Loewner form, and the signed-Hadamard preprocessing
  randomness is now composed with the iid uniform sampling randomness on the
  product law for the scoped coordinate-Hoeffding route.  The corresponding
  floating-point uniform-sketch transfer is also formalized with an explicit
  sample-dependent FP perturbation budget.  A deterministic-radius refinement
  is also formalized: any scalar \(\tau\) dominating that budget over all joint
  outcomes gives the same high-probability event with radius
  \(\varepsilon+\tau\), and a sampled-row norm cap gives a closed-form fixed
  budget.  The source-directed expected Euclidean row-norm step, deterministic
  row-norm convexity and Lipschitz inputs from Tropp's Lemma~3.3 proof, the
  affine scaling that turns Ledoux's \([0,1]^m\) concentration constant into
  Tropp's Rademacher constant, the finite MGF/Herbst/Laplace algebraic
  substrate, the corrected-log-MGF quotient monotonicity step, the right-limit
  at zero and final finite Herbst extraction to a log-Laplace bound, the
  entropy-bound-to-tail wrapper, the finite product-law entropy chain rule,
  the unbiased Bernoulli coordinate expectation and entropy formulas, the
  fair-Bernoulli coordinate log-Sobolev inequality, the first one-coordinate
  product peel-off for \(P\otimes\nu\), the finite \(L^2\) section-norm bridge,
  the abstract Bernoulli-product induction lift, the concrete
  \(\texttt{RademacherTrace}\ m\) cube entropy-gradient theorem, the
  conditional exponential-tilt reduction from that theorem, the conditional
  finite-cube Chernoff wrapper, scalar half-tilt exponential inequalities,
  uniform Rademacher flip invariance, non-sharp pointwise half-difference and
  coordinate-absolute-difference bridges into the tilt flip-gradient
  hypothesis, the concrete signed-Hadamard row-norm positive-flip
  self-bounding estimate
  \(\sum_k\max(F_i(\omega)-F_i(\omega^k),0)^2\le 4/m\), the positive-drop
  exponential-tilt bridge, the source-sharp one-row and all-row SRHT row-norm
  tails, the resulting leverage-probability cap, the logarithmic
  \(t=\sqrt{8\log(m/\delta)/m}\) preprocessing-budget corollary, and the exact
  source-sharp SRHT preprocessing plus iid uniform-sampling composition theorem,
  and the matching source-sharp rounded sample-Gram floating-point transfer are
  formalized in explicit \(t\)-parameter, logarithmic-preprocessing,
  tail-budget, and fixed FP-budget forms.  The remaining Algorithm~3
  distribution-specific claims concern non-SRHT Gaussian/FJLT/input-sparsity
  preprocessing;
\item the low-rank structural condition of equation~(9).  The exact
  rank-factorization vocabulary, the left/right basis-product rank-at-most
  certificates, the column-sketch projector rank bridge, and the explicit
  residual-certificate and coefficient-handoff surfaces are now formalized.
  Rectangular SVD, pseudoinverse, full-rank \(V_k^T Z\), best rank-\(k\),
  unitarily invariant norm, and computed SVD/projector routine certificates
  are still open;
\item the matrix-completion claims around equations~(10)--(11);
\item the effective-resistance sparsification and Laplacian-solver claims
  described later in the paper.
\end{itemize}
The deterministic transfer and sketch-objective theorems above are real Lean
theorems, but they should not be read as proofs of those still-open randomized
or structural results.  Consequently, the current full-paper formalization gate
is intentionally marked as failed until the open paper-level claims in the
ledger have matching Lean theorems.

\section{File Map}

\begin{description}
\item[Algorithm 1 update and trace stability.]
  \leanfit{RandNLA/ElementwiseSampling.lean}
\item[Algorithm 1 concentration.]
  \leanfit{RandNLA/HitCountConcentration.lean}
\item[Algorithm 1 one-step trace-MGF product-law adapters.]
  \leanfit{RandNLA/ElementwiseTraceMGF.lean}
\item[Algorithm 1 equation (2) spectral transfer.]
  \leanfit{RandNLA/ElementwiseSpectral.lean}
\item[Finite symmetric matrix spectral bridge.]
  \leanfit{Analysis/MatrixSpectral.lean}
\item[Finite-real to complex C\(^*\)-matrix bridge.]
  \leanfit{Analysis/CStarMatrixBridge.lean}
\item[Operator-log/exponential bridge.]
  \leanfit{Analysis/OperatorLog.lean}
\item[Complex C\(^*\)-matrix expectation/order/Jensen bridge.]
  \leanfit{Analysis/CStarMatrixExpectation.lean}
\item[Lieb trace and relative-entropy route vocabulary.]
  \leanfit{Analysis/LiebTrace.lean}
\item[Trace-exponential probability bridge.]
  \leanfit{Analysis/MatrixConcentration.lean}
\item[Algorithm 2 sampled rows.]
  \leanfit{RandNLA/RowSampling.lean}
\item[Algorithm 2 Gram analysis and FP equation (5).]
  \leanfit{RandNLA/RowSamplingGram.lean}
\item[Algorithm 2 leverage scores and equation (7).]
  \leanfit{RandNLA/RowSamplingLeverage.lean}
\item[Algorithm 2 row-trace and leverage log-CGF prerequisites.]
  \leanfit{RandNLA/RowSamplingTraceMGF.lean}
  \leanfit{RandNLA/RowSamplingLeverageMGF.lean}
\item[Algorithm 3 preconditioning.]
  \leanfit{RandNLA/UniformRowSampling.lean}
  \leanfit{RandNLA/UniformRowSamplingMGF.lean}
  \leanfit{RandNLA/UniformRowSamplingComposition.lean}
  \leanfit{RandNLA/UniformRowSamplingFP.lean}
  \leanfit{RandNLA/Preconditioning.lean}
\item[Equation (8) sketched least squares.]
  \leanfit{RandNLA/LeastSquaresSketch.lean}
\item[Equation (9) low-rank foundations.]
  \leanfit{RandNLA/LowRankApprox.lean}
\item[QR and Householder stability substrate.]
  \leanfit{Algorithms/InverseBounds.lean}
  \leanfit{Algorithms/QR/HouseholderSpec.lean}
  \leanfit{Algorithms/QR/HouseholderApply.lean}
  \leanfit{Algorithms/QR/HouseholderQR.lean}
  \leanfit{Algorithms/LeastSquares/LSQRSolve.lean}
\item[Full theorem ledger.]
  \leanfit{docs/RANDNLA_CACM_THEOREM_LEDGER.md}
\item[Open full-paper backlog.]
  \leanfit{docs/RANDNLA_CACM_NOT_PROVED_LEDGER.md}
\end{description}

\end{document}
